\documentclass[10pt,a4paper,twoside,bg=print,twocolumn,openany,nodeprecatedcode]{dndbook}
\usepackage[utf8]{inputenc}
\usepackage{tabulary}
\usepackage[normalem]{ulem}
\usepackage{color,soul}
\usepackage{float}
\usepackage{caption}
\usepackage{wrapfig}
\usepackage{tocloft}
\usepackage[hidelinks]{hyperref}
\raggedbottom
\extrafloats{2000}
\cftsetindents{section}{0em}{0em}
\cftsetindents{subsection}{0em}{0em}
\cftsetindents{subsubsection}{0.2em}{0em}
\setcounter{tocdepth}{1}
\definecolor {mytable} {HTML} {f4edd8}
\definecolor {mycomment} {HTML} {ffffff}
\definecolor {mysidebar} {HTML} {d0cfc2}
\definecolor {myreadaloud} {HTML} {d9e8f1}
\definecolor {myreadaloudborder} {HTML} {5a6173}
\definecolor {mypagenum} {HTML} {2a2d2e}
\DndSetComplexThemeColor[mytable][mycomment][mysidebar][myreadaloud][myreadaloudborder][titlered][titlegold][titlegold]\setlist[description]{nolistsep,listparindent=3pt,topsep=6pt}
\begin{document}
\tableofcontents
\newpage
\chapter{Tal'Dorei Campaign Setting Reborn}
\section{Preface}
\DndDropCapLine{I}{}nspiration can come from anywhere: a musician's song, the advice of a close friend, or just the trust that another being places in you. In my experience, few things can compare to the surge of joy and motivation that being inspired gives.


I have often found inspiration in fantasy worlds and books of adventure, mystery, and magic. I grew up devouring every novel I was handed, seeking to lose myself in the words of another mind, channeling the images they evoked in my mind into my sketchbook. But it wasn't until I discovered tabletop roleplaying games that I felt truly empowered to create my own worlds, places like the ones that I so often found myself lost in. Even better, it let me share all the ideas in my head with my friends, and they could do the same with me. I was spellbound.

It's been about twenty-five years since my first game. Since then, I've drawn endless inspiration from the imaginations of my friends and loved ones as we tell heroic stories around a table with dice and paper. I've found a circle of friends who share my love of dramatic storytelling, my belief that failure in games is just as exciting as success (maybe even more so), and my delight in being inspired by moments of improvisational brilliance. It's still surreal that we found a way to share that with the world; \textit{Critical Role} has become a phenomenon that's surpassed my wildest dreams. Over the past few years, we've built the magical world of Exandria together, warts and all, and somehow it's become a world that millions of others have engaged with and invested their imaginations in. I hope it's inspired you to create your own adventures, stories, and joyful memories. Through Vox Machina's epic tale, the land of Tal'Dorei became a landscape rife with intrigue, conflict, and endless possibility... and we wanted to share it with you.

Years ago, I had an opportunity to bring Tal'Dorei to the page and invite the world to tell their own stories in Exandria. The thrill of the opportunity was incredible, but the immense pressure and fear of failure (given my absolute lack of experience doing something like this) made it a strenuous journey on an incredibly short timeline. I was incredibly blessed to have become friends with the immeasurably brilliant James Haeck around that time, and I asked him to join me in this task. My good fortune found him agreeing, and with his help and inspiration, it somehow came together. With that, a limited run of the \textit{Tal'Dorei Campaign Setting} went into the world.

Time has passed, both in our world and within the world of Exandria. Vox Machina's greatest adventures came to a close. Looking two decades later (within Exandria), we embarked on a new adventure with the Mighty Nein in the land of Wildemount...but Vox Machina and Tal'Dorei remained deeply dear to my heart. The question arose: what would Tal'Dorei look like these two decades later? What changes would have taken place in the recovery following the Chroma Conclave and the ascension of the Whispered One? How would the legendary heroes of Vox Machina have adjusted to life after the end of their story? What dangers and threats, both old and new, would the next generation of heroes need to rise up against?

Inspiration struck once more, and now that \textit{Critical Role} had the means to do so...we needed to answer these questions for both ourselves and the numerous Critters who often conjectured as we did. "\textit{Who is on the Tal'Dorei Council?}" Well, we can finally answer that! I immediately called James Haeck back to the fray, and we were joined this time by the extraordinary mind and talents of Hannah Rose to elevate these ideas into something better than I could have hoped. Between us, Tal'Dorei was Reborn. We invite you to get lost in this world, learn of its intricacies and secrets, and, I hope, be inspired to make this world yours as well.

Matthew Mercer,

Game Master of \textit{Critical Role}

\chapter{Welcome to Tal'Dorei}
\DndDropCapLine{T}{}al'Dorei is a land of epic adventure, filled with heroes who do great deeds for love, for justice, and—maybe more often than not—for gold and their own egos. Villains are draped in shadow and draw upon the power of cruel gods—but they are just as motivated by love, their own twisted concepts of justice, and their hunger for gold, power, and glory.


Tal'Dorei is a land created by Matthew Mercer to introduce his friends to fantasy roleplaying games, and he filled it with the greatest high fantasy tropes and played them to the hilt—and then twisted them, defying his own clichés when it would hit the hardest.

This book is for Game Masters who wish to run fantasy roleplaying game campaigns set on the continent of \textit{Critical Role's} first campaign, and for the players who will create characters for those games. The game mechanics in this book use "fifth edition rules," but it also includes enough lore about the world to form the foundation for a game using the rules of any fantasy RPG.

\section{Lands of Tal'Dorei}
The name Tal'Dorei refers to three things throughout this book. Most often, Tal'Dorei means the continent that is home to all of the people, locations, monsters, magic items, and more within these pages. It is one of four major continents on the world of Exandria. Tal'Dorei can also mean the Republic of Tal'Dorei, the largest nation on the continent. And finally, Tal'Dorei was first the surname of a mighty hero who lent her name to the land that she helped liberate from evil.

See "\textit{A History of Tal'Dorei}" later in this chapter for more details about the heroic \textit{Zan Tal'Dorei}, the land she founded, and her descendants.

The continent of Tal'Dorei encompasses several nations and independent city-states, all of which are home to diverse peoples who view the world in ways influenced by their cultures and the lands they call home. The most significant nations of Tal'Dorei include:

\subparagraph{The Republic of Tal'Dorei}
From the great city of \textit{Emon} in the west, across the vast \textit{Dividing Plains}, and even as far as the isolated city-state of \textit{Whitestone} in the northeast, the Republic of Tal'Dorei spans the breadth of the continent that bears its name. It is ruled by a council of elected representatives and is cordially allied with all other nations on the continent.

\subparagraph{The Ashari}
The elementally attuned Ashari people are reclusive and dedicated to their mission of safeguarding elemental rifts across the world. Two of the greatest rifts exist in Tal'Dorei, and the Earth Ashari of \textit{Terrah} and the Air Ashari of \textit{Zephrah} keep watch over them from their secluded homes.

\subparagraph{The Deep Halls of Kraghammer}
Deep beneath the \textit{Cliffkeep Mountains} is \textit{Kraghammer}, the ancestral home of Tal'Dorei's dwarven people. Though \textit{Kraghammer} is home to people other than \textit{dwarves}, its dark and claustrophobic halls act as a deterrent to non-dwarven residents—as does the dwarf clans' famous inhospitality.

\subparagraph{The Verdant Enclave of Syngorn}
In the heart of the trackless \textit{Verdant Expanse} is a city of \textit{elves} who hid in the Fey Realm from a war that nearly annihilated all life on Exandria. \textit{Syngorn} is welcoming to all who brave the \textit{Verdant Expanse} to reach them, though there is an expectation that those who visit take pains not to flaunt the strict and ancient rules of \textit{Syngornian} society, lest their stay come to a swift and unceremonious end.

\subparagraph{The Iron Authority}
This imperial force rules the \textit{Beynsfal Plateau} in the southernmost reaches of the continent. There they train legions of soldiers in an endless campaign to conquer the entire \textit{Rifenmist Peninsula}. The other powers of Tal'Dorei are just beginning to take notice of the movements of the \textit{Iron Authority}, but most of their knowledge is based in rumor. The \textit{Tal'Dorei Council} has little interest in drawing the authority's ire while it is preoccupied with its conquest of \textit{Rifenmist}.

\section{What's in This Book?}
\textit{Chapter 1} is an introduction to Tal'Dorei's lands and history. It also describes certain fundamental details about the world of Exandria as a whole, including its calendar, the planes, and the cosmos. Finally, this chapter provides an overview of Tal'Dorei's society today—and the brewing tensions and secrets that Game Masters can use to plot their own campaigns.

\textit{Chapter 2} describes the allegiances player characters and non-player characters (NPCs) may have to deepen their connection to the world. Their allegiances may be to factions, such as governments or guilds, or to the gods. These divine beings can only affect the Material Plane through their mortal followers, or occasionally supernatural messengers like celestials or fiends.

\textit{Chapter 3} is the Tal'Dorei Gazetteer, a detailed look at all of Tal'Dorei's regions and cities, their points of interest, the people who reside within them, and the quests that heroic adventurers might undertake there.

\textit{Chapter 4} contains new options that players can use to build their characters, including story descriptions of the different ancestries; new subclasses, backgrounds, and feats; and a handful of optional rules for the whole table.

\textit{Chapter 5} contains advice to help Game Masters create exciting and compelling campaigns in Tal'Dorei, and magic items that you can give your players as treasure during their perilous adventures—including the unstoppably powerful \textit{Vestiges of Divergence}.

\textit{Chapter 6} presents game statistics for monsters and NPCs that you can use to challenge your characters or aid them in a pinch. Among them are the legendary heroes Vox Machina, some twenty years after the conclusion of their greatest adventures.

\section{Calendar, Time, and the Cosmos}
The world of Exandria is a planet with regular seasons and a wide variety of climes, ranging from frigid ice caps to arid deserts, temperate forests, windswept plains, sweltering jungles, and storm-tossed oceans. It orbits a single sun, is orbited in turn by two moons, and is surrounded by distant stars that twinkle in the evening skies.

The ancient people of Exandria told time by these celestial bodies, and the \textit{elves} created a calendar based on their movements.

\subsection{Calendar}
The calendar year of Exandria runs a total of 328 days, grouped into seven-day weeks over the course of eleven months. These months are outlined below in the order of their arrival within the calendar year, along with their respective number of days and holidays common throughout Tal'Dorei (many of which are detailed within the "\textit{Pantheon of Exandria}" section in \textit{chapter 2}).

\subsubsection{Counting Years}
In Exandria, the years are most commonly counted by a reckoning known as "\textit{Post Divergence}," with year 0 being the fateful year that the \textit{Prime Deities} ended the \textit{Calamity} by sealing all the gods behind the Divine Gate. This book describes the world as it is in the year 836 P.D. If you're familiar with the \textit{Critical Role} campaigns, this year is twenty-four years after the end of Campaign 1 (Vox Machina) and toward the end of Campaign 2 (The Mighty Nein).

\subsubsection{Months and Seasons}
The names of the months and days of the week vary between cultures, but the elven calendar described here is the standard means of reckoning dates in Tal'Dorei—as well as in most nations around the world, for ease of international trade and diplomacy.

Just as in our world, there are four seasons, and the year ends and begins just after the winter solstice.

The first blush of spring in Tal'Dorei is felt as new flowers bloom on Wild's Grandeur, a holy day on the 20th of Dualahei. The official start of the season, however, is observed a week earlier and celebrated with games, music, and just-ripened food at the Renewal Festival. Summer brings hotter days toward the middle of Unndilar. The noonday sun of the 26th day is called the Zenith, and this solstice is considered the first true moment of the summer.

The colors change and winds cool as autumn begins in the early days of the month of Fessuran, marked by the Harvest's Close on the 3rd day. The chill of winter arrives to bring longer nights and cleansing snow on Barren Eve, the 2nd of Duscar, a nighttime celebration and remembrance of those who fell in battle.

The seven days of the week are named Miresen, Grissen, Whelsen, Conthsen, Folsen, Yulisen, and Da'leysen. Each day is 24 hours long, and most city-dwellers in the Republic of Tal'Dorei are expected to work eight hours of that day—except for Yulisen and Da'leysen, which comprise the weekend, and on holidays.

These social rules don't necessarily apply to cultures where sleeping patterns are different (such as in \textit{Syngorn}, where most of its people are \textit{elves} who trance a mere four hours of the day), or to people outside of cities—such as farmers, who rise with the sun, and traveling merchants and adventurers, who work highly unusual hours.

\subsubsection{Using a Calendar}
The Exandrian calendar is intentionally quite different from our own. The creation of calendars is a messy and often illogical process, with strange artifacts of their making that linger on only because of tradition—and because it would be too much of a hassle to get the entire world to change their way of counting time and days.

You don't need to bother with the Exandrian calendar in your home game. It's perfectly reasonable for you to use a nice, simple seven-day week, thirty-day month, and twelve-month year. You may even use our names for the months and the days of the week so that when an NPC tells the party to meet them at noon on "Grissen, Misuthar 7th," they're basically talking about "Tuesday, February 7th."

However, if you want to immerse your players in the fantasy world of Exandria, using this unfamiliar calendar may help you on your quest.

\begin{DndTable}[header=Exandrian Calendar]{XcX}

Month & Day & Holidays\\
Horisal & 29 & New Dawn (1st)Hillsgold (27th)\\
Misuthar & 30 & Day of Challenging (7th)\\
Dualahei & 30 & Renewal Festival (13th)Wild's Grandeur (20th)\\
Thunsheer & 31 & Harvest's Rise (11th)Merryfrond's Day (31st)\\
Unndilar & 28 & Deep Solace (8th)Zenith (26th)\\
Brussendar & 31 & Artisan's Faire (15th)Elvendawn or Midsummer (20th)\\
Sydenstar & 32 & Highsummer (7th)Morn of Largesse (14th)\\
Fessuran & 29 & Harvest's Close (3rd)\\
Quen'pillar & 27 & The Hazel Festival (10th)Civilization's Dawn (22nd)\\
Cuersaar & 29 & Night of Ascension (13th)Zan's Cup (21st)\\
Duscar & 32 & Barren Eve (2nd)Embertide (5th)Winter's Crest (20th)\\
\end{DndTable}

\subsection{Catha and Ruidus}
When the people of Tal'Dorei gaze into the morning sky, they see the sun, a blazing disc the size of a gold piece, looking back at them. It is the domain of the The Dawnfather, a god of life and light honored politely by nearly all the realm's denizens.

At night, the skies are filled with countless distant stars, a gleaming silver-white moon called Catha, and occasionally, a dim, ruddy moon called Ruidus that is nearly half Catha's size. Both moons are the domain of the The Moonweaver, a capricious god of trickery and illusions, but there are sects throughout the world that believe only Catha represents the The Moonweaver cunning and grace. Catha's pearly glow is said to bless the just with cunning and caution, and to make hidden the goodhearted when they require stealth and subtlety.

Ruidus, on the other hand, is so surrounded by disquieting rumors and folkloric tales of misfortune that some believe another unknown god or power rules this small, reddish-brown moon. Cultures around the world tell countless legends of prideful rulers who made grand plans or attempted deeds under the moon's full light—when it shines a brilliant vermilion rather than its usual ruddy color—and were forced to watch in horror as their endeavors fell to unforeseen misfortune. It is said those who fall afoul of Ruidus failed to give it the deference it is due—and so superstitious folk rarely dare to make plans while the full light of Ruidus shines above, let alone enact them.

Worse yet, some tales forebode dark fortunes for those born under the light of a full Ruidus, a curse of ill luck that will follow them throughout their lives. Though the cruel practice of moon-sacrifice is no longer permitted in any of Tal'Dorei's cities, some far-flung settlements still secretly sacrifice children born under a full Ruidus to "save" them from a cursed life, and to appease the dark and unknowable appetites of the red moon.

Fortunately for the superstitious, Ruidus is rarely full. While its cousin, Catha, completes a full cycle approximately once a month, Ruidus's haunting, half-year orbit, combined with its eerie and unexplainable tendency to simply not appear in the sky on certain nights or glow with unknown light, creating an unexpected full moon, has only added to its mythic reputation as an omen of ill fortune.

\begin{DndReadAloud}{Prologue to the Song of Alyxian}
\DndQuote%
{}
{''O Ruidus, grant humble chorus leave''}
{''To sing the song which hails the zenith of''}
{''Your accurséd, thrice-blessed Apotheon;''}
{''Remember'd best by deeds in war, and yet''}
{''Whose acts were driven oft by fate most foul.''}
{'' ''}
{''His kindly brow bore gifts from gods of change,''}
{''And art, and moon, yet in his soul was pain,''}
{''The suffering of your vermilion light''}
{''Drove the Paragon to a desert realm''}
{''Bestrewn with blades and drenched in crimson blood.''}
{'' ''}
{''So hear, O moon of curséd deeds and fates!''}
{''The song of he who rose above your great''}
{''And mighty pow'r, to save Exandria.''}
{''From flames of war fanned by the Ruiner's blade.''}
{}
\end{DndReadAloud}

\subsection{Planes}
Though this book concerns itself mostly with the world of Exandria, and specifically the continent of Tal'Dorei, high-level adventurers may eventually gain magic that allows them to travel to other planes of existence beyond the material world in which Exandria resides.

The planes surrounding Exandria are similar to the planes presented in fifth edition lore. Below are brief descriptions of some of the planes that a party of powerful adventurers may travel to. These planes go by many poetic and fanciful names.

\subsubsection{Material Plane}
The Material Plane is a land where natural laws work as we expect. It is a place of all four elements, and of plants, animals, and minerals. Exandria is but one of many worlds on the Material Plane—though none but the most accomplished astronomers have become aware of this fact, much less devised a way of traveling between material worlds.

\subsubsection{Elemental Planes}
When the Protean Gods discovered Exandria, it was like a barren ship adrift in a sea of churning, chaotic elements. As a part of their creation, the elements were thrust away from the world. Some theorize that the Elemental Planes were created from this banished chaos, while others believe the elemental chaos seeped into the Material Plane from tears in the fabric of the planes.

The latter may be more likely, as rifts between the Elemental Planes can be found across Exandria. The largest of these rifts are tended to by the four Ashari civilizations, and any who gaze within find wild wastelands of pure elemental power: roiling seas of fire, towering mountains of gnashing stone, infinite skies filled with cloud-top settlements, and endless, stormy oceans of water.

\subsubsection{Faerie, Shadow, and Ether}
Parallel to Exandria are odd reflections of its reality. The Fey Realm is a plane of natural beauty and ferocity, inhabited by fey folk. The Plane of Shadow is a dour, sinister mirror of the world that is. And the Ethereal Plane is a hazy film across the material world, as if it were viewed through an old, cloudy mirror, inhabited not only by restless spirits, but also by predators that target those who cling to the world of the living.

\subsubsection{Astral Plane}
All the prior planes exist directly adjacent to the Material Plane, buttressing and occasionally co-mingling with its inhabitants. All the planes that follow are separated from the material world by a vast sea of starlight, barren asteroids, and misty nebulae of unpredictable magic. In the ancient past, arcanists had to craft vessels to sail the treacherous sea of stars to travel between their home plane and the mysterious Outer Planes.

These days, magic such as the plane shift spell has allowed planar travelers to bypass these methods, but there are still a number of strange creatures and forgotten, imprisoned demigods who float eternally through the Cosmic Sea.

\subsubsection{Outer Planes}
The \textit{Prime Deities} imprisoned the \textit{Betrayer Gods} in prison planes after the \textit{Founding}, and retired to planes of their own creation after the \textit{Divergence}. Over time, the divine beings on these planes reshaped their new homes in their own image, creating places where natural laws don't function as one native to the Material Plane might expect.

All of these planes are strongly aligned to one or more cosmic principles of good, evil, law, and chaos, based on the divine powers that reside there. Because of the Divine Gate established after the \textit{Divergence}, no divine powers can depart these planes, for the sake of all Creation. The gods' supernatural creations can pass through the Divine Gate, though the more powerful they grow—and thus, the closer they come to divinity—the harder it is for them to pass through. However, mortals and lesser supernatural beings who possess magic to travel between the planes can travel to and from Outer Planes like the Hells, the Abyss, Elysium, and Pandemonium unhindered by the Divine Gate.

\subsubsection{Realms Beyond}
Far beyond even the divine Outer Planes is a realm where all laws of nature break down completely. Mortal senses have no way of properly perceiving the nature of these far-distant realms, and most minds that try to comprehend them snap under the strain of so much unfamiliar stimulus. These maddening realities and the alien horrors that live within them are mercifully kept away from the known multiverse by the unimaginable distance of space.

Some creatures of the beyond have lurked on the Material Plane since ancient times, perhaps at the invitation of the The Chained Oblivion itself, or because the divine powers unleashed in the \textit{Calamity} punched holes in time and space that allowed such aberrations to slither into Exandria unimpeded.

Major known locations of aberrant activity in Tal'Dorei include \textit{Yug'Voril} beneath \textit{Kraghammer}, the \textit{Crystalfen Caverns} beneath \textit{Emon}, and the realm of \textit{Ruhn-Shak} beneath the \textit{Stormcrest Mountains}.

\subsubsection{Demiplanes}
Massive, infinite spaces such as the Material, Elemental, and Outer Planes are not the only types of reality. Through the use of powerful magic, arcanists can create stable demiplanes—finite pockets of extraplanar reality that they can shape to their whims. These "pocket dimensions" exist as bubbles of floating color within the Cosmic Sea, and can only be traced or penetrated by extremely powerful, lost magical arts.

\section{A History of Tal'Dorei}
\begin{DndReadAloud}{}
\DndQuote%
{}
{''...then, the vast empty heart of endless shadow befell the light of the Seeker. Cascading from the pits was the font of life. The first moon formed the endless oceans. The second sun brought the lush soil. The third wind carried breath of life. The fourth flame ignited the Heart of Exandria.''}
{}
\end{DndReadAloud}

Where did Tal'Dorei come from? The origin of this land is a time-shrouded question that everyone has pondered in some way, from the most exalted, world-traveling hero to the common farmer who has never traveled three miles beyond their village. Everyone has their own ideas on where Tal'Dorei, and all of the wide world of Exandria, came from. These varied thoughts all have some things in common, since they are all born of jumbled misunderstandings of ancient myths told orally from bard to bard, then passed from father to daughter, and eventually codified into religious re-writings of history that favor the teachings of a given deity.

To speak a full and accurate truth of Tal'Dorei is an impossible task, yet Tal'Dorei is known for creating folk who regularly do the impossible. What follows is the Myth of Exandria, an account based on the investigations and best efforts of Tal'Dorei's foremost scholars—the homegrown \textit{Alabaster Lyceum} of \textit{Emon}, and the international knowledge-seekers and anti-propagandists of the \textit{Cobalt Soul}. These two organizations have assembled this history by digging beneath the layers of dust and decay that hide battles long past, unearthing texts long thought burned and censored, and exhuming historical treasures from tombs of heroes long forgotten. The details are often debated, for the question of the land's origin remains—consuming the curious, calling those hungry for purpose, and fueling the business of adventuring to delve into the tantalizing unknown places of the world.

\subsection{Myth of Exandria}
Before the history of the lands of Tal'Dorei can be explored, its context in the history of Exandria itself must be thoroughly established. The world of Exandria is home to four major continents, dozens of nations, and untold tiny pockets of civilization amid a wondrous, dangerous wilderness. Every civilization, from the largest empire to the tiniest nomadic clan, has its own interpretation of where its story began. Even within the world of Exandria, different cultures have creation myths that eventually converge with history, but there is no known definitive story.

Even so, life ever seeks to understand its inception. In this report, the collected minds of the \textit{Cobalt Soul} and \textit{Emon's} \textit{Alabaster Lyceum} can say, with only the most trivial of doubts, that the city of Vasselheim in \textit{Othanzia} truly deserves its title of the "Dawn City." It is here, to our researchers' best understanding, where the stories of all mortal peoples begin. Vasselheim became Exandria's oldest surviving city by virtue of not just being built long ago, but by enduring the \textit{Calamity} (which shall be elucidated presently), a miraculous feat that few other charted settlements can claim to have accomplished.

In short, Vasselheim houses the earliest known temples to the gods, and the earliest known records of history that survived this catastrophe. It is these sources that this joint research initiative treats as its primary documents, though great effort has been made to corroborate their claims through later sources. It is in the chronicles of Vasselheim that the earliest account of the myth of the \textit{Founding} is recorded—and indeed, this history is the most commonly accepted origin of the world in the Republic of Tal'Dorei. Though countless regional variations muddy the historical record, the broad strokes of this mythic history are undeniable.

\subsubsection{The Founding}
Long ago, this world was one of tumultuous and chaotic forces, naught but unbridled fire, seething oceans, and churning, gnashing rock. Through the ashen skies of Primordial Creation, the gods arrived from an unknown realm located beyond the ether. These ambitious divinities—young and still formless—looked down upon this roiling realm and saw potential for great strength and outstanding beauty, and the chance to learn their own place in creation.

From the divine hands of the The Arch Heart sprung forth the First Children, the \textit{elves}. It is thought that they were made in their creator's image, yet it is just as likely that the The Arch Heart, and all the gods who shared their ideas with them, actually adopted forms inspired by the image of their creations. The First Children walked the lush Green and gazed upon the heavenly Blue as the gods tamed Primordial Creation around them.

Not long after, in the scale of the cosmos, a second mortal creation was wrought by the divine The All-Hammer: the \textit{dwarves}, a hearty people imbued with the desire to continue the gods' work—taming a primordial world filled with the craft and invention of the divinity that lay beyond the ashen void. In their turn, a third people were given life by the passionate lovers of Order and Chaos, the The Lawbearer and the The Wildmother. Their children were the \textit{humans}, endowed with hearts of passion to bring meaning to the world as it assumed an orderly form. Yet in their hearts burned a spark of chaos, cursing them with short life just as it blessed them with the unquenchable ambitions, mirth, and curiosities of the gods that crafted them.

As a note from the High Curator of \textit{Emon's} Cobalt Reserve, the chronicle of Vasselheim states that the gods were given form and name only when the peoples of Exandria began to worship them—yet many of the various races of the world ascribe their creation to a single god or group of gods. It would be remiss to not at least conjecture that while the Protean Gods lacked distinct form, they still had unique thoughts, emotions, and motivations that were later codified through the worship of the people.

As inspiration flowed from the Protean Gods, other creations followed, giving life to Exandria's many peoples. Here the mythic record of Vasselheim grows unclear, for while the \textit{elves}, \textit{dwarves}, and \textit{humans} were undeniably Exandria's first races, the identities of the many Children of Creation that followed were left unrecorded by the Dawn City's chroniclers.

These myriad young races explored their freshly founded world. As their knowledge grew, they attempted to build. But the land was fierce and treacherous, and the lives of the children were extinguished mercilessly by the still-chaotic world. Sorrow is said to have filled the hearts of the gods as they watched their children struggle against a land that rebuked them. The gods ultimately lent their children tiny fragments of their own power, gifts with which they might shape the world around them.

This small mercy of the gods was the birth of magic, and with it the people of Exandria learned to bend the angry earth to their will, to temper the wrath of its explosive mountains, to tame its unyielding floods, and to make the hard earth soft and bountiful. Wandering peoples settled, language became commonplace, survival gave way to art, and governance replaced anarchy. The Protean Creators, the divinity beyond the ashen skies, saw progress and saw that it was good—but fragile, and in need of guardians.

Thus were born the First Protectors: the Dragons Metallic of Tal'Dorei. These Protectors pledged themselves to defend the weak against tyrants. It is here that the mythic tales connect to the oldest, most fragmented passages of recorded history as civilization's reach stretched across the world. The people praised the gods for their newfound safety, and their prayers gave them forms, names, and purpose.

\subsubsection{Wrath of the Primordials}
Yet the realm refused to be tamed. Quaking cliffs roared in defiance of mortalkind's magical gifts. Seas swelled and seethed. Flames erupted from below with greater intensity than any had seen before. Beneath the elements, unknown to the Creators beyond the ashen skies, lived ancient beings who were born in the primeval chaos of the world the gods had found: the Primordials. These elemental titans that once dwelt deep within the land now rose from their unseen domain to sunder it once more. The gods watched as their children were crushed between gnashing rocks and fed to formless terrors unleashed in the wake of the destruction. Thus were the seeds of Exandria's greatest cataclysm set into motion.

Some gods were so full of grief and anger that they wished to abandon this world for another. They tried to convince their divine kindred to join the Primordials in reclaiming the realm for chaos, so that they could move on to start anew. Others wished to prove they could tame the world for their beloved children. This caused a divide among the gods. The schismatic deities, known today as the \textit{Betrayer Gods}, joined their songs and swords to those of chaos and destruction, taking and twisting their children into the image of their intent.

The cosmos itself seemed to recognize the magnitude of the battle on Exandria, and beings of Order and Chaos spilled forth from the still-young multiverse. Creatures of chaos emerged from a formless void that Exandria's peoples came to call the Abyss.

These demons fed upon the suffering of the wounded and the dying, and grew cruel. Sentinels of order emerged from a distant realm of celestial harmony. The \textit{Prime Deities} beseeched these celestials to aid them. In turn, the \textit{Betrayer Gods} forced order upon the demons and made them devils, perfectly lawful soldiers who lived only to cause harm.

As the clash of gods and Primordials spiraled out of control, the desperate mortal races sought their own way to defend themselves from the rebellious elements. Seeing that their creations still suffered, the The Arch Heart taught them the secrets of making magic themselves. The most intelligent and inquisitive mortal minds followed the path the gods had revealed to them, and created incantations and formulae to create new spells on their own terms, without the aid of divine power. These people called themselves arcanists, for they were masters of the arcane secrets of magic.

With their potent new magics, these resourceful mortals subdued the Primordials long enough for the \textit{Prime Deities} to banish their traitorous kin to secluded prison planes. With the \textit{Betrayer Gods} locked away, the gods were overjoyed to see the resourcefulness of their children. The defeated Primordials were destroyed, and from their ever-enduring essences, the gods channeled their wild and destructive powers into the Elemental Planes that surround the material world of Exandria.

For the first time since the Creation, Exandria finally settled into an idyllic peace. From this peace was born the first enduring mortal civilization, the heart of which was a grand city called Vasselheim, the Cradle of Creation and the Dawn City.

Culture developed anew, the races ventured beyond to explore and discover their own lands, and great music filled the air to give name to this world once and for all: Exandria.

\subsubsection{Age of Arcanum}
\begin{DndReadAloud}{}
\DndQuote%
{}
{''The foremost scholars of our day would make it seem an indisputable fact that the gods created Exandria. That they blessed us with the bounty that all peoples of the world now make use of. So too, they say, was magic their gift to us. But such zealots ignore the fact that there is another likelihood: that the gods did not create mortals, and their world—but the opposite. That the dreams of mortals mingled in the wild magic of the young world, and gave that magic form and order: divinity. Of course, the dogmatic zealots of our day will dismiss such ideas as heresy—but not you, reader! The avid mind feeds not on dogma, but on new ideas!''}
{Anonymous}
\end{DndReadAloud}

Ages passed, and society flourished. Great kingdoms sprung up. Castles were built in a day, accelerated by the arcanists' newfound power. Even though magic could be used to complete the most difficult tasks with hitherto unknown speed, magic-users strove always to innovate. As mages practiced and perfected their powers of creation, they soon unlocked the secrets of life itself, giving birth to wondrous, dangerous new forms of life and power.

Yet, with the benefit of hindsight, it is clear that the very magic which allowed mortalkind to prosper also instilled a deep rot within their civilizations. The arcanists of the world grew arrogant. They came to see their arcane gifts as proof that the gods held no sway over their fate, and that with a sophisticated enough command of magic, they could become as powerful as the gods themselves. Though this hurt and surprised the \textit{Prime Deities} who had imprisoned their kin for the sake of their children, they sought to sympathize with the willfulness of their creations, remaining out of love and hoping that the mortals would learn the error of their ways.

The advent of the arcane seemed to be the key to a bountiful age of plenty, but also proved to threaten it, as prosperity soon gave way to greed. Rather than universal goodwill and comfort for all, the \textit{Age of Arcanum} came to be defined by endless petty squabbles over wealth and power among the elite, while those without magic hunted for their scraps. Tantalizing rumors of a path to immortality slithered through the most decadent circles of magi and nobility. One mortal mage, her name lost and presumed to be struck from history, crafted now-forbidden rites to challenge the God of Death, felling him and taking his place among the pantheon, making of her the first and only living mortal to ascend. The name of this original death god was likewise lost to history. Only the title of his now-godly successor, the The Matron of Ravens, survives.

Her victory over divinity was a catalyst for many of the horrors of the \textit{Age of Arcanum}. An archmage named Vespin Chloras, renowned throughout ancient Vasselheim for his wealth, skill, and cruelty, was inspired by this display. He sought the guidance and power of the banished gods, rending open the gates of their prisons and releasing the Betrayers into the mortal world.

In their imprisonment, these Gods of Hatred and Despair had warped their prisons into reflections of their depravity. The Abyss, once a formless place of chaos, a byproduct of the gods' Creation, was twisted into a place of evil. Countless other planes of cruelty were born from the \textit{Betrayer Gods'} hatred, and those realms' evil depths endlessly churned forth horrors that lived only to transform peace into suffering, and righteousness into arrogance and greed.

Released unto the mortal world once more, the \textit{Betrayer Gods'} urge to ruin was supplanted by the desire to dominate, and they turned their sights first to the archmage who had freed them, making Vespin the first of their many mortal thralls. Records of diabolical texts safeguarded within the libraries of the \textit{Alabaster Lyceum} suggest that Vespin, now long dead, serves the \textit{Betrayer Gods} to this day—now as a devil at the left hand of the The Lord of the Hells. These corrupt gods sought out the shattered remnants of the devils they had crafted and the demons they conscripted, and secretly created with them a fearsome new kingdom on the far end of the world, the capital of which was named Ghor Dranas, the Gathering of Shadows.

In this land, where the twisted power of the \textit{Betrayer Gods'} corrupted planes seeped into the world, the lords of evil welcomed mortals whose heartlessness and lust for power made them susceptible to the gods' grand promises. The most fertile soil of all for these noxious seeds lay in mortal hearts obsessed with the infinite power of the arcane. With a legion of the damned behind them, the \textit{Betrayer Gods} soon made their presence known to the world with a surprise assault on Vasselheim itself.

Though much of the city was reduced to rubble, Vasselheim weathered the initial assault, saved by the intervention of their protectors, the metallic dragons, and even a number of the \textit{Prime Deities} themselves. The creators descended to trade blows with their former brethren. The battle, which pitted gods against mortals and heroes against demons, raged ceaselessly for twenty days and nights, until the dark forces were forced to retreat, their surprise attack thwarted.

Yet this victory was a tainted one. Evil had been repulsed momentarily, but the revelation of such a terrible foe incited an arcane arms race. Trust was shattered indefinitely: if mortals could fall under the sway of the \textit{Betrayer Gods}, who could be true allies? If ruin like this could be unleashed under the watchful eyes of divinity, what value did mortal lives hold? Fearing all powers but their own, the most self-interested and singular human arcanists warped their greatest creations, magical instruments of prosperity and joy, into arms and armor of horrific power. The \textit{dwarves'} fascination with rock and earth turned toward isolation as they burrowed further into the mountains, using their divine gifts to animate legions of autonomous golems to protect their ancestral halls. \textit{Elves} used their understanding of creation's beauty and intricacies to weave spells of unimaginable destructive force, the likes of which Exandria had never before seen.

For the first time since the awakening of the Primordials, the focus of magic was warfare. The gods themselves agreed to join their children on the field of battle, descending from the heavens to take up arms once more for the war now referred to as the \textit{Calamity}.

\subsection{The Calamity}
Few records remain of the terrible war that followed, but its effects are still felt today. Most tales that remain are of great champions who the \textit{Prime Deities} blessed with their power. Some deities created warriors of their own from light and holy fire. Others imbued fragments of their power into the weapons now called \textit{Vestiges of Divergence}. More still bestowed blessings upon mortals and called them champions. One tale even tells of a champion blessed by three of the gods in their times of greatest need. Now all but forgotten, this Apotheon was forced by fate into terrible battles across the world.

The sheer magnitude of the energies unleashed in the ensuing battles of gods and mortals frayed the boundaries holding back the elemental chaos, spilling unbridled destruction into the world and bringing utter annihilation. Even the ley lines that direct the flow of magic across Exandria like veins direct blood through the body show signs of being warped by the raw power that the Calamity unleashed. Some well-preserved ruins of the ancient civilizations remain, such as the Drowned City of Cael Morrow in the lands of \textit{Marquet}, but for the most part, all traces of the old world were erased from the face of Exandria.

So great was the loss of life during the war that even liberal estimates suggest no more than a third of Exandria's population survived it. This unthinkable devastation inspired some civilizations to flee Exandria entirely and seek refuge on other planes, scattering the Children of Creation across the multiverse. By the war's end, the only bastion of civilization left on the face of Exandria was the Dawn City itself: Vasselheim.

\subsection{The Divergence}
In the wake of the \textit{Calamity}, the victorious \textit{Prime Deities} once more banished their traitorous kin to their realms of deception and hate. The world entered a long, dark period of recovery, as history had to be recovered and purpose restored, and the threat of the \textit{Betrayer Gods} still loomed heavily upon the minds of all. Even the \textit{Prime Deities} felt guilt for their role in the conflict, for it was the unrestrained clash of divine power that had unleashed such horror upon the world.

The records of the Scalebearers assert that the The Platinum Dragon and the The Lawbearer descended upon Vasselheim and spoke a decision: the \textit{Prime Deities} would depart from the world and establish a Divine Gate that would forever prevent any god from ever acting directly upon Exandria again. This proclamation shocked the other \textit{Prime Deities} and the beleaguered survivors of the war in equal measure, but while mortals railed against this announcement, the other gods quickly realized that the decision of the two most dutiful and self-sacrificing of their number was unimpeachable and just.

Thus, in hopes of ensuring such ruin would not befall Exandria again, they left their children to rebuild civilization anew within the walls of Vasselheim and beyond. The Creators returned to their own realms, sealing all divine powers behind their newly constructed Divine Gate. Only in this way could they prevent their corrupted brethren from physically returning to the Material Plane.

Much time has passed since, and the world has been reborn once again. The gods still exhibit their influence and guidance from beyond the Divine Gate, bestowing their knowledge and power to their most devout worshipers, but the path of mortals is now their own to make. New cities, kingdoms, and cultures have retaken the world, building over the ashes of the old. New songs fill the air, and the hope of a brighter future drives people day after day, while buried ruins and forgotten relics remind all people of a darker time—and of mistakes that should never be repeated.

\begin{DndSidebar}[float=!b]{Names of the Divergence}
Of these many formative events, the tale of the \textit{Divergence} is the most oft-told among the many faiths and civilizations of the world. When seeking the truth of this event, the researchers of the \textit{Cobalt Soul} found that it was called by many titles around the world, for many different reasons.




\begin{itemize}
\item \textbf{The Second Spark.} This name is most often used by those who study the arcane, for it was the gods giving them a second chance to use magic for the betterment of the world.
\item \textbf{The Penance.} Priests, clerics, and religious scholars across Exandria refer to the departure of the gods by this name, for they see it as a self-imposed atonement for the gods' role in the ravaging of the world.
\item \textbf{The Divergence.} The most common name for this event, the \textit{Divergence}, is used by people of all classes and creeds. This descriptive name simply refers to the gods' final act of divergence from the mortal world.
\end{itemize}



\end{DndSidebar}

\subsection{Origins of Tal'Dorei}
Following the creation, and subsequent razing, of Exandria, a \textit{post-Divergence} world was now left to rise from the ashes and begin a new era. While every region had its own rebirth following the terrible destruction of the \textit{Calamity}, the continent that now bears the name of this setting shall remain the focus. The modern calendar began in the year 0 P.D. (\textit{Post-Divergence}), over eight hundred years before the writing of this text. Though the modern nation-state of Tal'Dorei did not emerge for some time after year 0, the history of Tal'Dorei truly began with the founding of \textit{Gwessar}.

\subsubsection{Gwessar}
At the dawning of modern history, what is now known as the continent of Tal'Dorei housed the germinating seeds of civilization. It was the hardy, dependable \textit{dwarves} who best weathered the war between gods and mortals, beneath the \textit{Cliffkeep Mountains}. Dozens of dwarven redoubts rose and fell in the tumult of the \textit{Calamity} and the uncertain era that followed it. For many years, \textit{dwarves} lived in the tunnels beneath the \textit{Cliffkeep Mountains}, leaderless and chaotic, until the various clans unified to form the subterranean city of \textit{Kraghammer}. Proud of their grand new home, the dwarf clans were content to remain beneath the earth. They saw the war-scarred surface world as a source of naught but misery and death, and their new underground domain as a place of undiscovered prosperity. This city was dug underneath the mountains, and its deep roots, though young, survived where more ancient ancestral halls were annihilated. Clan Jaggenstrike were the architects of \textit{Kraghammer's} hardy, unassailable redoubts, and easily became the first ruling clan of the \textit{Kraghammer} \textit{dwarves}.

While the \textit{dwarves} busied themselves excavating the world below, a group of \textit{elves} appeared suddenly in the south. At the outset of the \textit{Calamity}, a society of \textit{elves} used a powerful, obscure ritual to transport many of their people to safety on another plane: the Fey Realm, a plane of primordial beauty where elven legend says the The Arch Heart lovingly fashioned their people before placing them upon the face of Exandria. They returned at the turning of the ages under the guidance of an elven sorceress named Yenlara Alderwreath.

The \textit{elves} rallied around her both for her defiant strength and for her compassion in the face of adversity. It was under Yenlara's wise rule that elven society once again began to reform. Upon returning to Exandria, Yenlara led her people westward to the \textit{Verdant Expanse}, an untamed forest born from the residual elemental power that lingered after the end of the \textit{Calamity}.

So beautiful was the forest they settled in that Yenlara's people were the first \textit{elves} known to bear the now-widespread name of syn'alfen—wood \textit{elves}, in Common parlance. Likewise, the beauty of their forest home inspired the name of their first settlement: \textit{Syngorn}, the great elvenhome that stands proud within the \textit{Verdant Expanse} to this very day.

As Yenlara's people began to explore beyond the forest, they found vast fields of grass emerging from the ash-darkened ground, with tall, snowcapped mountains visible on the horizon. Overwhelmed by the beauty of the restored world, the syn'alfen called their new land \textit{Gwessar}, the Fields of Joy, a name by which \textit{elves} (and those enamored with elven culture) still call the continent of Tal'Dorei to this day.

The \textit{dwarves} and \textit{elves} are long-lived people, and when they struggled to rebuild their civilizations, there were those among them who still remembered the world that was. Humanity was not so fortunate. Human histories, written by wasteland warlords in fading ink on waterlogged parchment and vellum, did not survive the years. Yet, somehow, humanity endured. Several centuries after the Jaggenstrike \textit{dwarves} began this period of renewal, a clan of \textit{humans} braved the angry Ozmit Sea and sailed to \textit{Gwessar's} western coast from the continent of \textit{Issylra}, though whether or not they came from Vasselheim itself is unknown.

These people had the sea in their blood, and leapt from island to island for generations, but something called them to \textit{Gwessar}. The ruins of their first settlement—the port city of \textit{O'Noa}—still stand today. From \textit{O'Noa} the seafarers expanded outward, until all of the western shores flew their banner. In the north, they found fertile fields protected from the salt of the sea, and an inlet unmarred by looming rocks and dangerous reefs, and they began to build. They did not know their city would become the heart of a great empire. They did not know the glory and sorrow that would surround the city of \textit{Emon}.

\subsubsection{The Iron Rule of Drassig}
The rise of human colonies irritated the \textit{elves} of \textit{Syngorn}. Forests that had stood for centuries fell under the axes of rash, short-lived beings who sought to exploit, expand, and propagate thoughtlessly. These tensions did not build to war, for long were the memories of the \textit{elves}, but humanity remained a thorn in the side of elfkind, chafing at their hard-won efforts to protect their society and threatening to encroach upon their sacred woodlands.

As the first human civilization on this part of the world after the \textit{Calamity}, \textit{Emon} gave rise to a handful of self-entitled noble houses who went on to establish civilization in the name of the The Lawbearer—but those who wrote the game stacked the deck in their favor. Corruption spread through the upper echelons of \textit{Emon}, and power-hungry politicians seized each new and valuable resource that was discovered in their bountiful new kingdom. They turned their citizens against each other, forcing them to fight for scraps while they hoarded the lion's share.

\textit{Emon} was a political war zone, and the greatest warrior of them all was a loudmouthed braggart and cunning oligarch named Warren Drassig. The chaos and mistrust in \textit{Emon} allowed Drassig and his agents to seize power and transform the realm into the Kingdom of Drassig, with Warren himself as its supreme monarch. Drassig was quick to sever any remaining connections with the \textit{elves} of \textit{Syngorn} and make new alliances with the \textit{dwarves} of \textit{Kraghammer}, marrying Drassig's autocratic power with the \textit{dwarves'} immense material wealth.

The \textit{elves} were furious, but the ambassador from \textit{Syngorn} to \textit{Emon}, an idealistic grandson of the still-living Yenlara, hoped to resolve this diplomatically. Upon arrival, he was apprehended, tortured, and slain. This final act of treachery drove \textit{Syngorn} to arms, and the continent erupted into a terrible, protracted war between Yenlara's kin and Drassig's bloodline known as the \textit{Scattered War}.

\subsubsection{The Scattered War}
The Scattered War lasted for thirty-two years, and is recalled in bard-song as the Time of Shrouds. The war spanned the \textit{Cliffkeep Mountains} and parts of the \textit{Verdant Expanse}, with the human colonies spread throughout the soon-to-be-warring territories. \textit{King Drassig} encountered little resistance as he conquered each village, town, and city in turn, giving these isolated settlements no way of warning one another of the warlord's advance.

No common warmonger, \textit{Drassig} was a cunning, heartless master of psychological warfare. Weeks in advance of his armies, he seeded spies throughout the land. They traveled swiftly through the secret dwarf-tunnels that crisscrossed the \textit{Cliffkeep Mountains}, and infiltrated settlements throughout the region, turning the people of each settlement against their own leaders, eating at them from within, poisoning them against one another. By the time \textit{Drassig's} armies reached their targets, they were already poised to fall. The lights of human towns and elven groves were snuffed out with equal savagery.

Yet, a mere nine years after the war began, King Warren Drassig died.

It happened at \textit{Torthil}, an abandoned village in \textit{Gwessar's} heartlands. For years, rumors had grown that \textit{Torthil} was a haven for war refugees. The king and his soldiers fell upon the village with a vengeance, seeking survivors to slaughter. As they regrouped in the town square, \textit{Drassig} stood among them, blade raised, eyes wild, and voice booming, to rally their bloodthirsty spirits. It was at the height of his wicked speech that the first arrow struck, followed shortly by a volley that clouded the sky. After nine years of suffering, Yenlara's wood \textit{elves} and the rebellious \textit{humans} of the scattered colonies had joined forces and formed an alliance against their common enemy.

After Warren \textit{Drassig} fell, he was succeeded by his eldest son, Neminar \textit{Drassig}. Neminar shunned the crude, brutal methods of his father, instead finding his interests in more sinister powers. He became known as "Neminar the Black-Fingered," as his forays into necromancy left one arm withered and tainted. Nevertheless, his weaponization of forbidden magic elevated the threat of \textit{Drassig's} war machine. King Neminar, bent on bloody-minded vengeance, led his army back to \textit{Torthil}, where the corpses of the elven and human traitors were piled high upon the outer walls. Their pyres ignited the night as the entire city was reduced to rubble. And so the Scattered War continued, unabated by the death of Warren Drassig.

The rule of Neminar marked the darkest days of the war. He introduced tactics and magic that drove his soldiers beyond human limits, leaving their bodies altered and mutated by foul necromancy. Their warped minds craved bloodshed and domination, and their bodies needed no sustenance to fuel their bloodlust. \textit{Drassig's} forces became the perfect weapons of war, their ranks blessed by the divine touch of the The Strife Emperor—the first \textit{Betrayer God} to extend his power to mortals beyond the Divine Gate. All the gods soon learned that though they could not walk upon Exandria themselves, they could still grant magic to their faithful.

In the face of such terrible power, the alliance of \textit{humans} and \textit{elves} bolstered their ranks with vengeful orphans and eager heroes of all ancestries who suffered under \textit{Drassig's} tyranny. One such hero who came to be instrumental in the coming conflict was \textit{Zan Tal'Dorei}. A human who rose from the harried streets of \textit{Syngorn}, Zan quickly showed her mettle as both a warrior and an inspirational leader. Rallying the broken ranks of the resistance, Zan lured Neminar and his lead forces into the \textit{Verdant Expanse}.

Neminar's undoing was an arrogant attack upon the \textit{Shifting Keep}, a \textit{Syngornian} outpost protected by the slippery magic of the Fey Realm. The illusory fortress vanished, leaving the \textit{Drassig} army without a target and vulnerable to ambush. The forest itself seemed to lend its fury to the weapons of Zan's warriors, and Neminar and his blighted soldiers were crushed in one swift stroke.

Word of Neminar's defeat and death spread like wildfire across \textit{Gwessar}, and the banner of \textit{Zan Tal'Dorei's} rebel army became a symbol of hope for all of the oppressed peoples of the realm.

Even so, the last years of fighting were still ahead of them, for the late King Warren Drassig's youngest son still stood to take power, and swore to avenge his family.

\subsubsection{Battle of the Umbra Hills}
\textit{King Trist Drassig}, second-born son of the despotic Warren \textit{Drassig}, took the throne with unearned confidence. So enamored was he with arrogant visions of victory that he ignored the truth of his tenuous rule: he was neither as brilliant nor as charismatic as his father, his armies were stretched thin, and his people rioted. The rebels, led by the young warrior \textit{Zan Tal'Dorei}, won battle after battle, and before long, King Trist found his army backed against the imposing base of the \textit{Cliffkeep Mountains}.

In the wake of certain victory, Zan and her rebels, allied with the \textit{elves} of \textit{Syngorn}, pursued \textit{Drassig} to this valley, but were met with fiends amidst the ranks of their enemy. King Trist had a secret weapon—through his family's dealings with the The Strife Emperor, the spawn of the Betrayers had returned to the world. They spilled into the battlefield like a river of nightmares, and in minutes, the surrounding hills ran dark with blood and ichor, and the bodies of \textit{humans} and demons alike littered the battlefield. Yet, despite all odds, the hero Zan defeated King Trist, ending the \textit{Drassig} bloodline—and with it, the demonic pact the \textit{Drassigs} had made. It is said that the grass and flowers of the \textit{Umbra Hills} grow black and burnt as an echo of this battle, their sap coursing with the fiendish blood that was spilled that day.

\subsection{Tal'Dorei Ascendant}
Society had all but collapsed once more amidst the \textit{Scattered War}. Within the \textit{Verdant Expanse}, the triumphant rebels assembled a council of trusted and proven minds, but the people were accustomed to a singular leader, a king. In the hearts of the people, it was not the crown that had failed them, but the bloody-minded family that wore it. The council assented to the will of their people, and nominated the war hero Zan Tal'Dorei to take the throne. She humbly accepted the role, but refused to take the title of king or queen. Her name, she asserted, was Zan. After some debate, she eventually relented to being called Sovereign Tal'Dorei, if such formality was required.

Again in spite of Zan's protestations, the council unanimously agreed the realm should be renamed Tal'Dorei in her honor. Power was divided between the sovereign and the council, both of whom ruled from \textit{Emon}. From there, \textit{Emon's} leaders expunged any remnants of \textit{Drassig's} reign from their city, while the allied armies of Tal'Dorei and \textit{Syngorn} did the same throughout the realm. The leading clans of \textit{Kraghammer} pleaded for clemency, claiming they were coerced to serve \textit{Drassig}. They spent many years paying reparations and making amends, but the fractured trust between \textit{Gwessar's} peoples took generations to heal.

\begin{DndReadAloud}{}
\DndQuote%
{}
{''A king should not stand apart from their people. A ruler must know sacrifice, for how else are they to know those who till the fields, who tend to the cattle, and who raise our children? A ruler should know pain and sorrow, for war is bloody and cruel. A ruler who asks their people to fight is one who asks them to die, and no person should die for a selfish king or an unjust cause. A ruler should know justice, but also mercy, for not all crimes are committed with evil intent, but rather in grief and shame. I am no king. I am the peasant in the fields. I am the painter at the easel. I am the young soldier donning their father's rusted armor. You will not find me sitting upon a bejeweled throne disdaining my supplicants, but working beside your mothers and fathers, fighting beside your valorous children, and aiding those who have suffered in the name of unjust kings.''}
{''Henceforth, Tal'Dorei shall not be just a name: it is a principle. From this day, Tal'Dorei is a principle that stands for honor, an everlasting honor that binds its people. A principle that is a beacon for the lost, the helpless, and the forgotten. Tal'Dorei is a principle that you should declare with honor as you stand side-by-side with the people of Tal'Dorei as kin. Tal'Dorei is not a name; it is to be proud, it is to have honor, it is to be loyal, and loving, and kind. Tal'Dorei is not a name; it is a people. And you are the people of Tal'Dorei! I refuse the title of king. If you wish me to rule, I shall do so as a sovereign of the people—and if you permit me, if you accept me as such, the Sovereign of Tal'Dorei! For the Honor of Tal'Dorei!''}
{}
\end{DndReadAloud}

\subsubsection{The Icelost Years}
Not two years into Sovereign Zan's reign, a cataclysmic evil tested her mettle and the skill of the council. A creature from the Elemental Plane of Ice slunk into the Material Plane through a hitherto unknown rift between the planes. It appeared in a mystical forest of eternal winter known today as the \textit{Frostweald}. From its hibernal abode, this elemental spirit spied upon the untested realm of Tal'Dorei and saw an opportunity. The spirit dove back to its extraplanar home and informed its master of a land ripe for conquest. Errevon the Rimelord, a cruel elemental so mighty that he is speculated to have been a scion of the Primordials, one that escaped the gods' wrath in the \textit{Founding}, seized this opportunity.

On the next winter solstice, Errevon tore wide the \textit{Frostweald} rift and commanded the endless blizzards of his home plane to ravage the prosperous land around it. Tal'Dorei's southern reaches, thus weakened, were easy pickings for the Algid Legion, a mighty army of frost giants and animate blizzards led by the Rimelord himself. Their relentless advance covered all they touched in snow and death.

Countless brave warriors and civilians alike fell to the invaders' gelid weapons, for they were unprepared for such an onslaught so soon after the \textit{Drassigs'} protracted war. The storm of ice widened year by year, consuming the Tal'Dorei heartland as Errevon claimed all lands that fell beneath the ice as his own. North of the rift, he built a towering citadel of frost, forcing those he conquered to swear fealty in exchange for warmth and liquid water. The Rimelord's tyranny lasted for three long years as the nascent \textit{Council of Tal'Dorei} struggled to convince \textit{Syngorn} and \textit{Kraghammer} to ally for the sake of the realm. This monumental diplomatic mission completed, the combined might of Tal'Dorei, \textit{Syngorn}, and \textit{Kraghammer} assaulted Errevon's Citadel and forced the Rimelord back to the rift from whence he came. It was in this time of need that the druidic Ashari people first revealed themselves to the people of Tal'Dorei, and used their unrivaled elemental magic to seal the rift once and for all. The snow began to melt, the fortress toppled into the nearby \textit{Foramere Basin}, and the people celebrated their freedom from the oppressive cold. This victory is now celebrated annually throughout Tal'Dorei as the Winter's Crest festival.

\subsubsection{Thordak, the Cinder King}
\textit{Zan Tal'Dorei} passed peacefully after a long and happy rule, and was succeeded by a long line of her descendants. Many generations of peace ensued, and the Tal'Dorei Empire formed as dozens of tiny city-states across the \textit{Dividing Plains} flocked to the stable rule and heroic reputation of Tal'Dorei in the wake of the \textit{Icelost Years}.

The adroit statecraft of the \textit{Council of Tal'Dorei} maintained cordial alliances with both \textit{Syngorn} and \textit{Kraghammer}, but, despite their best efforts, couldn't keep the \textit{elves} and \textit{dwarves} from regressing to hostility over ancient quarrels. In the time of Sovereign Odellan Tal'Dorei, word of a shadow in the south came to the ears of the council. Contact with numerous outlying townships at the edges of the empire halted abruptly, and traders that made circuits from \textit{Emon} to the \textit{Rifenmist Peninsula} in the south vanished without a trace. The council was slow to act in defense of territories that routinely flaunted imperial law, but were finally driven to action when allied villages south of the \textit{Verdant Expanse} begged for protection against a "nightmare of fire and malice." The dour Sovereign Odellan blocked the council's relief efforts, claiming the economic cost of sending a regiment that far south to outsider communities was a misuse of resources, especially without proof of a threat. It wasn't until two years later, when reports of a powerful red dragon reached the ears of the sovereign, that the \textit{Council of Tal'Dorei} was able to act.

By then, much of the \textit{Mornset Countryside} had been reduced to ash by a self-obsessed and power-hungry red dragon named \textit{Thordak}, who preferred the self-styled moniker of "the Cinder King." According to \textit{Cobalt Soul} records from the Temple of the Mentor in the lands of \textit{Marquet}, the \textit{Cinder King} was presumed dead nearly two hundred years prior, killed by the brass dragon Devo'ssa—yet it could not be denied that \textit{Thordak} survived, and that he sought new peoples to enslave, ones with fewer defenses.

The mighty armies of Tal'Dorei marched finally south, but morale was low. They marched toward fiery, painful death, and soldiers grumbled of a heartless sovereign who luxuriated in his castle in \textit{Emon} while his people put their lives on the line for the realm. They finally clashed with \textit{Thordak} and his mercenary armies south of the \textit{Stormcrest Mountains}, but the forces of Tal'Dorei were soon forced into a panicked, disorganized retreat as the rocky, exposed terrain worked to the advantage of their airborne foe.

The valiant soldiers of Tal'Dorei would surely have been brought to utter ruin were it not for the heroic intervention of a band of young—yet able—adventurers. Led by \textit{Arcanist Allura Vysoren}, who sits today upon the \textit{Tal'Dorei Council}, they cornered the \textit{Cinder King} alone and defeated him—though his power was so great that they could not manage to kill the dragon outright. Instead, \textit{Allura} and her companions used powerful magic to imprison \textit{Thordak} within the Elemental Plane of Fire forever. The Fire Ashari of distant \textit{Issylra} vowed to watch over \textit{Thordak's} prison until age could consume even him.

\subsubsection{The Chroma Conclave}
Forever lasted a mere sixteen years. In that time, life returned to mundane quarrels over the price of bread and taxes. During this time of quietude, an old and unexpected ally of \textit{Thordak's}, a green dragon called Raishan the Diseased Deceiver, found \textit{Thordak} and conspired to release him from his prison so he could conquer the land that had once rebuked him. Raishan formed a plan and forged an unprecedented alliance with three other power-hungry chromatic dragons: Umbrasyl the Hope Devourer, Brimscythe the Iron Storm, and Vorugal the Frigid Doom.

After infiltrating the Fire Ashari of Pyrah for nearly five years, Raishan finally unlocked the seal between the planes that bound \textit{Thordak} and set him free. \textit{Thordak's} physical body was warped by the fiery energies of his prison. Now stronger than ever, he rallied his draconic allies, the Chroma Conclave, to destroy or conquer Tal'Dorei as they saw fit—starting with a retributive assault of catastrophic proportions on \textit{Emon}. Nearly three centuries of Tal'Dorei rule in the realm ended under the rule of Sovereign Uriel Tal'Dorei II, his life cut short when much of \textit{Emon} was reduced to rubble.

The attack not only killed the Sovereign Uriel, but also scattered the \textit{Council of Tal'Dorei}, throwing the realm into chaos. The remnants of the civilized lands were divided up and taken as trophies by the Conclave, while the capital \textit{Emon} was ruled by the \textit{Cinder King} himself, with Raishan secretly manipulating the course of his rule from the shadows.

The \textit{Alabaster Lyceum} need hardly introduce the heroes of this story. One by one, the members of the Conclave fell to the might and cleverness of a band of misfit warriors known as Vox Machina, their might bolstered by a number of the lost \textit{Vestiges of Divergence}, artifacts of incalculable power forged amidst the chaos of the \textit{Calamity}. After many adventures, these heroes gathered their allies and stormed the capital of \textit{Emon}, slaying \textit{Thordak the Cinder King} and freeing its people. Upon discovering the machinations of the real mastermind, Raishan the Diseased Deceiver, with whom Vox Machina had held a tenuous alliance, they gave chase and finally slew this last standing member of the Chroma Conclave, ending their reign over the land and restoring rule once more to the \textit{Council of Tal'Dorei}.

\subsection{A Fledgling Republic}
The final act of Sovereign Uriel Tal'Dorei II's life was to abdicate the throne and end the line of sovereigns—permanently. To analyze the historic Last Sovereign's final acts is the duty of biographers, not chroniclers, yet it can readily be surmised that he saw flaws inherent to his station. Uriel was the son of Odellan Tal'Dorei, the sovereign who singlehandedly prevented the council from swiftly handling the \textit{Cinder King} when first he rose in the wilderness south of the empire.

How could the Last Sovereign not feel that it was his office's immense power to overrule the council that ultimately cost thousands of loyal citizens their lives? How could he not compare the singular power of the sovereign to the abuses of power committed by \textit{King Drassig} in ancient times?

Regardless of his private feelings, Uriel II's last act as sovereign transformed his family's imperial dynasty into the Republic of Tal'Dorei. \textit{Zan Tal'Dorei} was reluctant to accept the singular role of rulership at the dawning of the nation, an ideal she had to compromise upon to ensure stability for her fledgling nation. Uriel's relinquishment of power is seen today as a fulfillment of that ideal. It will be the duty of future historians to judge if his decision stood the test of time.

\subsubsection{The Time of Regrowth}
The \textit{Tal'Dorei Council} has guided the fate of the nation since Vox Machina slew the \textit{Cinder King}. For most of the past two decades, the goal of the council has been to rebuild that which was lost—and to assure far-flung territories that the nascent republic is just as stable as the empire ever was.

Recovery has been slow in both Tal'Dorei's heartland and its most distant reaches. Even though cities like \textit{Kymal} and \textit{Stilben} were not direct targets of the dragons' rampage, the age of unrest that followed made these small cities easy targets for ransacking by roving bandits and corrupt magistrates alike. Even today, outlying settlements are slow to trust the good intentions of the \textit{Tal'Dorei Council}.

Major cities like \textit{Westruun} and \textit{Emon} have been the greatest beneficiaries of the Time of Regrowth, for even though they suffered the worst of the \textit{Conclave's} destruction, the dragons also kept their lairs close by, allowing vast amounts of plundered wealth from their treasure hoards to flow rapidly back into these cities' coffers. Over the last twenty years, these cities have hired legions of arcanists to magically restore much of what was lost.

However, perhaps the one city that can say it directly benefited from the \textit{Chroma Conclave's} reign of chaos is \textit{Whitestone}, a formerly independent city-state that is now a member of the republic. \textit{Whitestone} took in countless refugees from \textit{Emon} after the \textit{Conclave's} initial attack, and countless more in the aftermath. Sequestered in Tal'Dorei's cold northeastern reaches, \textit{Whitestone} is an isolated city, but it has rapidly become a powerhouse of political, cultural, artistic, and economic influence—to the chagrin of some traditionalists in \textit{Emon}.

\subsubsection{Apotheosis Thwarted}
The world was troubled by a disturbance in \textit{Issylra} a year after the \textit{Chroma Conclave's} defeat, as a lich attempted to achieve godhood and crush the Dawn City of Vasselheim in one fell swoop. His success would have meant catastrophe for the Material Plane, leaving him the sole god on this side of the world's impenetrable Divine Gate. However, beyond hearing news of Vox Machina's involvement in this conflict, few in Tal'Dorei were concerned by rumblings on the other side of the world.

Though the lich's defeat confined him to a remote demiplane, there is no way to strip him of his stolen divinity. Occult worshipers yet scurry throughout Tal'Dorei's darkest corners, seeking the favor of a lich-god known as The Whispered One.

\subsubsection{Seeds of Peril}
A generation has passed in relative peace. Yet all who have studied history know that each age's conflict sows the seeds of the next. The Republic of Tal'Dorei has proven it can prosper in peacetime, but it has yet to be tested by the pressures of war or a supernatural cataclysm.

After two decades of expensive reconstruction and princely investments in \textit{skyship} moorings, the Republic of Tal'Dorei is deeply indebted to arcanists across the land. Chief among them is the \textit{League of Miracles}, a consortium of thaumaturgists who have claimed untold fortunes from reconstruction bounties all over the country.

The simple fact is that the cost of rebuilding Tal'Dorei was far more than the council or local leaders could afford to pay. Far be it from the \textit{Alabaster Lyceum} to cast aspersions on so beneficial an organization, but there are fears among the Tal'Dorei people that the league has already begun to collect on what they are owed in forms other than gold, as more and more league mages receive preferential treatment from local margraves and even members of the \textit{Tal'Dorei Council} itself.

Beyond \textit{Emon}, the landscape of Tal'Dorei has been forever scarred by the presence of the \textit{Conclave}, especially the \textit{Cinder King}. His hellfire left marks on the land that will never heal—magical wounds that bleed fire and ooze chaotic magic. Smaller cities beyond the reach of the council have begun to succumb to sinister influences. Demons and devils crawl into the world, summoned by occult powers answering fell voices from beyond the Divine Gate. The \textit{Betrayer Gods} sense weakness within the world. The entropic fury of the The Chained Oblivion churns beneath \textit{Gatshadow}. The lost blade of Neminar the Black-Fingered calls from the \textit{Umbra Hills} for a new master. The The Cloaked Serpent raises armies south of \textit{Syngorn}. Cults of the now-divine The Whispered One spread like poison through the veins of Tal'Dorei's cities.

Yet despite these threats, the legendary heroes of Vox Machina have largely set down their blades, guns, and spells and taken up the mantle of leadership. More than twenty years have passed since their greatest adventures, and though the surviving members of Vox Machina are still canny rulers and cunning warriors, the \textit{Tal'Dorei Council} can no longer rely upon their skills as readily as they once did. Rumors of the heroes' permanent retirement swirl, inspiring fear and nostalgic sorrow in many of the people—but adventurers across the land see this as an opportunity to rise to the level of Vox Machina and beyond.

Tal'Dorei needs heroes—and this time, Vox Machina alone will not be enough.

\section{Running a Tal'Dorei Campaign}
Tal'Dorei is more than just its history; do not treat this book as an encyclopedia, nor as shackles that bind you to portray its characters and environments with perfect fidelity. As you read this manual, consider why you want to run an RPG campaign set in Tal'Dorei, and in the larger world of Exandria. What about \textit{Critical Role} inspires you as a Game Master? You might cherish the bonds of family and love between its many heroic characters. Perhaps you are drawn to the various factions that vie for power over Tal'Dorei's peoples. Or perhaps it's the way that evils from Exandria's mythic past inevitably rear their ugly heads, always when circumstances are already at their most dire, that most captures your attention.

No matter the initial spark that led you here, hold tight to it as you read, and search for ways to feed your personal flame of imagination. And when you are done, create. Create boldly, using these words as a launching point for your own tales.

This section is for Game Masters' eyes only, though GMs are welcome to share tidbits with their players to set the stage for the campaign, or to inspire them to create characters that will dovetail with the GM's story ideas.

\subsection{Secrets of Tal'Dorei}
The world of Exandria is steeped in familiar fantasy themes, with all the \textit{elves}, \textit{dwarves}, and wizards you might expect from a fantasy roleplaying game. But it also holds secrets that will challenge your expectations, and you may enjoy putting your own twists on these great mysteries of the world. Each of these secrets can form the backbone of an entire campaign, or can serve as flavor to create the impression of a world that is larger than just the player characters' story.

\subsubsection{All-Powerful, Imprisoned Gods}
Tal'Dorei was nearly destroyed when the gods fought one another in the \textit{Calamity}. Now, the \textit{Betrayer Gods} are sealed in otherworldly prison planes, and the victorious \textit{Prime Deities} live in self-imposed exile. An unbreakable barrier stands between the mortal and immortal realms. At best, a deity can grant their most skilled devotees a small measure of divine power to enact their will upon Exandria.

However, there are a number of cults who yearn to circumvent this barrier and summon their dark patron onto the Material Plane. Others still seek to unravel the forgotten rites of ascension, just as the The Matron of Ravens did in ancient times—just as The Whispered One did in living memory. What would happen if an omnipotent deity somehow managed to arise on the mortal side of the Divine Gate? Or even simply managed to fool the gods into thinking they had done so? Would the gods of good be forced to unleash their powers as one to break the Gate, unleashing another \textit{Calamity} upon the world?

\subsubsection{Dawn of a Republic}
It was the dying decree of Sovereign Uriel Tal'Dorei II that the continent-spanning Tal'Dorei Empire be reborn as the Republic of Tal'Dorei. No more would a sovereign rule the myriad peoples of this land, but instead the wise \textit{Council of Tal'Dorei} would enact the will of the people and protect them from the dangers of the world. But things don't always go according to plan. The nascent Republic of Tal'Dorei is suffering growing pains, and unless some heroes are able to help, the realm may fall into anarchy and chaos.

While fantasy politics isn't every gamer's cup of tea, the struggles of the fledgling republic can serve several purposes. Players looking for fantasy intrigue can meddle in the internal politics of the city of \textit{Emon} and root out corruption within the council, which may have been planted by any number of this land's power-hungry factions, from the invasive crime syndicate known as the \textit{Myriad} to the shadowy mages within the \textit{Arcana Pansophical} who splintered off into a new schismatic faction, the \textit{League of Miracles}.

Game Masters can also use political elements to flavor or provide a backdrop for their otherwise action-focused games. On a smaller scale, corrupt margraves may hire bandits to pillage their own villages, requiring heroes to step in. On a broader scale, a threat to the stability of the Republic of Tal'Dorei itself may force the characters to stop a megalomaniac bent on conquering the republic from within.

\subsubsection{Criminal Shadow War}
A criminal syndicate known as the \textit{Clasp} has lurked in Tal'Dorei's underworld since the foundation of the old Tal'Dorei Empire. It's almost become a cliché, for people across the land know the \textit{Clasp} as the stereotypical face of organized crime. Novels romanticizing their allure have been adapted into stage plays and bard-song in all corners of the republic.

And perhaps the \textit{Clasp} has bought its own good publicity. After the \textit{Cinder King's} famous conquest of \textit{Emon} some twenty years ago, it was the \textit{Clasp} that saved countless innocent people and led them to safety through their extensive network of smuggling tunnels beneath the city. It was the organized action of criminals, just as much as Vox Machina's heroic power, which saved Tal'Dorei—and the \textit{Tal'Dorei Council} knows it. The \textit{Clasp} has gone mainstream.

However, the \textit{Clasp} is under threat from the \textit{Myriad}, another criminal syndicate from the nearby continent of \textit{Wildemount}. The \textit{Myriad} sees Tal'Dorei's slow legitimization of the \textit{Clasp} as an opportunity to seize power in the nation's criminal underground. Violence is on the rise in Tal'Dorei's cities as the \textit{Myriad} uses local gangs as pawns in their bid to dethrone the \textit{Clasp}.

\subsubsection{A Land Thrice-Destroyed}
First there was the \textit{Calamity}, and the palaces of the gods were devoured by the earth. Then came the wars of \textit{Drassig}, and his defeat at the hands of \textit{Zan Tal'Dorei}. And then came \textit{Thordak}. The \textit{Cinder King} and the \textit{Chroma Conclave} shook the very foundations of Tal'Dorei. \textit{Thordak} "merely" reduced the city of \textit{Emon} to slag, but through fear, panic, and malice, the flames of destruction were spread across the continent by mortal hands. Today, a generation after the dawning of the Republic of Tal'Dorei, the ruins of three civilizations lie buried beneath your feet, and some secrets long to be unearthed.

\subsubsection{Mother of Invention}
\textbf{Percival de Rolo}—\textbf{Percival de Rolo}, to his friends—wanted to take the secrets of his inventions to his grave, especially the deadly power of his firearms. Unfortunately, because of his now-deceased archenemy Anna Ripley, the secret of gunpowder has begun to spread across Tal'Dorei. Now, isolated groups of tinkerers, both noble and unscrupulous are working to develop new technologies from \textbf{Percival de Rolo} research, and other black powder technologies from \textit{Marquet}. Some of these modern marvels will no doubt help the people of Exandria—but others, if placed in the wrong hands, are sure to enable the deaths of thousands.

Renaissance technologies, especially firearms, are not common in Tal'Dorei. If you want to keep your fantasy purely medieval, it's a simple matter to ignore the tinkerers, gunsmiths, and budding industrialists scattered across the continent.

\subsubsection{Uncharted Territory}
The corners of Tal'Dorei's maps are not all filled in. North of the \textit{Cliffkeep Mountains}, the trackless \textit{Neverfields} push up against the northern edge of the map. South of the \textit{Verdant Expanse} is the \textit{Rifenmist Peninsula}, where vast jungles and the brutal forces of the \textit{Iron Authority} prevent any explorer from filling in the map's southern tips.

To help you create stories that feel at home in Tal'Dorei, \textit{chapter 3} of this book contains a number of story hooks throughout its description of Tal'Dorei's landscape and settlements. See "\textit{Creating Adventures}" for more information on how to turn these plot seeds into full adventures that you can run for your gaming group.

While this book gives insight to one version of Exandria, its lore exists to serve you, not the other way around. Your campaign is your own creation—if you delight in adhering to \textit{Critical Role} canon and delving into its lore, then by all means, dig deeper! If you like subverting expectations, revealing that things believed to be true were nothing but myth and rumor all along, tell that story! And, if you love characters and places from \textit{Critical Role}, but want your stories to go to brand-new lands of your own creation, take inspiration from this book and steer your campaign in whatever direction you see fit. Never forget that you are the creator of your own version of Tal'Dorei. No matter who is playing at your table, even if it's Matthew Mercer himself, no one can tell you that your vision of the world is inaccurate—because it's yours now.

\chapter{Allegiances of Tal'Dorei}
\DndDropCapLine{T}{}he lands of Tal'Dorei are governed not just by their history, but by factions and faiths striving to achieve myriad goals, fighting for conflicting ideals, and vying for the power to come out on top. The gods of Exandria may be barred from walking upon the land in material form, but they can still put their thumbs on the scales of history through their mortal followers. Such followers include clerics with the power to cast mighty spells, paladins who wield blades alight with idealistic verve, and untold thousands of common folk whose voices can affect the decisions of even the most powerful leaders of the land.


A number of factions hold just as much sway over the daily lives of the people as the gods and their clergy—perhaps even more. These factions range in might and influence from the greatest of the great to the smallest of the small. The \textit{Tal'Dorei Council} institutes political policies that affect everyone in grand, sweeping, and often inscrutable ways, while the \textit{Golden Grin} is a group of folk heroes who seek to subvert evil across the land. Whether these factions are governments or just tiny organizations, they can add depth and conflict to your Tal'Dorei campaign.

\section{Pantheon of Exandria}
During the \textit{Divergence}, the \textit{Prime Deities} created a powerful barrier between the Material Plane and the supernatural realms of the Outer Planes. This gate sealed all gods, including the evil \textit{Betrayer Gods}, within their respective domains. If the Divine Gate were to be destroyed by the unanimous effort of the \textit{Prime Deities}, all powers would be unleashed, threatening armageddon. Thus, the gods patiently watch their creations from beyond the veil, aiding their faithful with what small power they can send through the barrier.

Some lesser idols, such as archdevils, demon lords, fiends and celestials of near-divine power, as well as a number of unaccounted-for demigods reside among the planes. They are affected by the Divine Gate, preventing them from crossing it of their own volition, yet are not so bound by it that they cannot be summoned by mortal magic.

The following gods are recommended as the existing Exandrian pantheon, but they are only a recommendation. As long as the GM and the players agree, you are welcome to tailor, alter, or completely change the gods in your Tal'Dorei campaign to fit your needs. The domains listed are the likely choices for followers of that deity, but these are not the only options.

\subsection{Divine Domains}
Most clerics can gain their magic from a single god, or even from a number of gods that share domains or divine provinces. Ultimately, how your cleric serves the gods or their core religious philosophies is for you and your Game Master to discuss. Each god in this chapter has two or three Divine Domains from which many of their clerics are granted power.

Nevertheless, almost all deities can have clerics who gain the power of any domain, if you and your Game Master agree.

\subsection{Prime Deities}
The Prime Deities are the leaders and luminary creators that battled the Primordial Titans and instigated the \textit{Founding}, forging the mortal races of Exandria. They hold sway over powers that represent a spectrum of sunlight, truth, moonlight, benevolent concealment, protection, love, death, and all other facets of freedom and life in the world. While these gods may disagree and squabble, they are united in an alliance dedicated to ensuring the survival of their creations.

\begin{DndTable}[header=Prime Deities]{XXX}

Deity & Alignment & Province\\
The Changebringer & Chaotic good & Change, freedom, luck\\
The Platinum Dragon & Lawful good & Honor, justice\\
The Arch Heart & Chaotic good & Art, beauty, \textit{elves}\\
The Lawbearer & Lawful neutral & Civilization, law, peace\\
The Knowing Mentor & Neutral & Knowledge, learning, teaching\\
The Stormlord & Chaotic neutral & Battle, competition, storms\\
The Wildmother & Neutral & Seas, wilderness\\
The All-Hammer & Lawful good & Craft, creation\\
The Dawnfather & Neutral good & Healing, sun\\
The Everlight & Neutral good & Atonement, compassion\\
The Matron of Ravens & Lawful neutral & Death, fate, winter\\
The Moonweaver & Chaotic good & Illusion, moonlight, night\\
\end{DndTable}

\subsection{Betrayer Gods}
The \textit{Betrayer Gods} are the deities who strayed from the ideals of the \textit{Founding} and embraced the destructive chaos of the Primordials. The \textit{Betrayer Gods} rarely work together, for they see one another as threats to their own goals. This very weakness allowed the \textit{Prime Deities} to defeat and banish them, ending the \textit{Calamity}.

\begin{DndTable}[header=Betrayer Gods]{XXX}

Deity & Alignment & Province\\
The Lord of the Hells & Lawful evil & God of devils and the Hells\\
The Strife Emperor & Lawful evil & Conquest, tyranny\\
The Ruiner & Chaotic evil & Slaughter, warfare\\
The Spider Queen & Chaotic evil & Deceit, spiders\\
The Chained Oblivion & Chaotic evil & Darkness, destruction\\
The Scaled Tyrant & Lawful evil & Dragon god of evil\\
The Crawling King & Neutral evil & Enslavement, torture\\
The Whispered One & Neutral evil & Necromancy, secrets\\
The Cloaked Serpent & Chaotic evil & Assassins, poison, snakes\\
\end{DndTable}

\subsubsection{Betrayer Curses}
The malign powers of some \textit{Betrayer Gods} linger in Exandria, both as myth and reality.

This book adds the following curses:


\begin{itemize}
\begin{multicols}{2}
\item Curse of Strife
\item Curse of Ruin
\end{multicols}

\end{itemize}

\subsection{Lesser Idols}
Though the gods still make their wills known upon the world by granting power to mortal clerics and by delivering edicts through their celestial or fiendish servitors, the Divine Gate inhibits them from exercising their full powers over the fate of Exandria. The same is not true for the lesser idols, beings of demi-godly power that have, in the absence of true gods, used their splendid and profane forms to amass their own cult followings.

Dozens of these powers are worshiped in relative secret across Exandria. Some are supernatural beings like archfey or archfiends that have taken up near-permanent residence on the Material Plane. Others are unknown fragments of the power of primordial Creation, given form. Most are mortal champions that have gained near-divine power through magic, the blessings of the gods, or through mythic effort and training. In all of these cases, the worship of mortals elevates these beings to something greater than they were before, imbuing them with fragments of divine power—and allowing them to grant scraps of that power back to the mortals in whom they take an interest.

\begin{DndTable}[header=Lesser Idols]{XXX}

Deity & Alignment & Province\\
Azgrah & Lawful neutral & Undying\\
Galdric & Lawful good & Celestial, Undying\\
Graz'tchar & Chaotic evil & Fiend, Hexblade\\
Khedive Xundi & Chaotic neutral & Genie\\
The Lost Beacon of Unknown Light & Neutral & Celestial\\
The Observer & Neutral good & Archfey, Celestial\\
The Sightless & Chaotic evil & Great Old One\\
The Traveler & Chaotic neutral & Archfey\\
Vesh & Neutral evil & Archfey, Undying\\
\end{DndTable}

\begin{DndSidebar}[float=!b]{Other Secluded Champions}
A handful of other mortals exist in Tal'Dorei who have risen to nearly divine power through destiny or sheer will. Much like Azgrah, the vast majority of these champions live out their lives in seclusion—and perhaps it is their seclusion that grants them the ability to lend their powers to others, while they await the day that their full strength is at last needed to defend their people.



All have their own stories, but no matter their tales, whether they are honored or despised by their kinsfolk, they are all mythic iconoclasts. Other secluded champions of Tal'Dorei that you may wish to develop into idols of your own include:




\begin{itemize}
\item Flavia, Queen of Unending Storms (\textbf{cloud giant})
\item Suthine, Eternal Princess of Vulkanon (\textbf{fire giant})
\item Valdemar, Draugr-Jarl of Farborg (\textbf{frost giant})
\item Ghaladron, Traitor to the Iron Crown (\textbf{hobgoblin})
\item Tevrosk, Whose Axe Was Wreathed in Flowers (\textbf{orc})
\item The First Favored Champion (\textbf{stone giant})
\item Typhoe, the Dreaming Shark (\textbf{storm giant})
\end{itemize}



\end{DndSidebar}

\section{Factions and Societies}
While the dominant cultures of the continent have molded structure and law from the rough clay left beyond the \textit{Divergence}, these very cultures are comprised of a number of smaller factions that together create the whole. Whether they be a union of opportunists, ever-vigilant to reap the benefits of an exploitable populace, or bands of like-minded politicians, striving to maintain control in a land beset by chaos and danger, each faction seeks to drive and manipulate the social and political direction of Tal'Dorei. The alliances and tensions that can arise between the following societies can guide your adventurers' destiny across the spectrum of heroism and villainy.

\subsection{Tal'Dorei Council}
When the new reign of Tal'Dorei was established three centuries ago, \textit{Zan Tal'Dorei} deposed a brutal tyrant. Though she was beloved by the people, she did not wish to be given unchecked or un-counseled rule over the realm. Nevertheless, her inner circle obeyed the will of the people and crowned her Sovereign of Tal'Dorei. She took to her new role well and swiftly expanded her council by appointing some of her trusted allies—and even some dissenting voices—to form an assembly of scholars, generals, and philosophers to govern at her side. This \textit{Council of Tal'Dorei} worked alongside their sovereign to rebuild the realm following the end of \textit{Drassig's} line, and were instrumental in the following prosperity of the nascent Tal'Dorei Empire.

Twenty-six years ago, Sovereign Uriel Tal'Dorei II, the latest and last of the Tal'Dorei line, publicly ended the monarchy and passed power over to the council, formally divesting his line of power. Tragically, this ceremony was interrupted by the arrival of the four ancient dragons known as the \textit{Chroma Conclave}. The dragons destroyed the city of \textit{Emon}, and the man who had until moments ago been sovereign was one of the first casualties. Uriel Tal'Dorei is survived by his wife Salda, and three children—\textit{Odessa}, Illiya, and Gren—who have used their family's still considerable wealth and power to become well-known public figures in \textit{Emon}. \textit{Odessa} herself won a seat on the council in a landslide victory.

In the generation since the fall of \textit{Emon}, the council has risen to the challenge of rebuilding and governing their new republic. This state's power is highly decentralized; each individual city-state in Tal'Dorei largely acts according to its own interests, with the council in place to ensure that pursuing those interests benefits the common good. Some critics of the republic say that the \textit{Tal'Dorei Council} is little more than a glorified city council for \textit{Emon}, though it is just as often praised for governing lightly as a rule, and for using the spy network helmed by the council's Master of Law to root out corruption as needed.

In addition to the historical Six Masters whose positions have existed on the \textit{Tal'Dorei Council} since Zan's time, a variable number of delegates also cast votes in matters of state. Each settlement in Tal'Dorei whose population numbers at least five thousand is offered the option to elect delegates to the council if they wish to engage in national politics. The settlement is permitted one delegate per five thousand people, and the means of their election is left to the settlement's local government. Additionally, any settlement with a population of over twenty thousand can appoint an ambassador to \textit{Emon}; these ambassadors also sit on the council and have voting power. Once in \textit{Emon}, each township's delegation is expected to negotiate with one another and find agreement. If an accord can't be reached, they must go through the slow process of drafting proposals to be voted on by seated members of the council.

Though \textit{Kraghammer} and \textit{Syngorn} remain sovereign city-states, the invasion of the \textit{Chroma Conclave} convinced lawmakers that even tighter alliances between the peoples of Tal'Dorei were necessary. One ambassador each from both \textit{Kraghammer} and \textit{Syngorn} now also holds a seat with voting power on the council. Each city-state has likewise sent a delegation to \textit{Emon} to vouch for their interests.

\subsubsection{Goals}
Following \textit{Zan Tal'Dorei's} coronation, the \textit{Tal'Dorei Council} existed to aid and inform the sovereign on how best to execute their duties. The council has evolved to manage the affairs of the city of \textit{Emon} and provide a private forum for delegates from across the land to shape the future of Tal'Dorei. The council is made up of six masters, whose roles have existed since the time of \textit{Zan Tal'Dorei}, and a variable number of ambassadors from all city-states across the continent who wish to participate in the coalition of power that is the Republic of Tal'Dorei.

The Six Masters are elected by a popular vote of all city-states who have sent an ambassador to the council. These masters are each beholden to a domain of responsibility, and they either inherit the existing infrastructure of civil servants, or restructure with the approval of the rest of the council. The domains are divided as such: Development, Arcana, Law, Commerce, Information, and Defense. Additionally, the first business of the council upon delivering or receiving a declaration of war is to elect a Seventh Master: the Master of War, who is appointed to the council by a vote of all sitting masters, ambassadors, and delegates. The council members are responsible for these domains throughout the city of \textit{Emon} and work closely with community leaders in other cities to maintain order.

\textbf{The Master of Development} works with the masonry and carpentry guilds to oversee all major infrastructure projects within \textit{Emon}, as well as allocating resources to projects needed in settlements that have a delegation in the council. In the past two decades, the current Master of Development has added an unprecedented eight Lieutenant Masters of Development to his staff to serve the reconstruction needs of each city in the republic.

\textbf{The Master of Arcana} keeps vigil over the lawful use of magics within Tal'Dorei, informs the council on matters beyond the mundane, and occasionally works with the \textit{Alabaster Lyceum's} Headmaster. This council member is currently embroiled in debates over the danger that the \textit{League of Miracles} poses to organizations like the \textit{Arcana Pansophical}.

\textbf{The Master of Commerce} works with all major merchant guilds and trading networks within \textit{Emon} and beyond. They are to maintain healthy trade at home and abroad, especially between squabbling delegates, and establish the government's annual budget.

\textbf{The Master of Information} manages foreign diplomacy, and also serves as the head of a network of spies that, these days, is supplemented by the skilled eyes of the \textit{Clasp}. Their work as spymaster and chief diplomat seeks to root out threats to the republic from within and without.

\textbf{The Master of Defense} is head of the \textit{Emon} Guard, High Commander of \textit{Fort Daxio}, and manages the allocation of federal resources to city guard groups throughout Tal'Dorei. They work closely with the Master of Law to mobilize armed troops throughout the country if supernatural or nonmilitary threats arise during peacetime.

\textbf{The Master of War} is a temporary position only instated during times of open or impending warfare. As Tal'Dorei has only a small standing army, this newly appointed master then assembles a military from city guards and town militias across the nation, mobilizing the elite warriors of \textit{Fort Daxio} as needed. When peace is declared, the position is dissolved until it is called for again.

\subsubsection{Relationships}
The council maintains free travel, free trade, and free communication between its constituent settlements, and between the republic and its allies—both on Tal'Dorei and abroad. The incursion of the \textit{Chroma Conclave} has united the oft-feuding city-states of \textit{Syngorn} and \textit{Kraghammer} like never before, and all of Tal'Dorei now enjoys free trade and open borders between even its most remote enclaves. Thanks to the actions of Vox Machina, the previously solitary city-state of \textit{Whitestone} has joined the republic, to the benefit of both.

The \textit{Clasp} was once considered a serious threat to the council's rule, but times have changed. These days, \textit{Clasp} Spirelings in \textit{Emon} frequently make public appearances before the council on behalf of their own organization and their often unscrupulous business partners. This is to say nothing of the \textit{Clasp's} deep involvement in the Master of Information's spy ring.

The \textit{Tal'Dorei Council} has comfortable trade and diplomatic agreements between the realm of \textit{Othanzia} in \textit{Issylra} and the loosely federated cities of the Clovis Concord in \textit{Wildemount}. The recent conclusion of a bloody war between the Dwendalian Empire and the Kryn Dynasty in \textit{Wildemount} has given the \textit{Tal'Dorei Council} cause to send ambassadors to both nations to establish embassies, in hopes of creating an historic arrangement with all of the sovereign nations of that oft-maligned continent. On the continent of \textit{Marquet} to the southwest, the great city of \textit{Ank'Harel} has been friendly and open when it comes to trade and commerce, but has shown no interest in any political relationships beyond that.

\subsubsection{Figures of Interest}
These figures are major members of the \textit{Tal'Dorei Council} that you can use as NPCs in your game. You should also take inspiration from these figures to make up new NPCs for your own campaign.

\subparagraph{Guardian Tofor Brotoras}
\textit{\textit{Dragonblood} Master of Defense}

Stern, honorable, and quick to anger, Tofor is proud to serve on the council. After losing many friends (and one of her own arms) in the struggles against the \textit{Chroma Conclave}, she took the vulnerability of the city to heart, swearing never to allow such a thing to happen again. She is a beloved hero of the war against the \textit{Chroma Conclave} and has served on the council since before the death of the Last Sovereign. Most folk in \textit{Emon} suspect she'll win reelection until the day she dies.

\subparagraph{Seeker Odessa Tal'Dorei}
\textit{Human Master of Information}

When the eldest daughter of the Last Sovereign campaigned for the Master of Information's seat, its current holder was Seeker Assum Emring, a loyal servant of her father's. He bowed out of the race, and his blithe quote made headlines in \textit{Emon}: "I taught her everything she knows." She won in a landslide. Odessa wields the dual blades of Assum's spy network and the mercenary agents of the \textit{Clasp} with practiced finesse. Nevertheless, her name and the legacy of her father's martyrdom weigh heavily upon her. As Seeker, Odessa is calm, methodical, and laconic, but she carries a dignified sorrow.

\subparagraph{Arcanist Allura Vysoren}
\textit{Human Master of Arcana}

\textit{Allura Vysoren} has held the position of Master of Arcana since before the \textit{Chroma Conclave's} attack, and is held in the esteem of Tal'Dorei's people—mages and common folk alike. As a member of the \textit{Arcana Pansophical}, she is also increasingly worried as more and more of Tal'Dorei's politicians fall in debt to the \textit{League of Miracles}. After decades of investigating supernatural abnormalities herself, age has finally forced her to limit her adventuresome excursions. She now is one of the council's foremost employers of adventurers.

\subparagraph{Arbiter Brom Goldhand}
\textit{Human Master of Law}

A stone-faced, gray-haired, flame-scarred cleric of the The Knowing Mentor, Brom is as mysterious as he is proficient with law. One thing is clear: he is dedicated to keeping the realm just and orderly, and there are few who would dare challenge his judgement.

\subparagraph{Hearthmaker Edelbern Cleareyes}
\textit{Goliath Master of Development}

As a goliath of the \textit{Rivermaw nomads}, no one expected Edelbern to have an interest in civic architecture. Yet, after growing up in \textit{Westruun} while it was occupied by the Herd of Storms some twenty years ago, he found that he preferred the beauty of stone towers to towering pines. He also possessed a remarkable gift of mathematical acuity, and used his skills to set \textit{Westruun}—and eventually all of Tal'Dorei—on the path of reconstruction. Edelbern is also idealistic to a fault, and has, in his dedication to restoration, been manipulated into handing numerous lucrative contracts to the \textit{League of Miracles}, allowing them to amass greater and greater power.

\subparagraph{Coinmistress Vex'ahlia de Rolo}
\textit{Half-elf Master of Commerce}

As a member of Vox Machina and hero of the realm, \textbf{Vex'ahlia} could help Tal'Dorei in any number of ways. She already held a ceremonial title in \textit{Whitestone} as \textit{Grand Mistress of the Grey Hunt}. Yet when she and her husband pushed for \textit{Whitestone} to be formally recognized as a constituent city of the Republic of Tal'Dorei, all who knew her name (and her legendary financial savvy) clamored for her to be elected to the council as Coinmistress. She took to the duty with aplomb, and fights hard to ensure that \textit{Emon} never gets the bad end of a deal when gold is involved.

\subparagraph{Ambassador Syldor Vessa}
\textit{Elf ambassador from \textit{Syngorn}}

\textbf{Vex'ahlia} elven father has served as his city's ambassador to \textit{Emon} for well over a century, and requested to be appointed a seated ambassador as soon as the office was created. Though he is otherwise a staunch traditionalist, Syldor was inspired by his brave children to move past his old racial prejudices and drive his homeland to be less bigoted and xenophobic—learning from his own mistakes as a father to two half-elven children. His relationship with his elder daughter is still strained, and it is a newsworthy occasion when they vote in accord.

\subparagraph{Ambassador Deirdrik Greyspine}
\textit{Dwarf ambassador from \textit{Kraghammer}}

The youngest son of the ruling House Greyspine of \textit{Kraghammer}, Deirdrik has little hope of being elected Ironkeeper. Instead, this ambitious young dwarf—a scant seventy years old—left his home to win political renown in the national arena. He is a romantic idealist at heart and always votes for bold new policies, but doesn't always think through the consequences of his actions.

\subparagraph{Ambassador Keyleth of the Ashari}
\textit{Half-elf ambassador from \textit{Zephrah}}

A legendary archdruid and member of Vox Machina, \textbf{Keyleth, Voice of the Tempest} is both leader of the Air Ashari and ambassador for all Tal'Dorei's \textit{Ashari} peoples. She despises politics, and is focused entirely on improving the wellbeing of as many people as possible—as well as involving \textit{the Ashari} on the national stage. She spends most of her time in \textit{Zephrah}, with her people—including her father, Korrin, and her newly returned mother, Vilya—but she can and will arrive in \textit{Emon} at the drop of a hat, if \textbf{Vex'ahlia} asks it of her.

\subparagraph{Delegate Kel'jaia Uloeh}
\textit{Gnome delegate from \textit{Syngorn}}

Kel'jaia is a confident gnome with incredible strength and resilience—both physical and emotional. She is thought by all to be the council's first choice for Master of War, should the time ever call for it. She was born in the \textit{Rifenmist Jungle} and grew up amidst mercenaries who fought against the \textit{Iron Authority}. After her commander was killed in battle, she traveled north to \textit{Syngorn} to seek aid in her fight. She dreams of forging bonds of unity between her homeland and Tal'Dorei.

\subsection{Arcana Pansophical}
All of Tal'Dorei's greatest mages fought at the \textit{Battle of the Umbra Hills}. They saw the incursion of demons and devils firsthand, and were rightly terrified that the spawn of the \textit{Betrayer Gods} could ever return to the world. It was clear that the \textit{Divergence} had not wholly prevented supernatural evil from infringing upon the world, and that the mages of Exandria must prevent future "surprises" of this nature. The mistakes of the \textit{Age of Arcanum} must never be repeated. In light of this, \textit{Yurek Windkeeper}, high enchanter of \textit{Syngorn} and a well-respected arcanist, reached out to the realm's most trusted and capable mages, and together they formed a private society called the \textit{Arcana Pansophical}.

\subsubsection{Goals}
The chief priority of this fellowship of powerful arcane specialists is to maintain the Truth of Magic, a strict code enacted to never again allow the rampant misuse of magic seen in the \textit{Age of Arcanum}. General members are referred to as Guides, while the higher circle of members who deal in judgment and oversight are called Makers. At least one member is established at each place of arcane learning.

Membership is largely kept secret unless out of necessity, though the existence of the order is relatively common knowledge. Anyone found to be practicing magic is logged and routinely checked up on, as misuse of magic is grounds for arrest, judgment by the Makers, and punishment outside of the local law. Even so, the reach and sight of the \textit{Arcana Pansophical} are not without limits, and many dark and terrible dealings manage to escape their notice.

\begin{DndSidebar}[float=!b]{The Truth of Magic}
The \textit{Pansophical's} myriad rules on the use and misuse of magic are innumerable—detailed at length in centuries of amendments made to the full text of the Truth of Magic. Though the Truth's bizarre intricacies are left up to the GM, the Truth's preamble establishes three unbreakable edicts:



\DndQuote%
{}
{''In the Eternal interest of the Preservation of Harmony within the realm of Exandria, we the Arcana Pansophical establish the Truth of Magic, and three Edicts from which none shall stray. Those Arcanists found in gross violation of the Truth will be punished by the full extent of this Order—outside of any local law.''}
{''Though the school of Necromancy shall be restricted to none, in the interest of magical study and Understanding, the animation of the dead is a Violation of the Truth.''}
{''The Arcane is a tool to be wielded for the good of the People. To use its Power in the pursuit of wanton Destruction or unjust Murder is a Violation of the Truth.''}
{''In addition, though the Jurisdiction of the Pansophical supersedes the power of local laws, a mage who willfully breaks the laws of the land is in Violation of the Truth.''}
{}


\end{DndSidebar}

\subsubsection{Relationships}
The \textit{Arcana Pansophical} is respected, if somewhat feared, throughout the Republic of Tal'Dorei. \textit{Syngorn} and other elf-founded settlements see the \textit{Pansophical} as allies, as the roots of the society began within \textit{Syngorn} itself. The people of \textit{Kraghammer} are suspicious of the \textit{Pansophical's} legalistic Truth of Magic, and they fear that its members hold more loyalty to the \textit{Pansophical} than to their homeland.

Recently, a group of profit-driven arcanists splintered from the \textit{Arcana Pansophical} to form the self-styled \textit{League of Miracles}. Their involvement in the for-profit reconstruction of Tal'Dorei has made them popular among mages who see the \textit{Pansophical} as a stodgy, elitist institution. The \textit{Pansophical}, however, sees the \textit{League of Miracles} as a distasteful upstart that will soon buckle under its own ambition.

\subsubsection{Figures of Interest}
These figures are major members of the \textit{Arcana Pansophical} that you can use as NPCs in your game. You should also take inspiration from these figures to make up your own NPCs.

\subparagraph{Enchanter Yurek Windkeeper}
\textit{Elf wizard (enchanter) and Maker}

Founder of the \textit{Arcana Pansophical}, the ancient archmage Yurek generally stays within \textit{Syngorn}, acting as the nexus of the order and a consultant to the \textit{Wardens of Syngorn}.

\subparagraph{Alchemist Oz Gruude}
\textit{Halfling alchemist and Maker}

Gruude is an ornery old halfling who is more interested in traveling the world to study plants and minerals with latent magical properties than doing alchemical lab work. He is rarely seen in the halls of the \textit{Pansophical}.

\subparagraph{Seer Gloria Ios}
\textit{Human wizard (diviner) and Maker}

Gloria is one of the great diviners of the age, and her duties have sent her to Vasselheim, where she acts as a foreign outreach for the \textit{Pansophical}.

\subparagraph{Magus Seanor Wiles}
\textit{Human wizard (enchanter), Maker, and tutor}

A humble-looking traveler, Wiles delights in finding fresh and emerging talent and helping his new students refine their arcane gifts.

\subparagraph{Planerider Ryn}
\textit{Tiefling wizard (conjurer), Maker, and planar studies specialist}

This former pupil of \textit{Yurek Windkeeper} has grown into a master of planar travel. She likewise has taken an interest in the flow of magic across the world by means of ley lines. She's constantly on the move, traveling from plane to plane to conduct her research.

\subparagraph{Tinkerer Quash Dentdruggle}
\textit{Gnome construct specialist and Guide}

Obsessed with the development of autonomous and living arcane constructs, Quash is usually holed up in his laboratory within \textit{Kraghammer}.

\subparagraph{Arcanist Allura Vysoren}
\textit{Human wizard (abjurer) and Guide}

As a member of the \textit{Tal'Dorei Council} and a retired adventurer, \textit{Allura} has helped maintain the political significance of the \textit{Arcana Pansophical} within \textit{Emon}. Even among her colleagues in the \textit{Pansophical}, her unpretentious strength commands great respect.

\subparagraph{Elementalist Drake Thunderbrand}
\textit{Dwarf wizard (evoker) and Guide}

Driven to bind the physical elements of the chaotic planes to his will, Drake is a prodigy of new incantations in the field. He and \textit{Allura Vysoren} are lifelong friends, and together they're renowned for achieving the impossible.

\subparagraph{Realmseer Eskil Ryndarien}
\textit{Human archmage and Maker emeritus}

Though he was a Maker for most of his life, Realmseer Ryndarien grew frustrated with the politics and limitations of the \textit{Pansophical}, eventually retiring to \textit{Westruun} to pursue his own interests. These days, he has found a way to rejuvenate his aged body and is now renowned as a handsome and aloof man about town—much to the annoyance of all who formerly knew him.

\subsection{The Ashari}
The Material Plane is surrounded on all sides by the churning chaos of the elements. This elemental power is invisible to most, but in some places, the veil between this world and the Elemental Planes grows thin—and sometimes is torn entirely asunder. These places where the Elemental Planes bleed into the world are called rifts, and they first appeared on a long-ago Celestial Solstice when the planes aligned. That day, elemental rifts of untold proportions burst open across Exandria. Wildfires, tsunamis, hurricanes, and earthquakes ravaged the lands, and the people of the world struggled to combat the elemental beings that poured forth from the rifts. The tide of rampaging elementals was eventually stemmed, but the world was gripped by fear—what would happen on the next solstice?

In the wake of the planar disturbance, a society of druids called the Ashari made a covenant to divide and form four separate orders to seal and watch over the sources of these unleashed energies. Most of the lesser rifts that opened in that time have been healed, but one rift to each of the four Elemental Planes remains in Exandria. The four orders of the Ashari settled near these rifts, watching over them and soothing them when they grow volatile. Of these orders, two are in Tal'Dorei, where they watch over the rift of earth and the rift of air. A dozen generations later, the Ashari still hold watch over these tumultuous locations.

\subsubsection{Ashari Culture}
Though the Ashari began as a single group of elementalists, their individual orders have each developed their own ways of living over the centuries.

The people of Pyrah value reactivity and coolness under pressure, for panic is death during an eruption. The Ashari of Vesrah have a close relationship with the marine life that surrounds their rift, and their actions always prioritize the health of Exandria's waters. \textit{Zephrah's} culture values the freedom of the open sky, and balances their sense of duty with a free-spirited sense of play. Finally, the society of \textit{Terrah} is rooted in planning for the future, and fixing what is broken with patience and humility.

\subsubsection{Goals}
Each Ashari order guards their respective rift between one of the elemental planes and Exandria. Here, they strive to understand, master, and ultimately subdue the ever-escaping elemental forces. They attempt to bring balance to the surrounding lands and heal the rift over time, though recurring cosmic events tend to reopen these wounds.

\textit{The Ashari} of Pyrah oversee the Fire Elemental Rift in the valley caldera of the Sunderpeak Mountain range in \textit{Othanzia}, upon the continent of \textit{Issylra}. \textit{The Ashari} of Vesrah watch over the Water Elemental Rift on the Islands of Anamn, within the Ozmit Sea. On the southwest edge of the \textit{Cliffkeep Mountains}, \textit{the Ashari} of \textit{Terrah} stand vigilant over the Earth Elemental Rift, while high atop the windy mountaintops of the \textit{Summit Peaks}, \textit{the Ashari} of \textit{Zephrah} keep the Air Elemental Rift. The Air Ashari also watch over the \textit{Frostweald}, where a being of elemental ice once broke through to cause untold destruction.

Due to their isolation, \textit{the Ashari} go to great lengths to maintain contact with their kin. Each order is led by a single elemental master, a title that is passed on by blood or honor for generations. In every new generation that comes of age, the young leader-to-be must embark upon an Aramente. This "noble odyssey" sends them across the world to meet with the other Ashari leaders and train until they have earned both the respect and an understanding of their hosts.

Upon returning, they ascend to the head of their order, while the elder retires. This journey often takes years, and should one such Ashari not return within a reasonable period of time, a next-of-kin or a suitable elemental adept is chosen to take up the Aramente.

\subsubsection{Relationships}
\textit{The Ashari} maintain peaceful, if sometimes strained, relationships between the orders, though inner conflict occasionally arises within them. Beyond this, much of life within \textit{the Ashari} is hidden from outsiders' prying eyes. In recent years, however, \textbf{Keyleth, Voice of the Tempest} of the \textit{Zephrah} has taken steps to unite \textit{the Ashari} with the rest of the people of Tal'Dorei. She has gone so far as to become an ambassador for \textit{the Ashari} people on the \textit{Tal'Dorei Council}. Part of this work has been expanding the crisis orbs system, which \textit{the Ashari} used to contact their distant kin in times of peril, to connect cities across Tal'Dorei in the event of another crisis on the level of the \textit{Chroma Conclave}.

\textit{The Ashari} of \textit{Terrah} nurture a seed of disrespect and mistrust for \textit{Kraghammer}, whose miners constantly disregard their warnings against reckless mining. This arrogance caused countless deaths in the most recent eruption of the \textit{Terrah} Rift—and led to the collapse of an entire mountain.

Each Ashari culture has its own favored deities. Many Ashari of Pyrah venerate the The Dawnfather, for his flames give life to Exandria. Many in Vesrah worship the The Moonweaver, for her work controls the tides. In \textit{Zephrah}, the faith of the The Wildmother is predominant, for her freedom is like that of the wind. And the people of \textit{Terrah} hold the The All-Hammer in great respect for his association with creation and the earth. There are many Ashari, as well, who prefer to revere the elements themselves, forgoing the worship of gods.

\subsubsection{Figures of Interest}
These figures are major members of \textit{the Ashari} that you can use as NPCs in your game. You should also take inspiration from these figures to make up new NPCs for your own campaign.

\subparagraph{Pa'tice, Heart of the Mountain}
\textit{Human leader of \textit{Terrah}}

Elderly leader of the Earth Ashari, Pa'tice still shows vigor and strength to resist his age, having led his people through many very tough challenges. His stubbornness is legendary, and it's said even death has been forced to wait until Pa'tice sees it his proper time. Many believe that his druidic training will keep him alive until he himself chooses to let go of this mortal coil.

\subparagraph{Korrin of Zephrah}
\textit{Half-elf elder of \textit{Zephrah}}

His wife Vilya was lost while she undertook the Aramente, and when the previous Voice of the Tempest died of illness, Korrin was called upon to act as Voice until his daughter \textbf{Keyleth, Voice of the Tempest} came of age and completed her Aramente. Korrin was shocked when Vilya suddenly returned, and now spends his days in disbelieving bliss, reunited with his long-lost wife. He is gentle with those who show honor, and merciless to those who threaten his people.

\subparagraph{Vilya, Who Escaped Death}
\textit{Half-elf elder of \textit{Zephrah}}

Vilya disappeared on her Aramente nearly thirty years ago, and was presumed dead. Her daughter, \textbf{Keyleth, Voice of the Tempest}, took up her own Aramente ten years after Vilya left, and discovered from the Water Ashari of Vesrah that Vilya was killed by a \textbf{kraken}, and only her leg was able to be recovered. Less than a year ago, Vilya suddenly strode out of a tree in \textit{Zephrah}, her leg replaced with a druidic prosthetic of her own creation, shocking all—though none were more overwhelmed with emotion than her daughter and husband. Though Vilya never earned the title of Voice of the Tempest, she now aids her daughter in leading her people without the faintest trace of jealousy.

\subparagraph{Keyleth, Voice of the Tempest}
\textit{Half-elf leader of \textit{Zephrah}}

Hero of Tal'Dorei and member of Vox Machina, \textbf{Keyleth, Voice of the Tempest} completed her Aramente and took leadership among her people following the destruction of the \textit{Chroma Conclave}. She seems to have not aged a day since Vox Machina's victory over The Whispered One. Since her beloved's death, she has kept busy to fill the void in her heart. From expanding the crisis orb network across the realm and building beautiful windmills in the \textit{Summit Peaks} with \textbf{Percival de Rolo} aid to training Air Ashari heroes at home, she has never allowed herself a free moment. She has rarely left \textit{Zephrah} in the months since her mother's emotional return.

\subsection{Brawler's League}
The brutal blood sports of the Brawler's League began, unsurprisingly, during the vicious \textit{Drassig} Dynasty. In that cruel era, many of \textit{Emon's} affluent and indolent social elite used their wealth to engage in entertainment of a more "visceral" variety. The kingdom's poor were already fighting each other for scraps; why not pay them to literally fight to the death for their amusement?

Some strong, charismatic ruffians used this opportunity to gain social power for themselves, and chief among them was a rough-and-tumble entrepreneur from \textit{Kraghammer} named Bruel Sunderchin. Reaching out to his network of lowlifes and cutthroats, Sunderchin became ringleader of an annual gladiatorial tournament that caters expressly to Exandria's wealthiest elite.

Those who have become invested members of the league's events are anonymously delivered a golden coin bearing their symbol in the weeks leading up to an event. These coins are required for entry into the events, and those who wish to join, whether as a patron or participant, require a coin of their own.

\subsubsection{Goals}
The \textit{Brawler's League} is always in search of hidden and well-guarded locations around Tal'Dorei (and sometimes other realms) to establish their semi-annual event. A combatant or team is required to have a known and cleared sponsor or manager to enter the event. There are two tiers of bouts; one tier has teams of six battling via single elimination until a victor is declared. The other features single combatants fighting one-on-one, single-elimination, until a victor is declared. Killing your opponents isn't encouraged, but isn't explicitly discouraged, either. Bets are placed in person at the event, or remotely (via arcane messages and \textit{sending stones}) from all around Exandria.

\subsubsection{Relationships}
The \textit{Brawler's League} is extremely illegal in all major cities across Tal'Dorei. Should law-keepers of any nation discover the location of the semiannual event, the venue would find itself swiftly strangled by the long arm of the law. As such, magical protections and sentinels are placed throughout the event to prevent any unapproved entry, or to wipe the memory of those that do infiltrate.

A common rumor states that numerous political figures who publicly decry the event secretly partake in the gambling, and that the location of the event is concealed not only by magic, but also impenetrable layers of bureaucracy. \textit{Kraghammer} has hosted the event a number of times, with most local guards and great houses turning a blind eye.

The \textit{Brawler's League} frequently hires local \textit{Clasp} agents (or \textit{Myriad} operatives, in a pinch) to spread word of their next event through the criminal underground. If all goes well, they get hired at the event itself to provide protection, run interference, and make money on the side by providing loans and taking bets.

\subsubsection{Figures of Interest}
These figures are major members of the \textit{Brawler's League} that you can use as NPCs in your game. You should also take inspiration from these figures to make up new NPCs for your own campaign.

\subparagraph{Oben "Broadstaff" Sunderchin}
\textit{Dwarf organizer and head of \textit{Brawler's League}}

Descendant of Bruel and inheritor of the league's top position, Oben secretly goes by the name Broadstaff and doesn't directly appear at events, as he avoids public connection with the league.

\subparagraph{Kradin Grimthorne}
Dwarf fight manager

A previously disgraced dwarf, exiled from \textit{Kraghammer} for murder, Grimthorne is the foremost brawler talent scout in Tal'Dorei.

\subsection{Chamber of Whitestone}
Generations ago, a team of intrepid explorers and fortune-seekers sailed from Port Damali in \textit{Wildemount} to northeastern Tal'Dorei through the \textit{Shearing Channel}, but wrecked upon the \textit{Alabaster Sierras'} rocky cliffs. The expedition was led by the de Rolo family, and these brave \textit{humans} led the survivors to safety after the wreck. They took refuge in a valley within the \textit{Alabaster Sierras}, and there they discovered the holy \textit{Sun Tree}, a radiant tree blessed by the The Dawnfather. A community grew around the \textit{Sun Tree}, and the respected and beloved de Rolos were asked by their shipwrecked allies to be the leaders of this new settlement. Pressed against the chalky cliffs of the \textit{Alabaster Sierras}, their new homeland was dubbed \textit{Whitestone}.

The sovereign city-state of \textit{Whitestone} enjoyed peace for many generations, secluded from the rest of Tal'Dorei and \textit{Wildemount}. This serenity was destroyed when Lord Sylas and Lady Delilah Briarwood arrived in \textit{Whitestone} and murdered most of the de Rolo family. Lady Delilah, a necromancer of The Whispered One and a disgraced member of the Cerberus Assembly in \textit{Wildemount}, and her husband, a power-hungry \textbf{vampire}, transformed the quiet city-state into a realm of death. They attempted to take control of all of Tal'Dorei by puppeteering the Last Sovereign, Uriel Tal'Dorei II, but were thwarted by Vox Machina.

Thirty years after \textit{Whitestone's} fall, the \textit{Briarwoods} and their sepulchral master are defeated, and \textit{Whitestone} has returned to the blessed peace of its past. Nevertheless, much has changed. The surviving de Rolo siblings, \textbf{Percival de Rolo} and Cassandra, reclaimed the city with the aid of Vox Machina. With the tyrants dethroned, \textbf{Percival de Rolo} established a council titled the Chamber of \textit{Whitestone} to rule in his frequent absence, and the city now sends a delegation to the \textit{Tal'Dorei Council} to help safeguard the republic against evil.

\subsubsection{Goals}
Under the watch of Lady Cassandra de Rolo, the chamber handles and surveys all aspects of political life in \textit{Whitestone}. Its members are appointed by the head of the chamber, who is likely to be a scion of the de Rolo family for generations to come. The head can also add or remove members from the council. This system is new and untested by a ruler less noble than \textbf{Percival de Rolo} or Cassandra. In many ways, this system of government is another one of \textbf{Percival de Rolo} tinkering projects.

Members of the Chamber are granted distinct areas of responsibility to the realm.

\textbf{The Guardian of Woven Stone} is the title given to the central figure on the council, responsible for higher civil justice, community unification, and oversight of internal government.

\textbf{The Steward of Sunblessed Gifts} is responsible for local farming and goods, overseeing distribution between citizens and exportation.

\textbf{The Curator of Fortune's Bounty} is responsible for all major commerce and economic business, both locally and with allied and foreign entities.

\textbf{The Keeper of Divine Virtue} is responsible for religious organization and spiritual guidance regarding multiple faiths within the region, both as a religious leader and a protector of practices.

\textbf{The Commissioner of Wardship} is responsible for the enforcement of local law, the organization of the Paleguard, acting as judge on smaller matters of civil justice, and defense of the city during times of war.

\textbf{The Architect of Enlightened Progress} is responsible for the construction and infrastructure of the city and castle.

\textbf{The Grand Mistress of the Grey Hunt} is responsible for diplomatic relations and for rooting out dangers beyond the city's walls.

\subsubsection{Relationships}
Despite its benefits, \textit{Whitestone's} previous solitary existence left it without allies in its time of need. The Chamber of \textit{Whitestone} has sent a delegation to the \textit{Tal'Dorei Council} for the betterment of both. These days, trade caravans are a common sight along the \textit{Silvercut Roadway}, transporting goods from \textit{Whitestone} to \textit{Westruun}, \textit{Kymal}, and ultimately \textit{Emon}. The most expensive goods travel via skyships, which land at \textit{Whitestone's} brand-new skyport.

\subsubsection{Figures of Interest}
These figures are major members of the \textit{Chamber of Whitestone} that you can use as NPCs in your game. You should also take inspiration from these figures to make up new NPCs for your own campaign.

\subparagraph{Lady Cassandra Johanna de Rolo, Guardian of Woven Stone}
\textit{Human head of the \textit{Chamber of Whitestone}}

Cassandra is \textbf{Percival de Rolo} younger sister, thought dead until he discovered her in the thrall of the \textit{Briarwoods}. She now manages \textit{Whitestone's} most pressing matters, with daily briefings from other members of the chamber. She is beloved throughout the city for her resilience and her decisive leadership.

\subparagraph{Lord Percival Fredrickstein von Musel Klossowski de Rolo III, Architect of Enlightened Progress}
\textit{Human member of the \textit{Chamber of Whitestone}}

\textbf{Percival de Rolo} is Cassandra's older brother, thought dead until Vox Machina discovered him in a prison cell. His life was filled with countless joys and regrets. He has taken well to retirement from adventuring, and is renowned across Tal'Dorei for his creations.

\subparagraph{Lady Vex'ahlia de Rolo, Baroness of the Third House of Whitestone and Grand Mistress of the Grey Hunt}
\textit{Half-Elf Member of the \textit{Chamber of Whitestone}}

\textbf{Vex'ahlia} is the Grand Mistress of the Grey Hunt, a title she has held since before she married \textbf{Percival de Rolo}. Over the past twenty-four years, she has become a talented and respected leader of the \textit{Grey Hunters}. And though the hunters use guns inspired by her husband's own designs, she still prefers to wield her trusty bow.

\subparagraph{Grog Strongjaw, Grand Poobah de Doink, In Charge of All This-and-That}
\textit{Goliath friend}

The noble, if simple, warrior of Vox Machina, no one really knows if \textbf{Grog Strongjaw} title is a real title at all, or if \textbf{Percival de Rolo} was just joking when he gave it to him. \textbf{Grog Strongjaw} likes it, and that's good enough for most people.

\subsection{Claret Orders}
\begin{DndReadAloud}{}
\DndQuote%
{}
{''As blood flows through and invigorates the body, so faith flows through and invigorates the soul. Ever remember this: your blood and your faith are one and the same. Spill not your blood nor that of another heedlessly, but do not hesitate to spill either when the cause is just.''}
{From The Crimson Canon}
\end{DndReadAloud}

Hundreds of years ago, a fell power exerted its dark will over \textit{Wildemount}, corrupting its outlying townships and spreading chaos among its people. Undeath washed over gravesites like a plague, shadowed beasts stalked the midnight woods, and the influence of fiends befouled even \textit{Wildemount's} purest souls.

In the eleventh hour, a priest of the The Matron of Ravens named Trence Orman prayed for a way to protect his flock. The The Matron of Ravens inspiration came in the form of long-hidden knowledge: the secrets to blood magic. Taking his gifts, Trence trained his most trusted warriors in these techniques, giving a portion of their humanity in exchange for the power to defend their people. This marked the origin of the \textit{Claret Orders}, and the first \textit{Blood Clerics}, \textit{Blood Mages}, and Blood Hunters.

The \textit{Claret Orders} work in relative secrecy, as the nature of their abilities is largely misunderstood, which often leads to persecution. Their numbers are small, as the price is great and the \textit{Orders} only accept those who are strong of heart and have little to lose. Small sects of the \textit{Orders} have begun work far enough West to wander the lands of Tal'Dorei in recent years, keeping out of the public eye and remaining vigilant to the signs of familiar darkness undermining the region.

Despite the \textit{Claret Orders'} best attempts at secrecy, the esoteric art of blood magic has escaped their clutches. Though rare, rogue blood mages and blood hunters have spread across both \textit{Wildemount} and Tal'Dorei, unbound by the tenets and strictures of the \textit{Claret Orders}.

\subsubsection{Goals}
The Creed of the \textit{Claret Orders} asks that its followers commit their lives to the hunt of entities and creatures that threaten the sanctity of life and joy. While the sacrifice is great, the reward is the continued existence of purity and good in the world. The \textit{Orders} strive to protect those who cannot protect themselves from the shadows, taking no credit for their services beyond the means to live and travel. Some have strayed, focusing on building a fortune for themselves, but should word of such intent make it back to the heads of the \textit{Orders}, and their actions begin to threaten the innocents they've sworn to protect, they may find themselves among the hunted.

\subsubsection{Relationships}
The \textit{Claret Orders} have no alliance with—nor allegiance to—any standing government or group. Solitary in action and discreet in display, they've not drawn the attention of many factions in Tal'Dorei as of yet. However, it is only a matter of time before curious eyes begin to seek answers to questions about these dark guardians and their true nature.

\subsubsection{Figures of Interest}
These figures are major members of the \textit{Claret Orders} that you can use as NPCs in your game. You should also take inspiration from these figures to make up new NPCs for your own campaign.

\subparagraph{Claret Director Alasterre de Vitrevos}
\textit{Human Blood Domain cleric of Ravens}

Calm and sanguine to a fault, \textit{Alasterre de Vitrevos} is the youngest-ever director of Tal'Dorei's \textit{Claret Order}. He sees his position as one he deserves without question, and is willing to use his command more boldly than his mentor, director emeritus Teresa Dulamar. He is frequently seen in the field, working with various adventuring parties and his own blood hunters to wipe out monsters wherever they emerge. The combination of his relaxed personality and his aggressive strategy may seem a contradiction at first glance, but those who know him see the truth of the immense self-assurance that lets him make hard decisions without hesitation. This quality is a strength when cool-headed leadership is needed in the thick of battle, but it can make him challenging to work with as an equal.

\subparagraph{Director Emeritus Teresa Dulamar}
\textit{Dwarf Blood Domain cleric of Ravens}

The former head of Tal'Dorei's \textit{Claret Orders} retired ten years ago after suffering a traumatic death in combat. Her soul accepted resurrection, but she took her defeat as a sign to let younger, fresher blood take her place. She remains a part of the \textit{Claret Orders} to advise the current director on matters he is too young to have experienced himself. Though she displays a chilly demeanor, she simply keeps her compassionate side hidden from the untrusted riffraff.

\subparagraph{Hunter Jorick LaMensh}
\textit{Human blood hunter}

Once among the greatest of all blood hunters, Jorick is now Director Emeritus Teresa Dulamar's personal bodyguard. He is still a capable member of the Order of the Profane Soul, and urges immediate, decisive action to destroy demons and devils whenever their occult activity is discovered.

\subsection{The Clasp}
The \textit{Clasp} is a well-established secret organization consisting of spies, thieves, fences, assassins, saboteurs, and smugglers. It was built upon the ruins of the Winksman Thieves' Guild, an organization that flourished during the reign of \textit{Drassig}, but crumbled under the rule of \textit{Zan Tal'Dorei}. Based largely in \textit{Emon}—and under it, via a network of subterranean tunnels—the \textit{Clasp} also runs smaller outfits in \textit{Westruun}, \textit{Stilben}, and \textit{Kymal}. Members are branded with the guild's symbol on their mid-back, signifying allegiance and commitment. The \textit{Clasp} operates like a business-minded family or mafia, celebrating loyalty and an ability to keep one's word. A member pays their debts, keeps their promises, and always seeks opportunity for the betterment of the \textit{Clasp}.

Leaders of a \textit{Clasp} sect bear the title of Spireling, and they oversee assigned elements of the organization's functions. Spirelings hold equal power within a sect, and there is always an odd number acting within that location to break any disagreements over business ventures and actions.

\subsubsection{Goals}
While most common folk assume that the \textit{The Clasp} is merely interested in profit at any cost (and that is a merit of the \textit{The Clasp}), they also wish to preserve the function and prosperity of civilization. It is commonly held that a thriving city is a profitable one, and an enemy of civilization is a foe of the \textit{The Clasp}. They largely consider themselves a necessary dark underbelly, willing to do the under-the-table deeds that more wealthy and powerful citizens cannot commit with their own hands. On the other hand, less scrupulous Spirelings also understand that discontent breeds profit, and have occasionally kindled the flames of unrest to keep things tense and profitable.

Were it not for the branch of the \textit{The Clasp} in \textit{Emon}, the \textit{Chroma Conclave's} attack would likely have destroyed the city's culture entirely. However, thanks to the \textit{Clasp's} role as a "shadow civilization," the people of \textit{Emon} were able to maintain their foothold at home, even under the watchful eye of \textit{Thordak, the Cinder King}.

\subsubsection{Relationships}
A generation ago, the \textit{Clasp} was considered by most to be an uncouth cesspool of cutthroats and criminals. These days, thousands of people know the \textit{Clasp} as the roguish heroes who saved their lives, or the lives of their parents, from the \textit{Chroma Conclave}. Nowhere is this truer than in \textit{Emon}, where its Spirelings make public appearances and even, on occasion, deal openly with the \textit{Tal'Dorei Council}. Nevertheless, \textit{Clasp} agents outside of \textit{Emon} keep their wits about them, for most local governments still view them as underhanded scoundrels—and not without reason. Most guards and town militias will arrest known \textit{Clasp} members on sight; however, most \textit{Clasp} members know well enough to have a friend on the force that will help them escape.

The \textit{Clasp} frowns upon dogmatic religion or cultish behavior, as it has conflicted with the organization's business in the past, and always ended messily.

Since the arrival of the \textit{Myriad} within Tal'Dorei from \textit{Wildemount}, the \textit{Clasp} has been on extremely high alert, seeking information on them and any means of banishing this new competition. Gang warfare has broken out sporadically between the \textit{Clasp} and the \textit{Myriad} in cities from \textit{Stilben} to \textit{Emon} over the past twenty years, and though the rival organizations are currently at peace, everyone knows the whole thing is just one poached client away from reigniting.

\subsubsection{Figures of Interest}
These figures are major members of the \textit{Clasp} that you can use as NPCs in your game. You should also take inspiration from these figures to make up new NPCs for your own campaign.

\subparagraph{Spireling Shenn}
\textit{Human Spireling of Secrets (\textit{Emon})}

A powerful illusionist, Shenn runs a far-reaching network of spies and informants. In his younger days, he was a canny negotiator who sought to spread the \textit{Clasp} beyond Tal'Dorei. The arrival of the \textit{Myriad} forced the \textit{Clasp} onto the defensive, and he will never forgive them for souring his ambitions.

\subparagraph{Spireling Warren}
\textit{Half-orc Spireling of Blades (\textit{Emon})}

Deadly duelist turned master assassin of the \textit{Clasp}, Warren handles all contracts that involve kidnapping, assassination, and intimidation. He hates public appearances.

\subparagraph{Spireling Zilloa}
\textit{Half-elf Spireling of Shadows (\textit{Emon})}

An accomplished burglar, Zilloa has risen to the head of all thieving, smuggling, and black-market dealings within the \textit{Clasp} of \textit{Emon}. She plays a dangerous game by enjoying folk-hero status as the public liaison between the \textit{Clasp} and the \textit{Tal'Dorei Council}.

\subparagraph{Spireling Gholesh}
\textit{Human Spireling of Shadows (\textit{Westruun})}

A crass and ambitious man, Gholesh's appearance is a mystery as his face is always hidden behind a leather mask that resembles an ogre. He has secretly thrown in his lot with the \textit{League of Miracles} and, with their backing, he aspires to take over the \textit{Clasp} and enjoy all its wealth and prestige.

\subparagraph{Spireling Melchior}
\textit{Tiefling Spireling of Blades (\textit{Westruun})}

This clean-faced, curly-haired tiefling with pale blue skin and stubby horns may have a soft, babyish visage, but it belies their powerful physique. They are equally skilled with a sword, axe, or bow, and made their name in underground brawling rings before joining the \textit{Clasp}.

\subparagraph{Spireling Fetch}
\textit{Halfling Spireling of Secrets (\textit{Westruun})}

Creepy and unkempt in appearance, Fetch's gaunt face only distracts from his cunning and skills of keen manipulation.

\subsection{Golden Grin}
In dark times, entertainment and good cheer are just as vital as resistance. During the despotic rule of King Warren \textit{Drassig} and his sons, a group of bards, dancers, and storytellers united to subvert his rule and inspire the people to rise against him. This band took the name "the \textit{Golden Grin}" and traveled across the realm to offer escapism and entertainment, all while planting the seeds of discontent, rebellion, and heroism.

Since the fall of \textit{Drassig}, the \textit{Golden Grin} has been ever-present, using song, comedy, and theatrics to inspire the people of Tal'Dorei to endure hardship and fight back against evil. These "Grinners" include among their ranks spiritual leaders, artists, musicians, innkeepers, and dozens of retired adventurers who still want to do good outside the public eye. Little is publicly known about the \textit{Golden Grin}, other than myths and rumors, and the Grinners gladly perpetuate these tall tales.

\subsubsection{Goals}
The \textit{Golden Grin} (or simply "the Grin") exists to enlighten and inspire, keeping ears and eyes out for the murmurs of tyranny and cruelty. The Grin's core philosophy is that each and every person is capable of great things alone, and even greater things together. Rather than taking direct action themselves, most Grinners use songs and tales of great heroes to galvanize people to action.

\subsubsection{Relationships}
The \textit{Golden Grin} claims no formal allies, nor do members reveal their secret affiliation to outsiders unless circumstances are dire indeed. They use their covert, grassroots power to empower local freedom fighters and activists. Some Grinners do occasionally find their way into government, granting them a higher vantage point and platform to disseminate the ideals behind their cause—but even these members of the Grin take a dim view of the slowly turning gears of bureaucracy.

Enemies of the \textit{Golden Grin} are those that oppress and abuse freedom of expression and living. The Grin likewise mistrusts people and institutions that control or regulate the populace too heavily, even if they seem to have the best of intentions.

\subsubsection{Figures of Interest}
These figures are major members of the \textit{Golden Grin} that you can use as NPCs in your game. You should also take inspiration from these figures to make up new NPCs for your own campaign.

\subparagraph{Theona Balmhand}
\textit{Gnome cleric of the Lawbearer}

A wandering healer hailing from \textit{Kymal}, Theona lives by donation and perpetually travels from place to place, seeking to aid who she can and to learn whatever the people whisper.

\subparagraph{Doctor Dranzel}
\textit{Half-orc traveling bard}

A jovial and lackadaisical violinist and frontman for Doctor Dranzel's Traveling Troupe, the renowned Doctor Dranzel exaggerates his stage persona for a number of reasons, not the least of which is to make him seem like a preening fool—the perfect cover for a master spy of the Grin.

\subsection{Houses of Kraghammer}
The fortified, subterranean city-state of \textit{Kraghammer} was carved out of the bedrock of the \textit{Cliffkeep Mountains} in the years following the \textit{Divergence}. It is the ancestral home of the native \textit{dwarves} of Tal'Dorei, and the five major bloodlines that founded the city in ages past still rule to this day. The ruling families are Houses Greyspine, Zuurthom, Bronzegrip, Thunderbrand, and Glorendar. These Great Houses elect an official to the office of Ironkeeper every ten years to keep the alliances between houses just and healthy.

Though the culture of \textit{Kraghammer} is highly traditional and as stoic as the stone it was carved from, it is not without unrest. It seems every dwarven generation has its own great social movement where disenfranchised workers or greedy, lesser noble houses attempt to unseat the ruling five. Yet for good or for ill, the Great Houses have time and time again shown that they can put aside their differences to preserve their mutual rulership.

\subsubsection{Goals}
Each ruling house has a specialty and holds responsibility over that domain within \textit{Kraghammer}. House Greyspine maintains the largest mine under \textit{Kraghammer}, as well as the Pools of Solace. House Zuurthom trains the chief architects of the city, and most masonry and building is organized through them. House Bronzegrip funds the smithing guilds and maintains the Bronzegrip Metalworks, \textit{Kraghammer's} largest blast furnace and ore refinery. House Thunderbrand prides itself on training \textit{Kraghammer's} foremost scholars and premiere arcanists. House Glorendar maintains the law and delivers punishment.

\subsubsection{Relationships}
Almost every dwarf living in \textit{Kraghammer} can remember a time when, "Kraghammer for the dwarves!" was a common cheer at the pubs. And when they said, "for the dwarves," they meant just the \textit{dwarves} born in \textit{Kraghammer}. Times have changed. The prevailing opinion is that a unified Tal'Dorei is best for all. Those \textit{dwarves} whose traditionalism is just a veil for bigotry have found that they aren't welcome in \textit{Kraghammer}—they've subsequently retreated to make their own isolated enclaves in the \textit{Cliffkeeps}.

The \textit{Houses of Kraghammer} facilitate a healthy trade of metals, weapons, and armor with \textit{Emon} and \textit{Westruun}. These bonds of commerce and communication have only grown stronger since an ambassador from \textit{Kraghammer} has taken a seat on the \textit{Tal'Dorei Council}.

Many members of House Thunderbrand also have ties to the \textit{Arcana Pansophical}. However, lesser clans with aspirations of supplanting the renowned Thunderbrands have thrown their lot in with the \textit{League of Miracles}, hoping to become the new masters of the arcane. They have been wooed by to this upstart faction with promises of disrupting the old ways of commerce and prestige, establishing something new.

\begin{DndReadAloud}{}
\DndQuote%
{}
{''Five the fingers that grasp hammer or blade,''}
{''Five Houses of Kraghammer rightly made:''}
{''Greyspine delves 'neath dwarven halls,''}
{''Zuurthom raises stony walls,''}
{''Bronzegrip masters of metal wrought,''}
{''Thunderbrand uncovers secrets sought.''}
{''And Glorenthal to judgements hear—''}
{''In Kraghammer, dwarf-folk need never fear.''}
{}
\end{DndReadAloud}

\subsubsection{Figures of Interest}
These figures are major members of the \textit{Houses of Kraghammer} that you can use as NPCs in your game. You should also take inspiration from these figures to make up new NPCs for your own campaign.

\subparagraph{Ironkeeper Palluda Zuurthom}
\textit{Dwarf Ironkeeper of House Zuurthom}

Even after weathering the \textit{Chroma Conclave}, uniting \textit{Kraghammer's} squabbling houses in support of the \textit{Tal'Dorei Council} was no mean feat. Ironkeeper Palluda of House Zuurthom proved up to the task. They made a name for themself by taking control of a trade deal that ultimately sold tens of thousands of gold pieces' worth of quarried stone to the \textit{League of Miracles}, which began the reconstruction of Tal'Dorei in earnest. Palluda became a hero of the working dwarf within \textit{Kraghammer}. They are a master at getting all sides to agree to a compromise, and do so with an iron-edged smile.

\subparagraph{Lord Nostoc Greyspine}
\textit{Dwarf overseer of House Greyspine}

The dour executive of dozens of quarries across the \textit{Cliffkeep Mountains}, Nostoc makes up for his unpleasant personality by being a formidable businessman. In recent years, he has become fabulously wealthy as the reconstruction of cities across Tal'Dorei devoured thousands of tons of stone. House Thunderbrand and the \textit{League of Miracles} made it possible to transport it all with teleportation spells.

\subparagraph{Elementalist Drake Thunderbrand}
One of the eldest of his House, Drake is known for his surprising empathy and fierce temper. He is also a member of the \textit{Arcana Pansophical}. His friendship with \textit{Allura Vysoren} made him a popular choice for \textit{Kraghammer's} ambassador to the council, but he politely refused, claiming it was a younger dwarf's duty.

\subsection{League of Miracles}
"Magic isn't cheap, but miracles are priceless." Such is the credo of the \textit{Wonderworker}, the mysterious thaumaturge whose words inspire and guide the \textit{League of Miracles}. This company of mercenary magi is the reason why Tal'Dorei has recovered so rapidly from the near-apocalyptic destruction of the \textit{Chroma Conclave}. \textit{Westruun} and \textit{Emon} rose from rubble to the height of their former glory in a matter of months.

The people of Tal'Dorei adore the league. Their mighty spellwrought \textbf{adranach} can be seen lifting thousand-pound bricks of limestone high into the sky to rebuild walls, towers, and even castles in a matter of days. Many of the league's spellwrights have become folk heroes and celebrities for their miraculous acts, which range from healing those maimed in the \textit{Conclave's} attack to transmuting raw timber, mud, and straw into full-fledged houses in minutes.

For most of Tal'Dorei's elite, the \textit{League of Miracles} is likewise a great ally. Trade and travel have resumed stronger than before, for the league has eagerly subcontracted their talented mages out as guards for caravans and persons of interest.

Yet, not everyone sees the league as an ally. In the early aftermath of \textit{Thordak's} defeat, the new and untested \textit{Tal'Dorei Council} put out thousands of bounties on reconstruction projects to get their towns and cities back on their feet, and local mayors and margraves across the nation agreed to these subsidies to reduce the cost of reconstruction for their already devastated populace. They expected Tal'Dorei's rebirth to be a multi-generational undertaking—yet the \textit{League of Miracles} snapped up these bounties like candy, and the league does not work their miracles cheaply. Bills for hundreds of thousands of gold pieces' worth of bounties found their way to the desks of local leaders and to the \textit{Tal'Dorei Council} by the end of the year. Such vast wealth for thirty years' worth of rebuilding in three years' time simply did not exist.

But this was all part of the \textit{Wonderworker's} plan, and their spellwrights executed it flawlessly. The league needs money to purchase their expensive spell components and to pay their talented mercenary membership, of course, but their true goal is the priceless currency of power, in the form of political favors from leaders ranging from small-town mayors, to city-ruling margraves, and even delegates to the \textit{Tal'Dorei Council}. Rumors even swirl that one or more of the members of the council itself are in the league's pocket.

The league's members are delineated into a simple hierarchy.

\subparagraph{Mercenary Mage}
Spellcasters drawn from across Tal'Dorei form the backbone of the \textit{League of Miracles}. The league works as an agency, handing mercenaries lucrative jobs, collecting the payment themselves, and paying out in a timely fashion—minus twenty percent. These mercenaries come from all walks of life. Some are aspiring adventurers seeking a less life-threatening first gig, while others are skilled arcanists with shady pasts, including \textbf{Remnant cultist} and double-dipping \textit{Myriad} agents with a talent for magic.

\subparagraph{Spellwright}
The most talented (or best-connected) members of the league are awarded the title of spellwright. Only two dozen such masters of the arcane exist across Tal'Dorei. These spellwrights know the secrets of how to create \textbf{adranach}, mighty constructs of magical energy that were used publicly to rebuild Tal'Dorei's ravaged infrastructure—and are also used in secret as hunter-destroyers of the league's enemies. These spellwrights sign a blood pact that connects their minds to the mysterious communications of the \textit{Wonderworker}, which all spellwrights are blood-bound to interpret and enact.

\subparagraph{Grand Thaumaturge}
The official title of the league's shadowy master is the Grand Thaumaturge, though most colloquially refer to them as the \textit{Wonderworker}. None know their true identity, or their ultimate goal. The spellwrights understand that the Grand Thaumaturge has a master plan that must be followed. Only these spellwrights have the privilege of hearing their master's psychic messages. Over the years, they have interpreted them to mean that the \textit{League of Miracles} must become the secret puppet-masters of all Tal'Dorei.

\subsubsection{Goals}
The \textit{League of Miracles} is comprised mostly of mercenary magi pulled from all corners of Tal'Dorei, with some even coming from abroad as news of the league's reputation spreads across Exandria. For most of these arcanists, their goal is simple: to make money. The goals of the league's spellwrights are more elaborate. They listen for the cryptic words of the mysterious Grand Thaumaturge and dedicate themselves to deciphering and enacting their grand design. The spellwrights' current aim is one of total control—once all of Tal'Dorei's social and political elite are in their pockets, they can shape the land as they see fit, all to the cheering adoration of the fawning people of Tal'Dorei.

\subsubsection{Relationships}
The \textit{League of Miracles} tries to keep cordial, if distant, relationships with every major power on Tal'Dorei. Their public image is shiny and unimpeachable. Many members of the \textit{Tal'Dorei Council} are already in their pocket and cast votes that further the league's self-interest. The \textit{Arcana Pansophical}—especially \textit{Allura Vysoren}, who also sits on the council—are deeply distrustful of this unregulated, uncommonly powerful organization.

The \textit{League of Miracles} has a tense relationship with the \textit{Chamber of Whitestone}, as they have long sought access to the \textit{residuum} unique to that territory. The spellwrights fear outright hostility with the legendary heroes of Vox Machina, and have supposedly abandoned their quest to mine \textit{Whitestone's} naturally occurring \textit{residuum}.

As "no questions asked" employers of mages, the league is in the good graces of many less scrupulous arcane factions, including the \textit{Remnants}. Necromancers, blood mages, Ruidus cultists, and other unsavory magi now exist under the league's protection.

\subsubsection{Figures of Interest}
These figures are major members of the \textit{League of Miracles} that you can use as NPCs in your game. You should also take inspiration from these figures to make up new NPCs for your own campaign.

\subparagraph{The Wonderworker}
\textit{Mysterious enactor of true miracles}

All spellwrights hear the authoritative—yet enigmatically distorted—voice of the \textit{Wonderworker}. The \textit{Wonderworker} has supposedly appeared in the flesh only three times throughout the league's twenty-four-year history, as a humanoid draped in robes of starlight and wearing a mithral mask. In those moments, they have produced true miracles—raising a castle from the ground in a matter of seconds, halting the rampaging silver dragon Saldarthoryn in her tracks and healing her twisted mind, and stopping the Center Slab of \textit{Kraghammer} from collapsing into the Pits in a catastrophic mining accident that would have cost countless lives. These legendary acts inspired utmost awe and faith in all who witnessed them, leading the spellwrights to create \textbf{adranach} in the \textit{Wonderworker's} image and enact their cryptic will.

\subparagraph{Blaine Kraverrogg}
\textit{Half-halfling mercenary necromancer}

Short and brimming with a malign lust for power, Blaine appears to be a short, skeletally skinny human who hides his lack of stature in billowing black-and-gold robes. He is a member of the \textit{Remnants}, but pays his bills by doing undercover grave-robbing for the league. He longs to reach the bottom of \textit{Shadebarrow}, a mysterious crypt somewhere on the \textit{Dividing Plains}. His obsession with \textit{Shadebarrow} stems from a hunger for the power of the Demon Prince of Undeath.

\subparagraph{Honor Kinnabari}
\textit{Tiefling blood mage and novice spellwright}

As the newest spellwright in the \textit{League of Miracles}, \textit{Honor Kinnabari} has a lot to prove. She was granted a chance to prove her quality in the eyes of the \textit{Wonderworker} when the league was contracted to rebuild the Market Ward of \textit{Westruun} after it was devastated by a herd of \textbf{magma landshark} that descended from the \textit{Cliffkeep Mountains}. Her red eyes and sheet-white skin and horns complement a traditional red-and-white mage's robe from the reclusive dwarven clans of the \textit{Alabaster Sierras}. She is fully dedicated to the league, but she is as honorable as her name implies; her fellow spellwrights fear that she could turn traitor, should the unsavory truth of the league be laid plain before her.

\subsection{Library of the Cobalt Soul}
The Cobalt Reserve in \textit{Westruun} is home to a monastic order so old, enduring, and tirelessly dedicated to the preservation of knowledge that most people in Tal'Dorei simply take it for granted. Just as birds have always flown and grass has always grown, the Library of the Cobalt Soul has always sent its agents across the world, taking minutes at public trials, archiving broadsheets, exploring ancient ruins, and copying scrolls. To all but the most dedicated scholars, the work of the Cobalt Soul seems like dreary, banal stuff.

Nothing could be further from the truth. Beneath the veneer of being a dopey order of bookish monks, the Cobalt Soul is one of the most dogged espionage organizations in the world. The secret agents of the Cobalt Soul, its expositors, are spies who wield the truth as a weapon more potent than any blade. The Cobalt Soul was born in \textit{Wildemount}, just after the end of the \textit{Calamity}, when the The Knowing Mentor was maimed by an evil god. Her most loyal followers went into hiding and vowed to continue her divine work in knowing all there is to be known—and most importantly, to use the truth as a holy weapon against lies and propaganda.

Each archive of the Cobalt Soul has a common hierarchy, and higher-ranking members of the Cobalt Soul can command their lower-ranking comrades from another archive.

\subparagraph{Monk}
The monks of the Cobalt Soul are their most visible members, for they travel the land in search of knowledge both modern and archaic. Not all monks are engaged in fieldwork, especially new recruits. Those who work entirely within an archive aid in the organization and preservation of information. Few know that these monks are dedicated to the The Knowing Mentor; the common opinion in Tal'Dorei is that they are ascetics, obsessively dedicated to the chronicling of knowledge, not devotees of an obscure god.

\subparagraph{Expositor}
Trained in the art of espionage, the expositors of the Cobalt Soul use their infiltration expertise to uncover the secrets of governments, world leaders, and obscure cults. In times of great crisis, undercover expositors have even incited revolutions by sharing the knowledge they discovered with people toiling under evil regimes. This true mission of the expositors is a secret even to low-ranking monks, who see them simply as well-trained field agents.

\subparagraph{Archivist}
Though they are not involved in field work, archivists play the vital role of organizing the overwhelming flow of information that the monks and expositors return to their archive. They also train new monks and perform administrative tasks, like negotiating with sellers of artifacts and with government officials to ensure the stability of their archive.

\subparagraph{Curator}
A small group of three to five curators manage the top-level operations of each archive. One of their number is given the special designation of High Curator, to whom all other members of their archive must ultimately answer.

\subsubsection{Goals}
The public mission statement of the Cobalt Soul is to archive knowledge for the betterment of all peoples of Exandria. The secret mission of the Cobalt Soul is to use their accumulated knowledge to root out tyranny and ensure the freedom of all people.

As instructed by their reclusive god, the The Knowing Mentor, the Cobalt Soul in Tal'Dorei makes their archives available to everyone, so that knowledge can be used to improve and enrich the lives of all. Despite this, the curators of the Cobalt Reserve privately determine what knowledge is too volatile to be freely disseminated. Typically, this privileged information is magical in nature, but it also includes secrets recovered by expositors that must be strategically employed as a weapon against unjust leaders, cults, and power-hungry factions.

The Cobalt Soul itself is not immune to corruption and complacency. Nowhere is this truer than in the Cobalt Soul's birthplace of \textit{Wildemount}, where it must work hand-in-hand with the authoritarian Dwendalian Empire. Its highest officials say that the Cobalt Soul has outlasted every tyrant on Exandria, and that they must be content with small acts of insurrection—subtly subverting propaganda and sharing the truth with any who will listen—while waiting for the right moment to strike a fatal blow. Nevertheless, this mantra of patience has caused some of its more radical members to abandon the order entirely. Typically, deserters of the Cobalt Soul are expositors, who see the ugliest side of the world firsthand again and again.

\subsubsection{Relationships}
The Cobalt Soul is an international organization with archives in the Dwendalian Empire and Clovis Concord of \textit{Wildemount}, the Republic of Tal'Dorei, the Dawn City of Vasselheim on \textit{Issylra}, and the great city of \textit{Ank'Harel} on \textit{Marquet}. Each of these archives acts independently under the orders of their commanding High Curator, though they must ultimately answer to the joint leadership of the Cobalt Soul: the High Curators of the Rexxentrum and Vasselheim archives.

Though Tal'Dorei is a freer society than the Dwendalian Empire, the Cobalt Soul knows better than to think that means it is free of corruption. From the Cobalt Reserve, their archive in the city of \textit{Westruun}, the Cobalt Soul keeps close tabs on members of the \textit{Tal'Dorei Council}, the leaders of \textit{Syngorn} and \textit{Kraghammer}, and other powerful individuals across the land.

The Cobalt Soul actively opposes the \textit{Remnants}—followers of The Whispered One, a tyrannical god of secrets who is diametrically opposed to the goals of the The Knowing Mentor. They are also concerned by the meteoric rise of a guild of arcanists called the \textit{League of Miracles}, whose political maneuverings have made puppets of dozens of delegates to the \textit{Tal'Dorei Council} and political leaders across the land. This worry has only deepened as rumors reach them of the league filling their ranks with \textbf{Remnant cultist} who seek to dominate Tal'Dorei.

\subsubsection{Figures of Interest}
These figures are major members of the \textit{Library of the Cobalt Soul} that you can use as NPCs in your game. You should also take inspiration from these figures to make up new NPCs for your own campaign.

\subparagraph{High Curator Greisalda Cassios}
\textit{Half-\textit{dragonblood} leader of the Cobalt Reserve}

The High Curator of the \textit{Westruun} archive of the \textit{Cobalt Soul} is a high-cheekboned, bald woman who would appear fully human if not for her slitted, silver eyes and the silver scales that run across her head, neck, and upper arms. She is regarded as one of the fastest minds in Tal'Dorei, for she has the rare gift of perfect recall. Monks tremble when her clear, unpretentious voice rings through an archive chamber, for Greisalda remembers acts of justice and personal slights with equal intensity, making her a powerful ally and a furious enemy. She is concerned that corruption is spreading throughout the \textit{Tal'Dorei Council}.

\subparagraph{Cressida Holt}
\textit{Orc expositor of the Cobalt Reserve}

The name \textit{Cressida Holt} is synonymous with style and effortless charisma—and also a blunt, merciless coldness. \textit{Cressida} is nonbinary, and uses both she and they pronouns interchangeably. Though she's fresh blood to the expositors of the Cobalt Reserve, she has already taken advantage of her license to kill on her espionage missions. They are slim for an orc, but still more muscular than many of their human colleagues. Their assignments take them across Tal'Dorei in search of political corruption. She is currently stationed in \textit{Emon} to observe a person of interest on the \textit{Tal'Dorei Council}, where she alternates between lurking on rooftops and striding through ballrooms in fine suits.

\subparagraph{Archivist Salomane Rothtree}
\textit{Half-elf archivist of the Cobalt Reserve}

A haggard, overworked administrator, Archivist Rothtree is known throughout the Cobalt Reserve as someone who tries too hard to prove himself and constantly bites off more than he can chew. Half-orc and half-elf, he has said that he hopes to accomplish as much in his lifetime as an elf could in theirs. He often is tasked with helping adventurers find information. Though he resents the task, he may slowly open up to heroes who treat him well.

\subsection{The Myriad}
\textit{Wildemount's} greatest organized crime ring is now becoming known throughout Tal'Dorei, but it came from humble beginnings. The Myriad began in the Dwendalian city of Yrrosa, as a small shipping company specializing in exotic goods and services, exploiting and selling curiosities from the distant continents of \textit{Othanzia} and \textit{Ank'Harel}. Over time, as legal trading became more difficult, the Myriad's illegal occupations became their best business. Through years establishing underworld contacts and tactical blackmail, the Myriad evolved into a powerful, well-veiled criminal faction. Over the past two decades, the Myriad's aspirations have spread beyond \textit{Wildemount}, and they are now deeply entrenched in Tal'Dorei's criminal underworld.

Known for displays of opulence and social grace, this misdirection enables the Myriad to seduce and control the lives and wills of whomever is deemed valuable to their plans. Dozens of truly amoral cutthroats and thieves noticed when the \textit{Clasp} began to work hand-in-hand with the \textit{Tal'Dorei Council}. They felt the dagger of accountability pressed against their backs, and broke with their former guild. Now, the Myriad has filled the vacuum of truly lawless criminality that the \textit{Clasp} left in its wake.

The Myriad operates as a loose network of gangs whose bosses run their own localized sects, without direct oversight from the mysterious heads of the society. Tithes and information are expected to be delivered to the leadership at a steady rate from each satellite of the organization, and should the success slow, threats of enslavement or destruction soon follow.

\subsubsection{Goals}
The \textit{Myriad} seeks wealth in hideous excess. The means are unimportant, as long as they're lucrative. No trade is too abhorrent, from bloody theft to kidnapping and slavery. And though they lack the political capital to bribe and puppeteer politicians the way they did in \textit{Wildemount}, the \textit{Myriad} are confident that they can get even the supposedly noble \textit{Tal'Dorei Council} in their pocket.

Most fronts for the \textit{Myriad} are businesses that publicly deal in antiquities and exotic textiles. Behind these masks is a ruthlessly efficient black market of slaves, magical beasts, weapons, and illicit substances. Overt violence tends to draw unwanted attention, and the \textit{Myriad} avoid it when possible, preferring to discredit, blackmail, or frame those who cross them. Should that fail, however, kidnapping or assassination can be a quick and traceless final alternative.

\subsubsection{Relationships}
Though they are relative newcomers to the criminal underworld of Tal'Dorei, the \textit{Myriad} is well-versed in the covert art of making money by any means necessary. The \textit{Clasp} is well aware of their presence in the broadest sense, and the two thieves' guilds have spilled blood in the streets of every city in Tal'Dorei at some point in the past twenty years. These days, the bloodshed has subsided into a cold war of espionage and counter-espionage, as each organization tries to sabotage the other.

The \textit{Myriad} has only been routed from cities a number of times—they no longer have any presence in \textit{Whitestone} or \textit{Syngorn}—and each time, it has been by a concentrated citizens' uprising. In both instances, when they sought the source, \textit{Myriad} leaders were furious to discover an empty safe house containing only a single gold coin, marked by a grinning face: the calling card of the \textit{Golden Grin}. Even across the \textit{Lucidian Ocean}, the old enmity between the \textit{Myriad} and the \textit{Golden Grin} remains today.

\subsubsection{Figures of Interest}
These figures are major members of the \textit{Myriad} that you can use as NPCs in your game. You should also take inspiration from these figures to make up new NPCs for your own campaign.

\subparagraph{Lustran "the Backstabber" Zeth}
\textit{Human boss of \textit{Stilben}}

Lustran was one of the first \textit{Clasp} members to turn double agent. By the time anyone realized his betrayal, he had granted the \textit{Myriad} their foothold in Tal'Dorei in the swampy port city of \textit{Stilben}. His lean, hungry face and graying hair give him a handsome, intimidating appearance. All \textit{Myriad} agents in Tal'Dorei respect him, and all members of the \textit{Clasp} hate him. He revels in his bad reputation as the renowned Backstabber of \textit{Stilben}.

\subparagraph{Tully "the Conclave of One" Dezzeram}
\textit{Gold \textit{dragonblood} boss of \textit{Emon}}

Tully's distasteful nickname delights her inner circle in \textit{Emon's} underground. The \textit{Myriad} boss of \textit{Emon} is a position that usually changes hands by the season, thanks to the diligent efforts of the \textit{Clasp}, but Tully has evaded capture for three years, and grown rich and powerful in her reign of smuggling, trafficking, and grand burglary.

\subparagraph{Sylker "the Millionaire" Uttolot}
\textit{Half-elf boss of \textit{Kymal}}

This gold-eyed half-elf, a former underling of the crime families of Shady Creek Run in \textit{Wildemount}, joined the \textit{Myriad} and became the first individual on that continent to hoard one million gold pieces. He has a masterful mind for money, and spearheaded the \textit{Myriad's} infiltration of \textit{Kymal}. "Gold is a disease," he often says, "and I'm bloody sick."

\subsection{The Remnants}
In a time long past, an archmage of unparalleled power known as The Whispered One attempted to ascend to godhood by conducting the Ritual of Seeding from high atop the citadel of Thar Ampala within the Plane of Shadow. The Whispered One surrounded himself with a cult of fanatical mages and empowered murderers to protect him during his apotheosis. Despite these protections, he failed. His cult, however, survived, and helped him return to power centuries later, after the defeat of the \textit{Chroma Conclave}. He succeeded in his bid for godhood, but was kept from utterly dominating the Material Plane at the last second by Vox Machina, who banished him beyond the Divine Gate, where he was shackled by all the same limitations as the other gods.

His cult, the Remnants, went into hiding once again. Their fanatical ranks, bolstered by his true apotheosis, now seethe in secret, convinced that there is a way to grant their god the ultimate power he sought. Many of them have wormed their way into positions of power within Tal'Dorei's political institutions, churches, arcane organizations, and criminal outfits. Many more have joined up with the dubiously moral, profit-driven arcanists of the \textit{League of Miracles}. They exist in shadows, lurking, waiting, longing for a new dark prophet to lead them to glory.

\subsubsection{Goals}
The Whispered One achieved an incomplete victory over twenty years ago, and now his cult is building a new generation of faithful to make his megalomaniacal dreams a reality. The cult's leadership believes that reuniting these malign artifacts will give them greater power to commune with their god and let them begin their quest in earnest to reach him through the Divine Gate. This leadership is known as the Circle of the Dismembered. Each cult leader in the circle commands a branch of the cult, named for a part of The Whispered One body. These factions are:

\subparagraph{The Unveiled Sight}
Members of this cult are tasked with scouring the planes for arcane lore that could restore their master to the Material Plane.

\subparagraph{The Mordant Voice}
Members of this cult serve their master by infiltrating other factions across Tal'Dorei and subverting their purposes with subtle propaganda.

\subparagraph{The Devouring Blood}
Members of this cult specialize in the art of silent death. They train in stealth and necromancy to become truly horrific assassins.

\subparagraph{The Unseen Touch}
Members of this cult are trained in magic that corrupts the mind and perverts the will, making them vile puppet-masters that prepare the world for their god's return.

\subparagraph{The Insatiable Heart}
Members of this cult do the vital ground-level work of seeking out new blood. They are trained to exploit the emotional weaknesses of their targets to convert them to the cult, without the need of magic.

\subsubsection{Relationships}
The \textit{Remnants} remain largely unknown, even after their master's grand ascension, and wish to keep it that way. Those who are aware of their presence are either cult members or people who have uncovered their insidious existence and wish to stamp it out wherever it may appear. Thus, there are rarely any alliances made by or with the \textit{Remnants}.

\subsubsection{Figures of Interest}
These figures are major members of the \textit{Remnants} that you can use as NPCs in your game. You should also take inspiration from these figures to make up new NPCs for your own campaign.

\subparagraph{Quinton Puck}
\textit{Human jeweler, Master of the Unveiled Sight}

While he runs a middling jewelry business within \textit{Emon}, Quinton is secretly the leader of the Unveiled Sight. He is a master of creating and piercing illusions, deciphering ancient script, and all things to do with the magical arts of obscurement and unveiling. He can be recognized by his blood-red spectacles.

\subparagraph{Ixrattu Khar}
\textit{Tiefling vampire, Master of the Devouring Blood}

This escapee from the \textit{Black Bastille} prison in \textit{Emon} is Tal'Dorei's most infamous mass murderer, who once slaughtered and consumed hundreds of innocents on a deadly killing spree from \textit{Kymal} to \textit{Emon}. During her imprisonment, Ixrattu was plagued by shadowy visions of a woman beckoning to her from beyond the bars. The visions pointed her towards weaknesses within the \textit{Black Bastille's} guard, and helped her achieve the unthinkable: escape. Once free, she vanished for a year as she scoured the countryside, puzzling over the cryptic whispers of her continuing visions.

At the end of her quest was a fateful discovery: within a lightless underground passage, she found the venomous crossbow \textit{Condemner}, a \textit{Vestige of Divergence}. Now she lurks in \textit{Kraghammer}, where she trains a network of elite assassins who gleefully dispatch anyone the Unveiled Sight identifies as an enemy of the \textit{Remnants}.

\subsection{Wardens of Syngorn}
The elven city-state of \textit{Syngorn} was established in the \textit{Verdant Expanse} by the sorceress Yenlara, the wood \textit{elves'} first leader after the \textit{Divergence}. Today, the physical and cultural wellbeing of \textit{Syngorn} is safeguarded by the three offices of elders called \textit{Wardens}, united by the High Warden. Their primary duty is to protect the city and the surrounding lands from intrusion.

\subsubsection{Goals}
The \textit{Wardens} believe that preserving the ancient culture of \textit{Syngorn} is of utmost importance. Even as their fate grows ever more closely tied to the fate of all Tal'Dorei, the traditions that make \textit{Syngorn} unique must be conserved. The three offices of the \textit{Wardens} are the Verdant Lord, the \textit{Voice of Memory}, and the Guildrunner. The office of Verdant Lord commands \textit{Syngorn's} formidable military, which is maintained and scattered throughout the \textit{Verdant Expanse}, ready to be called to \textit{Syngorn's} defense at a moment's notice. As preservers of culture, the \textit{Wardens} also encourage public respect for the arts, including the art of magic. It's not uncommon for a \textit{Syngornian} to rise in social status after being publicly praised by a warden for a masterful artistic creation or performance.

\subsubsection{Relationships}
\textit{Syngorn} has long been an isolationist society. As a mostly elven enclave, they have often been criticized as xenophobic, or even racist, toward outsiders. There is truth to these claims. Nevertheless, the rise and fall of the \textit{Chroma Conclave} has goaded even some of \textit{Syngorn's} most conservative voices into strengthening the bonds of trade, travel, and kinship with the other free peoples of Tal'Dorei. Even the people of \textit{Kraghammer}, who many \textit{elves} still remember siding with the tyrant \textit{Drassig} in centuries past, have been welcomed as friends and allies for the loyalty and courage they displayed in the battle against \textit{Thordak}.

\subsubsection{Figures of Interest}
These figures are major members of the \textit{Wardens of Syngorn} that you can use as NPCs in your game. You should also take inspiration from these figures to make up new NPCs for your own campaign.

\subparagraph{High Warden Tirelda}
\textit{Elf Warden of \textit{Syngorn}}

Though she is in her twilight years, even by elven standards, High Warden Tirelda is a trusted descendent of Yenlara, founder of \textit{Syngorn}. She is skilled with magic and blessed with calm wisdom. Thus, she was entrusted with the highest level of responsibility amongst the \textit{Wardens}—to maintain law, guide society, and preside over all other \textit{Wardens}.

\subparagraph{Verdant Lord Celindar}
\textit{Elf Warden of \textit{Syngorn}}

Head of the Verdant Guard of \textit{Syngorn}, and an accomplished warrior in their own right, Celindar handles all military strategy and instruction within the forest's bounds. They only return to \textit{Syngorn} when called by their fellow wardens, and can otherwise constantly be found prowling the \textit{Verdant Expanse}. They use any pronouns, and have relinquished any sort of gender presentation after receiving a vision from the The Arch Heart.

\subparagraph{Ouestra, the Voice of Memory}
\textit{Elf Warden of \textit{Syngorn}}

A renowned historian and a beloved vocalist, Ouestra is charged with the preservation of \textit{Syngorn's} history and culture. It's said her most magical songs can control moonlight and bring light even to the Abyss itself. These days, she has grown somewhat unpopular by also singing songs of \textit{Syngorn's} darker days, for she says true greatness can only be inspired by learning from failure.

\subparagraph{Guildrunner Rawndel}
\textit{Elf Warden of \textit{Syngorn}}

One of the oldest living \textit{elves} within \textit{Syngorn}, Rawndel is more cunning and insightful than most, allowing him to be a very effective head of commerce, business, and the treasury.

\chapter{Tal'Dorei Gazetteer}
\begin{DndReadAloud}{}
\DndQuote%
{}
{''I've wandered over many a shore and plain in my journeys, and my mind still finds itself drawn back to Tal'Dorei. Ain't no place in this world more vibrant; the land feels young. Unscarred. Mind you, that sort of calm causes my old bones to itch, too. Folk in that land are kind and honest, and its vistas are second to none, yet...well, I'll just say that cities are cities, and the wilds are the wilds. Tal'Dorei's got plenty of both. Keep a sword or a wand with you wherever you go, adventurer.''}
{Norad Firth}
\end{DndReadAloud}

\DndDropCapLine{T}{}his chapter contains a wealth of information about the lands of Tal'Dorei, the adventures to be found therein, and the many centers of civilization to be found across its vast wilds. Consider this chapter a compilation of a wide variety of assorted notes, research, and scattered publications of historians of the \textit{Library of the Cobalt Soul} and the \textit{Alabaster Lyceum}.


Boundless possibilities await eager adventurers of all kinds who explore Tal'Dorei, from the vast fields of the \textit{Dividing Plains} to the treacherous crags of Tal'Dorei's many mountains, from the storm-wracked \textit{Lucidian Coast} to the immemorial depths of the \textit{Rifenmist Jungle}. Countless nomadic civilizations exist within these realms of magic and monsters, as do dozens of tiny villages, thriving towns, and bustling cities. The people of Tal'Dorei do what they can to survive, and share the tales of treasures and adventure with any who will listen.  Consider this gazetteer a compilation of tales and fragmented histories, passed down from bards and scholars. Your campaign can follow its descriptions to the letter, or diverge as wildly as you see fit.

\section{Settlements}
Every permanent settlement in Tal'Dorei has a few basic statistics that impart a top-level impression of that place. This includes what type of settlement it is—typically a village, town, or city—its population, and its rough demographics. Don't feel restricted by these demographics if you want your character to come from a specific location but you don't see their race listed there. People of all races can be found in just about every community in Tal'Dorei, and your character can come from anywhere.

Additionally, cities and large towns often have a number of small rural communities outside their fortified walls and population center. The population numbers of these larger settlements include the people in the outlying areas. On the other hand, the official populations of a village or other small settlement just includes the people living in the village itself.

In fact, though there are a few villages and small settlements listed in this gazetteer, Tal'Dorei is home to hundreds, if not thousands, of tiny settlements too small to detail or place on the map. The Game Master can make up any number of villages or other communities wherever they wish in Tal'Dorei without changing any of the lore presented in this gazetteer. And of course, the GM can also create large settlements unique to their version of Tal'Dorei—the only difference is that they should be prepared to think about how their version of Tal'Dorei is different, now that they've added another hub of population, trade, politics, and culture.

\section{Story Hooks}
This chapter contains a variety of story hooks that you can use as a starting point for Tal'Dorei adventures of your own creation. \textit{Chapter 5} of this book contains advice for expanding these hooks into full adventures, and even more advice for creating RPG stories can be found in the fifth edition core rules.

These plot hooks are tied to specific locations, and all have a level range. These level recommendations aren't strictly mandated, but have generally been chosen either because of the monsters involved, because of the narrative weight of the location or characters involved, or to indicate how epic and world-affecting the adventure should be.


\begin{itemize}
\item For low-level characters: 1st to 4th level
\item For mid-level characters: 5th to 10th level
\item For high-level characters: 11th to 16th level
\item For epic-level characters: 17th to 20th level
\end{itemize}

Additionally, some plot hooks are marked as "any level." These hooks present stories broad enough to suit characters of any level, anywhere in their party's story. These hooks generally don't suggest specific monsters to use as villains. When creating adventures based off of these hooks, use the encounter-building advice in the fifth edition core rules to choose appropriate foes for your players.

\section{Lucidian Coast}
The stormy Lucidian Ocean crashes furiously against the rocky coasts of eastern Tal'Dorei, chewing inlets into the coastline and spurring vast banks of impenetrable fog across the waters. Most ships rightly avoid the volatile northern waters, instead docking at southern ports like \textit{Stilben}. Adventurous folk leave their seafaring vessels behind and brave the \textit{Mooren River Run} to \textit{Drynna}. Yet even there, the Lucidian Coast is not without its dangers; vicious wildlife stalks the wilderness, and flocks of \textbf{harpy} have harried the shores since time immemorial, drawing lost sailors to their doom upon the rocks.

\begin{DndReadAloud}{}

\begin{itemize}
\item{Majority Faiths:}
The Changebringer, The Wildmother, The Stormlord

\item{Minority Faiths:}
The Lawbearer, The Strife Emperor, The Scaled Tyrant

\item{Imports:}
Livestock, lumber

\item{Exports:}
\textbf{Fish}, ships, grain, \textit{iron}, \textit{silver}, \textit{gold}, trade goods from \textit{Wildemount}

\end{itemize}

\end{DndReadAloud}

\subsection{Drynna}

\begin{description}
\item[Town:]
Population 2,765 (75\% \textit{humans}, 8\% \textit{gnomes}, 8\% \textit{halflings}, 4\% \textit{elves}, 5\% \textit{other races})

\end{description}

Along the shore of \textit{Mooren Lake} rests the sleepy town of Drynna. The community thrives on the bountiful wildlife around the lake, and the Drynna fisherfolk are responsible for most of eastern Tal'Dorei's freshwater fish supply. The social guidance of the city revolves around the Sunrise Lodge (or simply "the Lodge"), a membership of community leaders who meet to discuss and vote on important social and legal matters. The Lodge has also sent a variety of increasingly frustrated petitions to the \textit{Tal'Dorei Council}, since their community's small size means they have traditionally lacked a delegate of their own on the council. However, the recent surge in prospectors has increased their population enough for a member of the Lodge to attend council sessions, and the Lodge is doing everything it can to make sure the boom lasts as long as possible before it falls into bust.

\subsubsection{Boomtown}
\textit{Drynna} has a small militia to protect the town from wild beasts and roving monsters, and to capture horse thieves and other petty criminals. This militia is presently overwhelmed by hundreds of prospectors who have swept into town over the past three years after hearing rumors that chunks of \textit{whitestone} have been found in \textit{Mooren Lake}, washed down the river from the \textit{Alabaster Sierras}.

Given the high price of \textit{residuum}, this new occurrence has made \textit{Drynna} into something of a boomtown. Thousands of small homesteads now surround the town center, and burglary and banditry are both on the rise.

\subsubsection{Drynna Adventures}
Game Masters who set their adventure in \textit{Drynna} can use these plot hooks for inspiration.

\subparagraph{Unfriendly Waters (mid-level)}
Recently, \textit{Drynna} has suffered a number of losses, with prospectors and fisherfolk being attacked. Worse still, the supply of fish itself seems to be declining rapidly; without this local food staple, villages across eastern Tal'Dorei will certainly starve. The party hears the call for mercenaries reaching as far as \textit{Westruun} and \textit{Stilben}. The culprit is a \textbf{hydra} within the lake that has gained intelligence by living for so long in the \textit{residuum}-infused waters. It has an Intelligence score of 12 (+1), can speak Common, and commands a tribe of \textbf{lizardfolk}.

\subparagraph{Bleak Inheritance (any level)}
Last year, Hareth Val Bardo, a younger member of the Sunrise Lodge, made a pact with the Moon Mistress (see "\textit{Mooren Lake}," below). Hareth's elder brother, Pauvren, died in his sleep shortly thereafter, leaving the entire Val Bardo family fortune to Hareth, the only surviving heir. However, nightmares now haunt Hareth every night, and violent urges surge through his mind. The Lodge, having noticed Hareth's paranoid eyes and twitchy hands, calls upon outsiders to investigate.

\subsection{K'Tawl Swamp}
The sprawling swamplands of \textit{K'Tawl} cover much of the southern \textit{Lucidian Coast}, enveloping \textit{Stilben} and the K'Tawl Bay. The saltwater swamp is thick with cypress and tupelo trees reaching up from muddy waters that range, in places, from one to ten feet deep—or more. The hot, muggy air and thick, sludge-covered landscape host all manner of terrible beasts, including venomous \textbf{flying snake} and man-eating \textbf{giant frog}.

Bands of frog-headed humanoids harry the edges of \textit{Stilben} from their deep-marsh encampments, while a vicious school of \textbf{sahuagin} called the Ghormauth command the shoreline of \textit{K'Tawl}. The bloodthirsty \textit{Ravagers} use the forested edges of the swamp to ambush prey—as do countless bandits and petty thieves. As such, caravans and messengers must remain on constant alert for the possibility of an ambush. With such dangers lurking in the swamp, travel to and from \textit{Stilben} is rarely made without an armed escort.

\subsubsection{K'Tawl Swamp Adventures}
Game Masters who set their adventure in \textit{K'Tawl Swamp} can use these plot hooks for inspiration.

\subparagraph{They Shamble (mid-level)}
People have been disappearing from the isolated swamp-burghs within the \textit{K'Tawl}. The characters, perhaps en route to \textit{Stilben}, spend the night in the village of Dunghill. After hearing tales of the disappearances, the characters awaken the next morning to screams—as a terrified human is carried off within the rotting orifice of a \textbf{shambling mound} of swamp vegetation.

\subparagraph{Monolith of the Hells (high-level)}
Why have members of the Torchraiser herd been found dead, their hearts exploded and their teeth plucked from their mouths? Why have the stagnant \textit{K'Tawl} waters suddenly begun to flow toward the dismal center of the swamp? A divine messenger of the The Everlight visits the characters, its celestial form torn and wavering, and dies before it can impart its request. On the angel's lips is a single dying phrase: "Destroy K'Tawl's Heart."

\subsection{Mooren Lake}
North of the \textit{Summit Peaks} is a shimmering lake fed by the snowmelt from the \textit{Alabaster Sierras} by the \textit{Mooren River Run}. \textit{Mooren Lake} is surrounded both by bountiful wildlife and the frightful creatures attracted to such a thriving ecosystem, including roving \textbf{lizardfolk}, flocks of \textbf{harpy}, and several \textbf{hydra} that lurk in the lake's depths.

Centuries of \textit{whitestone} runoff from the \textit{Alabaster Sierras} have allowed deposits of \textit{residuum} to collect at the bottom of the lake—and the \textit{residuum} has formed into a thrumming organ of raw magical energy, beginning to gain sentience. The twin islands in the center of the lake are blanketed by an ever-present mist, and folk tales tell of a beautiful, ageless woman, known only as the Moon Mistress, who watches over the island shores. It's said that any mortal who seeks an audience is judged by her, and is either granted a favor, a boon of unveiled knowledge, or a swift and terrible end.

\subsubsection{Mooren Lake Adventures}
Game Masters who set their adventure in \textit{Mooren Lake} can use this plot hook for inspiration.

\subparagraph{Renew the Mists (mid-level)}
The Moon Mistress, a \textbf{night hag} who takes the guise of a water nymph, uses the fading mists to shroud her location from a blood hunter named Perron Brill, a woman dedicated to hunting the Mistress down out of vengeance. The mists require a blasphemous mortal sacrifice to renew, and the Mistress is willing to reward those who would aid her in retrieving a victim—willing or otherwise. Perron Brill resides in nearby \textit{Drynna}, and though she is no saint herself, more honorable characters may think it best to help the hunter find her prey in \textit{Rootgarden Marsh}.

\subsection{Rootgarden Marsh}
The shallows of \textit{Mooren Lake} leak into the tender soil of the lake-bed and spread eastward, creating the cold and desolate \textit{Rootgarden Marsh}. Tall, thin trees lift five or more feet above the muck upon stilt-like roots, like freshwater mangroves. Their sparse placement gives them the appearance of a village of scattered huts. The fresh water mingles with salt here, leaving only adapted vegetation. The smell of moss and rot fills the air, while strange, feral creatures leap from tree to tree, or patiently wait from within their dome-like tree-root dens to snatch wary travelers. An organized group of gnoll marauders called the Moonsteeth keep their cavern hideout here, striking out at easy targets outside of \textit{Drynna} and \textit{Mooren Lake}.

Since the region is populated almost entirely by monsters, most humanoids avoid it whenever possible. Those few who explore the deadly marsh do so only out of necessity, either because they need to unlock Rootgarden's many secrets, or because "civilization" is even more dangerous to them than the primal dangers found within.

\subsubsection{Rootgarden Marsh Adventures}
Game Masters who set their adventure in \textit{Rootgarden Marsh} can use these plot hooks for inspiration.

\subparagraph{The Ley-Knot (mid-level)}
Without warning, arcane magic stops functioning within a hundred-mile radius of \textit{Rootgarden Marsh}. As panic spreads across northeastern Tal'Dorei, the characters or a patron of theirs discover that a major ley line, the source of arcane power in this part of the world, has been twisted and tied into a knot at the center of the marsh. This knot has been tied by High Moonspeaker Luneskon of the Moonsteeth \textit{gnolls}. Her ambitions are small, but whatever method she used to tangle the ley lines must be erased from history—knowledge of its existence could bring disaster to all of Tal'Dorei.

\subparagraph{Death's Garden (any level)}
Rootgarden's strange ecology has produced many distinctive and powerful forms of plant life. During the course of a campaign, a group of adventurers may need to create a unique potion requiring the essence of manawort, a sentient mushroom found only within Rootgarden. The manawort mushroom is fearful of being harvested, and has made close friends with a black dragon named Thardraxxus, who keeps the mushroom and its sporelings safe. This dragon can be of any age, depending on what level the characters are.

\subsection{Stilben}

\begin{description}
\item[Large Town:]
Population 10,660 (67\% \textit{humans}, 11\% \textit{elves}, 11\% \textit{halflings}, 5\% \textit{half-elves}, 6\% \textit{other races})

\end{description}

Stilben stinks to high heaven. Most sailors who pass through it—be they honest traders of the Clovis Concord in \textit{Wildemount} or pirates seeking their next conquest—call this dismal burgh the Rotted Lot. The stink and the crime that run through this town make the name an apt one, but things aren't all bad in Stilben. The largest port on this coast of Tal'Dorei, Stilben is the focal point of international trade with \textit{Wildemount}. Even though Stilben's permanent population is a scant ten thousand, its more transient population of sailors who stay for a few days or weeks at a time is just as large, and fully triples Stilben's total population in the summer months.

Predictably, the vast majority of Stilben's business caters to folk who won't stay for long. There are brothels aplenty, at least one inn and one tavern each on every street, and countless other comforts for the weary traveler. The town is nestled in the sticky depths of the \textit{K'Tawl Swamp}, and its putrescent stink mingles with the humidity and the ever-present buzz of insects to produce an air of misery around Stilben's outer districts. Pushed out of Stilben's cozy, cosmopolitan heart, its poor and disenfranchised citizens—derisively called muckdwellers by the town's bigshots—struggle to eke out a living in the \textit{K'Tawl's} murky waters.

Central Stilben houses most of its businesses, residences, and attractions. Those living in Stilben's heart like to present a clean, "civilized" view of their town, maintaining it as a thriving commercial hub. The cross-continental \textit{Silvercut Roadway} leads from Stilben all the way to \textit{Emon}, and its robust trade of rare goods, textiles, and spices has given rise to a formidable merchant ruling class. Though the aging Margrave Wendle Truss is Stilben's ruler in name, it is the guilds that hold true power. Their vast wealth and sizeable "protection" force have impelled the bombastic-yet-ineffectual margrave to acquiesce to the guilds' will in order to maintain peace, as well as to preserve his own station within the city.

Order is kept by the Waterwatch, Stilben's official constabulary, but corruption is common within their ranks. The past few decades have been hard for the \textit{Clasp} in Stilben. Though the local Spireling has kept a tight grip on the margrave and the Waterwatch, the \textit{Myriad} has used bribes, intimidation, and every dirty trick in the book to winnow away the \textit{Clasp's} power. Now, Stilben is little more than a crucible for the power struggle between the two criminal organizations.

\begin{DndSidebar}[float=!b]{Origins of Vox Machina}
The citizens of \textit{Stilben} do not know that some of Exandria's greatest heroes, Vox Machina, first met in their humble city and went on their first adventure in the muck of the \textit{K'Tawl Swamp}. Back then, even the exalted Vox Machina were just two-copper adventurers, quickly forgotten by \textit{Stilben's} haughty elite. One dwarf named Heinrich Runescribe, a junior member of the \textit{Alabaster Lyceum's} revered Lorekeeper Society, practically worships the heroes that defeated the \textit{Chroma Conclave}, and has made it his mission to record all of their travels. He has found a wealth of evidence that the heroes of Vox Machina first gathered in \textit{Stilben} under inauspicious circumstances and were dragged into an unexpected battle against the \textit{Myriad}.



\end{DndSidebar}

\subsubsection{Stilben Adventures}
Game Masters who set their adventure in \textit{Stilben} can use these plot hooks for inspiration.

\subparagraph{Lost Below (low-level)}
A \textit{Stilben} guild known as the Relic Seekers partnered with the \textit{Clasp} to smuggle a hoard of stolen \textit{Syngornian} artifacts through the heart of the \textit{K'Tawl Swamp} to the port of \textit{Stilben}. The smugglers were last seen by a \textit{Clasp} agent watching a bridge in the western swamp, but two weeks have passed since their last contact. They should have arrived by now. The characters are approached by one of several groups: an agent of the Relic Seekers in search of the relics, one of the \textit{Clasp} seeking to find their smugglers, or a \textit{Myriad} taskmaster who seeks to undermine the \textit{Clasp}. Their first contact: the bridge guard, a hulking tiefling \textbf{Clasp enforcer} named L'Arkhelle.

\subparagraph{Ashari Agitator (mid-level)}
Gang warfare between the \textit{Clasp} and the \textit{Myriad} has been a daily affair in \textit{Stilben} for the past thirty years, but something is changing. A cruel \textbf{Ashari waverider} named Galessi has joined the \textit{Myriad} and is using her elemental magic to turn the tide of this gang war. The \textit{Clasp}, the Waterwatch, and even the Air Ashari of \textit{Zephrah} want to see Galessi laid low, but she has a secret base within the swamp, protected by a charmed \textbf{giant crocodile} and putrid \textbf{water elemental}.

\subparagraph{Demonblood Vengeance (high-level)}
A \textbf{marilith} demon using the guise of a blonde human named Iselda has been lurking on the outskirts of \textit{Stilben} since her defeat at the hands of Vox Machina decades ago. Her demonic ichor has spawned a brood of \textbf{demonfeed spider} that now view her as their mother. Iselda, sensing an opportunity, dedicated herself to the The Spider Queen and attracted a small legion of dark elf fanatics to her side. Aided by all these evil creatures, she plans on conquering \textit{Stilben}, not for the \textit{Myriad}, but for herself and her evil god.

\subsection{Summit Peaks}
The Summit Peaks loom over the southeastern plains of Tal'Dorei, like great spires of rock reaching for the sky. The tall, thin mountains burst so aggressively from the surrounding fields and swampland that even travelers from distant lands recognize them as a landmark of eastern Tal'Dorei. They are a welcome sight for any ship's crew making their way to shore. The mountain valleys recess into marshy gorges that house all manner of creatures migrating from the \textit{K'Tawl}. Reclusive \textbf{hill giant} and \textbf{stone giant} also make their homes in the marshy valleys. As the mountains rise higher, strong winds blow across the battered rock, leaving strangely smooth peaks that harbor \textbf{griffon} lairs, \textbf{harpy} dens, and convocations of \textbf{giant eagle} hunting the lower lands.

\subsection{Zephrah}

\begin{description}
\item[Small Town:]
Population 1,310 (62\% \textit{half-elves}, 12\% \textit{elemental ancestry}, 10\% \textit{humans}, 9\% \textit{halflings}, 5\% \textit{tieflings}, 2\% \textit{other races})

\end{description}

High atop a mountain aerie within the \textit{Summit Peaks} sits Zephrah, the tribal home of the Air Ashari people, keeping vigil over a gate into the Elemental Plane of Air. \textit{The Ashari} people worship and live by the balance of the natural world, and the people of Zephrah have adapted to life among the clouds. Most of the tribe consists of skilled hunters and climbers who hunt \textbf{mountain goat}, large birds, and other cliff-dwelling creatures, along with subsisting on high-elevation vegetation. Utilizing the proximity of powerful elemental magic, the druidic masters of the Air Ashari craft skysails to ride the ever-present winds, allowing their people to traverse the sheer cliffs of the mountains. The leader of Zephrah is called the Voice of the Tempest; they oversee all travel through the rift, all major druidic rituals, and any visiting outsiders.

Though it's still a small settlement, Zephrah has grown greatly in the past several years. Many people with an innate connection to the element of air have petitioned the current \textbf{Keyleth, Voice of the Tempest} for her permission to live here among \textit{the Ashari}. Likewise, a number of \textbf{druid} and other spellcasters who wish to learn the ways of the Air Ashari have come to live in relative isolation here—after they have proven their good intentions and respect for the sanctity of the skies.

\subsubsection{Diplomacy and Developments}
\textit{Zephrah} has grown in more than just population. The close friendship shared between \textbf{Keyleth, Voice of the Tempest}, and the genius inventor \textbf{Percival de Rolo} of \textit{Whitestone} has led to a number of modern technologies being adopted in \textit{Zephrah}. Windmills now dot the cliffsides, milling grain purchased from farms on the \textit{Dividing Plains}. \textbf{Keyleth, Voice of the Tempest} continues to work with \textbf{Percival de Rolo} to improve the lives of her people through technology and innovation. The operation of two of these mills is also kept private from the rest of \textit{Zephrah}, and they're used to refine \textit{whitestone} into \textit{residuum} for the construction of magic items to arm the \textbf{Ashari skydancer}.

\textbf{Keyleth, Voice of the Tempest} confident leadership has also guided her to make further diplomatic ties between her people and the rest of Tal'Dorei. These attempts at breaking down the traditional isolationism of \textit{the Ashari} have been met with skepticism from within, slowing the process to a crawl. Allowing any outsiders into \textit{Zephrah} at all has been seen as a huge step. Beyond this, \textbf{Keyleth, Voice of the Tempest} has pushed for greater economic and political involvement, so that the needs of \textit{Zephrah} can be heard in the \textit{Tal'Dorei Council}. Resistance toward political integration has mostly been internal within \textit{the Ashari}, but \textbf{Keyleth, Voice of the Tempest} considers any progress at all a good start.

\subsubsection{Defense}
\textit{Zephrah} has never maintained a formal army, though past Voices of the Tempest have trained martial forces and druidic warriors to defend their home against the monsters of the \textit{Summit Peaks}. As \textbf{Keyleth, Voice of the Tempest} attempts to build bonds between her people and the rest of Tal'Dorei, she has expanded the training of new \textbf{ashari skydancer} and created a unit of druidic warriors capable of responding to any emergency across Tal'Dorei. If a threat like the \textit{Chroma Conclave} were to rise again in Tal'Dorei, the \textbf{ashari skydancer} of \textit{Zephrah} would be there to help save the nation's people.

Thanks to \textit{Zephrah's} mountainous terrain, \textit{the Ashari} enclave has never fallen under siege. Because of the dangerous beasts that lurk in the mountains, most Zephran adults have learned how to wield a \textit{spear} or a \textit{staff} to defend themselves. Roughly one-in-ten people have undergone a small amount of elemental training—enough to be able to cast the \textit{gust of wind} spell a number of times per day equal to their Wisdom modifier (minimum of once).

Were \textit{Zephrah} to fall under attack, the \textbf{ashari skydancer} would defend the town and its elemental rift for as long as possible while the rest of its people evacuated to hidden retreats in the mountain caves.

\subsubsection{Zephrah Adventures}
Game Masters who set their adventure in \textit{Zephrah} can use these plot hooks for inspiration.

\subparagraph{Poached Eggs (low-level)}
The local Ashari have noticed a terrible trend of poachers climbing the peaks to steal entire \textbf{griffon} nests full of eggs, leaving the parents slain. The efficiency, brutality, and growing rate of these attacks suggest the hunters are dangerous professionals, and \textit{the Ashari} are seeking aid in finding and punishing those responsible before the locally endangered creatures are all dead or taken. The culprits, human poachers (use \textbf{scout} statistics) led by a \textbf{quasit} demon named Ryzzix, are fencing the stolen eggs on the \textit{Stilben} black market.

\subparagraph{A Welcome Overstayed (mid-level)}
Shaktia, an overbearing \textbf{djinni} soldier from the Plane of Air, has recently discovered and come through the elemental rift into \textit{Zephrah}. However, the \textbf{djinni} refuses to leave, instead choosing to bother the people with pranks, debates, and invitations to become his well-paid servants. The \textbf{Keyleth, Voice of the Tempest}, at the end of her patience, seeks adventurers to capture Shaktia and return him through the rift to the Elemental Plane of Air. However, when the adventurers learn what caused the \textbf{djinni} to leave his home plane in the first place, they may find it hard to force him to leave.

\section{Alabaster Sierras}
In the northeastern reaches of the continent, where the waters of the \textit{Lucidian Ocean} break between the rocky shores of Tal'Dorei, a lush treeline fills the valleys between the chalk-white bluffs of the Alabaster Sierras. This gelid northern territory has for decades been beyond the reach of the \textit{Tal'Dorei Council}, and fallen under the rule of the \textit{Chamber of Whitestone}. However, over twenty years ago, thousands of \textit{Emonian} refugees fled to Whitestone to ride out the devastation of the \textit{Chroma Conclave}. These refugees included the bereaved family of the Last Sovereign. Since those dark days, the \textit{Chamber of Whitestone} has enjoyed a close relationship with the \textit{Tal'Dorei Council}. Lady \textbf{Vex'ahlia}, a member of the \textit{Chamber of Whitestone}, also sits on the \textit{Tal'Dorei Council} itself, further strengthening the bonds of unity between this once-isolated land and the rest of the realm.

Beyond \textit{Whitestone}, the region's vibrant-yet-gothic urban center, the Alabaster Sierras are framed by dangerously windy shoreside cliffs and dry, desolate mountains swarming with \textbf{harpy}, cold-adapted \textbf{wyvern}, and other hungry monstrosities. Beneath the peaks, the valley opens into the thickly wooded pine forest known as the \textit{Parchwood}. This dense, temperate rainforest is beautiful and lush, but many shadowed threats loom in the dark of the wood.

The northern center of the valley marks a single clearing, where the humble city of \textit{Whitestone} stands as the source of rule and law in this land. The people here have been hardened through strife and exalted through freedom, enough to know honor and appreciate loyalty like few other civilizations do.

\begin{DndReadAloud}{}

\begin{itemize}
\item{Majority Faiths:}
The Dawnfather

\item{Minority Faiths:}
The Wildmother, The Matron of Ravens, The Chained Oblivion

\item{Imports:}
Industrial and precious metals, cloth, livestock, grain

\item{Exports:}
Granite, limestone, \textit{whitestone}

\end{itemize}

\end{DndReadAloud}

\subsection{Mooren River Run}
While many smaller streams and rivers run through the \textit{Parchwood Timberland}, gathering rain and snow runoff from the nearby mountains, they nearly all empty into the massive Mooren River Run. A powerful river that tears through the heavy forests of the \textit{Alabaster Sierras}, the Mooren River Run flows fiercely southward, out of the woods and into \textit{Mooren Lake}, near the \textit{Lucidian Coast}. Freshwater fish are in ready supply along the river's path, and \textit{Whitestone} maintains a well-guarded dock for their fishermen to feed their people.

The small town surrounding the dock is simply called Mooren, and seasonally is home to some five thousand people of various ancestries in the spring and summer. When conditions grow too cold, mages from \textit{Whitestone} are called to place \textit{arcane lock} upon all establishments and to place the empty town in lockdown until next season.

Due to the plentiful nature of the river, many denizens of the \textit{Parchwood} wander its banks to reap its bounty, or to seek bounty from those who frequent it.

\subsubsection{Mooren River Run Adventures}
Game Masters who set their adventure in \textit{Mooren River Run} can use these plot hooks for inspiration.

\subparagraph{Bear of the White River (low-level)}
While traveling near \textit{Whitestone} or through the \textit{Parchwood Timberland}, the characters find the corpse of half-orc forester; a successful DC 10 Intelligence (Nature) check reveals that he was savaged by a large bear. This half-orc's forester clothing marks her as a cleric of the The Stormlord and a member of the \textit{Rivermaw herd}, a long way from the \textit{Dividing Plains}. Investigating the scene shows that bear tracks in the mud lead westward, toward the \textit{Mooren River Run}. These tracks, when followed, soon fade into the bare footprints of a humanoid—the killer was a \textbf{werebear}, but who could it be?

\subparagraph{Off-Season Suspicion (mid-level)}
The cold nights of autumn and winter force even the most stalwart fisherfolk away from Mooren. Some Mooren residents who have returned to \textit{Whitestone} for the off-season have seen lights in the abandoned town, and are spreading rumors that some of the mages contracted to "lock up" after everyone has left town are actually agents of the \textit{Myriad}. The truth is that these mages are mercenaries hired by the \textit{League of Miracles} to secretly conduct experiments on \textbf{cold snap spirit} to transform them into loyal elemental warriors.

\subsection{Parchwood Timberlands}
Dark, dense, and often assailed by snow and cold winds, the Parchwood Timberlands consume most of the valleys within the \textit{Alabaster Sierras} with pine trees and heavily overgrown pathways. The grim and forbidding interior of the Parchwood is home to many grim inhabitants, from the spectral remains of unlucky travelers past to wandering blights spawned by corrupt trees that consume the life around them. Howls can be heard in the darkest of nights, giving rumor to \textbf{werewolf} hunting parties that stalk the wood. A handful of nearby farmers swear they've encountered a community of territorial \textit{centaurs} in the southern region of the Parchwood, though such tales are largely met with derision. It's common knowledge that if you intend to travel through the thick fog of the Parchwood Timberlands, you stick to the roads, you bring some muscle, and you pray to whatever gods bring you comfort.

\subsubsection{Parchwood Timberlands Adventures}
Game Masters who set their adventure in the \textit{Parchwood Timberlands} can use these plot hooks for inspiration.

\subparagraph{Headless Horse-Man (mid-level)}
The \textit{centaurs} of the \textit{Parchwood} are no myth, and many are succumbing to a curse of undeath. One among them, a headless stallion-lord named Ichabarr, now leads the Herd of the Damned across the valleys of the \textit{Parchwood}. Several small settlements outside \textit{Whitestone} have told the \textit{Chamber} that these skeletal monsters have been spiriting away eligible young bachelors, and soon no men will be left in the Timberlands. These marauders are \textbf{centaur skeleton}. Their leader Ichabarr uses the same statistics, and can also use the optional Spirit Away feature.

\subparagraph{They Seek Only Rest (high-level)}
Those who travel the \textit{Parchwood} sometimes hear thunder roll through the valley, accompanied by the cracking of timbers, but no storm follows. Over twenty years ago, the \textit{Briarwoods} slaughtered and reanimated an entire clan of giants native to the \textit{Parchwood}. Some of these zombie giants survived the fall of the \textit{Briarwoods}. Most of these giants have been destroyed by the \textit{Grey Hunters}, but a few still remain. However, their mindless actions are not without purpose: each giant is tormented by being denied the halls of their ancestors, and quietly weeps for their lost afterlife as they roam.

Zombie giants can be represented by changing any giant's type to undead, decreasing its AC to 8, giving it resistance to bludgeoning damage from nonmagical attacks, and giving it immunity to poison damage and the poisoned condition. This does not affect its challenge rating.

\subsection{Salted Bluffs}
The tumultuous waters of the \textit{Lucidian Ocean} that encompass the northeast edge of Tal'Dorei end at the towering seaside cliffs called the Salted Bluffs. These cliffs mark the coastal boundary of the \textit{Alabaster Sierras} and are home to seemingly infinite cliffside caverns that overlook the ocean, sometimes hundreds of feet up. The waters along the base are rocky and dangerous to traverse by boat, and only the most skilled of captains even consider sailing through these treacherous waters. These factors make the Salted Bluffs an ideal hideaway for pirates, who have taken to these bluffs to hide away their riches or build a base of operations.

However, these outlaws must also deal with the colonies of \textbf{harpy} native to the region. They can never be too sure if a cave they're holing up in for the evening is abandoned—or if it's home to a beast of alluring song and shredding talons. Even empty caves house the remains of these \textbf{harpy} brutal kills, for they sweep up unsuspecting creatures and dash them against the cliffs and rocks, recovering the broken corpses to feed upon safely.

\subsubsection{Salted Bluffs Adventures}
Game Masters who set their adventure on the Salted Bluffs can use these plot hooks for inspiration.

\subparagraph{Crow's Nest (low-level)}
The Ashwing \textbf{harpy} tribe that roost among the Salted Bluffs have recently commandeered the galley owned by the group of penitent corsairs and started raiding coastal villages. They are led by Captain Silvercomb, an arrogant, dashing, and well-spoken \textbf{harpy matriarch}. The other \textbf{harpy} are just pillaging for fun, but Silvercomb seems to be searching for something, and she always wears a pendant with a black feather around her neck. What could it be?

\subparagraph{Start Believing in Ghost Stories (low-level)}
Pirates are one thing. Undead pirates are another. A group of \textbf{swashbuckler} have willingly turned themselves in to the \textit{Chamber of Whitestone}, begging for absolution from the The Dawnfather for their crimes. They refuse to return to the cove where they kept their stolen treasure since it's guarded by the \textbf{zombie} and \textbf{skeleton} of their risen comrades—as well as the \textbf{wight} that killed them. Perhaps a group of adventurers can save the day?

\subsection{Shearing Channel}
Long ago, a land bridge once connected the landmasses of Tal'Dorei and \textit{Wildemount}. However, the catastrophic forces of the \textit{Calamity} blasted it apart, leaving in its place a deathly strait of jagged rock and gale-force winds. Numerous ships have been ruined against the tides of the channel, and sea creatures that thrive on wreckage have taken up residence beneath the waters. It's said that no honest sailor knows a safe passage through the channel, but rumors abound of hidden coves and unknown routes that the Revelry pirates of \textit{Wildemount's} Menagerie Coast use to safely traverse these deadly waters.

\begin{DndReadAloud}{}
\DndQuote%
{}
{''Amidst fluted white bones, serenaded by siren's call.''}
{'''Twas there did I bury my treasures, one and all.''}
{''But deaf and blind must the finder be,''}
{''Lest a briny death be that which ye truly seek.''}
{}
\end{DndReadAloud}

\subsubsection{Shearing Channel Adventures}
Game Masters who set their adventure in the \textit{Shearing Channel} can use this plot hook for inspiration.

\subparagraph{Scion of Oblivion (epic-level)}
The The Chained Oblivion has no physical form, but some residue of its power was cast into the ocean when the The Dawnfather pursued it from the \textit{Alabaster Sierras} to \textit{Gatshadow}. This seed of evil has had eons to gestate undisturbed in the tempestuous channel—and now it rises. A \textbf{kraken}, mutated by the The Chained Oblivion aberrant power, explodes from the water. Its mere presence causes the flesh and minds of sailors to fuse to their ships, sea creatures to hideously mutate into humanoid abominations, and the clouds above to rain thick, ichorous black droplets. A storm now lingers permanently over the \textit{Shearing Channel}, and forces from \textit{Whitestone} and the Menagerie Coast of western \textit{Wildemount} relay frantic messages, seeking to somehow halt this evil. It seems to be building an army here for some larger attack. What is it planning?

\subsection{Whitestone}

\begin{description}
\item[Small City:]
Population 17,710 (78\% \textit{humans}, 7\% \textit{halflings}, 6\% \textit{dwarves}, 4\% \textit{tieflings}, 5\% \textit{other races})

\end{description}

Over two hundred years ago, the de Rolo family helmed an expedition from \textit{Wildemount} to the shores of Tal'Dorei. When their vessel wrecked upon the stones of the \textit{Shearing Channel}, the de Rolos and their remaining crew survived for weeks in the \textit{Alabaster Sierras}, despite the tempestuous weather and slavering beasts. After a time, the survivors found a clearing containing a single tree that glowed with the radiance of the sun itself. They called it the \textit{Sun Tree}—a literal blessing of the The Dawnfather himself—and built their home around it.

Since then, the de Rolos' pioneering settlement has grown into a robust, if isolated, city. The sights of the city range from the gothic majesty of \textit{Castle Whitestone} to the radiant glory of the \textit{Sun Tree}, and from the quaint townhouses of the city center to the humble homesteads of its outlying townships. Because Whitestone is limited in fertile farmland, locally grown produce is scarce, leading to the necessity to develop a good trade relationship with the rural communities of the \textit{Turst Fields}.

Though the region's scarcity of vital resources made for a harrowing first few years, the de Rolo family stumbled upon a miraculous mineral they called \textit{whitestone}. This geological marvel both gave their city its name and quickly became its most important and most sought-after export. Before long, the de Rolos' sale of \textit{whitestone} allowed them to quickly amass a significant fortune and construct a princely keep within their city.

The city is constantly improved (and unsettled) by the brilliant tinkering of \textbf{Percival de Rolo}. His greatest creation, the \textit{Heart of Whitestone}, overlooks the city. Most prominent these days, however, is a rudimentary electric light grid. Using acid pits in the relatively unexplored caverns beneath the city as a sort of battery, six dozen electric lamps throughout the city are given just enough power to gently illuminate the streets, the watch posts on the castle, and the city walls with a soft, greenish glow all through the night.

\subsubsection{Scars of the Briarwoods}
Lady Delilah and Lord Sylas Briarwood, former cult leaders of The Whispered One, were ousted from \textit{Whitestone} over twenty years ago. Though they are no more, the evil they committed still lingers. The nearly six years they ruled \textit{Whitestone} sapped life and hope from its people. Though homes have been rebuilt and the Briarwoods' undead monstrosities have been eradicated, the trauma of the vampires' reign of terror will linger for generations. Some people of \textit{Whitestone} were so deeply affected that they turned to necromancy and vampirism themselves, and have fled into the \textit{Alabaster Sierras} to learn the secrets of the Briarwoods' foul magic—and to emulate their takeover of \textit{Whitestone}. The enigmatic \textit{Grey Hunters}, the secretive elite warriors of the \textit{Chamber of Whitestone}, have worked tirelessly to pull out the \textit{Remnants} by their roots, but permanently eliminating them seems an impossible task.

\subsubsection{An Enchanting Export}
One of the climactic battles of the \textit{Calamity} was a clash between the stalwart The Dawnfather and the engine of callous destruction known as the The Chained Oblivion. This battle came to a dramatic and explosive conclusion in the \textit{Alabaster Sierras}. The unstable arcane forces unleashed in their battle created the valley that now houses the \textit{Parchwood}, thrusting the surrounding mountains upward and infusing them with incredible amounts of \textit{residuum}—the physical residue of arcane energy.

The sheer amount of leftover arcane power within the \textit{whitestone} of the mountains is incredibly receptive to enchantment. Magic items that incorporate at least an ounce of \textit{whitestone} into their construction require only one-quarter the creation time of other magic items. Some alchemists in \textit{Emon} have discovered that dissolving \textit{whitestone} with specific acids can leave behind pure \textit{residuum}, which can substitute for expensive spell components when such materials are not readily available.

If a spell has a component cost that consumes the component, its caster can substitute the required component for an amount of \textit{residuum} of an equal value. One pound of pure \textit{residuum} is worth approximately 500 gp in most markets, though it becomes extremely expensive and difficult to find outside of \textit{Emon} or the city of \textit{Whitestone} itself.

\textit{Residuum} can also be heated and blown into glass, forming a shimmering, greenish surface. Residuum slates or orbs can be used in stationary enchanting tables. Archmages willing to pay exorbitant rates—easily 20,000 gp or more—may acquire an enchanting slate that can be used to expedite the creation of all magic items. Using such a table allows anyone crafting a magic item to do so in one-quarter of the usual time.

\subsubsection{Government}
For many years, the de Rolo family ruled \textit{Whitestone} with dignity and justice. Their noble rule was broken by Delilah Briarwood, an exiled necromancer of the Cerberus Assembly from \textit{Wildemount}, and her \textbf{vampire} husband, Sylas. Together, they slaughtered the family, took over \textit{Whitestone}, and plunged the city into despair.

Nearly six years later, Vox Machina returned with \textbf{Percival de Rolo}, the only known de Rolo heir, and discovered that his sister Cassandra also still lived. The heroes defeated the \textit{Briarwoods} and expunged their evil from \textit{Whitestone}, and the de Rolo siblings founded a new system of governance: the \textit{Chamber of Whitestone}.

This council consists of members appointed by the Guardian of Woven Stone, who is to remain a member of the de Rolo bloodline as long as it exists, or until the rest of the Chamber unanimously agrees there are no suitable de Rolos to accept the position.

\subsubsection{Defense}
\textit{Whitestone's} first line of defense is made up of the obelisk wards. Hastily constructed during the time of the \textit{Chroma Conclave} to avoid detection by their roving \textbf{wyvern} armies, five \textit{whitestone} obelisks mark the perimeter of the city and castle grounds. These obelisks form a dome of illusory magic when powered by an arcane spellcaster, hiding the city under the guise of continued, undeveloped \textit{Parchwood}, and shielding everything within from arcane detection. The obelisks haven't been used since the attack of the \textit{Chroma Conclave} and are beginning to be covered by graffiti and overgrowth.

For each spell slot that is syphoned into an obelisk, two hours of illusion is powered per spell level sacrificed, and the cover of the illusion barrier is incomplete unless all five obelisks are powered simultaneously. This siphoning is a taxing experience, and for every 30 spell levels sacrificed within a 24-hour period, the caster gains 1 level of exhaustion.

In addition to this rarely used magical ward, three groups of defensive personnel protect \textit{Whitestone} from both within and without.

\subparagraph{Pale Guard}
The people of \textit{Whitestone} rely upon the Pale Guard to resolve any (admittedly rare) crimes and disturbances of the peace. The Pale Guard are a force of sentinels who answer to the Pale Lord of Wardship. Since criminal activity is fairly low these days, most of the Pale Guard are given other civic tasks to help \textit{Whitestone's} people.

These tasks range from helping newcomers navigate the city streets to escorting visiting dignitaries to diplomatic functions. While performing a duty that requires it, members of the Pale Guard are permitted to carry simple, three-barrel pepperbox revolvers.

\subparagraph{Whitestone Rifle Corps}
The defense of the city from outside invaders and threats from the darkened boughs of the \textit{Parchwood Timberlands} is attended to by the \textit{Whitestone Rifle Corps}, who are trained in the deadly long-range rifles created by \textbf{Percival de Rolo}. They are only permitted to wield the weapons while patrolling the city walls or castle battlements—except in times of emergency, or when granted special permission by the \textit{Chamber of Whitestone}.

\subparagraph{Grey Hunters}
An elite force called the Grey Hunters makes up the most secret of \textit{Whitestone's} defenses. Trained in the arts of hand-to-hand combat, marksmanship, espionage, and etiquette, the Hunters have a reach that extends across all Tal'Dorei, but they exist only as a rumor in the eyes of the public—and they answer only to the \textit{Grand Mistress of the Grey Hunt}.

The firearms wielded by members of the Grey Hunters are custom-made and named by their owner. Their rifles' stocks are made of \textit{Parchwood} inlaid with \textit{whitestone} and enchanted against wear with \textit{residuum}. The lock mechanism bears the \textit{Whitestone} crest. Hunters are trained to meticulously care for their weapons, and many learn enough to make personal modifications, or even build one from scratch. A Hunter must never lose their weapon. If they cannot prove it was destroyed, they aren't permitted to return to \textit{Whitestone} until it is recovered—although they may request help from other Hunters in retrieving it. If they come across a weapon of \textit{Whitestone} in the wrong hands, they are honor-bound to return it to the city.

This covert organization is headquartered in the ruins of a failed city expansion deep beneath the eastern \textit{Parchwood}. This base is used for training, for weapon storage, and as a repository for their mission logs, so that a complete record of every Hunter and their weapons can forever be kept. As with many of his inventions, \textbf{Percival de Rolo} fears that the Hunters are a powerful—and even necessary—tool in the hands of the righteous, but a terrible weapon, should they be perverted to suit evil.

\subsubsection{Crime}
The community of \textit{Whitestone} has endured much hardship, from its shipwrecked beginnings to the bloody reign of the \textit{Briarwoods}. Its hardy people are generally more concerned with personal duty, small comforts, and the creation of new art and technologies than with the pursuit of easy profit. The \textit{Myriad} has tried several times to gain a foothold in \textit{Whitestone}, but they've been foiled time and again by the consistent counterespionage efforts of the \textit{Grey Hunters}. Nevertheless, they continue to try, and their presence has inspired petty theft throughout the city.

The altruistic Lord \textbf{Percival de Rolo} and Lady \textbf{Vex'ahlia} of \textit{Whitestone} prefer to address crime at the source—poverty and suffering—rather than punishing criminals. Likewise, most petty criminals are seen in \textit{Whitestone} as worthy of rehabilitation, not punishment.

\subsubsection{Geography}
\textit{Whitestone} was originally constructed around the newly discovered \textit{Sun Tree}, and has slowly expanded outward ever since. \textit{Castle Whitestone} was built atop a small hill overlooking the city to the north. Over the years, farmers and homesteaders have cleared more and more of the \textit{Parchwood} to the south of the city. Even today, as \textit{Whitestone's} population has swelled, the \textit{Parchwood} remains defiant, dark, and threatening. It is in little danger of being overrun by mortals, and its eerie creatures strike back whenever the people of \textit{Whitestone} press against it too greedily. The \textit{Grey Hunt} was formed to keep the deadly denizens of the woods at bay, though now their purview reaches farther than just the \textit{Parchwood}.

Dawnfather Square is the central district of \textit{Whitestone}. The square surrounds a small hill, and the golden glow of the revered \textit{Sun Tree} that sits atop it can be seen throughout the city. Many of \textit{Whitestone's} original settlers already worshiped the The Dawnfather, and the discovery of the \textit{Sun Tree} was seen as a sign. Today, Dawnfather Square is filled with quaint, bustling shops and small town-homes.

\subsubsection{Points of Interest}
The City of \textit{Whitestone} has a number of famous landmarks. They are keyed to the map of \textit{Whitestone}.

\subparagraph{1. Castle Whitestone}
This gothic fortress was erected not long after the city's founding, and has long been the seat of power for the sovereign de Rolo family. \textbf{Percival de Rolo} and Cassandra de Rolo have opted to end their absolute power over the city by sharing it with the \textit{Chamber of Whitestone}—on which they each hold a seat.

The castle's interior has a certain stark regality to it. While the hunting trophies and tapestries that decorate its walls are clear signs of the de Rolos' ancestral wealth, there is a solemn, beautiful hollowness to its corridors. Cassandra de Rolo says that something of her brother's melancholy seeped into its \textit{whitestone} walls, and the castle will always carry the weight of his sorrow—and of the \textit{Briarwoods'} evil.

The castle itself is generally off-limits to the public, but its grounds are more accessible. The \textit{Whitestone Rifle Corps} train in a firing yard behind the castle, and many people of the town flock to watch—only to be politely rebuffed by the \textit{Pale Guard}. Open to all, however, are the gardens—except for \textbf{Percival de Rolo} latest project, the Widow's Garden. This walled garden was named for Lady Melanie de Rolo, who originally cultivated its deadly blooms to slowly murder her husband. The Widow's Garden used to be a macabre menagerie of poisonous plants from across \textit{Wildemount}. It fell fallow during the reign of the \textit{Briarwoods}, but \textbf{Percival de Rolo} has resurrected it with \textbf{Keyleth, Voice of the Tempest} horticultural aid.

\subparagraph{2. The Sun Tree}
After the The Chained Oblivion nearly consumed the Knowing Mistress during the \textit{Calamity}, a furious and vengeful The Dawnfather pursued the fleeing The Chained Oblivion across the mountaintops and over the sea, to the peak of \textit{Gatshadow}, where the mad god was eventually defeated and imprisoned once again. The land wept with pain as the two wrathful gods laid waste to Tal'Dorei.

After his victory, the The Dawnfather returned to \textit{Whitestone} and placed a single seed into the ground where the The Knowing Mentor was nearly destroyed. The god wept, and from the seed grew a symbol of protection—a tree with radiant orange foliage and a golden glow that surrounds it when in bloom. The Sun Tree became a beacon of light within the dark forest that enclosed it. While the extent of the tree's divine power remains a mystery, it stands as an emblem of unity and hope, both within the city of \textit{Whitestone}, and upon its official crest.

\subparagraph{3. Temples and the Greyfield}
A patch of land outside the eastern wall of the city sees generations of \textit{Whitestone's} dead interred within a graveyard known as the Greyfield. Custody of this cemetery is shared by the Temple of the Dawn and the Altar of the Raven, sanctuaries maintained by the clergy of the The Dawnfather and the The Matron of Ravens, respectively. Since most people in \textit{Whitestone} consider the The Dawnfather their patron deity, the Temple of the Dawn is the largest of these shrines. The smaller Altar of the Raven is right next door, and most well-to-do devotees of the The Dawnfather look upon the dour, dark-robed clerics of the The Matron of Ravens with a mixture of awe and unease.

The temples' grounds also house a number of mausoleums containing the remains of long-forgotten noble families, though the last generation of lost de Rolos were not recovered following the rule of the \textit{Briarwoods}. Traces of the \textit{Briarwoods'} defiling magic still linger within the Greyfield, and the \textit{Pale Guard} are posted here at night to quell the restless dead that occasionally claw their way to the surface.

\subparagraph{4. The Heart of Whitestone}
A majestic clock tower called the Heart of \textit{Whitestone} stands proud and tall over the city. It is \textbf{Percival de Rolo} favorite invention, one he dedicated much of the past two decades to constructing to honor his loyal friends, their sacrifices, and the trials they endured together. The area around it is known as Forlington Square, named for one of the first families to settle in \textit{Whitestone}. This bustling courtyard is lined with shops and street performers that attract locals and tourists alike.

The clock tower's construction was completed five years ago, but \textbf{Percival de Rolo} continues to tinker with it, adding both functionality and self-indulgent embellishments. The tower has individual movements for \textit{most major holidays}, including depictions of famous moments and persons from \textit{Whitestone} history. The clock tracks the movements of the moons, stars, and spheres. The mechanics of the clock are only truly understood by \textbf{Percival de Rolo} himself, though small clubs of scholars and esotericists take great joy in unraveling its mysteries. Further, it has long been rumored that \textbf{Percival de Rolo} secreted a vault within the tower's innards to safeguard his most dangerous inventions.

\subparagraph{5. Snowdrop Memorial Trail}
In the east of the city, winding between the \textit{Grey Hunt Estate} and the temple to the The Matron of Ravens, is a private trail through the woods. Only locals know of this trail, and even among them, few know that the path was made by \textbf{Keyleth, Voice of the Tempest} of \textit{Zephrah} on one of her many visits to the city. Countless snowdrop flowers line this secluded track, and they are tended to by Lady \textbf{Vex'ahlia} when she has time away from her duties to both the \textit{Chamber of Whitestone} and the \textit{Tal'Dorei Council}. The significance of this memorial garden isn't listed on any plaque, but most in \textit{Whitestone} have an inkling that it's dedicated to a beloved \textbf{champion of ravens}.

\subparagraph{6. Grey Hunt Manor}
The estate of the \textit{Grand Mistress of the Grey Hunt} is rarely used by the current holder of that title these days, as all her nights not spent away on \textit{Tal'Dorei Council} business are passed with her family in \textit{Castle Whitestone}. Members of the \textit{Grey Hunt} still use it as a formal meeting place—though meetings of true import take place in their secret lair in the eastern ruins (11).

\subparagraph{7. House of the Lawbearer}
This temple is both a place of worship and a stately courthouse. All who enter are sworn to tell the truth by the vigilant clerics of the The Lawbearer—for the entire building is under the protection of a vast, permanent \textit{zone of truth} spell. This protection is fallible, but it sends a clear message to all who enter.

\subparagraph{8. Slayer's Cake Bakery}
Founded by members of Vox Machina, the Slayer's Cake has quickly become one of the most beloved bakeries in Tal'Dorei, and it has a satellite shop in \textit{Emon}. It ships bread and pastries to Nicodranas in \textit{Wildemount} via \textit{teleportation circle} weekly, and the bakery's owners are thinking of expanding.

\subparagraph{9. Lord Trinket's Public Park}
A section of the \textit{Parchwood} to the west of the city was cleared in the aftermath of The Whispered One defeat in anticipation of a major expansion. However, more of the \textit{Emonian} refugees left the city than anticipated, leaving these plans untenable, but the forest already cleared. At the request of \textbf{Vex'ahlia}, Grand Mistress of the \textit{Grey Hunt}, the space was transformed into a public park that she uses to play with her now-venerable brown bear, \textbf{Trinket}—and a lifelike marble statue of that same bear, clad in magnificent armor, stands in its center.

Portions of this city expansion did eventually come to pass, and now several manor houses have been built near the edge of the park to house visiting dignitaries.

\subparagraph{10. Galdric's Gate}
The new, northwestern gate of \textit{Whitestone} opens onto the Timberland Trail, which grants hunters easy access to the city when they return from expeditions into the \textit{Parchwood}.

\subparagraph{11. Eastern Ruins}
\textit{Whitestone's} first expansion in recent times happened just after the fall of The Whispered One, to accommodate the city's large refugee population. However, a mix of social unrest and monstrous incursions from the \textit{Parchwood} stymied the project, and the bones of the expansion were left to fall into ruin. No attempt has been made to reclaim the ruins, but the \textit{Grey Hunt} have made a secret base of operations in the basement of a magnificent temple not yet dedicated to any god.

\subparagraph{12. Hearth of the Everlight}
Established by \textbf{Pike Trickfoot} in the years following Vox Machina's victory over The Whispered One, this humble temple stands as a bastion of welcoming light and shelter to those seeking solace and direction. Within its stone walls, the faithful of the The Everlight study her philosophy of compassion and fierce protection. Still, the temple is far from tranquil—the congregation also hosts monthly fêtes and revels as a part of their worship in the image of their exuberant founder.

\subsubsection{Whitestone Adventures}
Game Masters who set their adventure in \textit{Whitestone} can use these plot hooks for inspiration.

\subparagraph{Dad Scientist (any level)}
\textbf{Percival de Rolo} doesn't usually involve outsiders in his creative process. However, when his youngest daughter Gwendolyn takes one of his latest creations and disappears into the dungeons of \textit{Castle Whitestone}, he has no choice but to rally all adventurers in \textit{Whitestone} to join him, \textbf{Vex'ahlia}, and the city's elite warriors in scouring the dungeons for his missing girl.

\subparagraph{Matter of Some Importance (low-level)}
The \textit{Castle Whitestone} library recently acquired a new book containing unprecedented knowledge about the Outer Planes, but opening the tome also opens an infernal portal that summons a steady stream of \textbf{imp}. The characters are hired by Head Librarian JB Trickfoot to keep the \textbf{imp} out of her hair—and out of the library, please—while she scours the tome for information.

\subparagraph{Grey Hunt (high-level)}
The characters have been recommended to the \textit{Grand Mistress of the Grey Hunt} as formidable adventurers by an unknown ally—or adversary. \textbf{Vex'ahlia} de Rolo considers them cautiously, but offers them the chance to prove themselves worthy of the title of \textit{Grey Hunter}. She sets them a task to be completed within the next three months: track down a legendary \textbf{ember roc} in the \textit{Cliffkeep Mountains} and claim its head as a trophy. Of course, she knows that this quest is more than just a hunt; the question is, will the characters realize it takes more than just skill to join the hunters?

\section{Dividing Plains}
Tal'Dorei's vast heartland, framed by the \textit{Stormcrest Mountains} in the southwest, the \textit{Cliffkeeps} in the north, and the \textit{Summit Peaks} in the south, contains miles and miles of rolling hills, tall-grass prairies, sky-blue rivers, and fertile farmlands. Within these plains, the hardworking people of the realm tend to their fields and businesses while adventurers and village militias battle back the wild beasts and wandering bandits of the region. The safest way to avoid being beset by monsters is to travel the \textit{Silvercut Roadway}, which slices through the endless plains, connects the coasts, and shepherds caravans between them. Of course, the roads are coursing with highwaymen and bandits, so which route is safer is entirely up to the traveler to decide. Cults of the \textit{Betrayer Gods} also prey upon the people of the plains. The most fearsome of these death cults are the \textit{Ravagers}, a bloody band of warriors exiled from cities and nomadic clans alike. All people of the plains fear the \textit{Ravagers}, and even the smallest villages train people old enough to hold a spear how to defend their homes. Beyond the grassy plains, \textbf{owlbear} and other monstrosities stalk the forests and foothills that dot the plains and border it on many sides. Packs of \textit{centaurs} also gallop across the open plains, delivering the The Arch Heart justice as they see fit.

\begin{DndReadAloud}{}

\begin{itemize}
\item{Majority Faiths:}
The Wildmother, The Lawbearer, The Platinum Dragon

\item{Minority Faiths:}
The Stormlord, The Everlight, The Ruiner, The Strife Emperor, The Chained Oblivion

\item{Imports:}
\textbf{Fish}, \textit{gold}, diamonds, spell components

\item{Exports:}
Lumber, grain, produce, cobalt, \textit{iron}, \textit{silver}, livestock

\end{itemize}

\end{DndReadAloud}

\subsection{The Bramblewood}
Surrounding the base of the southernmost stretch of the \textit{Cliffkeep Mountains} is the Bramblewood, a dark and tangled forest named for its jagged, thorny-barked trees. The Bramblewood's impenetrable canopy and thick, gnarled roots give it an ominous appearance, but its eerie boughs are populated with countless deadly beasts. It's a grim warning to all who would dare approach \textit{Gatshadow}, the mountain of despair: come no closer if you value your life.

Yet despite this, without the abundant game and plentiful timber of the Bramblewood, the nearby city of \textit{Westruun} would likely not exist. A trade city with nothing to trade rarely lasts long, and the Bramblewood ensures that local carpenters, artisans, and tradesfolk always have plentiful resources. For those brave enough to venture into the Bramblewood, a pair of well-tread paths have resisted its overgrowth; one leads to the iron-rich Murdoon Mines at the base of the \textit{Cliffkeeps}, while the other winds upward to \textit{Gatshadow's} storm-wreathed peak.

\subsubsection{Bramblewood Adventures}
Game Masters who set their adventure in the \textit{Bramblewood} can use these plot hooks for inspiration.

\subparagraph{Haunted Road (mid-level)}
Local folklore speaks of a spectral young woman sometimes appearing along the road through the forest, sobbing and leading men deeper into the \textit{Bramblewood} to never be seen again. While these tales are generally considered childish ghost stories, a mining crew was recently found dead near the fabled spot, unwounded and with looks of terror stretched across their lifeless faces. A \textbf{banshee}, aided by the \textbf{specter} of those she has slain, hungers for the sweet taste of another adventurer's dying breaths.

\subparagraph{Blood of the Ruiner (high-level)}
The Shields, \textit{Westruun's} beleaguered soldiery, have tracked a cell of \textit{Ravagers} to a base in the \textit{Bramblewood}. They have conducted a foul ritual to the The Ruiner to turn their cult master, an elf named Kalydria Darkeye, into a \textbf{Ravager slaughter lord}. They have also transformed a number of the \textbf{giant spider} endemic to the \textit{Bramblewood} into hideous \textbf{demonfeed spider} through the same bloody ritual. The \textit{Ravagers} have created a makeshift fortification, and the Shields need to know what their next move is—an infiltration mission only powerful adventurers are capable of completing.

\subsection{Foramere Basin}
This great lake was once the ice fortress of \textit{Errevon the Rimelord}, the elemental behemoth that nearly toppled the fledgling realm of Tal'Dorei just two years after its founding. It was here that the combined forces of \textit{Kraghammer}, \textit{Syngorn}, Tal'Dorei, and \textit{the Ashari} hurled the Rimelord back into the Elemental Plane of Ice and brought an end to the \textit{Icelost Years}. The massive lake left behind by \textit{Errevon's'} melted citadel is now the largest source of fresh water in central Tal'Dorei.

Foramere Basin's shores are peppered with tiny huts and fishing communities that trade with \textit{Kymal} and \textit{Westruun}. The great lake is fed by three rivers that give life to the \textit{Dividing Plains}: the Byhills, the Tundrun, and the Wildpath. After especially hard winters, these rivers run hard and fast, and the overflows spill into the Foramere Waterway, which flows into Owlset Bay, near \textit{Stilben}.

Some of these fishing villages are home to the descendants of the fighters that warred against the Rimelord's tyranny. Even in times when \textit{Syngorn} and \textit{Kraghammer} were bitter rivals, the \textit{elves} and \textit{dwarves} here took no part in their homelands' hatred.

\subsubsection{Foramere Basin Adventures}
Game Masters who set their adventure in \textit{Foramere Basin} can use this plot hook for inspiration.

\subparagraph{Nightmare Before Winter's Crest (mid-level)}
The festival of Winter's Crest celebrates the anniversary of the Rimelord's defeat and the fall of Foramere Citadel. The characters, far from home on the holiday, are not only witness to the fishing communities on the banks of the basin gathering to celebrate Winter's Crest, but to the grim sight of Foramere's waters surging upward from the basin. The water has begun to freeze into towers of ice—someone or something is recreating the domain of \textit{Errevon the Rimelord} and will not stop until all of Tal'Dorei is buried in ice.

\subsection{Ironseat Ridge}
The \textbf{gnoll} and hillsfolk of the \textit{Dividing Plains} say the crooked, angular mountain that juts out of the plains and foothills was once the throne of a titan. From a distance, \textit{Ironseat Ridge} does resemble a crooked chair. Whether or not this legend is true hardly matters; the tales of \textit{Ironseat Ridge} permeate the folklore of the plains.

The mountains surrounding Ironseat itself have been hollowed out over the centuries, and other legends say that the ruins of an abandoned dwarven redoubt called Inverspire can be found by explorers savvy enough to navigate the labyrinthine tunnels—and that it's surely filled with monsters and vast riches. Legends aside, countless abandoned mineshafts do snake through the stone like the veins of a dead titan. When the winds grow strong, you can sometimes hear a terrible, growling moan from the air pressing through the empty channels. People from \textit{Kymal} rarely venture close enough to the mountains to hear the noise, but sometimes on clear, silent nights, the monstrous sound echoes through \textit{Kymal's} brightly lit streets.

\subsubsection{Ironseat Ridge Adventures}
Game Masters who set their adventure in \textit{Ironseat Ridge} can use these plot hooks for inspiration.

\subparagraph{The Beast Below (any level)}
It's one of those still nights where the howls of the Ironseat echo across the plains. The next time the characters spend a night in town, whether in \textit{Kymal} or any small village of the plains, they learn of disappearances—explorers that went into the caves, never to return—and of the hefty reward their grieving families are offering for their return. What could be making that awful wailing in the caves? \textbf{werewolf}? Kidnappers using the echoing caverns?

\subparagraph{Dead Fires of Inverspire (epic-level)}
The mythic redoubt of Inverspire contains a legendary creation of the The All-Hammer: a \textbf{forge guardian}. It was placed here by the The All-Hammer during the \textit{Calamity} to watch over a forge with the power to create arms and armor fit only for a war between gods and mortals. Over twenty years ago, the elemental power of the \textit{Cinder King} somehow ignited this divine creation's ashen heart and filled it with elemental life. It has turned to the forge it was made to guard, and has begun creating beings of steel and flame that will tear Tal'Dorei apart stone by stone in order to find its missing father, the The All-Hammer. Only definitive proof of the The All-Hammer imprisonment beyond the Divine Gate can quench the guardian's anguished flames.

\subsection{Ivyheart Thicket}
Clustered at the western side of the \textit{Ironseat Ridge}, this thicket contains the ruins of an ancient civilization that remains under study by the \textit{Alabaster Lyceum}. Among scattered relics and the rubble of forgotten villages lie references to the \textit{Founding}: the first age of the world. The excavation pushes onward under the watch of the \textit{Lyceum}, with the promise of discovery looming, but very little understanding of what may be found beneath.

The western edges of these ruins are the most dangerous—and the most alluring—part of the Ivyheart Thicket. A vast, well decorated with ivy-crept marble pillars, burns with eternal fire. Explorers in the employ of the \textit{Alabaster Lyceum} say that this place has become a hotbed of activity. \textbf{Azer} have flocked to it from the Elemental Plane of Fire, and have made it their home, forging beautiful bronze weapons, armor, and \textit{objets d'art}. The \textit{Lyceum} reports that elemental activity within the nearby redoubt of Inverspire is somehow related to this elemental incursion—and that it may be connected to the mythic fire giant stronghold of \textit{Vulkanon}.

\subsection{Kymal}

\begin{description}
\item[Small City:]
Population 29,440 (63\% \textit{humans}, 24\% \textit{dwarves}, 13\% \textit{other races})

\end{description}

Glitzy. Gilded. Grimy. The city of Kymal stands as the grumbling ghost of a gold-rush boomtown. These days, it's a haven for gamblers and criminals: a city of gold built atop a mountain of mud.

Over two centuries before, the discovery of rich gold veins within the \textit{Ironseat Ridge} sent a rush of eager folk to its base. Kymal began life as their base camp, and over the next century it grew into a lively, booming city. \textit{Gold} and minerals came out of Kymal, and all the pleasures of home followed. Music, theater, sex, food, liquor, drugs, trinkets—if you could want it, you could find it in Kymal. Anything you could spend money on was there, and people spent money like the world was ending. The rich binged anything they could get their hands on to pretend they were important, and poor folk indulged in whatever comfort they could find from their wretched existence in a town where money meant everything.

In time, of course, the veins ran dry. Ten years after Kymal boomed into being, it went bust. The miners packed up and left. Not many of them had struck it rich; in fact, most folk were worse off than they had been before. In the intervening centuries, Kymal has gone through too many cycles of boom and bust to count. New veins of gold or mithral or adamantine have been found in the mountains, and people rush in to claim it. Or a cunning entrepreneur comes up with a scheme to put Kymal on the map—and it works, for a few years, before the con runs its course and people are right back to where they were before.

Most folk here make a modest living in rural villages surrounding Kymal's city center, or trawling the nearby \textit{Foramere Basin} for fish, and are trying to build interest as a trade center like \textit{Westruun}, though without much luck. Kymal's most recent change of fortune, however, has yet to run dry. Some seventy years ago, an entrepreneur from \textit{Ank'Harel} named Calis Krishtan opened a small casino called "The Maiden's Wish." It started something of a gold rush of its own, one that continues to this day.

The city has grown over the past seventy years, becoming a hub of entertainment and games of chance in Tal'Dorei. So great was Calis's revival of Kymal that the \textit{Tal'Dorei Council} appointed his son (and heir to the Maiden's Wish), Jaktur Krishtan, the Margrave of Kymal.

\subsubsection{Government}
Now entering old age, Margrave Jaktur Krishtan is desperately trying to maintain the boom that his father created, while grooming his son to take over his gambling empire. Krishtan has recently made headlines across Tal'Dorei by relinquishing his position as margrave. It's an open secret that Jaktur Krishtan has never had a head for governance, simply enjoying the lavish, gold-plated life his casinos' profits bring him. He could easily have allowed his advisors to continue running \textit{Kymal} for him while he continued to throw magnificent galas, make nightly appearances at the Maiden's Wish, and put his stamp of approval on glitzy parades down the city's main street.

Investigators have flocked to \textit{Kymal} to uncover the truth of Margrave Krishtan's sudden change of heart. All suspect that the truth is a salacious secret, one which would shake \textit{Kymal} to its foundations. Krishtan himself seems unperturbed. He continues to live large as he always has, his detractors don't disparage him any more than usual, and the public that loves him hasn't changed their tune.

Jaktur's former majordomo, an austere tiefling woman named Khavasta "Prudence" Almiji, has taken up the mantle of Acting Margrave with quiet satisfaction. It is, after all, the title that she should have had for the past thirty years, and she hopes that the \textit{Tal'Dorei Council} will see fit to appoint her margrave permanently in very short order.

\subsubsection{Defense}
The acting margrave commands a force of about one thousand armed guards called the Public Defense. Despite the name, they are most often seen stationed around the casinos owned by ex-Margrave Krishtan, or working as bodyguards for \textit{Kymal's} politicians and wealthy cardsharps. Gold can buy anything in \textit{Kymal}, even the personal protection of the Public Defense.

\subsubsection{Crime}
\textit{Kymal} is no longer the den of crime it once was, but there is still scum beneath its glitzy and imperfect façade. Small-time cutpurses lurk in the shadows, waiting to make off with the small fortunes won by the casinos' luckiest patrons. The war between the \textit{Clasp} and the \textit{Myriad} smolders within \textit{Kymal} as well, but Acting Margrave Almiji has played her cards carefully and intelligently. Through an intricate web of bribes, promises, and threats, she has managed to suppress even the faintest trace of gang violence within her city.

Of course, crime still occurs in \textit{Kymal}. It's just not public and bloody. The \textit{Clasp} and \textit{Myriad} both own casinos in town—named the Wishing Well and the Dragon's Hoard, respectively. Not many know of these casinos' less-than-savory proprietorship. Even those who do fail to realize just how deeply rigged its tables are—or how many pickpockets and swindlers hide in plain sight in their gilded halls.

Moreover, the \textit{Clasp} and \textit{Myriad} alike have immense illicit drug operations within \textit{Kymal}. The \textit{Public Defense} is technically bound to arrest anyone caught selling or possessing these substances, but the bulk of these arrests fall on small-time crooks working for these criminal organizations, almost never the high-ranking members.

\subsubsection{Geography}
\textit{Kymal} is built on the rocky foothills of the \textit{Ironseat Ridge}. Exploring beyond the glamorous city leads one to nothing but a desolate expanse of boulders and scrub-brush.

\subsection{Ruins of Torthil}
While traveling south of \textit{Kymal} along the \textit{Wildwood Byway}, just outside of the Ivyheart Thicket, travelers can spot a mile-wide cluster of broken stone, salted earth, and a single obsidian monument to the death of Warren \textit{Drassig}, one of the most infamous and terrible figures in Tal'Dorei history. This scarred patch of land marks the ruins of the village of Torthil, which was the site of a pivotal battle of the \textit{Scattered War}—one in which King Warren \textit{Drassig} was ambushed and slain by a rebel army.

The city was destroyed in an act of vengeance by Warren's eldest son Neminar over two centuries ago. Now, its ruins harbor only memories—and occasionally outlaws and cultists who use the ghost stories that surround it to avoid suspicion. The now long-weathered wreckage and broken foundations of the town harbor only ghosts and unsavory creatures drawn to places of such dark history.

\subsubsection{Ruins of Torthil Adventures}
Game Masters who set their adventure in the Ruins of \textit{Torthil} can use this plot hook for inspiration.

\subparagraph{Forgotten War Machine (high-level)}
The Ruins of \textit{Torthil} hold one lost secret. Buried beneath a landslide in one of the nearby hills is a \textbf{platinum golem}. Any crests that may have indicated who it fought for—Tal'Dorei's rebels, the \textit{elves} of \textit{Syngorn}, or \textit{Drassig's} legions—have been scoured away by time. The characters are drawn to \textit{Torthil} when exiles and criminals living in the ruins flood into \textit{Kymal}, saying that a greedy treasure hunter awakened a metallic monster that laid waste to the already ruined town—and is now making its way steadily toward \textit{Kymal}!

\subsection{Shadebarrow}
No one who has visited the \textit{Shadebarrow} has ever returned alive—but neither do they come back entirely dead. The \textit{Shadebarrow} was once a monolithic henge used as a ritual site for sun-worshiping druids known as the Dawn Circle, but their shrine became their tomb after \textit{Trist Drassig} slew them for aiding \textit{Zan Tal'Dorei} in the \textit{Scattered War}, centuries ago.

After \textit{Trist Drassig's} death, the trade guilds of \textit{Westruun} laid claim to the unowned land, and from there, the Dawn Circle's abandoned and treasure-less burial tunnels were purchased from the unions at great cost by an eccentric \textit{Westruunian} baron named Sevil Howthess. His obsession with its history led him to be interred there upon his passing, the tomb outfitted with protections and, supposedly, the remnants of his fortune.

The exact location of the \textit{Shadebarrow} has been forgotten by all but the Baron Howthess's few surviving grandchildren, and they deny any involvement with the forgotten crypt. The surviving descendants of Sevil Howthess, Camilla Tenver and Buddleia Austan, live in \textit{Westruun}, but have changed their names through marriage. Only searching through the public records in \textit{Westruun's} \textit{Hall of Reason} or seeking lore about the \textit{Shadebarrow} in the \textit{Cobalt Reserve} will turn up their names.

Despite the barrow's location being lost to time, treasure hunters and historians still seek out—and occasionally find—the infamous crypt. If they ever return to civilization, it is as a wailing spirit, cursed to eternally torment whoever they thought of in their final moments, or as a shambling corpse focused only on murdering their loved ones.

\subsubsection{Shadebarrow Adventures}
The \textit{Shadebarrow} is a crypt that you can turn into a dungeon of any size in your campaign. It may have only a single layer filled with \textbf{ghost}, and perhaps a plot-important item. Or it could be a multi-level dungeon with interconnected passages and a mighty reward at the bottom, like a \textit{Vestige of Divergence} guarded by a terrible monster.

Game Masters who set their adventure in the \textit{Shadebarrow} can also use these plot hooks for inspiration.

\subparagraph{Forsaken Monolith (low-level)}
While traveling across the \textit{Dividing Plains}, the characters stumble upon the \textit{Shadebarrow}. The ancient druidic henges stand regally atop the grassy hill, and all seems peaceful—until the \textbf{specter} of dead explorers rise from the ground. They wail an incessant curse as they try to kill or scare intruders away: "Defiler! Howthess!"

\subparagraph{Undeath Unbound (epic-level)}
The Demon Prince of Undeath is one of the foul princes of the Abyss and a mighty servant of the \textit{Betrayer Gods}. The cruelty shown at the Dawn Circle so long ago drew his loathsome gaze, and for centuries he has bided his time, waiting for a mortal to disturb this nest of hate and undeath. Ever since Baron Howthess disturbed the \textit{Shadebarrow's} seal, this prince of evil has rallied an army of undeath in the accursed depths of this crypt. Now, the time has come for him to spill his plague of undeath across the \textit{Dividing Plains}. The undead are spreading to all corners of Tal'Dorei, and the characters must find the \textit{Shadebarrow} and reach its deepest sanctum to halt their spread for good.

\subsection{Silvercut Roadway}
Named for the shimmer of endless fields of windswept, sunlit grasses that accompany the way, the \textit{Silvercut Roadway} connects the heart of Tal'Dorei—\textit{Emon}—to its distant extremities, affording traders and travelers some measure of safety against the wilds. Parts of this ancient road predate the \textit{Tal'Dorei Council}, and it now runs from the \textit{Bladeshimmer Shoreline} to the \textit{Lucidian Coast}, connecting most major hubs of civilization across the country.

\subsubsection{Silvercut Crossroads}
The \textit{Silvercut Roadway} and the \textit{Wildwood Byway} meet in the shadow of the \textit{Ironseat Ridge}. During the rule of \textit{Drassig}, many innocents and war heroes were executed by hanging at a massive gallows that once stood at the crossroads, with bodies hung from signs as a warning to dissenters. The gallows were destroyed by \textit{Zan Tal'Dorei's} rebel army, and the rebels created a stone monument with an inscription to commemorate all who lost their lives to tyranny:

\DndQuote%
{}
{''The standard-bearers of revolution honor the innocent dead. Vengeance will not restore the lives of the lost, nor will it appease their spirits. We vow you will receive justice, for justice will ensure that no more guiltless blood shall be shed.''}
{}
Most travelers choose not to camp too closely to the site, for there are countless urban legends about hauntings that surround this grim monument.

\subsubsection{Silvercut Roadway Adventures}
Game Masters who set their adventure on the \textit{Silvercut Roadway} can use these plot hooks for inspiration.

\subparagraph{The Road Warriors (low-level)}
While traveling across the Silvercut, the characters come across the smoking wreckage of a massive wagon train, and investigation shows that it has been ransacked after a fierce battle. A team of mercenaries called the Road Warriors spring upon the characters, blaming them for the attack. Investigating the wreckage turns up unholy symbols of the The Ruiner, suggesting that the \textit{Ravagers} were responsible for the massacre, allowing the characters to easily clear their names.

But is it really that simple? Further investigation reveals that this convoy had a number of mercenary mages in their midst, and many of them had contracts from the \textit{League of Miracles}.

\subparagraph{Shard of the Shadebarrow (mid-level)}
A \textbf{wraith} calling herself the spirit of Archdruid Avandros has animated an army of \textbf{ghoul} and \textbf{ghast} from the graveyard. The undead have taken over the \textit{Silvercut Crossroads}, and they must be dealt with before a national trade crisis emerges. But what is the spirit of an archdruid of the Dawn Circle doing in this mass grave? Sneaking into the graveyard reveals a human skull with sinister amethyst growths protruding from the bone, and Avandros's spirit is desperate to protect it.

\subsection{Throne of the Arch Heart}
In the foothills of the \textit{Ironseat Ridge} is a small rise known as the Throne of the The Arch Heart: a gathering place for the many nomadic \textit{centaurs} who ride across the \textit{Dividing Plains}. The \textit{centaurs} were first born when the The Arch Heart showed mercy to a tribe of elven horse riders who gave their lives serving their god during the \textit{Calamity}. For centuries, the \textit{centaurs} have all returned to the Throne of the The Arch Heart once per decade to gather in a Herdsmeet to discuss the state of the herds, and to pay homage to their deity on the very hill where the first \textit{centaurs} were created. This event causes great chaos in the eastern plains; the massive migration of the centaur herds disrupts trade and stirs the \textit{Ravagers} to acts of terrible vengeance against the creations of the The Arch Heart—the The Ruiner divine nemesis.

\subsubsection{Throne of the Arch Heart Adventures}
Game Masters who set their adventure in the Throne of the Arch Heart can use this plot hook for inspiration.

\subparagraph{Fist of the Ruiner (mid-level)}
Many \textit{orcs} call this rise the Fist of the The Ruiner, for they too were created here in the image of their god. Legend tells that \textbf{orc} and \textit{centaurs} were both created at this rise in the aftermath of a battle between the The Arch Heart and the The Ruiner. \textit{Orcs} with a deep appreciation for ancient myths from across Tal'Dorei make a pilgrimage at the time of the Herdsmeet. Some are long-time friends of the \textit{centaurs}, and they pay homage to the The Arch Heart together in thanks for spilling the blood that created them. Other \textit{orcs} join with the \textit{Ravagers} to lay siege to the Herdsmeet—and the \textit{Ravagers'} bloodthirsty power is so great that \textbf{orc} and \textit{centaurs} alike have begun hiring humanoid mercenaries to defend the herds and pilgrims as they attempt to worship the The Arch Heart in peace.

\subsection{Torian Forest}
Cradling the southwestern edge of the Cliffkeep Mountain range, the Torian Forest supplies much of \textit{Emon's} timber, as well as \textit{Kraghammer's} firewood. Though its locals proudly proclaim the Torian to be the "most welcoming" and "least haunted" of Tal'Dorei's forests, this is more of a myth than people realize. The forest is welcoming and bright at its edges, with a beautiful wetland at its northern end, but within its deep and unexplored heart are spirits not understood since the \textit{Age of Arcanum}—and a cruel and hungry \textbf{wraithroot tree} guarding a magical artifact likewise not seen since time immemorial.

\subsubsection{Rivermaw Herd}
Born from the wild peoples of the \textit{Dividing Plains}, free folk who shunned the constraints of society, the Rivermaw is the greatest of the \textit{Dividing Plains'} herds. It now boasts a membership of nearly ten thousand, along with a number of satellite herds that owe their allegiance to the Rivermaw.

Though the Rivermaw is now a mighty force for good on the \textit{Dividing Plains}, they had humble beginnings. Just two generations ago, they numbered little more than three hundred \textit{half-giants}, \textit{orcs}, \textit{gnolls}, \textit{humans}, and nomads of other races that made their home in the Torian Forest.

It wasn't until they were conquered by the Herd of Storms, helmed by the tyrannical goliath Kevdak, that a more violent way of life was forced upon them for a number of years. Once Kevdak was slain, the remaining members of both the Rivermaw and the Herd of Storms recanted his brutal ways and become the \textit{Rivermaw herd} known across Tal'Dorei today. It is led by two \textit{goliaths}: Zanroar, son of Kevdak, and Skyhaal, daughter of Zanroar.

\subsection{Turst Fields}

\begin{description}
\item[Village:]
Population 2,420 (40\% \textit{humans}, 38\% \textit{halflings}, 15\% \textit{gnolls}, 8\% \textit{other races})

\end{description}

To the northwest of \textit{Drynna}, not far from the Dawnmist Pines, the farming community of Turst Fields toils over extremely fertile land fed by the \textit{Mooren River Run}. Providing much of Tal'Dorei and \textit{Whitestone's} grown food and produce, the community must import most all other commodities outside of wood, making it a lucrative stop for trade caravans traveling along the Parchwood Way to \textit{Drynna}—and an easy target for swindlers. As such, nearly two hundred Shields from \textit{Westruun} are permanently stationed here to keep the peace and prevent unlawful activity. Nearly as common as the arrival of a trader's wagon is an armored carriage transporting brigands, grifters, and horse thieves back to \textit{Westruun} for trial.

\subsubsection{Gnoll Mercy}
\textit{Turst Fields} has a significant \textbf{gnoll} population. As mutant beings created by demonic rituals, \textbf{gnoll} are shunned as monsters in most humanoid settlements on the \textit{Dividing Plains}. Nevertheless, without the Dustpaw \textbf{gnoll}, Turst would not have survived its first summer. When the village was founded, it was easy prey for the \textit{Ravagers}. Their leader sought blood tribute from the people of Turst, but the villagers refused to offer a sacrifice. This cell of \textit{Ravagers} had recently sought the favor of the Demon Lord of Slaughter, transforming a number of the plains-dwelling hyenas into \textbf{gnoll} in that pursuit.

When the villagers of Turst refused to give up their own, the \textit{Ravagers} unleashed the \textbf{gnoll} upon them. But something unusual happened. After the \textbf{gnoll} first drew blood, one of them began to weep. Then another, until all of the \textbf{gnoll} were howling in sorrow and fury. These beings, though infused with demon ichor, turned against their creators to defend their helpless prey.

To this day, the Dustpaw \textbf{gnoll} have lingered in Turst, slowly becoming trusted members of their community. Their unsettling fangs and canine faces upset many who travel to this village, but any who deride the \textbf{gnoll} do so at their own risk. Not at the risk of mauling from a \textbf{gnoll}—though any humanoid being will fight back if pushed too far—but because hateful agitators will soon find themselves surrounded by the spears of the Turst militia.

\subsubsection{Turst Fields Adventures}
Game Masters who set their adventure in \textit{Turst Fields} can use this plot hook for inspiration.

\subparagraph{Herd Mentality (mid-level)}
A half-elf, half-\textit{dragonblood} \textbf{Ravager slaughter lord} calling herself Grud the Great has taken control of an arm of the \textit{Ravagers}. Her highly organized regiment worships the The Strife Emperor. Now, they march on Turst to take an easy foothold in the northern plains. The army is too large to be stopped by just the characters, but a Dustpaw \textbf{gnoll} warrior named Akaksa knows of a number of nomadic groups wandering nearby that could help defend the town, if properly persuaded. The characters can seek out the Dawnmist Herd, led by a goliath \textbf{Rivermaw stormborn} named Gardessk, and convince him to aid the people of Turst.

\subsection{Westruun}

\begin{description}
\item[City:]
Population 51,980 (61\% \textit{humans}, 10\% \textit{gnomes}, 7\% \textit{tieflings}, 5\% \textit{orcs}, 4\% \textit{half-giants}, 1\% \textit{gnolls}, 12\% \textit{other races})

\end{description}

The city of Westruun is the center of Tal'Dorei, a crossroads of culture along the \textit{Silvercut Roadway} that welcomes adventurers, traders, and vagabonds of all kinds. Despite its name, Westruun is more east than west, and was in fact named after Palest Westruun, the trader and philosopher who founded the city. Westruun suffered under the rule of the goliath tyrant Kevdak's Herd of Storms. Worse, its outlying towns and great stone buildings were ravaged by the acidic breath of Umbrasyl the Hope Devourer, an \textbf{ancient black dragon} of the \textit{Chroma Conclave} who exacted tribute from his lair in the eerie mountain of \textit{Gatshadow}. Despite the damage inflicted by its evil rulers, Westruun was one of the first cities in Tal'Dorei to recover from the \textit{Conclave}. Though its government is deeply in debt to both the \textit{League of Miracles} and a number of adventuring parties, the city has long since reclaimed its place of importance as an economic, academic, and political power within Tal'Dorei.

\subsubsection{Government}
\textit{Westruun} underwent a political revolution in the wake of the \textit{Chroma Conclave's} defeat. The former margrave, Brandon Zimmerset, was a tyrant who claimed absolute power for himself following Kevdak's defeat at the hands of Vox Machina. The title of margrave is a holdover from the authoritarian rule of Warren \textit{Drassig} centuries ago, and those who hold that office are appointed for life by the council. A group of local idealists called the Retrievers banded with adventurers and mercenaries to overthrow Zimmerset a decade after he took power, and instated their leader as the lord mayor instead. Elections for the position of Lord Mayor of \textit{Westruun} are held every five years, with no term limits. All adult citizens of the city who have lived there for at least five years have a vote, though only a fraction of the people actually use it.

The Retrievers continue to be a political force in \textit{Westruun}, and several of \textit{Westruun's} delegates to the \textit{Tal'Dorei Council} are members of this faction. They are campaigning to abolish the office of margrave from all settlements in Tal'Dorei. Though they have broad support in the council, the council doesn't have the authority to unilaterally demand its constituent city-states to change their system of government. The revolution-minded Retrievers are somewhat stifled by the need to compromise with other cities.

Lord Mayor Lysandra Kallos is a middle-aged goliath and former leader of the Retrievers. She has never lost a vote for re-election, and holds executive power over the city's defenses, along with the ability to dissolve any business or guild, and to exile any person or group of up to ten people from the city at her discretion. She rules \textit{Westruun} with the aid of a council of ten advisors.

\subsubsection{Defense}
The lord mayor commands a city guard of about 1,200 trained soldiers known as the Shields of the Plain, though most folk just call them the Shields. \textit{Westruun's} soldiery is spread thin across the walls of the city itself, stationed at the gates of the dozens of outlying settlements and conducting raids against small cells of \textit{Ravagers} across the plains.

Many people in town would like to gather their forces and put an end to the \textit{Ravagers'} death cult once and for all, but the Shields are spread too thin to mount an effective offensive. Even in their current state, dozens—if not hundreds—of mercenaries, bounty hunters, and adventurers pass through \textit{Westruun} every year in search of bloody work that the Shields are too beleaguered to handle themselves.

\subsubsection{Debt}
Before he was ousted, Margrave Brandon Zimmerset authorized hundreds of reconstruction bounties—offers of gold to anyone who could rebuild the immense damage the city suffered during Umbrasyl's attack and the Herd of Storms' occupation of the city. Bounties were placed on the city's siege-breaking stone walls, majestic towers and parapets, and even its iconic post-and-beam houses, many of which were razed completely during the attack.

The margrave expected these bounties to be slowly completed over decades, but they were all snapped up and completed in a few short years by mercenary mages serving a group called the \textit{League of Miracles}. The margrave was saddled with a crushing debt, one that he was both unwilling and unable to pay before being ousted by the Retrievers. As his successor, Lord Mayor Lysandra Kallos inherited Zimmerset's immense debt to the \textit{League}, and is constantly seeking to find ways to alleviate their pressure—especially as the league continues to use the debt as leverage to buy the allegiance of her advisory council.

\subsubsection{Crime}
The \textit{Clasp} has become a major and public force in \textit{Westruun} since the criminal organization's dubiously virtuous makeover following the \textit{Cinder King's} defeat in \textit{Emon}. Many \textit{Clasp} agents enjoy the reputation of folk heroes across the \textit{Dividing Plains}, and make public appearances in taverns and town squares throughout \textit{Westruun}.

Nevertheless, the fact remains that the \textit{Clasp} is still a criminal organization. They have controlled and efficiently run their shadowy trades for decades right under the Shields' noses. Their network of fences runs from coast to coast, \textit{Stilben} to \textit{Emon}. Their hideout changes every few years, when the Shields manage to track them down. Their current hideout is secreted away behind the storefront of the Saddled Plainscow, a popular tavern owned and operated by \textit{Clasp} agents.

Folk in search of real folk heroes often seek the covert agents of the \textit{Golden Grin}, who tend to wait and watch in local taverns. The Grin sometimes gets a bad reputation in \textit{Westruun}, since people need help every day, but the Grinners prefer to wait for just the right moment to reveal themselves and strike against evil.

\subsubsection{Geography}
Miles of farmland and fields surround \textit{Westruun's} towering stone walls—except to the west, where the city is bordered by the \textit{Bramblewood} and dwarfed by the mist-shrouded peak of \textit{Gatshadow}. Nearly two dozen villages and towns surround the city's walls.

The city's interior is divided into six wards:

\subparagraph{Scholar Ward}
The city's small southern ward is replete with major temples to the \textit{Prime Deities} and homes for their clergy. The most important gods of the \textit{Dividing Plains} are the The Lawbearer, the The Wildmother, and the The Platinum Dragon, and it's made obvious by their huge, resplendent sanctuaries in the Temple Ward. The greatest temples here are the First Bastion, the Wild Arbor, and the The Platinum Dragon Rest.

\subparagraph{Opal Ward}
\textit{Westruun's} central district is also its municipal heart. The ward has no residential housing, but is home to the Lord Mayor's Residence, the \textit{Hall of Reason}, and the public square.

\subparagraph{Residential Ward}
The northeastern section of \textit{Westruun} is the city's most densely populated ward, housing over two-thirds of its population. There are no major landmarks in this ward, though every neighborhood in the ward has its own unique culture. Even though tourists might not visit the landmarks here, there are countless local bakeries, shops, and pubs beloved by its residents.

This ward is rife with crime. \textit{Westruun's} poorest citizens are pressed together in tight quarters, sometimes as little as a single city block away from the city's wealthiest manor houses. The \textit{Clasp} and the \textit{Myriad} prey upon the poor and affluent alike, often stoking the tensions affecting the city's downtrodden to create chaos that masks their crimes. The Shields are stretched too thin to properly deal with these agitators. Unlike in the Temple Ward, free expression of religion is not a guaranteed right here, even within the privacy of one's own home. Though the The Lawbearer, the The Wildmother, and the The Platinum Dragon are all benevolent deities, many of their \textit{Westruunian} followers are not particularly tolerant of other faiths.

\subparagraph{Market Ward}
\textit{Westruun's} southwestern ward contains the majority of its businesses, trade stands, and production warehouses. Some of these businesses are historic fixtures of the city, rebuilt by the \textit{League of Miracles} in years past, but many sell overpriced wares to cater to tourists, adventurers, and the city's wealthy elite.

It's currently in fashion for store owners to live in a second story above their shops, though less scrupulous businesspeople have turned their residences into gambling dens, black markets, and illegal brothels. The most significant building in this ward is the Exandrian Exchange, a marble-walled auction house where hundreds of vendors hawk rare and unusual items from across Tal'Dorei. A satellite location of \textit{Gilmore's Glorious Goods}, a famed magic item shop in \textit{Emon}, is a beloved attraction of this district.

\subparagraph{Underwalk Ward}
The city's newest ward is underground. It's an extension of the residential sprawl built from repurposed sections of the sewers, underground military bunkers, and the wine cellars of mansions annihilated by Umbrasyl and the Herd of Storms.

The Underwalk was frantically excavated during the invasion of the \textit{Chroma Conclave} in hopes of keeping the populace hidden. It has been largely abandoned for the past twenty-four years, save for those too poor to return even to the Residential Ward, but an effort has been made in the past two years to create a unique and vibrant ecosystem in the Underwalk. The optimistic entrepreneurs trying to popularize Underwalk housing hire adventurers in great quantities, for every day, new and unsettling rumors of many-headed monstrosities prowling the Underwalk's shadowed roads spread to the surface. The best-known entrance to this tangle of tunnels is the Underwalk Gates, a forty-foot-wide staircase in the Residential Ward that descends into a gaping opening lit only by torchlight. Another entrance known mostly to the \textit{Clasp} is a single unassuming manhole in the Opal Ward, topped with a metal cover emblazoned with the crest of \textit{Westruun}.

\subsubsection{Points of Interest}
The City of \textit{Westruun} has a number of noteworthy landmarks. They are keyed to the map of \textit{Westruun}.

\subparagraph{1. Cobalt Reserve}
This vast repository looks more like a house of worship than a library. Nevertheless, this landmark of the Scholar Ward is Tal'Dorei's largest house of knowledge, and it is maintained and fiercely defended by the monks of the \textit{Library of the Cobalt Soul}.

A character who spends 1 hour in the public section of the Cobalt Reserve can make a DC 15 Intelligence (History) check, discovering any piece of common knowledge in Tal'Dorei. On a failure, this check can be repeated after another hour of study. The GM can decide that the knowledge the character seeks can only be found in the restricted archives, which only members of the \textit{Cobalt Soul}, or characters in the \textit{Cobalt Soul's} good graces, can access.

\subparagraph{2. Westhall Academy}
No university in Tal'Dorei can hope to compete with the prestige of \textit{Emon's} \textit{Alabaster Lyceum}, but the Westhall Academy does its best. Its sand-colored stone walls and copper-patina dome are an unmistakable landmark of the Scholar Ward. They imply a stately and refined institution with an unpretentious, workaday attitude. Most of the Westhall Academy's students come from humble backgrounds, and the lord mayor does her level best to ensure that the academy is funded enough to let anyone from the plains who shows great promise attend, free of cost. Mages, architects, poets, artists, agriculturalists, and people of all disciplines can find an education at Westhall.

\subparagraph{3. Yuminor Observatory}
\textit{Westruun} is a fairly flat city, and one tower on the edge of the Scholar Ward looms high above the nearby \textit{Westhall Academy in Westruun}. It looks forever skyward, the arcanists inside constantly charting the movements of the heavens. A quasi-religious order of diviners and astromancers named the Scions of Yuminor study here, supported by the academy's headmistress, Estella Ladimar, who is herself a blood descendant of the heroic wizard Atz Yuminor. The observatory is open to all, and the diviners here will gladly give a star reading for 10 gp, once per day. Receiving this reading allows you to roll 1d20. Note the result—you can exchange it for any d20 roll made by a creature you can see. This divination fades after 24 hours. Despite the Scions of Yuminor's affability, something seems odd about them, like they have stared into the stars just a bit too long and seen something they wish they could forget.

\subparagraph{4. First Bastion}
The most magnificent of the temples within \textit{Westruun's} Temple Ward is the First Bastion. Dedicated to the The Lawbearer, this opulent sanctuary is a common attraction for visiting folk of all faiths. Its magnificent, even decadent, gold-and-marble exterior is a sight most people can't imagine exists upon the Material Plane—it seems like nothing less than a building from the realm of the gods themselves. The First Bastion was a secret home for refugees during Kevdak's occupation, and many \textit{Westruun} survivors have turned to the faith of the The Lawbearer after being protected by her grace.

\subparagraph{5. Lord Mayor's Residence}
The Lord Mayor's Residence is a squat manor in the Opal Ward. It was originally named the Margrave's Keep, and as such, it is a defensible and highly fortified compound buttressed against military barracks. The interior has been livened up significantly by the lord mayor, but the exterior is still a spare, military edifice.

\subparagraph{6. Hall of Reason}
This beautiful courthouse in the Opal Ward is constructed in sweeping Lyrengornian style and bears a marble statue of the The Platinum Dragon above its steps. The square before the courtroom steps is a forum for public gatherings and bears a fountain depicting Palest Westruun, the founder of the city, atop an onyx horse.

\subparagraph{7. The Black King}
Palest Westruun was not a king, but his new statue in the Opal Ward is often called the \textbf{Black King}, a subtle jab at the city's previous "Black King," the tyrannical dragon Umbrasyl. The high priest of the First Bastion blessed both the \textbf{Black King} and \textit{Westruun} itself when the statue was unveiled on the anniversary of \textit{Westruun's} reclamation, saying, "May the The Lawbearer grant us all strength of spirit when chaos next threatens our peaceful civilization. When that time comes, may the \textbf{Black King} unsheathe his sword to rally our people and lead \textit{Westruun} to victory."

When the safety of all of \textit{Westruun} is threatened, the \textbf{Black King} will animate and defend its people, becoming a \textbf{Black King}. Its personality is dour but determined—regal and single-minded in its desire to protect its city and its people.

\subparagraph{8. Survivors' Legacy}
The names of all those in \textit{Westruun} who survived the reign of Kevdak and Umbrasyl are carved on a wall in the first chamber of the Underwalk. Legend has it that the memorial was created by \textbf{Percival de Rolo} when the people of \textit{Westruun} first hid here from the wrath of the \textit{Chroma Conclave} after Umbrasyl's death, but it has been expanded in the years since, now including the names of all who perished in the struggle as well.

The Survivors' Legacy is a symbol of hope, warding away the sorrow and evil of the Underwalk beneath, and to this day common, folk make trips here to honor their lost loved ones. Some of the more generous folk give aid to the homeless living here as well, and some more adventurous travelers make prayers to the The Dawnfather and the The Matron of Ravens here before delving into the Underwalk's deeper reaches.

Once per week, a character related to one of the people named on the Survivors' Legacy can pray there and gain Advantage and Disadvantage on death saving throws for the next 24 hours.

\subsubsection{Westruun Adventures}
Game Masters who set their adventure in \textit{Westruun} can use these plot hooks for inspiration.

\subparagraph{They Can't Have Flown Away (low-level)}
Over the past week, over four dozen \textbf{plainscow}, exceptional livestock and beasts of burden, have disappeared from \textit{Westruun's} outlying villages. Left in their place were nothing but a series of five-foot holes in the ground and traces of acid scarring. Farmers and caravanners alike are frightened and confused; these recurring attacks threaten to destroy \textit{Westruun's} agriculture and disrupt commerce. Jeremiah Sook, head of the regional Rancher's Guild, is willing to pay a hefty bounty for the culprits.

\subparagraph{An Assistant's Plight (mid-level)}
Just outside \textit{Westruun} is Slaterock Tower, the abode of renowned archmage \textit{Realmseer Eskil Ryndarien}. Normally, none are invited to the tower, but the characters, wherever they may be, receive a missive from the Realmseer's assistant, Jekt. The letter claims the Realmseer has gone missing within the tower, and the building itself now bars Jekt's entry—and apparently, the Realmseer's exit. Upon arrival, and confirmation of Jekt's request, the tower seems to constantly shift its interior structure, almost acting like a sentient entity. Challenge after deadly challenge is thrown at the party as they traverse the labyrinthine paths of the tower, only to find that the Realmseer himself is actually trying to keep something far more dangerous imprisoned within.

\subparagraph{The Stuff of Nightmares (mid-level)}
Reports of missing children rise once again within \textit{Westruun}, with a seemingly unrelated scourge of nightmares plaguing people old and young alike. Mystic expertise is sought to acquire information regarding either issue, but the sadistic \textbf{oni}, secretly responsible for both, is deft at erasing its tracks. Nyx and Vryll, twin tiefling mages with only a single horn each, have also recently opened a dream-reading stall, and their business is booming. Are they responsible, or is their exploitative business just a red herring?

\subparagraph{Down Below (high-level)}
A massive sinkhole has emerged beneath the First Bastion, taking a portion of the temple into the darkness of the Underwalk Ward. With the depth seemingly bottomless, and initial investigators going missing, high priest Umentu Helmdoff seeks the assistance of brave and able-bodied adventurers to find answers and save the people still in the temple when it vanished. Only one thing is certain when descending into the unknown below—the terror only grows as the light from above fades.

\subsection{Wildwood Byway}
Known as the oldest road in Tal'Dorei, this path linked the adolescent cities of \textit{Kraghammer} and \textit{Syngorn} before the landfall of humankind. During the reign of \textit{Drassig}, the \textit{Silvercut Roadway} intersected with this ancient road. Warren \textit{Drassig} built countless fortifications along this road to defend his holdings against \textit{Syngorn}, a move that made him many allies within \textit{Kraghammer}.

While traveling the Wildwood in broad daylight isn't any more dangerous than riding along the Silvercut, it's common knowledge that the forest at the northern end of this road is fraught with danger. Wild beasts, cultists, and bandits make the Wildwood their home.

\section{Cliffkeep Mountains}
Spanning the northern reaches of Tal'Dorei, the \textit{Cliffkeep Mountains} are an impenetrable barrier between the Republic of Tal'Dorei and the frozen \textit{Neverfields}. The southern foothills are dotted with forests and pockets of civilization, standing resolute against the unwelcoming topography. Deep within the earth itself, the dwarven stronghold of \textit{Kraghammer} flourishes and expands, just a few scant leagues above ancient caverns that plunge into a domain of unknowable horrors.

The \textit{Cliffkeeps} are the largest mountain range in all of Tal'Dorei, and innumerable types of beings make their homes upon the crags, in mountainside caves, and in the miles upon miles of tunnels that descend into the earth below them. \textit{Dwarves}, \textit{half-giants}, \textit{goblins}, and kobolds are the best-known peoples of the \textit{Cliffkeeps}, but the mountains are also home to the fire giants of \textit{Vulkanon}, wandering \textbf{ettin} and \textbf{hill giant}, hungry \textbf{cyclops}, \textbf{frost giant} marauders from the \textit{Neverfields}, and cave-dwelling \textbf{troll}.

Near the earth elemental rift and \textit{the Ashari} village of \textit{Terrah}, the mountains themselves seem to rise up in defense as \textbf{earth elemental} spring from the living stones. Northward, peaks reach higher and winds grow colder as snow and ice take the lands to become the frozen waste of the \textit{Neverfields}. These challenges keep the common folk at bay, the secrets and spoils held within these crags calling to the brave, the clever, and the foolhardy.

\begin{DndReadAloud}{}

\begin{itemize}
\item{Majority Faiths:}
The All-Hammer, The Stormlord, The Dawnfather

\item{Minority Faiths:}
The Platinum Dragon, The Moonweaver, The Lord of the Hells, The Scaled Tyrant

\item{Imports:}
Lumber, spices, \textbf{fish}, grain, livestock

\item{Exports:}
Precious gems, industrial metals, \textit{quartz}, granite, \textit{gold}, \textit{silver}, \textit{platinum}, mithral, cobalt, jewelry, fungal produce, weapons, armor

\end{itemize}

\end{DndReadAloud}

\subsection{Fort Daxio}
Nestled within the \textit{Othendin Pass}, north of \textit{Emon}, \textit{Fort Daxio} is a massive stronghold and series of military structures that, under the rule of the \textit{Tal'Dorei Council}, trains Tal'Dorei's most elite warriors. The fortress was constructed following the \textit{Scattered War} as a strategic bulwark against opportunistic humanoid and giant bandit forces in the \textit{Cliffkeep Mountains}, and as a hidden reserve of military might in the event that \textit{Emon} was ever threatened.

Though Tal'Dorei doesn't maintain a standing army—rather, it assembles one in times of strife from mercenaries, adventurers, and local militias—the forces of \textit{Fort Daxio} are an exception. The legendary Daxio Outriders are a corps of cavaliers some four thousand strong. Half their forces remain stationed in \textit{Fort Daxio} and its surroundings at any given time, serving as a historic barrier between \textit{Emon} and the bandits and creatures of the \textit{Cliffkeep Mountains}. The other two thousand Outriders are spread across the continent, typically in battalions of up to one hundred riders, safeguarding villages, patrolling the \textit{Silvercut Roadway}, putting pressure on the \textit{Ravagers} of the \textit{Dividing Plains}, or protecting traveling diplomats.

\textit{Fort Daxio} itself is led by General Elle Gorgofon, a dour but diplomatic tactician. She served faithfully under the fort's previous leader, Warmaster Mikael Daxio, for nearly thirty-five years before he was killed at the hands of a \textbf{vampire} hoping to throw the fort into chaos. Daxio was resurrected, but is enjoying an extended holiday in \textit{Syngorn} rather than returning to arms.

The Daxio Outriders are divided into four different regiments: The Dusk Regiment, the Gale Regiment, the Aegis Regiment, and the Tide Regiment. Each regiment is led by their own assigned general, all answering directly to General Gorgofon (who in turn answers to the Master of War on the \textit{Tal'Dorei Council} in times of conflict).

The Dusk Regiment is responsible for covert operations across the country and is led by General Fei Yujian, a cheerful, affable man. The Gale Regiment is trained in skyborne combat in addition to horse-riding; after the \textit{Chroma Conclave} was defeated, \textit{Fort Daxio} trained a variety of \textbf{griffon}, \textbf{hippogriff}, and \textbf{wyvern} to ward against airborne threats. It is led by General Jillian Sylph, a calm, tactical human with air \textit{elemental ancestry}. The Aegis Regiment is charged with remaining in \textit{Fort Daxio} to protect \textit{Emon}, and is led by General Elle Gorgofon herself. The Tide Regiment is trained in spellcasting and ship-to-ship combat, in case \textit{Emon} is ever threatened from the western seas. This final regiment is led by General Kay Clearsight, a half-elf trained in the arcane arts at the renowned Soltryce Academy in the Dwendalian Empire of \textit{Wildemount}.

\subsubsection{Fort Daxio Adventures}
Game Masters who set their adventure in \textit{Fort Daxio} can use this plot hook for inspiration.

\subparagraph{Vampire Hunters (high-level)}
Years ago, Warmaster Mikael Daxio was killed by a \textbf{vampire} named Vrylaska Fellbranch. She disappeared without a trace, and the Daxio Outriders have been searching fruitlessly for her for years. While the characters are in \textit{Fort Daxio}, a scout returns with breathless news of vampire attacks in remote settlements in the \textit{Cliffkeep Mountains}. General Elle Gorgofon suspects a trap, but has no choice but to mobilize the Aegis Regiment—and the characters, if they accept her offer of 5,000 gp—to hunt the monster down.

Little do they know that, while the Aegis Regiment is away, the vampire's minions hidden in \textit{Fort Daxio} will come out to play.

\subsection{Gatshadow}
The mighty mountain known as \textit{Gatshadow} is a spire of death that looms over the \textit{Dividing Plains}. \textit{Westruun} sees it always upon the horizon, and the most vicious monsters of the \textit{Bramblewood} grow cruel in its fell shadow. Folkloric tales about the evil powers within \textit{Gatshadow} are told in all cultures of the \textit{Dividing Plains}: stories of \textbf{giant spider} in the forest, of beings of pure shadow that dwell within the mountain's tunnels, and of the tunnels themselves, which are so labyrinthine and shifting that no one who enters will ever escape. These days, most of these stories are told to keep children from wandering into the forest alone—but they are far from simply children's tales.

\textit{Gatshadow} was transformed into a place of evil in the \textit{Age of Arcanum}. Acek Orattim, a priest of the Betrayer God known as the The Chained Oblivion, made the mountain his seat of power. The The Chained Oblivion inexorable might corrupted the mountain, stretching its stone-like putty until it was a jagged needle that towered over all other mountains, and etching a mind-bending labyrinth of twisting passages into the mountain's heart.

Acek spread the The Chained Oblivion nihilistic despair across the \textit{Cliffkeep Mountains} and the \textit{Dividing Plains} below. All bowed before their influence their minds lost to unfathomable contemplation of the end of all things. Untold thousands died at the hands of Acek and his god, willingly entering \textit{Gatshadow}—which then became their tomb.

The bitter cold of oblivion threatened to consume the land, but it was repulsed by the The Dawnfather and his warriors of light. The The Dawnfather struck a telling blow against the Oblivion, and doggedly pursued as it retreated to \textit{Gatshadow}. He shackled and banished the god of the end within \textit{Gatshadow's} depths. It was said that Acek's mind and spirit were so saturated with his master's essence that the god's banishment tore his body asunder. Today, only the empty maze of deadly chambers and traps remain beneath the mountain, forgotten and awaiting discovery.

The high bluffs of \textit{Gatshadow} have in recent years been the home of many a villain, most recently Umbrasyl the Hope Devourer, the black dragon of the \textit{Chroma Conclave}. Legend says that \textit{Gatshadow's} interior tunnels spiral infinitely downwards, into a void of utter darkness. Immortal, white-eyed devotees of the The Chained Oblivion pace endlessly into \textit{Gatshadow's} tunnels, carrying lightless torches, whispering that the secret to freeing their god from the Divine Gate is sealed in the depths below. The stone doors that bar entrance to \textit{Gatshadow's} lowest reaches are sealed by magic so great, even the most inquisitive of the \textit{Arcana Pansophical} can't breach them—and it is well that they cannot, for who knows what terror would be unleashed if they were opened.

\subsubsection{Gatshadow Adventures}
Game Masters who set their adventure in \textit{Gatshadow} can use these plot hooks for inspiration.

\subparagraph{Blood of the Dragon (high-level)}
The body of the \textbf{ancient black dragon} Umbrasyl was never recovered. Some think that's a waste of a perfectly good business opportunity—after all, a black dragon's bones are a potent alchemical reagent, and its scales stronger than any steel.

However, there's a reason no one has dared to recover the Hope Devourer's corpse. When Umbrasyl died, the caustic acid and magical blood within his body seeped through the caverns of Upper \textit{Gatshadow}, devouring stone and transforming the mountain's upper levels into a vitriolic maze teeming with oozes, half-black dragon warriors, and the restless spirits of \textit{goliaths} from the Herd of Storms that failed to kill the dragon.

\subparagraph{No Lights in the Darkness (epic-level)}
The \textit{Betrayer Gods} were banished to their planes of eternal imprisonment, but something has upset the laws of the planes. A being called the Nameless—supposedly a cultist of The Whispered One—has stolen her master's secret knowledge and fled to \textit{Gatshadow} in search of the tortured spirit of Acek Orattim. Agents of the \textit{Arcana Pansophical} have sensed that the lower gates of \textit{Gatshadow} have been breached, and frantically call for Tal'Dorei's greatest heroes to rush to the evil mountain. As the characters hurry to the bottom of the mind-sundering labyrinth within, they discover that they race against not only the new Herald of Oblivion, but also the \textit{Remnants} she betrayed.

\subsection{Jorenn Village}

\begin{description}
\item[Large Village:]
Population 2,480 (64\% \textit{humans}, 9\% \textit{dwarves}, 7\% \textit{halflings}, 7\% \textit{gnomes}, 13\% \textit{other races})

\end{description}

For many years, Jorenn Village was just a tiny pioneer town trying to profit from the shadegrass of the \textit{Umbra Hills}. Over the past century, the village has exploded in size and population, after plentiful silver veins were discovered within the nearby rocky hills.

Hardworking common folk and opportunistic prospectors alike flocked to Jorenn, despite its unsettling proximity to the \textit{Umbra Hills}. Jorenn is far enough away from \textit{Emon} that the influence of the \textit{Tal'Dorei Council} rarely touches it, giving opportunists and outlaws freedom to subvert and control both the village government and the local economy. Thirty years ago, a small circle of vigilantes and mercenaries declared themselves protectors of Jorenn, and took the name Shadewatch, hunting and slaying creeping dangers in the night and returning to great fanfare. Over the past three decades, others seeking fame and glory joined the ranks of the Shadewatch, and the group ballooned into the unofficial town law enforcement.

The Shadewatch is deeply corrupt, and it has been almost since its first year. The vigilantes personally stoke the fires of fear, all while demanding increasingly galling tribute for their services—and breaking kneecaps if they don't get it. Today, the people of Jorenn are just as afraid of the Shadewatch as they are of the creatures the watch keeps at bay.

The Shadewatch is a minor organization in the grand scheme of things, but Shademaster Arhanna Lewyn has grand ambitions. She has sent messages to every faction in Tal'Dorei, from the \textit{Myriad} to the \textit{League of Miracles}, with an interest in cruel subterfuge. It's clear that she's planning something: every night, she consults a ball of polished white quartz, and every week she sends a trio of her agents to the \textit{Umbra Hills}—or to the base of \textit{Gatshadow}.

\subsubsection{Jorenn Village Adventures}
Game Masters who set their adventure in Jorenn Village can use this plot hook for inspiration.

\subparagraph{An Offer You Can't Refuse (low-level)}
The characters are attacked by a group of \textbf{worg}, and more just keep coming. Just as they're about to be completely overwhelmed, they're saved by a sudden volley of black-feathered bolts! Six thugs join the battle and "rescue" the characters by driving off the remaining \textbf{worg}. These Shadewatch "officers" ask for the characters' names first, then brusquely declare that they like their moxie. They offer the characters a position in their force back in Jorenn, with unbelievable pay: 50 gp a day each.

If this deal sounds too good to be true, that's because it is: the Shademaster needs strong-willed sacrifices if she's to open the gates of \textit{Gatshadow}, a feat she can accomplish by teleporting up to twelve creatures into the depths of the mountain with her whispering quartz orb.

\subsection{Kraghammer}

\begin{description}
\item[City:]
Population 73,550 (73\% \textit{dwarves}, 8\% \textit{gnomes}, 5\% \textit{humans}, 14\% \textit{other races})

\end{description}

Dozens of dwarven redoubts rose and fell in the tunnels beneath the \textit{Cliffkeep Mountains} during the \textit{Calamity}, and in the uncertain years that followed it. Legend holds that a great civilization called Uthtor once reigned supreme in the \textit{Cliffkeeps}, and that all of the deep settlements are fragments of Uthtor's glory, fractured by the \textit{Calamity}. Over the centuries, these fragments divided and coalesced dozens of times until, finally, those who remained came together to create Kraghammer.

Clan Jaggenstrike united thousands of \textit{dwarves}, \textit{gnomes}, and other dwellers of the deep halls in the creation of Kraghammer. Having quietly died out centuries ago, the clan now holds a position of mythic reverence within the city. Kraghammer's founding was dedicated to the The All-Hammer, divine creator of the \textit{dwarves} and patron of creation and ingenuity. This god's industrious spirit filled the people of Kraghammer and blessed them with rich veins of \textit{platinum}, mithral, adamantine, and other precious metals that have sustained them since the city's founding.

An ominous walk of a thousand stairs climbs the side of the mountain to the gates of the city—carved from the mountain itself, just wide enough for two \textit{dwarves} to walk abreast. With the path too steep and treacherous for carts or common beasts of burden to traverse, traders and trappers use \textbf{giant goat} to transport their goods to and from market, while wealthy dwarven merchant lords employ serfs and laborers to carry their goods. Kraghammer's impenetrable adamantine gates stand at the end of a perilous bridge overlooking a hundred-foot chasm, guarded by three watchtowers carved directly into the mountain.

Within those gates, the massive stronghold is made up of three cylindrical levels called "slabs" that descend deep beneath the mountains. Kraghammer's granite walls are lit by roaring fires, enchanted torches, and shimmering candlerock, which glows with internal fire. Even when the evening fires are extinguished, the amber light of the Bronzegrip Metalworks at the base of the city bathes Kraghammer in its warm glow.

\subsubsection{Government}
\textit{Kraghammer} law is written by a consortium of the five most powerful houses within the city, each appointing a delegate to serve on their behalf. That tenuous union of houses is kept by the city's executive officer, the Ironkeeper. The Ironkeeper is elected by the houses and carries a term of ten years, with no term limit, though most Ironkeepers contain their rule to three to five terms.

The current sitting Ironkeeper is \textit{Palluda of House Zuurthom}. She rose to prominence advocating for the rights of workers in the wake of the \textit{Chroma Conclave}, and is a strong advocate for uniting with the Republic of Tal'Dorei to strengthen their might against future supernatural threats. She is sitting comfortably in the midst of her second term and enjoys strong support from members of \textit{Kraghammer's} working class—and staunch opposition from traditionalists, who cling to the city's fading culture of \textit{isolationism}.

The laws of the city are enforced by a military class of elite warriors called Carvers that act as guards, soldiers, and jailers under the guidance of the Ironkeeper. Becoming a Carver is a lengthy, rigorous, and taxing pursuit, and many who choose to join fail to complete their training.

Most who join the Carvers are \textit{dwarves} from families with centuries of history within the Carvers. This isn't just because of the pressure to uphold their family's honor. Most hopefuls are also those who value \textit{Kraghammer's} ancient traditions and rigid social order—a set of beliefs largely held by dwarven clans with long family histories. Once officially named, a Carver is given a homestead built into the center slab of the city, a set of immaculate armor and weapons, and a steady income.

In the years following the defeat of the \textit{Chroma Conclave}, the famously \textit{isolationist} and xenophobic houses conceded to formally make an alliance with the Republic of Tal'Dorei. An ambassador from \textit{Kraghammer} now sits on the \textit{Tal'Dorei Council}, and rumors are beginning to circulate within the city that the heads of the Great \textit{Houses of Kraghammer} are in talks with the \textit{Tal'Dorei Council} to make \textit{Kraghammer} an official constituent city-state of the republic.

\subsubsection{Isolation of Kraghammer}
It's not easy being an outsider. \textit{Kraghammer} has long been a city of \textit{dwarves}, for \textit{dwarves}, and anyone else within the city is an interloper. Some traditionalists within the city still think that surface-born folk should spend as little time in \textit{Kraghammer} as possible, so as not to dilute the ancient splendor of their culture.

Yet times are changing. A coalition of people, ranging from \textit{dwarves} to \textit{gnomes} to traditionally surface-dwelling people who have made \textit{Kraghammer} their home, have demanded that the Great Houses amend the city's most xenophobic laws. Despite the angry protestations of the city's most bigoted \textit{dwarves}, the Great Houses have acquiesced to the requests of their people. This is due in no small part to these protestors having the support of the current Ironkeeper—\textit{Palluda Zuurthrom} has long been a proponent of \textit{Kraghammer} entering the Republic of Tal'Dorei and ending their centuries of relative isolation.

No non-dwarf has ever served as a delegate of the houses, let alone the Ironkeeper. Progress is slow, especially in \textit{Kraghammer}, but it is happening. The most despicable of the city's xenophobes have left of their own volition, creating other dwarven settlements throughout the \textit{Cliffkeeps} where they can keep to themselves and "their own kind" in peace.

\subsubsection{Crime}
Citizens of \textit{Emon} who visit \textit{Kraghammer} often marvel at how lawful and organized the dwarven enclave is. Crime is ostensibly low within \textit{Kraghammer} thanks to the Carvers' protection, but this peace comes at a cost. Crime is low because the courts of \textit{Kraghammer} almost always reach a guilty verdict. As a long-lived and durable people, most \textit{dwarves}, particularly those with a family reputation to uphold, would rather plead guilty and endure years, even decades, of hard labor than disturb the fabric of their society by fighting conviction.

Beyond this, corruption is rife within the Carvers, and they are able to brutalize even innocent people with relative impunity. This corruption is shielded from public scrutiny by the vast wealth and political influence of House Glorendar, their patron house. Because of Glorendar's close ties with the Carvers (and the mutual benefits thereof), corruption has been hard to uncover within either organization.

Over twenty years ago, however, a gnomish journalist named Ida Mudrake uncovered something damning within the vaults of House Glorendar, bringing some hope that change may yet move steadily forward. She has been forced to flee the city and now lives with a security detail in \textit{Emon}, but she works tirelessly to bring the crimes of House Glorendar to light. Many in \textit{Kraghammer} hope that this evidence, whatever it is, will shake both Glorendar and the Carvers to their foundations. Many more fear that even true justice is not worth the chaos that bringing Glorendar's misdeeds to light will cause.

\subsubsection{Geography and Climate}
Though the \textit{Cliffkeep Mountains} themselves are blanketed in thick snow, the bone-piercing chill does not extend into \textit{Kraghammer} itself. Heat from the Bronzegrip Metalworks and the other blast furnaces in the Bottom Slab rises and permeates the entire city, though the Top Slab still grows cold in the winter. The city itself is carved out of the granite core of Mount \textit{Kraghammer}; though the \textit{dwarves} rarely see the night sky, the inner walls still sparkle like starlight in the amber glow of the furnaces.

\subparagraph{Top Slab}
Also called the Arch, this uppermost tier is the entry level of the city, built on a massive ring. Like a hub with many spokes, countless neighborhoods and tunnels spread out from the Top Slab. A fine layer of soot tends to coat much of this slab, due to imperfect ventilation from the industrious center slab; despite this, this level is where most people make their homes, with thousands of stone-built abodes dotting the walls, while popular taverns like the trendy Firebrook Inn and the historic but austere Ironhearth Tavern fill alcoves pushed into the rock. For hundreds of years, most non-dwarves within \textit{Kraghammer} were relegated to living in a Top Slab neighborhood called the "Otherwalk." These days, non-dwarves live all over the city, and the Otherwalk has been renamed the Cavernwalk. It's still mostly inhabited by non-dwarves, and it has a vibrant culture all its own.

As members of the \textit{Kraghammer} working class, "Topslabbers" either make the long commute to labor in the Bottom Slab or work in the Toppers' farms above the mountain. Others are indentured servants to wealthy \textit{dwarves} in the Middle Slab, and only return to their families on the Top Slab on their monthly day of rest.

\subparagraph{Center Slab}
Also referred to as the Heart, this middle tier is the widest and most varied rung on the journey through \textit{Kraghammer}. It is stratified between the residences of \textit{Kraghammer's} old money (its long-established noble houses), and the new (the wealthy merchant class). At the exact center of the slab is the Pyrethrone, seat of the Ironkeeper. The geometric symbolism of the Pyrethrone being at the dead center of the city is a fact most visitors to \textit{Kraghammer} hear ad nauseum. Radiating out from the Pyrethrone are the various fortress-manors of the dwarven noble houses, with several Carver barracks in easy reach.

Below, the outer rings of the Center Slab are occupied by most of the city's non-mining businesses, from smithies, breweries, and jewelers, to tailors, butchers, and tinkerers like the \textit{Cracksackle Union}. The Hunter's Club is also well known as the chief provider for non-imported meats to the city, sending parties into the mountains or the nearby \textit{Torian Forest} for wild game. The extravagant marble temple known as The All-Hammer stands on the center slab of the city, the grandest of all shrines to the most worshipped deity under the mountain. Enormous in size and impeccable in both detail and architectural design, this hall calls to artisans and crafters throughout the city for inspiration.

Of the five noble houses, only two conduct their affairs entirely in the Center Slab. The magically talented Lord Steddos Thunderbrand and his family are known to rarely leave their mansion, except on business with the Ironkeeper. Wallera Glorendar used to oversee the combat training of all Carvers within \textit{Kraghammer}, but hasn't been seen in public for twenty years. Her daughter, Kazya, runs all her affairs within the public eye.

\subparagraph{Bottom Slab}
Best known as the Pit, this tier is the industrious powerhouse of \textit{Kraghammer}, boasting the most forges per square mile of any location in Tal'Dorei, courtesy of the Bronzegrip Metalworks and their competitors. Most of these forges are potent, highly industrialized, and mechanized blast furnaces. The expansive \textit{Keenstone Quarry} is chief among the numerous mining operations that have tunneled through the depths of the mountains—depths whose true dangers are obscured by centuries of legend, and therefore not fully understood.

A network of fungus-farming paths called the Glowgrove works its way around this level, passing through all sorts of caverns. The mushrooms are harvested by the Toppers, the dwarven agricultural league that got their name from the small topside farms they tend above the city.

Three of the five dwarven houses have business interests in the Bottom Slab, though they direct its operations from the next tier up. \textit{Nostoc Greyspine} is the grim overseer of \textit{Keenstone Quarry}. Haddi Bronzegrip is the foredwarf of the Bronzegrip Metalworks; she claims to be the wealthiest dwarf in the North, a claim no one has ever been able to successfully dispute. Blenton Zuurthom is an architectural savant, and one of the older and kinder noble \textit{dwarves}; he tries his best to oversee every major construction project in the Bottom Slab.

\subsubsection{Points of Interest}
\textit{Kraghammer} has a number of famous landmarks, as well as a few mythic locations said to exist in the caverns beneath the city, as depicted on its map.

\subparagraph{The Starshrine (Top Slab)}
Most \textit{dwarves} worship the The All-Hammer, and his teachings of community and ancestral piety have been all but completely subsumed into \textit{Kraghammer's} culture at large. For outlanders with different faiths, however, the monolithic religious culture can feel inescapable. Shrines to all the \textit{Prime Deities} have spread throughout \textit{Kraghammer} since the fall of the \textit{Chroma Conclave}, but the Starshrine was the first.

When the world trembled at the might of \textit{Thordak the Cinder King}, tremors rocked the upper layers of \textit{Kraghammer}, causing certain caverns to collapse—and in some cases, to reveal places of great beauty. The Starshrine is a vast grotto whose granite walls mysteriously sparkle with a perfect replica of the stars in the night sky. It can be found within the Top Slab's culturally diverse Cavernwalk district. Once per year, when a character prays to one of the \textit{Prime Deities} here, there is a Your prayer goes unanswered. chance per character level that their prayer is answered, as the cleric's Divine Intervention class feature.

\subparagraph{Cracksackle Guildhall (Middle Slab)}
The headquarters of the Cracksackle Union looks unlike anything in \textit{Kraghammer}. It's made of a dozen glistening steel domes, like sleek, metal igloos, and the light of unusual experiments flicker through its slitted windows at all hours of the night. Most within \textit{Kraghammer} consider the building an eyesore, but all interested in technological tinkering and ingenuity love it for its uniqueness. If nothing else, it plainly displays the oft-baffling ingenuity of its tinkerers.

The Cracksackle Union was founded by gnomish refugees from \textit{Wittebak} nearly four hundred years ago. While intrinsically tied to the Bronzegrip Metalworks for resources, contacts, and distribution, the Cracksacklers have made themselves indispensable to the dwarven elite with their inventive contributions to mining and masonry technology.

On any working day, a dozen tinkerers, many of them \textit{gnomes}, will set up shop in the courtyard, selling strange and untested devices they have made in their workshops. The GM determines which non-magical curiosities can be purchased here, but some include:

\begin{DndTable}[header=Cracksackle Wares]{XXX}

Item & Cost & Description\\
\textit{Dynamite} & 100gp & This bundle of explosives can be lit and thrown up to 30 feet as an action. It explodes at the start of your next turn. Creatures and objects within 20 feet of it must make a DC 12 Dexterity saving throws, taking 6d6 fire damage on a failed save, or half as much damage on a successful one.\\
\textit{Gluebomb} & 50gp & This unique explosive can be primed and thrown up to 30 feet as an action. It explodes on impact. Creatures within 20 feet of it must make a DC 12 Dexterity saving throws or be restrained for 1 minute. A restrained creature can make a DC 12 Strength check as an action, escaping on a success.\\
\textit{Stink bomb} & 25gp & This smelly explosive can be primed and thrown up to 30 feet as an action. It explodes at the start of your next turn. Creatures within 20 feet of it must make a DC 12 Constitution saving throws or be poisoned for 1 minute. A creature can repeat this save at the end of each of its turns, ending the effect on itself on a success.\\
\end{DndTable}

\subparagraph{Keenstone Quarry (Bottom Slab)}
Despite all of \textit{Kraghammer's} wondrous crafts and goods, its most significant export is raw stone—most of which is excavated from Keenstone Quarry, making the quarry's overseer, \textit{Nostoc Greyspine}, a very rich dwarf indeed. As the legend of Vox Machina spread across Tal'Dorei, \textit{Overseer Greyspine} got an idea: open the quarry to adventurers. Vox Machina famously entered the vast network of dangerous tunnels beneath \textit{Kraghammer} from a fissure at the depths of the quarry, and countless imitators have tried to follow in their footsteps—some emerging with fabulous treasure, but most never returning at all.

To take advantage of this opportunity, \textit{Greyspine} has advertised the quarry far and wide as the "Gateway to the Depths," and posted an entrance fee of 25 gp per adventurer. Of course, neither \textit{Greyspine} nor the staff of the quarry bears any liability for death, dismemberment, or emotional trauma sustained with the depths.

\subparagraph{Hall of Burning Mushrooms}
The Bronzegrip Metalworks just carved deeper into the mountain to expand their furnaces, only to discover the most magnificent cavern: a mile-wide cave filled with bioluminescent purple mushrooms. Industry had to move in, and the Bronzegrips hired Wyrmhide Thunderbrand and his pyromancer brigade to torch the cavern. The mushrooms never stopped burning, and the \textbf{myconid adult} that lived in the fungal forest never stopped fighting. Characters known within \textit{Kraghammer} may be requested by House Bronzegrip or Thunderbrand to cut through the inferno and destroy the \textbf{myconid Sovereign}.

\subsection{Emberhold}

\begin{description}
\item[City:]
Population 21,180 (73\% \textit{dwarves}, 15\% \textit{elves}, 12\% \textit{other races})

\end{description}

A passage from the \textit{Keenstone Quarry} in \textit{Kraghammer} leads below the earth. Deep. Deeper. Until the light of day is entirely gone, and then deeper still. Across chasms miles-deep and full of magma, beneath great hanging gardens of glowing \textit{malachite}, lies the \textit{Emberhold}. It stares imperiously over an underground kingdom of the \textbf{duergar}, \textit{dwarves} who long ago broke away from the rule of the Great \textit{Houses of Kraghammer}.

Millennia of fear and superstition have turned the \textit{dwarves} of \textit{Kraghammer} into bitter, fearful enemies of their estranged kin. Untold dozens of duergar villages span the glassy lava-plains that the \textit{Emberhold} is built upon, with mighty armies constantly coming and going, sometimes returning with just a few bloody survivors, other times returning with captives in chains—and other times returning not at all.

The \textit{dwarves} of the \textit{Emberhold} rarely rise to the surface; they have lived so long in the darkest depths of the earth that the light of the sun sears their eyes and prickles their skin. Few ever arise even to \textit{Kraghammer}, for they are taught that only death awaits them there. Nevertheless, many duergar doubt the narrative that has been fed to them since birth, and make attempts to escape to the surface. Those duergar that have left the ways of the \textit{Emberhold} behind and learned of the outside world have remarked how unbelievable it is that their people were so misled by lies about the barbaric cruelty of surface-dwellers.

\subsubsection{Government}
The \textit{Emberhold} and the towns that serve as its vassals are ruled by an absolute monarch. The current King of Embers is the wrathful King Thangrul II. He rules through fear and conquest, spreading his armies to the farthest corners of Tal'Dorei's underground caverns in the name of the The Crawling King.

Thangrul's rule is supported by a contingent of \textbf{drow} from their decadent, alien domain of \textit{Ruhn-Shak}. Most in Thangrul's court know well enough to mistrust these ambassadors of the far-flung caverns, for their minds have been twisted by the influence of not just the The Spider Queen, but also the aberrations with whom they have made strange alliances in order to sustain their power. Thangrul himself cares little for his advisors' mistrust of his elven allies and openly declares that he will make use of them as long as they are useful.

Vox Machina's assault of the \textit{Emberhold} over twenty years ago cost the duergar their king and queen, but rumors of a lost princess of Ember have begun to circulate through the realm, and the survivors of Vox Machina's raid have rebuilt their hierarchy and continued their endless war of conquest throughout the depths.

\subsubsection{Geography}
The \textit{Emberhold} itself is a fortress built of black stone and \textit{obsidian}, framed by rivers of molten rock and surrounded by a massive city of ramshackle dwellings and towers that mark the boundaries of its peoples' domain.

\subsubsection{Emberhold Adventures}
Game Masters who set their adventure in the \textit{Emberhold} can use this plot hook for inspiration.

\subparagraph{The Lost Heir (mid-level)}
King Murghol of the \textit{Emberhold} was killed in combat by \textbf{Scanlan Shorthalt}, gnome bard of Vox Machina, and his wife, Queen Ulara, was slain by the aberrant tyrant K'Varn. All of the \textit{Emberhold} knows that the King and Queen had a young daughter—only six years old at the time of their death—but none seem interested in seeking \textit{Emberhold's} missing heir. The new king, Thangrul II, grows particularly fiery whenever anyone mentions her name in his presence.

\subsection{Yug'Voril}
Where the \textit{dwarves} of \textit{Kraghammer} whisper fearful words of the \textbf{duergar}, so too do the \textbf{duergar} whisper of Yug'Voril. Beneath the deep \textit{dwarves'} fiery kingdom, ever deeper into the heart of Exandria, there is a miles-wide grotto coated in glowing crystal. A lake fills the grotto, its surface so smooth and placid it seems like a \textit{jade} mirror, perfectly reflecting the light of the crystals around it. In the middle of the lake is an island, atop which loom the grand, austere towers of Yug'Voril.

The clean lines and beautiful facade of the city belie its current, abominable denizens.

Thousands of years old, the builders and previous occupants have long vanished, leaving the unspoiled ruins to become home to a host of \textbf{mind flayer} from a plane beyond the boundaries of reality. The tyrannical society of these alien creatures flourishes here, far from their cosmic pursuers, giving them the freedom to build a legion of thralls: mind-slaves whose delicious brains are allowed to remain in their bodies...until those bodies fall apart from toil.

Yug'Voril and \textbf{mind flayer} fell recently under the control of the extraplanar aberration called K'Varn, before he was slain by Vox Machina. Now that the \textbf{mind flayer} have been freed from K'Varn's tyranny, they have reclaimed Yug'Voril as their own and have continued work on their strange and terrible plots.

\subsubsection{Yug'Voril Adventures}
Game Masters who set their adventure in \textit{Yug'Voril} can use this plot hook for inspiration.

\subparagraph{A Mind Is a Terrible Thing to Waste (high-level)}
Four of the foremost minds at \textit{Emon's} \textit{Alabaster Lyceum} have disappeared without a trace. All vanished on the same night, and there have been no confirmed sightings since. Investigations have been underway for several months, and some \textit{dwarves} of \textit{Kraghammer} have recently reported seeing four figures that could be the missing scholars entering the lower passages of the \textit{Keenstone Mine} under cover of night.

Something has mentally dominated these scholars. The alien inhabitants of \textit{Yug'Voril} hunger for a specific morsel of powerful knowledge within their brilliant minds.

\subsection{Lyrengorn, the Elvenpeaks}

\begin{description}
\item[Town:]
Population 5,430 (93\% \textit{elves}, 3\% \textit{dwarves}, 2\% \textit{humans}, 2\% \textit{other races})

\end{description}

Every year, dozens of travelers from across Tal'Dorei travel past the far northern reaches of the \textit{Cliffkeep Mountains} to see the The Moonweaver Ribbons twist around the Elvenpeaks, pulled through the sky by those mysterious figures on \textbf{wyvern}. The Elvenpeaks are a single mountain split perfectly in the center, with ice-sheeted forests topping both peaks. While the exterior of these forest canopies appears frozen and without life, the interior is instead a warm and humid wood, blooming with flowers of all colors and roamed by living plants and wild game. In the boughs of these thick tropical trees also lives a society of wood \textit{elves}, some of those wild folk who broke away from Yenlara's pact and the city of \textit{Syngorn}. Their city of Lyrengorn—\textit{lyren'gorn} literally meaning "mountain city" in Elvish—is the last bastion of safety before crossing into the endless ice of the \textit{Neverfields} to the north.

These lyren'alfen—"high \textit{elves}"—call themselves friends of none and enemies of fewer; they welcome all who manage to make it to their home, provided they come in peace, and take no sides in any conflict. In ages past, they were highly sought-after as mercenaries because their warriors had learned the rare art of \textbf{wyvern} riding. When Warren \textit{Drassig's} legions marched to conquer Tal'Dorei, the \textit{elves} of Lyrengorn refused to die for anyone else's war, retreating instead to their mountaintops and lush forests.

Today, few of the lyren'alfen are actually warriors, though many among them are hunters. In the intervening centuries, \textbf{wyvern} riding has transformed from a martial art into a ceremonial one; great athletes called skyswimmers train year-round for the annual coming of the The Moonweaver Ribbons, great streaks of multicolored light that illuminate the northern sky on the winter solstice. The skyswimmers and their \textbf{wyvern} Mounts and Vehicles seize the strands of light on their spears and paint radiant, esoteric portraits in the wintry sky, drawing intrepid outsiders from across Tal'Dorei to camp on remote mountains and marvel at their visual poetry.

Beyond the ceremony's artistic merit, it's also a ritual to the goddess known as the The Moonweaver, chief deity of these \textit{elves}. The high priestess of the moon claims that the annual rite purifies Lyrengorn and gives the wild \textit{elves} long life.

\subsubsection{Lyrengorn Adventures}
Game Masters who set their adventure in Lyrengorn can use these plot hooks for inspiration.

\subparagraph{The Impossible Art (low-level)}
While traveling through the \textit{Cliffkeep Mountains}, the cold becomes too much for the characters to bear, and they pass out in the snow. When they awaken, they find themselves being tended to by a trio of Lyrengorn \textbf{wyvern} riders, but their roles are suddenly reversed when a massive band of raiders, most of them \textbf{goblin} and human thugs, pincushions the riders in mid-flight. When the characters crash in the snowfields, they find their saviors dead, and must try to learn the impossible art of \textbf{wyvern} riding before they are claimed by the blizzard.

\subparagraph{The Incorruptible Light (high-level)}
While the characters are in Lyrengorn, they discover a small community of dark elf refugees. They fled here from a monstrous, nameless source of aberrant corruption that has been devouring their society for centuries. The dark \textit{elves} hope that the skyswimmers' ceremony will cleanse them of the corruption and fits of sadistic cruelty that still plague them from their days under the abberrations' terrible thrall. Strange accidents conspire to kill the skyswimmers and the priestesses, and on the day of the solstice, a pair of \textbf{beholder} rise from the depths of the earth and try to bend the The Moonweaver Ribbons to their evil purposes.

\subsection{The Neverfields}
The farther north one goes into the \textit{Cliffkeep Mountains}, the taller the peaks grow and the colder the weather becomes, until the snowstorms are so thick that no human unprepared for their chill can survive. The spiky peaks grow shallow, trading treacherous rock for glacial ice and the untempered fury of endless blizzards. No bastions of living society exist this far north, though some brave, foolish warriors endure the elements here to hunt rare and powerful game or to seek lost relics from before the \textit{Calamity}.

Legends tell of a time when this region was green and temperate, when trees and joyous elven enclaves covered the land. Moreover, these tales say that relics of this elysian era are still buried deep beneath the icy wastes. Many that go seeking answers to these rumors never return. The few that do, however, occasionally emerge with artifacts of intense power and mystery, reigniting interest in the region and sending further generations to their doom among the ancient, ice-rimed ruins and plains of endless glacial wasteland.

\subsubsection{Neverfields Adventures}
Game Masters who set their adventure in the \textit{Neverfields} can use this plot hook for inspiration.

\subparagraph{Seed of Life (mid-level)}
Heroes of Tal'Dorei who have made themselves known to \textit{Syngorn} may suddenly find themselves contacted by \textit{Verdant Lord Celindar}, leader of \textit{Syngorn's} armies. His report is grim: the warding trees of the \textit{Verdant Expanse} no longer protect the city; in fact, some have turned against it. The \textbf{treant} that once defended the \textit{elves} have grown pale and sickly, and are now attacking the elvenguard. \textit{Celindar} beseeches the characters to meet with his counterpart in Lyrengorn, Skyblade Hali Instralidar, who has been surveying the ice fields of the north for centuries. Why? Because elven legend speaks of the Seed of Life, the last bloom of an ancient \textbf{treant} named Cedargaunt, which may restore the withered treants' senses. The Seed of Life is worn on the neck of a \textbf{frost giant} jarl's pet \textbf{remorhaz} within the \textit{Neverfields}.

\subsection{Othendin Pass}
This massive valley carved into the southwest section of the \textit{Cliffkeep Mountains} was, in ancient times, the territory of \textbf{stone giant} clans that warred with the \textit{dwarves} of \textit{Kraghammer}. After the arrival of humanity on Tal'Dorei's shores, the people of \textit{Kraghammer} and the new settlers of \textit{O'Noa} banded together and pushed deep into the mineral-rich valley.

The \textbf{stone giant} fought back, but they were overwhelmed by the might of their combined foe. They retreated to their caves higher into the mountains to the northeast, and have rarely been seen by humanoids in the centuries since. This allowed the forces of \textit{Emon} to claim the resources and land of the Othendin Pass for its people, establishing \textit{Fort Daxio} to defend it.

However, the population of deadly landsharks—creatures called \textbf{bulette}—native to the region has been on the rise. Some say it's because the giants who hunted them were driven away, but since the giants have been gone for a century, few can understand why the landsharks are only taking over now. The miners, gatherers, and passing travelers do their best to cull the \textbf{bulette} population both on their own and with the help of hired adventurers, but they can only ever hope to keep the deadly monsters at bay.

\subsubsection{Othendin Pass Adventures}
Game Masters who set their adventure in \textit{Othendin Pass} can use this plot hook for inspiration.

\subparagraph{Bite the Bulette (mid-level)}
A druidic order called the Stoneherders, made up mostly of \textbf{stone giant} and \textit{goliaths}, have taken up residence in an important choke point of the \textit{Othendin Pass}. They claim that they are the Stoneherders, a great druidic order that keeps the balance between civilization and nature in the north. \textit{Fort Daxio's} military presence in the pass has kept them away for several generations, but an unusual rise in the \textbf{bulette} population has forced them to return.

The Stoneherders' plan to thin the \textbf{bulette} herd requires the characters' aid in rounding up the deadly creatures; once they have been assembled, the Stoneherders will tame as many as they can and move them elsewhere in the \textit{Cliffkeep Mountains}. Those that can't be tamed will be relocated to other parts of the mountains, where they can be hunted again.

\subsection{Palace of Wonder}
In the desolate and uninhabited northwestern reaches of the \textit{Cliffkeep Mountains} stands a ruined citadel. All that remains of it are crumbling walls, half-beaten to dust by the unrelenting winds of the mountains, and one forlorn tower, its turret torn asunder to reveal the mildewed, rain-fattened wooden beams within. This ruin, a relic of the early days of \textit{Issylran} settlement and a panoply of forgotten wars, is all but forgotten by the people of Tal'Dorei. Few explorers travel these wind-scarred peaks, or the storm-beaten northwestern coast of Tal'Dorei. The citadel itself is wracked by the elements, while the shoreline and the mountains beyond it are prowled by territorial elementals that brook no intrusions.

\subsubsection{Secrets of the Citadel}
The ruined citadel that overlooks this treacherous landscape is far more than meets the eye. In truth, it is the headquarters of the \textit{League of Miracles}. The power-hungry spellwrights who make up the league's upper echelons call this castle the Palace of Wonder. A crumbling but entirely illusory edifice has replaced the derelict fortress, maintained by the league to deter curiosity and avoid notice.

Beneath this arcane façade, the Palace of Wonder truly lives up to its name. The palace is made up of five wings, each of which has a dedicated \textit{teleportation circle}:

\subparagraph{Halls of Splendor}
The \textit{League of Miracles} is in the business of good first impressions. Though the Palace of Wonder is a secret to the people of Tal'Dorei, the spellwrights occasionally teleport particularly powerful mercenary mages to their home base in order to dazzle them with the castle's splendors. The palace's ground floor is made of a number of interconnected grand halls, awash in a tide of splendid, vibrant colors—bordering on garish—that crashes over all who enter.

These many grand dining, entertaining, and performance halls are gleaming spectacles of white marble, decorated by sparkling globes of magical light, and festooned with self-indulgent tapestries of the spellwrights' many dubious achievements. Any intruders here without the express permission of at least one spellwright will find themselves beset by a small army of \textbf{animated armor} and other constructs—including a \textbf{platinum golem} that watches over the front entrance.

\subparagraph{Bowers of Recovery}
The second floor of the palace is a verdant paradise, with vibrant plants curling around the castle's stately stone architecture. Most spellwrights have their quarters in this wing, which is protected by powerful enchantments that give life to the plants. All who enter here without a spellwright escort are viewed by the living plants as a tasty snack.

This wing of the palace also includes a three-tiered conference chamber where all twenty-four spellwrights can gather to discuss matters of import. In the center of the chamber is a pedestal, upon which sits a humanoid mask of gleaming mithral. When the \textit{Wonderworker}, the league's mysterious leader, wishes to deliver a proclamation to the spellwrights, their voice emanates from this mask.

\subparagraph{Towers of Farsight}
The pinnacle of the palace is made up of four sky-scraping towers that gaze out over the \textit{Cliffkeep Mountains}—and all of Exandria. These private studies are the personal quarters of four of the league's chief spellwrights, and they are equipped with all manner of \textit{scrying} crystals and tools of clairvoyance, clairaudience, and far-speech.

\subparagraph{Dungeons of Despair}
Woe betide the prisoners of the \textit{League of Miracles}, who are left to rot in the cells beneath their palace. The warden of this dungeon is an infallible \textbf{mage hunter golem}, and each cell is equipped with shackles that inhibit the spellcasting capabilities of anyone in their grasp. Instruments of unspeakable arcane torture litter these dungeons, capable of instilling endless pain, and even extracting the very soul from one's body. The shackled spirits of those who have died by this last vile torment serve eternally as the guardians of the palace dungeons.

\subparagraph{Cradle of the Adranach}
Winding about the palace dungeons are a network of twisting passages illuminated by glaring, magical lights. These forking tunnels shift daily, and the secret to navigating them is a mark visible only under the power of a \textit{true seeing} spell. Within these passages are a dozen laboratories stocked with the tools needed for the spellwrights to dream up new potions and spells. Some they sell at extravagant cost to the governments and militias of towns and cities across Tal'Dorei—but the best they keep for themselves, to wield against the factions they manipulate from the shadows, should they ever grow unruly.

At the deepest point of the tunnels is the Cradle of the Adranach, a unique arcane laboratory where all spellwrights are initiated. Here, they craft an \textbf{adranach}, a construct of mithral and astral energy that serves as a loyal servant—and that unmistakably marks its creator as a spellwright of the \textit{League of Miracles}.

\subsubsection{Gathering the Spellwrights}
The \textit{Palace of Wonder} is almost never empty, for at least one of its twenty-four spellwrights must always remain within its walls to protect it. Many more spellwrights frequently return by means of the palace's \textit{teleportation circle}, or upon the back of their personal \textbf{adranach}, to take advantage of its magical equipment and nigh-impenetrable defenses.

At times, all spellwrights are requested to return to the palace at the behest of one of the four chief spellwrights, but unless one's presence is demanded by name, these summons are not mandatory. The one summon which cannot be refused, however, is the annual calling of the \textit{Wonderworker}. The nameless, faceless leader of the \textit{League of Miracles} holds a near-prophetic sway over the spellwrights that cannot be understated. When all twenty-four spellwrights are gathered within the conference chamber, the \textit{Wonderworker} intones a slew of cryptic yet unquestionable commandments that all of their followers must interpret and obey to the best of their ability.

\subsection{Pools of Wittebak}
Wittebak was once a peaceful city of three thousand or more rock \textit{gnomes}, a safe haven for the small ones that was built into the mountain rock and heated by a series of geothermal pools. The Wittebak \textit{gnomes} weathered great wars by not getting involved, investing in contraptions and laughter rather than wars and empires. They found joy in inventing complicated ways to simplify their daily lives. Yet, for all their efforts, trouble eventually came to them.

Around four hundred years ago, a displaced band of \textbf{hill giant} found and luxuriated within Wittebak's hot springs and made a homestead on their banks. It is said that Wittebak was founded underneath the full light of Ruidus, for by pure misfortune, a flaw in the stonework of the gnomish city cracked under the force of the giants frolicking within the pools. Worse still, the cave-in did more than simply crush a few buildings beneath. The seismic force of thousands of tons of rock collapsed into the caverns, cracking open buried pockets of volcanic gas.

Over a thousand \textit{gnomes} suffocated that night, or were incinerated in isolated gas explosions. The survivors fled into the mountains, assailed by hungry giants as they ran. The \textit{gnomes} have found unlikely refuge in \textit{Kraghammer}, but Wittebak has been abandoned entirely, choked by caustic gas and locked in a terrible moment in time. The ten-times great-grandchildren of the giants that caused the \textit{gnomes} to flee now call the hot springs their ancestral home. To them, it is named Gundergrunch, and they flourish within their unchallenged territory.

Beneath them, the ruins of Wittebak are stratified into three distinct tiers of death and despair.

\subparagraph{Tinkerer's Tier}
This level is relatively unharmed, and the full glory of \textit{Wittebak's} tinkering shops can be seen here; some even include marvels of technology unseen for centuries. Creatures that need to breathe can survive the gases in the widest, uppermost level for a number of hours equal to their Constitution modifier, and take 3 (1d6) poison damage each minute after that time. The \textbf{hill giant} that live within have adapted to the gas and do not fear it.

\subparagraph{Homesteads}
The air in the middle level is heavy with gas, and the quaint gnomish homes on this level are filled with the decomposed remains of those too slow to escape. The entire area is lightly obscured with clouds of gas, and tiny streams of hot water leak from the walls, tumbling into the level below.

\subparagraph{Engine Core}
The bottommost tier is a mechanical marvel, a titanic wheel of rusted iron and shining steel with no known purpose. This level is completely flooded and filled with merrow. It is also the lair of a \textbf{sea hag} named Mauvlettir.

\subsubsection{Pools of Wittebak Adventures}
Game Masters who set their adventure in \textit{Wittebak} can use these plot hooks for inspiration.

\subparagraph{Getting into Gundergrunch (mid-level)}
Just getting into \textit{Wittebak} is an adventure unto itself. It's nearly impossible to fight through the hundreds of \textbf{hill giant} in the hot springs and the Tinkerer's Tier—and even if it weren't, such a fight would be a callous slaughter of countless innocents. Stealth is an option, for the giants take little heed of folk smaller than them. Diplomacy is another option; Gondalorr, prince of the giants, has a tender heart and loves the idea of repatriating \textit{Wittebak's} \textit{gnomes}. His sister Gondabrau, on the other hand, is cruel, and would like nothing more than to pop wriggling \textit{gnomes} into her mouth as a snack. Cunning adventurers could sneak in with Gondalorr's aid, or even cause chaos in the giant town by playing the siblings against each other.

\subparagraph{The Truth of the Disaster (mid-level)}
Was it truly just by chance that \textit{Wittebak} was destroyed? Perhaps there is truth to the rumors that Ruidus, the Moon of Ill Omen, shone brightly above \textit{Wittebak} on the night of its founding. The \textbf{merrow} that lurk in the Engine Core know nothing, but the \textbf{sea hag} who lairs there has an inkling. Delve deeper, and the \textbf{aboleth} and its aberrant servitors know more still—a secret that dates back to the \textit{Founding} itself.

\subsection{Reaching Bluff}

\begin{description}
\item[Village:]
Population 324 (53\% \textit{humans}, 27\% \textit{dwarves}, 3\% dark \textit{elves}, 3\% earth \textit{elemental ancestry}, 14\% \textit{other races})

\end{description}

In a small, fertile valley in the northwest of the \textit{Cliffkeep Mountains} is a small village founded by Tal'Dorei homesteaders just ten years ago. This simple settlement is inhabited by hardy folk seeking new places to mine \textit{gold} and \textit{iron} from the earth. Many unpleasant exiles from \textit{Kraghammer} have joined them, so while this village has a charming, rustic atmosphere, few travelers stay for long. The hostility toward outsiders is palpable to all.

This insular attitude has only grown more hostile following the discovery of a gleaming dodecahedron about the size of a human head. One of the miners who unearthed it found that her pickaxe barely scratched its surface—and when struck, it began emanating a radiant, internal golden glow. This mystical object was returned to the village foredwarf, and before the day was out, all the villagers had learned of it. The people of Reaching Bluff call it the The Lost Beacon of Unknown Light, and worship it as a god.

The settlers protect their mysterious beacon with cultish fervor, and Foredwarf Daz, a deceitful dwarf \textbf{veteran} with earth \textit{elemental ancestry}, is willing to murder anyone who tries to probe into their business.

\subsection{Serpent's Head}
\textit{Thordak} only flew to the \textit{Cliffkeep Mountains} once during his occupation of \textit{Emon}. No one living knows why the nearly invincible \textit{Cinder King} left his prized city to attack a single, isolated mountaintop village, but the ruins of \textit{Serpent's Head} prove that he did.

\textit{Serpent's Head} is a squat mountain in the southern reaches of the \textit{Cliffkeep Mountains}. Its people were mostly \textit{dwarves} and \textit{humans}, miners and farmers. They had no militia to speak of, and most had already evacuated to \textit{Kraghammer} when news of \textit{Emon's} fall reached them—but not everyone.

The village that surrounds \textit{Serpent's Head} Mountain is an ashen ruin. The immolated corpses of its people still lay buried beneath a thick layer of ash, and the mountain itself was transformed into a small volcano by \textit{Thordak's} magic.

\subsubsection{Serpent's Head Adventures}
Game Masters who set their adventure in \textit{Serpent's Head} can use this plot hook for inspiration.

\subparagraph{Factions of Flame (any level)}
Characters who wish to plumb the depths of \textit{Serpent's Head} for treasure—or perhaps to find an item or character connected to the greater story of your campaign—must contend with the creatures that have made it their home. Mindless \textbf{cinderslag elemental} prowl the shattered battlements of the fortress that once stood atop the mountain. Earth tremors have cracked the keep's foundation, revealing deep magma caverns filled with diamonds—precious and infused with chaotic magic. Occasionally, tiny rifts to the Elemental Plane of Fire flash open within the depths, allowing groups of \textbf{azer} and \textbf{salamander} to wriggle through. The two factions have made camp in the lower caverns, and both struggle for a foothold in the Material Plane.

\subparagraph{Hot-Headed Wanderers (high-level)}
The \textit{Cinder King's} incursion into Tal'Dorei attracted elementally infused monsters from deep beneath the \textit{Cliffkeep Mountains} to the surface. \textbf{Magma landshark} and their \textbf{young magma landshark} are the worst of these monsters, and even decades after the \textit{Cinder King's} defeat, they continue to prowl the mountains from their nests beneath \textit{Serpent's Head}, devouring entire villages unless brave adventurers stand in their way.

\subsection{Terrah}

\begin{description}
\item[Village:]
Population 838 (61\% \textit{humans}, 19\% \textit{dwarves}, 5\% \textit{gnomes}, 15\% \textit{other races})

\end{description}

Sheltered within a cauldron-like valley, miles deep into the northern mountain range, lies the home of the Earth \textit{Ashari}. The people of \textit{Terrah} have long stood watch over the exposed rift to the Elemental Plane of Earth, protecting the ever-quaking ground of this crumbling valley. Those without druidic power dig long trenches around their village to protect their people against rockslides that roll off the mountains, while a handful of mighty druids use their magic to reshape the valley and rebuild damaged structures.

Living off hunted game and scavenged fungi, the \textit{Terrah} people are stocky warriors of stubborn mind and immense pride and loyalty. Jewelry made of precious gemstones is regarded as a sign of station and respect. The Earth \textit{Ashari} supply \textit{Emon} with precious and semi-precious stones from the rift in exchange for essential goods and, occasionally, protection from the creatures of the mountains.

The peaks that surround the rift frequently seem to shake with rage, and from their angry slopes rise creatures of living stone, imbued with life by elemental spirits from the rift. The people of \textit{Terrah} sleep in shifts to keep watch for any sudden changes in seismic activity, lest the fury of the mountains undo all that \textit{the Ashari} have created.

\subsubsection{Government}
\textit{Terrah} is governed by an \textit{Ashari} archdruid given the honored title of Heart of the Mountain. This archdruid has completed their \textit{Aramente} and earned the respect of the four elemental enclaves of \textit{the Ashari}. Pa'tice, the current Heart of the Mountain, is a human man well over one hundred years old. Though tenacious and ornery, he is a venerable warrior with boundless appreciation and support for those who prove themselves honorable and self-reliant. Rumor has it that Pa'tice is seeking a worthy druid who can undertake the Aramente and succeed him as Heart of the Mountain.

\subsubsection{Terrah Adventures}
Game Masters who set their adventure in \textit{Terrah} can use these plot hooks for inspiration.

\subparagraph{Pestering Pa'tice (any level)}
Rumor has spread throughout the \textit{Cliffkeeps} that the Earth \textit{Ashari} are seeking a new Heart of the Mountain. In response, dozens of upstart druids with delusions of grandeur have flocked to \textit{Terrah}, seeking to impress the current leader of the Earth \textit{Ashari}. Pa'tice can't be bothered with this pageantry, and quietly hires a group of adventurers—the characters—to disperse these druids. If there is a druid in the party, perhaps their actions will impress the Heart of the Mountain and set them on the path of the Aramente.

\subparagraph{Dependable as Stone (high-level)}
The \textit{Terrah} people are insular and dutiful, and rarely take sides in any conflict. However, sometimes times of great peril require that even \textit{the Ashari} go to war. Whatever the conflict, the characters must convince the \textit{Terrah} to join their cause. Pa'tice meets with them and grimly states that the \textit{Terrah} cannot aid them because a terrible new threat has come through the portal. The Stonesight Council, four \textbf{medusa} aided by two \textbf{gorgon}, now command the \textbf{earth elemental} descending from the peaks around \textit{Terrah}. Pa'tice is a practical leader—he cannot lead his people to fight a threat to the outside world while danger lurks on their doorstep. "Eliminate this threat to Terrah," Pa'tice says, "and perhaps we can join your cause."

\subsection{Umbra Hills}
The \textit{Battle of the Umbra Hills}, the climactic battle of the \textit{Scattered War} and the death knell of the Kingdom of \textit{Drassig}, transformed the floral, sun-dappled Emerald Highlands into a blasted wasteland. A tide of demonic blood spilled across the highlands like fire, cursing the land forever. Generations have passed, but the heather that once grew on these hills is as black and burnt as the day of the battle. No animals live here, and the only plant that grows in the \textit{Umbra Hills} now is shadegrass. Its ash-gray stalks are dry and unfilling, but its unique, acrid flavor has drawn the interest of spice traders worldwide, propelling wealthy merchants to hire armed escorts for their spice-pickers.

Though no animals live here, danger still lurks within the \textit{Umbra Hills}. Undead soldiers rise from the grass when the silvery moon Catha is high, and the ruins of ancient \textit{Drassig} war camps seem to draw demons of all types, still bound by the contract between \textit{Trist Drassig} and the Demon Prince of Indulgence. Pathetic \textbf{dretch} gambol about in the ruins' long shadows, and rumors say that, on nights when the moon Ruidus is full, even mighty \textbf{balor} demons prowl the hills in search of mighty souls to corrupt.

\subsubsection{Umbra Hills Adventures}
Game Masters who set their adventure in the \textit{Umbra Hills} can use this plot hook for inspiration.

\subparagraph{Sword of the Demon Prince (high-level)}
In a dream, a skilled warrior within the party continually sees visions of a noble, kingly sword embedded in the ground, surrounded by shadows and demons. "Free me," it whispers. "I need your help...I am in the \textit{Umbra Hills}." The sword is not as it appears; though the Demon Prince of Indulgence was defeated at the \textit{Battle of the Umbra Hills} and banished to the Abyss, his accursed greatsword still lingers. The last shard of his malevolence in this world, \textit{Graz'tchar, the Decadent End} rests in the center of the valley, its tip embedded in the earth. The alluring, androgynous \textbf{marilith} demons that guard it pretend to be the ghosts of great knights, protecting the blade of a lost Sovereign of Tal'Dorei.

\subsection{Grey Valley}
North of the \textit{Umbra Hills} lies a wide valley of ill omen and dark history, named for the forest of petrified trees that fill the majority of the valley's forest floor. During the \textit{Scattered War}, \textit{Trist Drassig} obsessed over the possibility of losing and pledged his service to the Demon Prince of Indulgence in exchange for power to smite his foes.

The bargain accepted, the demon prince sent his forces to this valley to aid Trist in the \textit{Battle of the Umbra Hills}, but ultimately both fell to the incredible strength of the rebelling forces. The release of terrible energies exploded across the battlefield. Now left cursed and lifeless, the valley remains a warning against dealing with demons—and a dangerous locus of darker entities that are drawn to the ruined place.

\subsubsection{Grey Valley Adventures}
Game Masters who set their adventure in the \textit{Grey Valley} can use this plot hook for inspiration.

\subparagraph{Heart of the Wraithroot (high-level)}
At the dead center of the \textit{Grey Valley} is a petrified tree unlike any other. Knots in its stony bark gaze like eyes at its withered domain. This \textbf{wraithroot tree} is the eternal guardian of a dark treasure. Rumors abound about what could be secreted away behind this vile guardian—a relic of a demon prince, the crown of \textit{Drassig}, or an artifact of the \textit{Betrayer Gods} unknown to mortal minds? The true nature of this artifact is up to the Game Master, but whatever it is, rumors of its existence have lured countless adventurers to their deaths.

\section{Stormcrest Mountains}
Spanning Tal'Dorei's southern reaches, buttressed against the \textit{Verdant Expanse}, stand the massive \textit{Stormcrest Mountains}. The winds of the \textit{Lucidian Ocean} wail through the storm-battered peaks, and few dare make these treacherous mountains their home. The few who do live amidst the Stormcrests are ferocious monstrosities and the \textit{giantkin} strong enough to defeat them—and those few peoples cunning enough to evade storm, monster, and giant alike. The only friendly towns anywhere near these forsaken peaks lie in the \textit{Mornset Countryside} pushing along the Stormpoint range, far beyond the influence of the \textit{Tal'Dorei Council}. These rugged folk are used to fending for themselves and scoff at the thought of returning to so-called "civilization," preferring life among the swamps and shadowed forests of their homeland.

\begin{DndReadAloud}{}

\begin{itemize}
\item{Majority Faiths:}
The Arch Heart, The Dawnfather, The Wildmother

\item{Minority Faiths:}
The Moonweaver, The Stormlord, The Ruiner, The Scaled Tyrant, The Spider Queen

\item{Imports:}
Livestock, lumber, grain, \textit{silver}, spell components

\item{Exports:}
Lumber, \textit{gold}, grain, granite, \textit{quartz}, \textit{iron}

\end{itemize}

\end{DndReadAloud}

\subsection{Ashen Gorge}
One mountain stands out among the Stormcrests. Its twin peaks rumble with fury, breathing fire and smoke into the sky to mingle with the storm clouds that surround the mountains. Lightning tears through the ashen clouds with giddy excitement, making it clear to all who behold it that this is a place of supernatural power. Cartographers in \textit{Emon} call this crater the \textit{Ashen Gorge}, but the people of the low valleys call it the Dragon's Throne, for they know what creature once lived atop the mountain. They remember when they too were under the dominion of \textit{Thordak the Cinder King}.

Today, several tenuous alliances of \textbf{lizardfolk}, \textit{dragonblood}, \textbf{kobold}, and even a few young dragons trying to expand their treasure hoards battle for dominion over this ruined vale. When the mountain was the lair of the \textit{Cinder King} during his first incursion into Tal'Dorei, the \textbf{lizardfolk} and \textbf{kobold} flocked to his side, worshiping his might and cruelty. Though \textit{Thordak} is slain, his influence yet lingers, and his legend draws treasure hunters and would-be tyrants to this place. Gold, jewels, and magic items from all ages of the world lie in wait within the dozen or more secret vaults of the \textit{Cinder King's} lair. When the king of dragons still held dominion here, he had thousands of loyal scaled minions at his beck and call, moving invaluable treasures into deeper and deeper holds. Beyond the main cavern of the lair, even the \textit{Cinder King} himself had to magically take humanoid form to traverse the winding tunnels his servants carved through the rock.

The descendants of \textit{Thordak's} servitors now fight an endless war over his treasure. Every year, it seems a petty new "Cinder King" rises to power within the \textit{Ashen Gorge}, but they are always dethroned before their minions can enter the ancient vaults, still locked after all this time. To steal \textit{Thordak's} greatest treasures, an adventuring party would have to either broker a peace between the warring clans or somehow sneak into the mountain undetected.

The following major factions fight for control over the \textit{Ashen Gorge}. They make and break their alliances with one faction or another, and then break their new alliances just as readily. You can use these factions to inspire adventures set in the \textit{Ashen Gorge}.

\subparagraph{Scions of Flame}
This cult of black-scaled \textbf{lizardfolk} claim to have been \textit{Thordak's} elite guard in ancient times, and are the current rulers of the gorge. They're led by a five-hundred-year-old druid named Burning Oak, and while his awesome magic keeps the other tribes at bay, his frailty makes him vulnerable to attack—and the other tribes know it.

\subparagraph{Tinysoot}
Every ten years brings a new generation of \textbf{kobold}, and while the Tinysoot tribe's warriors are not the mightiest, their overwhelming numbers have allowed them to overwhelm even the greatest of the Flame Scions' champions. The Tinysoots' leader, a stubborn young Queen of Soot named Yabber Tinysoot XIX, wears around her neck the key to Everflame Crevasse, \textit{Thordak's} deepest treasure vault.

\subparagraph{Black Snow}
When \textit{Thordak's} ally, Umbrasyl, was slain atop \textit{Gatshadow}, his caustic blood seeped into the corpses of a dozen goliath scouts from the Herd of Storms. The black dragons' magic, combined with the foul power of \textit{Gatshadow} itself, transformed the \textit{goliaths} into a squadron of undead black \textit{dragonblood} warriors. Skeletal wings sprouted from their backs, and the Black Snow tribe flew instinctively to the \textit{Ashen Gorge}. They were drawn to the Ruby of Oblivion, a pitch-black ruby of immense unholy power—now stored in the Obsidian Geode, a treasure vault guarded by a legion of golems and elementals. This squadron of hateful revenants is small, but each one of them wields the strength of twenty lesser warriors.

\subparagraph{Tatterwing Kin}
Toxoshuul, Chloroxodon, and Snagglefang, a trio of \textbf{young green dragon}, long to plunder \textit{Thordak's} lost vaults for their own hoards, and demand tribute from the other factions of the \textit{Ashen Gorge} in exchange for their aid in battle. All three are the children of Tatterwing, a feared dragon of the Kirmont Valley, and even if they find themselves on opposite sides of a skirmish, they never fight one another, only annihilating the pathetic humanoids beneath them with their toxic breath and razor-sharp talons.

\subsection{Dreamseep Marshlands}
East of the Stormcrest Range, within the Kirmont Valley, is the sprawling, fetid Dreamseep Marsh. The perpetual rain rolling off the Stormcrests' enchanted slopes has transformed this once-lush forest into a morass of rotting vegetation, sulfurous mud, and gnarled, weeping trees. So many cruel murders have been committed within the Dreamseep that the very land is cursed, forsaken by the gods. Negative energy pools there like water, and the dead drink deep of it. The marsh is populated by more than just the walking zombies of killers and their victims—awful amalgamations of dozens of bodies, both humanoid and bestial, are birthed within the Dreamseep's putrid womb.

Somewhere in the middle of this accursed realm is a sinkhole that plunges deep beneath the earth. At its deepest point is the Tomb of Udah, a fabled necropolis of countless chambers, littered with traps and treasures that have claimed the lives and imaginations of untold hundreds. Every death within Udah's accursed walls only adds to the legions at its dread master's command; characters who explore its tunnels face not only undead warriors in armor from a bygone millennium, but steely warriors in armor from every epoch in Tal'Dorei history.

The small community of Bronbog holds the only lights of civilization within the shadows here, and those that continue to thrive against the oppressive swamp are hearty and stubborn folk..

\subsubsection{Dreamseep Marshlands Adventures}
Game Masters who set their adventure in the Dreamseep Marshlands can use this plot hook for inspiration.

\subparagraph{Skull of Udah (epic-level)}
Two rumors persist of the Tomb of Udah. First, it is said that within the tomb's deepest sanctum is a three-eyed golden skull, and that in each socket is inlaid a ruby that can grant any wish. Second, the locals say that, unlike the mindless undead of the \textit{Dreamseep} above, the abominations in the necropolis below are commanded by a hideous master. Little do these rumor-spreaders know that the skull and the master are one and the same. Udah the Undying is a \textbf{lich} of tremendous power, devouring the souls of all who die within the \textit{Dreamseep}. His disembodied spirit controls the dungeon itself, and only retreats into the skull—the last remnant of his physical form—when it is in danger.

\subsection{Bronbog}

\begin{description}
\item[Village:]
Population 683 (53\% \textit{humans}, 18\% \textit{half-elves}, 11\% \textit{goblinkin}, 9\% \textit{halflings}, 9\% \textit{other races})

\end{description}

If there ever was a civilization that called the \textit{Dreamseep} home, Bronbog is all that remains of it. Its buildings are made from planks of waterlogged bogwood, and the only sign of its previous greatness is the ring of stone pillars that once supported the Temple of the The Dawnfather. Even though the temple to their patron god has crumbled, the people of Bronbog keep the faith; their belief in the sun god's afterlife is their greatest comfort against their meager lives. Yet, for all the gloom that surrounds the \textit{Dreamseep Marshlands}, the Bronboggi maintain a determinedly optimistic disposition. There's a saying in the village: "If you don't feel the storm, you can't know there ain't sun."

The only reason anyone north of the \textit{Stormcrest Mountains} knows of Bronbog is because of queenscap, a rare swamp fungus that serves as a reagent in creating potent potions of superior healing. It's been harvested almost to extinction in the more accessible \textit{K'Tawl Swamp}, and Tal'Dorei's alchemists will pay good money for a shipment.

The biggest building in Bronbog is its tavern, the Queen's Knight, which has about six beds in drafty rooms to accommodate the town's very few visitors.

\subsubsection{Bronbog Adventures}
Game Masters who set their adventure in Bronbog can use this plot hook for inspiration.

\subparagraph{Morning Glory (mid-level)}
On a queenscap-gathering expedition to Bronbog, an unusual storm forces the characters to stay the night. As they spend time in the swamp town, vine-choked stone obelisks slowly emerge from the murky waters, emanating an unearthly light. Some villagers immediately fall under the obelisks' sway, claiming they are "heralds of the Dawnfather, the rising sun," but all others are disturbed by their neighbors' feverish worship of the stones.

The true power behind these obelisks is a misguided \textbf{deva} that calls himself Eclipse. He seeks to create a brainwashed, perfectly loyal Army of the Dawn to fight against the The Dawnfather enemies. He is posing as Father Larkylai, the town's affable priest of the The Dawnfather, and conducts services from within the town's tavern.

\subsection{The Frostweald}
Along the northern base of the Stormcrests is a forest locked in perpetual winter, cursed to eternal cold by the invasion of \textit{Errevon the Rimelord} during the \textit{Icelost Years}. Within the forest, ponds magically freeze into perfect mirrors, reflecting the snowy sky above, and the snow-draped trees shelter families of fey hiding from enemies in the Fey Realm. At first, the \textit{Frostweald} seems a wonderland of crisp snow and aromatic pines and firs, yet the serene landscape belies sinister danger. Herds of \textbf{basilisk} roam the woods, and travelers who encounter mysterious snow-covered statuaries or copses of petrified trees must think fast or join this grisly statuary—permanently.

\subsubsection{Frostweald Inhabitants}
The only humanoid peoples who live within the \textit{Frostweald} are hunter-gatherers called the Shivergut. Many of their people are \textit{orcs} who have lived within the \textit{Frostweald} for generations, but their ranks also include \textit{humans}, \textit{elves}, and \textit{dwarves} who adopted their ways centuries ago during the \textit{Icelost Years}. Most are helpful to travelers, if reserved. Others, however, are dour, selfish isolationists who have been apart from other folk for so long that they care more for the thrill of bloodshed than a warm conversation over a crackling fire.

The \textit{Frostweald} is also home to a host of benevolent fey, most of them pixies or dryads, using the forest as a safe haven away from the Fey Realm. If pressed for answers on why they left the Fey Realm, their answers are cryptic and dissatisfactory, but always ominously suggest a "Great Shadow" has descended upon their lands—perhaps even a host of warriors from the Plane of Shadow? A nymph named Arethusa watches over a cluster of three mirror-like pools that form a pathway into the Fey Realm, each to a different Archfey's forest. She is suspicious of all mortals, and both her trust and a favor are required to gain passage.

\subsubsection{Ruins of the Frostweald}
Many forgotten obelisks of the The Knowing Mentor rest secretly beneath the ice and snow. These lost beacons were created to guide the faithful of the The Knowing Mentor to the \textit{Cavern of Axiom}, a sacred place guarded since time immemorial by a \textbf{androsphinx} dedicated to the The Knowing Mentor. If any living creature knows why these ancient monoliths reside here when no other ruins of the The Knowing Mentor have been discovered, they have kept the knowledge secret from the other scholars of the world. Over thirty years ago, the half-orc \textit{Emonian} archaeologist Jorlund Vohr discovered a cache of Ioun stones buried here and brought them to the \textit{Alabaster Lyceum} of \textit{Emon}, but his research and findings were stolen the week after he returned. Vohr has since stated he "got over it," and is back to work on new research on \textit{Visa Isle}.

\begin{DndSidebar}[float=!b]{The Silver Tablet}
\DndQuote%
{}
{''Ye who seek my boundless wisdom,''}
{''Pass into the white trees' kingdom.''}
{''Return to the shrine Mentor's eyes,''}
{''And follow the light in the skies.''}
{''I am the Beast of the Heavens;''}
{''Thou shalt kneel before my presence,''}
{''Lest the wisdom thou seekest''}
{''Be lost to the cowardice of your weakest.''}
{}


This poem was discovered by Jorlund Vohr on his first expedition to the \textit{Frostweald}, on a silver tablet written in flowing Celestial script. Upon examination, scholars at the \textit{Alabaster Lyceum} discovered the language of the heavens is so sublime that its poetry rhymes even when translated into any language. The tablet was stolen later that week, along with the rest of Vohr's findings, but Jorlund keeps a parchment-and-charcoal copy of the tablet in a chest beneath his bed in \textit{Emon's} \textit{Erudite Quarter}.



\end{DndSidebar}

\subsection{Cavern of Axiom}
The Cavern of Axiom is a lost shrine to the The Knowing Mentor, spoken of only in riddles. Hidden by ever-shifting illusions, the entrance to this cavern opens only to those who are expected by fate. Its entrance looks different to every group that finds it, from the imperious to the humble, but it appears simply as a heavy snowdrift to those not fated to open its doors. Within its shifting facade are dangerous trials designed to test the will and mettle of those who seek the infinite knowledge of the The Knowing Mentor. The chambers plunge deeper into the rock beneath the mountain, while half-eroded murals and strange puzzles evaluate any wanderer who seeks an audience with the keeper of the cavern, an ageless \textbf{androsphinx} named Kamaljiori. The final challenges presented by the \textbf{androsphinx} alter from subject to subject, and failure banishes the petitioners from the cavern, barred from ever returning.

\subsection{Ruhn-Shak}

\begin{description}
\item[Small City:]
Population 6,670 (80\% \textit{elves}, 12\% \textit{dwarves}, 5\% \textit{gnomes}, 3\% \textit{other races})

\end{description}

If you find carved arches and steel gates in the mountain slopes, do not rest there, no matter the cold. There are no grand dwarf-halls among these forsaken peaks.

Deep below the surface world, countless caverns and tunnels wind into regions where light finds no purchase. It is here, beneath the \textit{Stormcrest Mountains}, that the largest subterranean society of the The Spider Queen devotees maintains its tyranny. Hidden entrances and tunnels riddle not just the mountain range, but the edges of the \textit{Dreamseep}, as well as darker regions of the \textit{Verdant Expanse}, granting raiding parties quick and easy access to the surface under cover of night—slaughtering many, capturing the rest, and taking the spoils as gifts to the The Spider Queen.

These tunnels, in their twisted network, are easily collapsed and reopened through the use of "Pit Witches," druids who master the art of rock and dirt manipulation, making it near impossible to give chase. The \textit{Wardens of Syngorn} are ever seeking a way to find and destroy this heart of the The Spider Queen power in Tal'Dorei, but the sly and cunning \textit{elves} of Ruhn-Shak have yet to meet their match.

\subsubsection{Ruhn-Shak Adventures}
Game Masters who set their adventure in Ruhn-Shak can use these plot hooks for inspiration.

\subparagraph{Dynastic Intervention (any level)}
While in \textit{Syngorn}, the characters are called by \textit{Ouestra, Voice of Memory}, to meet with guests from \textit{Wildemount}. They are emissaries of the Kryn Dynasty, a society founded by dark \textit{elves} who fled the The Spider Queen grip and now live on the surface. These Kryn envoys seek to undertake an expedition into \textit{Ruhn-Shak}, to try to use the radiant power of their mysterious god, the The Lost Beacon of Unknown Light, to free their kin from the The Spider Queen evil influence. The characters are offered a fine reward by the \textit{Wardens of Syngorn} to protect these Kryn priests on their mission.

\subparagraph{A Plea from the Dark (mid-level)}
The characters are drawn into the underground tunnels around \textit{Ruhn-Shak} by a psychic message: "I have displeased the Spider Queen; there is nothing I can do but beg for salvation—can anyone hear me?" The message comes from Ry'denyan, a \textbf{drow} of \textit{Ruhn-Shak} granted telepathic powers by an unknown aberration. This unexpected gift allowed him to break free of his society's brainwashing. He tried to escape, and was exiled into the tunnels, where he will surely be devoured by \textbf{drider} unless the characters save him.

\subsection{Wrettis}
A relic of the \textit{Age of Arcanum}, Wrettis is the ruined tower of a powerful mage, one who was driven mad by the seductive whispers of beings from the beyond. What few legends survive of the mage of Wrettis say he was known as Clemain Astural, the Sight Shephard, and that he was a powerful arcanist who peered into a realm beyond the planes in search of power to end the war that ravaged his world. He found it, and thought it would serve him.

He was wrong. Astural's sanity crumbled, but his power only grew, fueled by an entity he called the The Sightless. As devastation crept across southeastern Tal'Dorei, heroes of the land rode to Wrettis to end Astural's chaotic reign, destroy his tower, and bury his corpse in the rubble. Wrettis is now a moss-coated ruin, but some chambers within and below the tower still hold secrets of the mad mage, as well as some of his creations.

\begin{DndSidebar}[float=!b]{The Astural Scrolls}
Clemain Astural covered thousands of pages of parchment with his writings, ranging from alchemical formulas to powerful spells to esoteric star charts...but his later notes descend into paranoid, half-crazed ramblings. One of the heroes that killed Astural, a wizard named Atz Yuminor, took as many of the scrolls as he could find and brought them back to his house in the \textit{Dividing Plains}. Yuminor's bloodline continues in Estella Ladimar, headmistress of the \textit{Westhall Academy in Westruun} in \textit{Westruun}, and many of the Astural Scrolls are kept secretly by the \textit{Scions of Yuminor in the academy's observatory}.





\end{DndSidebar}

\section{Rifenmist Peninsula}
South of the \textit{Stormcrest Mountains} and the \textit{Verdant Expanse} are a number of autonomous communities that live outside the influence of the Republic of Tal'Dorei and its constituent city-states. A number of plans to expand into \textit{Rifenmist} have been brought to the \textit{Tal'Dorei Council}, but strong opposition from the \textit{Syngornian} delegation and the Master of Defense have consistently buried these plans before they can come to fruition.

A number of communities exist on the \textit{Rifenmist Peninsula}. Some are lost remnants of the \textit{Scattered War} that chose to flee \textit{Drassig's} rule to the isolated and unpopulated \textit{Mornset Countryside} rather than fight against his tyranny. These are the communities that some \textit{Emonian} merchants wish to bring into the fold of Tal'Dorei proper. Farther to the south (and occasionally roaming north into Mornset) are the \textit{Orroyen elves}, a collection of nomadic tribes. The \textit{Orroyen} and the \textit{elves} of \textit{Syngorn} see one another as family, and the \textit{Syngornian} delegation refuses to allow the council to interfere with their cousins' way of life.

\begin{DndReadAloud}{}

\begin{itemize}
\item{Majority Faiths:}
The Changebringer, The Lawbearer, The Strife Emperor

\item{Minority Faiths:}
The Wildmother, The Moonweaver, The Cloaked Serpent, The Lord of the Hells

\item{Imports:}
Arms, armor, stone, masonry, spell components

\item{Exports:}
Spell components, produce, rare metals

\end{itemize}

\end{DndReadAloud}

Beyond the Mornset people and the \textit{Orroyen} are a number of other small, individualistic settlements that live in fear of the imperial power that commands Rifenmist: the \textit{Iron Authority}. Tal'Dorei's Master of Defense warns the council of provoking this slumbering empire, for they are preoccupied with their unholy conquest of the \textit{Rifenmist Jungle} in the name of their patron deity, the The Strife Emperor. A number of \textit{Orroyen elves} and other people of \textit{Rifenmist} long for \textit{Emon's} aid in breaking the grip of the \textit{Iron Authority}, but fear that allowing the Republic of Tal'Dorei to gain a military foothold in \textit{Rifenmist} would simply be trading one imperial master for another.

Beyond the political struggles of this realm are terrors only true heroes can survive. Enormous, ancient beasts stalk the vine-twisted paths of the jungle that consumes the majority of the peninsula's coastline. An expanding morass of fungus spreads from a mysterious, corrupt source along the eastern depths of the underbrush, while a secret society of naga-worshippers bring bloody offerings to their snake-queen.

\subsection{Beynsfal Plateau}
Where below there is life and freedom, above there is naught but death and tyranny. In the high, southern plateaus of \textit{Rifenmist}, verdant jungle gives way to a cracked, volcanic landscape, its basalt fields interrupted only by towering structures of black iron. This hellish place is the home of the \textit{Iron Authority}, and one of the last battlegrounds of the \textit{Calamity}.

It is here that the Betrayer God known as the The Strife Emperor, aided by his legions of brainwashed goblinkin, clashed with the The Wildmother and her Free Children. Though the The Wildmother was victorious and the goblinkin were freed, the site of the The Strife Emperor defeat was reduced to rock and ash, a place where plants would never again grow.

Most of the goblinkin who survived the battle vanished into the jungle and traveled north into the rest of Tal'Dorei, leaving the shame of their connection to the The Strife Emperor behind. Centuries later, however, a number of goblinkin still remained among the towering iron of the The Strife Emperor fallen armor. One platoon of \textbf{hobgoblin}, the Iron Regiment, devoted themselves wholly to the The Strife Emperor, and pledged to revive his glorious legacy of conquest. They press-ganged as many of their kin as they could into armies and created the seed of an empire—one which they watered with blood.

When the Iron Regiment brought all of the Beyensfal Plateau under their rule and transformed themselves into the \textit{Iron Authority}, their empire only continued to grow. The Iron Tide never broke, and within a decade, it had engulfed all the other kingdoms along \textit{Rifenmist's} southern coast. They were united under the cruel banner of Tz'Jarr, the Iron Emperor.

While the entirety of Beynsfal is firmly controlled by the \textit{Iron Authority}, the empire is still reaching, now locked in a decades-long war to subjugate the native inhabitants of the jungles below: the \textit{Orroyen}.

Five major city-states comprise the core of the \textit{Iron Authority}, though all swear fealty to the Iron Emperor in its capital of \textit{Tz'Arrm}, the southernmost city of the plateau. From south to north, the other imperial city-states are Rybad-Kol, city of the Forge Lords; Hdar-Fye, the necropolis of Prince Hdar; Ezordam-Haar, eyrie of the red pegasi; and Ortem-Vellak, the petrified elventree.

\subsubsection{Beynsfal Plateau Adventures}
Game Masters who set their adventure on the \textit{Beynsfal Plateau} and the surrounding mountains can use this plot hook for inspiration.

\subparagraph{The Wayward Bloodline (mid-level)}
Tz'Jarr has a single heir to his throne: his \textbf{hobgoblin} child Zuun'dak, nearly of age and currently being groomed to become a great general. However, Zuun'dak has grown up with trepidation regarding the brutality of his society, and after being briefly captured by an outlying \textit{Orroyen} scouting party and learning of the gentler world abroad, he's begun to long for escape and reformation. When the party stumbles upon a disguised and fleeing Zuun'dak, they must decide whether to turn over such a valuable bargaining tool to the \textit{Orroyen}, or aid and groom the defecting heir as a rebellious force against his father.

\subsection{Mornset Countryside}
After a long journey south from \textit{Syngorn}, or north from the \textit{Rifenmist Jungle}, thick foliage gives way to rolling hills and winding rivers. No great cities are to be seen here, only an untouched wilderness pockmarked by fortified towns. The Mornset Countryside is home to many people whose ancestors fled the rule of \textit{Drassig}. Most of these people value their privacy—especially from \textit{Emonian} landscape painters who enrich themselves by romanticizing and exoticizing their idyllic way of life.

The Free Folk of Mornset, as they call themselves, include a wide variety of people who have spurned the ways of their forebears. The bulk of these people are the descendants of refugees from a long-ago war, but a number of Tal'Dorei people who long for freedom from the pressure of modern life have uprooted their lives and journeyed to Mornset to start anew. Likewise, some \textit{Orroyen elves} who dislike the nomadic lifestyle find homes in the small communities of Mornset.

Beyond those who follow their hearts to this realm, Mornset is also home to a number of criminals fleeing persecution. Most Free Folk don't ask questions about newcomers' pasts, and welcome them regardless. However, if these particular outlanders don't leave their past behind and continue to harm others, or bring bounty hunters (or worse) in their wake, every village in Mornset has a well-trained militia that will eagerly remove dangerous folk from their midst.

Bound by the necessity of survival, the people of this land share a common code of honor called the Fundaments. This code holds all to a baseline of respect for all travelers and outcasts, and codifies ancient principles of hospitality. Each denizen of the countryside looks out for their neighbor, and if a person in need arrives upon one's doorstep, the hospitality of food and shelter must be provided.

It is understood that the Fundaments may be taken advantage of by ruffians and power-hungry upstarts. Nevertheless, this code of honor has served the Free Folk well for centuries, and anyone who scorns it is looked upon with disdain and mistrust. Most of the Mornset Free Folk are humanoids, but many free settlements have alliances with the local \textbf{hill giant}, \textbf{cyclops}, and \textbf{ogre} that help protect them against the \textit{Iron Authority} to the south.

\subsubsection{Mornset Adventures}
Game Masters who set their adventure in Mornset can use this plot hook for inspiration.

\subparagraph{Anger Everburning (low-level)}
A number of residents of the uncharted village of Heldenfaire have been unexpectedly—and suddenly—flying into a rage both mindless and destructive. Many within the village are confused, fearing a plague of violent madness, and they seek adventurers (especially parties with clerics or skilled healers) to save them. Investigation by the party reveals that each individual driven to these fits is connected to the town alchemist's dark past and a failed herbalism experiment from ten years ago, when she lived in \textit{Kymal}. If the characters wish to find an antidote, they must explore the alchemist's sinister laboratory in the hills.

\subsection{Byroden}

\begin{description}
\item[Town:]
Population 8,050 (45\% \textit{humans}, 21\% \textit{half-elves}, 12\% \textit{goblinkin}, 8\% \textit{halflings}, 8\% \textit{gnomes}, 6\% \textit{other races})

\end{description}

The northernmost and largest town in the \textit{Mornset Countryside}, \textit{Byroden} borders the \textit{Gladepools} and marks the first stop for any traveler on their way to the \textit{Rifenmist Peninsula}. It's also a hub of trade for the entire \textit{Mornset Countryside}, and even some trappers and artisans of the \textit{Verdant Expanse}. The town has a dense town center surrounded by palisades and tall wooden walls, but most of the community lives in homesteads and farmhouses spread across the surrounding fields. \textit{Byroden} is mostly self-sufficient, and has a wealth of farmers, fishers, artisans, and miners, giving them everything they need to survive. Trade provides the opportunity for some luxuries, with raw materials serving as their primary currency.

While pleasant and welcoming on the outside, this lifestyle is protected by a fierce and intense drive to defend the community's way of life at the first sign of threat or danger. Most citizens are armed at all times, trained to rise as a communal militia at a moment's notice from the war ringer, a sentinel charged with the town's protection.

Despite (or perhaps because of) the ever-present dangers that threaten \textit{Byroden}, the town holds frequent festivals, feasts, and fairs, including the highly anticipated Gem of \textit{Byroden} pageant and the midsummer Peachiest Pie bake-off. On festival days, small children roam freely within the protective town walls as their parents enjoy the chance to let loose and celebrate life in all its ephemeral glory.

\subsubsection{Hillmaw Chasm}
A few miles north of the center of \textit{Byroden} is a gaping chasm. An earthquake thirty years ago split open Deercrest Hill. The chasm is referred to officially as Deercrest Ravine, but is known to most superstitious folk in town as Hillmaw Chasm. After the earthquake, locals stayed as far away from the chasm as they could—until four years later, when a group of adventurers found vast veins of precious \textit{gold} and sturdy \textit{iron} within. Soon, all people of \textit{Byroden} strong enough to swing a pick or lift a shovel made their way into the chasm.

Most miners hire an armed escort because the chasm is known to be home to grumpy \textbf{earth elemental}, as well as unpleasant, slimy creatures of the dark. Most serious beliefs about the chasm being haunted have dissipated over the past thirty years. Still, rumors persist, and as folklore becomes myth, some say that myth can become reality.

\subsubsection{Advent of the Cinder King}
About forty years ago, \textit{Byroden} was brought to its knees by \textit{Thordak the Cinder King}. This incursion was about two decades before the \textit{Cinder King's} destruction of \textit{Emon}, and it was stopped by an adventuring party led by \textit{Allura Vysoren}, now a senior member of the \textit{Tal'Dorei Council}. Many in \textit{Byroden} still see \textit{Allura}, Lady Kima of Vord, and \textit{Drake Thunderbrand} as heroes. The graves of the companions they lost in that fight, Sirus Kaldrem, Dohla Lorian, and Ghenn Talevesh, hold places of honor in the town cemetery.

While \textit{Byroden} has long-since been built anew, the terror of that day still lives on in its elders and adults. All but the youngest generation of \textit{Byroden} takes great caution against another such monstrous attack—as well as against any attack by the \textit{Iron Authority} to the south.

\subsubsection{Byroden Adventures}
Game Masters who set their adventure in \textit{Byroden} can use these plot hooks for inspiration.

\subparagraph{Ironbound Scouts (low-level)}
All are awoken in \textit{Byroden} by the ringing of the town's alarm bell, including the characters. As the town militia is armed for battle, the characters overhear that \textbf{goblin} and \textbf{hobgoblin} scouts of the \textit{Iron Authority} have been spotted twenty miles south of town—closer than they've been seen in nearly two decades.

Worried that this scouting party is an omen of invasion, the town's war ringer deputizes the characters and offers them a handsome sum of money to investigate the scouting camp. Their goal is simple: learn the interlopers' mission, either by interrogating their commander or by finding a written command note. And, if the characters can capture or eliminate the scouts before the town's upcoming festival, the locals will celebrate their success with all the free food and drink they can stomach.

\subparagraph{The Dreamgate (mid-level)}
After a mining expedition goes missing deep within the Hillmaw Chasm, locals begin to report vivid, communal dreams involving members of the missing party pleading for help. A half dozen mercenaries have been hired to investigate—only for them to disappear as well. Just a day before the characters are asked for help, the townsfolk of \textit{Byroden} discovered with horror that the vanished mercenaries have also appeared in their dreams. The adventurers are asked to delve into this unearthed pathway and to find both answers and the missing people.

\subsection{Niirdal-Poc}

\begin{description}
\item[City:]
Population 15,442 (18\% \textit{humans}, 17\% \textit{goblinkin}, 12\% \textit{elves}, 12\% Loxodon, 10\% \textit{half-elves}, 7\% \textit{elemental ancestry}, 24\% \textit{other races})

\end{description}

Long before the \textit{Calamity}, a great civilization called the Qoniira made \textit{Rifenmist} their home. During the \textit{Age of Arcanum}, the Qoniira Tetrarchy grew from a loose coalition of towns led by arcanists, clerics of the The Wildmother, and druids who spoke with the jungle around them. These towns were inhabited by \textit{humans} and the animal-bodied humanoids of the jungle. The first stones of the Qoniiran heart-city Niirdal-Poc were laid by the broad, towering Loxodon, and places for new settlements were scouted out by \textit{Catfolk} seers whose eyes glistened with starlight.

\subsubsection{Surviving the Calamity}
In time, the Qoniira civilization grew into a Tetrarchy, with four great satellite cities and the central heart-city of \textit{Niirdal-Poc}. The tetrarchs established their cities with mathematical precision. A constellation of the night sky guided them—Surrac's Shield, a four-starred diamond with a fifth star in its center. The constellation still shines over \textit{Rifenmist} and southern Exandria today. The four outer cities of Qoniira match the points of Surrac's shield precisely, and \textit{Niirdal-Poc} is aligned with its center star.

The \textit{Calamity} struck at the peak of Qoniiran civilization, yet rather than take up arms against gods and mortals, the tetrarchs opted to remove themselves from the world—not unlike the \textit{elves} of \textit{Syngorn} to the north. Something happened in \textit{Niirdal-Poc}, some sort of magic that even the Qoniirans of today do not understand.

\subsubsection{Geography}
Today, only the heart-city of \textit{Niirdal-Poc} survives. Its satellites are little more than crumbling ruins. Something shifted in the magic of the jungle when the Qoniirans chose not to fight in the \textit{Calamity}. The \textit{Rifenmist Jungle} closed its boughs against the world, shielding the city. When the The Strife Emperor died on the \textit{Beynsfal Plateau} to the south, the Qoniirans knew nothing of it. All the tumult of the world passed them by.

The \textit{Rifenmist Jungle} sunk \textit{Niirdal-Poc} into a deep basin, and its thick canopies shield it from above, while its intertwining boughs hold the heart-city tightly in its embrace. Few are permitted by the protective, ancient magic of the Qoniirans to enter the city. Its people live contentedly within \textit{Niirdal-Poc}, only occasionally traveling outside their basin to trade with \textit{Orroyen nomads} or, in mythically rare instances, with the town of \textit{Byroden}.

The four ruins which surround \textit{Niirdal-Poc} are:

\subparagraph{Niirdal-Gan}
This outer city to the north of \textit{Niirdal-Poc} was the agricultural hub of Qoniira. Some farmers have opted to live outside the basin and tend to these still-fertile lands. They are few, for magic is able to supply food for most people of \textit{Niirdal-Poc}.

\subparagraph{Niirdal-Teek}
This outer city, east of \textit{Niirdal-Poc}, was an academic center of the Qoniiran civilization. Most of the city is one vast university that once was populated by over 30,000 people. Some explorers venture into these ruins to find lost knowledge.

\subparagraph{Niirdal-Sarqet}
Located south of \textit{Niirdal-Poc}, this outer city was a center of faith for the Qoniirans. All that remains of the sacred city is a vast, floating cube called the Hexahedron of Sarqet. It is adorned with statues of the \textit{Prime Deities} and the \textit{Betrayer Gods} alike, and none living know what is inside its unbreachable walls. Even today, modern Qoniiran priests teach that balance is the true way of the world, telling their people to invoke the names of the \textit{Prime Deities} for blessings, and of the \textit{Betrayer Gods} for curses.

\subparagraph{Niirdal-Hup}
Precious little remains of this fortress-city, the heart of the Qoniiran military. A common parable reminds those who hear it to not fight fate—for of all the ruined cities, only the city of warriors was reduced to dust.

All of these ruins are filled with statues from the height of the Qoniiran Tetrarchy. Some statues have been known to move under the light of the stars, illuminating hidden runes upon their bodies as they patrol like sentinels of antiquity.

\subsubsection{Government}
The Tetrarchy of Qoniira, in ages past, was ruled by four tetrarchs, each of whom ruled one of the civilization's four outer cities. In times where the fate of Qoniira was at a crossroads, they gathered within the heart-city to guide the course of their civilization together.

The relative autonomy of the four outer cities and the looseness of the coalition that bound them seems at odds with the perfect precision with which the cities were built, so precisely matching the stars in the sky. The Qoniirans have always been an easygoing people who typically follow the flow of fate, rather than fight against it.

Today, the tetrarchy survives, even though the four outer cities are in ruins. They gather daily in the Hakredic Dome, an architecturally magnificent granite rotunda adorned with an outer layer of grayish-purple porphyry stone and lifelike, reddish porphyry statues. There, the tetrarchs and their advisors work to keep life simple and content for their people.

The tetrarchy is a hereditary office, but citywide elections are held between eligible heirs to determine which heir should take up their parent's mantle. Elections are also held when a tetrarch's line ends without heirs—often giving advisors a chance to gain power by virtue of their service to the city.

\subsubsection{Magic Beyond Reckoning}
Magic suffuses the city of \textit{Niirdal-Poc}. Its magic is something older than the gods, or the \textit{Calamity}, or anything wielded by mortals today. However, the mortal mages of \textit{Niirdal-Poc} wield the same arcane force gifted by the gods, and its druids and clerics touch upon the power of the gods behind the Divine Gate, just as Exandria's other mortal peoples do.

Both ancient and modern magic mingle in \textit{Niirdal-Poc}, for the patterns of the ancient powers are mysterious, and its methods are unknowable. Some call it the Light of Fate, others the Blessing of the Stars. The most common name is simply "the Gift." It was the Gift that protected \textit{Niirdal-Poc} in the \textit{Calamity}. It is the Gift that causes miracles to occur, like a person surviving a fall from the heights of Vedrim's Pinnacle to the streets of the city. It is the Gift that causes the Hexahedron of Sarqet to hover mysteriously; it is the Gift that surrounds the Statue of The Observer, a four-winged feline humanoid figure in the city's center, with a broken halo of ashen-gold light—and that causes that same light to flicker within the lamps of the city.

\subsection{Orroyen Tribes}

\begin{description}
\item[Nomadic Tribes:]
Population 11,850 (64\% \textit{elves}, 12\% \textit{half-elves}, 10\% \textit{humans}, 14\% \textit{other races})

\end{description}

\begin{DndReadAloud}{}
\DndQuote%
{}
{''Storms bring death, but rain brings life. Storms are power, the power of the distant gods, on display for all to see. Climb into the Stormpoints, young one, and shout your fury into the relentless wind and booming thunder. Your prayers may not be answered by the gods, but you will feel them echo in your heart.''}
{Cloudwriter Ix}
\end{DndReadAloud}

Born from a wandering colony fleeing the \textit{Calamity}, these wood \textit{elves} have not just survived within the jungles of \textit{Rifenmist}—they have flourished. Learning the dangers of the surrounding lands, refining their hunting techniques, and maintaining a healthy arms-length relationship with outsiders, the Orroyen call the jungle their domain and protect it furiously.

The Orroyen tribes have no central government, and their membership shifts as different clans move through the region and encounter other tribes, sometimes trading members as they go. Elders are given the most oversight and respect among a traveling tribe, the eldest given the title of dura, and the duty of leading their tribe. Each nomadic group maintains a population of anywhere between a hundred and a thousand people. Temporary lodging is constructed at each resting point, called a \textit{tomenda}.

Most Orroyen groups remain at a tomenda they've constructed for about two years, before depletion of resources or the arrival of dangerous creatures or soldiers forces them to move on. Countless abandoned tomenda can be found throughout the jungle, and Orroyen clans will frequently make camp in old tomenda made by a different tribe—it is considered a gift from the jungle.

To come into prominence within the tribe, or request admittance into the tribes as an ally or member, a hopeful entrant must complete a series of ceremonial trials. These trials change from dura to dura, but are designed to test strength of body and mind. It is not uncommon for those who attempt these trials to return maimed, or to not return at all, and many elders bear the marks and scars of their trials long past. Few outlanders from the north, known as orfindes to the Orroyen, have been allowed into any of the dozens of Orroyen tribes, and the elders have little intention of changing this, for they have heard distant tales of the ancient tyranny of \textit{Drassig} and the supernatural horrors the people of northern \textit{Gwessar} invite upon themselves.

Nevertheless, the Orroyen are practical, not prideful. They trade often with the Free Folk of Mornset, and sometimes even travel as far north as \textit{Byroden} to trade rare goods. There, some Orroyen meet people from \textit{Syngorn} or beyond and share stories of their lands. Love can bloom in these meetings of separate worlds, and sometime results in an orfinde being brought into an Orroyen clan. And sometimes curious Orroyen leave simply to see the world beyond the jungle.

\subsection{Rifenmist Jungle}
The jungles that cover Tal'Dorei's southern expanses are a vast wilderness. Many \textit{elves} call this lush landscape their home, particularly those of the nomadic \textit{Orroyen Tribes}. There is no doubt that this jungle is as magical and mysterious as Tal'Dorei's temperate forests, for the jungle's greenery is home to countless mystical spirits, ferocious monsters, and arcane mysteries. The vine-strangled floors of the jungle are clustered with swollen trunks, spine-laden ferns, and carnivorous plant life that waits for the unknowing wanderer to make an incorrect step.

The sweltering heat of the sun-baked peninsula grows even more oppressive beneath the jungle's humid canopy. Clusters of venomous insects buzz through the mist-shrouded trees, alongside large lizards, and wild beasts in search of carrion. Sunken valleys lead to pockets of marshland or quicksand pit traps, while others are wrought with a dire, spreading fungus. Those who have lived all their lives in this environment know its dangers well and have hundreds of years of precautions to take against them. On the other hand, those who venture into the jungle from northern Tal'Dorei rarely return, and those that do warn others to never make the same mistake.

The farther south one travels into \textit{Rifenmist}, the denser the jungle grows. Though some \textit{elves} and \textit{humans} of northern Tal'Dorei have attempted to colonize the jungles for their vast resources in decades past, they were quickly rebuffed by the \textit{Orroyen elves} of the jungle. Their message is clear: no cities are to be built south of the jungle's edge—and if the northern jungles are cleared or burned, the trees' vengeance will be swift and just. The \textit{Tal'Dorei Council} reacted quickly as well—the \textit{elves} of the jungle are to be heeded. The few settlers who have joined the Free Folk of the \textit{Mornset Countryside} are deeply superstitious about the southern jungles, and rarely travel deeper than they need to in order to parley with their \textit{Orroyen} allies.

Even the \textit{Iron Authority}, who has been at war with the trees of the jungle itself for decades, has not managed to destroy their foe. The \textit{Orroyen nomads} fear the conquerors of \textit{Beynsfal Plateau} and the massive armies supplied to them by their authoritarian leaders. Their raids ebb and flow as imperial leaders rise and fall. Each new era begins with a bold and cruel warlord rallying their people against their age-old enemy. War is waged, and countless innocents and Authority soldiers perish in a senseless fight. And then resources run low, morale dips, and the legions return to the \textit{Beynsfal Plateau} to lick their wounds.

Legend holds that the jungle will never die, so long as even one of the \textit{Orroyen} lives within it. So far, history has only proved the jungle's enduring strength.

\subsubsection{Rifenmist Jungle Adventures}
Game Masters who set their adventure in the \textit{Rifenmist Jungle} can use these plot hooks for inspiration.

\subparagraph{The Insatiable Sanctuary (mid-level)}
Long have rumors stirred regarding a vine-obscured temple from an epoch long past, but reports have mapped it at many different locations, with returning expeditions finding it missing. This temple has seemed to reveal itself once more in the northern region of the \textit{Rifenmist}, its shadowed passages and secret bounties ripe for the plucking. It isn't until a venturing band delves deeper into the structure that signs begin to reveal the shrine itself to be a living entity, and the treasure a lure to draw in prey.

\subparagraph{Seed of Immolation (high-level)}
Deep beneath the twisted floors of the jungle lies a long-slumbering relic from before the \textit{Founding}: an unborn primordial titan seed, its energy responsible for the fertile landscape and extreme overgrowth of the jungle. The \textit{Iron Authority} and their power-hungry clerics of the The Strife Emperor have tampered with the seed and unleashed powers beyond their control. The titan has begun to develop and wake from its eons-long slumber. If the process isn't halted, it will fully awaken and terrorize the continent. However, if the seed is destroyed, vast swathes of jungle may wither and die without its magical essence.

\subsection{Stormpoint Mountains}
The southern offshoot of the \textit{Stormcrest Mountains}, the Stormpoint Mountains push deep into \textit{Rifenmist Jungle} itself. These steep peaks are shrouded in low-sitting clouds and fog for most of the year, their rocky surface hidden from view by lush overgrowth that's fed by the near-perpetual rain.

This region is beautiful, and in spring, thousands of waterfalls spray mist and cast gorgeous rainbows across the crags. However, due to its distance from most of Tal'Dorei's population centers, few but adventurers and treasure hunters come here—seeking wealth, not beauty. Some Free Folk from the \textit{Mornset Countryside} and pilgrims of the The Wildmother do take great pains to travel here and revel in the Stormpoint Mountains' natural splendor.

\subsubsection{Stormpoint Mountains Adventures}
Game Masters who set their adventure in the Stormpoint Mountains can use this plot hook for inspiration.

\subparagraph{Allies of the Storm (high-level)}
The people of the \textit{Mornset Countryside} and the free people of the \textit{Rifenmist Jungle} long for an end to the tyranny of the \textit{Iron Authority}. A group of freedom fighters who have heard of the characters' legendary feats of heroism contact them with a request: help. They lack the funds to request the characters fight the Authority themselves, but they wish for the characters to recruit a clan of \textbf{cyclops} that live in the Stormpoint Mountains, including the mighty \textbf{cyclops stormcaller} that lead their people. Once there, the characters learn that their clan is beset by demons who have possessed their \textbf{cyclops stormcaller} and are using their mighty powers to conjure a deadly storm to keep the rest of the \textbf{cyclops} in submission. Only exorcising the demons and saving the \textbf{cyclops stormcaller} will earn the characters their aid.

\subsection{Tz'Arrm, Helm of the Emperor}

\begin{description}
\item[City:]
Population 47,400 (70\% \textit{hobgoblins}, 10\% \textit{goblins}, 20\% \textit{other races})

\end{description}

\begin{DndSidebar}[float=!b]{Content Warning: Slavery}
The \textit{Iron Authority} is a despicable society that embodies all the evils of ruthless authoritarianism. This includes the subjugation of innocents, racial supremacy, poverty, and slavery. People at your gaming table may be more sensitive to this type of realistic evil than villains who simply wish to destroy the world. It's worth discussing with your players in a \textit{session zero} if everyone will still have fun (and be comfortable) if the game contains this type of evil before introducing the \textit{Iron Authority} as a villain in your campaign.



If you and your players want to use the \textit{Iron Authority} in your game, but don't feel comfortable with the depth of its evil, you can strip out these elements. What remains is a classic "evil empire" that wants to conquer Tal'Dorei. Its conquest can serve as a backdrop to the other adventures that your players want to go on.



\end{DndSidebar}

When the The Strife Emperor fought the The Wildmother, he took the form of a giant of unimaginable size, clad from head to toe in armor of pitted iron. When he was defeated and cast back into his planar prison, his immense, divine armor remained and crashed to the ground. Centuries have passed, and the thirty-foot tall Helm of Strife is now the palace of the Iron Emperor. The rest of their imprisoned god's armor fuels the unstoppable war machine of the \textit{Iron Authority}, melted down to create walls, weapons, and worse.

\subsubsection{Government}
The armies of the \textit{Iron Authority} are strictly regimented, and its society is no less ruthlessly structured. Every \textbf{goblin}, \textbf{hobgoblin}, and \textbf{bugbear} knows their place within the imperial hierarchy. And every being of other races, including a number of \textit{elves} who have betrayed the \textit{Orroyen} in hopes of gaining power within the Authority, know their place beneath the goblinkin. Iron Emperor Tz'Jarr is the supreme military commander of the empire, loyally served by the four royal generals that rule the empire's individual city-states. Tz'Jarr is unimaginably old, and it is said that his clerics have used blood magic granted to them by their god to sustain him beyond his years.

The military government of the \textit{Iron Authority} is headquartered in Tz'Arrm. Its cultish generals are doggedly focused on maintaining the Authority's war machine—and their own personal power. Beneath these generals are innumerable propagandists, military officials, slave-owning nobles, and their cronies, all of whom ravenously seek a way to climb the ladder of power. No elections are held; positions are only vacated in the event of death or promotion, and all leaders are appointed by the emperor or an immediate superior.

\subsubsection{Society}
As the capital of an empire with an insatiable appetite for expansion and conquest, Tz'Arrm is awash with propaganda and crawling with secret police. In addition to being the headquarters of the military, Tz'Arrm is also the center of operations for the Gauntlet, the empire's intelligence and espionage division. As the children of the The Strife Emperor, loyal goblinkin are seen as the supreme citizens of the empire. The upper echelons of the \textit{Iron Authority} are filled almost exclusively by \textbf{hobgoblin}, for the Iron Emperor has declared them the strongest and most worthy children of their god. By the same token, goblinkin who betray the empire are seen as the ultimate traitors, and suffer a fate worse than death.

Those with wealth and status within the empire are constantly paranoid about falling out of favor. Even great generals see all those around them as hungry wolves waiting to gobble them up. At the bottom rungs of society are the serving class, who eke out a meager existence as indentured servants for the elite, or as merchants who are bound to whatever master owns their shop. Beneath even them are slaves, most of whom are prisoners of war.

\subsubsection{Slavery}
The armies of the \textit{Iron Authority} have captured tens of thousands of slaves during their wars of conquest across \textit{Rifenmist}. The slaves of Tz'Arrm are mostly used as laborers, sent to chisel iron scraps from the The Strife Emperor fallen armor, or to venture back into the lower jungles to forage for food. The elite of the \textit{Iron Authority} sometimes play at benevolence by granting comfortable jobs indoors to slaves skilled in the arts—such as personal portrait-painters, musicians, or sculptors. Such kindnesses are rarely a sign of true mercy; they are more likely a simple (and meager) expression of guilt.

Groups of slaves that pass in and out of the city to labor in the armor yards or furnaces are escorted by a military detachment whose primary purpose is to discourage escape. The powerful people of the city prize their elven slaves above all others, passing them down like heirlooms from generation to generation. \textit{Elves} are treasured possessions by the Authority's despicable elite, not just because of their beauty and long life, but because of how difficult it is to procure slaves from the sheltering boughs of the jungle.

\subsubsection{Geography and Climate}
Tz'Arrm is hellishly hot. The black stone of the \textit{Beynsfal Plateau} and the iron walls of the city trap the heat of the relentless \textit{Rifenmist} sun. The imperial elite enjoy spells that cool their homes and shelter them from the sun, but those who toil under the sun pay a heavy price. The Authority's \textbf{hobgoblin} elite are well suited to such a climate, but their \textbf{goblin} and \textbf{bugbear} slaves suffer as they toil—as do the \textit{Orroyen elves}, whose life in the shady, humid boughs of the \textit{Rifenmist Jungle} has left them unprepared for the blistering heat of the plateau.

\section{Verdant Expanse}
East of the \textit{Stormcrest Mountains}, hundreds of miles of massive, unbridled forest shroud the landscape in mystery and shadow. This dense greenwood was for centuries traversed by few but the \textit{elves} of Tal'Dorei, for it was known to be their domain, and theirs alone. Even in the time of \textit{Zan Tal'Dorei}, when her rebels and their elven allies saved the realm from tyranny, it was understood that the coasts were humanity's domain, and the woods were that of elvenkind. Times have changed, however. The alliances formed to oppose \textit{Thordak the Cinder King} united virtually all of Tal'Dorei's disparate peoples, and many more than simply \textit{elves} now walk the shaded paths of the \textit{Verdant Expanse}. Even so, much of the Expanse is protected by the watchful eyes of the \textit{Wardens of Syngorn}, as it has been for hundreds of years.

\begin{DndReadAloud}{}

\begin{itemize}
\item{Majority Faiths:}
The Arch Heart, The Wildmother, The Moonweaver

\item{Minority Faiths:}
The Matron of Ravens, The Knowing Mentor, The Stormlord, The Cloaked Serpent, The Spider Queen

\item{Imports:}
Precious gems, grain, \textit{gold}

\item{Exports:}
Lumber, elvish goods, weapons, jewelry, meat and vegetables, enchantments

\end{itemize}

\end{DndReadAloud}

Merchants travel from \textit{Emon} to \textit{Syngorn} along roads sanctioned for trade, hunters both human and elven alike seek glory in tracking the wild beasts that roam the untamed wilderness, and bands of dark elf raiders from \textit{Ruhn-Shak} prowl the forest under cover of night. This enchanted wood holds the decay of the seasons at bay, and its trees are green year-round.

The \textit{Verdant Expanse} is sustained by a massive confluence of ley lines—rivers of magical power that flow through the earth—after the devastation of the \textit{Calamity} shifted the flow of magic across the world. The region is saturated with magic. Magical creatures flock to the forest's supernaturally vibrant boughs, including wayward fey, \textbf{unicorn}, unusual arcane monstrosities like \textbf{owlbear}, and displaced aberrations that call the darker groves their home.

\subsection{The Gladepools}
Surrounding the southern edges of the Expanse lies a cluster of broken ponds and lakes that ever draw fresh water from the \textit{Stormcrest Mountains} to the east, along with saltwater from the Ozmit Sea to the west. The strange mix of habitats across this somewhat marshy grassland has led to unusual ecosystems and odd, dangerous denizens, which in turn has drawn the attention of many fishermen from the \textit{Rifenmist} region, and of \textit{Syngornian} hunters from within the \textit{Verdant Expanse}.

The clay and silt gathered along the shores of the saltwater lakes can be refined into fine ceramics and simple constructs that fetch fair prices in northern Tal'Dorei. \textit{Syngornian} tradition tells of an oracle's spirit that is bound to the lakes, and when given sufficient offering and respect, the oracle will emerge to grant sufficiently reverent supplicants a cryptic prophecy.

\subsubsection{Gladepools Adventures}
Game Masters who set their adventure in the \textit{Gladepools} can use this plot hook for inspiration.

\subparagraph{No Basis for a System of Government (any level)}
While traveling along the edge of the \textit{Gladepools}, an elf character, or a character with elven ancestry, notices a gleaming \textit{longsword} embedded in the bank of the lake. No one but that character can see it until they draw it from the mud. Flowing elven script upon the blade reads, "Whosoever draws this blade bears the mantle of the true scion of Yenlara."

Attempts to appraise the sword and discover if it's a true ancient \textit{Syngornian} artifact plunge the characters into a dark conspiracy that lurks at the heart of \textit{Syngorn}: The High Warden may be wise and just, but she is not the rightful ruler—at least, she is not Yenlara's heir. Centuries ago, someone strove to subvert Yenlara's bloodline. Why? And, if the High Warden is truly a good ruler, does the bloodline deserve to be restored?

\subsection{The Mirescar}
Deep within the western reaches of the \textit{Verdant Expanse} is a region where the trees grow so close together that sunlight is barely permitted to pass through their boughs. These ancient trees cluster and twist into a labyrinth of knotted roots and branches, while a tangled canopy of drooping foliage and gray moss blocks the sky from view. Strange powers dwell within the Mirescar, and beasts and travelers that wander into its depths often emerge with their hearts and flesh touched by a deep and corrupting evil—if they emerge at all. The woodlands surrounding the Mirescar are rife with dire beasts, cruel fey, and many-eyed aberrations.

The Mirescar is the subject of countless \textit{Syngornian} folk tales, many of which purport to tell the true story of how it was formed. Some say it contains the bones of a titan. Others say that a meteor from Ruidus itself landed in the forest. Another, that the The Spider Queen cursed it during the \textit{Calamity}. Any one of the tales may be true—or none of them. Only adventurers of great courage and skill have the power to penetrate deep enough into the Mirescar to learn the truth and retain their humanity.

More often than not, these tales serve to frighten the people of \textit{Syngorn} and \textit{Byroden} away from the region altogether.

\subsubsection{Mirescar Adventures}
Game Masters who set their adventure in the Mirescar can use these plot hooks for inspiration.

\subparagraph{A Sickness Spreads (mid-level)}
A plague has begun to spread from the Mirescar. People living in the woods around this evil realm are suffering as they fall blind, with milky-white eyes; or their bodies are wracked by fevers; or they are struck by a number of other mysterious symptoms. All are aware that the vile Mirescar is the source of the \textit{contagion}. The characters are hired to delve into its boughs and find a cure, or to at least stem the flow of disease.

In ancient times, the Mirescar was called the Heartwood. It was the domain of a wise and kindly \textbf{guardian naga} named \textbf{Sagacitous Erusaire}, who welcomed all with good hearts to share their bounty or receive alms. However, this \textbf{Sagacitous Erusaire} has been corrupted into a cruel shadow of his former self. He and his \textbf{diseased grick} minions have been twisted into hideous, pus-oozing spreaders of disease—and spreading it is now the only thought in their minds. The \textbf{Sagacitous Erusaire} can cast \textit{contagion} at will, and his \textbf{diseased grick} minions can do so once per day. Is it possible to save the mind of this corrupted guardian?

\subparagraph{Tree of Scarred Knowledge (high-level)}
A \textbf{wraithroot tree} dwells at the heart of the Mirescar. It is covered with deep gouges, all of which open and close like eyes and mouths as it watches and, on occasion, speaks in a hissing voice to its surroundings. It stands atop a fell artifact—the very thing that turned the Mirescar to evil in ancient times—and has drunk deep of this artifact's forbidden knowledge and its cruelty.

The tree is excited to bargain with any creature that meets it, trading secrets of unknowable origin for the memories held most dear. It fears no attack, for it is protected by a small legion of \textbf{invisible stalker} that enact its will throughout the Mirescar.

\subsection{The Shifting Keep}
The Verdant Guard has a secondary base of operations, beyond \textit{Syngorn's} walls, called the \textit{Shifting Keep}. It is built within and around a cluster of massive, living trees that form the heart of the \textit{Verdant Expanse}, and enchanted with the most powerful illusion magic that could be mustered by \textit{Syngorn's} military arcanists and their guild of Spellbenders. Enchantments woven into the keep's foundations allow it to cloak its presence, seemingly vanishing and reappearing elsewhere within the forest at will.

The \textit{Shifting Keep} can't teleport, but it can turn invisible and project an illusion of itself to any point within a mile of its true location. This remarkable enchantment has made it almost impossible for invaders to besiege the keep, and it played a major role in \textit{Syngorn's} victory against Neminar \textit{Drassig} in the \textit{Scattered War}—tales of the unassailable \textit{Shifting Keep} are told in \textit{Syngorn} to this day. Nearly a thousand trained elven warriors, hunters, and spies call the \textit{Shifting Keep} their home, though many spend their time patrolling and exploring the \textit{Verdant Expanse}. These rangers are vigilant, seeking signs of whatever foolish invaders would dare trifle with the might of \textit{Syngorn}.

\subsection{Syngorn}

\begin{description}
\item[City:]
Population 38,540 (85\% \textit{elves}, 9\% \textit{humans}, 3\% \textit{halflings}, 3\% \textit{other races})

\end{description}

Founded by the sorceress Yenlara in the wake of the \textit{Divergence}, \textit{Syngorn} became the heart of elven civilization on Tal'Dorei following its calamitous destruction. The remnants of the fallen Court of Ullusa fled to the Fey Realm for a generation following the end of the \textit{Age of Arcanum}, returning only when Yenlara found the realm safe to return. These four hundred survivors built a new home, taking inspiration from their temporary home among the fey, and formed the foundations of \textit{Syngorn}. Though \textit{elves} have spread all across Tal'Dorei, even as far as the frigid \textit{Neverfields}, there was once a time where nearly all of them resided in \textit{Syngorn}.

This beautiful city of curving stone is built in harmony with the trees of the \textit{Expanse} and sits against the western base of the \textit{Stormcrest Mountains}. It is all but impervious to external assault—not simply because of its 40-foot-high walls of ivy-covered \textit{jade}, and not only because it's also surrounded by living trees of the \textit{Verdant Expanse} and beacons of detection that keep constant vigil for intruders. No, \textit{Syngorn} is virtually unassailable because every entrance to the city is warded by a series of threshold crests: massive emblems of a crescent moon, flanked by two trees, over a deep cerulean stone. These enchanted stones act as an anchor to the Fey Realm, where \textit{Syngorn's} founders once recovered, allowing the city itself to vanish tracelessly into the Fey Realm as a final act of preservation.

\begin{DndReadAloud}{}
\DndQuote%
{}
{''All here will say to you that Syngorn is a city without equal on Exandria. Its towers and bowers are impeccable, yes. Its cloisters and arches so graceful and fine they make your heart ache, yes, it has been said. But our dreams and memories are the most precious treasure here. They flicker like fireflies through the night, and they are heard in the laughter of the youngest elven child, and in the reminiscing groans of the eldest of our kind.''}
{Ouestra}
\end{DndReadAloud}

\subsubsection{Government}
The city has been governed since its inception by the \textit{High Warden}, a hereditary monarch who appoints proven individuals to three other offices alongside them as the \textit{Wardens of Syngorn}. Each Warden carries a title and bears a responsibility to guide the city toward safety and prosperity. The Verdant Lord is the head of the city guard, though they delegate this responsibility to a Vice Protector when leading the armies of \textit{Syngorn} to war abroad.

The Guildrunner manages the city's treasury and oversees commerce within \textit{Syngorn's} borders. The \textit{Voice of Memory} is a heralded keeper of history and culture, and is often seen interacting with the elven people, gathering new memories for the archives. The High Warden has always been of Yenlara's bloodline, and is responsible for keeping order within \textit{Syngorn} and its territories, as well as within the \textit{Wardens} themselves.

\subsubsection{Society}
\textit{Syngorn} is steeped in elven tradition that dates back to before the \textit{Divergence}. The arts are lauded and revered, the pursuit of knowledge is respected and encouraged, and some training in refined martial techniques is culturally expected. The idea of trade between other cities and nations is understood to be both healthy and beneficial, but most foreign trade and travel is relegated to outer areas of the city.

There was a time that few foreigners ever saw \textit{Syngorn's} innermost reaches, but that time has passed. \textit{Syngorn's} gates are open to travelers from across Tal'Dorei—though all who visit the city are bound by \textit{Syngorn's} strict cultural laws to respect their millennia-old traditions. Any who disrespect the ways of \textit{Syngorn} are warned but once, and are thrown out of the city thereafter. Delicate crafts by \textit{Syngornian} hands are sought after by collectors around the world, so many \textit{Syngornians} take up the trade for both profit and honor.

In a populace made up of people who are extraordinarily long-lived, having children is a rare and highly-regulated process. Prospective parents must gain approval from the office of the \textit{Voice of Memory} to procreate, and any unapproved children are sent out of the city to be raised in outposts or foreign cities. All \textit{elves} call such an occasion "regrettable," and truly do treat it with sorrow—but the laws of \textit{Syngorn} are clear on this point, and few \textit{elves} are willing to speak out against rulings which have stood for a thousand years.

\subsubsection{Crime}
If \textit{Syngorn} appears free of crime, it's because elven criminals have hundreds of years of experience at lurking in the shadows. A thriving market for illegal goods runs invisibly through the city's guilds. Though time often seems to stand still in \textit{Syngorn}, crime moves with the same celerity as it does in other cities—it must, if the criminal underbelly wants to stay one step ahead of the \textit{Wardens of Syngorn}.

The most common contraband to pass through \textit{Syngorn} are shipments of \textit{suude}, stolen goods, and dwarven trinkets. Many come to \textit{Syngorn} to purchase these goods for a low price on the black market, but those unwilling to make a journey and risk being stopped by the Verdant Guard can wait for these goods to make their way up to \textit{Kymal}, where they can be purchased at an outrageous markup of five times their cost in \textit{Syngorn}.

Most people in \textit{Syngorn} have their basic needs taken care of by the city's government, to ensure the wellbeing of their people. Nevertheless, some people still fall into poverty, and these destitute individuals are often preyed upon by pickpockets, swindlers, and smugglers who seek to use them as patsies for a paltry sum of gold.

\subsubsection{Geography and Climate}
The perpetual shade of the thick, green canopy of the forest maintains cooler weather throughout the summer months, while this far south, the ice and snow of winter is present, but minimal.

\textit{Syngorn} is broken up into six main districts and has a variety of neighborhoods:

\subparagraph{Memory Ward}
The mind of \textit{Syngorn} is contained within a ring of marble walls in \textit{Syngorn's} northeast reaches. A winding stair, the One Thousand Steps, leads up to its lofty gate. Standing in the center of its grand court is the mystical Sequoia of Remembrance, a three-hundred-foot-tall redwood within which are stored the lives and memories of every elf in Tal'Dorei. At least, those whose bodies are brought before the sequoia and are granted the Rite of Remembrance, as performed by the \textit{Voice of Memory} and her disciples, have their memories preserved and added to the collective. Nearly all \textit{elves} in \textit{Syngorn} receive this rite, and a vast number of \textit{elves} who live elsewhere write into their wills that they must be subject to the rite after their death.

Wisps of violet light flit about the sequoia's ancient boughs, and those who look closely into their centers have been known to catch glimpses of eras past, troubles present, and things that yet may be. The great redwood and its spirits are watched and tended to by the Dreamweavers, \textit{elves} who have dedicated their waking and sleeping lives to the protection of their heritage. About five hundred Dreamweavers and five thousand scholars, merchants, and other \textit{elves} call this district their home.

The \textit{elves} are a culture obsessed with tradition and the past. As such, the Dreamweavers are among the most well-respected factions within \textit{Syngorn}. Their leader, \textit{Ouestra, the Voice of Memory}, is likewise the most revered and socially powerful of the \textit{Wardens}.

\subparagraph{Beryl Keep}
This fortress-district houses the martial might of \textit{Syngorn}. Situated on a hill in the northwest of the city, thick walls of leaf-green beryl separate the district from the rest of the city. Within these walls are rows of barracks, archery ranges, trance chambers, forges, mess halls, and all the other necessities to train an army of long-lived \textit{elves}. At the far northwest of the district is the Beryl Keep itself. The near-unassailable fortress proves the saying that even the most utilitarian of elven things are beautiful to mortal eyes. Six eucalyptus trees of unknown age mark the boundaries of the keep, and between them grows a thick curtain of ironbark. A dense canopy of impenetrable steelfern forms its roof. Legend says that when \textit{Syngorn's} need is greatest, the six ancient eucalyptus trees that guard the keep will uproot and march as mighty \textbf{treant} to defend the city.

Beryl Keep is commanded by \textit{Verdant Lord Celindar}, \textit{Syngorn's} master strategist. He possesses a \textit{crystal ball} of \textit{true seeing} within his war room, which he can use to \textit{scrying} upon any part of the \textit{Verdant Expanse} and to pierce any illusion. From this vantage point, he commands his forces like a chess master, always one step ahead of the threats within the enchanted forest. So far, \textit{Syngorn} remains safe because its enemies are wild and disorganized. Were they to unify, the might of the Verdant Guard would truly be put to the test.

The armies of \textit{Syngorn} are about five thousand \textit{elves} strong, though only about three thousand are within this district at any time. The rest are stationed at outposts within the \textit{Verdant Expanse} and the Fey Realm.

\subparagraph{Tarn Ward}
\textit{Syngorn's} central ward is a peaceful commerce district surrounding Lake Ywynnlas and split by the channels that feed it. From sunrise to sunset, no business is conducted here, save for food service at public houses and beds at inns. During the daylight hours, \textit{elves} socialize, play, sing songs, write poetry, paint great works of art, and meditate here. \textit{Elves} are long-lived, and if half the day is not spent in peace and self-improvement, it is thought that the day was truly wasted, no matter how productive one might have been in matters of business.

At nightfall, however, the Tarn Ward changes. The Tarn Thoroughfare opens at moonrise, lit only by the heavens and by magical, floating lanterns. The streets are silent but for the ethereal elvensong that guides those who know how to listen through the streets. \textit{Elves} and half-elves raised in \textit{Syngorn} know this musical language, but those unfamiliar with it must use a \textit{comprehend languages} spell to follow the melody. When the river of sound guides a customer to a vendor, they speak in hushed Elvish, like an audience whispering at the theatre.

Most shops along the Tarn Thoroughfare are affiliated with one of three guilds—the Spellbenders' Guild, the Elvencraft Alliance, or the Mithral Fellowship—and the sly \textit{Guildrunner Rawndel} sits at the head of all three. The Warden of this district, Rawndel lives in a permanent \textit{Mordenkainen's Magnificent Mansion} with a façade adorned with splendid fey grotesques, and is respected by elven merchants and nobles alike. The \textit{elves} don't seem to mind his monopoly on Thoroughfare trade, as their lives haven't been affected much, but \textit{High Warden Tirelda} worries of the long-term consequences of Rawndel's power-hungry behavior.

\begin{DndSidebar}[float=!b]{Syngornian Goods}
Any basic goods can be purchased in \textit{Syngorn} at twice their usual cost. These goods, however, are of fine elven make. Weapons have a +1 bonus to damage rolls, horses and other mounts increase their movement speed by 10 feet, and other items are possessed of unearthly beauty and preternatural durability. All common magic items can be easily found here, and any uncommon magic item has a The item is nowhere to be found. chance of being found here in the ward's many market stalls. Finding magic items of greater rarity requires a character to make a DC 21 Charisma (Investigation) check to find a rare item, or a DC 25 (Investigation) check to find a very rare item. These items are rarely for sale, and can only be bartered for—typically by completing a dangerous quest for the person who possesses them.



\end{DndSidebar}

\subparagraph{Emerald Citadel}
The great palace of \textit{Syngorn} looms high over the north of the city—a mighty castle of faded white marble, now covered with climbing ivy and sprouting plants. As the highest point in the gradually sloping city, its emerald-tipped spires can be seen from anywhere below. Visitors to the citadel first climb a grand set of stairs, curved like a flowing river of marble, before reaching its brass gates. The Verdant Guard aggressively protects the castle's main entrance, rejecting any commoner who does not have an invitation to the palace marked with the High Warden's seal. \textit{High Warden Tirelda's} monocled majordomo, a silver-haired elf named Ibbimas, screens all of the High Warden's supplicants.

The palace's high, vaulted ceilings are supported by pillars of gleaming marble carved in the image of tall elm trees, such that the ceilings are a canopy of polished stone. Beautiful portraits and busts of past \textit{Wardens} line the walls, along with opulent tapestries of Yenlara and the creation of \textit{Syngorn}. The citadel is made up of four main levels. The basement levels hold the castle dungeons, as well as grand vaults containing fabulous treasure, historical artifacts, and overwrought gifts—the latter given by emissaries uncertain of how to present a gift fine enough for an elven monarch. The first floor holds lush quarters for ambassadors and visiting dignitaries. The second floor's chambers are dedicated to business of state, including the High Warden's throne room. The towers above the citadel contain studies and quarters for the High Warden and her family.

\subparagraph{Feygrove}
A once-splendid manor house in southern \textit{Syngorn} is now completely overgrown with plant life. Massive mulberry trees sprout at odd angles out of windows and through gables, and fairy lights dance around the house at all hours of the day. When \textit{Syngorn} returned from the Fey Realm after the \textit{Conclave's} defeat, the \textit{elves} unintentionally brought a little of the Fey Realm with them. A number of fey creatures slipped through a gate from the Fey Realm to Exandria, then wriggled through \textit{Syngorn's} defenses and made their home in the mansion of an elven noblewoman named Lady Ladri Il'shavfa. She has tried for years to reclaim her home, but even the patience of \textit{elves} wears thin; the fey are simply too numerous and too tenacious to be kicked out for good.

The fey interlopers, however, are having the time of their lives. Far from merely throwing nightly parties—though they do that, too—they have made Il'shavfa Manor a place for them to experiment with wild new magic. Their leader, a \textbf{pixie} prince named Windybranch, has found great pleasure in planting strange fey plants in the house and watching as they integrate with the mansion. The house is now a living, breathing, thinking thing—a new friend for the fey! It keeps the \textit{elves} out and it's a great conversationalist! To the fey, this arrangement has no downsides! The fairies have named this manor after the Archfey who first allowed them into the city. To them, their new home is Artagan's Lodge (and it's said that the The Traveler occasionally stays within its walls when he visits Tal'Dorei).

Since the Feygrove was first "founded" over twenty years ago, the fey have spilled out of Il'shavfa Manor and made an entire district of the city into their home. By this point, only the most stubborn of elven nobility still gripe about "the fey menace" within their midst. Most of \textit{Syngorn's} residents are delighted to have a district filled with the fey, even if it causes the southern end of the city to be a little more chaotic than usual. The fey here seem to have kind hearts, and their chaos rarely crosses the line from churlish pranks into cruelty.

\subparagraph{Reverie Walks}
The artistic and spiritual center of \textit{Syngorn} is a winding labyrinth of living trees and stones. \textit{Elves} who seek a day-long (or days-long) journey of meditation may wander the Reverie Walks, entering a trance-like state as they do so. Those who seek true enlightenment may dedicate years of their near-immortal lives to wandering the ever-shifting labyrinth, finding peace in isolation or in search of the Stone of the The Arch Heart in the labyrinth's center. An elven monk named Lyssev Sorveline has walked the labyrinth for five hundred years, seeking answers to questions even the gods do not know. Some \textit{elves} leave gifts of food for the Wanderer within the labyrinth, allowing them to continue their eternal meditation in peace.

No Warden holds sway over the Reverie Walks. Legend holds that the living stones and trees that shift the labyrinth are devotees of the The Arch Heart that swore to be their god's eternal wardens.

\subsubsection{Points of Interest}
\textit{Syngorn} has a number of famous landmarks. They are keyed to the map of \textit{Syngorn}.

\subparagraph{1. Elvencraft Alliance Guildhall}
The Elvencraft Alliance is supposedly the oldest guild in \textit{Syngorn}. They are entrusted with creating works of art that spread the culture of \textit{Syngorn} far and wide across Tal'Dorei. The front rooms of their guildhall are a museum of fine elven crafts from ages past, and they're open to all who wish to appreciate \textit{Syngorn's} cultural heritage.

\subparagraph{2. Guildrunner's Mansion}
\textit{Guildrunner Rawndel}, Warden of the Tarn Ward, makes his home within a permanent \textit{Mordenkainen's Magnificent Mansion}. Its exterior is a magnificent façade of stone and leering fey statues, yet if one were to look inside, they would find it completely empty. The door, which is constantly locked to all but the Guildrunner and his guests, is actually a portal to an extradimensional space containing a manor of supernatural beauty and luxury.

\subparagraph{3. Mithral Fellowship Guildhall}
The master metallurgists and smiths of the Mithral Fellowship are tasked with arming the Verdant Guard and the armies of \textit{Syngorn}. In peacetime, they turn their talents to creating things of beauty from metal and jewels, sometimes in tandem with the artists of the Elvencraft Alliance. Due to their military involvement, none but authorized members of the fellowship are allowed within their guildhall.

\subparagraph{4. One Thousand Steps}
This winding stair is the only way in or out of the \textit{Memory Ward}. Each step is inscribed with an ancient elven idiom, aphorism, or impossible philosophical question. Anyone who enters the \textit{Memory Ward} is told to read the aphorism upon the first step and ponder it the entire time they are climbing the stairway. If they find peace in their answer, they are to read the inscription of the second step on their next journey, and so on. There are few, even among the \textit{elves}, who have pondered the inscriptions of all of the steps to their satisfaction.

\subparagraph{5. Spires of Yurek}
Roads lined with dormitories and vendors radiate out from the branches of the Sequoia of Remembrance to the walls of the \textit{Memory Ward}, providing essentials for not only the Dreamweavers, but the elven scholars who teach, study, and experiment with magic within the Spires of Yurek. This academy was named for \textit{Yurek Windkeeper}, founder of the \textit{Arcana Pansophical} and personal counselor to the \textit{Voice of Memory} herself. The school is small compared to \textit{Emon's} magnificent \textit{Alabaster Lyceum}, but its seven marble towers are no less awe-inspiring. Characters who visit the Spires of Yurek can learn any historical fact by studying its tomes for one hour and then making a successful DC 20 Intelligence (History) check. They can even uncover hints toward long-lost secrets by completing a day's study and then succeeding on a DC 25 Intelligence (History) check. New spells are constantly being invented here, and countless tomes filled with ancient spells can be found within the libraries.

\subparagraph{6. Tower of Moonlight}
Most visitors to \textit{Syngorn} doubt that \textit{Ouestra, the Voice of Memory}, truly dwells in a tower of moonlight. It is the truth. Her tower does not exist under the light of the sun, and then appears as a semi-solid pillar of Catha's silvery glow, occasionally veined by lines of Ruidus's ruddy light when the smaller moon is full. Its inside has supposedly never been seen by any but the other \textit{Wardens of Syngorn}, and a thousand rumors circulate throughout the realm of the magical and divine marvels that must be contained within.

\subparagraph{7. Spellbenders' Guildhall}
The \textit{Verdant Expanse} is saturated with magic. The ley energies that suffuse the greenwood make it easy for arcanists to create works of spellcraft by themselves, when it might take a half-dozen mages working in concert in other lands. Many of \textit{Syngorn's} most trustworthy magi are members of the Spellbenders, a guild dedicated to creating new spells and using them to create works of art, defend their homeland, and improve their peoples' lives. Their guildhall is only accessible to those with the express permission of \textit{Guildrunner Rawndel}.

\subparagraph{8. Stone of the Arch Heart}
In the heart of the Reverie Walks is a column of pure diamond, carved in the radiant likeness of an androgynous elf. Though the The Arch Heart is worshiped throughout Tal'Dorei as a god of magic and art, they are something more to the \textit{elves} of \textit{Syngorn}: a progenitor of their people, a deity of fatherhood and motherhood, yet also of neither—a liminal god that by their very nature both embraces and destroys binaries. \textit{Elves} that are "born of the Arch Heart" are honored within \textit{Syngorn}, and many who seek that honor wander the Reverie Ward in search of the Stone, longing for physical transformation, social power, or arcane might. Some search for centuries and never find it, but in times of peril, the The Arch Heart always makes their wisdom known to those who truly need it.

\subsubsection{Syngorn Adventures}
Game Masters who set their adventure in \textit{Syngorn} can use this plot hook for inspiration.

\subparagraph{The River Does Not Run Through It (low-level)}
Word arrives from Verdant Scouts that the Feshun River has stopped emptying into the Ozmit Sea without warning. The powerful flow from the \textit{Tormor Falls} shows no signs of dwindling river flow, so the party is sent to investigate, only to find a sudden sinkhole that has consumed a small section of the \textit{Expanse}, causing the river to pour into this massive, underground cavern. The rock and earth below appears to be crumbling at an accelerated pace, and the air below smells of decay. What could be the source of this rapid destruction, and how can it be remedied?

\subsection{Tormor Falls}
Feeding the mouth of the rushing Feshun River that carves through the body of the \textit{Expanse}, Tormor Falls is an incredible multi-level series of waterfalls that cascade down the eastern side of Orencleft Mountain for hundreds of feet. Swelling with every major rainfall, and beautiful to behold at all times of the year, Tormor Falls is also host to a number of caves that hide beneath the mist and spray, leading beyond the forest and under the \textit{Stormcrest Mountains}.

A squad of Verdant Guard patrols the base of the falls, and many have explored the caverns in the past. Some return with nothing; others discover old relics and trinkets from the earliest days of Tal'Dorei, or even from the time of the \textit{Calamity}. It's rare for a patrol to fail to return, but when they do, no search party has ever been able to make contact with survivors, or to reclaim their remains. The guard has begun to whisper that the cave entrances move, shifting with each rising sun. The superstitions surrounding the caves around Tormor Falls have grown so great that expeditions which fail to return are no longer searched for at all, and no patrols go inside the caves without a critically important reason.

\subsubsection{Tormor Falls Adventures}
Game Masters who set their adventure in the \textit{Tormor Falls} can use this plot hook for inspiration.

\subparagraph{This Cave Was Made for Me! (mid-level)}
Characters that travel the edge of the \textit{Verdant Expanse} hear distressing news that people from small villages, and even from \textit{Syngorn}, have disappeared in the night—and now, a friend of the party has joined the ranks of the missing. What they don't know is that these people are leaving the village themselves, hypnotically drawn to the caves behind \textit{Tormor Falls} by an \textbf{aboleth} living within a pool deep inside the cave system. The beings that emerge from those caves are no longer humanoid, but nightmarish mockeries of life.

\subsection{Vues'dal Waters}

\begin{description}
\item[Scattered Villages:]
Population 2,420 (52\% \textit{elves}, 28\% \textit{goblinkin}, 12\% \textit{lizardfolk}, 5\% \textit{humans}, 3\% \textit{other races})

\end{description}

The volcano known as Mt. Vues'dal erupted for the last time in the early days of the \textit{Age of Arcanum}. Its explosion shook the earth so terribly that the mountain itself was swallowed, creating the Vues'dal Basin at the edge of the \textit{Stormcrest Mountains}. Today, the basin has been completely filled by the Feshun River, and the land around the Vues'dal Basin is among the most fertile farmland in southern Tal'Dorei. The villages of the Vues'dal Waters are known for their prize produce, exporting tomatoes, cherries, pumpkins, lentils, sweet potatoes, and dozens of other foods, both staples and luxuries, to \textit{Syngorn}—and from there, to the rest of Tal'Dorei.

The people of the Vues'dal Waters are mostly \textit{elves}, but over the years they've been joined by some of the local \textbf{lizardfolk} that have dwelt in the marshes, as well as goblinkin fleeing the \textit{Iron Authority} to the south. All of these various people have banded together to create something beautiful around the caldera of this long-dead volcano—and to repel bandits and monsters from the nearby marshes. The \textbf{lizardfolk} hate fighting their kin, but human and \textbf{lizardfolk} bandits alike have made the wetlands their home, and prey upon these honest farmers from encampments that shift from week to week.

\subsubsection{Vues'dal Adventures}
Game Masters who set their adventure in the Vues'dal Waters can use this plot hook for inspiration.

\subparagraph{Legions of the Reptile God (mid-level)}
The people of Vues'dal have unwittingly been skirmishing with the vanguard of a force that wishes to bring \textit{Syngorn} to its knees. The \textbf{spirit naga} Maledicta Hexos has come from below the \textit{Stormcrest Mountains} to rally the \textbf{lizardfolk} tribe in the name of the The Cloaked Serpent. Vues'dal will soon be overrun, and the Verdant Guard has called for \textit{Emon's} aid in driving back the legions of this monster with delusions of godhood.

Maledicta's armies contain not just countless \textbf{lizardfolk}, but also several \textbf{hydra}, and a \textbf{night hag} that serves as her lieutenant. Maledicta is also in possession of one of the \textit{Astural Scrolls}, an artifact looted from \textit{Wrettis} that grants her uncanny cosmic powers.

\section{Bladeshimmer Shoreline}
The Bladeshimmer Shoreline, named for the distant glimmering of sun across the Ozmit Sea, stretches across the central western coast of Tal'Dorei. This coast is where \textit{humans}, \textit{halflings}, and \textit{gnomes} first set foot on the continent of Tal'Dorei, and it bears the marks of their first steps on this land. It's also home to the present heart of the Republic of Tal'Dorei.

The shoreline is a hub of international trade, for the western coast of Tal'Dorei is closest to the continents of \textit{Issylra} and \textit{Marquet}. Dozens of tallmasted trading ships set sail to and from its calm shores each day—as well as the majestic skyships that daily soar in and out of the resplendent capital city of \textit{Emon}.

Inland Bladeshimmer is mostly temperate grassland, intercut by cool, winding rivers. The bulk of western Tal'Dorei's produce comes from farms here, blessed with clean water and non-salinated soil, despite the ocean's proximity. \textit{Emon's} presence lends the region stability. Nevertheless, the yeomen living beyond the city's walls struggle to defend their small plots from burrowing \textbf{ankheg}, corrupting fiends, and the various folkloric spirits and monsters that lurk in the shadows of their minds.

\begin{DndReadAloud}{}

\begin{itemize}
\item{Majority Faiths:}
The Lawbearer, The Wildmother, The Platinum Dragon

\item{Minority Faiths:}
The Moonweaver, The Matron of Ravens, The Dawnfather, The Lord of the Hells, The Scaled Tyrant

\item{Imports:}
Precious and industrial metals, lumber

\item{Exports:}
Stone, lumber, ships, \textbf{fish}, grain, produce, cobalt, \textit{gold}, livestock

\end{itemize}

\end{DndReadAloud}

\subsection{Daggerbay}
Daggerbay, so named for the jagged Slumber Reef that flanks the bay, was where the first human colonists of Tal'Dorei made landfall. In that bygone time, it was a bustling port for the human city of Port \textit{O'Noa}. It has since fallen into disuse after the \textit{Scattered War} and the destruction of \textit{O'Noa}. Centuries later, the bay is now naught but a haunting reminder of darker days. It is home to hundreds of sunken ships and the bodies of their lost crew, deep below the depths of the reef. Those who know of their history say that the waters are cursed, and the nearby settlers spread rumors of ghost ships and \textbf{siren} calling looters to their grave should they wander too close to the ominous waves.

\subsubsection{Daggerbay Adventures}
Game Masters who set their adventure in \textit{Daggerbay} can use this plot hook for inspiration.

\subparagraph{Stormbringers (mid-level)}
A society of banished \textbf{merfolk} have claimed \textit{Daggerbay} as their domain. They harass any treasure hunters that seek to loot the wreckage at the bottom of the bay. While the characters are traveling nearby, a terrible storm forces them to take shelter in the ruins of \textit{Daggerbay}.

In truth, the storm was caused by the renegade \textbf{merfolk}. Some weeks ago, they discovered a relic of the The Stormlord with the power to summon and direct storms. The coast is now shrouded in an endless maelstrom. That night, when the characters take shelter in the ruins, they're assailed by the \textbf{merfolk} land-dwelling humanoid thugs. These brutes attempt to capture the characters and deliver them to their leaders.

\subsection{Daggerbay Mountains}
The pages of history do not remember the Molten Titan, or when this fierce entity of liquid metal was felled by primordial \textit{elves} at the end of the \textit{Founding}. The entity was sealed within and consumed by the earth, giving birth to a range of mountains that has stood for the ages since. Jutting upward against the northwestern border of the \textit{Verdant Expanse} and stretching to the Ozmit Sea, these stormy peaks functioned as the first boundary between the long-established \textit{elves} of \textit{Syngorn} and the burgeoning human colonies on Tal'Dorei during their arrival.

For most of history, the \textit{Daggerbay Mountains}, or the lyren'alsi in Elvish, were known to be desolate and devoid of material worth. Prospectors from \textit{O'Noa} came home empty-handed, if they came back at all. The survivors returned with tales of bands of bloodthirsty, one-eyed giants that roamed the peaks: \textbf{cyclops}. \textit{Issylran} settlers and \textit{Syngornian} \textit{elves} both learned to leave the lyren'alsi well enough alone.

Decades later, the \textit{Emerald Outpost} was established as a major trading post for the people of \textit{Emon} and \textit{Syngorn}. With such a valuable settlement nearby, the \textit{Daggerbay's} hazardous peaks suddenly garnered the interest of adventurers and fortune-seekers. Small hunts for undiscovered riches grew into vast expeditions. These incursions into the mountains aggravated the mountains' territorial inhabitants; elven scouts reported \textbf{cyclops} raiders and \textbf{cyclops stormcaller} conducting strange rituals within the eyes of thunderstorms at night, and hungry \textbf{bulette} hunting beyond their normal domain, as if guided by some unknown intelligence.

\subsubsection{Daggerbay Mountains Adventures}
Game Masters who set their adventure in the \textit{Daggerbay Mountains} can use these plot hooks for inspiration.

\subparagraph{Eye of the Storm (mid-level)}
Groups of \textbf{hill giant} are migrating into the mountains, leaving their foothill settlements empty of warriors, civilians, and the infirm alike. A half-elf \textbf{scout} named Thunderchaser has caused a stir in the \textit{Emerald Outpost} by bringing back the lightning-charred corpse of a \textbf{hill giant} with one eye ritually gouged out. What are the \textbf{cyclops stormcaller} planning beneath the concealing stormclouds of the \textit{Daggerbay Mountains}?

\subparagraph{A Metal of Memory (high-level)}
The ancient molten beast that was sundered beneath these rocky peaks left behind pockets, located far below the surface, of an incredibly rare metal with mysterious properties, called orichalcum. The sudden discovery has sent local entrepreneurs into a frenzy, and the ensuing rush has already sparked violence among the bands of treasure hunters. As the adventuring parties dig deeper, they discover that some essence of the ancient creature still lives, and it is seeking vengeance.

\subsection{Emerald Outpost}

\begin{description}
\item[Town:]
Population 1,470 (76\% \textit{elves}, 19\% \textit{humans}, 3\% \textit{halflings}, 2\% \textit{other races})

\end{description}

Originally established by the \textit{elves} of \textit{Syngorn} as a hidden outpost to protect their homeland from the growing empire of \textit{Drassig}, the \textit{Emerald Outpost} fell out of use after the \textit{Scattered War}. This former military outpost is now an intermediary trading post and a place for merchants traveling between \textit{Emon} and \textit{Syngorn} to rest. The Emerald Chambers, once used as a war room, and more recently used as a diplomatic chamber for the allied governments of \textit{Emon} and \textit{Syngorn}, have these days been converted into a stylish bazaar, in which \textit{Emonian} and \textit{Syngornian} traders can sell their wares without making the full journey to either city.

\subsection{Emon, the City of Fellowship}

\begin{description}
\item[Metropolis:]
Population 365,026 (63\% \textit{humans}, 8\% \textit{dwarves}, 7\% \textit{elves}, 22\% \textit{other races})

\end{description}

\textit{Emon} stands defiantly against all who would threaten Tal'Dorei and its people. It is the cultural and political heart of the Republic of Tal'Dorei, as well as a nexus of commerce, entertainment, travel, education, and adventure.

\textit{Emon} is accessible only through its heavily patrolled gates and by \textit{skyship}. The denizens of the city are generally well protected from outside attackers and sieges.

\subsubsection{Government}
\textit{Emon} is the seat of the \textit{Tal'Dorei Council}, the nation's highest governing body. Though the council originally served under the Sovereign of Tal'Dorei, the nation's final sovereign was killed when the \textit{Chroma Conclave} attacked \textit{Emon}. Following the death of \textit{Thordak the Cinder King} and the rest of the \textit{Conclave}, \textit{Emon} was rebuilt and the council reformed as the backbone of the new Republic of Tal'Dorei. The council and its members are \textit{described in detail starting on page 42.}

One unintended consequence of relying on magic to rapidly rebuild \textit{Emon} is that the mages of Tal'Dorei now hold incredible influence over the fledgling council. Some fear that without a sovereign, the council will be unable to keep the arcanists in line, and Tal'Dorei will dissolve into magocracy.

\subsubsection{Current Events}
The \textit{Tal'Dorei Council} is best known for their policies that affect the whole of Tal'Dorei. It is often ignored that the council also governs \textit{Emon} and resolves the day-to-day issues of the city. These are some current events both great and small affecting \textit{Emon} today; these topics can add flavor to the discussions and rumors the characters hear around town, and could even inspire new adventures for them to take on.

\subparagraph{Elemental Disturbances}
Over twenty years have passed since \textit{Thordak the Cinder King} was killed in the caverns beneath \textit{Emon}. His wreckage in the \textit{Cloudtop District}, as well as the \textit{Scar of the Cinder King}, a swath of elemental destruction outside the city walls, have been tended to by members of the Fire \textit{Ashari}, gradually healing the stubborn wounds of his influence. Yet the primordial fires that seeped from his body continue to linger in the caverns beneath the city where he died. Every year or two, an incursion of wild fire and hateful \textbf{cinderslag elemental} burst through the city streets, causing mass chaos, expensive damages, and life-or-death battles for the city guard.

\subparagraph{Mercenary Mages}
A generation ago, the fledgling \textit{League of Miracles} hired hundreds, if not thousands, of mercenary mages to aid in the reconstruction of \textit{Emon} and other cities across Tal'Dorei. These days, Tal'Dorei is rebuilt, and the few infrastructure problems aren't enough to employ all the mercenaries left behind by the league. These out-of-work mercenaries now linger around dismal dive bars and make trouble for the locals. Most people would rather see them gone, but it's hard not to feel sorry for their situation.

\subparagraph{Walls and Slums}
The city's outer walls have fallen under heavy criticism from politicians and civilians alike in recent years. Many people are concerned that the walls' true purpose these days is to separate the slums outside the walls from the "respectable" districts within them. Popular opinion is now that the great \textit{City of Fellowship} should do more to help its less fortunate, but council politicians continue to stall. In private, Master of Commerce \textbf{Vex'ahlia} de Rolo fumes that she wants to help, but can't cut through political red tape with her legendary bow and arrow.

\subsubsection{Society}
While in antiquity \textit{Emon} was the heart of a human empire, the city has grown to encompass people from all over Tal'Dorei, and all over the world. Its citizens include people of nearly all known races from Exandria's many nations, thriving on the new innovations and ideas of its citizenry. It has been this way since the sovereignty of \textit{Zan Tal'Dorei}, when the alliances she formed invited dozens of displaced \textit{elves}, \textit{dwarves}, \textit{orcs}, and \textit{goliaths} into her city. Since then, \textit{Emon} has developed into a city whose people valued innovation and collaboration, especially in small, cohesive groups—virtues that some historians within the \textit{Alabaster Lyceum} believe have given rise to the prominence of the modern adventuring party.

While \textit{Emon} is becoming more commercial and the use of \textit{gold} as currency has been ubiquitous throughout Tal'Dorei's history, many communities within the city are still close-knit enough to use the barter system. During the reign of \textit{Drassig}, the \textit{humans} of \textit{Emon} adopted a dwarven oath called rudraz, an intimate promise between two people to repay a deed or trade. Though the rudraz is not a contract, the \textit{dwarves} believed that an oathbreaker would be forever barred from passing beyond the Brightguard Gates of Hilmaire, the gates that allow \textit{dwarves} to pass to the afterlife. \textit{Humans} in \textit{Emon} treat the rudraz more lightly, often using it to seal matters of business or politics rather than personal promises, but breaking this oath still carries massive social repercussions. Few look kindly upon a person with the epithet "Oathbreaker".

\subsubsection{Rebirth}
\textit{Emon} has mostly recovered from its destruction at the talons of the \textit{Chroma Conclave} and \textit{Thordak's} subsequent occupation. A number of Tal'Dorei's most significant factions played a role in \textit{Emon's} rebirth, and their members enjoy positions of prominence throughout the city. \textit{Emon} was rapidly rebuilt to its former glory, due to the expensive but effective intervention of mercenary mages serving the \textit{League of Miracles}. Thanks also to the \textit{Clasp's} underground networks and established hierarchies, \textit{Emonian} society and culture was able to bounce back more swiftly than any other major settlement in Tal'Dorei.

And the \textit{Tal'Dorei Council}, formerly an advisory council to the sovereign, is now the ruling council of the land. Folk heroes and career politicians alike sit on the council. Over the past two decades, the council has, with many missteps, enacted the will of the people—allowing Tal'Dorei to transition from a monarchical power to a republic without war or insurrection.

However, the swiftness of \textit{Emon's} rebirth has put the \textit{Tal'Dorei Council} in a precarious position. Not only are they socially indebted to a criminal faction—one whose members now command immense celebrity among the common folk—the council now owes vast amounts of gold to mages of the \textit{League of Miracles}. League conjurers created thousands of tons of stone and steel, and each one of their transmuters rebuilt everything from taverns to townhouses and sewers to stately manors as quickly as a hundred laborers. \textit{Thordak's} hoard of plundered wealth has dwindled to nothing, repaying the league, and they have done the same in cities across the land. Some members of the council have already caved to pressure from factions within the \textit{Clasp} and the \textit{League of Miracles}, and subsequently turn a blind eye to the crime and magical abuse that run rampant throughout the city. Barely a single human generation old, the young republic is already struggling under the weight of ruling in a just and noble manner.

\subsubsection{Defenses}
The Republic of Tal'Dorei has no standing army, and doesn't elect a Master of War—a military commander-in-chief—to its council outside of wartime. In peacetime, the Master of Defense is responsible for the acquisition of mercenaries and adventurers for tasks which require the use of force. They also have the ability to conscript existing militias and town guards into military service and call them to \textit{Emon's} defense at any time, though typically only while the Master of War is being elected.

The only standing forces pledged exclusively to the defense of the realm are the elite forces of \textit{Fort Daxio}. Mobilizing the armies of \textit{Fort Daxio} to \textit{Emon's} aid generally requires the city to withstand a week-long siege. However, in case of internal strife, there are stables positioned every fifteen miles along the Othen Trail, allowing Daxio's legendary outriders to arrive at \textit{Emon} in a mere two days' time by switching out their exhausted steeds.

\textit{Emon} is encircled by 60-foot-high stone walls that stretch from the eastern fields to the western shore—though a number of discussions in the \textit{Tal'Dorei Council} have revolved around the strategic irrelevance of great stone walls. They didn't protect \textit{Emon} against the \textit{Chroma Conclave} over twenty years ago, and they wouldn't have done much against an army of \textbf{cultist} and winged \textbf{gloomstalker} had The Whispered One led an army against Tal'Dorei. As skyships grow more popular and affordable around the world, even the council's Master of Defense is considering demolishing the city's walls—making the strange case that such fortifications are prized more for their historic and cultural value than for their strategic importance.

These walls are patrolled by lightly armed and armored members of the Arms of Emon, the city's official watch and constabulary. They are divided unofficially into the Arms upon the Wall and the Arms on the Streets, the latter of which are assigned to patrol the city streets to dissuade petty criminals and respond to violent crime.

\subsubsection{Crime}
Hundreds of miles of tunnels of varied provenance run beneath \textit{Emon}. Most are the forgotten remnants of paved-over neighborhoods and secret passageways made by thieves' guilds during the reign of \textit{Drassig}. Others are the lairs of tunneling monsters. The origins of others still are lost forever, for they have already long since been turned into sewer passages to whisk away the refuse of \textit{Emon's} citizens.

The vast majority of the tunnels closer to the surface are known to the \textit{Clasp} and used to rapidly move about the city unseen. Somewhere deep in this network of unlabeled roads is the \textit{Clasp's} secret headquarters, where all manner of thieves, \textbf{clasp cutthroat}, fences, and spies gather beneath their banner. Though the \textit{Clasp} operates in major cities across Tal'Dorei, \textit{Emon} has always been their home, and saving their neighbors from \textit{Thordak's} reign of terror gave the \textit{Clasp} the dubious benefit of a heroic reputation.

However, the \textit{Myriad} crime syndicate from \textit{Wildemount} has spent over two decades relentlessly clawing their way into the \textit{Clasp's} territory across Tal'Dorei—dismantling their networks, poaching their clients, and assassinating their informants. The \textit{Myriad} hasn't managed to completely eliminate the \textit{Clasp's} hold in any of the cities they've targeted, but that hasn't stopped them from bringing the ongoing power struggle to the \textit{Clasp's} doorstep. \textit{Myriad} agents bargain in shadowed alleys, gamble in upscale taverns, and lurk in the very tunnels that the \textit{Clasp} relies on, waiting for their chance to unseat the \textit{Clasp} once and for all.

\subsubsection{Geography and Climate}
\textit{Emon} is not a tourist destination for its climate—the city is known for sudden, pounding rainstorms just as much as for its beautiful, sunny days. The city is blessed with cool summers and warm winters, thanks to its proximity to the cool Ozmit Sea. Though snow rarely falls on the city itself, it relies on springtime snowmelt from the \textit{Cliffkeep Mountains} to fill its reservoirs throughout the year.

\textit{Emon} is made up of the following districts:

\subparagraph{Abdar's Promenade}
The Promenade is an open marketplace district and massive bazaar that dominates eastern \textit{Emon}, named for the legendary spicemonger from \textit{Marquet} who helped fund the construction of \textit{Emon}. To this day, the name of Abdar is synonymous with both generosity and business savvy. Nearly everything a person could want can be found within the tents, carts, warehouses, and shops of the promenade. Its intertwining roads stretch for miles, and it has dozens of tiny neighborhoods, each with their own distinct cultures and identities. The promenade is heavily patrolled by the Arms of Emon, for its sunny reputation is shrouded by an equally significant reputation for literal back-alley criminal deals.

\subparagraph{Central District}
This is the largest residential district within \textit{Emon's} walls. It houses most of the city's merchant class, including traveling traders, ship captains, and guild journeymen. The Central District is a patchwork of thousands of personal homes, tenements, and guildhalls of all shapes and sizes, peppered with small taverns and inns on nearly every street corner.

These neighborhoods lay clustered together among tightly set streets, occasionally broken up by park grasses or the Ozmit Waterways that snake through the region between the promenade and the port. Visitors to the Central District are advised not to go out at night; the wide disparity of wealth from home to home here has seen a recent rise in crime, and the city watch seems reluctant to find a constructive solution. Residents who know the lay of the land have an easier time at night and can help visitors avoid the most dangerous streets.

\subparagraph{Cloudtop District}
Out of all of \textit{Emon's} districts, the Cloudtop has, in some ways, changed the most over the decades. In \textit{Drassig's} time, a monstrous fortress loomed over the manors of his sycophants. In the age of \textit{Zan Tal'Dorei's} descendants, countless mages transfigured that decadent citadel into a sleek, elegant palace. By the time the \textit{Chroma Conclave} came, decadence had once more crept up on \textit{Emon}, and though the Palace of the Sovereign remained much the same, the Cloudtop District was once again a nest of luxury and the throne of the social elite. The latest additions in that era were \textit{skyship} platforms, so that the wealthiest and most connected in Tal'Dorei could travel in exclusive luxury—soaring through the skies rather than enduring arduous journeys over land or sea.

All of that was reduced to slag by the advent of the \textit{Cinder King}. \textit{Thordak} made the Cloudtop District his home, and the palace itself his throne. His very presence caused magma to surge from the caverns beneath \textit{Emon}, where previously no fires burned. \textit{Thordak} ruled there for only a matter of months, but those few months were enough to annihilate unknown millions of gold pieces' worth of art and creations of untold value—though the greedy dragon took care to command his minions to save the most beautiful works for his treasure hoard.

Today, the Cloudtop District is rebuilt, though the most stubborn elemental influences of the \textit{Cinder King} linger in \textit{Thordak's Crater}. A new castle called \textit{Cloudwatch Palace} has been built atop the city of \textit{Emon}, and it contains the chambers of the \textit{Tal'Dorei Council}, as well as grand dining halls and guest quarters for visiting dignitaries. Much of the castle is a museum, displaying what regalia and objets d'art survived the wrath of the \textit{Cinder King}. The private chambers of the palace also contain a number of \textit{teleportation circle} that leading councilmembers can use to return to their homes at night—particularly if they live in far-off cities such as \textit{Whitestone}.

\subparagraph{Temple District}
Some say that this northern district contains more temples than in the rest of Tal'Dorei's cities combined. This is a flagrant exaggeration, but at first glance, most people believe it. \textit{Emon} is a city of many people, from many different places, with many different faiths. Even folk who worship the same gods worship them differently, and have built their own temples. Legend has it that shrines and temples used to be built all across \textit{Emon} (and there are still some small sanctuaries elsewhere in the city), but a cleric of the The Dawnfather declared her temple the greatest in all \textit{Emon}. Her boast was met with scorn from a nearby church of the The Stormlord. Their clerics, always eager to rise to a challenge, built an even greater temple across the street, and soon, people of all faiths were building the most magnificent structures they could, all in one tiny corner of \textit{Emon}.

Today, cathedrals to the The Lawbearer and the The Platinum Dragon stand at apex among smaller temples to the The Dawnfather, the The Stormlord, the The Wildmother, and the The Matron of Ravens. All \textit{Prime Deities} have at least a moderately sized house of worship here, and even some of the \textit{Betrayer Gods} are worshiped in secrecy.

Anyone who seeks shelter here, be they rich or poor, can find sanctuary in the Temple District. Even the faithful of gods of strength, such as the The Stormlord and the The Platinum Dragon, have mercy for those who need the aid of others to help them stand on their own.

Travelers and sailors come to leave tokens and gifts at shrines to bless their journeys, while people of all walks of life seek the graces of the The Changebringer before undertaking risky endeavors. The meek search for wisdom in the great halls of the The Knowing Mentor, as the artists and wistful folk give praise to the The Arch Heart. However, not all folk are drawn toward religion, nor trusting of it, and temples past have been exposed as frauds—or worse, cults who venerate power.

Recently, an off-kilter faith to a lesser idol called the The Traveler has found purchase in the Temple Ward—primarily in the form of puerile graffiti in back alleys. Other quasi-religious spiritual entities have small cult followings in the Temple Ward. They are generally allowed to flourish to whatever degree they can, so long as their beliefs don't advocate for any behavior that would break the laws of \textit{Emon}.

\subparagraph{Erudite Quarter}
This district is home to members of the educated upper-middle class and destitute academics alike. Several institutions of higher learning have their stately campuses here, ranging from well-to-do boarding schools to Tal'Dorei's most premiere academic and arcane institutions.

Beautiful towers in ancient \textit{Syngornian} style stretch to the sky across a cityscape of centuries-old halls of learning and student boardinghouses. Here, the finest schools and colleges draw the wealthy, the gifted, and the brilliant—and none more so than the \textit{Alabaster Lyceum}, often lauded as the largest and most accomplished institute of arcane study on the continent. Though the district's walls were shattered and its libraries crushed during the invasion of the \textit{Chroma Conclave}, the Erudite Quarter has thankfully been mostly reconstructed, thanks to the efforts of the \textit{Lyceum's} staff and students.

Members of the \textit{Lyceum}, under the direction of Headmaster Thurmond Adlam, performed the reconstruction themselves not simply because they sought to project an image of good work ethic and self-sufficiency, but because the heads of their colleges distrust the \textit{League of Miracles}. Their motives and sources of funding are, in the eyes of the \textit{Lyceum}, too opaque to be trusted. Their suspicion has often been derided either as elitism over the diverse array of mercenary mages the league hires, or as sour grapes over losing out on countless lucrative reconstruction contracts.

\subparagraph{Military District}
Since \textit{Emon} has no standing army, the Military District is one of the city's smallest wards. Nevertheless, it has the capacity to house and train a vast number of soldiers if an army needs to be levied from the people of Tal'Dorei. The Military District is mostly filled with several prisons and jails, including the infamous \textit{Black Bastille}, as well as tightly locked barracks and instructional facilities for the Arms of Emon.

\subparagraph{Cemetery District}
This smaller section of the city is filled with graveyards and small religious sanctuaries so that people with means to purchase a gravesite or a mausoleum can bury their dead in whatever luxury they can afford. Countless gravestones line the rolling hills of the district. They are sectioned off from the roads by tall iron fences, with mausoleums of more affluent families dotting the grassy hillsides. As surface space grew limited in centuries past, the city excavated a network of catacombs called the Undervaults. Upkeep, expansion, and general safeguarding of the sites are all overseen by the Gravewatchers, a guild of gravekeepers that have held political and social hegemony over the district for generations.

Their attempts to expand the Undervaults have had mixed results over the centuries. Most new excavations meet without incident (likely because the \textit{Clasp} has learned to keep their tunnels far from this district). However, excavations in recent decades have broken into the \textit{Crystalfen Caverns}, an eerie cave system filled with mind-bending inscriptions and horrific beings. The breach was promptly sealed off.

Yet the sealing of the caverns did not stop some members of the Gravewatchers from seeing an opportunity to gain fame and glory. Most who ventured into the caverns were never seen again, and those who have returned shudder unspeakably of the horrors they escaped. Still, passage is occasionally granted to the \textit{Clasp} or adventurers willing to pay the Gravewatchers for entry to the dangerous realm below.

\subparagraph{Port of Emon}
\textit{Emon} is bordered on the west by the Ozmit Sea and is generally sheltered from the worst of its storms. The vast \textit{Port of Emon} allows hundreds of thousands of gold pieces' worth of goods to flow through the city every single day. Well over a hundred ships fill this port at any given time, and the dock crews are ever working, carrying crates and goods away to \textit{Abdar's Promenade}, or to empty vessels for exportation. The northern sector of the Port is mostly allocated to the Everline, a powerful fisherfolk's guild. They frequently feud with the Onyx Banner, \textit{Emon's} most powerful shipping guild, over which docks they are permitted to use. Both guilds keep their hands clean, but it's an open secret that both hire adventurers and mercenaries to trash the other's ships and warehouses to intimidate them. Adventurers in need of fast cash who don't mind dodging the Arms of \textit{Emon} can find good work as muscle for either of these two guilds.

\subparagraph{Outwall}
During the reign of \textit{Drassig}, a massive city of tents and hovels grew outside of \textit{Emon}, filled with those too poor to pay \textit{Drassig's} extortative taxes. By the end of the \textit{Scattered War}, the slum had grown to a third of the city's size, and nearly half its population. Even during the reign of \textit{Zan Tal'Dorei}, these slums continued to grow, as refugees of the \textit{Scattered War} tried to earn enough money to find stability within \textit{Emon's} walls.

Though their conditions were still meager when the magnanimous \textit{Zan Tal'Dorei} ended \textit{Drassig's} iron rule, the city's denizens chose to remain, having created a community with its own culture of inexpensive living and brotherhood in poverty, away from the bustle of the city's inner streets. A microcosm of \textit{Emonian} society now exists within the city of Outwall, including its own trade square, shrines for worship, and makeshift farms on the outskirts. As early as two decades ago, Outwall was called the Upper Slums—a derogatory name given to it by \textit{Emon's} Cloudtop elite. People within Outwall tried to transform the term "Upper Slums" into a symbol of civic pride, but it never stuck. The name Outwall eventually emerged after community leaders demanded an audience with the \textit{Tal'Dorei Council} in an event that came to be called the Outer Wall Debates—for they demanded the council convene in the slums' own town square.

These days, Outwall is becoming a trendy place to live, thanks to its vibrant culture and relatively inexpensive housing. Its locals are fighting hard to keep young people with money to burn from the \textit{Central District} from gentrifying their home. It hasn't come to violence yet, but the angrier members of the Outwall community clandestinely meet with adventurers to burn down the stylish houses and mansions that are driving out local residents and businesses.

\subparagraph{Southgate Farms}
Inspired by the civic action of the people of \textit{Outwall}, the farmers of Southgate were able to break from the ugly name of "slum" in the past decade. This community of farmers and ranchers, formerly known as the Lower Slum, exists outside \textit{Emon's} southern gates. It first formed when a large refugee band from \textit{Othanzia} was barred entry into the city in the midst of the \textit{Scattered War}. Instead of dispersing, they decided to lay down roots while \textit{Drassig's} armies were out fighting a losing battle against \textit{Zan Tal'Dorei}. These squatters eventually formed a community of farmers who petitioned \textit{Zan Tal'Dorei} for land in exchange for providing produce and grown goods to the people of the city.

Thanks to the aid of the First Sovereign, the descendants of those \textit{Othanzian} refugees have become one of the most significant sources of fresh food for the people of \textit{Emon}. Frustrated by being saddled with the name of "slum," and inspired by the Outer Wall Debates, the leading families of the Lower Slums demanded recognition by the \textit{Tal'Dorei Council} as well. They were met with little debate, and these humble farmers are now a respectable part of \textit{Emonian} society—even if they live outside the city's defensive walls.

This southern region of the city is also home to a famous site of modern history: Greyskull Keep, the former base of operations for Vox Machina in \textit{Emon}. This fortress was gifted to them by the Last Sovereign as thanks for saving his life. The various heroes of Vox Machina have moved on to other homes (and now typically gather in \textit{Whitestone} for holidays), and they allowed the keep to be turned into a museum of their heroic exploits.

\subparagraph{Shoreline Farms}
Less respectable than the \textit{Southgate Farms} are the communities along a stretch of murky shore north of \textit{Emon}. These people of \textit{Outwall} have tried for decades to turn terrible, salty land into effective farmland, to mixed results. They have successfully cultivated the few vegetables that grow in saltwater-soaked soil, leading to bountiful harvests of savory quattet gourds and the hard-to-grow brineroot. However, the farms' distance from the city gives them little protection from dangers that city folk think little of. Adventurers are always wanted to protect the farms from petty criminals and wandering monsters, but the farmers have little to pay. Generally, only utterly novice adventurers find their start in a place like the Shoreline Farms—but all heroes have to start somewhere.

\subparagraph{The Grotto}
This sprawling series of underground chambers composes the \textit{Clasp's} headquarters beneath the city of \textit{Emon}. It's rumored that over a dozen entrances to the Grotto are concealed amidst the labyrinthine tunnels, which are both naturally occurring and a part of \textit{Emon's} sewers. That doesn't even begin to get into the hundreds of other entrances and access shafts abandoned and filled in for fear of discovery. The very nature of the \textit{Clasp's} business is to keep shifting out of sight, and as such, the actual core of the organization moves from base to base between multiple subterranean structures, ever building further underground, or repurposing a long-abandoned hideout whenever necessary.

\subsubsection{Points of Interest}
A number of both popular and historically significant landmarks can be found in \textit{Emon}. They are keyed to the map of \textit{Emon}.

\subparagraph{1. Laughing Lamia Inn}
While the city is home to countless taverns and inns for the weary traveler or practiced performer, few are as grand and well known as the Laughing Lamia. Almost as old as the city itself, this fine establishment has traded hands numerous times over the years, and has expanded with each transition, leaving the massive, four-story institution a hodgepodge mass of odd rooms and themes, all centered around a raucous central tavern room filled with dozens of tables. Currently under the ownership of a brash and cheerful lass called Luthania Wells, this hot spot for traveling traders and adventurers alike rarely sees an hour in the day (or night) where there isn't something interesting going on within.

\subparagraph{2. Azalea Street Park}
Azalea Street, within \textit{Abdar's Promenade}, is one of \textit{Emon's} oldest neighborhoods, and is full to bursting with small businesses, quaint homes, and hole-in-the-wall restaurants, such as the \textit{Laughing Lamia}. Amidst the chaos of daily life, Azalea Street is also home to a peaceful park overlooking the Ozmit Waterways. When the stress of constant battles becomes too much to bear, adventurers in the know take a moment to visit the Azalea Street Park and recover, maybe even swapping stories with other relaxing heroes. Legend has it that heroes who meditate among the flowers and drink in the sea air return invigorated by the spirit of the The Wildmother. A statue of the The Wildmother and The Lawbearer dancing together graces a fountain in the center of the park. Once per week, a character can gain inspiration (inspiration) by taking a resting in Azalea Street Park.

\subparagraph{3. Anvilgate}
Made of stark black stone and adorned with intricate brass metalwork, this grand blacksmithy is one of the most striking buildings in \textit{Abdar's Promenade}. Decades ago, a restless craftsman from \textit{Kraghammer} of the name Grahf Tiltkettle grew weary of solitary pursuits and worried that interest in blacksmithing would give way to the emerging trend of mass-production. Tiltkettle retired in \textit{Emon} and dedicated his respectable fortune to ignite the spark of inspiration in those who wished to learn the ways of the hammer and anvil, hoping to instill the joy of the The All-Hammer in a new generation. He funded the creation of Anvilgate, a massive, public blacksmithy where any folk who wished to try their hand at the smithing of metals could do so under the The All-Hammer guidance.

Cheaper ores and materials were often donated by friends of Tiltkettle to aid the volunteer mentors in teaching their craft to the less fortunate, while folks who wished to bring their own ingots were welcome to do so. The location has become popular enough to begin expanding its facility to include kilns for ceramics and looms for weaving. While the tools are open for anyone to use on the premises, all are warned of the curse that will beset any wanton thief who steals from under the The All-Hammer gaze.

\subparagraph{4. Traverse Junction}
The most magically active structure in Tal'Dorei shines like a pyramid of pure sapphire in the bustling center of the \textit{Erudite Quarter}, its walls thrumming with arcane energy. It is the Traverse Junction, and all visitors to the \textit{Erudite Quarter} stop to see it when they visit.

In nearly every major Exandrian city, there are \textit{teleportation circle} that link mages to other such \textit{teleportation circle} around the world. Each and every one of these major \textit{teleportation circle} has a twin in the Traverse Junction, a travel nexus for approved mages and world leaders. Any characters renowned in \textit{Emon} or honored by the \textit{Alabaster Lyceum} may make use of the Junction (for 50 gp, if they can't cast \textit{teleportation circle}), allowing them to travel to major cities such as \textit{Westruun}, \textit{Syngorn}, \textit{Kraghammer}, and \textit{Whitestone} in Tal'Dorei; Port Damali and Rexxentrum in \textit{Wildemount}; \textit{Ank'Harel} in \textit{Marquet}; and Vasselheim in \textit{Issylra}.

The telemagi who curate the \textit{teleportation circle} are always in search of new teleportation sigils to different lands; they are willing to pay a handsome sum for the services of powerful and magically skilled adventurers willing to seek out and catalog new circle sigils for their records.

\subparagraph{5. Alabaster Lyceum}
The gleaming white halls and elaborate gardens of the \textit{Alabaster Lyceum} signal to all who see them that they gaze upon the greatest institute of higher learning in Tal'Dorei—though numerous scholars in \textit{Kraghammer} and \textit{Syngorn} would heatedly dispute that claim. Regardless, this place is the heart of homegrown academia in the Republic of Tal'Dorei, boldly striving to make a name for itself alongside the ancient scholarly traditions of \textit{Issylra}, \textit{Marquet}, and \textit{Wildemount}. The \textit{Lyceum} (for those who can afford the impressive tuition) is an unparalleled center for study and research of history, economics, alchemy, art, and—most of all—magic and the arcane. Many of Tal'Dorei's greatest scholars and bards paid their dues in these halls, and the history of the \textit{Lyceum} and its graduates are featured in all manner of superstitions and folktales.

\subparagraph{6. Cloudwatch Palace}
In the wake of the \textit{Chroma Conclave's} destruction of \textit{Emon}, the towering palace that stood as home to the sovereign and the center of governance within the city fell to fire and ruin. In the decades since, construction began with the aid of the \textit{League of Miracles} to construct a new palace to help house the \textit{Tal'Dorei Council} and be a symbolic center for all major political affairs within the region. The Cloudwatch Palace now sits atop the once-smoldering hill where \textit{Thordak} once claimed his roost, a symbol of \textit{Emon's} enduring will and optimistic hope for the future. Within these high-arched halls, countless chambers hold space for diplomats, meetings, and in the case of holidays, grandiose celebrations that draw the attention of powerful figures all across Exandria.

\subparagraph{7. Thordak's Crater}
One last unsettling element of the \textit{Cloudtop District} is the elementally scarred remains of \textit{Thordak's Crater}. Even after two decades of healing from \textit{the Ashari} of Pyrah, \textit{Thordak's} fires still burn beneath the Cloudtop. The smoldering entrance to the caverns remains cordoned off from the public because of its unsightly appearance and sulfurous stench, and because fiery creatures occasionally try to claw their way out into the fresh air. The crater is often forgotten, but every few months an adventuring party makes their way inside to seek treasure or glory within the volatile caverns—and such expeditions are a favorite topic of conversation (and clandestine betting) among the politicians, political aides, remaining city nobility, and nouveau riche of \textit{Emon}.

Why \textit{Thordak's} magic lingers enough to create these mindless \textbf{cinderslag elemental} is a mystery; but it seems to be limited to the crater, as the elementals cannot go beyond it without collapsing into inert slurry.

\subparagraph{8. Emon Skyport}
This elevated port is complete with docking platforms solely for the use of skyfaring vessels. The original \textit{Emon Skyport} was the first of its kind in Tal'Dorei, and its successor is every bit as revolutionary as the original.

The \textit{Emon Skyport} was destroyed when the \textit{Cinder King} claimed the \textit{Cloudtop District} as his fiery throne, but like so much of \textit{Emon}, it has since been rebuilt. Unlike other parts of \textit{Emon}, this reconstruction was a joint venture between the \textit{Arcana Pansophical} and the Alsfarin Union, a \textit{Marquesian} guild responsible for the engineering and distribution of all skyships in Exandria.

Booking passage on a \textit{skyship} isn't cheap (see the "Skyships" sidebar), but the privilege is available to everyone that can pay. Generally, one passenger \textit{skyship} leaves \textit{Emon} per week to each of the following locations:


\begin{itemize}
\item Vasselheim, \textit{Issylra} (passenger ticket costs 2,500 gp)
\item \textit{Whitestone}, Tal'Dorei (passenger ticket costs 900 gp)
\item Port Damali, \textit{Wildemount} (layover in \textit{Whitestone}; passenger ticket costs 2,000 gp)
\item Ank'Harel, \textit{Marquet} (passenger ticket costs 2,500 gp)
\end{itemize}

Members of the \textit{Tal'Dorei Council} can also usually pull strings to grant adventurers they hire free passage when on vital business.

\subparagraph{9. Tomb of the Last Sovereign}
Uriel Tal'Dorei II, the Last Sovereign, was slain by the noxious breath of the green dragon Raishan during the attack of the \textit{Chroma Conclave}. Uriel's wife, Salda, survived him and petitioned the reformed \textit{Council of Tal'Dorei} to build their last sovereign a tomb befitting his benevolence and magnanimity. They complied, and the monumental Tomb of the Last Sovereign is now home to not just Uriel's ashes, but to tribute from all who loved him. Those with wealth gave gold, while those whose hearts outweighed their purses gave more personal tribute.

Over twenty years ago, the Gravewatchers closed the tomb to all, supposedly to prevent vandalism, but have not reopened it after more than two decades. The Last Sovereign's eldest daughter, \textit{Odessa Tal'Dorei}, is furious. Despite her position of power on the \textit{Tal'Dorei Council}, she is unable to dispute the act because of the power the Gravewatchers have as keepers of the dead. It has come to the point where she seeks to secretly employ adventurers to uncover the Gravewatchers' true intentions—which involve the psychic influence of an aberrant creature of the \textit{Crystalfen Caverns}.

\subparagraph{10. Godsbrawl Ring}
Though the Temple of the The Stormlord always has a fighting ring in its center, the annual Godsbrawl transforms the earthen, torchlit sanctum into one of Tal'Dorei's most unusual tournament grounds. On the Day of Challenging, the god of athleticism's holy day, the brawny priests of the The Stormlord invite warriors and worshipers of the entire pantheon to the Godsbrawl, asking that each temple offer forth their greatest warrior to act as their god's proxy in the tournament. The clergy of the The Dawnfather and the The Lawbearer send their champions, but rarely take it seriously, for they scorn the storm-priests' notion of "might makes right." Conversely, the champions of the The Stormlord and the The Platinum Dragon have a fierce—though friendly—rivalry, trading the title of Supreme Champion back and forth each year after a bloody (but rarely fatal) final round.

\subparagraph{11. The Black Bastille}
Named for its ash-blackened walls, the ominous Black Bastille is a single-story prison within the \textit{Military District} that holds Tal'Dorei's most dangerous and irredeemable criminals. The compound has no windows, no open courtyards, and only one entrance: two imposing steel doors, flanked by watchtowers. During the reign of the \textit{Chroma Conclave}, the prison was one of the first buildings attacked, setting hundreds of Tal'Dorei's worst criminals loose to sow chaos across \textit{Emon}. The Black Bastille has since been rebuilt, but dozens of its most deplorable inmates still run free, and the Arms of \textit{Emon} are eager to recover them. Some include the human demon-summoner Felrinn Derevar, betrayer of the \textit{Arcana Pansophical}; the half-elf pyromancer Illaman Falconsong, exile of the Fire Ashari; and the vampiric tiefling \textit{Ixrattu Khar}, cultist of The Whispered One and Tal'Dorei's foulest mass murderer.

\subparagraph{12. Walls of Tribute}
Anyone who wishes to uphold justice and keep the peace can apply to join the Arms of \textit{Emon}. But before a recruit is given their armor and takes a vow to protect the city and its people, they must endure a grueling training regimen within the tall, featureless stone barricades of the Walls of Tribute, referred to wryly by recruits as the Quarry. There, they are trained in the physical aspects of peacekeeping, instructed on the habits of the city's criminal underbelly, and drilled in the laws and ideals of \textit{Emon}.

\subparagraph{13. House of Discipline}
\textit{Emon} does not have a standing army of its own. The skills imparted to the Arms of Emon are sufficient to maintain a rowdy populace, but they are carefully trained to use nonlethal force to restrain, rather than kill, all those whom they apprehend. Thus, when war comes to Tal'Dorei, the council appoints a Master of War and levies an army from the Arms of Emon, as well as the militias of its constituent settlements. They must report to the House of Discipline, where they are taught martial combat alongside the harsh realities of warfare. Those who endure their time in the House of Discipline are forged into an army ready to protect \textit{Emon} and the rest of Tal'Dorei from the greater dangers of Exandria.

\subparagraph{14. Gilmore's Glorious Goods}
The broad, single-story façade of Shaun Gilmore's wondrous magical emporium belies the extravagance of its interior. Within, the air is thick with a dozen competing perfumes, each more pungent than the last. The interior is impossibly large—certainly larger than its exterior walls would have you believe—and filled with seemingly endless rows of arcane curiosities and artifacts, many of whose cryptic functions have been lost to time. Everburning candles in an array of unnatural colors light the shop, casting tantalizing shadows over every bubbling phial and mystical orb. Those who enter the shop are first met by Gilmore's long-time assistant Sherri, a half-elf draped in deep purple robes. If they are lucky, or particularly convincing, they may even meet the hero Gilmore himself.

Even rarer a sight than Sherri or Gilmore is the latter's husband, a handsome elf named Darius. He is shy in public, a perfect counterbalance to Gilmore's bombast, but is said to have a singing voice and skill with the \textit{lyre} that can make any mortal weep. Whether or not his husband is present, Gilmore is ever eager to recount his part in the fall of the \textit{Chroma Conclave}, often winkingly embellishing his role in the final battle with \textit{Thordak} himself.

Gilmore's Glorious Goods sells all manner of magic items. Its stock varies, but typically includes magic items of very rare rarity and lower for the price range listed in the fifth edition core rules. Additionally, there are always unidentified items in Gilmore's vast back-room storage that he says he'll "get to eventually;" these may include one or two long-lost legendary items. Artifact-quality items are beyond even Shaun Gilmore's usual fare, but he might know clues to a specific artifact's location, whether through his own extensive knowledge or his rumored relationship with the \textit{Tal'Dorei Council} as an expert consultant on matters of arcane objects. Magic items are expensive, and those strapped for cash may inquire about sponsorships and quests done in Gilmore's service.

\subsubsection{Emon Adventures}
Every district, point of interest, and scrap of current events described earlier in this section could be a plot seed you can turn into an adventure. In addition to those, you can use these adventure hooks to involve your players in an \textit{Emonian} adventure.

\subparagraph{Primordial Children (any level)}
Rumors of terrible creatures of fire and teeth stalking the sewers spread throughout \textit{Emon's} taverns. Such mutterings have gone largely ignored; however, an investigator has recovered what looks to be the broken remains of a large, leathery red egg, not far from a collapsed tunnel beneath the Cloudtop. It seems one of \textit{Thordak's} primordial dragon spawn survived its father's demise and is lurking in the city's tunnels.

\subparagraph{Council Business (any level)}
The \textit{Tal'Dorei Council}—particularly \textit{Tofor Brotoras}, Master of Defense, and \textit{Allura Vysoren}, Master of Arcana—frequently hire adventurers to undertake dangerous quests for the benefit of the realm. Once the characters do a deed great enough to catch the council's attention, they are contacted the next time they're in \textit{Emon} to undertake a quest; choose from any of the other story hooks in this book or make a new one of your own. The council offers a reward equal to 1,000 gp × the average character level of the party.

\subparagraph{Miraculous Escape (high-level)}
Master of Arcana \textit{Allura Vysoren} covertly contacts the characters with a worrisome claim. In the middle of the night, a former \textbf{Remnant cultist} named Hammond Kraith escaped the \textit{Black Bastille}, apparently with the aid of an invisible monster. She provides the party with the magic-inhibiting manacles he wore within the prison, in hopes that they can use them to follow his trail. What she doesn't know is that the \textit{League of Miracles} broke Kraith out of prison using an invisible \textbf{master adranach}. He is hiding out on an island in the \textit{Lucidian Ocean}, aided by a number of \textit{Myriad} mercenaries and a captured \textbf{mage hunter golem} he has reprogrammed to serve him.

\begin{DndSidebar}[float=!b]{The Sisters Grast}
Mythic tales of a pair of ancient hags swirl through learned circles within \textit{Emon}, typically to frighten \textit{Lyceum} students off of wandering forbidden vaults at night. Some folk, however, take these whispers seriously. A number of explorers and treasure seekers wandering the subterranean ruins of \textit{Crystalfen} have encountered one, or both, of these wretched witches, returning with tales of fright, dark dealings, and unholy pacts.



In reality, these creatures are elven scholars who stumbled upon the \textit{Ruins of Salar} long, long ago. Their discovery and research of the profoundly magical location twisted their minds and bodies, infused them with power, extended their lives, and granted a glimmer of knowledge of some immense, unknowable purpose to the lost city. They've become obsessed with the mysteries of the ruins. The elder, Trysta, remains below to continue to excavate and study, while her sister Forscythia travels back and forth from \textit{Emon} under the veil of a half-dozen illusions to acquire goods and gold in exchange for her unsettlingly insightful fortune-telling. Anyone who seeks the sisters with ill intent is rarely seen again, but those who stumble upon them often find themselves temporarily in service to them, a pact made out of fear and self-preservation.



\end{DndSidebar}

\subsection{Crystalfen Caverns}
Deep beneath the \textit{Bladeshimmer Coast} is a sprawling network of natural caverns and rivers that predate the \textit{Age of Arcanum}. These tunnels reach even below the ocean, and have roots in early colonization by denizens of the Realms Beyond.

An ancient society of psychically powerful \textit{aboleths} and other aberrations sprung up from a door between worlds and conquered this subterranean web, drove their slaves to construct \textit{Salar}, the Unseeable City, and slowly pressed upward toward the surface. Their plot was unintentionally ended when the final battle of the \textit{Calamity} sent powerful waves of magical force throughout Exandria, causing much of the underground caverns to collapse and the aboleth civilization to fall under rock and rubble. Fragments of the City Which Eyes Cannot See still remain, however, with surviving terrors slowly rebuilding from the dust within their once-great capital.

The ruins of this sprawling city are now known as the Crystalfen Caverns. These caves are home to scattered veins of \textit{azuremite}, a gorgeous blue crystal that formed from millennia of psychic energies existing near element deposits. Curious explorers who discovered the veins found that, when mined and refined into a fine dust, the \textit{azuremite} powder is a strong mind-altering agent and induces temporary visions and other psychic phenomena.

While a number of small, isolated entrances to this labyrinth of tunnels have been discovered and adventuring parties have attempted to chart the mines, the caverns are so vast and deep that either they gave up for fear of becoming lost, or they were assailed by the terrible denizens of the caverns, never to return. Most who now brave the caverns are foolish treasure hunters or criminals seeking to mine more \textit{azuremite} to sell on the black market.

\subsubsection{Crystalfen Caverns Adventures}
Game Masters who set their adventure in the \textit{Crystalfen Caverns} can use this plot hook for inspiration.

\subparagraph{Mining The Past (any level)}
When the characters finish a resting in the \textit{Crystalfen Caverns}, they must each make a DC 15 Wisdom saving throws. On a failure, a woman's voice fills their minds, saying, "I sense someone...are you truly here, too? Please...I am lost, and alone in the dark." The voice guides them down the twisting caves to a 120-foot-deep mine shaft. The stones at the bottom shimmer with faint azure light. Should they descend into the \textit{azuremite} mine, they travel through an uncharted tunnel and find a lone, dying \textbf{aboleth} with five levels of exhaustion that has wandered too far from its ruined domain.

The \textbf{aboleth} doesn't seek to kill them—yet. Its only priority is self-preservation, and it needs help fending off the \textbf{kobold} that have made the mine their home in order to reach a subterranean lake where it can recover. Do the characters kill it or save it? If it reaches the lake, or dies, they are rewarded with a haunting premonition of the \textit{Ruins of Salar}, and a murky vision revealing the depths and expanse of the \textit{Crystalfen Caverns}.

\subsection{Ruins of Salar}
In the lightless nadir of the \textit{Crystalfen Caverns}, directly beneath the splendid \textit{Port of Emon}, rest the eternally slumbering ruins of Salar. No mortal has beheld these ruins—once the capital of a mighty \textbf{aboleth} dominion—and returned to the surface with wits to tell the tale. Salar is known only in myth and legend as the Unseeable City, because mortal eyes cannot look upon its alien architecture without destroying their fragile minds. Salar is treated as an urban legend by \textit{Emon's} scholars and archeologists. Nevertheless, many believe in its existence, for numerous scholars have returned from expeditions into the deep caverns with their minds shattered by otherworldly power.

Those who have returned from Salar brought back journals filled with increasingly crazed notes. The most intelligible of their writings indicate that all that survives of the ruined city are emerald marble structures of unfathomable design, seeming to shift and shimmer as the eye tries to comprehend, moving through time as they remain static in space.

The once-renowned archaeologist Karaline von Ethro claimed to have visions of Salar, describing a subterranean city with twisting pillars that stretched higher than any tower, kissing the seat of the surface world from the bedrock of the earth. At the peak of the central pillar floats a giant pearlescent disc. This so-called Moon Disc bathes the entire city in unsettling, milky light that refracts eerily through the emerald spires.

Karaline von Ethro is long since dead, and her journal—by far the most coherent record of the Unseeable City—is an apocryphal artifact. The \textit{Arcana Pansophical} kept it in a closely guarded vault, which only added to its mystique. However, the journal was stolen long ago, and since then, dozens of manuscripts said to be the Von Ethro Journal have been sold in black markets across Tal'Dorei. Anyone lucky—or unlucky—enough to find the genuine journal would learn a terrible secret that the \textit{Pansophical} has held onto in silence for decades: dozens, perhaps hundreds, of \textit{aboleths} survive within Salar. Their master, Durrom of the Emerald Eyes, lives on in secrecy while expanding its army of thralls and learning all it can about the ignorant surface dwellers. Durrom's ultimate goal is to shatter the forgotten Pillars of the Earth beneath \textit{Emon} and cause the city to collapse into its subterranean domain.

\begin{DndSidebar}[float=!b]{Karaline von Ethro}
Archeologist, anthropologist, and explorer extraordinaire; Karaline von Ethro was a human adventurer and inventor. It was her designs that created the Sunbeam Compass, a device that always directs the user to the sun by emitting a small light at the tip of the needle in a glass orb, allowing it to rotate in all directions—a very useful tool when exploring subterranean territory. However, despite her discoveries and inventions, she is most remembered for going mad after allegedly having visions of the Ruins of Salar. Karaline's madness eventually led to her tragic demise at the young age of 37.



\end{DndSidebar}

\subsubsection{Ruins of Salar Adventures}
Game Masters who set their adventure in the \textit{Ruins of Salar} can use this plot hook for inspiration.

\subparagraph{Secrets Within Secrets (mid-level)}
The characters are approached by Jomen Tash, an earnest half-elf mage of the \textit{Arcana Pansophical}, who claims to have obtained pages from the lost Von Ethro Journal. He shares his research with the characters and hires them to accompany him into the ruins, promising that he will guard their minds from any alien influence. But in truth, though Jomen was once a member of the \textit{Pansophical}, he has recently aligned his interests with the mysterious goals of the \textit{League of Miracles}. Is he leading the characters into danger with only counterfeit scrawlings to guide them, or does he actually hold part of Karaline von Ethro's description of \textit{Salar} and its dangers?

\subsection{Scar of the Cinder King}
It was thought that the lands surrounding \textit{Emon} would never fully recover from \textit{Thordak's} overwhelming elemental power. His presence not only incinerated his enemies—it wounded the land itself. After the \textit{Cinder King's} defeat, the \textit{Tal'Dorei Council} beseeched the Fire Ashari to travel from the distant lands of \textit{Issylra} to \textit{Emon} and lend their wisdom.

The frightful \textit{Scar of the Cinder King} and \textit{Thordak's Crater}, both twisted by the red dragon's cancerous, fiery magic, were healed by the druids of the Fire Ashari over the course of about fifteen years—an act that the brightest minds of the \textit{Arcana Pansophical} thought would take an epoch to achieve. Clouds of white ash still occasionally blow across the fields surrounding \textit{Emon}, but the last lingering reminders of the \textit{Cinder King's} terror are beginning to vanish.

Nevertheless, the Fire Ashari left behind a group of four druids, led by a female half-orc named Lorkathar. They dwell in a small stone fortress known as Flamereach Outpost—a defensive fortification used when the scar was still spawning beings of flame and hatred. They vehemently warn all who enter to avoid using elemental magic of any sort—especially fire magic—as these evocations could cause the still tender division between the Material Plane and the Elemental Plane of Fire to resonate and tear.

\subsubsection{Scar of the Cinder King Adventures}
Game Masters who set their adventure in the \textit{Scar of the Cinder King} can use this plot hook for inspiration.

\subparagraph{Reopening Wounds (any level)}
\textit{The Ashari} of Pyrah have claimed that the \textit{Scar of the Cinder King} is healed—yet rumors have spread through \textit{Emon's} \textit{Erudite Quarter} that a representative of Pyrah has been seeking counsel with the \textit{Arcana Pansophical} regarding the scar. While the characters travel through the once-scarred lands east of \textit{Emon} to find their next adventure, they are shocked by a sudden earthquake, as the ground erupts with a serpentine, fiery crack. Magma and \textbf{cinderslag elemental} slowly seep from the crevasse, turning all they touch to ash.

\subsection{Ruins of O'Noa (New O'Noa)}

\begin{description}
\item[Village:]
Population: 310 (83\% \textit{humans}, 17\% \textit{other races})

\end{description}

The first human colonizers of Tal'Dorei journeyed across the Ozmit Sea, weathering terrible storms and deadly waves, and made landfall in \textit{Daggerbay}. From there, they founded the port city of O'Noa, and though the city quickly grew in wealth and influence, the early settlers were greeted with mistrust by the \textit{elves} of \textit{Syngorn}. When King Warren \textit{Drassig} took the throne of \textit{Emon}, no city suffered greater devastation than the contested region of O'Noa. It was the site of one of the \textit{Scattered War's} bloodiest battles, and the arcane energies unleashed reduced the city to ruins and permanently scarred the land. Built within a cursed bay, and deprived of fertile land, it was agreed after the war that neither the new nation of Tal'Dorei nor \textit{Syngorn} would attempt to rebuild the city, instead letting it stand as a reminder of the horrors of the \textit{Scattered War}.

Over centuries of dereliction, the ruins became a haven for outlaws and exiles, who named their refuge New O'Noa. They took this place of destitution and dust and made it into a refuge for all rejected by society. Though some of New O'Noa's people were exiled from their homes due to fear and prejudice—like people of ill fate born under the \textit{full light of Ruidus}—most are thieves and delinquents. New O'Noa is ruled by a charismatic tiefling called Waken, who rose to prominence through honeyed words and selective murders.

New O'Noa still appears abandoned to the eyes of outsiders. Its inhabitants live and hide away from prying eyes, amidst ancient stone ruins and their catacombs. Those who venture too deeply into the ruins are rarely seen or heard from again—unless their loved ones can offer up a princely ransom.

\subsubsection{Ruins of O'Noa Adventures}
Game Masters who set their adventure in the Ruins of \textit{O'Noa} can use this plot hook for inspiration.

\subparagraph{This Too Solid Flesh (mid-level)}
New \textit{O'Noa} is a seedy town filled with strange magic. A \textit{suude} dealer named Shef Silverleaf is in major debt to Waken, the leader of the outlaw village. The characters enter the picture when Waken surreptitiously hires them to collect on his debt—perhaps under false pretenses to hide his illegal operation.

When the party tracks down Shef, they find more than they bargained for: his body has been turned to solid \textit{silver}. A \textbf{Remnant cultist} lurks in the shadows of New \textit{O'Noa}, and has been purchasing \textit{suude} from Shef for nearly three months. Shef uncovered his secret identity and confronted him—but the mage's magic, made chaotic by the \textit{suude}, transformed Shef into solid \textit{silver}. Waken will take Shef's solid \textit{silver} corpse as payment—but what if the silver-skinned man were still alive and able to magically walk away? And what of the cultist who transformed him?

\subsection{Seashale Mountains}
A curious range of mountains runs along the \textit{Bladeshimmer Shoreline} north of \textit{Emon}. The Seashale Mountains are named for their unique shape, akin to a tidal wave about to crash on the inland territories, frozen in dynamic stone.

Crashing waves and salty sea breezes have carved miles of twisted valleys between these jagged peaks, creating an unusual home for the creatures lurking within the cavernous rock, shoreside pools, and towering bluffs. High tide leads the briny waters into the Nightwash, a miles-long lagoon that wends its way through the base of the bluffs, carrying hosts of dangerous ocean life with each pass. Conversely, the tidal phenomenon also brings a valuable bevy of deep-sea fish on occasion, which gave birth to the fishing community of the Shalesteps. The mountains' strange peaks are decorated by copses of scraggly trees, stunted by the gravelly soil and stripped nearly bare by salty ocean breezes.

\subsection{The Shalesteps}

\begin{description}
\item[Village:]
Population 770 (40\% \textit{humans}, 35\% \textit{halflings}, 32\% \textit{dwarves}, 3\% \textit{other races})

\end{description}

The Shalesteps are a series of small, allied villages whose inhabitants chose to live a more simple life outside of the bustle of the capital city, but close enough to enjoy imports and protection should the need arise. Most of their food is sourced from the ocean and the Nightwash tide, while a small hunting community wanders the peaks for avian quarry and other sources of meat and natural resources. Those who live here are sturdy and ready to defend themselves against the dangerous wilds surrounding them—and against the occasional monstrosity that washes in with the tide.

\subsubsection{Shalesteps Adventures}
Game Masters who set their adventure in the Shalesteps can use this plot hook for inspiration.

\subparagraph{From Hell's Heart (low-level)}
The Shalestep village of Puddlefoot has never had anything remarkable happen to it. People there are quiet and simple, with no interest in tales of grand heroics. No great adventures have ever come to their rocky shore—until the whale. The dead whale, bloated with corruption, its white hide covered in oozing green lesions, is a source of fear for the nearby townsfolk. What happened to this behemoth? And why is Rill the Hermit ranting about seeing tall shadows spying on Puddlefoot from the mountain peaks?

\subsection{Visa Isle}
Forty miles off the coast of the \textit{Bladeshimmer Shoreline}, west of the \textit{Verdant Expanse}, two islands belch clouds of black smoke into the open sky. \textit{Visa Isle} and its smaller sister island have been surrounded by rumors and sailors' legends since the dawn of civilization. Some say that every night, the island burns to ash—giving it a distinctive orange glow—and that every day new trees grow to maturity by sundown. It is an island of fire and death.

Those who have explored its dense jungles and returned claim only to have found common wild pigs and—at worst—unusually large fire beetles and mantises. They were the lucky ones. They didn't wander unprepared into the Ruins of Vos'sykriss, and did not encounter the serpentine spirits that haunt its elder halls. The only expedition to have ventured into the Ruins of Vos'sykriss and lived to speak of it returned to \textit{Emon} just as the \textit{Chroma Conclave} descended upon the city. Their leader, a half-orc archaeologist and spiritspeaker named Jorlund Vohr, survived the destruction of \textit{Emon}, but all his notes and the artifacts he recovered were stolen. Decades later, the aging professor is still hard at work creating endless proposals for expeditions to \textit{Visa Isle}.

\subsubsection{Visa Isle Adventures}
Game Masters who set their adventure in \textit{Visa Isle} can use this plot hook for inspiration.

\subparagraph{Awakening the Sleepers (mid-level)}
In ancient times, Vos'sykriss was the seat of a serpentfolk empire that dominated southern Tal'Dorei. Before the \textit{Calamity}, their seers had visions of the coming destruction and desperately sought a way to survive it. The serpents constructed a magical stasis field beneath their city and chose their strongest and most powerful to slumber and rebuild the empire when the danger had passed.

They are now awoken. The \textbf{Vos'skyriss serpentfolk} now prowl the island, being prepared by the lingering \textbf{Vos'skyriss serpentfolk ghost} to claim the world for their own. When Jorlund Vohr's expedition delved into Vos'sykriss's deepest sanctum, they unwittingly ended the serpents' ancient spell and awakened the Sleepers. Soon, a people long thought extinct will return to claim the world as their own.

\section{Other Lands of Exandria}
The land of \textit{Gwessar}, now known far and wide as the young nation of Tal'Dorei, is a sprawling, curious land full of adventure, danger, and glory. Though Tal'Dorei is a land rife with grand tales, the world of Exandria is far larger and more ancient than just the lands between the \textit{Lucidian Coast} and the \textit{Bladeshimmer Shoreline}. Many of the people who live here either descend from or are immigrants to this realm.

Some of these other lands have been explored by Vox Machina and the Mighty Nein in \textit{Critical Role}, but are beyond the scope of this book. You can populate these lands with dungeons, cities, characters, and stories of your own creation—or be inspired by the history of Exandria in \textit{chapter 1} of this book, and the stories in the campaigns of \textit{Critical Role}, to create a version of Exandria entirely your own.

\subsection{Issylra}
Northwest of Tal'Dorei, across the freezing waters of the Ozmit Sea, lies Issylra, home to Vasselheim, commonly called the oldest surviving city on Exandria. Issylra is considered by most to be the land where all civilization emerged and spread in the time of the \textit{Founding}, making it the religious core for many, and the focal point of historical study for others. It was also where most of the flying cities of the \textit{Age of Arcanum} first rose from the ground. Its people rose high—and fell hard from those dizzying heights, and to this day, many there still hate and fear the power of arcane magic.

The territory known as \textit{Othanzia} encompasses a majority of the continent's developed regions, keeping most domains protected and governed by the might of the theocracy of Vasselheim. Nestled among eastern taiga forests of the Vesper Timberland, the city of Vasselheim stands as the oldest known seed of civilization and religion, and the only city to withstand every catastrophe that has rocked Exandria since the gods brought life to it. As the spiritual center of most of the world, Vasselheim is ruled by a coalition of high priests of the many \textit{Prime Deities}. These hierophants are viewed by their people as the defenders of their virtuous society, and of all of \textit{Othanzian} culture.

Beyond Vasselheim's ancient walls are frigid, untamed wilds, some of which are untouched since the \textit{Founding} itself. These lands are populated by ancient monsters and dire beasts. Smaller camps, communities, and townships are scattered across the surrounding snowy plains, and those who call \textit{Othanzia} home have become hardy and self-sufficient.

Elsewhere in \textit{Othanzia}, the lands of Issylra are sparsely populated by scattered republics of varying morality. Traveling between these city-states is difficult, due to the lingering beasts of the \textit{Calamity} that stalk the lands beyond Vasselheim.

\subsection{Marquet}
Following the winds and waves to the southwest of Tal'Dorei across the Ozmit Sea eventually leads one to the chain of islands called the Hespet Archipelago and to the shores of the continent of \textit{Marquet}. The pebble beaches of the Bay of Gifts greet travelers with the charming sights of untouched nature, and the port city of Shamal invites all to engage in the pleasures of civilization, with beautiful ballrooms and gilded gambling halls.

Beyond the port of Shamal lies the rocky edge of the Aggrad Mountains, whose canyons and caverns house untold numbers of predatory creatures, before giving way to the immense Marquesian Desert. This sandy wasteland is rumored to have once been verdant jungle brought to ruin by the spear of the The Ruiner in the \textit{Calamity}. Some stories tell of a great hero that sacrificed himself to save \textit{Marquet} from utter annihilation, but only scattered fragments of this tale endure.

Near the center of this intense desert, surrounded by a sprawling network of small villages, is the grand city of \textit{Ank'Harel}, known as the Jewel of Hope. This city, built atop the Drowned City of Cael Morrow, is the center of culture, history, and power for all of \textit{Marquet}. Ruled by the mysteriously undying J'mon Sa Ord, \textit{Ank'Harel} is a magnificent haven of music and culture. It's not without its dangers, but it is a paradise compared to the monster-infested sands that surround it.

The Hands of Ord, the city's guards, also protect nearby villages and townships within the desert, but they can only reach so far. Further south, beyond the reach of \textit{Ank'Harel} and the Hands of Ord, the sands give way to scrubland and marshes tangled with the volatile mountains surrounding the smoking Suuthan Volcano and the dangerous clans that worship it. Westward, beyond the arid desert, lies a vibrant landscape of verdant jungles, craggy mountains, and magical mysteries that the people of \textit{Marquet} have spent eons unraveling. Within the jungles of \textit{Marquet} are dozens of cultures, some nomadic, some seafaring, and some with great cities of their own. One such city, founded by orcish survivors of the The Ruiner wrath during the \textit{Calamity}, is known across Exandria for its mathematically immaculate architecture. This realm welcomes scholars and travelers from across Exandria, who are eager to learn from the minds of some of the greatest architects the world has ever seen.

\subsection{The Shattered Teeth}
Few sail into the dangerous waters southeast of Tal'Dorei, past the ever-present fog bank called the Fool's Curtain, and into the Shattered Teeth. This cluster of broken lands is made up of forty-three islands, varying in size from small, mile-wide reef-toppers to the large, city-holding isles like Ruukva. Two societies dwell within the Shattered Teeth. The first is the Ossended Host, a vast fishing culture that values their independence—even to the point of isolationism—and self-empowerment through the worship of the power of dreams and nightmares. The second is the Wanderman Assembly, a centuries-old trade company that, upon being stranded across the Shattered Teeth after losing their way in a typhoon, grew into a brutal, ruthless society dedicated to gaining profit at all costs. It masquerades poorly under the banner of honor and brotherhood.

This tension has led to bloodshed in recent years, and each group is paranoid that the other is bent on destroying their way of life. What brims beneath the surface of this conflict is the nature and source of the dreams the Ossended Host have been harnessing for generations, and what that answer may mean for the struggles of the people of the Shattered Teeth.

\subsection{Wildemount}
Due east of Tal'Dorei, the lands of Wildemount are highly geologically varied, from the tropical Menagerie Coast, to the dour and rainy valleys of Western Wynandir, to the blasted wastes of Eastern Wynandir—better known as Xhorhas. The people of Tal'Dorei trade frequently with Wildemount, though the waters of the \textit{Lucidian Ocean} are prowled by pirates, sea monsters, and forgotten aquatic demigods. Such a journey leads to the Clovis Concord, a federation of city-states that govern the Menagerie Coast, where a lively culture of trade, entertainment, and the arts flourishes amidst the illicit activities of local smugglers and the ever-present \textit{Myriad} syndicate.

The Dwendalian Empire, an authoritarian power that seeks to expand its reach across all of Wildemount, reigns over the center of the continent. The empire is ruled by King Bertrand Dwendal, but tensions boil between the Crown and the Cerberus Assembly, the empire's league of arcanists and covert operatives. Though the assembly's power waxes and wanes, they are constantly maneuvering to become the true, shadowy power behind the Dwendalian Crown.

The empire's control over the resource-rich valleys in the center of Wildemount has helped make it a military superpower. Likewise, the natural defenses of the surrounding mountains make their lands easily defensible. For centuries, the idea that the empire is invincible has fueled recruitment propaganda for its military, the sanctimoniously titled Righteous Brand. King Dwendal's edicts have created a highly nationalistic society—one which permits only the worship of Crown-approved deities, submits to mandatory military service, and jingoistically endures immense taxes to support the Crown's wars.

Across the mountains east of the Dwendalian Empire are the wastes of Xhorhas, ruled by the Kryn Dynasty. This society was founded by dark \textit{elves} who escaped the grasp of the The Spider Queen through the grace of a supernatural being of light called the The Lost Beacon of Unknown Light that granted them the ability to reincarnate into new bodies, life after life. Those who are reincarnated seek to broaden the knowledge of the The Lost Beacon of Unknown Light by living bold lives and gaining new experiences. Today, the Kryn Dynasty is filled with not just dark \textit{elves}, but goblinkin, \textit{orcs}, ogres, kobolds, and all manner of other peoples considered to be "monsters" by the Dwendalian Empire.

Beyond the Kryn Dynasty and the Dwendalian Empire, and the tensions that burn between them, are the lands of the Miskath Strand and the islands of Eiselcross. Both are realms twisted by forbidden magic, and inhabited by eldritch abominations and other creatures that are anthema to mortal life. Yet, despite the danger, the most courageous adventurers of Wildemount—and indeed, all of Exandria—see exploration as the pinnacle of heroism. At the fringes of the known world, where monsters from the edge of reality lurk...this is where true adventure begins.

\chapter{Character Options}
\DndDropCapLine{T}{}he sprawling lands and endless challenges of Exandria forge heroes and villains of curious vocations and abilities. The dangers of the chaotic world urge stalwart heroes to defend innocents in need of protection, while uncovered fragments of forgotten powers inspire power-hungry villains to seek new ways to bend reality to their will. The material in this chapter offers a number of new options for characters to define who they are and what they wish to become in a Tal'Dorei campaign.


\section{Playing in a Tal'Dorei Campaign}
This section is for players creating characters for a Tal'Dorei campaign that another player is running as the Game Master. It suggests a number of tips and tricks for how to create a character, and how to make that character's story fit smoothly into the world of Exandria.

\subsection{Standard Procedure}
Start by creating a character as you normally would for any fifth edition RPG campaign. This process typically involves choosing your character's race, lineage, or ancestry; selecting your class, background, starting equipment, and spells (if applicable); and giving the character a name and backstory. If your Game Master has any house rules or special requests that change this process, they'll let you know.

\subsubsection{Ancestries}
All of the fantasy peoples commonly found in Tal'Dorei, including \textit{humans}, \textit{elves}, \textit{dwarves}, \textit{goblins}, \textit{tieflings}, and \textit{dragonblood}, are described in this chapter. These descriptions don't include game statistics, which can be found in the fifth edition rules, but rather detail the lore of each of these ancestries and how characters of those ancestries fit into the tapestry of Exandrian life.

\subsubsection{Classes}
All the classes in the fifth edition core rules are fair game in Tal'Dorei, and classes and subclasses found in other supplements might be available with your Game Master's permission. This book also contains the following new subclasses, all detailed in this chapter:


\begin{itemize}
\begin{multicols}{3}
\item barbarian
\item bard
\item cleric
\item cleric
\item Druid
\item monk
\item paladin
\item sorcerer
\item wizard
\end{multicols}

\end{itemize}

\subsubsection{Backgrounds}
The backgrounds found in the fifth edition core rules are all well suited to life in Tal'Dorei. This fantasy realm is teeming with soldier, Guild Artisan, sailor, and noble, all of whom might make interesting adventurers. This book also presents six new backgrounds, detailed later in this chapter:


\begin{itemize}
\begin{multicols}{3}
\item Ashari
\item Clasp Member
\item Variant Clasp Member (Myriad Operative)
\item Lyceum Scholar
\item Reformed Cultist
\item Whitestone Rifle Corps
\item Variant Whitestone Rifle Corps (Grey Hunter)
\end{multicols}

\end{itemize}

\subsubsection{Session Zero}
The cast of \textit{Critical Role} creates great stories by communicating openly with their Game Master before the campaign begins, and maintaining that communication between game sessions. Doing so ensures that everyone's stories are being told in a fun way, and that everyone is having fun at the table. To get that same experience, be sure to talk with your Game Master about what's important to you while you build your character and flesh out their backstory. That way, your GM will be able to build a story that centers you and your allies as main characters, rather than just bystanders.

Session zero is a catchall term for a pregame meetup between the players (including the GM), where everyone can lay out expectations and boundaries for the campaign. Lots of guides designed to give your campaign a strong start with session zero can be found on the internet and in other fifth edition books. Session zero often includes discussions of the following:


\begin{itemize}
\item What each player and the GM like about fifth edition campaigns, and what they're hoping for in this campaign.
\item Likewise, what aspects of play people dislike or are uncomfortable with, or simply want to avoid in this campaign for any reason.
\item If this campaign were a movie, what would its rating be? PG? PG-13? R?
\item Which topics are fun to explore in game, which are absolutely off-limits due to personal discomfort, and which are okay in moderation? Everyone should clearly agree to playing a game that includes topics like extreme violence; romantic or sexual situations; any gender-, race-, disability-, or age-based discrimination or violence; or themes that involve real-life politics or religion.
\item What amount of combat, exploration, and social encounters the players want to experience.
\item What kind of character each player wants to create, and how the characters might fit together into a cohesive party.
\item What type of story everyone wants to experience together, whether that's a meandering episodic tale or a lengthy epic, and whether the story will be driven by player actions and their consequences or by GM-directed plot events.
\end{itemize}

Additionally, after each game, many people find it helpful to check in with each other and discuss what everyone enjoyed, wants less of, or wants more of, so the GM and players can prepare to bring more of what worked to the next session, and less of what didn't.

\subsection{Flavorful Details}
In addition to discussing your character's backstory with your GM and the other players in your group, you can think about other flavorful details to help ground your character in the world of Exandria.

\subsubsection{Gods}
A vast pantheon of deities created Exandria in primeval times. Which of these deities does your character worship, if any? Do they worship an obscure god, or one commonly venerated throughout Tal'Dorei? How devout are they in their faith?

The gods of Exandria are described in more detail in \textit{chapter 2}, though your GM might want to keep some of the details of those gods secret.

\subsubsection{Factions}
The politics of Tal'Dorei are influenced by a number of factions. Not all are actively hostile toward each other, and many are actually allies.

But all jockey constantly for power—and frequently come into conflict as a result. Do you belong to one of these factions? Are you a card-carrying member? Do you simply sympathize with their cause, or perhaps have friends or family who are part of the faction? Are there factions you actively despise, perhaps because of some slight they committed against you or those close to you in the past?

Tal'Dorei's factions are described in more detail in \textit{chapter 2}, though your GM might want to keep some of the details of those factions secret.

\section{Races and Cultures}
Tal'Dorei is home to manifold humanoid peoples, from the \textit{elves} who emerged from the Fey Realm into the untouched lands of \textit{Gwessar} in the wake of the \textit{Calamity}, to the \textit{humans} who first came to Tal'Dorei on stout ships from the distant shores of \textit{Issylra}. The familiar fantasy races of \textit{elves}, \textit{dwarves}, \textit{humans}, and \textit{halflings} are the most populous beings on Tal'Dorei, but they are far from the only peoples to make their mark in city and wilderness alike.

Over the past thousand years of history, the peoples described in this chapter have left an indelible mark on the world. And though many of these races appear monolithic at first glance, all are incredibly varied. A first-time visitor to Tal'Dorei might initially perceive all \textit{elves} as haughty and fascinated with heritage, because that's the dominant culture in \textit{Syngorn}, the largest elf-majority city on the continent. However, the \textit{elves} living in villages and cities across the land are often exceptions to that apparent rule, and it's dangerous—not to mention rude—for travelers to assume that any of Tal'Dorei's varied peoples fit into simple boxes.

This section is a resource for you to use to inspire your character's backstory—or to inspire your NPCs if you're a Game Master. Whether you enjoy leaning into fantasy archetypes, intentionally breaking them, or ignoring them entirely, this section has something for you. The game statistics of each of the races of Tal'Dorei are identical to those in the fifth edition rules.

\subsection{People and Monsters}
For generations of roleplaying gamers in many different RPG campaign settings, certain of the fantasy peoples in the world of Exandria have long been considered "monsters" or "evil races." \textit{Goblins}, \textit{orcs}, and \textit{bugbears} are examples of peoples whose very nature was supposedly tainted by inherent evil, and were thus morally unambiguous targets for adventurers to kill and loot. This mindset treats these people as no different from monsters like \textbf{basilisk} or \textbf{giant spider}.

The humanoid people of Tal'Dorei aren't set up that way. Any creature with the humanoid type is a thinking, feeling person. Creatures with the dragon, giant, or fey type are also unambiguously sapient, and thus can have any alignment. It's just as likely for a human to be a neutral evil \textbf{bandit} as it is for an orc to be a chaotic good \textbf{swashbuckler}, or for a \textbf{hobgoblin} to be a lawful good \textbf{knight}. And if you like, you can take this even further in your games and say that any fantasy creature with an Intelligence score of 6 or higher is capable of free thought, and should have just as wide a range of alignments as other sapient beings.

The stories in \textit{Critical Role} are filled with shades of moral gray, but Tal'Dorei still has plenty of unambiguously evil creatures. Most demons, devils, and undead are suffused with cosmic evil, and always try to subvert righteousness. Likewise, unthinking creatures such as oozes and constructs, and creatures of instinctual intelligence such as beasts, many monstrosities, and plant creatures, all make morally unambiguous foes.

Likewise, there are plenty of evil humanoids in the world. Cultists of evil gods, including the \textit{Remnants} who worship The Whispered One, and merciless, blood-soaked bandits such as the \textit{Ravagers} are clear and immediate threats to both the safety of the player characters and all free peoples of Tal'Dorei. Depending on the tone of your story, the characters might cut such foes down without mercy—or they might try to rehabilitate them. The difference is that these villains' choices, not their ancestries, determine whether or not they are evil.

\subsection{Cultures across Tal'Dorei}
The various peoples of Tal'Dorei don't have racial cultures. Rather, someone's culture is based on the cities and regions in which they spend their formative years. For most of Tal'Dorei's history, members of the different fantasy races lived in largely monolithic enclaves. However, the more recent past has instilled in people the need to break down old barriers, work together, and live together.

All cities and large towns have a majority population of certain races, typically based on the peoples who founded those settlements. This includes such notable sites as the elf-majority city-state of \textit{Syngorn} and the dwarf-majority city-state of \textit{Kraghammer}. Some people in Tal'Dorei might assume that certain races have common cultural traits that unite them, but many others know that this isn't the whole truth. A dwarf raised in \textit{Kraghammer} and a dwarf raised in \textit{Stilben} are likely to have very different attitudes about family, food, faith, and all the other factors that make up a culture.

Even when people of the same race gather in neighborhoods within a city, those people don't necessarily share the common culture of their best-known homelands. In many neighborhoods in the great city of \textit{Emon}, for example, \textit{elves} and half-elves are a majority. But since not all those elf-kin come from \textit{Syngorn}, they don't all share a monolithic \textit{Syngornian} "elf culture."

That said, beings with distinctly magical traits often form cultures that include their unique abilities. This includes folk such as \textit{dragonblood}, with their ability to channel elemental power. And because \textit{goliaths} are known for their powerful, muscular bodies, goliath-majority societies often worship the The Stormlord, god of strength, and engage in festivals that center around feats of athleticism.

\subsection{Dragonblood}
The dragonblood of Exandria are a people of unknown origin, and countless myths surround their ancestral creation. One particularly pervasive tale states that dragonblood were created by the The Scaled Tyrant and the The Platinum Dragon during the \textit{Calamity}, serving as loyal soldiers who eventually earned their freedom. Regardless of how they came into being, though, the dragonblood lived for many centuries in relative isolation in the southern wastes of Xhorhas, on the continent of \textit{Wildemount}.

Today, dragonblood live all over Exandria, though they remain a scarcer sight in Tal'Dorei than in \textit{Wildemount}. Their spread across the world was hastened by the assault against the free peoples of Exandria by the \textit{Chroma Conclave}. During that conflict, the ancient white dragon known as Vorugal the Frigid Doom destroyed the magical city-state of Draconia, the ancestral homeland of most dragonblood.

The destruction of Draconia and the loss of many of its people was a tragedy. But to roughly half the dragonblood of Exandria, it also marked the end of an oppressive regime. Dragonblood of Exandria can either be born with a tail or without. In Draconia, those with tails were called draconians or draconbloods, for they were seen as closer in form and power to true dragons. Those without tails were given the name of ravenites and enslaved. Their blood and toil built the mighty towers of Draconia, yet the ravenites enjoyed none of their city's splendor.

The ruin of Draconia thus marked the start of an era of freedom and self-determination for the ravenites. In Tal'Dorei, it's apparent that any differences between the ravenites and the draconians have nothing to do with bloodline or draconic power. Those differences are cultural, and rooted in a past whose conflicts still haven't been resolved. The oppression that existed in Draconia remains largely unknown within Tal'Dorei to this day, and outside of major cities such as \textit{Emon}, dragonblood can go for months, even years, without seeing another of their kind.

\subsubsection{Draconians}
The tail-bearing draconians who survived the destruction of Draconia have been scattered across the world. In the time that's passed since the fall of the \textit{dragonblood} city-state, a generation of draconbloods have grown up in Tal'Dorei's towns and cities, disconnected from their parents' privileged upbringing. Only a scant few draconians—typically ones able to teleport away from Vorugal's attack—escaped with any sort of wealth. Most other refugees were lucky simply to escape with their lives.

\subparagraph{Clinging to the Past}
A social divide continues to grow between draconians who have wealth and those who don't. Many of those with wealth greedily cling to the shadows of their former tyrannical glory, but other draconians seek instead to bury the past and create their own destiny, on their own merits.

\subsubsection{Ravenites}
Following the end of the Draconian order, the ravenites ventured into the world eager to make the most of their newfound freedom. In the generation since, many have formed a new homeland in Xhorhas known as Xarzith Kitril, content to reclaim the lands of their birth. Those ravenites who departed \textit{Wildemount} have found a new home in the cities and the icy north of Tal'Dorei, where they enjoy the freedom and privileges of equality that this land provides.

\subparagraph{Rebuilding from Scratch}
Most ravenite children grew up poor, either in slum towns at the edges of cities like \textit{Emon} and \textit{Westruun}, or helping their parents laboriously carve out rural lives in the \textit{Dividing Plains} or \textit{Cliffkeep Mountains}.

\subsection{Dwarves}
When the world was nearly rent asunder by the \textit{Calamity}, the \textit{dwarves} of Tal'Dorei disappeared underground to wait out the worst of the conflict. Many dwarven citadels once dotted the subterranean world, but over the centuries, they dwindled, either merging together or falling into decadence and ruin. In the centuries following the \textit{Calamity}, these dwarf clans, isolated beneath the surface of Exandria, gathered into a single, steadfast redoubt: \textit{Kraghammer}. In one way or another, just about every dwarf in Tal'Dorei can trace their family back to that great city-state.

The mountainous stronghold still stands as the heart of settlement in the \textit{Cliffkeep Mountains}, and many \textit{dwarves} consider it a place all their folk should visit at least once a decade. \textit{Kraghammer's} legendary status gives it a highly conservative, unchanging culture, but even so, the dwarf people are far from monolithic.

Some \textit{dwarves} like to distinguish themselves as hill or mountain \textit{dwarves}, based on whether or not they grew up outside or inside of \textit{Kraghammer}. Gurdhe, a Dwarvish word translated into Common as "hill," colloquially means "beyond the mountains" in most Dwarvish dialects. And many \textit{dwarves}, even the so-called mountain-dwellers renowned for living their entire lives underground, make journeys out into the "hills" of Tal'Dorei to trade, gain artistic inspiration, or visit the children of long-dead friends they bonded with decades or centuries prior.

\subsubsection{Kraghammer Dwarves}
The culture of \textit{Kraghammer} is one of isolationism and tradition—even to the point of xenophobia. As a common saying goes: "Those who live beneath the mountain are as unmoving as the mountain itself." Whether they're \textit{dwarves}, \textit{gnomes}, or folk of other ancestries, the people who dwell in \textit{Kraghammer} typically embody this stubborn, dogmatic stoicism.

\subparagraph{Tradition}
\textit{Kraghammer} society is steeped in a culture of \textit{dwarves} protecting their own, and those \textit{dwarves'} traditions of meticulous chronicling, geometrically precise artistry, dynastic great houses, and resilience against even the greatest of threats has proved itself a central pillar of their society. For long centuries, new peoples and new ideas in \textit{Kraghammer} meant potential danger, threatening the death of monarchs and incursion by nefarious subterranean infiltrators. The first time that \textit{Kraghammer} proudly opened its gates to a new people, the oligarch Warren \textit{Drassig} led all of Tal'Dorei into the greatest war since the \textit{Age of Arcanum}. The city's relative safety during the assault of the \textit{Chroma Conclave} only cemented its people's belief that isolation is strength.

\subparagraph{Change}
In the present day, many of \textit{Kraghammer's} residents feel stifled by its dominant culture, which continues to be insular and suspicious of outsiders and new ideas. In the past few decades, a tide of refugees fleeing crises such as the \textit{Wittebak disaster} and the \textit{Cinder King's} attack on \textit{Emon} have brought new voices and perspectives to the formerly unyielding halls of \textit{Kraghammer}. These people now call the \textit{dwarves'} ancestral realm their home, and many of \textit{Kraghammer's} \textit{dwarves} have seen the innovation and transformation that new residents can bring. Change comes slowly in \textit{Kraghammer}, but it does come in time.

\subsubsection{Hill Travelers}
\textit{Dwarves} might leave their underground realm for myriad reasons, with many making homesteads in the foothills of the \textit{Cliffkeep Mountains} or venturing even farther afield. Hill \textit{dwarves} who left the mountain themselves, or who had parents or more distant ancestors who did so, have a certain shared camaraderie in their connection to the open sky above, the sun and wind around them, and the soft loam beneath their feet.

\subparagraph{Exploration}
It's common for \textit{dwarves} who live upon the hills and fields of Tal'Dorei to wander throughout their lives, never putting down roots for more than a decade or two at a time out of their centuries-long lives. \textit{Dwarves} with tanned skin and sun-bleached hair are so common a sight in taverns across the \textit{Dividing Plains} that the Dwarvish word gurdheledr, translated as "hill travelers," has come to mean "mercenary" or "bounty hunter" as commonly as it refers to simple wanderers.

\subparagraph{Community}
\textit{Dwarves} are a common sight in cities such as \textit{Emon} and \textit{Westruun}, and even in smaller towns across central Tal'Dorei. Not unlike mountain-dwelling \textit{dwarves}, these city dwellers are often as unmoving as a granite boulder on a rock shelf—unable to be budged unless given a tremendous push. \textit{Dwarves} typically adopt the culture of wherever they live, though their long lives sometimes make them more resistant to social change than shorter-lived peoples.

\subsubsection{Duergar of Emberhold}
Deep beneath Tal'Dorei is an underground realm lit only by magelight, magma flows, luminescent creatures and fungi, and the roaring fires of duergar forges. It is thought that a group of \textit{dwarves} delved far into the depths of Exandria countless centuries ago, predating even the chronicles of \textit{Kraghammer}. Whether by choice or by necessity, they made a home in the caverns they explored there. Something about the magic that suffuses their eerie realm changed the deep \textit{dwarves}—who are called the duergar now. That name was originally an epithet, bestowed by the mountain \textit{dwarves} who rediscovered and battled those deep \textit{dwarves}, but the duergar now wear it with pride.

\subparagraph{Survival}
The deep caverns of Tal'Dorei are an inhospitable place, filled with creatures whose aberrant forms defy imagination and whose alien minds reject mortal morality. The duergar became hardened by their struggle for survival, and their rejection at the hands of the \textit{Kraghammer} \textit{dwarves} hardened their hearts forever. The smiths of the \textit{Emberhold} do not forge art or works of beauty—only weapons and armor to turn back the onslaught of those who wish to see them dead.

\subparagraph{Alliance}
The duergar have suffered alone for so long, hating all and hated by all. Nevertheless, in the undercurrents of the \textit{Emberhold}, the idea spreads that there must be more to life than fire and death. Many duergar nurture a secret hope of finding some way to ally with their other dwarf kin—if only someone were brave enough to reach out a hand without fear of it being cut off.

\subsection{Elemental Ancestry}
Exandria's lands are rife with elemental power. Primordial energy constantly flows into the world from elemental rifts, and such primal power inevitably alters its surroundings. This is also true of peoples who spend long years in proximity to such sources of power, and of those who travel the planes long enough to find the seed of elemental power take root within their blood. Such folk are the beginning of an elemental ancestry, harboring the potential to give birth to children who are suffused with latent magic.

Children of an elemental ancestry bear the features of their parents' ancestry, but have clearly become something new. Their appearance is altered by their power, granting them skin that shines like the sun on the surface of the sea, hair that crackles like flame, or eyes that glitter like gemstones. Then comes the manifestation of elemental powers, which can make a child the envy of their peers—or someone to be feared.

While people of elemental ancestry are rare in Tal'Dorei, they are most common among \textit{the Ashari} due to the elemental affinity of that folk. Some possessed of elemental power enjoy the attention their uniqueness garners them. But for others, even the most benign interest draws painful attention to their otherness.

\subsubsection{Belonging}
People with an elemental bloodline are often raised by parents in a regular community, and so grow up with all the aspirations and prejudices inherent to their culture. For many such people, the traits that make them unique also make them a pillar of their community—for example, a baker who instinctually understands the will of her oven's fire. Within superstitious lands and cultures, however, a person's elemental powers might make them the target of scrutiny, fearmongering, or jealousy, eventually driving them to strike out into the world to seek a place where they truly belong.

\subsubsection{Called to the Ashari}
Those elemental-touched people who leave their communities often seek out \textit{the Ashari}, the isolated keepers of rifts to the Elemental Planes. Elementally attuned children are born with great frequency to \textit{the Ashari}, who understand the utmost importance of giving those children the education and training that will let them understand their gifts.

All people of \textit{elemental ancestry} are welcome among \textit{the Ashari}, including the Air Ashari and Earth Ashari found in Tal'Dorei. However, even those whose \textit{elemental ancestry} connects them to water or fire can seek out Tal'Dorei's Ashari peoples to help them on their journeys of discovery. In time, such a journey might lead them to the Fire Ashari on the continent of \textit{Issylra}, to the Water Ashari on the Islands of Anamn in the Ozmit Sea, or to Ashari who have journeyed from those lands to Tal'Dorei.

Elementally attuned folk who don't seek or can't find \textit{the Ashari} might instead find their way to the reclusive \textit{giants of Tal'Dorei}. Though the giants do their best to keep their distance from the lives and problems of humanoids, they have been known to take a shine to people whose shared elemental affinities give them common ground.

\subsection{Elves}
When the \textit{Calamity} threatened to annihilate all life on Tal'Dorei, the \textit{elves} of this land gathered in their last remaining city, \textit{Syngorn}, and used long-forgotten magic to transport themselves into the Fey Realm. Time runs strangely in that realm, and for an unknown stretch of history, the \textit{elves} of \textit{Syngorn} lived with their fey kindred in peace. When they eventually returned to Exandria, they found a land devoid of all but animals and the grandeur of nature. The rotten majesty of the \textit{Age of Arcanum}, burned away by the wrath of the gods, had given way to new growth.

To the \textit{elves}, the lands of Tal'Dorei are still called \textit{Gwessar}, and \textit{elves} who dwell here are called the gwes'alfen—the \textit{elves} of the Fields of Joy. For centuries, this folk watched with delight as other peoples populated \textit{Gwessar}, beginning with \textit{dwarves}, then \textit{humans} from across the sea, and so on. Then they looked on in horror as mortal ambition was perverted into greed, and war followed in its wake.

\textit{Syngorn} and the \textit{Verdant Expanse} that surrounds it are the ancestral homeland of Tal'Dorei's \textit{elves}, many of whom never venture beyond that homeland. Young \textit{elves} are most likely to journey afield to see what else the world has to offer, and while most of them return to \textit{Syngorn}, some make new homes in the many towns and cities of Tal'Dorei—or they become hermits, their status as mystical beings with a thousand-year life span commanding an almost mythic reverence from other folk.

The \textit{elves'} tendency to focus on the past has given rise to many stereotypes of those folk being haughty, aloof, and unconcerned with the day-to-day joys and sorrows of the shorter-lived races. But though this is certainly true for some \textit{elves}, it's hard for other folk to consider the perspective of an elf who has personally witnessed whole cultures and civilizations rise and fall. Most \textit{elves} are born with the same passion for life as any other being. But as centuries pass, life becomes a stream of historic events to be witnessed and contemplated, and not to be interfered with unless absolutely necessary.

\subsubsection{Syngorn Elves}
The syn'alfen, as the \textit{elves} of the \textit{Verdant Expanse} call themselves, almost universally revere the ancient history of \textit{Syngorn}, and the story of how its escape to the Fey Realm saved their culture from extinction. No elf who lives in \textit{Syngorn} can forget this past, and \textit{Syngornian} culture can sometimes feel overbearing to those who see it only from the outside. More than a few woodland \textit{elves} also feel stifled by their homeland's unrelenting focus on history, magic, the fey, and the preservation of elvenkind. Such \textit{elves} are most likely to leave \textit{Syngorn} and make their homes in other parts of Tal'Dorei.

\subparagraph{Departure}
Many \textit{elves} who leave \textit{Syngorn} (or who had parents who left) still consider themselves syn'alfen, whether they live in \textit{Emon} or other cities, or even the ancestral homes of other folk, such as \textit{Kraghammer}. The cultural imprint of \textit{Syngorn} is hard to escape, and most of these expatriate wood \textit{elves} hold a certain amount of love for their home even as part of a complicated relationship.

\subparagraph{Reunited Family}
South of the \textit{Stormcrest Mountains} lies the massive jungle of \textit{Rifenmist}, home to the nomadic elf \textit{Orroyen tribes}. The \textit{Orroyen} believe that they split off from the \textit{Syngorn} \textit{elves} centuries ago, with both peoples meeting each other again only within the past century. Though initially suspicious, both elf lines soon came to see one another as longlost family.

The onslaught of the \textit{Chroma Conclave} disrupted communication with the \textit{Orroyen elves}, and news from the south has grown more dire since. \textit{Orroyen} refugees continue to flow into \textit{Syngorn}, fleeing the tyrannical \textit{Iron Authority} that rules \textit{Rifenmist}.

\subsubsection{Lyrengorn Elves}
The lyren'alfen, known alternately as Lyrengorn \textit{elves} or high \textit{elves}, are a people who broke off from \textit{Syngorn} long ago, and who journeyed from southern Tal'Dorei to settle in the frigid \textit{Neverfields}, north of the \textit{Cliffkeep Mountains}. The city of Lyrengorn that they founded is every bit as majestic as \textit{Syngorn} in the south, but for better or worse, Lyrengorn lacks its sister city's millennia of culture—and its cultural baggage.

\subparagraph{Art and Expression}
Lyrengorn is nestled in the spires of the \textit{Elvenpeaks}, spreading under the mysterious northern lights known as the The Moonweaver Ribbons. The city is always aglow with magical light, and its people are taught spellcraft from a young age, for magic is the foundation upon which Lyrengorn was built. The spells of heat and comfort woven into its residents' clothing and architecture are essential to living comfortably in such relentlessly cold lands.

The \textit{elves} who make their home within the pines of the \textit{Elvenpeaks} are best known for riding through the wintry skies on \textbf{wyvern}. Every year, people from across Tal'Dorei travel to the \textit{Elvenpeaks} to watch skyswimmers, elite \textbf{wyvern} riders who shape the glow of the The Moonweaver Ribbons from the air.

\subparagraph{Magic and Transformation}
High elf culture values magical prowess above almost all else, but the high \textit{elves} lack the dogma of many of their \textit{Syngornian} cousins. Their society encourages experimentation, both in terms of magic and self-expression. Likewise, lyren'alfen who become adventurers often do so in pursuit of inspiration or self-discovery. High elf societies are filled with elf magi who use magic to transform their bodies, exploring different gender expressions and social roles—and they encourage all who visit Lyrengorn to join them in discovering all the different facets of one's self.

Many high \textit{elves} become religious sorcerers, exploring their latent powers through their faith. Others develop their arcane skill by seeking admission to the \textit{Alabaster Lyceum} in \textit{Emon}, where they study the theory of magic with the best minds from across Tal'Dorei.

\subsubsection{Dark Elves}
The myrk'alfen, also known as drow or dark \textit{elves}, have a complicated relationship with the other elven societies of Tal'Dorei. Though the \textit{elves} are one people by heritage, the dark \textit{elves} refused to join with the \textit{elves} of \textit{Syngorn} during the \textit{Calamity}, instead escaping into the trackless caverns that stretched beneath Tal'Dorei. There, they became the unwitting prey of the The Spider Queen, a Betrayer God of immense cruelty.

\subparagraph{Betrayed by the Gods}
In times long past, the drow were a wise, beautiful folk with long, silvery hair and radiant ashen or violet skin. But their underground enclaves grew decadent and cruel, and their leaders fell to the alluring whispers of the The Spider Queen. The myrk'alfen allied with their venomous goddess against the The Arch Heart and their \textit{elves} in the \textit{Calamity}. When the The Spider Queen was defeated, the myrk'alfen of Tal'Dorei were exiled permanently from the surface world, confined to their imperious realm within the earth.

\subparagraph{Struggle to Survive}
More evil powers than just the The Spider Queen lurk beneath the surface of the world. With the drow in exile, the whispers of the The Crawling King and the The Chained Oblivion first reached the ears of their nobility as the dark \textit{elves} struggled to stave away the aberrations that encircled their cities. Now, generations later, the dark \textit{elves} are a people seemingly on the brink of annihilation. Neighbors slaughter one another in the streets as they succumb to paranoia. For when aberrations can take or shape any form, who can be trusted? Unable to stop their citizens from rioting, the drow elite have grown ever more authoritarian, commanding their royal guards to keep order by violently suppressing their people.

Against this ongoing chaos, some dark \textit{elves} who fall deep into the abominable thrall of the The Spider Queen become truly monstrous, their skin and eyes turning deathly pale as they become little more than puppets for their tormenting goddess. Others willingly offer themselves to the aberrations to end their suffering, and are transformed into shapeshifting \textbf{doppelganger}. And the most power-hungry drow nobles often succumb entirely to their devotion to the The Spider Queen, willingly drinking her abyssal blood to become \textbf{drider} in her service.

\subparagraph{New Light for the Shadow}
The people of Tal'Dorei have recently received emissaries from dark \textit{elves} who live in the eastern reaches of \textit{Wildemount}, in a nation called the Kryn Dynasty. Their Bright Queen, Leylas Kryn, broke free of the The Spider Queen corrupting influence and led her people to glory and self-determination in the distant lands of Xhorhas. This rebellion was inspired by the discovery of the The Lost Beacon of Unknown Light, a mysterious being of light and rebirth that the Kryn drow worship as a god.

Though the dark \textit{elves} of the Kryn Dynasty have problems enough of their own, some of the Kryn seek the aid of \textit{Syngorn} and Lyrengorn on expeditions into the depths beneath Tal'Dorei. Their goals are noble, but none can deny that these are missionary expeditions in search of new followers—and in search of new The Lost Beacon of Unknown Light beacons, the scattered fragments of that new drow god.

Many adventuring drow harbor a violent hatred of aberrations, and have abandoned the The Spider Queen faith to become paladins in service of the The Dawnfather or the The Lost Beacon of Unknown Light, setting aside their affinity for the shadow to search for hope in the light of the sun.

\subsection{Firbolgs}
Fur-covered folk with floppy ears, flat noses, and long faces, firbolgs are a humanoid people with a unique, somewhat bovine appearance. They aren't insular by nature, but they feel most at ease when surrounded by the natural world, leading the vast majority of firbolgs to never leave their forest homes except under circumstances beyond their control. As a result, many folk in Tal'Dorei don't even know that firbolgs exist. Those firbolgs who do live in the realm's cosmopolitan cities either grow annoyed or accepting of people constantly questioning who—or what—they are.

\subsubsection{Dwellers of Cold Forests}
\textit{Firbolgs'} fuzzy bodies are acclimated to cold weather, and they often overheat in warm climes. They are most commonly seen in \textit{Kraghammer}, \textit{Zephrah}, \textit{Whitestone}, and \textit{Stilben}, both as residents and as traveling bards and merchants. But they are otherwise a rare sight south of the \textit{Cliffkeep Mountains}, though some \textit{firbolgs} enjoy a comfortable life in \textit{Emon} after learning a few handy spells to keep cool in the summer.

The largest and densest population of \textit{firbolgs} in Tal'Dorei is found in the eternally winter-bound forest of the \textit{Frostweald}. These \textit{firbolgs} have a friendly but cautious relationship with the capricious fey who also dwell there.

\subsection{Gnomes}
Gnomes and gnomish culture are relatively uncommon in Tal'Dorei, for although \textit{humans} and rock gnomes both hail from \textit{Issylra}, relatively few gnomes have the chance or the desire to leave their homeland. Moreover, their short stature has often proved a social hindrance in the lands of taller folk. But where they have settled, gnomes easily prove their mettle through deed and humor, becoming welcome members of the settlements where they choose to make their homes.

\subsubsection{Rock Gnomes}
The first \textit{gnomes} to enter Tal'Dorei migrated from \textit{Issylra} during the \textit{Age of Arcanum}, and endured the \textit{Calamity} within the deep caverns of the \textit{Cliffkeep Mountains}. These folk called themselves rock \textit{gnomes}, and founded a marvelous settlement called \textit{Wittebak}. Their society prized cleverness above all other things, eschewing the arts of magic for a focus on tinkering and an understanding of natural phenomena.

\subparagraph{Dizzying Heights}
By the end of the \textit{Calamity}, the city of \textit{Wittebak} was the most technologically advanced place in all of Tal'Dorei, even as the \textit{elves} and \textit{dwarves} were emerging from their isolation. However, such greatness could not last. The rock \textit{gnomes} boasted of their superiority to the "ignorant" stone giants who also called the mountains home, and incurred their elemental wrath. In the earthquakes that followed, \textit{Wittebak} and most of its people were destroyed.

\subparagraph{Eyes Ahead}
Most of those who survived the \textit{Wittebak} disaster made their way to \textit{Kraghammer}, relying upon the generosity of the \textit{dwarves} to survive. A few others traveled to \textit{Emon} and \textit{Kymal}, and helped build those cities in their earliest days. To this day, most \textit{gnomes} with any connection to the heady era of \textit{Wittebak} still call themselves rock \textit{gnomes} in its honor. These \textit{gnomes} carry a certain cultural memory of their city's fall from grace, which drives many of them to strive constantly toward greatness.

\subparagraph{Eyes Behind}
A fair few rock \textit{gnomes} dwelling in \textit{Kraghammer}—elders in particular—take a different view of their people's fall from grace and glory. These \textit{gnomes} work bitterly and tirelessly, fueled by spite and pride, with the intention of either reclaiming \textit{Wittebak} or proving that the city's lost technology is superior to any magic or other force on Exandria.

\subsubsection{Forest Gnomes}
Forest \textit{gnomes} are not native to Tal'Dorei, but unlike rock \textit{gnomes}, they are believed to hail from the Fey Realm, much like the \textit{elves}. A half dozen small forest gnome communities can be found in the \textit{Verdant Expanse}, hidden from the prying eyes of \textit{elves} and human explorers alike. But how \textit{gnomes} ended up in the Fey Realm to begin with is a mystery that even they can't begin to answer.

\subparagraph{Magic, Not Tech}
The biggest difference between the \textit{gnomes} of the forest and their rock gnome cousins are their philosophies toward magic and technology. After spending untold centuries in the Fey Realm, forest \textit{gnomes} have a great respect for magic, especially capricious, unpredictable magic like that of the fey. Forest gnome culture sees technology as predictable and boring, even if its usefulness must be conceded from time to time.

Forest \textit{gnomes} respect the rock \textit{gnomes'} dedication to science, and are more than happy to help their kin in their search to understand how nature works and how to manipulate it. But in the end, knowledge is not as important as canniness in forest gnome enclaves.

\subparagraph{Worldly Explorers}
Not all forest \textit{gnomes} are as close to their fey identity as the \textit{gnomes} of the \textit{Verdant Expanse}, and a number of more worldly clans have spread to other forests across Tal'Dorei. Most notably, a number of families, including the Trickfoot clan, have made the \textit{Bramblewood} outside of \textit{Westruun} their home for generations. These clans are more practical than their fey kin and less uptight than their distant rock gnome cousins, giving these \textit{gnomes} a curious personality perfect for the life of a traveling adventurer.

\subsection{Goblinkin}
\textit{Goblins}, \textit{hobgoblins}, and \textit{bugbears} are three distinct races with a united past. They live across Tal'Dorei, most often in isolated communities, though many people in the realm have become used to the sight of goblin children playing catch in the street, hobgoblin mercenaries hanging around job boards, and bugbear apothecaries running shops in the hills.

The origin of the goblinkin is shrouded in myth and legend, but it's generally accepted that these peoples were created by the The Strife Emperor during the \textit{Calamity} to be his perfect soldiers. The few records of that time hint that the goblinkin were once a single folk called the dranassar, a humanoid people of golden skin and green, violet, or yellow eyes who lived in eastern \textit{Wildemount}. They dwelled in great numbers in the city of Ghor Dranas, which came to be the seat of the \textit{Betrayer Gods'} power in the \textit{Calamity}.

For centuries, goblinkin were viewed as monsters by most other people of Tal'Dorei—particularly explorers and adventurers who regularly invaded and destroyed goblinkin settlements. But this ugly sentiment has undergone a radical shift in the past generation, during which countless \textit{goblins}, \textit{hobgoblins}, and \textit{bugbears} have fought alongside the other folk of Tal'Dorei in defense of their homeland against the \textit{Chroma Conclave} and other threats. Today, goblinkin are a welcome sight in Tal'Dorei's cities, and those folk who nurture old hostilities are invariably taken to task by the people around them.

\subsubsection{Goblins}
\textit{Goblins} are a green-skinned people, short of stature and typically having thin, dexterous fingers and strong, nimble legs. They fit easily into settlements with significant gnome or halfling populations, whose residences and furniture are already appropriately sized. With long traditions of living in isolated family units in forests and mountain caves, many city-dwelling \textit{goblins} dwell in large houses holding at least three generations of family and many different branches of their family tree.

\subsubsection{Hobgoblins}
Most \textit{hobgoblins} are alike in stature to \textit{humans} and have similar variations in build, ranging from thin and scrawny to broad shouldered and muscular. Their large ears, gleaming yellow eyes, and powerful incisors give them a vaguely feline appearance.

\textit{Hobgoblins} are often stereotyped as having a warlike, martial nature—a sentiment that arose from rumors surrounding the \textit{Iron Authority} in the \textit{Rifenmist Jungle}. That empire is majority \textbf{hobgoblin}, and its cruel leaders have press-ganged the land's entire population into compulsory military service. In ages past, most \textit{hobgoblins} that became known to the other peoples of Tal'Dorei were scouts or captured legionnaires of the \textit{Iron Authority}. Today, though, many escapees from the authority's influence live in cities such as \textit{Emon}, \textit{Kymal}, or \textit{Byroden}. Most of these free \textit{hobgoblins} know only their parents' stories of that distant empire, and are able to live happily free of its influence.

\subsubsection{Bugbears}
\textit{Bugbears} are tall humanoids covered with brown, black, or blue-gray fur, which combines with large ears and yellow eyes to give them an even more feline appearance than their \textbf{hobgoblin} cousins. Though \textit{bugbears} are tall, their posture is generally stooped, so that their heads are still roughly at human height.

Many \textit{bugbears} are solitary folk, living in forests and mountains where their thick fur protects them against the cold. They typically gather in large groups one to three times a year to socialize and play games with friends, reunite with family, and find love—a tradition that has been reworked by those who live in permanent communities. Such \textit{bugbears} often live with a number of romantic partners under the same roof, or "adopt" \textbf{goblin} families and dwell in their loving, chaotic, multigenerational households.

\subsection{Half-Giants}
Four peoples were known to exist on Tal'Dorei when the world was reborn out of the ashes of the \textit{Calamity}. The \textit{elves}, who emerged from the Fey Realm. The \textit{dwarves}, who emerged from the earth. The \textit{orcs}, who endured the \textit{Calamity} through strength and determination. And the half-giants, who safely rode out that war in mountaintop retreats before eventually descending to the windswept plains below.

Averaging seven feet tall, with hairless, craggy bodies and skin tough as stone, half-giants are often colloquially known as stoneborn or \textit{goliaths}. Descended from \textbf{stone giant}, these folk have long been viewed as exemplars of rugged endurance and individualism. They are found in settlements across Tal'Dorei, where they are known for their innate strength and their commonly held belief—passed down from the giants of old—that all must take fate into their own hands.

In ancient times, \textit{goliaths} and \textit{orcs} lived off the same land, and their nomadic clans frequently warred and intermingled in equal measure. As such, it is common for \textit{orcs} and \textit{goliaths} of Tal'Dorei to discover common ancestors somewhere in their family trees.

\subsubsection{Nomadic Herds}
Stone giants believe that physical perfection is the world's greatest virtue, and this belief has been passed down culturally to the \textit{goliaths} who live in the mountains and upon the \textit{Dividing Plains}. These plains-dwelling \textit{half-giants} call their clans "herds," and unite behind the stone giants' philosophy of "might makes right." Though at one time these herds were mostly made up of stoneborn almost exclusively, recent generations have come to include \textit{humans}, \textit{orcs}, \textit{dwarves}, and even physically diminutive people such as \textit{halflings}—as long as they agree with the philosophy of strength above all.

\subparagraph{Rivermaw}
The greatest of the half-giant herds is the \textit{Rivermaw}, whose members use their strength of arms and political cunning to protect their ancestral lands from incursion. The \textit{Rivermaw herd} includes many \textit{goliaths} from the now-broken Herd of Storms, whose folk suffered at the hands of the tyrant Kevdak many years ago, and who now seek to protect others from ever suffering as they did.

\subparagraph{Ravagers}
Opposing the \textit{Rivermaw} are the Ravagers—a group that the \textit{Rivermaw} have stripped of the title of "herd" in disgust. The Ravagers are more a death cult dedicated to the Betrayer God called the The Ruiner than a true clan. Just like the contemporary Rivermaw, the ranks of the Ravagers are filled with people of many different ancestries, from \textit{goliaths} to \textit{humans} to bloody-minded \textit{firbolgs}—all of them rejecting order to embrace slaughter and cruelty.

\subsubsection{Towers of Stone}
Tal'Dorei's most famous goliath, the legendary hero \textbf{Grog Strongjaw}, is renowned for being a warrior first and a thinker second. But \textit{goliaths} are no less intelligent or canny than Tal'Dorei's other peoples. Despite being known as nomadic folk, many \textit{half-giants} understand that stone is sturdiest when formed into towers, walls, and bridges, and cities such as \textit{Kraghammer} and \textit{Emon} have long been filled with the works of goliath masons and architects not given the historical recognition they deserved.

Great advances in mathematics and civil engineering have been made by other city-dwelling stoneborn, and scholars of the \textit{Cobalt Soul} are working to confirm which of the beloved \textit{Emonian} and \textit{Westruunian} structures that survived the \textit{Chroma Conclave} are actually the result of pioneering goliath minds. \textit{Westruun} has long been home to vast numbers of \textit{half-giants} dwelling there to be close to their former herds, but many others have set down more recent roots after the Battle of \textit{Westruun}, helping rebuild in the wake of the destruction wrought by the dragon Umbrasyl and the Herd of Storms.

\subsection{Halflings}
No one really knows how halflings appeared on Exandria, nor is it known how they first arrived in Tal'Dorei. Some halfling scholars are eager to trace the origin of their people, but many of those who live out in the countryside are more keen on inventing increasingly tall tales regarding the roots of their ancestry. One popular creation myth is that when the The All-Hammer made the \textit{dwarves}, a little bit of leftover clay slipped from his hands and landed on the rolling fields beyond the \textit{dwarves'} mountains, where it became the first halflings. Though most people dismiss this story as nothing more than folklore, a sense of the halflings having humble beginnings fits with the place they've found for themselves in Tal'Dorei.

The historical records of the \textit{Cobalt Soul} assert that halflings journeyed across the Ozmit Sea from \textit{Issylra} with Tal'Dorei's first \textit{humans}, with both peoples coexisting in peaceful harmony in this realm's earliest cities and idyllic agrarian towns. \textit{Humans} and halflings likewise fought together—and against each other—when the Iron Rule of \textit{Drassig} swept across the land.

The \textit{Scattered War} divided halflings across the continent in two ways. Thousands of \textit{humans} and halflings were impressed by \textit{Drassig's} strength and joined his armies, even as many of those who refused to join turned away from their old allegiances. One group of halflings decided that \textit{humans} had betrayed them and that Warren \textit{Drassig} was emblematic of all humanity's evils, and formed the Stoutheart Clan in response. They stockpiled weapons and began a campaign of guerilla warfare against \textit{Drassig's} armies—even battalions that held former friends and family among their ranks.

Another clan, the Lightfoots, asserted that all \textit{humans} couldn't be condemned for the evil actions of one king, and pointed to the thousands of halflings and other non-humans who had allied themselves with \textit{Drassig's} evil regime as proof. The Lightfoots advocated protecting as many people as possible while waiting the war out. In the end, it became clear to both halfling clans that in a conflict against a tyrant such as \textit{Drassig}, neutrality was not an option. They banded together behind \textit{Zan Tal'Dorei's} banner of rebellion.

These two clans dissolved after the \textit{Scattered War}, but many contemporary halfling families can trace their genealogies back to one or the other, and proudly proclaim that their great-grandparents were there when history was being made. Most halflings treat this sense of history only as idle fun, though, and few feel any real animosity over old ills and grudges.

\subsubsection{Lightfoot Clan}
The Lightfoot Clan dissolved into its many component families after the war, but the spirit that created the clan still survives. Those \textit{halflings} who strongly identify with their Lightfoot identity look out for number one first, and are more likely to try to evade conflict than to meet an oncoming storm head-on. Yet a contradiction lingers within the Lightfoot spirit, for despite all their love of the comforts of home, many Lightfoot \textit{halflings} feel a certain unquenchable hunger to explore. Many adventuring \textit{halflings} are Lightfoots who have given in to this compulsion.

\subsubsection{Stoutheart Clan}
The Stoutheart Clan also dissolved in the wake of the \textit{Scattered War}, but its ideals and identity still endure. Brashness and courage even in the face of absurd odds embodies the Stoutheart spirit. Recalling the guerilla army that stood against the legions of \textit{Drassig}, any halfling with the fortitude to stand up to a warrior three times their size has more than earned the honor of calling themself a Stoutheart—even if their family tree doesn't have true roots in that historic clan.

Stouthearts often surprise other people who know of \textit{halflings} only from experiences with more laidback Lightfoot folk—and many use this surprise to their advantage to become truly cunning diplomats. The proud Stoutheart named Assum Emring, a former member of the \textit{Tal'Dorei Council}, once boisterously claimed that he was so rude, so brash and so defiant in his first meeting with the late Sovereign Uriel Tal'Dorei II, that he was very nearly arrested—before being offered a position on the sovereign's staff the very next day!

\subsection{Humans}
Humans are the most populous race in Tal'Dorei today, though at one point in this land's history, all of humankind's holdings were naught but a speck on the western coastline.

If \textit{humans} ever lived upon Tal'Dorei in the \textit{Age of Arcanum}, they were wiped out entirely by the end of the \textit{Calamity}. But after the world had begun to recover from the devastation wrought by that war of gods and mortals, an expedition of human explorers departed from Vasselheim in \textit{Issylra}, thought to be the only bastion of humanoid culture on the surface of Exandria to survive the \textit{Calamity}. Sailing eastward, they reached the land that would eventually come to be called Tal'Dorei, and were shocked to find \textit{elves}, \textit{dwarves}, \textit{orcs}, and \textit{half-giants} already living there—shattering their belief that the \textit{Calamity} had washed away all mortal life except the people of Vasselheim.

Over the years, \textit{humans} and their ambition have been responsible for bloody wars, weapons of unimaginable devastation, and spells that have reduced entire cultures to ruin. Yet \textit{humans} have also driven great advances in art, music, poetry, and theater, and have shown love and compassion to the other mortal peoples of the world. The course of human history is never straight, and can turn at any time. After all, \textit{humans'} lives are short and their minds are chaotic—a fact that troubles many of the long-lived \textit{elves} and \textit{dwarves} of Tal'Dorei. Still, despite these weaknesses—or perhaps because of them—\textit{humans} have produced some of the most innovative and forward-thinking figures Tal'Dorei has ever known.

\subsubsection{A Far-Reaching People}
\textit{Humans} can hail from anywhere in Tal'Dorei, for although their presence on the continent began as a small colony in the west, that presence quickly spread. \textit{Humans} were eager to trade, to forge alliances, and to find love with the \textit{elves}, \textit{dwarves}, \textit{orcs}, and \textit{goliaths} of Tal'Dorei. Their lands quickly grew in size and power. But then a human despot by the name of Warren \textit{Drassig} broke the peace of trade and coalition in favor of conquest.

\textit{Drassig's} betrayal nearly condemned humankind to an eternal reputation for villainy. However, one human rebel named \textit{Zan Tal'Dorei} forged alliances of renewed trust between her rebellious faction and the other beleaguered peoples of \textit{Drassig's} kingdom. Her forthright idealism and steadfast determination won the \textit{Scattered War}, and laid the foundations for a new society that enshrined the alliances she created.

Though human adventurers can hail from anywhere in Tal'Dorei, \textit{humans} are most commonly found in the lands connected by the \textit{Silvercut Roadway}: the \textit{Bladeshimmer Shoreline} in the west, the \textit{Lucidian Coast} in the east, and the \textit{Dividing Plains} that separate them. A number of \textit{humans} from \textit{Wildemount} also sailed across the \textit{Shearing Channel} long ago and settled in the \textit{Alabaster Sierras}, where they founded the city-state of \textit{Whitestone}.

\textit{Humans} living along the \textit{Silvercut} provide the infrastructure for trade across the continent, and as such are often involved in some form of craft, business, or production. Human-dominated cities in Tal'Dorei are typically highly diverse, because they were founded with the aid of many allies of other ancestries. Today, members of all Exandria's other peoples come together in great cities such as \textit{Emon}, even as \textit{humans} from the distant continents of \textit{Issylra}, \textit{Marquet}, and \textit{Wildemount} travel to exchange goods and ideas with the vast and bountiful lands of Tal'Dorei. Beyond the cities, dozens of humanfounded outposts spread along the fringes of lands better known by Tal'Dorei's older ancestries—the \textit{elves}, \textit{dwarves}, \textit{orcs}, \textit{goliaths}, and even the reclusive \textit{giants}. Outlanders, explorers, and hermits take refuge here from the hectic pace of city life, but still find danger enough in the wilderness to keep them on their toes.

Many \textit{humans} likewise live in settlements founded by non-humans, ranging in size from tiny gnome-majority villages in \textit{Torian Forest} to the ancient and stately cities of \textit{Syngorn} and \textit{Kraghammer}. These \textit{humans} absorb the cultures of their homes, just as people of other ancestries raised in cities such as \textit{Emon} can grow up with a distinctly human outlook on life.

\subsection{Orcs}
Orcs survived on Tal'Dorei during the \textit{Calamity} through sheer determination and strength. Unlike the \textit{elves} and \textit{dwarves} who hid in their secret retreats, and the \textit{goliaths} who watched the destruction from their mountaintops, great clans of orcs fought endless bloody battles on the \textit{Dividing Plains} as the gods and their loyal minions clashed, wielding magic sufficient to rip the world asunder. And the orcs endured—perhaps in part because of the divine power behind that people's creation.

The myth surrounding the orcs' origin is a grim one. It is said that they were created by accident when the The Arch Heart, god of the \textit{elves}, shot out the eye of the war god known as the The Ruiner. The The Ruiner blood poured from his wound and rained upon the armies of \textit{elves} who served the The Arch Heart—and those unsuspecting mortals were transformed. Their muscles bulged, tusks grew from their jaws, and their minds burned with a thirst for blood.

Yet after that battle's end, the fury within these new humanoids burned away, and the orcs survived as a hardy, physically powerful people. In the early days of \textit{Gwessar}, orc clans warred with anyone who invaded their territory, especially the \textit{goliaths} who descended from the mountains to the plains, and who rivaled the orcs in both strength and stature. But once more, fury cooled in time, and reasonable voices won out over those of even the most passionate warmongers.

Today, orcs and other folk with orcish heritage can be found all across Tal'Dorei. Many orcs find their natural athleticism an asset in specific trades, but they can be found in every guild and profession. Though orcs and \textit{humans} have an ancient legacy of conflict that lives on in a prejudiced few, the realm is filled with people who denounce and ostracize any who discriminate against their orc neighbors.

\subsubsection{From Near and Far}
Though many of Tal'Dorei's \textit{orcs} are related to clans that have lived nomadically on the continent since time immemorial, the major cities of the land are also home to \textit{orcs} who've traveled to Tal'Dorei from across the sea. Many of those \textit{orcs} originated in the lands of Xhorhas, in \textit{Wildemount}, and the continent of \textit{Marquet}—places where \textit{orcs} have never been denied the opportunity to become great scholars or leaders.

\subsection{Tieflings}
Though Tal'Dorei was once awash with fell legends claiming to detail the evil origin of \textit{tieflings}, the more progressive parts of the realm have long since stomped out such ugly myths. Still, many folk can't help but show surprise the first time they see a person with horns curling from their brow, gleaming eyes of a single solid color, and lustrous, brightly hued skin of red, purple, or blue.

Scholars have scoured thousands of years of \textit{Issylran} history in search of the origin of \textit{tieflings}, to little avail. The most complete extant records suggest that during the \textit{Age of Arcanum}, a cabal of power-hungry \textit{Issylran} warlocks consorted with extraplanar entities, and that many of these unions resulted in children who shared the qualities of human and fiend. These progeny of the Lower Planes eventually came to break away from that origin to become a people unto themselves.

The first \textit{tieflings} to walk on Tal'Dorei sailed from \textit{Issylra} during the \textit{Age of Arcanum}, fleeing religious zealots who believed their very existence was an abomination. Many \textit{tieflings} in Tal'Dorei refused to fight during the \textit{Calamity}, embracing a nihilism which said that if the world could reject their kind, perhaps it was better to just let it burn. Yet the world and the \textit{tieflings} survived.

In the centuries after the \textit{Divergence}, \textit{tieflings} who had grown weary of gods and demons rejoiced, feeling that the gods' imprisonment had effectively freed them from their cursed past. As such, when Warren \textit{Drassig} and his sons conquered the elven lands of \textit{Gwessar}, countless \textit{tieflings} rose in rebellion against the tyrant. They had seen his kind before, they knew his methods, and they knew others would suffer as they once had.

\subsubsection{Tiefling Children}
\textit{Tieflings} who form a union always have tiefling children, and most \textit{tieflings} belong firmly to that ancient ancestry. Yet to this day, any humanoid—whether famous hero or common farmer—who makes a deal with fiends might give rise to children, grandchildren, or even more distant descendants born with the horns, eyes, and colorful skin of a tiefling.

Perhaps the most famous tiefling born in the past generation is Gwendolyn de Rolo, the youngest daughter of the renowned heroes \textbf{Percival de Rolo} and \textbf{Vex'ahlia} de Rolo of \textit{Whitestone}. It is an open secret in the city-state that \textbf{Percival de Rolo} had dealings with fiends as a young man. Yet though none of the rumors have precisely pinpointed the nature of \textbf{Percival de Rolo} supernatural bargain, all who know Gwendolyn agree that she is the sweetest and most innocent de Rolo to roam the halls of \textit{Castle Whitestone} in many years.

\subsection{Mixed Ancestry}
Even the oldest records of the \textit{Cobalt Soul} are unclear as to whether the ancestries of \textit{elves}, \textit{dwarves}, \textit{orcs}, and \textit{goliaths} routinely mixed before the arrival of humankind on Tal'Dorei's shores. But it's easy to assume that explorers and ambassadors would inevitably find love or companionship among the folk of unfamiliar lands—and it is well known that the arrival of \textit{humans} changed Tal'Dorei forever.

From the earliest days of human exploration on the continent, love and passion began to blossom between all of Tal'Dorei's people. The children that spring from the unions of folk of different heritage bear the traits of each of their parents, so that a human child might gain patches of scales and reptilian eyes from a \textit{dragonblood} parent, or an orc and a dwarf might raise a child of squat, stout stature and olive-green skin, and bearing wisps of a coppery beard even as an infant.

In the modern age, love unites all folk. People of mixed ancestry can be found all across Tal'Dorei, though they are most common in the realm's cosmopolitan cities. Some ancestries are biologically improbable (such as a goliath and a gnome having a child), but sufficient application of magic makes even those unions possible. However, money can be a limiting factor in such things, as having a skilled mage on hand during conception and childbirth is an expensive process.

Non-humanoid creatures with the magical ability to assume humanoid size and form, such as dragons, can also sire or give birth to mixed-ancestry children.

\subsection{Other Races}
In addition to humanoids, countless other sapient beings walk the lands of Tal'Dorei. Many of these are rare creatures, never building communities of folk who share their appearance and abilities. Some are wholly unique creations of unpredictable magical effects. Others are mighty beings such as dragons and giants, who can take on the role of hero or villain with equal ease but who intentionally separate themselves from the affairs of Tal'Dorei's other mortal peoples.

\section{Hemocraft}
\section{Subclasses}
This section features new subclass options for some of the familiar fifth edition character classes. These subclasses grant you new options not available to all members of your class. Information on your full class features, as well as other subclasses, can be found in the fifth edition rules.

\subsection{Barbarian}
You are a wild warrior whose emotions and instincts guide you in combat far more than logic or tactics. When your barbarian character chooses a primal path at 3rd level, the barbarian is a new option for use in a Tal'Dorei campaign.

\begin{DndSidebar}[float=!b]{Kel'jaia Uloeh}
\textit{Female gnome Path of the Juggernaut barbarian}



When Kel'jaia sets her eyes on something, nothing can stop her from achieving it. Her small size has no bearing on her unbelievable strength, which can overpower creatures many times her size. She came to Tal'Dorei after the village on the \textit{Rifenmist Peninsula} she was defending was razed by the \textit{Iron Authority}. She first arrived in \textit{Syngorn}, and through the connections and friendships she made there, was chosen to be a \textit{Syngornian} delegate to the \textit{Tal'Dorei Council}. She strongly believes that Tal'Dorei should take a stand against the cruelty of the Authority. Should strife come to Tal'Dorei, Kel'jaia is thought by all to be the council's first choice for the Master of War, due to her valor, tactical cunning, and battle prowess.





\end{DndSidebar}

\begin{DndSidebar}[float=!b]{Rules Tip: Forced Movement}
Usually when one creature moves out of a hostile creature's reach, the hostile creature can use its reaction to make an opportunity attack. However, forced movement—such as being pushed by a Path of the Juggernaut barbarian's Thunderous Blows feature—doesn't provoke opportunity attack.



Likewise, a juggernaut barbarian's Hurricane Strike feature allows an ally to make a melee weapon attack as a reaction only if the foe ends its forced movement within 5 feet of the ally. If a foe is pushed through other spaces within 5 feet of your allies, those allies can't make normal opportunity attack against the foe.



\end{DndSidebar}

\subsection{Bard}
You are a charming performer whose wit and words can inspire heroism in your allies and a rising sense of dread in your foes. When your bard character chooses a bard college at 3rd level, the bard is a new option available in a Tal'Dorei campaign.

\begin{DndSidebar}[float=!b]{Balthasar Bleakskull}
\textit{Male half-giant College of Tragedy bard}



Born in \textit{Westruun} and raised by two loving goliath parents—a towering warrior and an equally massive astronomer of the \textit{Yuminor Observatory}—Balthasar Bleakskull was exposed to the adventuring lifestyle and academia from an early age. He tumbled into the art of theater when he was but a child, and followed his passion for drama into adulthood—when he was entangled in a bloodcurdling academic conspiracy that left eight dead, including his dramatic mentor and his lover. Since then, he has devoted his dramatic skill to the art of tragedy, to grant his audiences the catharsis that he one day hopes to give himself.





\end{DndSidebar}

\subsection{Cleric}
You are a wise disciple of one or more of Exandria's gods, whose divine hands can no longer shape the world on their own. When your cleric character chooses a divine domain at 1st level, you can consult your god's entry in chapter 2 for suggestions on which domains they might grant. The cleric and cleric are new options available in a Tal'Dorei campaign.

\begin{DndSidebar}[float=!b]{Alasterre de Vitrevos}
\textit{Male human Blood Domain cleric}



The director of Tal'Dorei's \textit{Claret Order} is a meticulous, orderly man. He believes that he deserves all the good that comes to him in the world—and, in fairness, all the bad as well. He wields the unsettling magic of hemocraft created by his order in service of both his organization's interests, and the divine interests of the The Matron of Ravens. His goals are simple: purge the world of undead, particularly \textbf{vampire}, and make his name one sung across Exandria in the process.





\end{DndSidebar}

\begin{DndSidebar}[float=!b]{Nalys Ildareth}
\textit{Nonbinary (they/them) elf Moon Domain cleric}



The moons of Exandria show their faces for all to see, yet hide a secret face from the world. Nalys Ildareth, leader of an obscure \textit{Syngornian} sect of the The Moonweaver, believes that true reverence requires all faithful to have a face they show to the world—represented by a mask of silver they never remove—and one they keep to themselves, beneath the mask. Nalys believes that it is their duty to enact the The Moonweaver will upon Exandria, and keep the world safe within the balance of Catha's calm order and the ill-fated chaos of Ruidus. They are thoughtful and calculating, seeming an impassive observer while they wait for the proper moment to strike, but then transforming into a furious avenger when that moment arrives.





\end{DndSidebar}

\subsection{Druid}
As a druid, you are a keeper of the old ways, attuned to nature and channeling power that grants you spells and allows you to transform your body. When your druid character chooses a druid circle at 2nd level, the Druid is a new option available in a Tal'Dorei campaign.

\begin{DndSidebar}[float=!b]{Camellia Springshower}
\textit{Female firbolg Circle of the Blighted druid}



Life was beautiful under the light of the sun for Camellia Springshower. The sun caused her flowers to flourish, and the rains nourished them so they could grow when next the sun arose. One day, however, she wandered too far from her family while they traveled the \textit{Dividing Plains}, and tumbled into the bleak depths of \textit{Shadebarrow}, a forgotten, accursed tomb. She ran for the light, but a specter reached out from the walls and touched a ghastly finger to her chest. She felt her blood run cold, and the flowers within her basket wither away. Though she escaped the \textit{Shadebarrow}, she was never able to find her family again, and she now wanders the \textit{Dividing Plains} like an untethered shadow—hoping to use the blighted powers granted to her by the shade to do some semblance of good in the world, even if all life seems to wither beneath her touch.





\end{DndSidebar}


\begin{DndMonster}[float=t]{Blighted Sapling}
\DndMonsterType{Medium plant, U}
\DndMonsterBasics[
armor-class = {10 + PB (natural armor)},
hit-points = {twice your druid level},
speed = {30 ft.}
]
\DndMonsterAbilityScores[
 str = 8,
 dex = 13,
 con = 12,
 int = 4,
 wis = 8,
 cha = 3
]
\DndMonsterDetails[
damage-vulnerabilities = {fire},
damage-resistances = {necrotic, poison},
senses = {blindsight 60 ft. (blind beyond this radius), passive perception 9},
languages = {understands the languages you speak},
]
\DndMonsterAction{Blighted Resilience (10th level or higher)}
The creature has immunity to necrotic and poison damage and to the poisoned condition.

\DndMonsterAction{Toxic Demise (10th level or higher)}
When the creature is reduced to 0 hit points, it explodes in a burst of toxic mulch or fetid viscera. Each creature within 5 feet of the exploding creature must succeed on a Constitution saving throws against your spell save DC or take necrotic damage based on the creature's challenge rating (see the table below). As an action, you can also cause a summoned creature to explode, immediately killing it.

\begin{DndTable}[header=Toxic Demise Damage]{cX}
Creature CR & Damage\\
1/4 or lower & 1d4 necrotic damage\\
1/2 & 1d6 necrotic damage\\
1 or higher &  CR \texttimes{} d8 necrotic damage\\
No challenge rating & Prof. bonus \texttimes{} d6 necrotic damage\\
\end{DndTable}

\par
\DndMonsterSection{Actions}
\DndMonsterAction{Multiattack}
When you reach 14th level in this class, the blighted sapling makes two claw attacks.

\DndMonsterAction{Claws}
\textsl{Melee Weapon Attack:}  to hit, reach 5 ft., one target. \textsl{Hit:}  2d4 + PB piercing damage.

\end{DndMonster}


\subsection{Monk}
You are a canny martial artist whose connection to the unseen energy of life grants you supernatural insights and the ability to manipulate your foes. When your monk character chooses a monastic tradition at 3rd level, the monk is a new option available in a Tal'Dorei campaign.

\begin{DndSidebar}[float=!b]{Cressida Holt}
\textit{Nonbinary (she/they) orc Way of the Cobalt Soul monk}



A freshly titled expositor of the \textit{Cobalt Soul}, Cressida Holt has already made a name for themself as a suave, stylish secret agent. Their assignments as an expositor take them across Tal'Dorei, though they are most often seen in \textit{Emon}, \textit{Kymal}, and \textit{Westruun}, where they infiltrate political galas, high-stakes games of two-card, and other dens of high-society corruption.



Though most \textit{Cobalt Soul} expositors are expected to keep a low profile and gather intelligence, Cressida has found success with a flashier method. She strides into a room in a fine cobalt-blue suit, makes sensual eye contact with the most beautiful person in the room, and casually makes herself someone from whom no one can keep a secret.





\end{DndSidebar}

\subsection{Paladin}
You are an honorable warrior whose devotion to a cause or ideal empowers you with righteous magic. When your paladin character chooses a sacred oath at 3rd level, the paladin is a new option available in a Tal'Dorei campaign.

\begin{DndSidebar}[float=!b]{Shirkuh Marín}
\textit{Male human Oath of the Open Sea paladin}



Though he is only in his early 30s, famed sailor Shirkuh Marín is well on the way to having his life become an epic sung in taverns all over the Lucidean Ocean. Born a girl to Marquesian immigrant parents in Nicodranas on \textit{Wildemount's} Menagerie Coast, Shirkuh took his present name from his great-granduncle, an infamous Marquesian pirate. He swore to wear that name with pride, and has followed his own code of honor and ethics ever since. He has worked hand-in-hand with both honest sailors and pirates of the Revelry—and clashed swords with both, too.



Of the many tales that surround Shirkuh, it's unclear which are true—if he truly outran a \textbf{juvenile kraken} in the \textit{Shearing Channel}, if he indeed cheated the Plank King in a game of two-card and lived to tell the tale, and if he actually managed to steal the heart of a \textit{Clasp} Spireling in \textit{Kymal}. But most agree that even if the stories are made up, he possesses both the talent and the recklessness to have accomplished them if he set his mind to it.





\end{DndSidebar}

\subsubsection{Oath of the Open Sea Spells}
Paladins of the paladin automatically gain two new spells through their Oath Spells. These spells can also be learned by druids, rangers, and sorcerers of the appropriate level, at the GM's discretion.

This book adds the following new spells:


\begin{itemize}
\begin{multicols}{2}
\item \textit{Freedom of the Waves}
\item \textit{Freedom of the Winds}
\end{multicols}

\end{itemize}

\begin{DndSidebar}[float=!b]{Rules Tip: Visibility}
Fog and other effects can obscure vision for you, your enemies, and your allies. When you heavily obscure an area using your Channel Divinity option, all creatures within the area have their vision completely blocked, and creatures outside the area can't see in. Creatures that can't see automatically fail ability checks that require sight. Also, attack rolls against creatures that can't see have Advantage and Disadvantage, and their own attack rolls have Advantage and Disadvantage.



Creatures in a lightly obscured area have Advantage and Disadvantage only on Wisdom (Perception) checks that rely on sight. The rules for when your vision is obscured are described completely in the fifth edition core rules.



\end{DndSidebar}

\subsection{Sorcerer}
You are a spellcaster to whom magic comes as naturally as breathing, possessing an instinctive grasp of its power—for that power flows through your very veins. When your sorcerer character chooses a sorcerous origin at 1st level, the sorcerer is a new option available in a Tal'Dorei campaign.

\begin{DndSidebar}[float=!b]{Karzadrawn "Karzy" Sunscale}
\textit{Male brass \textit{dragonblood} Runechild sorcerer}



When you're a model in a city as big as \textit{Emon}, you've got to find a way to stand out. Being a broad-shouldered, heavyset \textit{dragonblood} with brass scales that twinkle in the sun is a good start, but Karzy Sunscale has another ace up his sleeve: he's a runechild. Ever since he was an infant, he has manifested strange magical abilities, and when he expresses intense emotions through his magic, mysterious runes flare to life upon his scales. All of these traits have always made Karzy the center of attention, and he basks in other peoples' interest like a lizard in the sun. He goes out of his way to dress in outrageous fashions and emote as vividly as he can. He works as a fashion model in \textit{Emon}, but the idea has worked its way into his head that he could make an even bigger splash if he were to also become a celebrity adventurer by completing a truly monumental quest.





\end{DndSidebar}

\subsection{Wizard}
You are a studious spellcaster who has earned your influence over arcane power through long hours of rigorous learning. When your wizard character chooses an arcane tradition at 2nd level, wizard is a new option available in a Tal'Dorei campaign.

\begin{DndSidebar}[float=!b]{Honor Kinnabari}
\textit{Female tiefling Blood Magic wizard}



It's not easy being the youngest spellwright in the \textit{League of Miracles}, especially when you've got a conscience. Honor Kinnabari was enamored with the league since she was a child in \textit{Westruun}, when she saw the spellwrights and their \textbf{adranach} constructs rebuild the entire Market Ward in a matter of days after it was reduced to a slag pit by a hoard of \textbf{magma landshark}. She was raised by dwarf parents from the \textit{Alabaster Sierras}, and wears wizardly robes in the style of their clan in honor of them. Honor has just constructed an \textbf{adranach} of her own, and has begun leading small-scale operations for the league, but she's starting to worry that something cultish and sinister lurks beneath her compatriots' practiced smiles.





\end{DndSidebar}

\section{Backgrounds}
Every player character in a fifth edition fantasy roleplaying game has a background. While your class represents your role in an adventuring party, your background represents who you were and what you learned before you began your adventuring career. Your background doesn't necessarily have to be confined to the past, though, as the way it marks your connections to the world might involve ongoing commitments to an organization or even an entire people.

The new backgrounds in this section each grant your character new proficiencies, languages, and equipment, as well as a unique background feature. Tables of characteristics can also help you quickly generate personality traits, ideals, character flaws, and bonds to tie your character to the people, places, and causes of a Tal'Dorei campaign.

This book adds the following new backgrounds:


\begin{itemize}
\begin{multicols}{3}
\item Ashari
\item Clasp Member
\item Variant Clasp Member (Myriad Operative)
\item Lyceum Scholar
\item Reformed Cultist
\item Whitestone Rifle Corps
\item Variant Whitestone Rifle Corps (Grey Hunter)
\end{multicols}

\end{itemize}

\section{Supernatural Blessing: Fate-Touched}
Most lives and souls find their ways easily along the paths of fate and destiny, oblivious and eager to be a part of the history meant for them. Some souls, however, carry themselves with mysterious influence, the golden threads of their fate subtly guiding history as they pass. Events both magnificent and catastrophic often follow such souls, which take the form of great rulers and terrible tyrants, powerful mages, religious leaders, and legendary heroes. These are the fate-touched.

Few fate-touched ever become aware of their prominent nature. Without the confirmation of the gods themselves, few can know for certain that they bear this blessing. And this is as it should be, for people of renown and power grow easily jealous of those whose gifts outstrip their own. The surviving chronicles of the \textit{Age of Arcanum} are filled with tales of magical hierarchs who enslaved those they took for fate-touched—or worse, attempted to extract the essence of fate from them by the most horrid means. Being touched by fate can be a blessing or a curse. It is most often those who bear the blessings of fortune who are remembered by history, while those cursed by fate are forgotten. Most cultures of Exandria have legends of those marked by ill fate simply by being born under the full light of Ruidus, the vermilion moon around which countless fell legends swirl.

Player characters and NPCs gain this blessing by having the GM grant it to them, gaining the Fortune's Grace feature as a result. A character might receive the blessing as a gift from the gods, or as a temporary and surprising phenomenon when the campaign shifts into a story arc focused on that character. An NPC ally might bear this blessing to help them survive as the characters protect them from a terrible threat. A previously benign NPC might even become a power-hungry villain as they test the limits of their fate-touched gift.

\section{New Feats}
The following feats are appropriate for a Tal'Dorei campaign.


\begin{itemize}
\begin{multicols}{3}
\item Cruel
\item Flash Recall
\item Mystic Conflux
\item Remarkable Recovery
\item Spelldriver
\item Thrown Arms Master
\item Vital Sacrifice
\end{multicols}

\end{itemize}

\chapter{Game Master's Toolkit}
\DndDropCapLine{G}{}ame Masters can make use of many different tools to run enjoyable games. A GM is the starting point of the storytelling process, coming up with narratively compelling adventure frameworks that they and the players then turn into a campaign. To do so, a GM works up locations and scenarios that their players will have fun with. They stock adventures with loot and magic items to reward daring characters, villains and monsters to defend those valuable treasures, and NPCs of all kinds who can give clues, help as allies, or oppose the characters in their goals.


This chapter contains storytelling advice for GMs running campaigns in Tal'Dorei. It also contains a number of new magic items, including the legendary \textit{Vestiges of Divergence}, plus optional rules that GMs can use to fine-tune the tone of a Tal'Dorei campaign. (\textit{Chapter 6} of this book presents monsters and NPCs unique to the world of Exandria, which can likewise help give a campaign a distinctly \textit{Critical Role} feeling.)

\section{Creating Adventures}
Before reading this section, you should familiarize yourself with "\textit{Running a Tal'Dorei Campaign}" in \textit{chapter 1}, which provides information about Tal'Dorei's lore and conflicts that are for the Game Master's eyes only. \textit{Chapter 3} of this book also contains dozens of specific story hooks, many of which tie into the big-picture secrets introduced in \textit{chapter 1}. Some hooks are simple, giving you just enough to inspire your next game session. Some are more complex, and can cascade into an entire campaign's worth of story.

You'll notice that these story hooks present a beginning, and sometimes a hint of a middle, but never a conclusion. That's because they're designed to spark your imagination, but in a way that still leaves you to fill out the rest of the story framework in your own way. The goal of this book is to provide the building blocks with which you can create your own Tal'Dorei, and these open-ended plot seeds will help you set up adventures that feel right at home in the world of \textit{Critical Role}—whether you're a longtime fan of the show or are brand-new to Tal'Dorei.

If you're an experienced GM, you might already have your own approach to creating adventures from story hooks. But especially if you're a beginning GM, this section presents a straightforward approach you can try out for expanding hooks into full adventures. And even as an experienced GM, if you ever feel stuck while planning an adventure, it can help to create a "fill-in-the-blanks" outline of your narrative. Start by dividing a page of your GM journal (whether it's a physical page or a digital one) into three sections labeled "Beginning," "Middle," and "End." Then think about what happens in each of those parts of your adventure as a means of opening up your creative process.

\subsection{Beginning an Adventure}
The beginning of most adventures comes when the characters receive a quest or are plunged headfirst into an unexpected scenario. In the Beginning section of your GM journal, write one sentence about how you want the characters to learn of the adventure. For example: "A member of the \textit{Tal'Dorei Council} posts an urgent notice about a dangerous creature in every tavern in \textit{Emon}, promising a 1,000 gp reward." Then write down little notes about other things you want in the beginning scene. Where is this opening scene set? What other NPCs are there? Are there environmental notes or clues you want to communicate to the players?

\subsection{Creating Conflict in the Middle}
The middle of an adventure is filled with conflict, whether episodic events during travel, a dungeon exploration filled with deadly traps and territorial monsters, or intense diplomatic negotiations. In the Middle section, write down notes about what sort of conflict you want your adventure to have. Most story hooks in this book contain a clear antagonist or threat—so start by writing down that antagonist and then brainstorming secondary threats that will challenge the characters on their way to facing their main foe. In addition to the monsters in \textit{chapter 6} of this book, the fifth edition rules are filled with useful monsters that are perfectly at home in Exandria.

\subsection{End with a Reward}
The end of an adventure is where the characters gain their reward and return to safety. In the Ending section of your journal, write down ideas for treasure or other rewards that you want to give the characters. You can make use of the new treasures featured in this book, as well as the varied rewards described in the fifth edition rules.

Remember that treasure doesn't have to be gold or magic items, and it doesn't have to be found in a dungeon. When the characters return to town carrying the head of the monster they've slain, their quest-giver might grant them an audience with the \textit{Tal'Dorei Council}, give them special use of the \textit{Traverse Junction} or some other magical marvel of the \textit{Alabaster Lyceum}, or bestow any number of other immaterial rewards.

\subsection{Going with the Flow}
Don't worry if your plans go awry—they will. No RPG session ever goes totally as planned. The most important skill for any GM is the ability to improvise, because you can never predict everything your players will do. Make your plans loose and variable, and be willing to have fun and change course. Your players are surprised by every challenge you throw at them—so give yourself a chance for the same experience! It can be fun to be surprised, especially when you're the GM who supposedly holds all the important story cards.

\section{Tal'Dorei Treasures}
This section presents magic items for use in your Tal'Dorei campaign, including specific items used by \textit{the Ashari} people of Exandria, and the legendary \textit{Vestiges of Divergence}.

\subsection{Magic Items}
Tal'Dorei is a land awash with magic. Numerous magical organizations, both stately and shadowy, create enchanted items to aid their members in mysterious quests. Shops in major cities, including \textit{Gilmore's Glorious Goods} in \textit{Emon}, openly sell magic items to anyone with the gold to pay their extravagant (some might say exorbitant) fees.

GMs can use the magic items in this section—some of which have appeared in \textit{Critical Role} before; some of which are all new—as treasure to reward the characters during their adventures. Unless otherwise stated, these items abide by the magic item rules laid down in the fifth edition core rules.


\begin{itemize}
\begin{multicols}{3}
\item \textit{Boots of Haste}
\item \textit{Boots of the Vigilant}
\item \textit{Cataclysm Bolts}
\item \textit{Coat of the Crest}
\item \textit{Corecut Dagger}
\item \textit{Dagger of Denial}
\item \textit{Doublet of Dramatic Demise}
\item \textit{Echo Stone}
\item \textit{Graz'tchar, the Decadent End}
\item \textit{Inescapable Lash}
\item \textit{Magician's Judge}
\item \textit{Mirror of Infinite Transpondence}
\item \textit{Raven's Slumber}
\item \textit{Rod of Mercurial Form}
\item \textit{Stormrider Boots}
\item \textit{Summer's Dance}
\item \textit{Tinkertop Boltblaster 1000}
\end{multicols}

\end{itemize}


\begin{itemize}
\begin{multicols}{3}
\item \textit{+2 Ring of Protection}
\item \textit{+2 Studded Leather Armor of Acid Resistance}
\item \textit{+2 Studded White Dragon Leather Armor of Cold Resistance}
\item \textit{Azuremite}
\item \textit{Circlet of Wisdom}
\item \textit{Dynamite}
\item \textit{Gluebomb}
\item \textit{Oloore Root Teabag}
\item \textit{Plainscow}
\item \textit{Residuum}
\item \textit{Residuum Enchanting Slate}
\item \textit{Skyship}
\item \textit{Stink Bomb}
\item \textit{Sunbeam Compass}
\item \textit{Suude}
\item \textit{Suude (Blue)}
\item \textit{Suude (Brown)}
\item \textit{Suude (Red)}
\item \textit{Whitestone}
\item \textit{Writing Kit}
\item \textit{Zeal}
\end{multicols}

\end{itemize}

\subsection{Tools of the Ashari}
\textit{The Ashari} people have long used their connection and attunement to elemental power to create unique magic items, helping them in their quest to further understand and control the elements. These items are kept secret by their creators and are rarely given to people not of \textit{the Ashari}. As such, they can be acquired only as a gift from an Ashari for rendering their people a great service, from the body of a dead Ashari explorer, or for an exorbitant price from someone who has acquired the item through dubious means.

At the GM's discretion, the unique nature of these items might mean that a character not of \textit{the Ashari} and unfamiliar with elemental magic can't automatically attune to them. Such characters can attune to one of the tools of \textit{the Ashari} with 24 hours worth of special training from an Ashari familiar with the item's use.


\begin{itemize}
\begin{multicols}{2}
\item \textit{Earthboard}
\item \textit{Flamefriend Lantern}
\item \textit{Oceanic Weapon}
\item \textit{Skysail}
\end{multicols}

\end{itemize}

\subsection{Vestiges of Divergence}
During the wars of the \textit{Calamity}, destruction rained across the world. Fragmented chronicles of this time recount battles between gods that leveled mountains and sank whole civilizations beneath the waves. Chosen heroes rose up to accept weapons, armor, and other magic items blessed by their patron gods to staunch the inexorable tide of death. Powerful archmages wove the overflowing arcane powers of the war into legendary magic items that could both thwart death and mete it out in equal measure.

These relics of the \textit{Age of Arcanum} armed mortals with god-like powers, and escalated the \textit{Calamity} to a fever pitch. Soon after, the war climaxed in a wave of total destruction that left much of Exandria in ruins. Thousands of years of history were reduced to ash—and in that ash were buried many of these selfsame relics. Some of those not lost were passed down through bloodlines for generations, becoming potent symbols of authority and rulership. Others were sealed away by those who feared their power. These mighty echoes of the war that nearly shattered Exandria and the divine banishment that followed are known as the Vestiges of Divergence.

The Vestiges of Divergence are unique even among legendary magic items, for each one touches the defiant soul of its wielder and grows in power alongside them. The base state of a Vestige of Divergence is its dormant state, which the item returns to if sealed away or otherwise kept from being attuned by a worthy individual for a long period of time. A dormant state still imbues a modicum of power to an attuned wielder. But with time, perseverance, and personal growth, the Vestige of Divergence can regain access to some of its lost power in its awakened state. Eventually, if its bearer overcomes personal hardship and transforms their own soul, a Vestige of Divergence reaches its full potential in its exalted state. This shows the true power of the item, as it was when first utilized in the \textit{Calamity} against gods armed for war.

The point at which a character finds a Vestige of Divergence, as well as the point at which that item starts to unlock its full power, is entirely in the hands of the Game Master—though it is a topic that players and GMs can discuss and mutually scheme over between game sessions. The acquisition of such an item—or the discovery of its true identity—has the potential to be a scene of momentous narrative weight. Those who bear a Vestige of Divergence are respected and feared by those who understand these items' true power. But such an honor can also paint a target on the bearer's back, for many covet such power, and would do anything to obtain it.

Many of the Vestiges of Divergence detailed in this section can currently be found in Tal'Dorei, while others are scattered elsewhere throughout Exandria. Some of these relics remain lost to time, waiting for someone to find them among the dust and rubble. Others might be in use by powerful heroes ready to pass the item's power onto a worthy successor who fights for the proper cause. Still others might be in the hands of twisted souls using them to terrorize the weak and defenseless. The Vestiges of Divergence are both magic items and narrative devices, and it's the GM's choice as to where any one of them might be found or earned.

\subsubsection{Advancement of a Vestige of Divergence}
The conditions by which a \textit{Vestige of Divergence} progresses to its next level of power can vary, and often revolve around the nature of the specific item and its magic. An attuned bearer of a \textit{Vestige of Divergence} must often symbolically progress to a new state of self-discovery or accomplishment. Alternatively, a \textit{Vestige of Divergence} might bequeath new power to a character in a state of extreme personal duress or desperation.

Such triggering moments are entirely up to the Game Master to identify and unveil, and as such, will likely manifest in unexpected ways. Players should allow the organic narrative moments that they feel might exemplify such an advancement to occur, and GMs should tailor the triggering moment accordingly to be important, impactful, and memorable.

Examples of the kinds of moments that can trigger the advancement of a \textit{Vestige of Divergence} include the following:


\begin{itemize}
\item A character finally surmounts one of their greatest fears, bravely overcoming a traumatic event to save a party member.
\item A character is mercilessly defeated by a long-hated foe, and in the face of such defeat, they feel a deep, dormant strength grow from within.
\item A character loses a close ally in battle, and the anguish and fury stirs the power within their \textit{Vestige of Divergence}.
\item A character discovers a facet of their destiny that calls their actions toward a dangerous cause—and against their own fear, they accept their fate and responsibility.
\item A character claims vengeance against a rival who has tormented them for ages.
\item A character known for restraint gives in to the amoral violence that the relic they bear was dedicated to.
\end{itemize}

If you want to keep the Vestiges "balanced," in addition to being narratively significant, then characters may find a dormant Vestige before 11th level—perhaps unaware that it's even a Vestige! That dormant Vestige is awakened somewhere between 11th and 17th level, and is exalted at any time after that. Even using these level ranges as a rule of thumb, you shouldn't rely too heavily on those mechanical benchmarks. A \textit{Vestige of Divergence's} advancement is triggered by character development, not simply by leveling up.

\subsubsection{To Suit the Wielder}
The \textit{Vestiges of Divergence} were all created by humanoids or gifted by the gods to their humanoid champions. As such, all of these legendary magic items are sized to suit Medium creatures. However, the mighty magic that thrums within each item allows it to grow or shrink in size to suit the creature attuned to it.

The game statistics of a \textit{Vestige of Divergence} don't change when it resizes unless the item is a weapon. When a Large or larger creature makes an attack with such a weapon, the weapon deals one additional damage die for each size category its wielder is larger than Medium. For example, if \textit{Fenthras} is used by a Large fey, attacks made using the longbow deal 2d8 piercing damage. If \textit{Pyremaul} is used by a Huge giant, attacks made using the maul deal 4d6 bludgeoning damage.

\begin{DndReadAloud}{}
\DndQuote%
{}
{''When the god-war spread across Exandria, no creature was safe, no matter how small and insignificant, no matter how mighty and powerful. For the battles between gods, would-be gods, and their divine monstrosities took place on such a scale that mountains were sundered and whole cities perished. How could any mortal hope to sway the fate of their own world?''}
{''The answer—perhaps ironically, given how the arcane arts and a hunger for godhood brought about the Calamity—lay in magic of both arcane and divine provenance. Not paltry charms or trinkets, but the kind of magic that suffuses the ley lines of this world and plucks at the very fabric of reality. True, a mortal army may be naught but a cloud of gnats to the gods, but a singular hero—or villain—wielding an artifact imbued with magic beyond reckoning...such a champion could bring their strength and cunning to bear amidst the devastation of the god-war.''}
{''Thus, the gods and even the greatest mortal mages wrought relics of a might never seen before, nor, I think, since. Moreover, the full potential of these items could only used by those with souls worthy of wielding such legendary power....''}
{Monograph on the Vestiges of Divergence from the Cobalt Reserve archives circa 547 P.D., attrib. High Curator Jorum Irrelios}
\end{DndReadAloud}

\begin{DndTable}[header=Vestiges of Divergence by Advancement]{XXXX}

Item Group & Dormant & Awakened & Exalted\\
\textit{Agony} & \textit{Agony (Dormant)} & \textit{Agony (Awakened)} & \textit{Agony (Exalted)}\\
\textit{Armor of the Valiant Soul} & \textit{Armor of the Valiant Soul (Dormant)} & \textit{Armor of the Valiant Soul (Awakened)} & \textit{Armor of the Valiant Soul (Exalted)}\\
\textit{Cabal's Ruin} & \textit{Cabal's Ruin (Dormant)} & \textit{Cabal's Ruin (Awakened)} & \textit{Cabal's Ruin (Exalted)}\\
\textit{Circlet of Barbed Vision} & \textit{Circlet of Barbed Vision (Dormant)} & \textit{Circlet of Barbed Vision (Awakened)} & \textit{Circlet of Barbed Vision (Exalted)}\\
\textit{Condemner} & \textit{Condemner (Dormant)} & \textit{Condemner (Awakened)} & \textit{Condemner (Exalted)}\\
\textit{Deathwalker's Ward} & \textit{Deathwalker's Ward (Dormant)} & \textit{Deathwalker's Ward (Awakened)} & \textit{Deathwalker's Ward (Exalted)}\\
\textit{Fenthras} & \textit{Fenthras (Dormant)} & \textit{Fenthras (Awakened)} & \textit{Fenthras (Exalted)}\\
\textit{Honor's Last Stand} & \textit{Honor's Last Stand (Dormant)} & \textit{Honor's Last Stand (Awakened)} & \textit{Honor's Last Stand (Exalted)}\\
\textit{Kiss of the Changebringer} & \textit{Kiss of the Changebringer (Dormant)} & \textit{Kiss of the Changebringer (Awakened)} & \textit{Kiss of the Changebringer (Exalted)}\\
\textit{Mythcarver} & \textit{Mythcarver (Dormant)} & \textit{Mythcarver (Awakened)} & \textit{Mythcarver (Exalted)}\\
\textit{Plate of the Dawnmartyr} & \textit{Plate of the Dawnmartyr (Dormant)} & \textit{Plate of the Dawnmartyr (Awakened)} & \textit{Plate of the Dawnmartyr (Exalted)}\\
\textit{Pyremaul} & \textit{Pyremaul (Dormant)} & \textit{Pyremaul (Awakened)} & \textit{Pyremaul (Exalted)}\\
\textit{Spire of Conflux} & \textit{Spire of Conflux (Dormant)} & \textit{Spire of Conflux (Awakened)} & \textit{Spire of Conflux (Exalted)}\\
\textit{Star Razor} & \textit{Star Razor (Dormant)} & \textit{Star Razor (Awakened)} & \textit{Star Razor (Exalted)}\\
\textit{Titanstone Knuckles} & \textit{Titanstone Knuckles (Dormant)} & \textit{Titanstone Knuckles (Awakened)} & \textit{Titanstone Knuckles (Exalted)}\\
\textit{Whisper} & \textit{Whisper (Dormant)} & \textit{Whisper (Awakened)} & \textit{Whisper (Exalted)}\\
\textit{Wraps of Dyamak} & \textit{Wraps of Dyamak (Dormant)} & \textit{Wraps of Dyamak (Awakened)} & \textit{Wraps of Dyamak (Exalted)}\\
\end{DndTable}

\section{Optional Campaign Rules}
Every campaign embarks on its heroic tale in its own way. Simple beginnings for a ragtag lot of sellswords. Students of the world falling in with the wrong crowd. Unwitting pawns of destiny thrust into plots much larger than themselves. However you wish to formulate your story and world—whether as player or GM—the rules of the game can help set a clear vision of how to run a campaign with consistency and intended design.

That being said, sometimes the experienced group wishes to add some new flavor to the existing structure of the gaming world. You can bring in new options for old rules, or adjust to a more advanced level of difficulty in a mystical world where danger—and even death—can be easily circumvented. This section outlines a number of homebrew rules and guidelines that any campaign can adopt to further customize the gaming experience, and which can fit in especially well with a Tal'Dorei campaign.

As with any homebrew or nonstandard rules and options, any of the options presented here that appeal to the players should be approved and implemented by the Game Master. Likewise, Game Masters should \textit{discuss the possible inclusion of any of these options with players beforehand}, to ensure that everyone will have fun with the intended changes. Talk these options over and make sure everyone is aware of and fully understands the new rules.


\begin{itemize}
\begin{multicols}{3}
\item Accelerated Rests
\item Rapid Quaffing
\item Alternative Resurrection Rules: Harrowing Return
\item Alternative Resurrection Rules: Fading Spirits
\item Alternative Resurrection Rules: Taxing Return
\end{multicols}

\end{itemize}


\begin{itemize}
\begin{multicols}{3}
\item Illegal Drugs
\item Mixed Ancestry Statistics
\item Hemocraft
\item Skyships
\end{multicols}

\end{itemize}

\chapter{Allies and Adversaries of Tal'Dorei}
\DndDropCapLine{T}{}he lands of Tal'Dorei are filled with people and monsters prepared to ambush adventurers, barter with entrepreneurial explorers, and forge alliances with charismatic characters—or at least with those willing to pay.


The first part of this chapter is a discussion of just some of the many types of creatures a GM might use as NPCs in a Tal'Dorei campaign (including creatures of the races presented in \textit{chapter 4}). The second part of the chapter is a bestiary with game statistics for new NPCs and monsters unique to Tal'Dorei and other parts of Exandria.

\section{Nonplayer Creatures}
Tal'Dorei is filled with non-humanoid beings that leave their mark on the world even if they typically don't go on adventures in dungeons or journey across the world on epic quests. The types of creatures in this section are those who the player characters can encounter on their adventures as either friends or foes. Many of these creatures are fully described in the fifth edition rules. The lore here simply covers how they function in the world of Exandria.

The creatures in this section include many not meant to be used as player characters in a typical campaign, including those with non-humanoid physiology (such as \textit{centaurs}), large size (including \textit{giants}), or alien minds (\textit{aboleths} and the like). But many other creatures of Exandria are suitable for player character use even if they don't have statistics in the fifth edition core rules. Creatures such as minotaurs, \textit{gnolls}, and the many humanoids with animal-like traits ranging from eagles to Loxodon can be used as player characters if players and Game Masters are willing to put in the effort to make them work in a campaign. The same is technically true for creatures like dragons, fey, and giants, but their size and innate powers can make it difficult for them to be members of a balanced adventuring party.

\subsection{Aboleths}
Vile, intelligent, and imbued with corrupting psychic abilities, \textbf{aboleth} are three-eyed piscine beings that lurk in watery grottos deep beneath the surface of Exandria. These aberrations are irredeemably evil, and all who attempt diplomacy with an \textbf{aboleth} end up as their food—or worse, their mindless thralls. Most \textbf{aboleth} grandiloquently claim that they have no notion of morality, but only an inexorable sense of purpose that is utterly inscrutable to the lesser minds of humanoid creatures. But even such "lesser minds" can understand that an \textbf{aboleth} purpose is the psychic domination of all life, for the lives of other creatures have no value to it.

\textbf{aboleth} can be found throughout the underground waterways of Tal'Dorei, typically living in isolation with only their thralls for company. These aberrations are loath to work together, for their perfect memories never permit them to forget a grudge. It is fortunate that these arrogant creatures are all but incapable of cooperation, for if their genius-level intellects were united, they might well be unstoppable.

During the \textit{Age of Arcanum}, an \textbf{aboleth} empire covered the length and breadth of Tal'Dorei's subterranean world. This dominion's heart was \textit{Salar}, the Unseeable City. Though \textit{Salar} is naught but a ruin in the \textit{Crystalfen Caverns} beneath \textit{Emon} now, \textbf{aboleth} still journey there in hopes of finding some remnant of their lost glory. \textit{Salar} is little more than legend to the world above—but even though few have seen it and survived, the Unseeable City is no myth. Some arcanists even speculate that the lingering psychic footprint of the ruined city is what drew the first \textit{Issylran} settlers to the area, and what inspired them to build \textit{Emon} above it.

\subsection{Catfolk}
The tabaxi of Tal'Dorei are as widespread and as variable in shape and color as the cats they resemble. They usually live solitary lives, claiming small territories to call their own, though some catfolk live as small family units with little conflict. Moreover, their feline instinct for solitude doesn't get in the way of catfolk interacting with other people around them. Indeed, many catfolk are highly social around others, safe in the knowledge that they can retreat to a solitude promising safety and calm whenever they need to.

No consistent creation myth exists for these people, who can be found living as cave dwellers in the \textit{Cliffkeep Mountains}, as city folk with their own private flats in \textit{Emon}, as cunning trackers and community builders in the \textit{Rifenmist Jungle}, and more. Generally, catfolk believe that their kind were all once normal animals, transformed by incredible magic in the ancient past. This transformation might have been a gift from the gods when magic washed over the world during the \textit{Age of Arcanum}, or when Exandria trembled beneath the divine powers unleashed in the \textit{Calamity}.

\subsection{Centaurs}
During the \textit{Calamity}, many \textit{elves} fled into the Fey Realm to escape that terrible war. One elven tribe that remained on Exandria to fight alongside their gods were known as the Horselords of Laphithas. The \textit{elves} of Laphithas rode alongside the The Arch Heart against the mortal servants of the The Ruiner, but the entire tribe was slain in a burst of chaotic magic when the The Arch Heart stabbed the The Ruiner through one eye. The overflowing power of the wounded god merged these \textit{elves'} mortal bodies with the bodies of their steeds even as they died.

The The Arch Heart rewarded the \textit{elves} of Laphithas by reviving them as \textbf{centaur}, who would ride across the \textit{Dividing Plains} forever free of war. \textbf{centaur} have long since diverged from their elven ancestry, viewing themselves as a people all their own. Still, some \textbf{centaur} throw their arms wide to the \textit{elves} of \textit{Syngorn} as long-lost cousins, even as others nurture an ancient grudge against them for fleeing when their people's need was greatest. Many of the centaur herds of the plains are incredibly faithful, worshiping the The Arch Heart with a fervor unknown even among the \textit{elves}.

\textbf{centaur} share an origin with \textit{orcs}, who were also originally \textit{elves} before being changed by the conflict between the The Arch Heart and the The Ruiner. Many \textbf{centaur} and \textit{orcs} of the \textit{Dividing Plains} see one another as kindred, but just as with real family, this bond can engender feelings of intense affection or equally intense hatred—or simply indifference. The circumstances of their lives determine the way \textit{orcs} and \textbf{centaur} feel about one another, though their kinship brings them into frequent contact either way.

\subsection{Dragons}
Dragons have existed in Tal'Dorei since the \textit{Founding}, when these mighty beings were created by the The Platinum Dragon and the The Scaled Tyrant to protect their interests on the Material Plane. All folk of Tal'Dorei know the power of dragons, and many people have witnessed it firsthand. Even the youngest humanoid victims of the \textit{Chroma Conclave's} attack on \textit{Emon} are old enough now to be parents, and to have passed down tales of the terror that filled the city's streets as doom rained down from above.

\subsubsection{Metallic Dragons}
Though countless popular stories tell of heroic warriors slaying evil, greedy, and monstrous dragons, the oldest legends of dragonkind speak of their roles as benevolent protectors of the weak. In the time of the world's beginning, the The Platinum Dragon created beings in his image, with scales of shining metal and kindly hearts. Those metallic dragons helped the mortal people of Exandria survive the turbulence of the \textit{Founding}. But since those days, most metallic dragons have vanished from the world, leaving their kindness and virtue to linger on in myth and legend rather than the chronicles of history.

On clear days, a lucky traveler might look up to the blue skies of the \textit{Dividing Plains} and feel their heart thump faster at the sight of sun glinting off silvery wings high above the world. For the most part, though, those metallic dragons who remain in the world use their shapechanging magic to live among humanoids, watching quietly from the background and gently nudging people in the direction of prosperity. After all, the The Platinum Dragon creations still look upon mortals as their charges—but times have changed. Mortals are now too numerous for the scant remaining metallic dragons to guide them as they once did. But a watchful observer can see the ills of society and whisper a word of caution into the ear of a mortal who might do something about it.

Lest any forget, the metallic dragons are not all friends to mortalkind. For some of these protectors of the small, time has caused them to grow selfish and disillusioned with their creator's commandment. These venal metallic dragons use their silver tongues and sterling reputations to manipulate mortals.

\subsubsection{Chromatic Dragons}
Fire, acid, poison, ice, and lightning rained devastation upon the world when the Betrayer God known as the The Scaled Tyrant created the unparalleled warriors that would serve her whims. The greed and malice of chromatic dragons are the stuff of countless legends—as are the mountains of gold and other treasures amassed within their lairs.

Dozens of cruel chromatic dragons lurk in the darkest depths of the forests, the highest mountain peaks, and the most arid, barren wastelands of Tal'Dorei. Though the devastation of the \textit{Chroma Conclave} is not yet three decades past, scores of adventurers have chosen to forget the horror of seeing flesh melt from bone beneath Umbrasyl's acid breath, or beholding scores choke in Raishan's poison fumes. Those adventurers now seek glory by hunting young and ambitious chromatic dragons, and claiming their nascent treasure hoards.

At the same time, not all chromatic dragons cling to the ways of their creator. A number of such dragons witnessed the horrors unleashed by the \textit{Cinder King} and his minions, and recoiled in disgust. Tales are told of how \textit{Westruun} was the beneficiary of a donation of over fifty thousand gold pieces from a mysterious young red dragon who supposedly admired the \textit{Cinder King's} ambition, but was appalled by his flagrant disregard for life.

\subsection{Driders}
The Betrayer God known as the The Spider Queen was one of the first casualties of the \textit{Calamity}—impaled by the The Stormlord thunderspear before her dark elf soldiers ever saw battle. Her blood, thick and silvery, seeped into the earth and pooled in the sunless caverns beneath Exandria. In the centuries that followed her defeat, these pools of divine blood were discovered by the The Spider Queen most loyal dark elf servants, and they drank deeply.

The blood of the The Spider Queen granted these drow the goddess's terrible power, and warped their bodies into forms that are half-spider, half-elf. Called \textbf{drider}, these beings grew in number as more and more drow accepted this gift. In time, the pools of divine blood began to run dry, and the \textbf{drider} found themselves tormented by its absence. They now roam the caverns and tunnels beneath Exandria, desperately searching for hidden pools of the silver nectar.

\subsection{Fey}
People don't discover the fey; rather, the fey discover people. Though enchanted forests are their favored domain, fey folk can be found wherever the power of nature is strong—even a flower garden in the middle of \textit{Emon}. Of course, \textit{Emonian} scholars endlessly debate whether the fey typically come to an enchanted forest because of its magic, or whether the magic of such sites blooms because of the fey. The answers to such mysteries are known to the Archfey, of course—but they don't like spilling secrets.

The Fey Realm is a place of intense emotion. There, happiness becomes elation, anger becomes fury, and passions become obsessions. The very shapes of the fey are defined by their innermost feelings, and even fey siblings can take different forms based on what they feel most strongly about. Two daughters of \textbf{pixie} parents are just as likely to be a \textbf{sprite} and a \textbf{satyr} if one sister is drawn to war and the other to revelry. And though uncommon, it is not unheard of for a fey to completely and spontaneously change form after a life-altering event. Since most fey live for upward of three centuries, a particularly unstable fey might take a dozen forms over the course of their life, wildly changing their physical form, voice, gender, likes and dislikes, and more.

Though the fey are beings of chaos and change, one immutable factor of fey life is that all are beholden to the Archfey. The immortal and imperious rulers of feykind, the Archfey rarely travel to Exandria, instead spending their time scheming against each other and plotting elaborate political intrigues. When mortals enter their realm, the Archfey take notice and often spy on them in person, using clever disguises and misdirection to keep those mortals from disrupting the Archfey's own maneuvering. Vox Machina was once pursued by Lord Artagan of the Morncrown, a jovial Archfey who finds the incursions of mortals more amusing than distressing. And there are rumors that Artagan has made his presence known on the Material Plane as well, as a mysterious being known as the The Traveler.

The mages of the \textit{Arcana Pansophical} have confirmed the identities of three other Archfey, though little else is known about them. They include Lord Saundor the Forsaken, the reportedly deceased Master of the Shademurk; Lady Elmenore the Unforgiving, High Warqueen of the Burning Vale and Matriarch of the Seelie Court; and Potentate Sammanar, They-Who-Walk-Unseen, keeper of the Sun's Shadow and Master of the Unseelie Court. Additionally, a spy of the \textit{Arcana Pansophical} traveling in the Fey Realm is said to have uncovered rumors of the existence of an Archfey known as the Keeper of the Moontides, or simply the Keeper. Since delivering her first report, however, this spy has not been heard from again.

\subsection{Giantkin}
Some \textit{Issylran} myths and legend speak of the giants as the precursors to humanity, calling them the gods' failed attempts to create what would become humankind. Other tales suggest that the giants are what became of \textit{humans} who tried to ascend to godhood, making bargains with fiends and imbibing the elemental essences of primordial spirits.

Regardless of their connection to humanoids, giants have longed dwelled far from humanoid cities and settlements. The so-called highborn giantkin stay isolated because they believe humanoids to be warlike, foolish, and untrustworthy. The lowborn giants who hold a place below the highborn are reclusive because humanoid legends paint them as terrible monsters to be hunted and slain.

Giants have a natural life span of about 300 years, though lives of hardship and toil mean that lowborn giants rarely live past 200 years. Highborn giants in their isolated and idyllic cities usually live full lives and die of old age more than any other cause—though highborn politics can be cutthroat, and assassination has ended the long life of many a giant over the centuries.

\subsubsection{Lowborn Giants}
\textbf{ogre}, \textbf{cyclops}, and \textbf{hill giant} occupy the lowest rungs of giant society. The highborn keep them at arm's length, using them as a sort of barrier between humanoid and giant civilizations. The lowborn band together, forming tight-knit social groups that must fight and forage to survive. \textbf{cyclops} typically rule lowborn settlements, and some \textbf{hill giant} warriors ritualistically tear out one eye to emulate their rulers. Zealous \textbf{hill giant} worshipers of the The Ruiner sometimes do the same, with that god granting them power to survive in exchange for ever-increasing sacrifices.

\textbf{cyclops}, ogres, and \textbf{hill giant} who turn to corrupting magic to help them survive lives of hunger and exile sometimes tumble into a cycle of cruelty and devotion to evil powers from which there is no escape. Their minds and bodies are warped beyond recognition, turning them into \textbf{troll} or \textbf{oni}. These giants gain tremendous power but lose any sense of compassion. Such corrupted giants typically live selfishly, and try to transform lowborn enclaves into their own personal armies of conquest.

\subsubsection{Highborn Giants}
The humanoid peoples of Tal'Dorei divide the highborn giants into five types: \textbf{cloud giant}, \textbf{fire giant}, \textbf{frost giant}, \textbf{stone giant}, and \textbf{storm giant}. These giants are considered one people divided into different cultures, having adopted different magics, fighting styles, and codes of ethics and law. Highborn giants are proud of their ability to look their kin in the eye as equals, making it ironic how few of them acknowledge the caste divide between highborn and lowborn.

The \textbf{hill giant} were once considered highborn, but were stripped of their nobility and exiled by the \textit{Council of Seven Scepters}, and their culture has long since fallen into ruin.

\subparagraph{Cloud Giants}
These lofty giants are denizens of Jovatthon, the Castle of Thunder. This mighty citystate flies invisibly above Tal'Dorei, enshrouded in thick storm clouds and tracing a clockwise circular route around the continent that avoids its cloudless, arid heartland. The cloud giants are ruled by two wedded monarchs, currently a City-King in charge of domestic affairs, and a War-King tasked with the defense of Jovatthon against dragons and other highborn giants. The current City-King is a stern, fearsome ruler named Ardokles, while the War-King is a much more charismatic and beloved general named Vasilios.

\subparagraph{Fire Giants}
These subterranean giants dwell in the same tunnels as \textbf{drow} and \textbf{duergar}, typically beneath volcanically active mountain ranges whose natural magma fuels their forges. The heart of fire giant civilization is the city-state of Vulkanon, located within the deep magma tubes beneath the \textit{Cliffkeep Mountains}. Though their citadels are spacious and vast, many \textbf{fire giant} have a perpetually stooped posture due to the low clearance of their underground homes. They are masters of the forge, and have been known to occasionally ally with the \textbf{duergar} of the \textit{Emberhold} to produce some of the finest arms ever seen in Tal'Dorei. Weapons of Vulkanon steel are prized by mercenaries across the continent, as much for their value as status symbols as their worth in combat.

The ruler of Vulkanon, the Grand Legate Phaestor, has long dreamed of expanding his empire, and the \textbf{fire giant} \textbf{duergar} allies have granted them a strategic foothold in Tal'Dorei's darkest depths. \textbf{fire giant} might prove to be a dangerous enemy to the folk of the surface world—but they might also be a fierce ally in the fight against the aberrations lurking in the depths, including the \textit{aboleths} of \textit{Salar}.

\subparagraph{Frost Giants}
Few beings can survive the \textit{Neverfields}, the frozen wasteland north of the \textit{Cliffkeep Mountains}. But the \textbf{frost giant}, riding on the backs of chitinous \textbf{remorhaz} and \textbf{mammoth}, have risen to be the undisputed rulers of this frigid land. Though their people are mostly nomadic reavers led by petty jarls, they all swear fealty to High King Jorskymmar, who rules from the fortress-city of Jotunborg.

The \textbf{frost giant} of Jotunborg are renowned for their ability to tame the fearsome beasts of the \textit{Neverfields}—including the white dragons that haunt the skies of that land. This skill and their knack for surviving in an inhospitable realm has made the giants of Jotunborg a pious people, and most pray to the The Wildmother for safety before crossing the ice fields. Some giants who have suffered near-death experiences, whether being trapped beneath an icy lake or smothered by an avalanche, clandestinely join one of the cults of the The Chained Oblivion, seeking solace in the inevitability of eternal cold.

\subparagraph{Stone Giants}
Few things are more striking than the sight of a \textbf{stone giant} atop a mountain peak, silhouetted against the sun. Some mountain climbers fortunate enough to see such a sight become convinced that they have gazed upon the The Dawnfather himself. \textbf{stone giant} are happy to encourage such rumors, for it makes their existence seem more like myth than fact. These reclusive highborn giantkin live in fortresses carved from the peaks of mountains. Their greatest stronghold is Skyanchor Citadel, a city hidden within a peak of the \textit{Alabaster Sierras}, its halls winding from the mountain's deepest roots to its highest summit.

\textbf{Stone giant} culture holds physical perfection as the greatest of all virtues, and \textbf{stone giant} are obsessed with physical appearance and feats of strength. But unlike the lowborn giantkin who have grown strong and resilient out of struggle, \textbf{stone giant} see physical power as an abstract virtue, best used in sporting events. A quadrennial tournament of prowess called the Sky Queen's Favor is held in coliseums atop Skyanchor, and highborn giants from across Tal'Dorei gather there to compete in tests of strength, agility, and endurance. Even the \textbf{stone giant} humanoid kin, the half-giant \textit{goliaths}, are occasionally present at the Sky Queen's Favor.

\subparagraph{Storm Giants}
The most mysterious of all the highborn, the \textbf{storm giant} are rarely seen even by other giantkin. Their flying city-state of Tempestar was destroyed by an unknown force centuries ago, and all the glories of their realm were cast into the \textit{Lucidian Ocean}, hundreds of miles off the coast. \textbf{storm giant} now live solitary lives, calling no other giants their neighbors, and caring only for themselves and their families.

Centuries have passed since Tempestar was last the seat of the \textit{Council of Seven Scepters}, yet the otherwise-reclusive \textbf{storm giant} leader, Sea Dreamer Galadawna, has never missed a meeting of the council. Though she takes little interest in the day-to-day minutiae of highborn society, she ensures that the other giant leaders never forget about her people—or the thunderous might of their sorcerers.

\subparagraph{Council of Seven Scepters}
The highborn are united by a democratic council that convenes annually for one week in the tallest spire of Jovatthon, the \textbf{cloud giant} citadel. The ruler of each giant civilization—including the two rulers of the \textbf{cloud giant}—have a seat on the council, though that seat is sometimes taken by a highly placed ambassador or representative of the ruler. When the council meets, its members deliberate on how best to maintain peace and prosperity. For even as they do their best to maintain a safe distance from the humanoid peoples of Tal'Dorei, the last thing the highborn giants desire is infighting among themselves.

Two of the seven thrones of the council have remained empty for generations. Hillqueen Ovam'mura refused to relinquish her scepter when her \textbf{hill giant} people were banished from the ranks of the highborn, and the Scepter of Wodensdottr has gathered dust in Jovatthon's vaults since Woden, demigod of the storm, abandoned the \textbf{storm giant} after the fall of Tempestar. Though the Sea Dreamer of Tempestar has never missed a gathering of the council, she has thus far refused to reclaim her throne or scepter, and the responsibilities that come with both.

\subsection{Gnolls}
As their canine appearance suggests, \textit{gnolls} are descended from \textbf{hyena} that once roamed the \textit{Dividing Plains}. A pack of these creatures was caught in the divine eminence of a Demon Prince that did battle in the \textit{Calamity}. When the nimbus of that demon's unholy power faded, the beasts were transformed.

Today, most \textit{gnolls} are nomadic groups led by a matriarch and her mate. Many people of the \textit{Dividing Plains} see \textit{gnolls} as demonic monsters, and are driven to acts of cruelty against them by their fear. Nevertheless, many \textbf{gnoll} clans live side by side with the other humanoid peoples of Tal'Dorei in small towns across the plains, including the Dustpaw \textit{gnolls} and the people of \textit{Turst Fields}.

\subsection{Lizardfolk}
The \textbf{lizardfolk} of Tal'Dorei emerged from the \textit{K'Tawl Swamp} in the \textit{Age of Arcanum}. Mages from an unnamed flying city descended into the swamp and conducted intense magical experiments on the local fauna. Some mutated into \textbf{giant lizard}, others into slithering swamp serpents—and others into humanoid beings. The mages took a few of these \textbf{lizardfolk} as specimens, and left the other fruits of their experiments behind.

In time, these abandoned \textbf{lizardfolk} used their cunning minds and skilled hands to create tools. They called themselves the Kuul'tevir—noble masters of the earth and water. These folk advanced rapidly, and though they never matched the heights of the flying cities of the magi who created them, they built magnificent domains filled with towering spires and glimmering magical waterfalls.

When the decadence of the \textit{Age of Arcanum} gave way to the chaos of the \textit{Calamity}, the Kuul'tevir delved deep into corrupting magic to save their grandeur. For a time, their people prospered, bolstered by the magic of demons and evil gods. But in the end, the desire to survive at any cost turned the Kuul'tevir's rulers against one another. Cabals of petty warlocks tore the \textbf{lizardfolk} magnificent societies apart, and the evil magic they wielded began to devour even the magic that the ancient mages had used to turn them from simple reptiles into a majestic people.

Today, the majesty of the Kuul'tevir is little more than distant myth. Most \textbf{lizardfolk} live in small communities in swamps, jungles, and fiery mountains across Tal'Dorei. Some live in cities, especially those with hot, humid climates, but no \textbf{lizardfolk} today possess the supernatural poise and power that they did in the \textit{Age of Arcanum}. Today, the \textbf{lizardfolk} are a people like any other—just trying to survive in modest comfort. But many still cling to the mythic memory of their ancient civilization and have allied themselves with terrible tyrants such as \textit{Thordak the Cinder King} in an attempt to resurrect their lost glory.

The \textbf{Vos'skyriss Serpentfolk} still dwell on \textit{Visa Isle}, south of \textit{Daggerbay}. They are more serpentine than lizard-like, but it is suspected that their origins are similar to that of the \textbf{lizardfolk}.

\subsection{Sahuagin}
Known by sailors as "sea devils," the shark-like \textbf{sahuagin} are both enemy and ally to \textit{humans} who ride the waves, and are rumored to have spawned from a bottomless trench in the ocean floor. Pods of \textbf{sahuagin} have no de facto leader, and usually have no more than six members. However, several pods might band together and collaborate on tasks they could not accomplish alone, forming schools. These schools follow the will of the majority, and disobedience is rare.

The few sea devils who dominate multiple schools through personal strength are known as \textbf{sahuagin baron}, who are as feared as they are revered. \textbf{Sahuagin} schools that fall under the sway of a \textbf{sahuagin baron} often become known for their piracy and hunger for slaughter, as the members of the school are directed to steal treasure to appease their tyrant's greed. Other schools might offer their services as mercenaries, aiding those humanoids willing to pay for their ferocity.

\textbf{Sahuagin} pods sometimes fall under the corrupting influence of \textit{aboleths} or other aquatic aberrations, with these aberrations taking on roles as demigods and the \textbf{sahuagin} who obey them becoming cultists. Instead of plundering ships for treasure, these \textbf{sahuagin} engineer shipwrecks and drag the survivors before their cruel overlord, who adds those humanoids to the ranks of their brainwashed thralls. The most feared of all these unholy figures is one that truly is the progeny of one of the \textit{Betrayer Gods}: Uk'otoa, a many-eyed \textbf{leviathan} spawn of the The Cloaked Serpent that is supposedly imprisoned deep beneath the \textit{Lucidian Ocean}.

\section{New Creatures}
Tal'Dorei is filled with wondrous creatures and characters. Whether they prove to be allies or enemies to the heroes, these beings' game statistics are provided here. These creatures are designed to supplement the monsters found in the fifth edition core rules, not replace them. Tal'Dorei is also filled with countless classic monsters, and adventures here will incorporate both familiar mythic creatures and the unique creatures of Exandria.

\begin{DndSidebar}[float=!b]{Magic Items in Stat Blocks}
Some creatures and NPCs in this section wear or wield magic items that increase their power, or that give them unique traits usable in or out of combat. New magic items are identified as such, and can be found in \textit{chapter 5} of this book. All other magic items can be found in the fifth edition core rules.



\end{DndSidebar}


\begin{itemize}
\begin{multicols}{3}
\subparagraph{New Creatures}
\item \textbf{Adranach}
\item \textbf{Master Adranach}
\item \textbf{Ashari Firetamer}
\item \textbf{Ashari Skydancer}
\item \textbf{Ashari Stoneguard}
\item \textbf{Ashari Waverider}
\item \textbf{Cinderslag Elemental}
\item \textbf{Clasp Cutthroat}
\item \textbf{Clasp Enforcer}
\item \textbf{Cold Snap Spirit}
\item \textbf{Cyclops Stormcaller}
\item \textbf{Demonfeed Spider}
\item \textbf{Demonfeed Spiderling}
\item \textbf{Ember Roc}
\item \textbf{Cobalt Golem}
\item \textbf{Forge Guardian}
\item \textbf{Mage Hunter Golem}
\item \textbf{Platinum Golem}
\item \textbf{Jourrael, the Caedogeist}
\item \textbf{Kraghammer Goat-Knight}
\item \textbf{Goat-Knight Steed}
\item \textbf{Magma Landshark}
\item \textbf{Young Magma Landshark}
\item \textbf{Plainscow}
\item \textbf{Ravager Slaughter Lord}
\item \textbf{Ravager Stabby-Stabber}
\item \textbf{Remnant Chosen}
\item \textbf{Remnant Cultist}
\item \textbf{Rivermaw Brawler}
\item \textbf{Rivermaw Stormborn}
\item \textbf{Centaur Skeleton}
\item \textbf{Flaming Skeleton}
\item \textbf{Vos'skyriss Serpentfolk}
\item \textbf{Vos'skyriss Serpentfolk Ghost}
\item \textbf{Wraithroot Tree}
\end{multicols}

\end{itemize}


\begin{itemize}
\begin{multicols}{3}
\subparagraph{Vox Machina:}
\item \textbf{Grog Strongjaw}
\item \textbf{Keyleth, Voice of the Tempest}
\item \textbf{Percival de Rolo}
\item \textbf{Trinket}
\item \textbf{Vex'ahlia}
\item \textbf{Pike Trickfoot}
\item \textbf{Scanlan Shorthalt}
\item \textbf{Doty X}
\item \textbf{Taryon Darrington}
\item \textbf{Champion of Ravens}
\end{multicols}

\end{itemize}


\begin{itemize}
\begin{multicols}{3}
\subparagraph{Adventure Hook NPCs:}
\item \textbf{Drynna Hydra}
\item \textbf{Black King}
\item \textbf{Diseased Grick}
\item \textbf{Sagacitous Erusaire}
\end{multicols}

\end{itemize}


\begin{itemize}
\begin{multicols}{3}
\item \textbf{Blighted Sapling}
\end{multicols}

\end{itemize}

\chapter{Index of Artists}

\begin{itemize}
\begin{multicols}{3}
\item{.}
\DndDropCapLine{2}{}16


\textbf{Aaron J. Riley}

aaronjriley.com

\item{.}
17, 25, 86

\textbf{Adrián Ibarra Lugo}

@Ailustrar

ailustrar.com

\item{.}
235, 247

\textbf{Allie Briggs}

@AllieBriggsArt

alliebriggs.com

\item{.}
8, 146, 187–189, 234

\textbf{Ameera Sheikh}

@Mikandii

mikandii.com

\item{.}
5, 142

\textbf{Andrey Vasilchenko}

artstation.com/artmage

\item{.}
66, 69, 72, 74, 79, 81, 92, 94, 97, 108, 113, 121, 126, 139

\textbf{Andy Law}

@Hapimeses

artstation.com/andylaw1

\item{.}
226–227, 229–230

\textbf{Anna Grinenko}

@tshortik

annagrin.squarespace.com

\item{.}
153

\textbf{Ariana Orner}

@ornerine

ornerine.tumblr.com

\item{.}
160–161, 163

\textbf{Azra Wheeler}

@Azraillu

artstation.com/azraelion

\item{.}
87, 252

\textbf{Biago D'Alessandro}

artstation.com/blas-t

\item{.}
99, 122, 219, 256

\textbf{Brian Syme}

artstation.com/bryansyme

\item{.}
240, 242–243

\textbf{Caio Santos}

@BlackSalander

artstation.com/caiosantos

\item{.}
42, 112, 116, 195, 197

\textbf{Clara Daly}

@EldritchBlep

tinyraptors.com

\item{.}
18, 155, 184, 199, 253

\textbf{Claudia Ianniciello}

@ClaudiaSG4

artstation.com/claudiasg

\item{.}
28–38, 44, 46, 48, 50–55, 57, 60-61, 63

\textbf{Conceptopolis, LLC}

conceptopolis.com

\item{.}
12, 14, 107

\textbf{Cyarna Trim}

@Cyarna

artstation.com/cyarna

\item{.}
58, 261

\textbf{De La Rosa Design, LLC}

delarosadesign.com

\item{.}
62, 236, 238, 249, 259

\textbf{Elisa Serio}

@ellinainthesky

elisaserioart.com

\item{.}
272, 274

\textbf{Elliott Berggren}

@planarbindings

\item{.}
cover, 3, 224–225

\textbf{Genel Jumalon}

@GenelJumalon

artstation.com/geneljumalon

\item{.}
244

\textbf{Hunter Bonyun}

@deerlordhunter

\item{.}
246

\textbf{Ilich Henriquez}

deviantart.com/ilacha

\item{.}
140

\textbf{Isabel Gibney}

@greyopals

isabelgibney.com

\item{.}
1, 202–212

\textbf{Jessica Nguyen}

@Jessketchin

jessketchin.com

\item{.}
265

\textbf{Jessica Scates}

@jessmightwork

\item{.}
82, 182, 223, 255

\textbf{John Anthony di Giovanni}

@jadillustrated

artstation.com/jad

\item{.}
262, 268, 270

\textbf{Jonah Baumann}

@GalacticJonah

galacticjonah.com

\item{.}
261

\textbf{Kennef Riggles}

@kennefriggles

\item{.}
22, 64, 70, 77, 91, 102, 118, 125, 133, 147, 191, 200

\textbf{Kent Davis}

@iDrawBagman

artstation.com/idrawbagman

\item{.}
166–168, 170–171, 173, 175, 177, 179, 242, 278–279

\textbf{Lauren Walsh}

@LaurenWalshArt

laurenwalshart.com

\item{.}
6

\textbf{Lea Bichlmaier}

@Jeleynai

artstation.com/jeleynai

\item{.}
11, 148

\textbf{Linda Lithén}

@LindaLithen

lindalithenart.com

\item{.}
158, 241, 257

\textbf{Nguyen Hieu}

@wth153

artstation.com/wthieu153

\item{.}
45, 154, 157, 159, 162, 164

\textbf{Nikki Dawes}

@nikkidawesdraws

nikkidawes.artstation.com

\item{.}
109, 137, 151, 192, 231

\textbf{Svetoslav Petrov}

artstation.com/svetoslavpetrov

\item{.}
195

\textbf{Stanislav Dikolenko}

artstation.com/stsdklnk

\item{.}
42, 223, 260

\textbf{Wesley Griffith}

@justwesley

artstation.com/wesleygriffith

\item{.}
26, 157

\textbf{Zuzanna Wuzyk}

@Zuzartii

artstation.com/zuzartii

\end{multicols}

\end{itemize}

\chapter{Continent Map}
\chapter{Credits}

\begin{description}
\begin{multicols}{2}

\begin{description}
\item[Exandria Created By.]
\DndDropCapLine{M}{}atthew Mercer


\item[Lead Designers.]
James J. Haeck, Matthew Mercer, Hannah Rose

\item[Additional Design.]
John Stavropoulos

\item[Additional Contributions by Critical Role Cast.]
Laura Bailey, Taliesin Jaffe, Ashley Johnson, Liam O'Brien, Marisha Ray, Sam Riegel, Travis Willingham

\item[Additional Worldbuilding.]
Aabria Iyengar

\item[Managing Editor.]
Hannah Rose

\item[Editors.]
Scott Fitzgerald Gray, Marcie Wood

\item[Proofreader.]
Matt Walloch-Key

\item[Lore Keeper.]
Dani Carr

\item[Cultural Consultant.]
Basheer Ghouse

\item[Safety Consultant.]
John Stavropoulos

\item[Accessibility Consultants.]
Accessible Games, Chris Hopper, Deven Rue

\item[Layout and Graphic Design.]
Christopher J. De La Rosa, Gordon McAlpin

\item[Art Directors.]
James J. Haeck, Hannah Rose

\item[Producers.]
James J. Haeck, Hannah Rose, Ivan Van Norman

\item[Cover Illustrator.]
Genel Jumalon

\item[Lead Character Illustrator.]
Lauren Walsh

\item[Cartographer.]
Andy Law

\item[Illustrators.]
Jonah Baumann, Elliott Berggren, Lea Bichlmaier, Hunter Bonyun, Allie Briggs, Conceptopolis, Clara Daly, Kent Davis, Nikki Dawes, Biagio D'Alessandro, Stanislav Dikolenko, Isabel Gibney, Anna Grinenko, John Anthony di Giovanni, Wesley Griffith, Ilich Henriquez, Claudia Ianniciello, Genel Jumalon, Linda Lithén, Adrián Ibarra Lugo, Nguyen Hieu, Jessica Nguyen, Ariana Orner, Svetoslav Petrov, Kennef Riggles, Aaron J. Riley, Caio Santos, Jessica Scates, Elisa Serio, Ameera Sheikh, Cyarna Trim, Andrey Vasilchenko, Lauren Walsh, Azra Wheeler, Zuzanna Wuzyk

\item[Digital Edition Alt Text.]
Mysty Vander

\item[Head of Darrington Press.]
Ivan Van Norman

\item[Creative Director.]
Matthew Mercer

\item[Head of Business Development.]
Ben Van Der Fluit

\item[Marketing Manager.]
Darcy L. Ross

\item[Playtest Coordinator.]
Christina Farber

\item[Director of Retail Partnerships.]
Brittany Walloch-Key

\item[Special Thanks.]
CritRoleStats, Ajit George, James Introcaso, Arthur Loftis, Sadie Lowry, F. Wesley Schneider

\item[Playtesters.]
Michelle Nguyen Bradley, Dani Carr, Adrienne Cho, Shaunette DeTie, Dominique Dickey, Christina Farber, Matthew Gilbert, Mark Hulmes, LaTia Jacquise, Taliesin Jaffe, Matt Walloch-Key, Christopher Lockey, Sadie Lowry, Carlos Luna, Surena Marie, Stevie Morley, Caroline Pitt, Darcy L. Ross, Kyle Shire, Lauren Walsh, Marcie Wood, Critical Role Community

\item[Critical Role Team.]
Stephanie Benjamin, Diana Jeanne Calalo, Sarah Marie Campbell, Dani Carr, Niki Chi, Adrienne Cho, Shaunette DeTie, Nadia Dilbert, Steve Failows, Christina Farber, Maxwell James, Will Lamborn, Sarah Leeper, Tal Levitas, Christopher Lockey, Ed Lopez, Surena Marie, Aaron Monroy, Khoa Nguyen, Jeremiah Rivas, Rachel Romero, Max Schapiro, Vinnie Singh, Kyle Shire, Spenser Starke, Jordyn Torrence

\end{description}

\subsection{First Edition Credits}

\begin{description}
\item[Publisher.]
Green Ronin Publishing

\item[Writing and Design.]
Matthew Mercer and James Haeck

\item[Additional Development.]
Joseph Carriker and Steve Kenson

\item[Editing.]
Jennifer Lawrence and Evan Sass

\item[Proofreading.]
Caroline Pitt

\item[Art Direction and Graphic Design.]
Hal Mangold

\end{description}

\end{multicols}

\end{description}

\chapter{Spells}
\paragraph{Warning} These spells are produced by parsing 5e.tools data for any spell mentioned in the adventure. There may be erorrs in the reproduction. Double check important spells with the source material.
\addcontentsline{toc}{section}{Arcane Lock}


\DndSpellHeader%
{Arcane Lock}
{2nd level Abjuration}
{1 action}
{Touch}
{V, S, gold dust worth at least 25 gp, which the spell consumes}
{Until dispelled}
You touch a closed door, window, gate, chest, or other entryway, and it becomes locked for the duration. You and the creatures you designate when you cast this spell can open the object normally. You can also set a password that, when spoken within 5 feet of the object, suppresses this spell for 1 minute. Otherwise, it is impassable until it is broken or the spell is dispelled or suppressed. Casting \textit{knock} on the object suppresses \textit{arcane lock} for 10 minutes.

While affected by this spell, the object is more difficult to break or force open; the DC to break it or pick any locks on it increases by 10.

\addcontentsline{toc}{section}{Comprehend Languages}


\DndSpellHeader%
{Comprehend Languages}
{1st level Divination}
{1 action}
{Self}
{V, S, a pinch of soot and salt}
{1 hour}
For the duration, you understand the literal meaning of any spoken language that you hear. You also understand any written language that you see, but you must be touching the surface on which the words are written. It takes about 1 minute to read one page of text.

This spell doesn't decode secret messages in a text or a glyph, such as an arcane sigil, that isn't part of a written language.

\addcontentsline{toc}{section}{Contagion}


\DndSpellHeader%
{Contagion}
{5th level Necromancy}
{1 action}
{Touch}
{V, S}
{7 days}
Your touch inflicts disease. Make a melee spell attack against a creature within your reach. On a hit, the target is poisoned.

At the end of each of the poisoned target's turns, the target must make a Constitution saving throw. If the target succeeds on three of these saves, it is no longer poisoned, and the spell ends. If the target fails three of these saves, the target is no longer poisoned, but choose one of the diseases below. The target is subjected to the chosen disease for the spell's duration.

Since this spell induces a natural disease in its target, any effect that removes a disease or otherwise ameliorates a disease's effects apply to it.

\subparagraph{Blinding Sickness}
Pain grips the creature's mind, and its eyes turn milky white. The creature has disadvantage on Wisdom checks and Wisdom saving throws and is blinded.

\subparagraph{Filth Fever}
A raging fever sweeps through the creature's body. The creature has disadvantage on Strength checks, Strength saving throws, and attack rolls that use Strength.

\subparagraph{Flesh Rot}
The creature's flesh decays. The creature has disadvantage on Charisma checks and vulnerability to all damage.

\subparagraph{Mindfire}
The creature's mind becomes feverish. The creature has disadvantage on Intelligence checks and Intelligence saving throws, and the creature behaves as if under the effects of the \textit{confusion} spell during combat.

\subparagraph{Seizure}
The creature is overcome with shaking. The creature has disadvantage on Dexterity checks, Dexterity saving throws, and attack rolls that use Dexterity.

\subparagraph{Slimy Doom}
The creature begins to bleed uncontrollably. The creature has disadvantage on Constitution checks and Constitution saving throws. In addition, whenever the creature takes damage, it is stunned until the end of its next turn.

\addcontentsline{toc}{section}{Freedom of the Waves}


\DndSpellHeader%
{Freedom of the Waves}
{3rd level Conjuration}
{1 action}
{120 feet}
{V, S, a strand of wet hair}
{Instantaneous}
You conjure a deluge of seawater in a 15-foot-radius, 10-foot-tall cylinder centered on a point within range. This water takes the form of a tidal wave, a whirlpool, a waterspout, or another form of your choice. Each creature in the area must succeed on a Strength saving throws against your spell save DC or take 2d8 bludgeoning damage and fall prone. You can choose a number of creatures equal to your spellcasting modifier (minimum of 1) to automatically succeed on this saving throws.

If you are within the spell's area, as part of the action you use to cast the spell, you can vanish into the deluge and teleport to an unoccupied space that you can see within the spell's area.

\addcontentsline{toc}{section}{Freedom of the Winds}


\DndSpellHeader%
{Freedom of the Winds}
{5th level Abjuration}
{1 action}
{Self}
{V, S, a scrap of sailcloth}
{Concentration, up to 10 minutes}
Wind wraps around your body, tugging at your hair and clothing as your feet lift off the ground. You gain a flying speed of 60 feet. Additionally, you have Advantage and Disadvantage on ability checks to avoid being grappled, and on saving throws against being restrained or paralyzed.

When you are targeted by a spell or attack while this spell is in effect, you can use a reaction to teleport up to 60 feet to an unoccupied space you can see. If this movement takes you out of range of the triggering spell or attack, you are unaffected by it. This spell then ends when you reappear.

\addcontentsline{toc}{section}{Gust of Wind}


\DndSpellHeader%
{Gust of Wind}
{2nd level Evocation}
{1 action}
{Self (60 feet line)}
{V, S, a legume seed}
{Concentration, up to 1 minute}
A line of strong wind 60 feet long and 10 feet wide blasts from you in a direction you choose for the spell's duration. Each creature that starts its turn in the line must succeed on a Strength saving throw or be pushed 15 feet away from you in a direction following the line.

Any creature in the line must spend 2 feet of movement for every 1 foot it moves when moving closer to you.

The gust disperses gas or vapor, and it extinguishes candles, torches, and similar unprotected flames in the area. It causes protected flames, such as those of lanterns, to dance wildly and has a No effect chance to extinguish them.

As a bonus action on each of your turns before the spell ends, you can change the direction in which the line blasts from you.

\addcontentsline{toc}{section}{Mordenkainen's Magnificent Mansion}


\DndSpellHeader%
{Mordenkainen's Magnificent Mansion}
{7th level Conjuration}
{1 minute}
{300 feet}
{V, S, a miniature portal carved from ivory, a small piece of polished marble, and a tiny silver spoon, each item worth at least 5 gp}
{24 hours}
You conjure an extradimensional dwelling in range that lasts for the duration. You choose where its one entrance is located. The entrance shimmers faintly and is 5 feet wide and 10 feet tall. You and any creature you designate when you cast the spell can enter the extradimensional dwelling as long as the portal remains open. You can open or close the portal if you are within 30 feet of it. While closed, the portal is invisible.

Beyond the portal is a magnificent foyer with numerous chambers beyond. The atmosphere is clean, fresh, and warm.

You can create any floor plan you like, but the space can't exceed 50 cubes, each cube being 10 feet on each side. The place is furnished and decorated as you choose. It contains sufficient food to serve a nine course banquet for up to 100 people. A staff of 100 near-transparent servants attends all who enter. You decide the visual appearance of these servants and their attire. They are completely obedient to your orders. Each servant can perform any task a normal human servant could perform, but they can't attack or take any action that would directly harm another creature. Thus the servants can fetch things, clean, mend, fold clothes, light fires, serve food, pour wine, and so on. The servants can go anywhere in the mansion but can't leave it. Furnishings and other objects created by this spell dissipate into smoke if removed from the mansion. When the spell ends, any creatures or objects left inside the extradimensional space are expelled into the open spaces nearest to the entrance.

\addcontentsline{toc}{section}{Scrying}


\DndSpellHeader%
{Scrying}
{5th level Divination}
{10 minute}
{Self}
{V, S, a focus worth at least 1,000 gp, such as a crystal ball, a silver mirror, or a font filled with holy water}
{Concentration, up to 10 minutes}
You can see and hear a particular creature you choose that is on the same plane of existence as you. The target must make a Wisdom saving throw, which is modified by how well you know the target and the sort of physical connection you have to it. If a target knows you're casting this spell, it can fail the saving throw voluntarily if it wants to be observed.

\begin{DndTable}[header=Knowledge of Target]{Xc}

Knowledge & Save Modifier\\
Secondhand (you have heard of the target) & +5\\
Firsthand (you have met the target) & +0\\
Familiar (you know the target well) & -5\\
\end{DndTable}

\begin{DndTable}[header=Connection to Target]{Xc}

Connection & Save Modifier\\
Likeness or picture & -2\\
Possession or garment & -4\\
Body part, lock of hair, bit of nail, or the like & -10\\
\end{DndTable}

On a successful save, the target isn't affected, and you can't use this spell against it again for 24 hours.

On a failed save, the spell creates an invisible sensor within 10 feet of the target. You can see and hear through the sensor as if you were there. The sensor moves with the target, remaining within 10 feet of it for the duration. A creature that can see invisible objects sees the sensor as a luminous orb about the size of your fist.

Instead of targeting a creature, you can choose a location you have seen before as the target of this spell. When you do, the sensor appears at that location and doesn't move.

\addcontentsline{toc}{section}{Teleportation Circle}


\DndSpellHeader%
{Teleportation Circle}
{5th level Conjuration}
{1 minute}
{10 feet}
{V, rare chalks and inks infused with precious gems worth 50 gp, which the spell consumes}
{1 round}
As you cast the spell, you draw a 10-foot-diameter circle on the ground inscribed with sigils that link your location to a permanent teleportation circle of your choice whose sigil sequence you know and that is on the same plane of existence as you. A shimmering portal opens within the circle you drew and remains open until the end of your next turn. Any creature that enters the portal instantly appears within 5 feet of the destination circle or in the nearest unoccupied space if that space is occupied.

Many major temples, guilds, and other important places have permanent teleportation circles inscribed somewhere within their confines. Each such circle includes a unique sigil sequence—a string of magical runes arranged in a particular pattern. When you first gain the ability to cast this spell, you learn the sigil sequences for two destinations on the Material Plane, determined by the DM. You can learn additional sigil sequences during your adventures. You can commit a new sigil sequence to memory after studying it for 1 minute.

You can create a permanent teleportation circle by casting this spell in the same location every day for one year. You need not use the circle to teleport when you cast the spell in this way.

\addcontentsline{toc}{section}{True Seeing}


\DndSpellHeader%
{True Seeing}
{6th level Divination}
{1 action}
{Touch}
{V, S, an ointment for the eyes that costs 25 gp; is made from mushroom powder, saffron, and fat; and is consumed by the spell}
{1 hour}
This spell gives the willing creature you touch the ability to see things as they actually are. For the duration, the creature has truesight, notices secret doors hidden by magic, and can see into the Ethereal Plane, all out to a range of 120 feet.

\addcontentsline{toc}{section}{Zone of Truth}


\DndSpellHeader%
{Zone of Truth}
{2nd level Enchantment}
{1 action}
{60 feet}
{V, S}
{10 minutes}
You create a magical zone that guards against deception in a 15-foot-radius sphere centered on a point of your choice within range. Until the spell ends, a creature that enters the spell's area for the first time on a turn or starts its turn there must make a Charisma saving throw. On a failed save, a creature can't speak a deliberate lie while in the radius. You know whether each creature succeeds or fails on its saving throw.

An affected creature is aware of the spell and can thus avoid answering questions to which it would normally respond with a lie. Such creatures can be evasive in its answers as long as it remains within the boundaries of the truth.

\chapter{Creatures}
\section{Printable Statblocks}
\paragraph{Warning} These creatures are produced by parsing 5e.tools data. There may be errors in the reproduction. Double check important creatures with the source material.\clearpage

\begin{DndMonster}[float*=b,width=\textwidth + 8pt]{Aboleth}
\begin{multicols}{2}
\DndMonsterType{Large aberration, lawful evil}
\DndMonsterBasics[
armor-class = {17 (natural armor)},
hit-points = {\DndDice{18d10 + 36}},
speed = {10 ft., swim 40 ft.}
]
\DndMonsterAbilityScores[
 str = 21,
 dex = 9,
 con = 15,
 int = 18,
 wis = 15,
 cha = 18
]
\DndMonsterDetails[
saving-throws = {Con +6, Int +8, Wis +6},
skills = {History +12, Perception +10},
senses = {darkvision 120 ft., passive perception 20},
languages = {Deep Speech, telepathy 120 ft.},
challenge = 10,
]
\DndMonsterAction{Amphibious}
The aboleth can breathe air and water.

\DndMonsterAction{Mucous Cloud}
While underwater, the aboleth is surrounded by transformative mucus. A creature that touches the aboleth or that hits it with a melee attack while within 5 feet of it must make a DC 14 Constitution saving throw. On a failure, the creature is diseased for 1d4 hours. The diseased creature can breathe only underwater.

\DndMonsterAction{Probing Telepathy}
If a creature communicates telepathically with the aboleth, the aboleth learns the creature's greatest desires if the aboleth can see the creature.

\par
\DndMonsterSection{Actions}
\DndMonsterAction{Multiattack}
The aboleth makes three tentacle attacks.

\DndMonsterAction{Tentacle}
\textsl{Melee Weapon Attack:} +9 to hit, reach 10 ft., one target. \textsl{Hit:} 12 (2d6 + 5) bludgeoning damage. If the target is a creature, it must succeed on a DC 14 Constitution saving throw or become diseased. The disease has no effect for 1 minute and can be removed by any magic that cures disease. After 1 minute, the diseased creature's skin becomes translucent and slimy, the creature can't regain hit points unless it is underwater, and the disease can be removed only by \textit{heal} or another disease-curing spell of 6th level or higher. When the creature is outside a body of water, it takes 6 (1d12) acid damage every 10 minutes unless moisture is applied to the skin before 10 minutes have passed.

\DndMonsterAction{Tail}
\textsl{Melee Weapon Attack:} +9 to hit, reach 10 ft., one target. \textsl{Hit:} 15 (3d6 + 5) bludgeoning damage.

\DndMonsterAction{Enslave (3/Day)}
The aboleth targets one creature it can see within 30 feet of it. The target must succeed on a DC 14 Wisdom saving throw or be magically charmed by the aboleth until the aboleth dies or until it is on a different plane of existence from the target. The charmed target is under the aboleth's control and can't take reactions, and the aboleth and the target can communicate telepathically with each other over any distance.

Whenever the charmed target takes damage, the target can repeat the saving throw. On a success, the effect ends. No more than once every 24 hours, the target can also repeat the saving throw when it is at least 1 mile away from the aboleth.

\par
\DndMonsterSection{Legendary Actions}
Aboleth can take 3 legendary actions, choosing from the options below. Only one legendary action can be used at a time and only at the end of another creature's turn. Aboleth regains spent legendary actions at the start of its turn.

\begin{DndMonsterLegendaryActions}
\DndMonsterLegendaryAction{Detect}{The aboleth makes a Wisdom (Perception) check.
}
\DndMonsterLegendaryAction{Tail Swipe}{The aboleth makes one tail attack.
}
\DndMonsterLegendaryAction{Psychic Drain (Costs 2 Actions)}{One creature charmed by the aboleth takes 10 (3d6) psychic damage, and the aboleth regains hit points equal to the damage the creature takes.
}
\end{DndMonsterLegendaryActions}
\par \DndMonsterSection{Lair Actions}
When fighting inside its lair, an aboleth can invoke the ambient magic to take lair actions. On initiative count 20 (losing initiative ties), the aboleth takes a lair action to cause one of the following effects:


\begin{itemize}
\item The aboleth casts \textit{phantasmal force} (no components required) on any number of creatures it can see within 60 feet of it. While maintaining concentration on this effect, the aboleth can't take other lair actions. If a target succeeds on the saving throw or if the effect ends for it, the target is immune to the aboleth's phantasmal force lair action for the next 24 hours, although such a creature can choose to be affected.
\item Pools of water within 90 feet of the aboleth surge outward in a grasping tide. Any creature on the ground within 20 feet of such a pool must succeed on a DC 14 Strength saving throw or be pulled up to 20 feet into the water and knocked prone. The aboleth can't use this lair action again until it has used a different one.
\item Water in the aboleth's lair magically becomes a conduit for the creature's rage. The aboleth can target any number of creatures it can see in such water within 90 feet of it. A target must succeed on a DC 14 Wisdom saving throw or take 7 (2d6) psychic damage. The aboleth can't use this lair action again until it has used a different one.
\end{itemize}

\par \DndMonsterSection{Regional Effects}
The region containing an aboleth's lair is warped by the creature's presence, which creates one or more of the following effects:


\begin{itemize}
\item Underground surfaces within 1 mile of the aboleth's lair are slimy and wet and are difficult terrain.
\item Water sources within 1 mile of the lair are supernaturally fouled. Enemies of the aboleth that drink such water vomit it within minutes.
\item As an action, the aboleth can create an illusory image of itself within 1 mile of the lair. The copy can appear at any location the aboleth has seen before or in any location a creature charmed by the aboleth can currently see. Once created, the image lasts for as long as the aboleth maintains concentration, as if concentration on a spell. Although the image is intangible, it looks, sounds, and can move like the aboleth. The aboleth can sense, speak, and use telepathy from the image's position as if present at that position. If the image takes any damage, it disappears.
\end{itemize}

If the aboleth dies, the first two effects fade over the course of 3d10 days.

\end{multicols}
\end{DndMonster}



\begin{DndMonster}[float=t]{Adranach}
\DndMonsterType{Huge construct, U}
\DndMonsterBasics[
armor-class = {17 (natural armor)},
hit-points = {\DndDice{20d12 + 40}},
speed = {40 ft., fly 60 ft.}
]
\DndMonsterAbilityScores[
 str = 23,
 dex = 16,
 con = 14,
 int = 8,
 wis = 14,
 cha = 11
]
\DndMonsterDetails[
saving-throws = {Str +11, Dex +8},
skills = {Athletics +11, Perception +7},
damage-resistances = {damage from spells},
damage-immunities = {poison; bludgeoning, piercing, slashing from non-magical attacks},
senses = {darkvision 120 ft., passive perception 17},
languages = {understands the languages spoken by its creator but can't speak},
challenge = 12,
]
\DndMonsterAction{Arcane Leak}
When the adranach is reduced to half its hit point maximum (85 hit points), its mithral mask cracks and its arcane form begins to waver, creating a field of unstable magic around it. Any creature that starts its turn within 10 feet of the adranach or enters that area for the first time on a turn takes 10 (3d6) radiant damage. Additionally, if a creature casts a spell within this area, it must make a DC 14 ability check using its spellcasting ability. On a failure, the spell backfires, consuming the spell slot and dealing 3 (1d6) force damage to the caster for each level of the spell slot that was consumed.

\DndMonsterAction{Immutable Form}
The adranach is immune to any spell or effect that would alter its form.

\DndMonsterAction{Magic Resistance}
The adranach has Advantage and Disadvantage on saving throws against spells and other magical effects.

\DndMonsterAction{Magic Weapons}
The adranach's weapon attacks are magical.

\DndMonsterAction{Powerful Build}
The adranach counts as one size larger when determining its carrying capacity and the weight it can push, drag, or lift.

\par
\DndMonsterSection{Actions}
\DndMonsterAction{Multiattack}
The adranach makes two attacks with its claws.

\DndMonsterAction{Claws}
\textsl{Melee Weapon Attack:} +11 to hit, one target, reach 5 ft. \textsl{Hit:} 32 (4d12 + 6) slashing damage.

\DndMonsterAction{Force Bolts (Recharge 5\textendash 6)}
The adranach targets up to four creatures it can see within 60 feet of it. Each target must succeed on a DC 18 Dexterity saving throws or take 28 (8d6) force damage.

\end{DndMonster}



\begin{DndMonster}[float*=b,width=\textwidth + 8pt]{Ancient Black Dragon}
\begin{multicols}{2}
\DndMonsterType{Gargantuan dragon, chaotic evil}
\DndMonsterBasics[
armor-class = {22 (natural armor)},
hit-points = {\DndDice{21d20 + 147}},
speed = {40 ft., fly 80 ft., swim 40 ft.}
]
\DndMonsterAbilityScores[
 str = 27,
 dex = 14,
 con = 25,
 int = 16,
 wis = 15,
 cha = 19
]
\DndMonsterDetails[
saving-throws = {Dex +9, Con +14, Wis +9, Cha +11},
skills = {Perception +16, Stealth +9},
damage-immunities = {acid},
senses = {blindsight 60 ft., darkvision 120 ft., passive perception 26},
languages = {Common, Draconic},
challenge = 21,
]
\DndMonsterAction{Amphibious}
The dragon can breathe air and water.

\DndMonsterAction{Legendary Resistance (3/Day)}
If the dragon fails a saving throw, it can choose to succeed instead.

\par
\DndMonsterSection{Actions}
\DndMonsterAction{Multiattack}
The dragon can use its Frightful Presence. It then makes three attacks: one with its bite and two with its claws.

\DndMonsterAction{Bite}
\textsl{Melee Weapon Attack:} +15 to hit, reach 15 ft., one target. \textsl{Hit:} 19 (2d10 + 8) piercing damage plus 9 (2d8) acid damage.

\DndMonsterAction{Claw}
\textsl{Melee Weapon Attack:} +15 to hit, reach 10 ft., one target. \textsl{Hit:} 15 (2d6 + 8) slashing damage.

\DndMonsterAction{Tail}
\textsl{Melee Weapon Attack:} +15 to hit, reach 20 ft., one target. \textsl{Hit:} 17 (2d8 + 8) bludgeoning damage.

\DndMonsterAction{Frightful Presence}
Each creature of the dragon's choice that is within 120 feet of the dragon and aware of it must succeed on a DC 19 Wisdom saving throw or become frightened for 1 minute. A creature can repeat the saving throw at the end of each of its turns, ending the effect on itself on a success. If a creature's saving throw is successful or the effect ends for it, the creature is immune to the dragon's Frightful Presence for the next 24 hours.

\DndMonsterAction{Acid Breath (Recharge 5\textendash 6)}
The dragon exhales acid in a 90-foot line that is 10 feet wide. Each creature in that line must make a DC 22 Dexterity saving throw, taking 67 (15d8) acid damage on a failed save, or half as much damage on a successful one.

\par
\DndMonsterSection{Legendary Actions}
Ancient Black Dragon can take 3 legendary actions, choosing from the options below. Only one legendary action can be used at a time and only at the end of another creature's turn. Ancient Black Dragon regains spent legendary actions at the start of its turn.

\begin{DndMonsterLegendaryActions}
\DndMonsterLegendaryAction{Detect}{The dragon makes a Wisdom (Perception) check.
}
\DndMonsterLegendaryAction{Tail Attack}{The dragon makes a tail attack.
}
\DndMonsterLegendaryAction{Wing Attack (Costs 2 Actions)}{The dragon beats its wings. Each creature within 15 feet of the dragon must succeed on a DC 23 Dexterity saving throw or take 15 (2d6 + 8) bludgeoning damage and be knocked prone. The dragon can then fly up to half its flying speed.
}
\end{DndMonsterLegendaryActions}
\end{multicols}
\end{DndMonster}



\begin{DndMonster}[float*=b,width=\textwidth + 8pt]{Androsphinx}
\begin{multicols}{2}
\DndMonsterType{Large monstrosity, lawful neutral}
\DndMonsterBasics[
armor-class = {17 (natural armor)},
hit-points = {\DndDice{19d10 + 95}},
speed = {40 ft., fly 60 ft.}
]
\DndMonsterAbilityScores[
 str = 22,
 dex = 10,
 con = 20,
 int = 16,
 wis = 18,
 cha = 23
]
\DndMonsterDetails[
saving-throws = {Dex +6, Con +11, Int +9, Wis +10},
skills = {Arcana +9, Perception +10, Religion +15},
damage-immunities = {psychic; bludgeoning, piercing, slashing from nonmagical attacks},
senses = {truesight 120 ft., passive perception 20},
languages = {Common, Sphinx},
challenge = 17,
]
\DndMonsterAction{Inscrutable}
The sphinx is immune to any effect that would sense its emotions or read its thoughts, as well as any divination spell that it refuses. Wisdom (Insight) checks made to ascertain the sphinx's intentions or sincerity have disadvantage.

\DndMonsterAction{Magic Weapons}
The sphinx's weapon attacks are magical.

\DndMonsterAction{Spellcasting}
The sphinx is a 12th-level spellcaster. Its spellcasting ability is Wisdom (spell save DC 18, +10 to hit with spell attacks). It requires no material components to cast its spells. The sphinx has the following cleric spells prepared:
\begin{DndMonsterSpells}
  \DndMonsterSpellLevel{sacred flame, spare the dying, thaumaturgy}
   \DndMonsterSpellLevel[1][4]{command, detect evil and good, detect magic}
   \DndMonsterSpellLevel[2][3]{lesser restoration, zone of truth}
   \DndMonsterSpellLevel[3][3]{dispel magic, tongues}
   \DndMonsterSpellLevel[4][3]{banishment, freedom of movement}
   \DndMonsterSpellLevel[5][2]{flame strike, greater restoration}
   \DndMonsterSpellLevel[6][1]{heroes' feast}
\end{DndMonsterSpells}
\par
\DndMonsterSection{Actions}
\DndMonsterAction{Multiattack}
The sphinx makes two claw attacks.

\DndMonsterAction{Claw}
\textsl{Melee Weapon Attack:} +12 to hit, reach 5 ft., one target. \textsl{Hit:} 17 (2d10 + 6) slashing damage.

\DndMonsterAction{Roar (3/Day)}
The sphinx emits a magical roar. Each time it roars before finishing a long rest, the roar is louder and the effect is different, as detailed below. Each creature within 500 feet of the sphinx and able to hear the roar must make a saving throw.


\begin{description}
\item[First Roar.]
Each creature that fails a DC 18 Wisdom saving throw is frightened for 1 minute. A frightened creature can repeat the saving throw at the end of each of its turns, ending the effect on itself on a success.
\item[Second Roar.]
Each creature that fails a DC 18 Wisdom saving throw is deafened and frightened for 1 minute. A frightened creature is paralyzed and can repeat the saving throw at the end of each of its turns, ending the effect on itself on a success.
\item[Third Roar.]
Each creature makes a DC 18 Constitution saving throw. On a failed save, a creature takes 44 (8d10) thunder damage and is knocked prone. On a successful save, the creature takes half as much damage and isn't knocked prone.
\end{description}

\par
\DndMonsterSection{Legendary Actions}
Androsphinx can take 3 legendary actions, choosing from the options below. Only one legendary action can be used at a time and only at the end of another creature's turn. Androsphinx regains spent legendary actions at the start of its turn.

\begin{DndMonsterLegendaryActions}
\DndMonsterLegendaryAction{Claw Attack}{The sphinx makes one claw attack.
}
\DndMonsterLegendaryAction{Teleport (Costs 2 Actions)}{The sphinx magically teleports, along with any equipment it is wearing or carrying, up to 120 feet to an unoccupied space it can see.
}
\DndMonsterLegendaryAction{Cast a Spell (Costs 3 Actions)}{The sphinx casts a spell from its list of prepared spells, using a spell slot as normal.
}
\end{DndMonsterLegendaryActions}
\end{multicols}
\end{DndMonster}



\begin{DndMonster}[float=t]{Animated Armor}
\DndMonsterType{Medium construct, unaligned}
\DndMonsterBasics[
armor-class = {18 (natural armor)},
hit-points = {\DndDice{6d8 + 6}},
speed = {25 ft.}
]
\DndMonsterAbilityScores[
 str = 14,
 dex = 11,
 con = 13,
 int = 1,
 wis = 3,
 cha = 1
]
\DndMonsterDetails[
damage-immunities = {poison, psychic},
senses = {blindsight 60 ft. (blind beyond this radius), passive perception 6},
challenge = 1,
]
\DndMonsterAction{Antimagic Susceptibility}
The armor is incapacitated while in the area of an \textit{antimagic field}. If targeted by \textit{dispel magic}, the armor must succeed on a Constitution saving throw against the caster's spell save DC or fall unconscious for 1 minute.

\DndMonsterAction{False Appearance}
While the armor remains motionless, it is indistinguishable from a normal suit of armor.

\par
\DndMonsterSection{Actions}
\DndMonsterAction{Multiattack}
The armor makes two melee attacks.

\DndMonsterAction{Slam}
\textsl{Melee Weapon Attack:} +4 to hit, reach 5 ft., one target. \textsl{Hit:} 5 (1d6 + 2) bludgeoning damage.

\end{DndMonster}



\begin{DndMonster}[float=t]{Ankheg}
\DndMonsterType{Large monstrosity, unaligned}
\DndMonsterBasics[
armor-class = {14 (natural armor) (11 while prone)},
hit-points = {\DndDice{6d10 + 6}},
speed = {30 ft., burrow 10 ft.}
]
\DndMonsterAbilityScores[
 str = 17,
 dex = 11,
 con = 13,
 int = 1,
 wis = 13,
 cha = 6
]
\DndMonsterDetails[
senses = {darkvision 60 ft., tremorsense 60 ft., passive perception 11},
challenge = 2,
]
\par
\DndMonsterSection{Actions}
\DndMonsterAction{Bite}
\textsl{Melee Weapon Attack:} +5 to hit, reach 5 ft., one target. \textsl{Hit:} 10 (2d6 + 3) slashing damage plus 3 (1d6) acid damage. If the target is a Large or smaller creature, it is grappled (escape DC 13). Until this grapple ends, the ankheg can bite only the grappled creature and has advantage on attack rolls to do so.

\DndMonsterAction{Acid Spray (Recharge 6)}
The ankheg spits acid in a line that is 30 feet long and 5 feet wide, provided that it has no creature grappled. Each creature in that line must make a DC 13 Dexterity saving throw, taking 10 (3d6) acid damage on a failed save, or half as much damage on a successful one.

\end{DndMonster}



\begin{DndMonster}[float=t]{Ashari Firetamer}
\DndMonsterType{Medium humanoid (any), U}
\DndMonsterBasics[
armor-class = {17 (\textit{red dragon scale mail})},
hit-points = {\DndDice{16d8 + 20}},
speed = {30 ft.}
]
\DndMonsterAbilityScores[
 str = 8,
 dex = 15,
 con = 14,
 int = 12,
 wis = 18,
 cha = 11
]
\DndMonsterDetails[
skills = {Arcana +4, Nature +7},
damage-resistances = {fire},
languages = {Common, Druidic, Ignan},
challenge = 7,
]
\DndMonsterAction{Spellcasting}
The firetamer is a 9th-level spellcaster. Its spellcasting ability is Wisdom (spell save DC 15, +7 to hit with spell attacks). It has the following druid spells prepared:
\begin{DndMonsterSpells}
  \DndMonsterSpellLevel{druidcraft, mending, produce flame}
   \DndMonsterSpellLevel[1][4]{cure wounds, faerie fire, jump}
   \DndMonsterSpellLevel[2][3]{flame blade, heat metal, lesser restoration}
   \DndMonsterSpellLevel[3][3]{daylight, dispel magic, protection from energy}
   \DndMonsterSpellLevel[4][3]{freedom of movement, wall of fire}
   \DndMonsterSpellLevel[5][1]{conjure elemental}
\end{DndMonsterSpells}
\par
\DndMonsterSection{Actions}
\DndMonsterAction{Scimitar}
\textsl{Melee Weapon Attack:} +5 to hit, reach 5 ft., one target. \textsl{Hit:} 5 (1d6 + 2) slashing damage plus 14 (4d6) fire damage.

\DndMonsterAction{Flamecharm (Recharges after a resting)}
The firetamer can cast the \textit{dominate monster} spell (save DC 15) on a \textbf{fire elemental} or other elemental fire creature. If the creature has 150 or more hit points, it has Advantage and Disadvantage on the saving throws.

\par
\DndMonsterSection{Bonus Actions}
\DndMonsterAction{Flameform (Recharges after a resting)}
The firetamer can transform into a \textbf{fire elemental}. Its game statistics are replaced by the statistics of the elemental, but the firetamer retains its alignment, personality, and Intelligence, Wisdom, and Charisma scores, and all of its skill and saving throws proficiencies. While in this form, the firetamer can't cast spells. When the firetamer is reduced to 0 hit points, falls unconscious, or dies in this form, it reverts to its humanoid form. It can remain in flameform for up to 5 hours or until it reverts to its humanoid form as a bonus action.

\DndMonsterAction{Druidic Recovery}
If the firetamer is in \textbf{fire elemental} form using Flameform, it can expend a spell slot to regain 1d8 hit points per level of the spell slot expended.

\end{DndMonster}



\begin{DndMonster}[float=t]{Ashari Skydancer}
\DndMonsterType{Medium humanoid (any), U}
\DndMonsterBasics[
armor-class = {14},
hit-points = {\DndDice{14d8}},
speed = {30 ft., fly 60 ft. \textit{(with \textit{skysail})}}
]
\DndMonsterAbilityScores[
 str = 10,
 dex = 18,
 con = 10,
 int = 12,
 wis = 16,
 cha = 11
]
\DndMonsterDetails[
saving-throws = {Dex +7},
skills = {Acrobatics +7, Perception +6},
languages = {Auran, Common, Druidic},
challenge = 5,
]
\DndMonsterAction{Evasion}
If the skydancer is subjected to an effect that allows it to make a Dexterity saving throws to take only half damage, the skydancer instead takes no damage if it succeeds on the saving throws, and only half damage if it fails.

\DndMonsterAction{Flyby}
The skydancer doesn't provoke an opportunity attack when it flies out of an enemy's reach.

\DndMonsterAction{Skysail}
The skydancer flies with a special Ashari magic item called a \textit{skysail}. While the skysail's wings are extended, the skydancer can glide through the air at a speed of 60 feet, but must descend by at least 10 feet at the end of its turn and can't gain altitude.

\DndMonsterAction{Spellcasting}
The skydancer is a 3rd-level spellcaster. Its spellcasting ability is Wisdom (spell save DC 14, +6 to hit with spell attacks). It has the following druid spells prepared:
\begin{DndMonsterSpells}
  \DndMonsterSpellLevel{guidance, shillelagh}
   \DndMonsterSpellLevel[1][4]{entangle, fog cloud, jump}
   \DndMonsterSpellLevel[2][2]{gust of wind, pass without trace}
\end{DndMonsterSpells}
\par
\DndMonsterSection{Actions}
\DndMonsterAction{Multiattack}
The skydancer makes two \textit{skysail} staff attacks.

\DndMonsterAction{Skysail Staff}
\textsl{Melee Weapon Attack:} +7 to hit, reach 5 ft., one target. \textsl{Hit:} 7 (1d6 + 4) bludgeoning damage. If the skydancer makes this attack while flying, the target must make a DC 14 Dexterity saving throws, taking 21 (6d6) lightning damage on a failed save, or half as much damage on a successful one.

\DndMonsterAction{Fly (1/Day)}
The skydancer uses its \textit{skysail} to cast \textit{fly} on itself for up to 1 minute (no concentration required).

\par
\DndMonsterSection{Reactions}
\DndMonsterAction{Slow Fall}
When the skydancer takes falling, it can reduce the damage by half.

\end{DndMonster}



\begin{DndMonster}[float=t]{Ashari Stoneguard}
\DndMonsterType{Medium humanoid (any), U}
\DndMonsterBasics[
armor-class = {15 (granite half plate)},
hit-points = {\DndDice{16d8 + 80}},
speed = {30 ft.}
]
\DndMonsterAbilityScores[
 str = 18,
 dex = 10,
 con = 20,
 int = 10,
 wis = 14,
 cha = 9
]
\DndMonsterDetails[
saving-throws = {Str +7, Con +8, Wis +5},
skills = {Athletics +7, Intimidation +2},
damage-resistances = {; bludgeoning, piercing, slashing from non-magical attacks, when using skin-to-stone},
senses = {tremorsense 30 ft., passive perception 12},
languages = {Common, Druidic, Terran},
challenge = 7,
]
\DndMonsterAction{Spellcasting}
The stoneguard is a 3rd-level spellcaster. Its spellcasting ability is Wisdom (spell save DC 13, +5 to hit with spell attacks). It has the following druid spells prepared:
\begin{DndMonsterSpells}
  \DndMonsterSpellLevel{druidcraft, resistance}
   \DndMonsterSpellLevel[1][4]{goodberry, speak with animals, thunderwave}
   \DndMonsterSpellLevel[2][2]{hold person, spike growth}
\end{DndMonsterSpells}
\par
\DndMonsterSection{Actions}
\DndMonsterAction{Multiattack}
The stoneguard makes three granite maul attacks.

\DndMonsterAction{Granite Maul}
\textsl{Melee Weapon Attack:} +7 to hit, reach 5 ft., one target. \textsl{Hit:} 11 (2d6 + 4) bludgeoning damage. If the attack hits, the stoneguard can also immediately cast \textit{thunderwave} as a bonus action. This casting uses a spell slot but requires no components.

\par
\DndMonsterSection{Reactions}
\DndMonsterAction{Sentinel}
When a creature within 5 feet of the stoneguard attacks a target other than the stoneguard, the stoneguard can make one attack against that creature.

\end{DndMonster}



\begin{DndMonster}[float=t]{Ashari Waverider}
\DndMonsterType{Medium humanoid (any), U}
\DndMonsterBasics[
armor-class = {14 (\textit{hide armor of cold resistance})},
hit-points = {\DndDice{14d8 + 14}},
speed = {30 ft.}
]
\DndMonsterAbilityScores[
 str = 15,
 dex = 14,
 con = 12,
 int = 10,
 wis = 16,
 cha = 13
]
\DndMonsterDetails[
saving-throws = {Con +4, Wis +6},
skills = {Athletics +8, Nature +3},
damage-resistances = {cold},
languages = {Aquan, Common, Druidic},
challenge = 5,
]
\DndMonsterAction{Healing Tides}
Whenever the waverider casts a spell of 1st level or higher that affects a non-hostile creature, that creature regains 3 hit points (in addition to any healing the spell may provide).

\DndMonsterAction{Marine Empathy}
The waverider can speak with and understand aquatic plants and animals.

\DndMonsterAction{Spellcasting}
The waverider is a 7th-level spellcaster. Its spellcasting ability is Wisdom (spell save DC 14, +6 to hit with spell attacks). It has the following druid spells prepared:
\begin{DndMonsterSpells}
  \DndMonsterSpellLevel{druidcraft, poison spray, resistance}
   \DndMonsterSpellLevel[1][4]{create or destroy water, cure wounds, healing word}
   \DndMonsterSpellLevel[2][3]{animal messenger, lesser restoration, moonbeam}
   \DndMonsterSpellLevel[3][3]{conjure animals, water breathing, water walk}
   \DndMonsterSpellLevel[4][1]{control water}
\end{DndMonsterSpells}
\par
\DndMonsterSection{Actions}
\DndMonsterAction{Harpoon}
\textsl{Melee or Ranged Weapon Attack:} +5 to hit, reach 5 ft., range 20/60 ft., one target. \textsl{Hit:} 7 (1d10 + 2) piercing damage. If the target is a Large or smaller creature, it must succeed on a Strength contest against the waverider or be pulled up to 20 feet toward the waverider. Attacks with this weapon made while Underwater Combat do not have Advantage and Disadvantage.

\par
\DndMonsterSection{Reactions}
\DndMonsterAction{Fishform (Recharges after a resting)}
The waverider can transform into a \textbf{hunter shark} or a \textbf{giant octopus}. Its game statistics are replaced by the statistics of the chosen creature, but the waverider retains its alignment, personality, and Intelligence, Wisdom, and Charisma scores. It also retains its skill proficiencies, in addition to gaining those of the creature it transforms into. When the waverider is reduced to 0 hit points, falls unconscious, or dies in this form, it reverts to its humanoid form. It can remain in fishform for up to 3 hours or until it reverts to its humanoid form as a bonus action.

\end{DndMonster}



\begin{DndMonster}[float=t]{Azer}
\DndMonsterType{Medium elemental, lawful neutral}
\DndMonsterBasics[
armor-class = {17 (natural armor)},
hit-points = {\DndDice{6d8 + 12}},
speed = {30 ft.}
]
\DndMonsterAbilityScores[
 str = 17,
 dex = 12,
 con = 15,
 int = 12,
 wis = 13,
 cha = 10
]
\DndMonsterDetails[
saving-throws = {Con +4},
damage-immunities = {fire, poison},
languages = {Ignan},
challenge = 2,
]
\DndMonsterAction{Heated Body}
A creature that touches the azer or hits it with a melee attack while within 5 feet of it takes 5 (1d10) fire damage.

\DndMonsterAction{Heated Weapons}
When the azer hits with a metal melee weapon, it deals an extra 3 (1d6) fire damage (included in the attack).

\DndMonsterAction{Illumination}
The azer sheds bright light in a 10-foot radius and dim light for an additional 10 feet.

\par
\DndMonsterSection{Actions}
\DndMonsterAction{Warhammer}
\textsl{Melee Weapon Attack:} +5 to hit, reach 5 ft., one target. \textsl{Hit:} 7 (1d8 + 3) bludgeoning damage, or 8 (1d10 + 3) bludgeoning damage if used with two hands to make a melee attack, plus 3 (1d6) fire damage.

\end{DndMonster}



\begin{DndMonster}[float=t]{Balor}
\DndMonsterType{Huge fiend (demon), chaotic evil}
\DndMonsterBasics[
armor-class = {19 (natural armor)},
hit-points = {\DndDice{21d12 + 126}},
speed = {40 ft., fly 80 ft.}
]
\DndMonsterAbilityScores[
 str = 26,
 dex = 15,
 con = 22,
 int = 20,
 wis = 16,
 cha = 22
]
\DndMonsterDetails[
saving-throws = {Str +14, Con +12, Wis +9, Cha +12},
damage-resistances = {cold, lightning; bludgeoning, piercing, slashing from nonmagical attacks},
damage-immunities = {fire, poison},
senses = {truesight 120 ft., passive perception 13},
languages = {Abyssal, telepathy 120 ft.},
challenge = 19,
]
\DndMonsterAction{Death Throes}
When the balor dies, it explodes, and each creature within 30 feet of it must make a DC 20 Dexterity saving throw, taking 70 (20d6) fire damage on a failed save, or half as much damage on a successful one. The explosion ignites flammable objects in that area that aren't being worn or carried, and it destroys the balor's weapons.

\DndMonsterAction{Fire Aura}
At the start of each of the balor's turns, each creature within 5 feet of it takes 10 (3d6) fire damage, and flammable objects in the aura that aren't being worn or carried ignite. A creature that touches the balor or hits it with a melee attack while within 5 feet of it takes 10 (3d6) fire damage.

\DndMonsterAction{Magic Resistance}
The balor has advantage on saving throws against spells and other magical effects.

\DndMonsterAction{Magic Weapons}
The balor's weapon attacks are magical.

\par
\DndMonsterSection{Actions}
\DndMonsterAction{Multiattack}
The balor makes two attacks: one with its longsword and one with its whip.

\DndMonsterAction{Longsword}
\textsl{Melee Weapon Attack:} +14 to hit, reach 10 ft., one target. \textsl{Hit:} 21 (3d8 + 8) slashing damage plus 13 (3d8) lightning damage. If the balor scores a critical hit, it rolls damage dice three times, instead of twice.

\DndMonsterAction{Whip}
\textsl{Melee Weapon Attack:} +14 to hit, reach 30 ft., one target. \textsl{Hit:} 15 (2d6 + 8) slashing damage plus 10 (3d6) fire damage, and the target must succeed on a DC 20 Strength saving throw or be pulled up to 25 feet toward the balor.

\DndMonsterAction{Teleport}
The balor magically teleports, along with any equipment it is wearing or carrying, up to 120 feet to an unoccupied space it can see.

\end{DndMonster}



\begin{DndMonster}[float=t]{Bandit}
\DndMonsterType{Medium humanoid (any race), any non-lawful alignment}
\DndMonsterBasics[
armor-class = {12 (\textit{leather armor})},
hit-points = {\DndDice{2d8 + 2}},
speed = {30 ft.}
]
\DndMonsterAbilityScores[
 str = 11,
 dex = 12,
 con = 12,
 int = 10,
 wis = 10,
 cha = 10
]
\DndMonsterDetails[
languages = {any one language (usually Common)},
challenge = 1/8,
]
\par
\DndMonsterSection{Actions}
\DndMonsterAction{Scimitar}
\textsl{Melee Weapon Attack:} +3 to hit, reach 5 ft., one target. \textsl{Hit:} 4 (1d6 + 1) slashing damage.

\DndMonsterAction{Light Crossbow}
\textsl{Ranged Weapon Attack:} +3 to hit, range 80/320 ft., one target. \textsl{Hit:} 5 (1d8 + 1) piercing damage.

\end{DndMonster}



\begin{DndMonster}[float=t]{Banshee}
\DndMonsterType{Medium undead, chaotic evil}
\DndMonsterBasics[
armor-class = {12},
hit-points = {\DndDice{13d8}},
speed = {0 ft., fly 40 ft. \textit{(hover)}}
]
\DndMonsterAbilityScores[
 str = 1,
 dex = 14,
 con = 10,
 int = 12,
 wis = 11,
 cha = 17
]
\DndMonsterDetails[
saving-throws = {Wis +2, Cha +5},
damage-resistances = {acid, fire, lightning, thunder; bludgeoning, piercing, slashing from nonmagical attacks},
damage-immunities = {cold, necrotic, poison},
senses = {darkvision 60 ft., passive perception 10},
languages = {Common, Elvish},
challenge = 4,
]
\DndMonsterAction{Detect Life}
The banshee can magically sense the presence of living creatures up to 5 miles away that aren't undead or constructs. She knows the general direction they're in but not their exact locations.

\DndMonsterAction{Incorporeal Movement}
The banshee can move through other creatures and objects as if they were difficult terrain. She takes 5 (1d10) force damage if she ends her turn inside an object.

\par
\DndMonsterSection{Actions}
\DndMonsterAction{Corrupting Touch}
\textsl{Melee Spell Attack:} +4 to hit, reach 5 ft., one target. \textsl{Hit:} 12 (3d6 + 2) necrotic damage.

\DndMonsterAction{Horrifying Visage}
Each non-undead creature within 60 feet of the banshee that can see her must succeed on a DC 13 Wisdom saving throw or be frightened for 1 minute. A frightened target can repeat the saving throw at the end of each of its turns, with disadvantage if the banshee is within line of sight, ending the effect on itself on a success. If a target's saving throw is successful or the effect ends for it, the target is immune to the banshee's Horrifying Visage for the next 24 hours.

\DndMonsterAction{Wail (1/Day)}
The banshee releases a mournful wail, provided that she isn't in sunlight. This wail has no effect on constructs and undead. All other creatures within 30 feet of her that can hear her must make a DC 13 Constitution saving throw. On a failure, a creature drops to 0 hit points. On a success, a creature takes 10 (3d6) psychic damage.

\end{DndMonster}



\begin{DndMonster}[float=t]{Basilisk}
\DndMonsterType{Medium monstrosity, unaligned}
\DndMonsterBasics[
armor-class = {15 (natural armor)},
hit-points = {\DndDice{8d8 + 16}},
speed = {20 ft.}
]
\DndMonsterAbilityScores[
 str = 16,
 dex = 8,
 con = 15,
 int = 2,
 wis = 8,
 cha = 7
]
\DndMonsterDetails[
senses = {darkvision 60 ft., passive perception 9},
challenge = 3,
]
\DndMonsterAction{Petrifying Gaze}
If a creature starts its turn within 30 feet of the basilisk and the two of them can see each other, the basilisk can force the creature to make a DC 12 Constitution saving throw if the basilisk isn't incapacitated. On a failed save, the creature magically begins to turn to stone and is restrained. It must repeat the saving throw at the end of its next turn. On a success, the effect ends. On a failure, the creature is petrified until freed by the  \textit{greater restoration} spell or other magic.

A creature that isn't Surprise can avert its eyes to avoid the saving throw at the start of its turn. If it does so, it can't see the basilisk until the start of its next turn, when it can avert its eyes again. If it looks at the basilisk in the meantime, it must immediately make the save.

If the basilisk sees its reflection within 30 feet of it in bright light, it mistakes itself for a rival and targets itself with its gaze.

\par
\DndMonsterSection{Actions}
\DndMonsterAction{Bite}
\textsl{Melee Weapon Attack:} +5 to hit, reach 5 ft., one target. \textsl{Hit:} 10 (2d6 + 3) piercing damage plus 7 (2d6) poison damage.

\end{DndMonster}



\begin{DndMonster}[float*=b,width=\textwidth + 8pt]{Beholder}
\begin{multicols}{2}
\DndMonsterType{Large aberration, lawful evil}
\DndMonsterBasics[
armor-class = {18 (natural armor)},
hit-points = {\DndDice{19d10 + 76}},
speed = {0 ft., fly 20 ft. \textit{(hover)}}
]
\DndMonsterAbilityScores[
 str = 10,
 dex = 14,
 con = 18,
 int = 17,
 wis = 15,
 cha = 17
]
\DndMonsterDetails[
saving-throws = {Int +8, Wis +7, Cha +8},
skills = {Perception +12},
senses = {darkvision 120 ft., passive perception 22},
languages = {Deep Speech, Undercommon},
challenge = {'cr': '13', 'lair': '14'},
]
\DndMonsterAction{Antimagic Cone}
The beholder's central eye creates an area of antimagic, as in the \textit{antimagic field} spell, in a 150-foot cone. At the start of each of its turns, the beholder decides which way the cone faces and whether the cone is active. The area works against the beholder's own eye rays.

\par
\DndMonsterSection{Actions}
\DndMonsterAction{Bite}
\textsl{Melee Weapon Attack:} +5 to hit, reach 5 ft., one target. \textsl{Hit:} 14 (4d6) piercing damage.

\DndMonsterAction{Eye Rays}
The beholder shoots three of the following magical eye rays at random (reroll duplicates), choosing one to three targets it can see within 120 feet of it:


\begin{description}
\item[1. Charm Ray.]
The targeted creature must succeed on a DC 16 Wisdom saving throw or be charmed by the beholder for 1 hour, or until the beholder harms the creature.
\item[2. Paralyzing Ray.]
The targeted creature must succeed on a DC 16 Constitution saving throw or be paralyzed for 1 minute. The target can repeat the saving throw at the end of each of its turns, ending the effect on itself on a success.
\item[3. Fear Ray.]
The targeted creature must succeed on a DC 16 Wisdom saving throw or be frightened for 1 minute. The target can repeat the saving throw at the end of each of its turns, ending the effect on itself on a success.
\item[4. Slowing Ray.]
The targeted creature must succeed on a DC 16 Dexterity saving throw. On a failed save, the target's speed is halved for 1 minute. In addition, the creature can't take reactions, and it can take either an action or a bonus action on its turn, not both. The creature can repeat the saving throw at the end of each of its turns, ending the effect on itself on a success.
\item[5. Enervation Ray.]
The targeted creature must make a DC 16 Constitution saving throw, taking 36 (8d8) necrotic damage on a failed save, or half as much damage on a successful one.
\item[6. Telekinetic Ray.]
If the target is a creature, it must succeed on a DC 16 Strength saving throw or the beholder moves it up to 30 feet in any direction. It is restrained by the ray's telekinetic grip until the start of the beholder's next turn or until the beholder is incapacitated.

If the target is an object weighing 300 pounds or less that isn't being worn or carried, it is moved up to 30 feet in any direction. The beholder can also exert fine control on objects with this ray, such as manipulating a simple tool or opening a door or a container.

\item[7. Sleep Ray.]
The targeted creature must succeed on a DC 16 Wisdom saving throw or fall asleep and remain unconscious for 1 minute. The target awakens if it takes damage or another creature takes an action to wake it. This ray has no effect on constructs and undead.
\item[8. Petrification Ray.]
The targeted creature must make a DC 16 Dexterity saving throw. On a failed save, the creature begins to turn to stone and is restrained. It must repeat the saving throw at the end of its next turn. On a success, the effect ends. On a failure, the creature is petrified until freed by the  \textit{greater restoration} spell or other magic.
\item[9. Disintegration Ray.]
If the target is a creature, it must succeed on a DC 16 Dexterity saving throw or take 45 (10d8) force damage. If this damage reduces the creature to 0 hit points, its body becomes a pile of fine gray dust.

If the target is a Large or smaller nonmagical object or creation of magical force, it is disintegrated without a saving throw. If the target is a Huge or larger object or creation of magical force, this ray disintegrates a 10-foot cube of it.

\item[10. Death Ray.]
The targeted creature must succeed on a DC 16 Dexterity saving throw or take 55 (10d10) necrotic damage. The target dies if the ray reduces it to 0 hit points.
\end{description}

\par
\DndMonsterSection{Legendary Actions}
The beholder can take 3 legendary actions, using the Eye Ray option below. It can take only one legendary action at a time and only at the end of another creature's turn. The beholder regains spent legendary actions at the start of its turn.

\begin{DndMonsterLegendaryActions}
\DndMonsterLegendaryAction{Eye Ray}{The beholder uses one random eye ray.
}
\end{DndMonsterLegendaryActions}
\par \DndMonsterSection{Lair Actions}
When fighting inside its lair, a beholder can invoke the ambient magic to take lair actions. On initiative count 20 (losing initiative ties), the beholder can take one lair action to cause one of the following effects:


\begin{itemize}
\item A 50-foot-square area of ground within 120 feet of the beholder becomes slimy; that area is difficult terrain until initiative count 20 on the next round.
\item Walls within 120 feet of the beholder sprout grasping appendages until initiative count 20 on the round after next. Each creature of the beholder's choice that starts its turn within 10 feet of such a wall must succeed on a DC 15 Dexterity saving throw or be grappled. Escaping requires a successful DC 15 Strength (Athletics) or Dexterity (Acrobatics) check.
\item An eye opens on a solid surface within 60 feet of the beholder. One random eye ray of the beholder shoots from that eye at a target of the beholder's choice that it can see. The eye then closes and disappears.
\end{itemize}

The beholder can't repeat an effect until they have all been used, and it can't use the same effect two rounds in a row.

\par \DndMonsterSection{Regional Effects}
A region containing a beholder's lair is warped by the creature's unnatural presence, which creates one or more of the following effects:


\begin{itemize}
\item Creatures within 1 mile of the beholder's lair sometimes feel as if they're being watched when they aren't.
\item When the beholder sleeps, minor warps in reality occur within 1 mile of its lair and then vanish 24 hours later. Marks on cave walls might change subtly, an eerie trinket might appear where none existed before, harmless slime might coat a statue, and so on. These effects apply only to natural surfaces and to nonmagical objects that aren't on anyone's person.
\end{itemize}

If the beholder dies, these effects fade over the course of 1d10 days.

\end{multicols}
\end{DndMonster}



\begin{DndMonster}[float=t]{Stone Golem}
\DndMonsterType{Large construct, unaligned}
\DndMonsterBasics[
armor-class = {17 (natural armor)},
hit-points = {\DndDice{17d10 + 85}},
speed = {30 ft.}
]
\DndMonsterAbilityScores[
 str = 22,
 dex = 9,
 con = 20,
 int = 3,
 wis = 11,
 cha = 1
]
\DndMonsterDetails[
damage-immunities = {poison, psychic; bludgeoning, piercing, slashing from nonmagical attacks that aren't adamantine},
senses = {darkvision 120 ft., passive perception 10},
languages = {understands the languages of its creator but can't speak},
challenge = 10,
]
\DndMonsterAction{Immutable Form}
The golem is immune to any spell or effect that would alter its form.

\DndMonsterAction{Magic Resistance}
The golem has advantage on saving throws against spells and other magical effects.

\DndMonsterAction{Magic Weapons}
The golem's weapon attacks are magical.

\par
\DndMonsterSection{Actions}
\DndMonsterAction{Multiattack}
The golem makes two slam attacks.

\DndMonsterAction{Slam}
\textsl{Melee Weapon Attack:} +10 to hit, reach 5 ft., one target. \textsl{Hit:} 19 (3d8 + 6) bludgeoning damage.

\DndMonsterAction{Slow (Recharge 5\textendash 6)}
The golem targets one or more creatures it can see within 10 feet of it. Each target must make a DC 17 Wisdom saving throw against this magic. On a failed save, a target can't use reactions, its speed is halved, and it can't make more than one attack on its turn. In addition, the target can take either an action or a bonus action on its turn, not both. These effects last for 1 minute. A target can repeat the saving throw at the end of each of its turns, ending the effect on itself on a success.

\end{DndMonster}



\begin{DndMonster}[float=t]{Blighted Sapling}
\DndMonsterType{Medium plant, U}
\DndMonsterBasics[
armor-class = {10 + PB (natural armor)},
hit-points = {\DndDice{twice your druid level}},
speed = {30 ft.}
]
\DndMonsterAbilityScores[
 str = 8,
 dex = 13,
 con = 12,
 int = 4,
 wis = 8,
 cha = 3
]
\DndMonsterDetails[
damage-vulnerabilities = {fire},
damage-resistances = {necrotic, poison},
senses = {blindsight 60 ft. (blind beyond this radius), passive perception 9},
languages = {understands the languages you speak},
]
\DndMonsterAction{Blighted Resilience (10th level or higher)}
The creature has immunity to necrotic and poison damage and to the poisoned condition.

\DndMonsterAction{Toxic Demise (10th level or higher)}
When the creature is reduced to 0 hit points, it explodes in a burst of toxic mulch or fetid viscera. Each creature within 5 feet of the exploding creature must succeed on a Constitution saving throws against your spell save DC or take necrotic damage based on the creature's challenge rating (see the table below). As an action, you can also cause a summoned creature to explode, immediately killing it.

\begin{DndTable}[header=Toxic Demise Damage]{cX}

Creature CR & Damage\\
1/4 or lower & 1d4 necrotic damage\\
1/2 & 1d6 necrotic damage\\
1 or higher & (#$prompt\\_number:min=1,title=Enter the creature's CR!,default=1$#)d8\\
No challenge rating & (PB)d6\\
\end{DndTable}

\par
\DndMonsterSection{Actions}
\DndMonsterAction{Multiattack}
When you reach 14th level in this class, the blighted sapling makes two claw attacks.

\DndMonsterAction{Claws}
\textsl{Melee Weapon Attack:}  to hit, reach 5 ft., one target. \textsl{Hit:}  2d4 + PB piercing damage.

\end{DndMonster}



\begin{DndMonster}[float=t]{Bugbear}
\DndMonsterType{Medium humanoid (goblinoid), chaotic evil}
\DndMonsterBasics[
armor-class = {16 (\textit{hide armor})},
hit-points = {\DndDice{5d8 + 5}},
speed = {30 ft.}
]
\DndMonsterAbilityScores[
 str = 15,
 dex = 14,
 con = 13,
 int = 8,
 wis = 11,
 cha = 9
]
\DndMonsterDetails[
skills = {Stealth +6, Survival +2},
senses = {darkvision 60 ft., passive perception 10},
languages = {Common, Goblin},
challenge = 1,
]
\DndMonsterAction{Brute}
A melee weapon deals one extra die of its damage when the bugbear hits with it (included in the attack).

\DndMonsterAction{Surprise Attack}
If the bugbear surprises a creature and hits it with an attack during the first round of combat, the target takes an extra 7 (2d6) damage from the attack.

\par
\DndMonsterSection{Actions}
\DndMonsterAction{Morningstar}
\textsl{Melee Weapon Attack:} +4 to hit, reach 5 ft., one target. \textsl{Hit:} 11 (2d8 + 2) piercing damage.

\DndMonsterAction{Javelin}
\textsl{Melee or Ranged Weapon Attack:} +4 to hit, reach 5 ft. or range 30/120 ft., one target. \textsl{Hit:} 9 (2d6 + 2) piercing damage in melee or 5 (1d6 + 2) piercing damage at range.

\end{DndMonster}



\begin{DndMonster}[float=t]{Bulette}
\DndMonsterType{Large monstrosity, unaligned}
\DndMonsterBasics[
armor-class = {17 (natural armor)},
hit-points = {\DndDice{9d10 + 45}},
speed = {40 ft., burrow 40 ft.}
]
\DndMonsterAbilityScores[
 str = 19,
 dex = 11,
 con = 21,
 int = 2,
 wis = 10,
 cha = 5
]
\DndMonsterDetails[
skills = {Perception +6},
senses = {darkvision 60 ft., tremorsense 60 ft., passive perception 16},
challenge = 5,
]
\DndMonsterAction{Standing Leap}
The bulette's long jump is up to 30 feet and its high jump is up to 15 feet, with or without a running start.

\par
\DndMonsterSection{Actions}
\DndMonsterAction{Bite}
\textsl{Melee Weapon Attack:} +7 to hit, reach 5 ft., one target. \textsl{Hit:} 30 (4d12 + 4) piercing damage.

\DndMonsterAction{Deadly Leap}
If the bulette jumps at least 15 feet as part of its movement, it can then use this action to land on its feet in a space that contains one or more other creatures. Each of those creatures must succeed on a DC 16 Strength or Dexterity saving throw (target's choice) or be knocked prone and take 14 (3d6 + 4) bludgeoning damage plus 14 (3d6 + 4) slashing damage. On a successful save, the creature takes only half the damage, isn't knocked prone, and is pushed 5 feet out of the bulette's space into an unoccupied space of the creature's choice. If no unoccupied space is within range, the creature instead falls prone in the bulette's space.

\end{DndMonster}


\clearpage

\begin{DndMonster}[float=t]{Centaur}
\DndMonsterType{Large monstrosity, neutral good}
\DndMonsterBasics[
armor-class = {12},
hit-points = {\DndDice{6d10 + 12}},
speed = {50 ft.}
]
\DndMonsterAbilityScores[
 str = 18,
 dex = 14,
 con = 14,
 int = 9,
 wis = 13,
 cha = 11
]
\DndMonsterDetails[
skills = {Athletics +6, Perception +3, Survival +3},
languages = {Elvish, Sylvan},
challenge = 2,
]
\DndMonsterAction{Charge}
If the centaur moves at least 30 feet straight toward a target and then hits it with a pike attack on the same turn, the target takes an extra 10 (3d6) piercing damage.

\par
\DndMonsterSection{Actions}
\DndMonsterAction{Multiattack}
The centaur makes two attacks: one with its pike and one with its hooves or two with its longbow.

\DndMonsterAction{Pike}
\textsl{Melee Weapon Attack:} +6 to hit, reach 10 ft., one target. \textsl{Hit:} 9 (1d10 + 4) piercing damage.

\DndMonsterAction{Hooves}
\textsl{Melee Weapon Attack:} +6 to hit, reach 5 ft., one target. \textsl{Hit:} 11 (2d6 + 4) bludgeoning damage.

\DndMonsterAction{Longbow}
\textsl{Ranged Weapon Attack:} +4 to hit, range 150/600 ft., one target. \textsl{Hit:} 6 (1d8 + 2) piercing damage.

\end{DndMonster}



\begin{DndMonster}[float=t]{Centaur Skeleton}
\DndMonsterType{Large undead, neutral evil}
\DndMonsterBasics[
armor-class = {12},
hit-points = {\DndDice{6d10 + 12}},
speed = {50 ft.}
]
\DndMonsterAbilityScores[
 str = 17,
 dex = 14,
 con = 14,
 int = 6,
 wis = 9,
 cha = 5
]
\DndMonsterDetails[
damage-vulnerabilities = {bludgeoning},
damage-immunities = {poison},
senses = {darkvision 60 ft., passive perception 9},
languages = {understands Common but can't speak},
challenge = 2,
]
\DndMonsterAction{Charge}
If the centaur moves at least 30 feet straight toward a target and then hits it with a pike attack on the same turn, the target takes an extra 10 (3d6) piercing damage.

\par
\DndMonsterSection{Actions}
\DndMonsterAction{Multiattack}
The centaur skeleton makes two attacks: one with its pike and one with its hooves, or two with its longbow.

\DndMonsterAction{Pike}
\textsl{Melee Weapon Attack:} +5 to hit, reach 10 ft., one target. \textsl{Hit:} 8 (1d10 + 3) piercing damage.

\DndMonsterAction{Hooves}
\textsl{Melee Weapon Attack:} +5 to hit, reach 5 ft., one target. \textsl{Hit:} 10 (2d6 + 3) bludgeoning damage.

\DndMonsterAction{Longbow}
\textsl{Ranged Weapon Attack:} +4 to hit, range 150/600 ft., one target. \textsl{Hit:} 6 (1d8 + 2) piercing damage.

\end{DndMonster}



\begin{DndMonster}[float=t]{Champion of Ravens}
\DndMonsterType{Medium celestial, U}
\DndMonsterBasics[
armor-class = {20 (deathwalker's ward)},
hit-points = {\DndDice{20d10 + 40}},
speed = {30 ft., fly 60 ft.}
]
\DndMonsterAbilityScores[
 str = 14,
 dex = 20,
 con = 14,
 int = 16,
 wis = 14,
 cha = 17
]
\DndMonsterDetails[
saving-throws = {Str +5, Dex +15, Con +5, Int +13, Wis +5, Cha +6},
skills = {Acrobatics +19, Intimidation +10, Investigation +10, Perception +16, Persuasion +10, Sleight Of Hand +12, Stealth +19},
damage-resistances = {radiant, necrotic, poison; bludgeoning, piercing, slashing from nonmagical attacks},
senses = {darkvision 60 ft., passive perception 26},
languages = {Abyssal, Celestial, Common, Druidic, Elvish, Thieves' cant},
challenge = 21,
]
\DndMonsterAction{Assassinate}
During his first turn, the Champion has Advantage and Disadvantage on attack rolls against any creature that hasn't taken a turn. Any hit the Champion scores against a surprised creature is a critical hit.

\DndMonsterAction{Aura of Protection}
While the Champion is conscious, he and friendly creatures within 10 feet of him have a +3 bonus to saving throws (included above).

\DndMonsterAction{Eternal Champion}
When the Champion is reduced to 0 hit points or dies, his body is destroyed but his spirit returns to the The Matron of Ravens side, and he gains a new body after 1d4 days.

\DndMonsterAction{Evasion}
If the Champion is subjected to an effect that allows him to make a Dexterity saving throws to take only half damage, he instead takes no damage if he succeeds on the saving throws, and only half damage if he fails.

\DndMonsterAction{Fate-Touched (4/Day)}
When the Champion makes an attack roll, ability check, or saving throws, he can reroll the d20 and use either roll. Alternatively, when a creature the Champion can see makes an attack roll against him, or a saving throws against one of his spells or features, he can force that creature to reroll the attack or save and use the new roll.

\DndMonsterAction{Shadow Smite (1/Turn)}
The Champion deals an extra 49 (14d6) radiant or necrotic damage (his choice) when he hits a target with a weapon attack and he doesn't have Advantage and Disadvantage on the attack roll.

\DndMonsterAction{Unerring Sight}
The Champion's ranged weapon attacks ignore cover and cover and have no range penalty.

\DndMonsterAction{Special Equipment}
The Champion wears \textit{boots of haste} and the Deathwalker's Ward, and wields Whisper.

\DndMonsterAction{Innate Spellcasting}
The Champion's innate spellcasting ability is Charisma (+10 to hit with spell attacks, save DC 18). He can innately cast the following spells, requiring no material components:
\begin{DndMonsterSpells}
  \DndInnateSpellLevel{bane, bless, druidcraft}
  \DndInnateSpellLevel[1]{greater invisibility}
\end{DndMonsterSpells}
\par
\DndMonsterSection{Actions}
\DndMonsterAction{Multiattack}
The Champion makes three attacks with Whisper.

\DndMonsterAction{Whisper}
\textsl{Melee or Ranged Weapon Attack:} +15 to hit, reach 5 ft. or range 60 ft., one target. \textsl{Hit:} 10 (1d4 + 8) piercing damage plus 4 (1d8) psychic damage. On a critical hit, the target must succeed on a DC 18 Wisdom saving throws or be frightened of the Champion for 1 minute. On a ranged attack, the Champion can teleport with Whisper, appearing within 5 feet of the target on a hit, or in a random space within 30 feet of the target on a miss.

\DndMonsterAction{Touch of Life and Death (Recharges after a resting)}
The Champion touches a creature within 5 feet of him, and chooses to either heal or harm that creature, restoring 55 (10d10) hit points or dealing 55 (10d10) necrotic damage. If the creature is an unwilling target, it can avoid either effect with a successful DC 20 Dexterity saving throws.

\par
\DndMonsterSection{Bonus Actions}
\DndMonsterAction{Cunning Action}
The Champion takes the Dash, Disengage, or Hide action.

\DndMonsterAction{Boots of Haste (1/Day)}
The Champion clicks his heels together, casting \textit{haste} on himself. He doesn't suffer from lethargy when the spell ends.

\par
\DndMonsterSection{Reactions}
\DndMonsterAction{Uncanny Dodge}
The Champion halves the damage that he takes from an attack that hits him. He must be able to see the attacker.

\end{DndMonster}



\begin{DndMonster}[float=t]{Cinderslag Elemental}
\DndMonsterType{Large elemental, U}
\DndMonsterBasics[
armor-class = {15 (natural armor)},
hit-points = {\DndDice{12d10 + 36}},
speed = {20 ft., burrow 40 ft.}
]
\DndMonsterAbilityScores[
 str = 15,
 dex = 12,
 con = 17,
 int = 1,
 wis = 10,
 cha = 3
]
\DndMonsterDetails[
damage-resistances = {; bludgeoning, piercing, slashing from non-magical attacks},
damage-immunities = {fire, poison},
senses = {darkvision 60 ft., passive perception 10},
challenge = 8,
]
\DndMonsterAction{Searing Presence}
Any creature that starts its turn within 5 feet of the elemental takes 5 (1d10) fire damage.

\DndMonsterAction{Molten Form}
The elemental can move through a space as narrow as 1 inch wide without squeezing. In addition, the elemental can enter a hostile creature's space and stop there. The first time it enters a creature's space on a turn, that creature takes 5 (1d10) fire damage and must succeed on a DC 14 Strength saving throws or be restrained. A creature that starts its turn restrained in this way takes 5 (1d10) fire damage. It can repeat the saving throws at the end of each of its turns, freeing itself on a success.

\DndMonsterAction{Water Susceptibility}
For every 5 feet the elemental moves in water, or for every gallon of water splashed on it, it takes 1 cold damage.

\par
\DndMonsterSection{Actions}
\DndMonsterAction{Multiattack}
The elemental makes two slam attacks.

\DndMonsterAction{Slam}
\textsl{Melee Weapon Attack:} +5 to hit, reach 5 ft., one target. \textsl{Hit:} 12 (3d6 + 2) bludgeoning damage plus 5 (1d10) fire damage. If the target is a creature or a flammable object, it ignites. Until a creature takes an action to douse the fire, the target takes 5 (1d10) fire damage at the start of each of its turns.

\DndMonsterAction{Molten Gaze}
\textsl{Ranged Spell Attack:} +8 to hit, range 30 ft., one target. \textsl{Hit:} 21 (6d6) fire damage, and the target must succeed on a DC 14 Constitution saving throws or have one nonmagical item of the elemental's choice that the target is carrying instantly melt or burn to cinders.

\end{DndMonster}



\begin{DndMonster}[float=t]{Clasp Cutthroat}
\DndMonsterType{Medium humanoid (any), U}
\DndMonsterBasics[
armor-class = {15 (leather)},
hit-points = {\DndDice{8d8 + 8}},
speed = {30 ft.}
]
\DndMonsterAbilityScores[
 str = 10,
 dex = 18,
 con = 12,
 int = 8,
 wis = 14,
 cha = 9
]
\DndMonsterDetails[
saving-throws = {Dex +6},
skills = {Deception +3, Stealth +8},
languages = {Common, Thieves' cant},
challenge = 3,
]
\DndMonsterAction{Sneak Attack (1/Turn)}
The cutthroat deals an extra 14 (4d6) damage when it hits a target with a weapon attack and has Advantage and Disadvantage on the attack roll, or when the target is within 5 feet of an ally of the cutthroat that isn't incapacitated and the cutthroat doesn't have Advantage and Disadvantage on the attack roll.

\par
\DndMonsterSection{Actions}
\DndMonsterAction{Multiattack}
The cutthroat makes two shortsword or dagger attacks.

\DndMonsterAction{Shortsword}
\textsl{Melee Weapon Attack:} +6 to hit, reach 5 ft., one target. \textsl{Hit:} 7 (1d6 + 4) piercing damage.

\DndMonsterAction{Dagger}
\textsl{Melee or Ranged Weapon Attack:} +6 to hit, reach 5 ft., range 20/60 ft., one target. \textsl{Hit:} 6 (1d4 + 4) piercing damage.

\par
\DndMonsterSection{Bonus Actions}
\DndMonsterAction{Cunning Action}
The cutthroat can take the Dash, Disengage, or Hide action.

\par
\DndMonsterSection{Reactions}
\DndMonsterAction{Uncanny Dodge}
The cutthroat halves the damage that it takes from an attack that hits it. The cutthroat must be able to see the attacker.

\end{DndMonster}



\begin{DndMonster}[float=t]{Clasp Enforcer}
\DndMonsterType{Medium humanoid (any), U}
\DndMonsterBasics[
armor-class = {16 (half plate)},
hit-points = {\DndDice{12d8 + 48}},
speed = {30 ft.}
]
\DndMonsterAbilityScores[
 str = 16,
 dex = 12,
 con = 18,
 int = 8,
 wis = 11,
 cha = 14
]
\DndMonsterDetails[
skills = {Athletics +6, Intimidation +8},
languages = {Common, Thieves' cant},
challenge = 5,
]
\DndMonsterAction{Intimidating Presence}
Whenever the enforcer hits a creature with a melee attack, the target must succeed on a DC 15 Wisdom saving throws or be frightened of the enforcer until the end of the target's next turn. The enforcer and its allies have Advantage and Disadvantage on attack rolls against any creature frightened in this way. If a creature's saving throws is successful or the effect ends for it, the creature is immune to the enforcer's Intimidating Presence for the next 24 hours.

\par
\DndMonsterSection{Actions}
\DndMonsterAction{Multiattack}
The enforcer makes three warhammer attacks.

\DndMonsterAction{Warhammer}
\textsl{Melee Weapon Attack:} +6 to hit, reach 5 ft., one target. \textsl{Hit:} 8 (1d10 + 3) bludgeoning damage.

\par
\DndMonsterSection{Bonus Actions}
\DndMonsterAction{Second Wind (Recharges after a resting)}
The enforcer regains 12 hit points

\end{DndMonster}



\begin{DndMonster}[float=t]{Cloud Giant}
\DndMonsterType{Huge giant, neutral good (50%) neutral evil (50%)}
\DndMonsterBasics[
armor-class = {14 (natural armor)},
hit-points = {\DndDice{16d12 + 96}},
speed = {40 ft.}
]
\DndMonsterAbilityScores[
 str = 27,
 dex = 10,
 con = 22,
 int = 12,
 wis = 16,
 cha = 16
]
\DndMonsterDetails[
saving-throws = {Con +10, Wis +7, Cha +7},
skills = {Insight +7, Perception +7},
languages = {Common, Giant},
challenge = 9,
]
\DndMonsterAction{Keen Smell}
The giant has advantage on Wisdom (Perception) checks that rely on smell.

\DndMonsterAction{Innate Spellcasting}
The giant's innate spellcasting ability is Charisma. It can innately cast the following spells, requiring no material components:
\begin{DndMonsterSpells}
  \DndInnateSpellLevel{detect magic, fog cloud, light}
  \DndInnateSpellLevel[3e]{feather fall, fly, misty step, telekinesis}
  \DndInnateSpellLevel[1e]{control weather, gaseous form}
\end{DndMonsterSpells}
\par
\DndMonsterSection{Actions}
\DndMonsterAction{Multiattack}
The giant makes two morningstar attacks.

\DndMonsterAction{Morningstar}
\textsl{Melee Weapon Attack:} +12 to hit, reach 10 ft., one target. \textsl{Hit:} 21 (3d8 + 8) piercing damage.

\DndMonsterAction{Rock}
\textsl{Ranged Weapon Attack:} +12 to hit, range 60/240 ft., one target. \textsl{Hit:} 30 (4d10 + 8) bludgeoning damage.

\end{DndMonster}



\begin{DndMonster}[float=t]{Cobalt Golem}
\DndMonsterType{Huge construct, U}
\DndMonsterBasics[
armor-class = {21 (natural armor)},
hit-points = {\DndDice{24d12 + 144}},
speed = {40 ft.}
]
\DndMonsterAbilityScores[
 str = 28,
 dex = 9,
 con = 22,
 int = 3,
 wis = 11,
 cha = 1
]
\DndMonsterDetails[
damage-immunities = {fire, lightning, poison, psychic; bludgeoning, piercing, slashing from nonmagical weapons that aren't adamantine},
senses = {darkvision 120 ft., passive perception 10},
languages = {understands the languages of its creator but can't speak},
challenge = 18,
]
\DndMonsterAction{Conductive Body}
Whenever the golem is subjected to lightning damage, it takes no damage. Instead, its weapon attacks deal an extra 11 (2d10) lightning damage until the end of its next turn. Also until the end of the golem's next turn, any creature that touches it or hits it with a melee attack while within 5 feet of it takes 11 (2d10) lightning damage.

\DndMonsterAction{Immutable Form}
The golem is immune to any spell or effect that would alter its form.

\DndMonsterAction{Magic Resistance}
The golem has Advantage and Disadvantage on saving throws against spells and other magical effects.

\DndMonsterAction{Magic Weapons}
The golem's weapon attacks are magical.

\par
\DndMonsterSection{Actions}
\DndMonsterAction{Multiattack}
The golem makes two attacks: one with its slam and one with its sword, or it uses its Magnetic Pulse and then its Spinning Cleave.

\DndMonsterAction{Slam}
\textsl{Melee Weapon Attack:} +15 to hit, reach 5 ft., one target. \textsl{Hit:} 25 (3d10 + 9) bludgeoning damage.

\DndMonsterAction{Sword}
\textsl{Melee Weapon Attack:} +15 to hit, reach 10 ft., one target. \textsl{Hit:} 28 (3d12 + 9) slashing damage.

\DndMonsterAction{Spinning Cleave}
The golem makes a sword attack against each creature within 10 feet of it.

\DndMonsterAction{Magnetic Pulse (Recharge 6)}
The golem raises its arms and unleashes a magnetic pulse. All unattended metal objects within 30 feet of it are pulled to an unoccupied space within 5 feet of the golem. Each creature wearing metal armor or made of metal within 30 feet of the golem must succeed on a DC 20 Strength saving throws or be pulled to an unoccupied space within 5 feet of the golem, then fall prone.

\end{DndMonster}



\begin{DndMonster}[float=t]{Cold Snap Spirit}
\DndMonsterType{Medium elemental, U}
\DndMonsterBasics[
armor-class = {18 (natural armor)},
hit-points = {\DndDice{12d8 + 24}},
speed = {30 ft., fly 30 ft.}
]
\DndMonsterAbilityScores[
 str = 10,
 dex = 20,
 con = 14,
 int = 2,
 wis = 7,
 cha = 3
]
\DndMonsterDetails[
saving-throws = {Dex +8, Wis +1},
skills = {Perception +1, Stealth +8},
damage-vulnerabilities = {fire},
damage-resistances = {; slashing, piercing from nonmagical attacks},
damage-immunities = {cold, poison},
senses = {darkvision 120 ft., passive perception 11},
challenge = 5,
]
\DndMonsterAction{Sense Heat}
The cold snap spirit can sense heat within 60 feet of it.

\DndMonsterAction{Frozen Aura}
Any creature that starts its turn within 5 feet of a cold snap spirit takes 5 (1d10) cold damage.

\DndMonsterAction{Jealous Spirit}
If the spirit attacks a creature that is within 5 feet of another cold snap spirit, the attack deals an extra 10 (3d6) cold damage on a hit.

\par
\DndMonsterSection{Actions}
\DndMonsterAction{Consume Warmth}
\textsl{Melee Spell Attack:} +8 to hit, reach 5 ft., one target. \textsl{Hit:} 21 (6d6) cold damage, and the target's speed is halved until the end of the spirit's next turn.

\DndMonsterAction{Subzero Wave (1/Day)}
Each creature within 30 feet of the cold snap spirit must make a DC 15 Constitution saving throws, taking 18 (4d8) cold damage on a failed save, or half as much damage on a successful one.

\end{DndMonster}



\begin{DndMonster}[float=t]{Cultist}
\DndMonsterType{Medium humanoid (any race), any non-good alignment}
\DndMonsterBasics[
armor-class = {12 (\textit{leather armor})},
hit-points = {\DndDice{2d8}},
speed = {30 ft.}
]
\DndMonsterAbilityScores[
 str = 11,
 dex = 12,
 con = 10,
 int = 10,
 wis = 11,
 cha = 10
]
\DndMonsterDetails[
skills = {Deception +2, Religion +2},
languages = {any one language (usually Common)},
challenge = 1/8,
]
\DndMonsterAction{Dark Devotion}
The cultist has advantage on saving throws against being charmed or frightened.

\par
\DndMonsterSection{Actions}
\DndMonsterAction{Scimitar}
\textsl{Melee Weapon Attack:} +3 to hit, reach 5 ft., one creature. \textsl{Hit:} 4 (1d6 + 1) slashing damage.

\end{DndMonster}



\begin{DndMonster}[float=t]{Cyclops}
\DndMonsterType{Huge giant, chaotic neutral}
\DndMonsterBasics[
armor-class = {14 (natural armor)},
hit-points = {\DndDice{12d12 + 60}},
speed = {30 ft.}
]
\DndMonsterAbilityScores[
 str = 22,
 dex = 11,
 con = 20,
 int = 8,
 wis = 6,
 cha = 10
]
\DndMonsterDetails[
languages = {Giant},
challenge = 6,
]
\DndMonsterAction{Poor Depth Perception}
The cyclops has disadvantage on any attack roll against a target more than 30 feet away.

\par
\DndMonsterSection{Actions}
\DndMonsterAction{Multiattack}
The cyclops makes two greatclub attacks.

\DndMonsterAction{Greatclub}
\textsl{Melee Weapon Attack:} +9 to hit, reach 10 ft., one target. \textsl{Hit:} 19 (3d8 + 6) bludgeoning damage.

\DndMonsterAction{Rock}
\textsl{Ranged Weapon Attack:} +9 to hit, range 30/120 ft., one target. \textsl{Hit:} 28 (4d10 + 6) bludgeoning damage.

\end{DndMonster}



\begin{DndMonster}[float=t]{Cyclops Stormcaller}
\DndMonsterType{Large giant, U}
\DndMonsterBasics[
armor-class = {18 (\textit{chain mail})},
hit-points = {\DndDice{12d10 + 48}},
speed = {30 ft., fly 60 ft. \textit{(stormy conditions only)}}
]
\DndMonsterAbilityScores[
 str = 16,
 dex = 10,
 con = 18,
 int = 15,
 wis = 8,
 cha = 20
]
\DndMonsterDetails[
saving-throws = {Str +5, Dex +2, Con +6, Int +4, Wis +5, Cha +11},
skills = {Arcana +6},
languages = {Common, Elvish, Giant},
challenge = 10,
]
\DndMonsterAction{One-Eyed}
The cyclops has Advantage and Disadvantage on any attack roll against a target more than 30 feet away.

\DndMonsterAction{Storm Wings}
While outdoors in stormy conditions, the cyclops has a flying speed of 60 feet.

\DndMonsterAction{Supernatural Focus}
The cyclops has Advantage and Disadvantage on Constitution saving throws made to maintain concentration on spells, and cannot lose concentration because of turbulent weather.

\DndMonsterAction{Innate Spellcasting}
The cyclops's innate spellcasting ability is Charisma (spell save DC 17, +9 to hit with spell attacks). The cyclops can innately cast the following spells, requiring no material components:
\begin{DndMonsterSpells}
  \DndInnateSpellLevel{ray of frost, water walk}
  \DndInnateSpellLevel[3e]{ice storm, sleet storm, wind wall}
  \DndInnateSpellLevel[1e]{control weather, storm of vengeance}
\end{DndMonsterSpells}
\par
\DndMonsterSection{Actions}
\DndMonsterAction{Multiattack}
The cyclops makes two ice claw attacks.

\DndMonsterAction{Ice Claw}
\textsl{Melee Weapon Attack:} +7 to hit, reach 5 ft., one target. \textsl{Hit:} 13 (3d6 + 3) slashing damage plus 3 (1d6) cold damage.

\end{DndMonster}



\begin{DndMonster}[float=t]{Demonfeed Spider}
\DndMonsterType{Large fiend, chaotic evil}
\DndMonsterBasics[
armor-class = {16 (natural armor)},
hit-points = {\DndDice{10d10 + 20}},
speed = {40 ft., climb 40 ft.}
]
\DndMonsterAbilityScores[
 str = 16,
 dex = 18,
 con = 15,
 int = 6,
 wis = 10,
 cha = 6
]
\DndMonsterDetails[
saving-throws = {Dex +7},
skills = {Perception +3, Stealth +7},
damage-resistances = {cold, fire, lightning},
damage-immunities = {poison},
senses = {darkvision 120 ft., passive perception 13},
challenge = 8,
]
\DndMonsterAction{Spider Climb}
The spider can climb difficult surfaces, including upside down on ceilings, without needing to make an ability check.

\DndMonsterAction{Web Walker}
The spider ignores movement restrictions caused by webbing.

\par
\DndMonsterSection{Actions}
\DndMonsterAction{Multiattack}
The spider makes two melee attacks: one with its bite and one with its stinger; or the spider can make one attack with its web spinner followed by a melee attack.

\DndMonsterAction{Bite}
\textsl{Melee Weapon Attack:} +7 to hit, reach 5 ft., one target. \textsl{Hit:} 8 (1d10 + 3) piercing damage, and the target must make a DC 14 Constitution saving throws. On a failed save, the target takes 27 (6d8) poison damage and is poisoned for 1 minute. On a successful save, the target takes half as much damage and isn't poisoned. If this poison damage reduces the target to 0 hit points, the target is stable but poisoned for 1 hour, even after regaining hit points, and is paralyzed while poisoned in this way.

\DndMonsterAction{Stinger}
\textsl{Melee Weapon Attack:} +7 to hit, reach 5 ft., one target. \textsl{Hit:} 9 (1d12 + 3) piercing damage, and the target must succeed on a DC 14 Constitution saving throws or be paralyzed for 1 minute. The target can repeat the saving throws at the end of each of its turns, ending the effect on itself on a success.

\DndMonsterAction{Web Spinner}
\textsl{Ranged Weapon Attack:} +7 to hit, reach 40 ft., one target. \textsl{Hit:}  The target is pulled 40 feet toward the spider and is grappled (escape DC 14). The spider can grapple only one creature at a time in this way. As an action, the spider can cocoon a creature it has grappled in webbing, causing it to be restrained and ending the grappled condition.

\par
\DndMonsterSection{Reactions}
\DndMonsterAction{Arterial Spray}
When the spider takes damage while below half its hit point maximum, it can spray poisonous blood in a 15-foot cone. Each creature in the area must make a DC 14 Constitution saving throws. On a failed save, a creature takes 13 (3d8) poison damage and is poisoned for 1 minute. On a successful save, the creature takes half as much damage and isn't poisoned. If this poison damage reduces a creature to 0 hit points, the creature is stable but poisoned for 1 hour, even after regaining hit points, and is paralyzed while poisoned in this way.

\end{DndMonster}



\begin{DndMonster}[float=t]{Demonfeed Spiderling}
\DndMonsterType{Medium fiend, U}
\DndMonsterBasics[
armor-class = {13 (natural armor)},
hit-points = {\DndDice{3d8 + 3}},
speed = {30 ft., climb 30 ft.}
]
\DndMonsterAbilityScores[
 str = 10,
 dex = 12,
 con = 12,
 int = 2,
 wis = 7,
 cha = 3
]
\DndMonsterDetails[
saving-throws = {Dex +3},
skills = {Stealth +5, Perception +0},
damage-resistances = {cold, fire, lightning},
damage-immunities = {poison},
senses = {darkvision 120 ft., passive perception 10},
challenge = 1,
]
\DndMonsterAction{Spider Climb}
The spiderling can climb difficult surfaces, including upside down on ceilings, without needing to make an ability check.

\DndMonsterAction{Web Walker}
The spiderling ignores movement restrictions caused by webbing.

\par
\DndMonsterSection{Actions}
\DndMonsterAction{Bite}
\textsl{Melee Weapon Attack:} +3 to hit, reach 5 ft., one target. \textsl{Hit:} 5 (1d8 + 1) piercing damage, and the target must make a DC 13 Constitution saving throws, taking 10 (3d6) poison damage on a failed save, or half as much damage on a successful one.

\end{DndMonster}



\begin{DndMonster}[float=t]{Deva}
\DndMonsterType{Medium celestial, lawful good}
\DndMonsterBasics[
armor-class = {17 (natural armor)},
hit-points = {\DndDice{16d8 + 64}},
speed = {30 ft., fly 90 ft.}
]
\DndMonsterAbilityScores[
 str = 18,
 dex = 18,
 con = 18,
 int = 17,
 wis = 20,
 cha = 20
]
\DndMonsterDetails[
saving-throws = {Wis +9, Cha +9},
skills = {Insight +9, Perception +9},
damage-resistances = {radiant; bludgeoning, piercing, slashing from nonmagical attacks},
senses = {darkvision 120 ft., passive perception 19},
languages = {all, telepathy 120 ft.},
challenge = 10,
]
\DndMonsterAction{Angelic Weapons}
The deva's weapon attacks are magical. When the deva hits with any weapon, the weapon deals an extra 4d8 radiant damage (included in the attack).

\DndMonsterAction{Magic Resistance}
The deva has advantage on saving throws against spells and other magical effects.

\DndMonsterAction{Innate Spellcasting}
The deva's spellcasting ability is Charisma (spell save DC 17). The deva can innately cast the following spells, requiring only verbal components:
\begin{DndMonsterSpells}
  \DndInnateSpellLevel{detect evil and good}
  \DndInnateSpellLevel[1e]{commune, raise dead}
\end{DndMonsterSpells}
\par
\DndMonsterSection{Actions}
\DndMonsterAction{Multiattack}
The deva makes two melee attacks.

\DndMonsterAction{Mace}
\textsl{Melee Weapon Attack:} +8 to hit, reach 5 ft., one target. \textsl{Hit:} 7 (1d6 + 4) bludgeoning damage plus 18 (4d8) radiant damage.

\DndMonsterAction{Healing Touch (3/Day)}
The deva touches another creature. The target magically regains 20 (4d8 + 2) hit points and is freed from any curse, disease, poison, blindness, or deafness.

\DndMonsterAction{Change Shape}
The deva magically polymorphs into a humanoid or beast that has a challenge rating equal to or less than its own, or back into its true form. It reverts to its true form if it dies. Any equipment it is wearing or carrying is absorbed or borne by the new form (the deva's choice).

In a new form, the deva retains its game statistics and ability to speak, but its AC, movement modes, Strength, Dexterity, and special senses are replaced by those of the new form, and it gains any statistics and capabilities (except class features, legendary actions, and lair actions) that the new form has but that it lacks.

\end{DndMonster}



\begin{DndMonster}[float=t]{Grick}
\DndMonsterType{Medium monstrosity, neutral}
\DndMonsterBasics[
armor-class = {14 (natural armor)},
hit-points = {\DndDice{6d8}},
speed = {30 ft., climb 30 ft.}
]
\DndMonsterAbilityScores[
 str = 14,
 dex = 14,
 con = 11,
 int = 3,
 wis = 14,
 cha = 5
]
\DndMonsterDetails[
damage-resistances = {; bludgeoning, piercing, slashing from nonmagical attacks},
senses = {darkvision 60 ft., passive perception 12},
challenge = 2,
]
\DndMonsterAction{Stone Camouflage}
The grick has advantage on Dexterity (Stealth) checks made to hide in rocky terrain.

\par
\DndMonsterSection{Actions}
\DndMonsterAction{Multiattack}
The grick makes one attack with its tentacles. If that attack hits, the grick can make one beak attack against the same target.

\DndMonsterAction{Tentacles}
\textsl{Melee Weapon Attack:} +4 to hit, reach 5 ft., one target. \textsl{Hit:} 9 (2d6 + 2) slashing damage.

\DndMonsterAction{Beak}
\textsl{Melee Weapon Attack:} +4 to hit, reach 5 ft., one target. \textsl{Hit:} 5 (1d6 + 2) piercing damage.

\end{DndMonster}



\begin{DndMonster}[float=t]{Djinni}
\DndMonsterType{Large elemental, chaotic good}
\DndMonsterBasics[
armor-class = {17 (natural armor)},
hit-points = {\DndDice{14d10 + 84}},
speed = {30 ft., fly 90 ft.}
]
\DndMonsterAbilityScores[
 str = 21,
 dex = 15,
 con = 22,
 int = 15,
 wis = 16,
 cha = 20
]
\DndMonsterDetails[
saving-throws = {Dex +6, Wis +7, Cha +9},
damage-immunities = {lightning, thunder},
senses = {darkvision 120 ft., passive perception 13},
languages = {Auran},
challenge = 11,
]
\DndMonsterAction{Elemental Demise}
If the djinni dies, its body disintegrates into a warm breeze, leaving behind only equipment the djinni was wearing or carrying.

\DndMonsterAction{Innate Spellcasting}
The djinni's innate spellcasting ability is Charisma (spell save DC 17, +9 to hit with spell attacks). It can innately cast the following spells, requiring no material components:
\begin{DndMonsterSpells}
  \DndInnateSpellLevel{detect evil and good, detect magic, thunderwave}
  \DndInnateSpellLevel[3e]{create food and water, tongues, wind walk}
  \DndInnateSpellLevel[1e]{conjure elemental, creation, gaseous form, invisibility, major image, plane shift}
\end{DndMonsterSpells}
\par
\DndMonsterSection{Actions}
\DndMonsterAction{Multiattack}
The djinni makes three scimitar attacks.

\DndMonsterAction{Scimitar}
\textsl{Melee Weapon Attack:} +9 to hit, reach 5 ft., one target. \textsl{Hit:} 12 (2d6 + 5) slashing damage plus 3 (1d6) lightning or thunder damage (djinni's choice).

\DndMonsterAction{Create Whirlwind}
A 5-foot-radius, 30-foot-tall cylinder of swirling air magically forms on a point the djinni can see within 120 feet of it. The whirlwind lasts as long as the djinni maintains concentration (as if concentration on a spell). Any creature but the djinni that enters the whirlwind must succeed on a DC 18 Strength saving throw or be restrained by it. The djinni can move the whirlwind up to 60 feet as an action, and creatures restrained by the whirlwind move with it. The whirlwind ends if the djinni loses sight of it.

A creature can use its action to free a creature restrained by the whirlwind, including itself, by succeeding on a DC 18 Strength check. If the check succeeds, the creature is no longer restrained and moves to the nearest space outside the whirlwind.

\end{DndMonster}



\begin{DndMonster}[float=t]{Doppelganger}
\DndMonsterType{Medium monstrosity (shapechanger), neutral}
\DndMonsterBasics[
armor-class = {14},
hit-points = {\DndDice{8d8 + 16}},
speed = {30 ft.}
]
\DndMonsterAbilityScores[
 str = 11,
 dex = 18,
 con = 14,
 int = 11,
 wis = 12,
 cha = 14
]
\DndMonsterDetails[
skills = {Deception +6, Insight +3},
senses = {darkvision 60 ft., passive perception 11},
languages = {Common},
challenge = 3,
]
\DndMonsterAction{Shapechanger}
The doppelganger can use its action to polymorph into a Small or Medium humanoid it has seen, or back into its true form. Its statistics, other than its size, are the same in each form. Any equipment it is wearing or carrying isn't transformed. It reverts to its true form if it dies.

\DndMonsterAction{Ambusher}
In the first round of a combat, the doppelganger has advantage on attack rolls against any creature it Surprise.

\DndMonsterAction{Surprise Attack}
If the doppelganger surprises a creature and hits it with an attack during the first round of combat, the target takes an extra 10 (3d6) damage from the attack.

\par
\DndMonsterSection{Actions}
\DndMonsterAction{Multiattack}
The doppelganger makes two melee attacks.

\DndMonsterAction{Slam}
\textsl{Melee Weapon Attack:} +6 to hit, reach 5 ft., one target. \textsl{Hit:} 7 (1d6 + 4) bludgeoning damage.

\DndMonsterAction{Read Thoughts}
The doppelganger magically reads the surface thoughts of one creature within 60 feet of it. The effect can penetrate barriers, but 3 feet of wood or dirt, 2 feet of stone, 2 inches of metal, or a thin sheet of lead blocks it. While the target is in range, the doppelganger can continue reading its thoughts, as long as the doppelganger's concentration isn't broken (as if concentration on a spell). While reading the target's mind, the doppelganger has advantage on Wisdom (Insight) and Charisma (Deception, Intimidation, and Persuasion) checks against the target.

\end{DndMonster}



\begin{DndMonster}[float=t]{Doty X}
\DndMonsterType{Medium construct, U}
\DndMonsterBasics[
armor-class = {17 (natural armor)},
hit-points = {\DndDice{9d8 + 18 + 20}},
speed = {40 ft.}
]
\DndMonsterAbilityScores[
 str = 14,
 dex = 12,
 con = 14,
 int = 4,
 wis = 10,
 cha = 6
]
\DndMonsterDetails[
saving-throws = {Dex +4, Con +5},
skills = {Athletics +7, Perception +8},
damage-immunities = {poison},
senses = {darkvision 60 ft., passive perception 18},
languages = {understands the languages that \textbf{Taryon Darrington} speaks},
challenge = 6,
]
\DndMonsterAction{Arcane Surge (4/Day)}
When Doty hits a target, \textbf{Taryon Darrington} can channel magical energy to deal an extra 14 (4d6) force damage to the target as long as \textbf{Taryon Darrington} is not incapacitated.

\DndMonsterAction{Dedicated Servant}
Doty has an additional 20 hit points, and has a +2 modifier to saving throws, ability checks using skill proficiencies, and attack rolls.

\DndMonsterAction{Limited Vocabulary}
Doty can say the following words in Common: "Tary," "Yes," "Correct," "Absolutely," "Soon," and "Handsome."

\DndMonsterAction{Vigilant}
Doty can't be surprised.

\par
\DndMonsterSection{Actions}
\DndMonsterAction{Empowered Fist}
\textsl{Melee Weapon Attack:} +7 to hit, reach 5 ft., one target. \textsl{Hit:} 8 (1d8 + 4) force damage.

\DndMonsterAction{Repair (3/Day)}
Doty regains 15 (2d8 + 6) hit points.

\par
\DndMonsterSection{Reactions}
\DndMonsterAction{Deflect}
If a creature within 5 feet of Doty that he can see makes an attack against a creature other than Doty, he can impose Advantage and Disadvantage on the attack roll. The attacker then takes 6 (1d4 + 4) force damage.

\end{DndMonster}



\begin{DndMonster}[float=t]{Dretch}
\DndMonsterType{Small fiend (demon), chaotic evil}
\DndMonsterBasics[
armor-class = {11 (natural armor)},
hit-points = {\DndDice{4d6 + 4}},
speed = {20 ft.}
]
\DndMonsterAbilityScores[
 str = 11,
 dex = 11,
 con = 12,
 int = 5,
 wis = 8,
 cha = 3
]
\DndMonsterDetails[
damage-resistances = {cold, fire, lightning},
damage-immunities = {poison},
senses = {darkvision 60 ft., passive perception 9},
languages = {Abyssal, telepathy 60 ft. (works only with creatures that understand Abyssal)},
challenge = 1/4,
]
\par
\DndMonsterSection{Actions}
\DndMonsterAction{Multiattack}
The dretch makes two attacks: one with its bite and one with its claws.

\DndMonsterAction{Bite}
\textsl{Melee Weapon Attack:} +2 to hit, reach 5 ft., one target. \textsl{Hit:} 3 (1d6) piercing damage.

\DndMonsterAction{Claws}
\textsl{Melee Weapon Attack:} +2 to hit, reach 5 ft., one target. \textsl{Hit:} 5 (2d4) slashing damage.

\DndMonsterAction{Fetid Cloud (1/Day)}
A 10-foot radius of disgusting green gas extends out from the dretch. The gas spreads around corners, and its area is lightly obscured. It lasts for 1 minute or until a strong wind disperses it. Any creature that starts its turn in that area must succeed on a DC 11 Constitution saving throw or be poisoned until the start of its next turn. While poisoned in this way, the target can take either an action or a bonus action on its turn, not both, and can't take reactions.

\end{DndMonster}


\clearpage

\begin{DndMonster}[float=t]{Drider}
\DndMonsterType{Large monstrosity, chaotic evil}
\DndMonsterBasics[
armor-class = {19 (natural armor)},
hit-points = {\DndDice{13d10 + 52}},
speed = {30 ft., climb 30 ft.}
]
\DndMonsterAbilityScores[
 str = 16,
 dex = 16,
 con = 18,
 int = 13,
 wis = 14,
 cha = 12
]
\DndMonsterDetails[
skills = {Perception +5, Stealth +9},
senses = {darkvision 120 ft., passive perception 15},
languages = {Elvish, Undercommon},
challenge = 6,
]
\DndMonsterAction{Fey Ancestry}
The drider has advantage on saving throws against being charmed, and magic can't put the drider to sleep.

\DndMonsterAction{Spider Climb}
The drider can climb difficult surfaces, including upside down on ceilings, without needing to make an ability check.

\DndMonsterAction{Sunlight Sensitivity}
While in sunlight, the drider has disadvantage on attack rolls, as well as on Wisdom (Perception) checks that rely on sight.

\DndMonsterAction{Web Walker}
The drider ignores movement restrictions caused by webbing.

\DndMonsterAction{Innate Spellcasting}
The drider's innate spellcasting ability is Wisdom (spell save DC 13). The drider can innately cast the following spells, requiring no material components:
\begin{DndMonsterSpells}
  \DndInnateSpellLevel{dancing lights}
  \DndInnateSpellLevel[1e]{darkness, faerie fire}
\end{DndMonsterSpells}
\par
\DndMonsterSection{Actions}
\DndMonsterAction{Multiattack}
The drider makes three attacks, either with its longsword or its longbow. It can replace one of those attacks with a bite attack.

\DndMonsterAction{Bite}
\textsl{Melee Weapon Attack:} +6 to hit, reach 5 ft., one creature. \textsl{Hit:} 2 (1d4) piercing damage plus 9 (2d8) poison damage.

\DndMonsterAction{Longsword}
\textsl{Melee Weapon Attack:} +6 to hit, reach 5 ft., one target. \textsl{Hit:} 7 (1d8 + 3) slashing damage, or 8 (1d10 + 3) slashing damage if used with two hands.

\DndMonsterAction{Longbow}
\textsl{Ranged Weapon Attack:} +6 to hit, range 150/600 ft., one target. \textsl{Hit:} 7 (1d8 + 3) piercing damage plus 4 (1d8) poison damage.

\end{DndMonster}



\begin{DndMonster}[float=t]{Drow}
\DndMonsterType{Medium humanoid (elf), neutral evil}
\DndMonsterBasics[
armor-class = {15 (\textit{chain shirt})},
hit-points = {\DndDice{3d8}},
speed = {30 ft.}
]
\DndMonsterAbilityScores[
 str = 10,
 dex = 14,
 con = 10,
 int = 11,
 wis = 11,
 cha = 12
]
\DndMonsterDetails[
skills = {Perception +2, Stealth +4},
senses = {darkvision 120 ft., passive perception 12},
languages = {Elvish, Undercommon},
challenge = 1/4,
]
\DndMonsterAction{Fey Ancestry}
The drow has advantage on saving throws against being charmed, and magic can't put the drow to sleep.

\DndMonsterAction{Sunlight Sensitivity}
While in sunlight, the drow has disadvantage on attack rolls, as well as on Wisdom (Perception) checks that rely on sight.

\DndMonsterAction{Innate Spellcasting}
The drow's spellcasting ability is Charisma (spell save DC 11). It can innately cast the following spells, requiring no material components:
\begin{DndMonsterSpells}
  \DndInnateSpellLevel{dancing lights}
  \DndInnateSpellLevel[1e]{darkness, faerie fire}
\end{DndMonsterSpells}
\par
\DndMonsterSection{Actions}
\DndMonsterAction{Shortsword}
\textsl{Melee Weapon Attack:} +4 to hit, reach 5 ft., one target. \textsl{Hit:} 5 (1d6 + 2) piercing damage.

\DndMonsterAction{Hand Crossbow}
\textsl{Ranged Weapon Attack:} +4 to hit, range 30/120 ft., one target. \textsl{Hit:} 5 (1d6 + 2) piercing damage, and the target must succeed on a DC 13 Constitution saving throw or be poisoned for 1 hour. If the saving throw fails by 5 or more, the target is also unconscious while poisoned in this way. The target wakes up if it takes damage or if another creature takes an action to shake it awake.

\end{DndMonster}



\begin{DndMonster}[float=t]{Druid}
\DndMonsterType{Medium humanoid (any race), any alignment}
\DndMonsterBasics[
armor-class = {11 (16 with \textit{barkskin})},
hit-points = {\DndDice{5d8 + 5}},
speed = {30 ft.}
]
\DndMonsterAbilityScores[
 str = 10,
 dex = 12,
 con = 13,
 int = 12,
 wis = 15,
 cha = 11
]
\DndMonsterDetails[
skills = {Medicine +4, Nature +3, Perception +4},
languages = {Druidic plus any two languages},
challenge = 2,
]
\DndMonsterAction{Spellcasting}
The druid is a 4th-level spellcaster. Its spellcasting ability is Wisdom (spell save DC 12, +4 to hit with spell attacks). It has the following druid spells prepared:
\begin{DndMonsterSpells}
  \DndMonsterSpellLevel{druidcraft, produce flame, shillelagh}
   \DndMonsterSpellLevel[1][4]{entangle, longstrider, speak with animals, thunderwave}
   \DndMonsterSpellLevel[2][3]{animal messenger, barkskin}
\end{DndMonsterSpells}
\par
\DndMonsterSection{Actions}
\DndMonsterAction{Quarterstaff}
\textsl{Melee Weapon Attack:} +2 to hit (+4 to hit with shillelagh), reach 5 ft., one target. \textsl{Hit:} 3 (1d6) bludgeoning damage, 4 (1d8) bludgeoning damage if wielded with two hands, or 6 (1d8 + 2) bludgeoning damage with \textit{shillelagh}.

\end{DndMonster}



\begin{DndMonster}[float=t]{Hydra}
\DndMonsterType{Huge monstrosity, unaligned}
\DndMonsterBasics[
armor-class = {15 (natural armor)},
hit-points = {\DndDice{15d12 + 75}},
speed = {30 ft., swim 30 ft.}
]
\DndMonsterAbilityScores[
 str = 20,
 dex = 12,
 con = 20,
 int = 2,
 wis = 10,
 cha = 7
]
\DndMonsterDetails[
skills = {Perception +6},
senses = {darkvision 60 ft., passive perception 16},
challenge = 8,
]
\DndMonsterAction{Hold Breath}
The hydra can hold its breath for 1 hour.

\DndMonsterAction{Multiple Heads}
The hydra has five heads. While it has more than one head, the hydra has advantage on saving throws against being blinded, charmed, deafened, frightened, stunned, and knocked unconscious.

Whenever the hydra takes 25 or more damage in a single turn, one of its heads dies. If all its heads die, the hydra dies.

At the end of its turn, it grows two heads for each of its heads that died since its last turn, unless it has taken fire damage since its last turn. The hydra regains 10 hit points for each head regrown in this way.

\DndMonsterAction{Reactive Heads}
For each head the hydra has beyond one, it gets an extra reaction that can be used only for opportunity attacks.

\DndMonsterAction{Wakeful}
While the hydra sleeps, at least one of its heads is awake.

\par
\DndMonsterSection{Actions}
\DndMonsterAction{Multiattack}
The hydra makes as many bite attacks as it has heads.

\DndMonsterAction{Bite}
\textsl{Melee Weapon Attack:} +8 to hit, reach 10 ft., one target. \textsl{Hit:} 10 (1d10 + 5) piercing damage.

\end{DndMonster}



\begin{DndMonster}[float=t]{Duergar}
\DndMonsterType{Medium humanoid (dwarf), lawful evil}
\DndMonsterBasics[
armor-class = {16 (\textit{scale mail})},
hit-points = {\DndDice{4d8 + 4}},
speed = {25 ft.}
]
\DndMonsterAbilityScores[
 str = 14,
 dex = 11,
 con = 14,
 int = 11,
 wis = 10,
 cha = 9
]
\DndMonsterDetails[
damage-resistances = {poison},
senses = {darkvision 120 ft., passive perception 10},
languages = {Dwarvish, Undercommon},
challenge = 1,
]
\DndMonsterAction{Duergar Resilience}
The duergar has advantage on saving throws against poison, spells, and illusions, as well as to resist being charmed or paralyzed.

\DndMonsterAction{Sunlight Sensitivity}
While in sunlight, the duergar has disadvantage on attack rolls, as well as on Wisdom (Perception) checks that rely on sight.

\par
\DndMonsterSection{Actions}
\DndMonsterAction{Enlarge (Recharges after a Short or Long Rest)}
For 1 minute, the duergar magically increases in size, along with anything it is wearing or carrying. While enlarged, the duergar is Large, doubles its damage dice on Strength-based weapon attacks (included in the attacks), and makes Strength checks and Strength saving throws with advantage. If the duergar lacks the room to become Large, it attains the maximum size possible in the space available.

\DndMonsterAction{War Pick}
\textsl{Melee Weapon Attack:} +4 to hit, reach 5 ft., one target. \textsl{Hit:} 6 (1d8 + 2) piercing damage, or 11 (2d8 + 2) piercing damage while enlarged.

\DndMonsterAction{Javelin}
\textsl{Melee or Ranged Weapon Attack:} +4 to hit, reach 5 ft. or range 30/120 ft., one target. \textsl{Hit:} 5 (1d6 + 2) piercing damage, or 9 (2d6 + 2) piercing damage while enlarged.

\DndMonsterAction{Invisibility (Recharges after a Short or Long Rest)}
The duergar magically turns invisible until it attacks, casts a spell, or uses its Enlarge, or until its concentration is broken, up to 1 hour (as if concentration on a spell). Any equipment the duergar wears or carries is invisible with it.

\end{DndMonster}



\begin{DndMonster}[float=t]{Earth Elemental}
\DndMonsterType{Large elemental, neutral}
\DndMonsterBasics[
armor-class = {17 (natural armor)},
hit-points = {\DndDice{12d10 + 60}},
speed = {30 ft., burrow 30 ft.}
]
\DndMonsterAbilityScores[
 str = 20,
 dex = 8,
 con = 20,
 int = 5,
 wis = 10,
 cha = 5
]
\DndMonsterDetails[
damage-vulnerabilities = {thunder},
damage-resistances = {; bludgeoning, piercing, slashing from nonmagical attacks},
damage-immunities = {poison},
senses = {darkvision 60 ft., tremorsense 60 ft., passive perception 10},
languages = {Terran},
challenge = 5,
]
\DndMonsterAction{Earth Glide}
The elemental can burrow through nonmagical, unworked earth and stone. While doing so, the elemental doesn't disturb the material it moves through.

\DndMonsterAction{Siege Monster}
The elemental deals double damage to objects and structures.

\par
\DndMonsterSection{Actions}
\DndMonsterAction{Multiattack}
The elemental makes two slam attacks.

\DndMonsterAction{Slam}
\textsl{Melee Weapon Attack:} +8 to hit, reach 10 ft., one target. \textsl{Hit:} 14 (2d8 + 5) bludgeoning damage.

\end{DndMonster}



\begin{DndMonster}[float=t]{Ember Roc}
\DndMonsterType{Gargantuan monstrosity, U}
\DndMonsterBasics[
armor-class = {14 (natural armor)},
hit-points = {\DndDice{16d20 + 64}},
speed = {20 ft., fly 120 ft.}
]
\DndMonsterAbilityScores[
 str = 26,
 dex = 10,
 con = 18,
 int = 4,
 wis = 11,
 cha = 9
]
\DndMonsterDetails[
saving-throws = {Dex +5, Con +9, Wis +5, Cha +4},
skills = {Perception +5},
damage-immunities = {fire},
languages = {understands Ignan but can't speak},
challenge = 14,
]
\DndMonsterAction{Keen Sight}
The roc has Advantage and Disadvantage on Wisdom (Perception) checks that rely on sight.

\DndMonsterAction{Searing Presence}
Any creature that starts its turn within 5 feet of the roc takes 7 (2d6) fire damage.

\DndMonsterAction{Illumination}
The roc sheds bright light in a 40-foot radius and dim light for an additional 40 feet.

\par
\DndMonsterSection{Actions}
\DndMonsterAction{Multiattack}
The roc makes two attacks: one with its beak and one with its talons.

\DndMonsterAction{Beak}
\textsl{Melee Weapon Attack:} +13 to hit, reach 10 ft., one target. \textsl{Hit:} 26 (4d8 + 8) piercing damage and 7 (2d6) fire damage.

\DndMonsterAction{Talons}
\textsl{Melee Weapon Attack:} +13 to hit, reach 5 ft., one target. \textsl{Hit:} 22 (4d6 + 8) slashing damage and 7 (2d6) fire damage.

\DndMonsterAction{Inferno (Recharge 5\textendash 6)}
The roc beats its wings, creating a swirling conflagration. Each creature within 30 feet of the roc must succeed on a DC 18 Dexterity saving throws or take 21 (6d6) fire damage.

\end{DndMonster}



\begin{DndMonster}[float=t]{Ettin}
\DndMonsterType{Large giant, chaotic evil}
\DndMonsterBasics[
armor-class = {12 (natural armor)},
hit-points = {\DndDice{10d10 + 30}},
speed = {40 ft.}
]
\DndMonsterAbilityScores[
 str = 21,
 dex = 8,
 con = 17,
 int = 6,
 wis = 10,
 cha = 8
]
\DndMonsterDetails[
skills = {Perception +4},
senses = {darkvision 60 ft., passive perception 14},
languages = {Giant, Orc},
challenge = 4,
]
\DndMonsterAction{Two Heads}
The ettin has advantage on Wisdom (Perception) checks and on saving throws against being blinded, charmed, deafened, frightened, stunned, and knocked unconscious.

\DndMonsterAction{Wakeful}
When one of the ettin's heads is asleep, its other head is awake.

\par
\DndMonsterSection{Actions}
\DndMonsterAction{Multiattack}
The ettin makes two attacks: one with its battleaxe and one with its morningstar.

\DndMonsterAction{Battleaxe}
\textsl{Melee Weapon Attack:} +7 to hit, reach 5 ft., one target. \textsl{Hit:} 14 (2d8 + 5) slashing damage.

\DndMonsterAction{Morningstar}
\textsl{Melee Weapon Attack:} +7 to hit, reach 5 ft., one target. \textsl{Hit:} 14 (2d8 + 5) piercing damage.

\end{DndMonster}



\begin{DndMonster}[float=t]{Fire Giant}
\DndMonsterType{Huge giant, lawful evil}
\DndMonsterBasics[
armor-class = {18 (\textit{plate armor})},
hit-points = {\DndDice{13d12 + 78}},
speed = {30 ft.}
]
\DndMonsterAbilityScores[
 str = 25,
 dex = 9,
 con = 23,
 int = 10,
 wis = 14,
 cha = 13
]
\DndMonsterDetails[
saving-throws = {Dex +3, Con +10, Cha +5},
skills = {Athletics +11, Perception +6},
damage-immunities = {fire},
languages = {Giant},
challenge = 9,
]
\par
\DndMonsterSection{Actions}
\DndMonsterAction{Multiattack}
The giant makes two greatsword attacks.

\DndMonsterAction{Greatsword}
\textsl{Melee Weapon Attack:} +11 to hit, reach 10 ft., one target. \textsl{Hit:} 28 (6d6 + 7) slashing damage.

\DndMonsterAction{Rock}
\textsl{Ranged Weapon Attack:} +11 to hit, range 60/240 ft., one target. \textsl{Hit:} 29 (4d10 + 7) bludgeoning damage.

\end{DndMonster}



\begin{DndMonster}[float=t]{Quipper}
\DndMonsterType{Tiny beast, unaligned}
\DndMonsterBasics[
armor-class = {13},
hit-points = {\DndDice{1d4 - 1}},
speed = {0 ft., swim 40 ft.}
]
\DndMonsterAbilityScores[
 str = 2,
 dex = 16,
 con = 9,
 int = 1,
 wis = 7,
 cha = 2
]
\DndMonsterDetails[
senses = {darkvision 60 ft., passive perception 8},
challenge = 0,
]
\DndMonsterAction{Blood Frenzy}
The quipper has advantage on melee attack rolls against any creature that doesn't have all its hit points.

\DndMonsterAction{Water Breathing}
The quipper can breathe only underwater.

\par
\DndMonsterSection{Actions}
\DndMonsterAction{Bite}
\textsl{Melee Weapon Attack:} +5 to hit, reach 5 ft., one target. \textsl{Hit:} 1 piercing damage.

\end{DndMonster}



\begin{DndMonster}[float=t]{Flaming Skeleton}
\DndMonsterType{Medium undead, lawful evil}
\DndMonsterBasics[
armor-class = {12},
hit-points = {\DndDice{3d8 + 6}},
speed = {30 ft.}
]
\DndMonsterAbilityScores[
 str = 10,
 dex = 14,
 con = 15,
 int = 6,
 wis = 8,
 cha = 5
]
\DndMonsterDetails[
damage-vulnerabilities = {bludgeoning},
damage-immunities = {fire, poison},
senses = {darkvision 60 ft., passive perception 9},
languages = {understands all languages it knew in life but can't speak},
challenge = 1/2,
]
\DndMonsterAction{Heated Weapons}
When the skeleton hits with a metal melee weapon, the weapon deals an extra 3 (1d6) fire damage (included in the attack).

\DndMonsterAction{Illumination}
The skeleton sheds bright light in a 10-foot radius and dim light for an additional 10 feet.

\par
\DndMonsterSection{Actions}
\DndMonsterAction{Shortsword}
\textsl{Melee Weapon Attack:} +4 to hit, reach 5 ft., one target. \textsl{Hit:} 5 (1d6 + 2) piercing damage and 3 (1d6) fire damage.

\DndMonsterAction{Ignite}
The skeleton flings itself at a Large or smaller creature within 5 feet of it. The target must succeed on a DC 12 Dexterity saving throws or be grappled (escape DC 12). On a failed save, the target also takes 3 (1d6) fire damage at the start of each of its turns, even after the grapple ends. If the target or a creature within 5 feet of it uses an action to put out the flames, or if some other effect douses the flames (such as the target being submerged in water), the effect ends.

\end{DndMonster}



\begin{DndMonster}[float=t]{Flying Snake}
\DndMonsterType{Tiny beast, unaligned}
\DndMonsterBasics[
armor-class = {14},
hit-points = {\DndDice{2d4}},
speed = {30 ft., fly 60 ft., swim 30 ft.}
]
\DndMonsterAbilityScores[
 str = 4,
 dex = 18,
 con = 11,
 int = 2,
 wis = 12,
 cha = 5
]
\DndMonsterDetails[
senses = {blindsight 10 ft., passive perception 11},
challenge = 1/8,
]
\DndMonsterAction{Flyby}
The snake doesn't provoke opportunity attacks when it flies out of an enemy's reach.

\par
\DndMonsterSection{Actions}
\DndMonsterAction{Bite}
\textsl{Melee Weapon Attack:} +6 to hit, reach 5 ft., one target. \textsl{Hit:} 1 piercing damage plus 7 (3d4) poison damage.

\end{DndMonster}



\begin{DndMonster}[float*=b,width=\textwidth + 8pt]{Forge Guardian}
\begin{multicols}{2}
\DndMonsterType{Gargantuan construct, U}
\DndMonsterBasics[
armor-class = {24 (natural armor)},
hit-points = {\DndDice{24d20 + 168}},
speed = {50 ft.}
]
\DndMonsterAbilityScores[
 str = 30,
 dex = 10,
 con = 24,
 int = 3,
 wis = 11,
 cha = 1
]
\DndMonsterDetails[
saving-throws = {Str +18, Con +15},
damage-immunities = {fire, poison, psychic; bludgeoning, piercing, slashing from nonmagical weapons that aren't adamantine},
senses = {truesight 120 ft., passive perception 10},
challenge = 27,
]
\DndMonsterAction{Fire Absorption}
Whenever the guardian is subjected to fire damage, it takes no damage and instead regains a number of hit points equal to the fire damage dealt.

\DndMonsterAction{Immutable Form}
The guardian is immune to any spell or effect that would alter its form.

\DndMonsterAction{Magic Resistance}
The guardian has Advantage and Disadvantage on saving throws against spells and other magical effects.

\DndMonsterAction{Magic Weapons}
The guardian's weapon attacks are magical.

\par
\DndMonsterSection{Actions}
\DndMonsterAction{Multiattack}
The guardian makes two melee attacks.

\DndMonsterAction{Runeblade}
\textsl{Melee Weapon Attack:} +18 to hit, reach 15 ft., one target. \textsl{Hit:} 36 (4d12 + 10) slashing damage, and the target must succeed on a DC 23 Constitution saving throws or be blinded until it deals damage to the guardian.

\DndMonsterAction{Kick}
\textsl{Melee Weapon Attack:} +18 to hit, reach 10 ft., one target. \textsl{Hit:} 23 (2d12 + 10) bludgeoning damage, and the target is pushed up to 15 feet away from the guardian and knocked prone.

\DndMonsterAction{Bladestorm (Recharge 5\textendash 6)}
The guardian makes a Runeblade attack against every creature within 15 feet of it. On a hit, a target is also pushed up to 15 feet away from the guardian and knocked prone.

\par
\DndMonsterSection{Legendary Actions}
Forge Guardian can take 3 legendary actions, choosing from the options below. Only one legendary action can be used at a time and only at the end of another creature's turn. Forge Guardian regains spent legendary actions at the start of its turn.

\begin{DndMonsterLegendaryActions}
\DndMonsterLegendaryAction{Runeblade}{The guardian makes one Runeblade attack.
}
\DndMonsterLegendaryAction{Kick}{The guardian makes one Kick attack.
}
\DndMonsterLegendaryAction{Trample (Costs 2 Actions)}{The guardian moves up to 50 feet without triggering opportunity attack, and can pass through other creature's spaces. Each creature whose space the guardian passes through must make a DC 23 Dexterity saving throws. On a failure, the creature takes 23 (2d12 + 10) bludgeoning damage and is knocked prone. On a success, the creature takes half as much damage and isn't knocked prone. Any creature that shares the guardian's space when it stops this movement is pushed out of its space into the nearest unoccupied space.
}
\end{DndMonsterLegendaryActions}
\end{multicols}
\end{DndMonster}



\begin{DndMonster}[float=t]{Frost Giant}
\DndMonsterType{Huge giant, neutral evil}
\DndMonsterBasics[
armor-class = {15 (patchwork armor)},
hit-points = {\DndDice{12d12 + 60}},
speed = {40 ft.}
]
\DndMonsterAbilityScores[
 str = 23,
 dex = 9,
 con = 21,
 int = 9,
 wis = 10,
 cha = 12
]
\DndMonsterDetails[
saving-throws = {Con +8, Wis +3, Cha +4},
skills = {Athletics +9, Perception +3},
damage-immunities = {cold},
languages = {Giant},
challenge = 8,
]
\par
\DndMonsterSection{Actions}
\DndMonsterAction{Multiattack}
The giant makes two greataxe attacks.

\DndMonsterAction{Greataxe}
\textsl{Melee Weapon Attack:} +9 to hit, reach 10 ft., one target. \textsl{Hit:} 25 (3d12 + 6) slashing damage.

\DndMonsterAction{Rock}
\textsl{Ranged Weapon Attack:} +9 to hit, range 60/240 ft., one target. \textsl{Hit:} 28 (4d10 + 6) bludgeoning damage.

\end{DndMonster}



\begin{DndMonster}[float=t]{Ghast}
\DndMonsterType{Medium undead, chaotic evil}
\DndMonsterBasics[
armor-class = {13},
hit-points = {\DndDice{8d8}},
speed = {30 ft.}
]
\DndMonsterAbilityScores[
 str = 16,
 dex = 17,
 con = 10,
 int = 11,
 wis = 10,
 cha = 8
]
\DndMonsterDetails[
damage-resistances = {necrotic},
damage-immunities = {poison},
senses = {darkvision 60 ft., passive perception 10},
languages = {Common},
challenge = 2,
]
\DndMonsterAction{Stench}
Any creature that starts its turn within 5 feet of the ghast must succeed on a DC 10 Constitution saving throw or be poisoned until the start of its next turn. On a successful saving throw, the creature is immune to the ghast's Stench for 24 hours.

\DndMonsterAction{Turn Defiance}
The ghast and any ghouls within 30 feet of it have advantage on saving throws against effects that turn undead.

\par
\DndMonsterSection{Actions}
\DndMonsterAction{Bite}
\textsl{Melee Weapon Attack:} +3 to hit, reach 5 ft., one creature. \textsl{Hit:} 12 (2d8 + 3) piercing damage.

\DndMonsterAction{Claws}
\textsl{Melee Weapon Attack:} +5 to hit, reach 5 ft., one target. \textsl{Hit:} 10 (2d6 + 3) slashing damage. If the target is a creature other than an undead, it must succeed on a DC 10 Constitution saving throw or be paralyzed for 1 minute. The target can repeat the saving throw at the end of each of its turns, ending the effect on itself on a success.

\end{DndMonster}



\begin{DndMonster}[float=t]{Ghost}
\DndMonsterType{Medium undead, any alignment}
\DndMonsterBasics[
armor-class = {11},
hit-points = {\DndDice{10d8}},
speed = {0 ft., fly 40 ft. \textit{(hover)}}
]
\DndMonsterAbilityScores[
 str = 7,
 dex = 13,
 con = 10,
 int = 10,
 wis = 12,
 cha = 17
]
\DndMonsterDetails[
damage-resistances = {acid, fire, lightning, thunder; bludgeoning, piercing, slashing from nonmagical attacks},
damage-immunities = {cold, necrotic, poison},
senses = {darkvision 60 ft., passive perception 11},
languages = {any languages it knew in life},
challenge = 4,
]
\DndMonsterAction{Ethereal Sight}
The ghost can see 60 feet into the Ethereal Plane when it is on the Material Plane, and vice versa.

\DndMonsterAction{Incorporeal Movement}
The ghost can move through other creatures and objects as if they were difficult terrain. It takes 5 (1d10) force damage if it ends its turn inside an object.

\par
\DndMonsterSection{Actions}
\DndMonsterAction{Withering Touch}
\textsl{Melee Weapon Attack:} +5 to hit, reach 5 ft., one target. \textsl{Hit:} 17 (4d6 + 3) necrotic damage.

\DndMonsterAction{Etherealness}
The ghost enters the Ethereal Plane from the Material Plane, or vice versa. It is visible on the Material Plane while it is in the Border Ethereal, and vice versa, yet it can't affect or be affected by anything on the other plane.

\DndMonsterAction{Horrifying Visage}
Each non-undead creature within 60 feet of the ghost that can see it must succeed on a DC 13 Wisdom saving throw or be frightened for 1 minute. If the save fails by 5 or more, the target also ages 1d4 × 10 years. A frightened target can repeat the saving throw at the end of each of its turns, ending the frightened condition on itself on a success. If a target's saving throw is successful or the effect ends for it, the target is immune to this ghost's Horrifying Visage for the next 24 hours. The aging effect can be reversed with a  \textit{greater restoration} spell, but only within 24 hours of it occurring.

\DndMonsterAction{Possession (Recharge 6)}
One humanoid that the ghost can see within 5 feet of it must succeed on a DC 13 Charisma saving throw or be possessed by the ghost; the ghost then disappears, and the target is incapacitated and loses control of its body. The ghost now controls the body but doesn't deprive the target of awareness. The ghost can't be targeted by any attack, spell, or other effect, except ones that turn undead, and it retains its alignment, Intelligence, Wisdom, Charisma, and immunity to being charmed and frightened. It otherwise uses the possessed target's statistics, but doesn't gain access to the target's knowledge, class features, or proficiencies.

The possession lasts until the body drops to 0 hit points, the ghost ends it as a bonus action, or the ghost is turned or forced out by an effect like the \textit{dispel evil and good} spell. When the possession ends, the ghost reappears in an unoccupied space within 5 feet of the body. The target is immune to this ghost's Possession for 24 hours after succeeding on the saving throw or after the possession ends.

\end{DndMonster}



\begin{DndMonster}[float=t]{Ghoul}
\DndMonsterType{Medium undead, chaotic evil}
\DndMonsterBasics[
armor-class = {12},
hit-points = {\DndDice{5d8}},
speed = {30 ft.}
]
\DndMonsterAbilityScores[
 str = 13,
 dex = 15,
 con = 10,
 int = 7,
 wis = 10,
 cha = 6
]
\DndMonsterDetails[
damage-immunities = {poison},
senses = {darkvision 60 ft., passive perception 10},
languages = {Common},
challenge = 1,
]
\par
\DndMonsterSection{Actions}
\DndMonsterAction{Bite}
\textsl{Melee Weapon Attack:} +2 to hit, reach 5 ft., one creature. \textsl{Hit:} 9 (2d6 + 2) piercing damage.

\DndMonsterAction{Claws}
\textsl{Melee Weapon Attack:} +4 to hit, reach 5 ft., one target. \textsl{Hit:} 7 (2d4 + 2) slashing damage. If the target is a creature other than an elf or undead, it must succeed on a DC 10 Constitution saving throw or be paralyzed for 1 minute. The target can repeat the saving throw at the end of each of its turns, ending the effect on itself on a success.

\end{DndMonster}



\begin{DndMonster}[float=t]{Giant Crocodile}
\DndMonsterType{Huge beast, unaligned}
\DndMonsterBasics[
armor-class = {14 (natural armor)},
hit-points = {\DndDice{9d12 + 27}},
speed = {30 ft., swim 50 ft.}
]
\DndMonsterAbilityScores[
 str = 21,
 dex = 9,
 con = 17,
 int = 2,
 wis = 10,
 cha = 7
]
\DndMonsterDetails[
skills = {Stealth +5},
challenge = 5,
]
\DndMonsterAction{Hold Breath}
The crocodile can hold its breath for 30 minutes.

\par
\DndMonsterSection{Actions}
\DndMonsterAction{Multiattack}
The crocodile makes two attacks: one with its bite and one with its tail.

\DndMonsterAction{Bite}
\textsl{Melee Weapon Attack:} +8 to hit, reach 5 ft., one target. \textsl{Hit:} 21 (3d10 + 5) piercing damage, and the target is grappled (escape DC 16). Until this grapple ends, the target is restrained, and the crocodile can't bite another target.

\DndMonsterAction{Tail}
\textsl{Melee Weapon Attack:} +8 to hit, reach 10 ft., one target not grappled by the crocodile. \textsl{Hit:} 14 (2d8 + 5) bludgeoning damage. If the target is a creature, it must succeed on a DC 16 Strength saving throw or be knocked prone.

\end{DndMonster}



\begin{DndMonster}[float=t]{Giant Eagle}
\DndMonsterType{Large beast, neutral good}
\DndMonsterBasics[
armor-class = {13},
hit-points = {\DndDice{4d10 + 4}},
speed = {10 ft., fly 80 ft.}
]
\DndMonsterAbilityScores[
 str = 16,
 dex = 17,
 con = 13,
 int = 8,
 wis = 14,
 cha = 10
]
\DndMonsterDetails[
skills = {Perception +4},
languages = {Giant Eagle, understands Common and Auran but can't speak them},
challenge = 1,
]
\DndMonsterAction{Keen Sight}
The eagle has advantage on Wisdom (Perception) checks that rely on sight.

\par
\DndMonsterSection{Actions}
\DndMonsterAction{Multiattack}
The eagle makes two attacks: one with its beak and one with its talons.

\DndMonsterAction{Beak}
\textsl{Melee Weapon Attack:} +5 to hit, reach 5 ft., one target. \textsl{Hit:} 6 (1d6 + 3) piercing damage.

\DndMonsterAction{Talons}
\textsl{Melee Weapon Attack:} +5 to hit, reach 5 ft., one target. \textsl{Hit:} 10 (2d6 + 3) slashing damage.

\end{DndMonster}



\begin{DndMonster}[float=t]{Giant Frog}
\DndMonsterType{Medium beast, unaligned}
\DndMonsterBasics[
armor-class = {11},
hit-points = {\DndDice{4d8}},
speed = {30 ft., swim 30 ft.}
]
\DndMonsterAbilityScores[
 str = 12,
 dex = 13,
 con = 11,
 int = 2,
 wis = 10,
 cha = 3
]
\DndMonsterDetails[
skills = {Perception +2, Stealth +3},
senses = {darkvision 30 ft., passive perception 12},
challenge = 1/4,
]
\DndMonsterAction{Amphibious}
The frog can breathe air and water.

\DndMonsterAction{Standing Leap}
The frog's long jump is up to 20 feet and its high jump is up to 10 feet, with or without a running start.

\par
\DndMonsterSection{Actions}
\DndMonsterAction{Bite}
\textsl{Melee Weapon Attack:} +3 to hit, reach 5 ft., one target. \textsl{Hit:} 4 (1d6 + 1) piercing damage, and the target is grappled (escape DC 11). Until this grapple ends, the target is restrained, and the frog can't bite another target.

\DndMonsterAction{Swallow}
The frog makes one bite attack against a Small or smaller target it is grappling. If the attack hits, the target is swallowed, and the grapple ends. The swallowed target is blinded and restrained, it has total cover against attacks and other effects outside the frog, and it takes 5 (2d4) acid damage at the start of each of the frog's turns. The frog can have only one target swallowed at a time. If the frog dies, a swallowed creature is no longer restrained by it and can escape from the corpse using 5 feet of movement, exiting prone.

\end{DndMonster}


\clearpage

\begin{DndMonster}[float=t]{Giant Goat}
\DndMonsterType{Large beast, unaligned}
\DndMonsterBasics[
armor-class = {11 (natural armor)},
hit-points = {\DndDice{3d10 + 3}},
speed = {40 ft.}
]
\DndMonsterAbilityScores[
 str = 17,
 dex = 11,
 con = 12,
 int = 3,
 wis = 12,
 cha = 6
]
\DndMonsterDetails[
challenge = 1/2,
]
\DndMonsterAction{Charge}
If the goat moves at least 20 feet straight toward a target and then hits it with a ram attack on the same turn, the target takes an extra 5 (2d4) bludgeoning damage. If the target is a creature, it must succeed on a DC 13 Strength saving throw or be knocked prone.

\DndMonsterAction{Sure-Footed}
The goat has advantage on Strength and Dexterity saving throws made against effects that would knock it prone.

\par
\DndMonsterSection{Actions}
\DndMonsterAction{Ram}
\textsl{Melee Weapon Attack:} +5 to hit, reach 5 ft., one target. \textsl{Hit:} 8 (2d4 + 3) bludgeoning damage.

\end{DndMonster}



\begin{DndMonster}[float=t]{Giant Lizard}
\DndMonsterType{Large beast, unaligned}
\DndMonsterBasics[
armor-class = {12 (natural armor)},
hit-points = {\DndDice{3d10 + 3}},
speed = {30 ft., climb 30 ft.}
]
\DndMonsterAbilityScores[
 str = 15,
 dex = 12,
 con = 13,
 int = 2,
 wis = 10,
 cha = 5
]
\DndMonsterDetails[
senses = {darkvision 30 ft., passive perception 10},
challenge = 1/4,
]
\par
\DndMonsterSection{Actions}
\DndMonsterAction{Bite}
\textsl{Melee Weapon Attack:} +4 to hit, reach 5 ft., one target. \textsl{Hit:} 6 (1d8 + 2) piercing damage.

\end{DndMonster}



\begin{DndMonster}[float=t]{Giant Spider}
\DndMonsterType{Large beast, unaligned}
\DndMonsterBasics[
armor-class = {14 (natural armor)},
hit-points = {\DndDice{4d10 + 4}},
speed = {30 ft., climb 30 ft.}
]
\DndMonsterAbilityScores[
 str = 14,
 dex = 16,
 con = 12,
 int = 2,
 wis = 11,
 cha = 4
]
\DndMonsterDetails[
skills = {Stealth +7},
senses = {blindsight 10 ft., darkvision 60 ft., passive perception 10},
challenge = 1,
]
\DndMonsterAction{Spider Climb}
The spider can climb difficult surfaces, including upside down on ceilings, without needing to make an ability check.

\DndMonsterAction{Web Sense}
While in contact with a web, the spider knows the exact location of any other creature in contact with the same web.

\DndMonsterAction{Web Walker}
The spider ignores movement restrictions caused by webbing.

\par
\DndMonsterSection{Actions}
\DndMonsterAction{Bite}
\textsl{Melee Weapon Attack:} +5 to hit, reach 5 ft., one creature. \textsl{Hit:} 7 (1d8 + 3) piercing damage, and the target must make a DC 11 Constitution saving throw, taking 9 (2d8) poison damage on a failed save, or half as much damage on a successful one. If the poison damage reduces the target to 0 hit points, the target is stable but poisoned for 1 hour, even after regaining hit points, and is paralyzed while poisoned in this way.

\DndMonsterAction{Web (Recharge 5\textendash 6)}
\textsl{Ranged Weapon Attack:} +5 to hit, range 30/60 ft., one creature. \textsl{Hit:} The target is restrained by webbing. As an action, the restrained target can make a DC 12 Strength check, bursting the webbing on a success. The webbing can also be attacked and destroyed (AC 10; hp 5; vulnerability to fire damage; immunity to bludgeoning, poison, and psychic damage).

\end{DndMonster}



\begin{DndMonster}[float=t]{Gloomstalker}
\DndMonsterType{Large monstrosity, neutral evil}
\DndMonsterBasics[
armor-class = {15 (natural armor)},
hit-points = {\DndDice{12d10 + 24}},
speed = {40 ft., fly 80 ft.}
]
\DndMonsterAbilityScores[
 str = 22,
 dex = 16,
 con = 14,
 int = 5,
 wis = 17,
 cha = 14
]
\DndMonsterDetails[
saving-throws = {Str +9, Dex +6},
skills = {Athletics +9, Intimidation +5, Perception +6, Stealth +6},
damage-vulnerabilities = {radiant},
senses = {darkvision 240 ft, passive perception 16},
languages = {understands Common but can't speak},
challenge = 6,
]
\DndMonsterAction{Shadowstep}
As a bonus action, the gloomstalker can teleport up to 40 feet to an unoccupied space it can see.

\DndMonsterAction{Sunlight Sensitivity}
While in sunlight, the gloomstalker has disadvantage on attack rolls, as well as on Wisdom (Perception) checks that rely on sight.

\par
\DndMonsterSection{Actions}
\DndMonsterAction{Multiattack}
The gloomstalker makes two attacks: one with its bite and one with its claws.

\DndMonsterAction{Bite}
\textsl{Melee Weapon Attack:} +9 to hit, reach 5 ft., one creature. \textsl{Hit:} 15 (2d8 + 6) piercing damage plus 7 (2d6) necrotic damage.

\DndMonsterAction{Claws}
\textsl{Melee Weapon Attack:} +9 to hit, reach 5 ft., one target. \textsl{Hit:} 13 (2d6 + 6) slashing damage plus 7 (2d6) necrotic damage.

\DndMonsterAction{Snatch}
\textsl{Melee Weapon Attack:} +9 to hit, reach 5 ft., one Medium or smaller creature. \textsl{Hit:} 13 (2d6 + 6) slashing damage plus 7 (2d6) necrotic damage, and the target is grappled (escape DC 17). While grappled in this way, the target is restrained.

\DndMonsterAction{Shriek (Recharge 6)}
The gloomstalker emits a terrible shriek. Each enemy within 60 feet of the gloomstalker that can hear it must succeed on a DC 13 Constitution saving throw or be paralyzed until the end of the enemy's next turn.

\end{DndMonster}



\begin{DndMonster}[float=t]{Gnoll}
\DndMonsterType{Medium humanoid (gnoll), chaotic evil}
\DndMonsterBasics[
armor-class = {15 (\textit{hide armor})},
hit-points = {\DndDice{5d8}},
speed = {30 ft.}
]
\DndMonsterAbilityScores[
 str = 14,
 dex = 12,
 con = 11,
 int = 6,
 wis = 10,
 cha = 7
]
\DndMonsterDetails[
senses = {darkvision 60 ft., passive perception 10},
languages = {Gnoll},
challenge = 1/2,
]
\DndMonsterAction{Rampage}
When the gnoll reduces a creature to 0 hit points with a melee attack on its turn, the gnoll can take a bonus action to move up to half its speed and make a bite attack.

\par
\DndMonsterSection{Actions}
\DndMonsterAction{Bite}
\textsl{Melee Weapon Attack:} +4 to hit, reach 5 ft., one creature. \textsl{Hit:} 4 (1d4 + 2) piercing damage.

\DndMonsterAction{Spear}
\textsl{Melee or Ranged Weapon Attack:} +4 to hit, reach 5 ft. or range 20/60 ft., one target. \textsl{Hit:} 5 (1d6 + 2) piercing damage, or 6 (1d8 + 2) piercing damage if used with two hands to make a melee attack.

\DndMonsterAction{Longbow}
\textsl{Ranged Weapon Attack:} +3 to hit, range 150/600 ft., one target. \textsl{Hit:} 5 (1d8 + 1) piercing damage.

\end{DndMonster}



\begin{DndMonster}[float=t]{Giant Goat}
\DndMonsterType{Large beast, unaligned}
\DndMonsterBasics[
armor-class = {11 (natural armor)},
hit-points = {\DndDice{3d10 + 3}},
speed = {40 ft.}
]
\DndMonsterAbilityScores[
 str = 17,
 dex = 11,
 con = 12,
 int = 3,
 wis = 12,
 cha = 6
]
\DndMonsterDetails[
challenge = 1/2,
]
\DndMonsterAction{Charge}
If the goat moves at least 20 feet straight toward a target and then hits it with a ram attack on the same turn, the target takes an extra 5 (2d4) bludgeoning damage. If the target is a creature, it must succeed on a DC 13 Strength saving throw or be knocked prone.

\DndMonsterAction{Sure-Footed}
The goat has advantage on Strength and Dexterity saving throws made against effects that would knock it prone.

\par
\DndMonsterSection{Actions}
\DndMonsterAction{Ram}
\textsl{Melee Weapon Attack:} +5 to hit, reach 5 ft., one target. \textsl{Hit:} 8 (2d4 + 3) bludgeoning damage.

\end{DndMonster}



\begin{DndMonster}[float=t]{Goblin}
\DndMonsterType{Small humanoid (goblinoid), neutral evil}
\DndMonsterBasics[
armor-class = {15 (\textit{leather armor})},
hit-points = {\DndDice{2d6}},
speed = {30 ft.}
]
\DndMonsterAbilityScores[
 str = 8,
 dex = 14,
 con = 10,
 int = 10,
 wis = 8,
 cha = 8
]
\DndMonsterDetails[
skills = {Stealth +6},
senses = {darkvision 60 ft., passive perception 9},
languages = {Common, Goblin},
challenge = 1/4,
]
\DndMonsterAction{Nimble Escape}
The goblin can take the Disengage or Hide action as a bonus action on each of its turns.

\par
\DndMonsterSection{Actions}
\DndMonsterAction{Scimitar}
\textsl{Melee Weapon Attack:} +4 to hit, reach 5 ft., one target. \textsl{Hit:} 5 (1d6 + 2) slashing damage.

\DndMonsterAction{Shortbow}
\textsl{Ranged Weapon Attack:} +4 to hit, range 80/320 ft., one target. \textsl{Hit:} 5 (1d6 + 2) piercing damage.

\end{DndMonster}



\begin{DndMonster}[float=t]{Gorgon}
\DndMonsterType{Large monstrosity, unaligned}
\DndMonsterBasics[
armor-class = {19 (natural armor)},
hit-points = {\DndDice{12d10 + 48}},
speed = {40 ft.}
]
\DndMonsterAbilityScores[
 str = 20,
 dex = 11,
 con = 18,
 int = 2,
 wis = 12,
 cha = 7
]
\DndMonsterDetails[
skills = {Perception +4},
senses = {darkvision 60 ft., passive perception 14},
challenge = 5,
]
\DndMonsterAction{Trampling Charge}
If the gorgon moves at least 20 feet straight toward a creature and then hits it with a gore attack on the same turn, that target must succeed on a DC 16 Strength saving throw or be knocked prone. If the target is prone, the gorgon can make one attack with its hooves against it as a bonus action.

\par
\DndMonsterSection{Actions}
\DndMonsterAction{Gore}
\textsl{Melee Weapon Attack:} +8 to hit, reach 5 ft., one target. \textsl{Hit:} 18 (2d12 + 5) piercing damage.

\DndMonsterAction{Hooves}
\textsl{Melee Weapon Attack:} +8 to hit, reach 5 ft., one target. \textsl{Hit:} 16 (2d10 + 5) bludgeoning damage.

\DndMonsterAction{Petrifying Breath (Recharge 5\textendash 6)}
The gorgon exhales petrifying gas in a 30-foot cone. Each creature in that area must succeed on a DC 13 Constitution saving throw. On a failed save, a target begins to turn to stone and is restrained. The restrained target must repeat the saving throw at the end of its next turn. On a success, the effect ends on the target. On a failure, the target is petrified until freed by the  \textit{greater restoration} spell or other magic.

\end{DndMonster}



\begin{DndMonster}[float=t]{Griffon}
\DndMonsterType{Large monstrosity, unaligned}
\DndMonsterBasics[
armor-class = {12},
hit-points = {\DndDice{7d10 + 21}},
speed = {30 ft., fly 80 ft.}
]
\DndMonsterAbilityScores[
 str = 18,
 dex = 15,
 con = 16,
 int = 2,
 wis = 13,
 cha = 8
]
\DndMonsterDetails[
skills = {Perception +5},
senses = {darkvision 60 ft., passive perception 15},
challenge = 2,
]
\DndMonsterAction{Keen Sight}
The griffon has advantage on Wisdom (Perception) checks that rely on sight.

\par
\DndMonsterSection{Actions}
\DndMonsterAction{Multiattack}
The griffon makes two attacks: one with its beak and one with its claws.

\DndMonsterAction{Beak}
\textsl{Melee Weapon Attack:} +6 to hit, reach 5 ft., one target. \textsl{Hit:} 8 (1d8 + 4) piercing damage.

\DndMonsterAction{Claws}
\textsl{Melee Weapon Attack:} +6 to hit, reach 5 ft., one target. \textsl{Hit:} 11 (2d6 + 4) slashing damage.

\end{DndMonster}



\begin{DndMonster}[float=t]{Grog Strongjaw}
\DndMonsterType{Medium humanoid (half-giant), U}
\DndMonsterBasics[
armor-class = {19},
hit-points = {\DndDice{20d12 + 120 + 40}},
speed = {50 ft.}
]
\DndMonsterAbilityScores[
 str = 26,
 dex = 15,
 con = 22,
 int = 8,
 wis = 10,
 cha = 13
]
\DndMonsterDetails[
saving-throws = {Str +14, Con +12},
skills = {Animal Handling +6, Athletics +14, Intimidation +7, Survival +6},
damage-resistances = {; cold, fire, lightning, thunder while enlarged; bludgeoning, piercing, slashing while raging},
senses = {darkvision 60 ft., passive perception 10},
languages = {Common, Dwarvish, Giant, with limited reading and writing},
challenge = 18,
]
\DndMonsterAction{Brutal Critical}
When Grog scores a critical hit, he rolls double weapon damage dice as usual, then rolls 3 additional weapon damage dice.

\DndMonsterAction{Feral Instinct}
Grog has Advantage and Disadvantage on initiative rolls.

\DndMonsterAction{Melee Master}
Grog can take a −5 penalty to any melee attack roll to gain a +10 bonus to the damage roll.

\DndMonsterAction{Reckless}
At the start of his turn, Grog can gain Advantage and Disadvantage on all melee weapon attack rolls during that turn, but attack rolls against him have Advantage and Disadvantage until the start of his next turn.

\DndMonsterAction{Relentless Rage}
If Grog is reduced to 0 hit points while raging, he drops to 1 hit point instead if he succeeds on a DC 10 Constitution saving throws.

\DndMonsterAction{Special Equipment}
Grog wears a \textit{belt of dwarvenkind}. He wields the \textit{Bloodaxe} (a +2 magic weapon), a \textit{dwarven thrower}, and the Titanstone Knuckles.

\DndMonsterAction{Spellbreaker}
When Grog damages a creature that is concentrating on a spell, that creature has Advantage and Disadvantage on the saving throws to maintain its concentration. In addition, Grog has Advantage and Disadvantage on saving throws against spells cast by creatures within 5 feet of him.

\DndMonsterAction{Toughness}
Grog has an additional 40 hit points.

\par
\DndMonsterSection{Actions}
\DndMonsterAction{Multiattack}
Grog makes two attacks (three when raging) with his \textit{Bloodaxe}, \textit{dwarven thrower}, or Titanstone Knuckles.

\DndMonsterAction{\textit{Bloodaxe}}
\textsl{Melee Weapon Attack:} +16 to hit, reach 5 ft., one target. \textsl{Hit:} 20 (1d12 + 14) slashing damage and 3 (1d6) necrotic damage. If the target is reduced to 0 hit points, Grog gains 10 temporary hit points.

\DndMonsterAction{\textit{Dwarven Thrower}}
\textsl{Melee or Ranged Weapon Attack:} +17 to hit, reach 5 ft. or range 20/60, one target. \textsl{Hit:} 24 (2d8 + 15) bludgeoning damage.

\DndMonsterAction{Titanstone Knuckles}
\textsl{Melee Weapon Attack:} +14 to hit, reach 5 ft., one target. \textsl{Hit:} 14 (1d4 + 12) bludgeoning damage.

\DndMonsterAction{Enlarge (Recharges after a resting)}
Grog casts \textit{enlarge/reduce} on himself to grow in size for 10 minutes. While enlarged in this way, he has resistance to cold, fire, lightning, and thunder damage.

\par
\DndMonsterSection{Bonus Actions}
\DndMonsterAction{Rage (Recharges after a resting)}
Grog enters a rage that lasts until he falls unconscious or chooses to end it (no action required). While raging, Grog gains the following benefits:


\begin{itemize}
\item He has Advantage and Disadvantage on Strength checks and Strength saving throws.
\item He deals an extra 4 damage when he hits a target with a melee weapon attack (included in attacks)
\item He has resistance to bludgeoning, piercing, and slashing damage
\item He can't be charmed or frightened
\item He can make an additional melee weapon attack (included in the Multiattack action)
\end{itemize}

When Grog's rage ends, he gains one level of exhaustion.

\par
\DndMonsterSection{Reactions}
\DndMonsterAction{Opportunist}
When a creature within 5 feet of Grog casts a spell or deals damage to him, Grog can make a melee weapon attack against that creature.

\end{DndMonster}



\begin{DndMonster}[float=t]{Guardian Naga}
\DndMonsterType{Large monstrosity, lawful good}
\DndMonsterBasics[
armor-class = {18 (natural armor)},
hit-points = {\DndDice{15d10 + 45}},
speed = {40 ft.}
]
\DndMonsterAbilityScores[
 str = 19,
 dex = 18,
 con = 16,
 int = 16,
 wis = 19,
 cha = 18
]
\DndMonsterDetails[
saving-throws = {Dex +8, Con +7, Int +7, Wis +8, Cha +8},
damage-immunities = {poison},
senses = {darkvision 60 ft., passive perception 14},
languages = {Celestial, Common},
challenge = 10,
]
\DndMonsterAction{Rejuvenation}
If it dies, the naga returns to life in 1d6 days and regains all its hit points. Only a \textit{wish} spell can prevent this trait from functioning.

\DndMonsterAction{Spellcasting}
The naga is an 11th-level spellcaster. Its spellcasting ability is Wisdom (spell save DC 16, +8 to hit with spell attacks), and it needs only verbal components to cast its spells. It has the following cleric spells prepared:
\begin{DndMonsterSpells}
  \DndMonsterSpellLevel{mending, sacred flame, thaumaturgy}
   \DndMonsterSpellLevel[1][4]{command, cure wounds, shield of faith}
   \DndMonsterSpellLevel[2][3]{calm emotions, hold person}
   \DndMonsterSpellLevel[3][3]{bestow curse, clairvoyance}
   \DndMonsterSpellLevel[4][3]{banishment, freedom of movement}
   \DndMonsterSpellLevel[5][2]{flame strike, geas}
   \DndMonsterSpellLevel[6][1]{true seeing}
\end{DndMonsterSpells}
\par
\DndMonsterSection{Actions}
\DndMonsterAction{Bite}
\textsl{Melee Weapon Attack:} +8 to hit, reach 10 ft., one creature. \textsl{Hit:} 8 (1d8 + 4) piercing damage, and the target must make a DC 15 Constitution saving throw, taking 45 (10d8) poison damage on a failed save, or half as much damage on a successful one.

\DndMonsterAction{Spit Poison}
\textsl{Ranged Weapon Attack:} +8 to hit, range 15/30 ft., one creature. \textsl{Hit:} The target must make a DC 15 Constitution saving throw, taking 45 (10d8) poison damage on a failed save, or half as much damage on a successful one.

\end{DndMonster}



\begin{DndMonster}[float=t]{Harpy}
\DndMonsterType{Medium monstrosity, chaotic evil}
\DndMonsterBasics[
armor-class = {11},
hit-points = {\DndDice{7d8 + 7}},
speed = {20 ft., fly 40 ft.}
]
\DndMonsterAbilityScores[
 str = 12,
 dex = 13,
 con = 12,
 int = 7,
 wis = 10,
 cha = 13
]
\DndMonsterDetails[
languages = {Common},
challenge = 1,
]
\par
\DndMonsterSection{Actions}
\DndMonsterAction{Multiattack}
The harpy makes two attacks: one with its claws and one with its club.

\DndMonsterAction{Claws}
\textsl{Melee Weapon Attack:} +3 to hit, reach 5 ft., one target. \textsl{Hit:} 6 (2d4 + 1) slashing damage.

\DndMonsterAction{Club}
\textsl{Melee Weapon Attack:} +3 to hit, reach 5 ft., one target. \textsl{Hit:} 3 (1d4 + 1) bludgeoning damage.

\DndMonsterAction{Luring Song}
The harpy sings a magical melody. Every humanoid and giant within 300 feet of the harpy that can hear the song must succeed on a DC 11 Wisdom saving throw or be charmed until the song ends. The harpy must take a bonus action on its subsequent turns to continue singing. It can stop singing at any time. The song ends if the harpy is incapacitated.

While charmed by the harpy, a target is incapacitated and ignores the songs of other harpies. If the charmed target is more than 5 feet away from the harpy, the target must move on its turn toward the harpy by the most direct route. It doesn't avoid opportunity attacks, but before moving into damaging terrain, such as lava or a pit, and whenever it takes damage from a source other than the harpy, a target can repeat the saving throw. A creature can also repeat the saving throw at the end of each of its turns. If a creature's saving throw is successful, the effect ends on it.

A target that successfully saves is immune to this harpy's song for the next 24 hours.

\end{DndMonster}



\begin{DndMonster}[float=t]{Harpy Matriarch}
\DndMonsterType{Medium monstrosity, chaotic evil}
\DndMonsterBasics[
armor-class = {14 (natural armor)},
hit-points = {\DndDice{16d8 + 16}},
speed = {20 ft., fly 40 ft.}
]
\DndMonsterAbilityScores[
 str = 13,
 dex = 16,
 con = 12,
 int = 9,
 wis = 10,
 cha = 16
]
\DndMonsterDetails[
saving-throws = {Dex +6, Cha +6},
skills = {Perception +3},
senses = {darkvision 60 ft., passive perception 13},
languages = {Common},
challenge = 5,
]
\DndMonsterAction{Luring Maestro}
While within 60 feet of the matriarch, creatures have disadvantage on saving throws against the matriarch's Luring Song.

\DndMonsterAction{Magic Resistance}
The matriarch has advantage on saving throws against spells and other magical effects.

\par
\DndMonsterSection{Actions}
\DndMonsterAction{Multiattack}
The matriarch makes two claws attacks.

\DndMonsterAction{Claws}
\textsl{Melee Weapon Attack:} +6 to hit, reach 5 ft., one target. \textsl{Hit:} 13 (3d6 + 3) slashing damage.

\DndMonsterAction{Fleeting Form}
The matriarch can magically disguise itself to resemble a humanoid of roughly similar size and shape for up to 1 hour. It can revert to its true form as a bonus action. This illusion does not hold up to close scrutiny.

\DndMonsterAction{Luring Song}
The matriarch sings a magical melody. Every humanoid and giant within 300 feet of the matriarch that can hear the song must succeed on a DC 14 Wisdom saving throw or be charmed until the song ends. The matriarch must take a bonus action on its subsequent turns to continue singing. It can stop singing at any time. The song ends if the matriarch is incapacitated.

While charmed by the matriarch, a target is incapacitated and ignores the songs of other harpies. If the charmed target is more than 5 feet away from the matriarch, the target must move on its turn toward the matriarch by the most direct route, trying to get within 5 feet. It doesn't avoid opportunity attacks, but before moving into damaging terrain, such as lava or a pit, and whenever it takes damage from a source other than the matriarch, the target can repeat the saving throw. A charmed target can also repeat the saving throw at the end of each of its turns. If the saving throw is successful, the effect ends on it.

A target that successfully saves is immune to this matriarch's song for the next 24 hours.

\DndMonsterAction{Visage of Desire (1/Day)}
The matriarch projects a vision into the minds of creatures within 30 feet of it that aren't constructs or undead, showing each creature achieving whatever it most desires. An affected creature must succeed on a DC 14 Wisdom saving throw or drop whatever it is holding and become paralyzed until the end of its next turn.

\end{DndMonster}



\begin{DndMonster}[float=t]{Hill Giant}
\DndMonsterType{Huge giant, chaotic evil}
\DndMonsterBasics[
armor-class = {13 (natural armor)},
hit-points = {\DndDice{10d12 + 40}},
speed = {40 ft.}
]
\DndMonsterAbilityScores[
 str = 21,
 dex = 8,
 con = 19,
 int = 5,
 wis = 9,
 cha = 6
]
\DndMonsterDetails[
skills = {Perception +2},
languages = {Giant},
challenge = 5,
]
\par
\DndMonsterSection{Actions}
\DndMonsterAction{Multiattack}
The giant makes two greatclub attacks.

\DndMonsterAction{Greatclub}
\textsl{Melee Weapon Attack:} +8 to hit, reach 10 ft., one target. \textsl{Hit:} 18 (3d8 + 5) bludgeoning damage.

\DndMonsterAction{Rock}
\textsl{Ranged Weapon Attack:} +8 to hit, range 60/240 ft., one target. \textsl{Hit:} 21 (3d10 + 5) bludgeoning damage.

\end{DndMonster}



\begin{DndMonster}[float=t]{Hippogriff}
\DndMonsterType{Large monstrosity, unaligned}
\DndMonsterBasics[
armor-class = {11},
hit-points = {\DndDice{3d10 + 3}},
speed = {40 ft., fly 60 ft.}
]
\DndMonsterAbilityScores[
 str = 17,
 dex = 13,
 con = 13,
 int = 2,
 wis = 12,
 cha = 8
]
\DndMonsterDetails[
skills = {Perception +5},
challenge = 1,
]
\DndMonsterAction{Keen Sight}
The hippogriff has advantage on Wisdom (Perception) checks that rely on sight.

\par
\DndMonsterSection{Actions}
\DndMonsterAction{Multiattack}
The hippogriff makes two attacks: one with its beak and one with its claws.

\DndMonsterAction{Beak}
\textsl{Melee Weapon Attack:} +5 to hit, reach 5 ft., one target. \textsl{Hit:} 8 (1d10 + 3) piercing damage.

\DndMonsterAction{Claws}
\textsl{Melee Weapon Attack:} +5 to hit, reach 5 ft., one target. \textsl{Hit:} 10 (2d6 + 3) slashing damage.

\end{DndMonster}



\begin{DndMonster}[float=t]{Hobgoblin}
\DndMonsterType{Medium humanoid (goblinoid), lawful evil}
\DndMonsterBasics[
armor-class = {18 (\textit{chain mail})},
hit-points = {\DndDice{2d8 + 2}},
speed = {30 ft.}
]
\DndMonsterAbilityScores[
 str = 13,
 dex = 12,
 con = 12,
 int = 10,
 wis = 10,
 cha = 9
]
\DndMonsterDetails[
senses = {darkvision 60 ft., passive perception 10},
languages = {Common, Goblin},
challenge = 1/2,
]
\DndMonsterAction{Martial Advantage}
Once per turn, the hobgoblin can deal an extra 7 (2d6) damage to a creature it hits with a weapon attack if that creature is within 5 feet of an ally of the hobgoblin that isn't incapacitated.

\par
\DndMonsterSection{Actions}
\DndMonsterAction{Longsword}
\textsl{Melee Weapon Attack:} +3 to hit, reach 5 ft., one target. \textsl{Hit:} 5 (1d8 + 1) slashing damage, or 6 (1d10 + 1) slashing damage if used with two hands.

\DndMonsterAction{Longbow}
\textsl{Ranged Weapon Attack:} +3 to hit, range 150/600 ft., one target. \textsl{Hit:} 5 (1d8 + 1) piercing damage.

\end{DndMonster}



\begin{DndMonster}[float=t]{Hydra}
\DndMonsterType{Huge monstrosity, unaligned}
\DndMonsterBasics[
armor-class = {15 (natural armor)},
hit-points = {\DndDice{15d12 + 75}},
speed = {30 ft., swim 30 ft.}
]
\DndMonsterAbilityScores[
 str = 20,
 dex = 12,
 con = 20,
 int = 2,
 wis = 10,
 cha = 7
]
\DndMonsterDetails[
skills = {Perception +6},
senses = {darkvision 60 ft., passive perception 16},
challenge = 8,
]
\DndMonsterAction{Hold Breath}
The hydra can hold its breath for 1 hour.

\DndMonsterAction{Multiple Heads}
The hydra has five heads. While it has more than one head, the hydra has advantage on saving throws against being blinded, charmed, deafened, frightened, stunned, and knocked unconscious.

Whenever the hydra takes 25 or more damage in a single turn, one of its heads dies. If all its heads die, the hydra dies.

At the end of its turn, it grows two heads for each of its heads that died since its last turn, unless it has taken fire damage since its last turn. The hydra regains 10 hit points for each head regrown in this way.

\DndMonsterAction{Reactive Heads}
For each head the hydra has beyond one, it gets an extra reaction that can be used only for opportunity attacks.

\DndMonsterAction{Wakeful}
While the hydra sleeps, at least one of its heads is awake.

\par
\DndMonsterSection{Actions}
\DndMonsterAction{Multiattack}
The hydra makes as many bite attacks as it has heads.

\DndMonsterAction{Bite}
\textsl{Melee Weapon Attack:} +8 to hit, reach 10 ft., one target. \textsl{Hit:} 10 (1d10 + 5) piercing damage.

\end{DndMonster}



\begin{DndMonster}[float=t]{Hyena}
\DndMonsterType{Medium beast, unaligned}
\DndMonsterBasics[
armor-class = {11},
hit-points = {\DndDice{1d8 + 1}},
speed = {50 ft.}
]
\DndMonsterAbilityScores[
 str = 11,
 dex = 13,
 con = 12,
 int = 2,
 wis = 12,
 cha = 5
]
\DndMonsterDetails[
skills = {Perception +3},
challenge = 0,
]
\DndMonsterAction{Pack Tactics}
The hyena has advantage on an attack roll against a creature if at least one of the hyena's allies is within 5 feet of the creature and the ally isn't incapacitated.

\par
\DndMonsterSection{Actions}
\DndMonsterAction{Bite}
\textsl{Melee Weapon Attack:} +2 to hit, reach 5 ft., one target. \textsl{Hit:} 3 (1d6) piercing damage.

\end{DndMonster}



\begin{DndMonster}[float=t]{Imp}
\DndMonsterType{Tiny fiend (devil), lawful evil}
\DndMonsterBasics[
armor-class = {13},
hit-points = {\DndDice{3d4 + 3}},
speed = {20 ft., fly 40 ft.}
]
\DndMonsterAbilityScores[
 str = 6,
 dex = 17,
 con = 13,
 int = 11,
 wis = 12,
 cha = 14
]
\DndMonsterDetails[
skills = {Deception +4, Insight +3, Persuasion +4, Stealth +5},
damage-resistances = {cold; bludgeoning, piercing, slashing from nonmagical attacks not made with silvered weapons},
damage-immunities = {fire, poison},
senses = {darkvision 120 ft., passive perception 11},
languages = {Infernal, Common},
challenge = 1,
]
\DndMonsterAction{Shapechanger}
The imp can use its action to polymorph into a beast form that resembles a rat (speed 20 ft.), a raven (20 ft., fly 60 ft.), or a spider (20 ft., climb 20 ft.), or back into its true form. Its statistics are the same in each form, except for the speed changes noted. Any equipment it is wearing or carrying isn't transformed. It reverts to its true form if it dies.

\DndMonsterAction{Devil's Sight}
Magical darkness doesn't impede the imp's darkvision.

\DndMonsterAction{Magic Resistance}
The imp has advantage on saving throws against spells and other magical effects.

\par
\DndMonsterSection{Actions}
\DndMonsterAction{Sting (Bite in Beast Form)}
\textsl{Melee Weapon Attack:} +5 to hit, reach 5 ft., one target. \textsl{Hit:} 5 (1d4 + 3) piercing damage, and the target must make a DC 11 Constitution saving throw, taking 10 (3d6) poison damage on a failed save, or half as much damage on a successful one.

\DndMonsterAction{Invisibility}
The imp magically turns invisible until it attacks, or until its concentration ends (as if concentration on a spell). Any equipment the imp wears or carries is invisible with it.

\end{DndMonster}



\begin{DndMonster}[float=t]{Invisible Stalker}
\DndMonsterType{Medium elemental, neutral}
\DndMonsterBasics[
armor-class = {14},
hit-points = {\DndDice{16d8 + 32}},
speed = {50 ft., fly 50 ft. \textit{(hover)}}
]
\DndMonsterAbilityScores[
 str = 16,
 dex = 19,
 con = 14,
 int = 10,
 wis = 15,
 cha = 11
]
\DndMonsterDetails[
skills = {Perception +8, Stealth +10},
damage-resistances = {; bludgeoning, piercing, slashing from nonmagical attacks},
damage-immunities = {poison},
senses = {darkvision 60 ft., passive perception 18},
languages = {Auran, understands Common but doesn't speak it},
challenge = 6,
]
\DndMonsterAction{Invisibility}
The stalker is invisible.

\DndMonsterAction{Faultless Tracker}
The stalker is given a quarry by its summoner. The stalker knows the direction and distance to its quarry as long as the two of them are on the same plane of existence. The stalker also knows the location of its summoner.

\par
\DndMonsterSection{Actions}
\DndMonsterAction{Multiattack}
The stalker makes two slam attacks.

\DndMonsterAction{Slam}
\textsl{Melee Weapon Attack:} +6 to hit, reach 5 ft., one target. \textsl{Hit:} 10 (2d6 + 3) bludgeoning damage.

\end{DndMonster}


\clearpage

\begin{DndMonster}[float=t]{Jourrael, the Caedogeist}
\DndMonsterType{Medium fiend (drow), chaotic evil}
\DndMonsterBasics[
armor-class = {19 (\textit{studded leather armor})},
hit-points = {\DndDice{16d8 + 80}},
speed = {80 ft., fly 40 ft.}
]
\DndMonsterAbilityScores[
 str = 13,
 dex = 24,
 con = 20,
 int = 14,
 wis = 17,
 cha = 15
]
\DndMonsterDetails[
saving-throws = {Dex +12, Con +10, Wis +8},
skills = {Acrobatics +12, Perception +13, Stealth +17},
damage-resistances = {acid, fire, lightning, cold, thunder; bludgeoning, piercing, slashing from nonmagical attacks},
damage-immunities = {poison},
senses = {darkvision 120 ft., passive perception 23},
languages = {Abyssal, Common, Infernal, Undercommon},
challenge = 15,
]
\DndMonsterAction{Incorporeal Movement}
Jourrael can move through other creatures and objects as if they were difficult terrain. They take 5 (1d10) force damage if they end their turn inside an object.

\DndMonsterAction{Sneak Attack (1/Turn)}
Jourrael deals an extra 35 (10d6) damage when they hit a target with a weapon attack and have Advantage and Disadvantage on the attack roll, or when the target is within 5 feet of an ally of Jourrael that isn't incapacitated and Jourrael doesn't have Advantage and Disadvantage on the attack roll.

\DndMonsterAction{Inevitable}
If Jourrael is reduced to 0 hit points, they do not fall unconscious but becomes an incorporeal spirit for 24 hours. During this time, Jourrael can't return to their original form until they returns to full hit points. In this form, they are invisible, can't attack or interact with corporeal creatures or objects, and can't be targeted by spells. After 24 hours have passed, Jourrael regains 1 hit point and returns to their original form with four levels of exhaustion.

\DndMonsterAction{Innate Spellcasting}
Jourrael's spellcasting ability is Charisma (spell save DC 15). She can innately cast the following spells, requiring no material components:
\begin{DndMonsterSpells}
  \DndInnateSpellLevel{dancing lights}
  \DndInnateSpellLevel[1e]{darkness, faerie fire}
\end{DndMonsterSpells}
\par
\DndMonsterSection{Actions}
\DndMonsterAction{Multiattack}
Jourrael makes two melee attacks or two ranged attacks.

\DndMonsterAction{Spectral Blade}
\textsl{Melee Weapon Attack:} +12 to hit, reach 5 ft., one target. \textsl{Hit:} 10 (1d6 + 7) slashing damage, and the target must succeed on a DC 14 Constitution saving throws or be poisoned for 1 hour. If the saving throws fails by 5 or more, the target is also unconscious while poisoned in this way. The target wakes up if it takes damage or if another creature takes an action to shake it awake.

\DndMonsterAction{Void Dagger}
\textsl{Ranged Weapon Attack:} +12 to hit, range 60/120 ft., one target. \textsl{Hit:} 9 (1d4 + 7) piercing damage.

\par
\DndMonsterSection{Bonus Actions}
\DndMonsterAction{Cunning Action}
Jourrael can take the Dash, Disengage, or Hide action.

\DndMonsterAction{Mirrored Form (Recharges after a resting)}
Jourrael manifests three illusory duplicates that hover in her space. Each time a creature targets Jourrael with an attack, roll a d20 to see if it hits one of the duplicates. If there are three duplicates, a duplicate is targeted on a roll of 6 or higher. With two duplicates, a duplicate is targeted on a roll of 8 or higher. A single duplicate is targeted on a roll of 11 or higher. A duplicate's AC is 17, and it is destroyed if hit. Jourrael can hide the duplicates as a bonus action.

\par
\DndMonsterSection{Reactions}
\DndMonsterAction{Parry}
Jourrael adds 5 to her AC against one melee attack that would hit her. To do so, she must see the attacker

\end{DndMonster}



\begin{DndMonster}[float*=b,width=\textwidth + 8pt]{Juvenile Kraken}
\begin{multicols}{2}
\DndMonsterType{Huge monstrosity (titan), chaotic evil}
\DndMonsterBasics[
armor-class = {16 (natural armor)},
hit-points = {\DndDice{18d12 + 90}},
speed = {20 ft., swim 50 ft.}
]
\DndMonsterAbilityScores[
 str = 24,
 dex = 11,
 con = 20,
 int = 19,
 wis = 15,
 cha = 17
]
\DndMonsterDetails[
saving-throws = {Str +12, Dex +5, Con +10, Int +9, Wis +7},
damage-resistances = {; bludgeoning, piercing, slashing from nonmagical attacks},
damage-immunities = {lightning},
senses = {truesight 120 ft., passive perception 12},
languages = {Abyssal, Celestial, Infernal, Primordial, telepathy 60 ft. but can't speak},
challenge = 14,
]
\DndMonsterAction{Amphibious}
The kraken can breathe air and water.

\DndMonsterAction{Freedom of Movement}
The kraken ignores difficult terrain, and magical effects can't reduce its speed or cause it to be restrained. It can spend 5 feet of movement to escape from nonmagical restraints or being grappled.

\par
\DndMonsterSection{Actions}
\DndMonsterAction{Multiattack}
The kraken makes two tentacle attacks, each of which it can replace with a use of Fling.

\DndMonsterAction{Bite}
\textsl{Melee Weapon Attack:} +12 to hit, reach 5 ft., one target. \textsl{Hit:} 20 (3d8 + 7) piercing damage. If the target is a Medium or smaller creature grappled by the kraken, that creature is swallowed and the grapple ends. While swallowed, the creature is blinded and restrained, it has total cover against attacks and other effects outside the kraken, and it takes 21 (6d6) acid damage at the start of each of the kraken's turns. One Medium or two smaller creatures can be swallowed at the same time.

If the kraken takes 35 damage or more on a single turn from a creature inside it, the kraken must succeed on a DC 23 Constitution saving throw at the end of that turn or regurgitate all swallowed creatures, which fall prone in spaces within 10 feet of the kraken. If the kraken dies, a swallowed creature is no longer restrained by it and can escape from the corpse using 10 feet of movement, exiting prone.

\DndMonsterAction{Tentacle}
\textsl{Melee Weapon Attack:} +12 to hit, reach 20 ft., one target. \textsl{Hit:} 17 (3d6 + 7) bludgeoning damage, and the target is grappled (escape DC 20). Until the grapple ends, the target is restrained. The kraken has ten tentacles, each of which can grapple one target.

\DndMonsterAction{Fling}
One Medium or smaller object held or creature grappled by the kraken is thrown up to 40 feet in a random direction and knocked prone. If a thrown target strikes a solid surface, the target takes 3 (1d6) bludgeoning damage for every 10 feet it was thrown. If the target is thrown at another creature, that creature must succeed on a DC 13 Dexterity saving throw or take the same damage and be knocked prone.

\DndMonsterAction{Lightning Strike}
The kraken magically create a bolt of lightning, which can strike a target the kraken can see within 90 feet of it. The target must make a DC 18 Dexterity saving throw, taking 22 (4d10) lightning damage on a failed save, or half as much damage on a successful one.

\par
\DndMonsterSection{Legendary Actions}
Juvenile Kraken can take 3 legendary actions, choosing from the options below. Only one legendary action can be used at a time and only at the end of another creature's turn. Juvenile Kraken regains spent legendary actions at the start of its turn.

\begin{DndMonsterLegendaryActions}
\DndMonsterLegendaryAction{Tentacle Attack (Costs 2 Actions)}{The kraken makes one tentacle attack.
}
\DndMonsterLegendaryAction{Fling}{The kraken uses Fling.
}
\DndMonsterLegendaryAction{Ink Cloud (Costs 3 Actions)}{While underwater, the kraken expels an ink cloud in a 40-foot radius. The cloud spreads around corners, and that area is heavily obscured to creatures other than the kraken. Each creature other than the kraken that ends its turn there must succeed on a DC 18 Constitution saving throw, taking 11 (2d10) poison damage on a failed save or half as much damage on a successful one. A strong current disperses the cloud, which otherwise disappears at the end of the kraken's next turn.
}
\end{DndMonsterLegendaryActions}
\end{multicols}
\end{DndMonster}



\begin{DndMonster}[float=t]{Keyleth, Voice of the Tempest}
\DndMonsterType{Medium humanoid (half-elf), U}
\DndMonsterBasics[
armor-class = {17 (\textit{+2 leather armor})},
hit-points = {\DndDice{20d8 + 60}},
speed = {30 ft.}
]
\DndMonsterAbilityScores[
 str = 14,
 dex = 15,
 con = 16,
 int = 15,
 wis = 22,
 cha = 15
]
\DndMonsterDetails[
saving-throws = {Str +4, Dex +4, Con +5, Int +10, Wis +14, Cha +4},
skills = {Athletics +8, Insight +12, Intimidation +8, Nature +12, Perception +12, Persuasion +8, Stealth +8, Survival +12},
senses = {darkvision 60 ft., passive perception 22},
languages = {Auran, Common, Druidic, Elvish},
challenge = 18,
]
\DndMonsterAction{Fey Ancestry}
Keyleth has Advantage and Disadvantage on saving throws against being charmed, and magic can't put her to sleep.

\DndMonsterAction{Focused}
Keyleth has Advantage and Disadvantage on Constitution saving throws made to maintain concentration on spells.

\DndMonsterAction{Special Equipment}
Keyleth wears a \textit{circlet of wisdom} (which increases her Wisdom score by 2), \textit{+2 leather armor}, and a ring of protection with a +2 bonus. She wields the Spire of Conflux.

\DndMonsterAction{Spellcasting}
Keyleth is a 20th-level spellcaster. Her spellcasting ability is Wisdom (spell save DC 22, +14 to hit with spell attacks). She has the following druid spells prepared:
\begin{DndMonsterSpells}
  \DndMonsterSpellLevel{druidcraft, guidance, mending}
   \DndMonsterSpellLevel[1][4]{cure wounds, faerie fire, fog cloud, healing word, thunderwave}
   \DndMonsterSpellLevel[2][3]{alter self, flaming sphere, gust of wind, pass without trace}
   \DndMonsterSpellLevel[3][3]{call lightning, wind wall}
   \DndMonsterSpellLevel[4][3]{blight, control water, grasping vine, polymorph, stone shape}
   \DndMonsterSpellLevel[5][3]{conjure elemental, greater restoration, scrying, wall of stone}
   \DndMonsterSpellLevel[6][2]{heroes' feast, move earth, transport via plants}
   \DndMonsterSpellLevel[7][2]{fire storm, plane shift}
   \DndMonsterSpellLevel[8][1]{control weather}
\end{DndMonsterSpells}
\par
\DndMonsterSection{Actions}
\DndMonsterAction{Shapechange}
Keyleth casts \textit{shapechange} to transform into a \textbf{planetar}, an \textbf{adult bronze dragon}, or another creature of CR 20 or lower that is not a construct or undead.

\DndMonsterAction{Meteor (\textbf{Earth Elemental} Form Only)}
When Keyleth falls at least 20 feet, she lands with meteoric force. Each creature within 20 feet of her must succeed on a DC 19 Dexterity saving throws or take 1d6 bludgeoning damage for every 10 feet Keyleth fell. Keyleth still takes falling while in \textbf{earth elemental} form.

\DndMonsterAction{Chain Lightning}
Keyleth casts \textit{chain lightning} (spell save DC 22) from the Spire of Conflux at a creature she can see within 150 feet of her.

\par
\DndMonsterSection{Bonus Actions}
\DndMonsterAction{Wild Shape}
Keyleth magically \textit{polymorph} into a beast or elemental of CR 6 or lower. Her game statistics are replaced by the statistics of the new form, but she retains her alignment, personality, and Intelligence, Wisdom, and Charisma scores. She also retains all of her skill and saving throws proficiencies, in addition to gaining those of the new form. She assumes the new form's hit points and Hit Dice and returns to the number of hit points she had when she reverts to her normal form. Any equipment she is wearing or carrying is absorbed or borne by her new form (her choice). She reverts to her humanoid form as a bonus action, or when she falls unconscious. Her attacks in Wild Shape form are magical, and she can cast spells while transformed.

\DndMonsterAction{Healing Word (4th Level)}
Keyleth casts \textit{healing word}, restoring 16 (4d4 + 6) hit points to herself or another creature she can see within 60 feet of her.

\end{DndMonster}



\begin{DndMonster}[float=t]{Knight}
\DndMonsterType{Medium humanoid (any race), any alignment}
\DndMonsterBasics[
armor-class = {18 (\textit{plate armor})},
hit-points = {\DndDice{8d8 + 16}},
speed = {30 ft.}
]
\DndMonsterAbilityScores[
 str = 16,
 dex = 11,
 con = 14,
 int = 11,
 wis = 11,
 cha = 15
]
\DndMonsterDetails[
saving-throws = {Con +4, Wis +2},
languages = {any one language (usually Common)},
challenge = 3,
]
\DndMonsterAction{Brave}
The knight has advantage on saving throws against being frightened.

\par
\DndMonsterSection{Actions}
\DndMonsterAction{Multiattack}
The knight makes two melee attacks.

\DndMonsterAction{Greatsword}
\textsl{Melee Weapon Attack:} +5 to hit, reach 5 ft., one target. \textsl{Hit:} 10 (2d6 + 3) slashing damage.

\DndMonsterAction{Heavy Crossbow}
\textsl{Ranged Weapon Attack:} +2 to hit, range 100/400 ft., one target. \textsl{Hit:} 5 (1d10) piercing damage.

\DndMonsterAction{Leadership (Recharges after a Short or Long Rest)}
For 1 minute, the knight can utter a special command or warning whenever a nonhostile creature that it can see within 30 feet of it makes an attack roll or a saving throw. The creature can add a d4 to its roll provided it can hear and understand the knight. A creature can benefit from only one Leadership die at a time. This effect ends if the knight is incapacitated.

\par
\DndMonsterSection{Reactions}
\DndMonsterAction{Parry}
The knight adds 2 to its AC against one melee attack that would hit it. To do so, the knight must see the attacker and be wielding a melee weapon.

\end{DndMonster}



\begin{DndMonster}[float=t]{Kobold}
\DndMonsterType{Small humanoid (kobold), lawful evil}
\DndMonsterBasics[
armor-class = {12},
hit-points = {\DndDice{2d6 - 2}},
speed = {30 ft.}
]
\DndMonsterAbilityScores[
 str = 7,
 dex = 15,
 con = 9,
 int = 8,
 wis = 7,
 cha = 8
]
\DndMonsterDetails[
senses = {darkvision 60 ft., passive perception 8},
languages = {Common, Draconic},
challenge = 1/8,
]
\DndMonsterAction{Sunlight Sensitivity}
While in sunlight, the kobold has disadvantage on attack rolls, as well as on Wisdom (Perception) checks that rely on sight.

\DndMonsterAction{Pack Tactics}
The kobold has advantage on an attack roll against a creature if at least one of the kobold's allies is within 5 feet of the creature and the ally isn't incapacitated.

\par
\DndMonsterSection{Actions}
\DndMonsterAction{Dagger}
\textsl{Melee Weapon Attack:} +4 to hit, reach 5 ft., one target. \textsl{Hit:} 4 (1d4 + 2) piercing damage.

\DndMonsterAction{Sling}
\textsl{Ranged Weapon Attack:} +4 to hit, range 30/120 ft., one target. \textsl{Hit:} 4 (1d4 + 2) bludgeoning damage.

\end{DndMonster}



\begin{DndMonster}[float=t]{Kraghammer Goat-Knight}
\DndMonsterType{Medium humanoid (dwarf), U}
\DndMonsterBasics[
armor-class = {20 (plate)},
hit-points = {\DndDice{8d8 + 16}},
speed = {25 ft.}
]
\DndMonsterAbilityScores[
 str = 17,
 dex = 8,
 con = 14,
 int = 10,
 wis = 11,
 cha = 16
]
\DndMonsterDetails[
saving-throws = {Str +5, Dex +1, Con +4, Int +2, Wis +2, Cha +5},
skills = {Nature +2, Religion +4},
damage-resistances = {poison},
senses = {darkvision 60 ft., passive perception 10},
languages = {Common, Dwarvish},
challenge = 3,
]
\DndMonsterAction{Aura of Protection}
Whenever the goat-knight or a creature friendly to it within 10 feet of it makes a saving throws, that creature gains a +2 bonus (included in the goat-knight's saving throws above). The goat-knight must be conscious to grant and gain this bonus.

\DndMonsterAction{Divine Health}
The goat-knight is immune to disease.

\DndMonsterAction{Dwarven Resilience}
The goat-knight has resistance to poison damage and Advantage and Disadvantage on saving throws against poison.

\DndMonsterAction{Spellcasting}
The goat-knight is a 6th-level spellcaster. Its spellcasting ability is Charisma (spell save DC 13, +5 to hit with spell attacks). It has the following paladin spells prepared:
\begin{DndMonsterSpells}
   \DndMonsterSpellLevel[1][4]{bless, cure wounds, protection from evil and good, sanctuary, shield of faith}
   \DndMonsterSpellLevel[2][2]{branding smite, find steed, lesser restoration, zone of truth}
\end{DndMonsterSpells}
\par
\DndMonsterSection{Actions}
\DndMonsterAction{Multiattack}
The goat-knight makes two warhammer attacks.

\DndMonsterAction{Warhammer}
\textsl{Melee Weapon Attack:} +5 to hit, reach 5 ft., one target. \textsl{Hit:} 7 (1d8 + 3) bludgeoning damage, or 8 (1d10 + 3) bludgeoning damage if used with two hands.

\end{DndMonster}



\begin{DndMonster}[float*=b,width=\textwidth + 8pt]{Kraken}
\begin{multicols}{2}
\DndMonsterType{Gargantuan monstrosity (titan), chaotic evil}
\DndMonsterBasics[
armor-class = {18 (natural armor)},
hit-points = {\DndDice{27d20 + 189}},
speed = {20 ft., swim 60 ft.}
]
\DndMonsterAbilityScores[
 str = 30,
 dex = 11,
 con = 25,
 int = 22,
 wis = 18,
 cha = 20
]
\DndMonsterDetails[
saving-throws = {Str +17, Dex +7, Con +14, Int +13, Wis +11},
damage-immunities = {lightning; bludgeoning, piercing, slashing from nonmagical attacks},
senses = {truesight 120 ft., passive perception 14},
languages = {Abyssal, Celestial, Infernal, Primordial, telepathy 120 ft. but can't speak},
challenge = 23,
]
\DndMonsterAction{Amphibious}
The kraken can breathe air and water.

\DndMonsterAction{Freedom of Movement}
The kraken ignores difficult terrain, and magical effects can't reduce its speed or cause it to be restrained. It can spend 5 feet of movement to escape from nonmagical restraints or being grappled.

\DndMonsterAction{Siege Monster}
The kraken deals double damage to objects and structures.

\par
\DndMonsterSection{Actions}
\DndMonsterAction{Multiattack}
The kraken makes three tentacle attacks, each of which it can replace with one use of Fling.

\DndMonsterAction{Bite}
\textsl{Melee Weapon Attack:} +17 to hit, reach 5 ft., one target. \textsl{Hit:} 23 (3d8 + 10) piercing damage. If the target is a Large or smaller creature grappled by the kraken, that creature is swallowed, and the grapple ends. While swallowed, the creature is blinded and restrained, it has total cover against attacks and other effects outside the kraken, and it takes 42 (12d6) acid damage at the start of each of the kraken's turns. If the kraken takes 50 damage or more on a single turn from a creature inside it, the kraken must succeed on a DC 25 Constitution saving throw at the end of that turn or regurgitate all swallowed creatures, which fall prone in a space within 10 feet of the kraken. If the kraken dies, a swallowed creature is no longer restrained by it and can escape from the corpse using 15 feet of movement, exiting prone.

\DndMonsterAction{Tentacle}
\textsl{Melee Weapon Attack:} +17 to hit, reach 30 ft., one target. \textsl{Hit:} 20 (3d6 + 10) bludgeoning damage, and the target is grappled (escape DC 18). Until this grapple ends, the target is restrained. The kraken has ten tentacles, each of which can grapple one target.

\DndMonsterAction{Fling}
One Large or smaller object held or creature grappled by the kraken is thrown up to 60 feet in a random direction and knocked prone. If a thrown target strikes a solid surface, the target takes 3 (1d6) bludgeoning damage for every 10 feet it was thrown. If the target is thrown at another creature, that creature must succeed on a DC 18 Dexterity saving throw or take the same damage and be knocked prone.

\DndMonsterAction{Lightning Storm}
The kraken magically creates three bolts of lightning, each of which can strike a target the kraken can see within 120 feet of it. A target must make a DC 23 Dexterity saving throw, taking 22 (4d10) lightning damage on a failed save, or half as much damage on a successful one.

\par
\DndMonsterSection{Legendary Actions}
Kraken can take 3 legendary actions, choosing from the options below. Only one legendary action can be used at a time and only at the end of another creature's turn. Kraken regains spent legendary actions at the start of its turn.

\begin{DndMonsterLegendaryActions}
\DndMonsterLegendaryAction{Tentacle Attack or Fling}{The kraken makes one tentacle attack or uses its Fling.
}
\DndMonsterLegendaryAction{Lightning Storm (Costs 2 Actions)}{The kraken uses Lightning Storm.
}
\DndMonsterLegendaryAction{Ink Cloud (Costs 3 Actions)}{While underwater, the kraken expels an ink cloud in a 60-foot radius. The cloud spreads around corners, and that area is heavily obscured to creatures other than the kraken. Each creature other than the kraken that ends its turn there must succeed on a DC 23 Constitution saving throw, taking 16 (3d10) poison damage on a failed save, or half as much damage on a successful one. A strong current disperses the cloud, which otherwise disappears at the end of the kraken's next turn.
}
\end{DndMonsterLegendaryActions}
\par \DndMonsterSection{Lair Actions}
On initiative count 20 (losing initiative ties), the kraken takes a lair action to cause one of the following magical effects:


\begin{itemize}
\item A strong current moves through the kraken's lair. Each creature within 60 feet of the kraken must succeed on a DC 23 Strength saving throw or be pushed up to 60 feet away from the kraken. On a success, the creature is pushed 10 feet away from the kraken.
\item Creatures in the water within 60 feet of the kraken have vulnerability to lightning damage until initiative count 20 on the next round.
\item The water in the kraken's lair becomes electrically charged. All creatures within 120 feet of the kraken must succeed on a DC 23 Constitution saving throw, taking 10 (3d6) lightning damage on a failed save, or half as much damage on a successful one.
\end{itemize}

\par \DndMonsterSection{Regional Effects}
The region containing a kraken's lair is warped by the creature's blasphemous presence, creating the following magical effects:


\begin{itemize}
\item The kraken can alter the weather at will in a 6-mile radius centered on its lair. The effect is identical to the \textit{control weather} spell.
\item Water elementals coalesce within 6 miles of the lair. These elementals can't leave the water and have Intelligence and Charisma scores of 1 (-5).
\item Aquatic creatures within 6 miles of the lair that have an Intelligence score of 2 or lower are charmed by the kraken and aggressive toward intruders in the area.
\end{itemize}

When the kraken dies, all of these regional effects fade immediately.

\end{multicols}
\end{DndMonster}



\begin{DndMonster}[float*=b,width=\textwidth + 8pt]{Leviathan}
\begin{multicols}{2}
\DndMonsterType{Gargantuan elemental, neutral}
\DndMonsterBasics[
armor-class = {17},
hit-points = {\DndDice{16d20 + 160}},
speed = {40 ft., swim 120 ft.}
]
\DndMonsterAbilityScores[
 str = 30,
 dex = 24,
 con = 30,
 int = 2,
 wis = 18,
 cha = 17
]
\DndMonsterDetails[
saving-throws = {Wis +10, Cha +9},
damage-resistances = {; bludgeoning, piercing, slashing from nonmagical attacks},
damage-immunities = {acid, poison},
senses = {darkvision 60 ft., passive perception 14},
challenge = 20,
]
\DndMonsterAction{Legendary Resistance (3/Day)}
If the leviathan fails a saving throw, it can choose to succeed instead.

\DndMonsterAction{Partial Freeze}
If the leviathan takes 50 cold damage or more during a single turn, the leviathan partially freezes; until the end of its next turn, its speeds are reduced to 20 feet, and it makes attack rolls with disadvantage.

\DndMonsterAction{Siege Monster}
The leviathan deals double damage to objects and structures (included in Tidal Wave).

\DndMonsterAction{Water Form}
The leviathan can enter a hostile creature's space and stop there. It can move through a space as narrow as 1 inch wide without squeezing.

\par
\DndMonsterSection{Actions}
\DndMonsterAction{Multiattack}
The leviathan makes two attacks: one with its slam and one with its tail.

\DndMonsterAction{Slam}
\textsl{Melee Weapon Attack:} +16 to hit, reach 20 ft., one target. \textsl{Hit:} 15 (1d10 + 10) bludgeoning damage plus 5 (1d10) acid damage.

\DndMonsterAction{Tail}
\textsl{Melee Weapon Attack:} +16 to hit, reach 20 ft., one target. \textsl{Hit:} 16 (1d12 + 10) bludgeoning damage plus 6 (1d12) acid damage.

\DndMonsterAction{Tidal Wave (Recharge 6)}
While submerged, the leviathan magically creates a wall of water centered on itself. The wall is up 250 feet long, up to 250 feet high, and up to 50 feet thick. When the wall appears, all other creatures within its area must each make a DC 24 Strength saving throw. A creature takes 33 (6d10) bludgeoning damage on failed save, or half as much damage on a successful one.

At the start of each of the leviathan's turns after the wall appears, the wall, along with any other creatures in it, moves 50 feet away from the leviathan. Any Huge or smaller creature inside the wall or whose space the wall enters when it moves must succeed on a DC 24 Strength saving throw or take 27 (5d10) bludgeoning damage. A creature takes this damage no more than once on a turn. At the end of each turn the wall moves, the wall's height is reduced by 50 feet, and the damage creatures take from the wall on subsequent rounds is reduced by 1d10. When the wall reaches 0 feet in height, the effect ends.

A creature caught in the wall can move by swimming. Because of the force of the wave, though, the creature must make a successful DC 24 Strength (Athletics) check to swim at all during that turn.

\par
\DndMonsterSection{Legendary Actions}
Leviathan can take 3 legendary actions, choosing from the options below. Only one legendary action can be used at a time and only at the end of another creature's turn. Leviathan regains spent legendary actions at the start of its turn.

\begin{DndMonsterLegendaryActions}
\DndMonsterLegendaryAction{Slam (Costs 2 Actions)}{The leviathan makes one slam attack.
}
\DndMonsterLegendaryAction{Move}{The leviathan moves up to its speed.
}
\end{DndMonsterLegendaryActions}
\end{multicols}
\end{DndMonster}



\begin{DndMonster}[float*=b,width=\textwidth + 8pt]{Lich}
\begin{multicols}{2}
\DndMonsterType{Medium undead, any evil alignment}
\DndMonsterBasics[
armor-class = {17 (natural armor)},
hit-points = {\DndDice{18d8 + 54}},
speed = {30 ft.}
]
\DndMonsterAbilityScores[
 str = 11,
 dex = 16,
 con = 16,
 int = 20,
 wis = 14,
 cha = 16
]
\DndMonsterDetails[
saving-throws = {Con +10, Int +12, Wis +9},
skills = {Arcana +19, History +12, Insight +9, Perception +9},
damage-resistances = {cold, lightning, necrotic},
damage-immunities = {poison; bludgeoning, piercing, slashing from nonmagical attacks},
senses = {truesight 120 ft., passive perception 19},
languages = {Common plus up to five other languages},
challenge = {'cr': '21', 'lair': '22'},
]
\DndMonsterAction{Legendary Resistance (3/Day)}
If the lich fails a saving throw, it can choose to succeed instead.

\DndMonsterAction{Rejuvenation}
If it has a phylactery, a destroyed lich gains a new body in 1d10 days, regaining all its hit points and becoming active again. The new body appears within 5 feet of the phylactery.

\DndMonsterAction{Turn Resistance}
The lich has advantage on saving throws against any effect that turns undead.

\DndMonsterAction{Spellcasting}
The lich is an 18th-level spellcaster. Its spellcasting ability is Intelligence (spell save DC 20, +12 to hit with spell attacks). The lich has the following wizard spells prepared:
\begin{DndMonsterSpells}
  \DndMonsterSpellLevel{mage hand, prestidigitation, ray of frost}
   \DndMonsterSpellLevel[1][4]{detect magic, magic missile, shield, thunderwave}
   \DndMonsterSpellLevel[2][3]{detect thoughts, invisibility, Melf's acid arrow, mirror image}
   \DndMonsterSpellLevel[3][3]{animate dead, counterspell, dispel magic, fireball}
   \DndMonsterSpellLevel[4][3]{blight, dimension door}
   \DndMonsterSpellLevel[5][3]{cloudkill, scrying}
   \DndMonsterSpellLevel[6][1]{disintegrate, globe of invulnerability}
   \DndMonsterSpellLevel[7][1]{finger of death, plane shift}
   \DndMonsterSpellLevel[8][1]{dominate monster, power word stun}
\end{DndMonsterSpells}
\par
\DndMonsterSection{Actions}
\DndMonsterAction{Paralyzing Touch}
\textsl{Melee Spell Attack:} +12 to hit, reach 5 ft., one creature. \textsl{Hit:} 10 (3d6) cold damage. The target must succeed on a DC 18 Constitution saving throw or be paralyzed for 1 minute. The target can repeat the saving throw at the end of each of its turns, ending the effect on itself on a success.

\par
\DndMonsterSection{Legendary Actions}
Lich can take 3 legendary actions, choosing from the options below. Only one legendary action can be used at a time and only at the end of another creature's turn. Lich regains spent legendary actions at the start of its turn.

\begin{DndMonsterLegendaryActions}
\DndMonsterLegendaryAction{Cantrip}{The lich casts a cantrip.
}
\DndMonsterLegendaryAction{Paralyzing Touch (Costs 2 Actions)}{The lich uses its Paralyzing Touch.
}
\DndMonsterLegendaryAction{Frightening Gaze (Costs 2 Actions)}{The lich fixes its gaze on one creature it can see within 10 feet of it. The target must succeed on a DC 18 Wisdom saving throw against this magic or become frightened for 1 minute. The frightened target can repeat the saving throw at the end of each of its turns, ending the effect on itself on a success. If a target's saving throw is successful or the effect ends for it, the target is immune to the lich's gaze for the next 24 hours.
}
\DndMonsterLegendaryAction{Disrupt Life (Costs 3 Actions)}{Each non-undead creature within 20 feet of the lich must make a DC 18 Constitution saving throw against this magic, taking 21 (6d6) necrotic damage on a failed save, or half as much damage on a successful one.
}
\end{DndMonsterLegendaryActions}
\par \DndMonsterSection{Lair Actions}
On initiative count 20 (losing initiative ties), the lich can take a lair action to cause one of the following magical effects; the lich can't use the same effect two rounds in a row:


\begin{itemize}
\item The lich rolls a d8 and regains a spell slot of that level or lower. If it has no spent spell slots of that level or lower, nothing happens.
\item The lich targets one creature it can see within 30 feet of it. A crackling cord of negative energy tethers the lich to the target. Whenever the lich takes damage, the target must make a DC 18 Constitution saving throw. On a failed save, the lich takes half the damage (rounded down), and the target takes the remaining damage. This tether lasts until initiative count 20 on the next round or until the lich or the target is no longer in the lich's lair.
\item The lich calls forth the spirits of creatures that died in its lair. These apparitions materialize and attack one creature that the lich can see within 60 feet of it. The target must succeed on a DC 18 Constitution saving throw, taking 52 (15d6) necrotic damage on a failed save, or half as much damage on a success. The apparitions then disappear.
\end{itemize}

\end{multicols}
\end{DndMonster}



\begin{DndMonster}[float=t]{Lizardfolk}
\DndMonsterType{Medium humanoid (lizardfolk), neutral}
\DndMonsterBasics[
armor-class = {15 (natural armor)},
hit-points = {\DndDice{4d8 + 4}},
speed = {30 ft., swim 30 ft.}
]
\DndMonsterAbilityScores[
 str = 15,
 dex = 10,
 con = 13,
 int = 7,
 wis = 12,
 cha = 7
]
\DndMonsterDetails[
skills = {Perception +3, Stealth +4, Survival +5},
languages = {Draconic},
challenge = 1/2,
]
\DndMonsterAction{Hold Breath}
The lizardfolk can hold its breath for 15 minutes.

\par
\DndMonsterSection{Actions}
\DndMonsterAction{Multiattack}
The lizardfolk makes two melee attacks, each one with a different weapon.

\DndMonsterAction{Bite}
\textsl{Melee Weapon Attack:} +4 to hit, reach 5 ft., one target. \textsl{Hit:} 5 (1d6 + 2) piercing damage.

\DndMonsterAction{Heavy Club}
\textsl{Melee Weapon Attack:} +4 to hit, reach 5 ft., one target. \textsl{Hit:} 5 (1d6 + 2) bludgeoning damage.

\DndMonsterAction{Javelin}
\textsl{Melee or Ranged Weapon Attack:} +4 to hit, reach 5 ft. or range 30/120 ft., one target. \textsl{Hit:} 5 (1d6 + 2) piercing damage.

\DndMonsterAction{Spiked Shield}
\textsl{Melee Weapon Attack:} +4 to hit, reach 5 ft., one target. \textsl{Hit:} 5 (1d6 + 2) piercing damage.

\end{DndMonster}



\begin{DndMonster}[float=t]{Mage Hunter Golem}
\DndMonsterType{Large construct, U}
\DndMonsterBasics[
armor-class = {20 (natural armor)},
hit-points = {\DndDice{18d10 + 144}},
speed = {30 ft., fly 20 ft.}
]
\DndMonsterAbilityScores[
 str = 24,
 dex = 9,
 con = 26,
 int = 3,
 wis = 11,
 cha = 1
]
\DndMonsterDetails[
saving-throws = {Wis +5, Cha +0},
skills = {Perception +5},
damage-immunities = {poison, psychic; bludgeoning, piercing, slashing from nonmagical weapons that aren't adamantine},
senses = {darkvision 120 ft., truesight 30 ft., passive perception 15},
languages = {understands the languages of its creator but can't speak},
challenge = 15,
]
\DndMonsterAction{Antimagic Cone}
The golem creates an area of antimagic, as in the \textit{antimagic field} spell, in a 60-foot cone originating from it. At the start of each of its turns, the golem decides which way the cone faces and whether the cone is active.

\DndMonsterAction{Immutable Form}
The golem is immune to any spell or effect that would alter its form.

\DndMonsterAction{Magic Absorption}
Whenever the golem is subjected to a damage-dealing spell of 4th level or lower, it takes no damage and instead regains a number of hit points equal to the damage dealt.

\DndMonsterAction{Magic Resistance}
The golem has Advantage and Disadvantage on saving throws against spells and other magical effects.

\DndMonsterAction{Magic Weapons}
The golem's weapon attacks are magical.

\par
\DndMonsterSection{Actions}
\DndMonsterAction{Multiattack}
The golem uses its Antimagic Jolt and makes two melee attacks, or it makes three melee attacks.

\DndMonsterAction{Slam}
\textsl{Melee Weapon Attack:} +12 to hit, reach 5 ft., one target. \textsl{Hit:} 20 (3d8 + 7) bludgeoning damage.

\DndMonsterAction{Claw (Recharge 5\textendash 6)}
\textsl{Melee Weapon Attack:} +12 to hit, reach 5 ft., one target. \textsl{Hit:} 11 (1d8 + 7) bludgeoning damage, and the target must succeed on a DC 16 Dexterity saving throws or have a magic collar placed upon its neck. While the target wears the collar, it is magically prevented from speaking and cannot cast any spell. A successful DC 25 Dexterity check using \textit{thieves' tools} removes the collar.

\DndMonsterAction{Antimagic Jolt}
Each creature within the golem's Antimagic Cone that has the ability to cast a spell of 1st level or higher must make a DC 18 Wisdom saving throws, taking 21 (6d6) psychic damage on a failed save, or half as much damage on a successful one.

\end{DndMonster}



\begin{DndMonster}[float=t]{Magma Landshark}
\DndMonsterType{Huge elemental, U}
\DndMonsterBasics[
armor-class = {18 (natural armor)},
hit-points = {\DndDice{12d12 + 72}},
speed = {50 ft., burrow 50 ft., swim 50 ft. \textit{(lava only)}}
]
\DndMonsterAbilityScores[
 str = 22,
 dex = 11,
 con = 22,
 int = 2,
 wis = 14,
 cha = 5
]
\DndMonsterDetails[
skills = {Perception +7},
damage-immunities = {fire},
senses = {darkvision 60 ft., tremorsense 120 ft., passive perception 17},
challenge = 14,
]
\DndMonsterAction{Searing Presence}
Any creature that starts its turn within 10 feet of the landshark takes 11 (2d10) fire damage.

\DndMonsterAction{Standing Leap}
The landshark's jumping is up to 40 feet and its jumping is up to 20 feet, with or without a running start.

\DndMonsterAction{Illumination}
The landshark sheds bright light in a 15-foot radius and dim light for an additional 15 feet.

\par
\DndMonsterSection{Actions}
\DndMonsterAction{Bite}
\textsl{Melee Weapon Attack:} +11 to hit, reach 5 ft., one target. \textsl{Hit:} 45 (6d12 + 6) piercing damage, and 14 (4d6) fire damage. If the target is a Large or smaller creature, it must succeed on a DC 18 Dexterity saving throws or be swallowed by the landshark. A swallowed creature is blinded and restrained, it has cover against attacks and other effects outside the landshark, and it takes 42 (12d6) fire damage at the start of each of the landshark's turns.

If the landshark takes 30 damage or more on a single turn from a creature inside it, the landshark must succeed on a DC 23 Constitution saving throws at the end of that turn or regurgitate all swallowed creatures, which fall prone in a space within 10 feet of the landshark. If the landshark dies, a swallowed creature is no longer restrained by it and can escape from the corpse by using 15 feet of movement, exiting prone.

\DndMonsterAction{Deadly Leap}
If the landshark jumping at least 15 feet as part of its movement, it can then use this action to land on its feet in a space that contains one or more other creatures. Each of those creatures must succeed on a DC 18 Strength or Dexterity saving throws (target's choice) or be knocked prone and take 27 (6d6 + 6) bludgeoning damage plus 21 (6d6) fire damage. On a successful save, the creature takes only half the damage, isn't knocked prone, and is pushed 5 feet out of the landshark's space into an unoccupied space of the creature's choice. If no unoccupied space is within range, the creature instead falls prone in the landshark's space.

\DndMonsterAction{Lava Pool (Recharge 6)}
If the landshark is on solid ground, it melts the area around it, creating a pool of lava in a 10-foot radius. Any creature that starts its turn in the area or enters it for the first time on a turn takes 44 (8d10) fire damage.

\end{DndMonster}



\begin{DndMonster}[float=t]{Mammoth}
\DndMonsterType{Huge beast, unaligned}
\DndMonsterBasics[
armor-class = {13 (natural armor)},
hit-points = {\DndDice{11d12 + 55}},
speed = {40 ft.}
]
\DndMonsterAbilityScores[
 str = 24,
 dex = 9,
 con = 21,
 int = 3,
 wis = 11,
 cha = 6
]
\DndMonsterDetails[
challenge = 6,
]
\DndMonsterAction{Trampling Charge}
If the mammoth moves at least 20 feet straight toward a creature and then hits it with a gore attack on the same turn, that target must succeed on a DC 18 Strength saving throw or be knocked prone. If the target is prone, the mammoth can make one stomp attack against it as a bonus action.

\par
\DndMonsterSection{Actions}
\DndMonsterAction{Gore}
\textsl{Melee Weapon Attack:} +10 to hit, reach 10 ft., one target. \textsl{Hit:} 25 (4d8 + 7) piercing damage.

\DndMonsterAction{Stomp}
\textsl{Melee Weapon Attack:} +10 to hit, reach 5 ft., one prone creature. \textsl{Hit:} 29 (4d10 + 7) bludgeoning damage.

\end{DndMonster}



\begin{DndMonster}[float=t]{Marilith}
\DndMonsterType{Large fiend (demon), chaotic evil}
\DndMonsterBasics[
armor-class = {18 (natural armor)},
hit-points = {\DndDice{18d10 + 90}},
speed = {40 ft.}
]
\DndMonsterAbilityScores[
 str = 18,
 dex = 20,
 con = 20,
 int = 18,
 wis = 16,
 cha = 20
]
\DndMonsterDetails[
saving-throws = {Str +9, Con +10, Wis +8, Cha +10},
damage-resistances = {cold, fire, lightning; bludgeoning, piercing, slashing from nonmagical attacks},
damage-immunities = {poison},
senses = {truesight 120 ft., passive perception 13},
languages = {Abyssal, telepathy 120 ft.},
challenge = 16,
]
\DndMonsterAction{Magic Resistance}
The marilith has advantage on saving throws against spells and other magical effects.

\DndMonsterAction{Magic Weapons}
The marilith's weapon attacks are magical.

\DndMonsterAction{Reactive}
The marilith can take one reaction on every turn in combat.

\par
\DndMonsterSection{Actions}
\DndMonsterAction{Multiattack}
The marilith can make seven attacks: six with its longswords and one with its tail.

\DndMonsterAction{Longsword}
\textsl{Melee Weapon Attack:} +9 to hit, reach 5 ft., one target. \textsl{Hit:} 13 (2d8 + 4) slashing damage.

\DndMonsterAction{Tail}
\textsl{Melee Weapon Attack:} +9 to hit, reach 10 ft., one creature. \textsl{Hit:} 15 (2d10 + 4) bludgeoning damage. If the target is Medium or smaller, it is grappled (escape DC 19). Until this grapple ends, the target is restrained, the marilith can automatically hit the target with its tail, and the marilith can't make tail attacks against other targets.

\DndMonsterAction{Teleport}
The marilith magically teleports, along with any equipment it is wearing or carrying, up to 120 feet to an unoccupied space it can see.

\par
\DndMonsterSection{Reactions}
\DndMonsterAction{Parry}
The marilith adds 5 to its AC against one melee attack that would hit it. To do so, the marilith must see the attacker and be wielding a melee weapon.

\end{DndMonster}



\begin{DndMonster}[float*=b,width=\textwidth + 8pt]{Master Adranach}
\begin{multicols}{2}
\DndMonsterType{Gargantuan construct, U}
\DndMonsterBasics[
armor-class = {19 (natural armor)},
hit-points = {\DndDice{20d20 + 80}},
speed = {60 ft., fly 120 ft.}
]
\DndMonsterAbilityScores[
 str = 25,
 dex = 16,
 con = 19,
 int = 10,
 wis = 14,
 cha = 11
]
\DndMonsterDetails[
saving-throws = {Str +14, Dex +10, Wis +9},
skills = {Athletics +14, Perception +16},
damage-resistances = {damage from spells},
damage-immunities = {poison; bludgeoning, piercing, slashing from non-magical attacks},
senses = {darkvision 360 ft., passive perception 26},
languages = {the languages spoken by its creator},
challenge = 21,
]
\DndMonsterAction{Arcane Leak}
When the adranach is reduced to half its hit point maximum (145 hit points), its mithral mask cracks and its arcane form begins to waver, creating a field of unstable magic around it. Any creature that starts its turn within 20 feet of the adranach or enters that area for the first time on a turn takes 21 (6d6) radiant damage. Additionally, if a creature casts a spell within this area, it must make a DC 18 ability check using its spellcasting ability. On a failure, the spell backfires, consuming the spell slot and dealing 7 (2d6) force damage to the caster for each level of the spell slot that was consumed.

\DndMonsterAction{Immutable Form}
The adranach is immune to any spell or effect that would alter its form.

\DndMonsterAction{Magic Resistance}
The adranach has Advantage and Disadvantage on saving throws against spells and other magical effects.

\DndMonsterAction{Magic Weapons}
The adranach's weapon attacks are magical.

\DndMonsterAction{Colossal Build}
The adranach's carrying capacity is doubled, as is the weight it can push, drag, or lift.

\DndMonsterAction{Slipstream}
The adranach can cast \textit{invisibility} on itself at will. While it is invisible, its speed is doubled.

\DndMonsterAction{Legendary Resistance (3/Day)}
If the adranach fails a saving throws, it can choose to succeed instead.

\par
\DndMonsterSection{Actions}
\DndMonsterAction{Multiattack}
The adranach makes three attacks: two with its claws and one with its tail.

\DndMonsterAction{Claws}
\textsl{Melee Weapon Attack:} +14 to hit, one target, reach 5 ft. \textsl{Hit:} 33 (4d12 + 7) slashing damage.

\DndMonsterAction{Tail}
\textsl{Melee Weapon Attack:} +14 to hit, one target, reach 15 ft. \textsl{Hit:} 20 (2d12 + 7) bludgeoning damage.

\DndMonsterAction{Force Bolts (Recharge 5\textendash 6)}
The adranach targets up to six creatures it can see within 120 feet of it. Each target must succeed on a DC 22 Dexterity saving throws or take 35 (10d6) force damage.

\par
\DndMonsterSection{Legendary Actions}
Master Adranach can take 3 legendary actions, choosing from the options below. Only one legendary action can be used at a time and only at the end of another creature's turn. Master Adranach regains spent legendary actions at the start of its turn.

\begin{DndMonsterLegendaryActions}
\DndMonsterLegendaryAction{Claw Attack}{The adranach makes one claw attack.
}
\DndMonsterLegendaryAction{Vanish (Costs 2 Actions)}{The adranach casts \textit{invisibility} and moves up to its walking or flying speed. (This movement doesn't benefit from its Slipstream trait.)
}
\DndMonsterLegendaryAction{Starry Burst (Costs 3 Actions)}{A burst of astral energy flares from the adranach's body. Each creature within 60 feet of it must make a DC 22 Constitution saving throws. On a failed save, a creature takes 27 (5d10) force damage and is blinded until the end of its next turn. On a successful save, the creature takes half as much damage and is not blinded. The adranach then turns invisible until the start of its next turn or until it makes an attack.
}
\end{DndMonsterLegendaryActions}
\end{multicols}
\end{DndMonster}



\begin{DndMonster}[float=t]{Medusa}
\DndMonsterType{Medium monstrosity, lawful evil}
\DndMonsterBasics[
armor-class = {15 (natural armor)},
hit-points = {\DndDice{17d8 + 51}},
speed = {30 ft.}
]
\DndMonsterAbilityScores[
 str = 10,
 dex = 15,
 con = 16,
 int = 12,
 wis = 13,
 cha = 15
]
\DndMonsterDetails[
skills = {Deception +5, Insight +4, Perception +4, Stealth +5},
senses = {darkvision 60 ft., passive perception 14},
languages = {Common},
challenge = 6,
]
\DndMonsterAction{Petrifying Gaze}
When a creature that can see the medusa's eyes starts its turn within 30 feet of the medusa, the medusa can force it to make a DC 14 Constitution saving throw if the medusa isn't incapacitated and can see the creature. If the saving throw fails by 5 or more, the creature is instantly petrified. Otherwise, a creature that fails the save begins to turn to stone and is restrained. The restrained creature must repeat the saving throw at the end of its next turn, becoming petrified on a failure or ending the effect on a success. The petrification lasts until the creature is freed by the  \textit{greater restoration} spell or other magic.

Unless Surprise, a creature can avert its eyes to avoid the saving throw at the start of its turn. If the creature does so, it can't see the medusa until the start of its next turn, when it can avert its eyes again. If the creature looks at the medusa in the meantime, it must immediately make the save.

If the medusa sees itself reflected on a polished surface within 30 feet of it and in an area of bright light, the medusa is, due to its curse, affected by its own gaze.

\par
\DndMonsterSection{Actions}
\DndMonsterAction{Multiattack}
The medusa makes either three melee attacks—one with its snake hair and two with its shortsword—or two ranged attacks with its longbow.

\DndMonsterAction{Snake Hair}
\textsl{Melee Weapon Attack:} +5 to hit, reach 5 ft., one creature. \textsl{Hit:} 4 (1d4 + 2) piercing damage plus 14 (4d6) poison damage.

\DndMonsterAction{Shortsword}
\textsl{Melee Weapon Attack:} +5 to hit, reach 5 ft., one target. \textsl{Hit:} 5 (1d6 + 2) piercing damage.

\DndMonsterAction{Longbow}
\textsl{Ranged Weapon Attack:} +5 to hit, range 150/600 ft., one target. \textsl{Hit:} 6 (1d8 + 2) piercing damage plus 7 (2d6) poison damage.

\end{DndMonster}



\begin{DndMonster}[float=t]{Merfolk}
\DndMonsterType{Medium humanoid (merfolk), neutral}
\DndMonsterBasics[
armor-class = {11},
hit-points = {\DndDice{2d8 + 2}},
speed = {10 ft., swim 40 ft.}
]
\DndMonsterAbilityScores[
 str = 10,
 dex = 13,
 con = 12,
 int = 11,
 wis = 11,
 cha = 12
]
\DndMonsterDetails[
skills = {Perception +2},
languages = {Aquan, Common},
challenge = 1/8,
]
\DndMonsterAction{Amphibious}
The merfolk can breathe air and water.

\par
\DndMonsterSection{Actions}
\DndMonsterAction{Spear}
\textsl{Melee or Ranged Weapon Attack:} +2 to hit, reach 5 ft. or range 20/60 ft., one target. \textsl{Hit:} 3 (1d6) piercing damage, or 4 (1d8) piercing damage if used with two hands to make a melee attack.

\end{DndMonster}



\begin{DndMonster}[float=t]{Merrow}
\DndMonsterType{Large monstrosity, chaotic evil}
\DndMonsterBasics[
armor-class = {13 (natural armor)},
hit-points = {\DndDice{6d10 + 12}},
speed = {10 ft., swim 40 ft.}
]
\DndMonsterAbilityScores[
 str = 18,
 dex = 10,
 con = 15,
 int = 8,
 wis = 10,
 cha = 9
]
\DndMonsterDetails[
senses = {darkvision 60 ft., passive perception 10},
languages = {Abyssal, Aquan},
challenge = 2,
]
\DndMonsterAction{Amphibious}
The merrow can breathe air and water.

\par
\DndMonsterSection{Actions}
\DndMonsterAction{Multiattack}
The merrow makes two attacks: one with its bite and one with its claws or harpoon.

\DndMonsterAction{Bite}
\textsl{Melee Weapon Attack:} +6 to hit, reach 5 ft., one target. \textsl{Hit:} 8 (1d8 + 4) piercing damage.

\DndMonsterAction{Claws}
\textsl{Melee Weapon Attack:} +6 to hit, reach 5 ft., one target. \textsl{Hit:} 9 (2d4 + 4) slashing damage.

\DndMonsterAction{Harpoon}
\textsl{Melee or Ranged Weapon Attack:} +6 to hit, reach 5 ft. or range 20/60 ft., one target. \textsl{Hit:} 11 (2d6 + 4) piercing damage. If the target is a Huge or smaller creature, it must succeed on a Strength contest against the merrow or be pulled up to 20 feet toward the merrow.

\end{DndMonster}



\begin{DndMonster}[float=t]{Mind Flayer}
\DndMonsterType{Medium aberration, lawful evil}
\DndMonsterBasics[
armor-class = {15 (\textit{breastplate})},
hit-points = {\DndDice{13d8 + 13}},
speed = {30 ft.}
]
\DndMonsterAbilityScores[
 str = 11,
 dex = 12,
 con = 12,
 int = 19,
 wis = 17,
 cha = 17
]
\DndMonsterDetails[
saving-throws = {Int +7, Wis +6, Cha +6},
skills = {Arcana +7, Deception +6, Insight +6, Perception +6, Persuasion +6, Stealth +4},
senses = {darkvision 120 ft., passive perception 16},
languages = {Deep Speech, Undercommon, telepathy 120 ft.},
challenge = 7,
]
\DndMonsterAction{Magic Resistance}
The mind flayer has advantage on saving throws against spells and other magical effects.

\DndMonsterAction{Innate Spellcasting (Psionics)}
The mind flayer's innate spellcasting ability is Intelligence (spell save DC 15). It can innately cast the following spells, requiring no components:
\begin{DndMonsterSpells}
  \DndInnateSpellLevel{detect thoughts, levitate}
  \DndInnateSpellLevel[1e]{dominate monster, plane shift}
\end{DndMonsterSpells}
\par
\DndMonsterSection{Actions}
\DndMonsterAction{Tentacles}
\textsl{Melee Weapon Attack:} +7 to hit, reach 5 ft., one creature. \textsl{Hit:} 15 (2d10 + 4) psychic damage. If the target is Medium or smaller, it is grappled (escape DC 15) and must succeed on a DC 15 Intelligence saving throw or be stunned until this grapple ends.

\DndMonsterAction{Extract Brain}
\textsl{Melee Weapon Attack:} +7 to hit, reach 5 ft., one incapacitated humanoid grappled by the mind flayer. \textsl{Hit:} The target takes 55 (10d10) piercing damage. If this damage reduces the target to 0 hit points, the mind flayer kills the target by extracting and devouring its brain.

\DndMonsterAction{Mind Blast (Recharge 5\textendash 6)}
The mind flayer magically emits psychic energy in a 60-foot cone. Each creature in that area must succeed on a DC 15 Intelligence saving throw or take 22 (4d8 + 4) psychic damage and be stunned for 1 minute. A creature can repeat the saving throw at the end of each of its turns, ending the effect on itself on a success.

\end{DndMonster}



\begin{DndMonster}[float=t]{Mountain Goat}
\DndMonsterType{Medium beast, unaligned}
\DndMonsterBasics[
armor-class = {11},
hit-points = {\DndDice{2d8 + 4}},
speed = {40 ft., climb 30 ft.}
]
\DndMonsterAbilityScores[
 str = 14,
 dex = 12,
 con = 14,
 int = 2,
 wis = 10,
 cha = 5
]
\DndMonsterDetails[
challenge = 1/8,
]
\DndMonsterAction{Charge}
If the goat moves at least 20 feet straight toward a target and then hits it with a ram attack on the same turn, the target takes an extra 3 (1d6) bludgeoning damage. If the target is a creature, it must succeed on a DC 12 Strength saving throw or be knocked prone.

\DndMonsterAction{Sure-Footed}
The goat has advantage on Strength and Dexterity saving throws made against effects that would knock it prone.

\par
\DndMonsterSection{Actions}
\DndMonsterAction{Ram}
\textsl{Melee Weapon Attack:} +4 to hit, reach 5 ft., one target. \textsl{Hit:} 5 (1d6 + 2) bludgeoning damage.

\end{DndMonster}


\clearpage

\begin{DndMonster}[float=t]{Myconid Adult}
\DndMonsterType{Medium plant, lawful neutral}
\DndMonsterBasics[
armor-class = {12 (natural armor)},
hit-points = {\DndDice{4d8 + 4}},
speed = {20 ft.}
]
\DndMonsterAbilityScores[
 str = 10,
 dex = 10,
 con = 12,
 int = 10,
 wis = 13,
 cha = 7
]
\DndMonsterDetails[
senses = {darkvision 120 ft., passive perception 11},
challenge = 1/2,
]
\DndMonsterAction{Distress Spores}
When the myconid takes damage, all other myconids within 240 feet of it can sense its pain.

\DndMonsterAction{Sun Sickness}
While in sunlight, the myconid has disadvantage on ability checks, attack rolls, and saving throws. The myconid dies if it spends more than 1 hour in direct sunlight.

\par
\DndMonsterSection{Actions}
\DndMonsterAction{Fist}
\textsl{Melee Weapon Attack:} +2 to hit, reach 5 ft., one target. \textsl{Hit:} 5 (2d4) bludgeoning damage plus 5 (2d4) poison damage.

\DndMonsterAction{Pacifying Spores (3/Day)}
The myconid ejects spores at one creature it can see within 5 feet of it. The target must succeed on a DC 11 Constitution saving throw or be stunned for 1 minute. The target can repeat the saving throw at the end of each of its turns, ending the effect on itself on a success.

\DndMonsterAction{Rapport Spores}
A 20-foot radius of spores extends from the myconid. These spores can go around corners and affect only creatures with an Intelligence of 2 or higher that aren't undead, constructs, or elementals. Affected creatures can communicate telepathically with one another while they are within 30 feet of each other. The effect lasts for 1 hour.

\end{DndMonster}



\begin{DndMonster}[float=t]{Myconid Sovereign}
\DndMonsterType{Large plant, lawful neutral}
\DndMonsterBasics[
armor-class = {13 (natural armor)},
hit-points = {\DndDice{8d10 + 16}},
speed = {30 ft.}
]
\DndMonsterAbilityScores[
 str = 12,
 dex = 10,
 con = 14,
 int = 13,
 wis = 15,
 cha = 10
]
\DndMonsterDetails[
senses = {darkvision 120 ft., passive perception 12},
challenge = 2,
]
\DndMonsterAction{Distress Spores}
When the myconid takes damage, all other myconids within 240 feet of it can sense its pain.

\DndMonsterAction{Sun Sickness}
While in sunlight, the myconid has disadvantage on ability checks, attack rolls, and saving throws. The myconid dies if it spends more than 1 hour in direct sunlight.

\par
\DndMonsterSection{Actions}
\DndMonsterAction{Multiattack}
The myconid uses either its Hallucination Spores or its Pacifying Spores, then makes a fist attack.

\DndMonsterAction{Fist}
\textsl{Melee Weapon Attack:} +3 to hit, reach 5 ft., one target. \textsl{Hit:} 8 (3d4 + 1) bludgeoning damage plus 7 (3d4) poison damage.

\DndMonsterAction{Animating Spores (3/Day)}
The myconid targets one corpse of a humanoid or a Large or smaller beast within 5 feet of it and releases spores at the corpse. In 24 hours, the corpse rises as a spore servant. The corpse stays animated for 1d4 + 1 weeks or until destroyed, and it can't be animated again in this way.

\DndMonsterAction{Hallucination Spores}
The myconid ejects spores at one creature it can see within 5 feet of it. The target must succeed on a DC 12 Constitution saving throw or be poisoned for 1 minute. The poisoned target is incapacitated while it hallucinates. The target can repeat the saving throw at the end of each of its turns, ending the effect on itself on a success.

\DndMonsterAction{Pacifying Spores}
The myconid ejects spores at one creature it can see within 5 feet of it. The target must succeed on a DC 12 Constitution saving throw or be stunned for 1 minute. The target can repeat the saving throw at the end of each of its turns, ending the effect on itself on a success.

\DndMonsterAction{Rapport Spores}
A 30-foot radius of spores extends from the myconid. These spores can go around corners and affect only creatures with an Intelligence of 2 or higher that aren't undead, constructs, or elementals. Affected creatures can communicate telepathically with one another while they are within 30 feet of each other. The effect lasts for 1 hour.

\end{DndMonster}



\begin{DndMonster}[float*=b,width=\textwidth + 8pt]{Night Hag}
\begin{multicols}{2}
\DndMonsterType{Medium fiend, neutral evil}
\DndMonsterBasics[
armor-class = {17 (natural armor)},
hit-points = {\DndDice{15d8 + 45}},
speed = {30 ft.}
]
\DndMonsterAbilityScores[
 str = 18,
 dex = 15,
 con = 16,
 int = 16,
 wis = 14,
 cha = 16
]
\DndMonsterDetails[
skills = {Deception +6, Insight +5, Perception +5, Stealth +5},
damage-resistances = {cold, fire; bludgeoning, piercing, slashing from nonmagical attacks that aren't silvered},
senses = {darkvision 120 ft., passive perception 16},
languages = {Abyssal, Common, Infernal, Primordial},
challenge = {'cr': '5', 'coven': '7'},
]
\DndMonsterAction{Magic Resistance}
The hag has advantage on saving throws against spells and other magical effects.

\DndMonsterAction{Night Hag Items}
A night hag carries two very rare magic items that she must craft for herself. If either object is lost, the night hag will go to great lengths to retrieve it, as creating a new tool takes time and effort.

Heartstone: This lustrous black gem allows a night hag to become ethereal while it is in her possession. The touch of a heartstone also cures any disease. Crafting a heartstone takes 30 days.

Soul Bag: When an evil humanoid dies as a result of a night hag's Nightmare Haunting, the hag catches the soul in this black sack made of stitched flesh. A soul bag can hold only one evil soul at a time, and only the night hag who crafted the bag can catch a soul with it. Crafting a soul bag takes 7 days and a humanoid sacrifice (whose flesh is used to make the bag).

\DndMonsterAction{Innate Spellcasting}
The hag's innate spellcasting ability is Charisma (spell save DC 14, +6 to hit with spell attacks). She can innately cast the following spells, requiring no material components:
\begin{DndMonsterSpells}
  \DndInnateSpellLevel{detect magic, magic missile}
  \DndInnateSpellLevel[2e]{plane shift, ray of enfeeblement, sleep}
\end{DndMonsterSpells}
\par
\DndMonsterSection{Actions}
\DndMonsterAction{Claws (Hag Form Only)}
\textsl{Melee Weapon Attack:} +7 to hit, reach 5 ft., one target. \textsl{Hit:} 13 (2d8 + 4) slashing damage.

\DndMonsterAction{Change Shape}
The hag magically polymorphs into a Small or Medium female humanoid, or back into her true form. Her statistics are the same in each form. Any equipment she is wearing or carrying isn't transformed. She reverts to her true form if she dies.

\DndMonsterAction{Etherealness}
The hag magically enters the Ethereal Plane from the Material Plane, or vice versa. To do so, the hag must have a heartstone in her possession.

\DndMonsterAction{Nightmare Haunting (1/Day)}
While on the Ethereal Plane, the hag magically touches a sleeping humanoid on the Material Plane. A \textit{protection from evil and good} spell cast on the target prevents this contact, as does a magic circle. As long as the contact persists, the target has dreadful visions. If these visions last for at least 1 hour, the target gains no benefit from its rest, and its hit point maximum is reduced by 5 (1d10). If this effect reduces the target's hit point maximum to 0, the target dies, and if the target was evil, its soul is trapped in the hag's soul bag. The reduction to the target's hit point maximum lasts until removed by the  \textit{greater restoration} spell or similar magic.

\end{multicols}
\end{DndMonster}



\begin{DndMonster}[float=t]{Ogre}
\DndMonsterType{Large giant, chaotic evil}
\DndMonsterBasics[
armor-class = {11 (\textit{hide armor})},
hit-points = {\DndDice{7d10 + 21}},
speed = {40 ft.}
]
\DndMonsterAbilityScores[
 str = 19,
 dex = 8,
 con = 16,
 int = 5,
 wis = 7,
 cha = 7
]
\DndMonsterDetails[
senses = {darkvision 60 ft., passive perception 8},
languages = {Common, Giant},
challenge = 2,
]
\par
\DndMonsterSection{Actions}
\DndMonsterAction{Greatclub}
\textsl{Melee Weapon Attack:} +6 to hit, reach 5 ft., one target. \textsl{Hit:} 13 (2d8 + 4) bludgeoning damage.

\DndMonsterAction{Javelin}
\textsl{Melee or Ranged Weapon Attack:} +6 to hit, reach 5 ft. or range 30/120 ft., one target. \textsl{Hit:} 11 (2d6 + 4) piercing damage.

\end{DndMonster}



\begin{DndMonster}[float=t]{Oni}
\DndMonsterType{Large giant, lawful evil}
\DndMonsterBasics[
armor-class = {16 (\textit{chain mail})},
hit-points = {\DndDice{13d10 + 39}},
speed = {30 ft., fly 30 ft.}
]
\DndMonsterAbilityScores[
 str = 19,
 dex = 11,
 con = 16,
 int = 14,
 wis = 12,
 cha = 15
]
\DndMonsterDetails[
saving-throws = {Dex +3, Con +6, Wis +4, Cha +5},
skills = {Arcana +5, Deception +8, Perception +4},
senses = {darkvision 60 ft., passive perception 14},
languages = {Common, Giant},
challenge = 7,
]
\DndMonsterAction{Magic Weapons}
The oni's weapon attacks are magical.

\DndMonsterAction{Regeneration}
The oni regains 10 hit points at the start of its turn if it has at least 1 hit point.

\DndMonsterAction{Innate Spellcasting}
The oni's innate spellcasting ability is Charisma (spell save DC 13). The oni can innately cast the following spells, requiring no material components:
\begin{DndMonsterSpells}
  \DndInnateSpellLevel{darkness, invisibility}
  \DndInnateSpellLevel[1e]{charm person, cone of cold, gaseous form, sleep}
\end{DndMonsterSpells}
\par
\DndMonsterSection{Actions}
\DndMonsterAction{Multiattack}
The oni makes two attacks, either with its claws or its glaive.

\DndMonsterAction{Claw (Oni Form Only)}
\textsl{Melee Weapon Attack:} +7 to hit, reach 5 ft., one target. \textsl{Hit:} 8 (1d8 + 4) slashing damage.

\DndMonsterAction{Glaive}
\textsl{Melee Weapon Attack:} +7 to hit, reach 10 ft., one target. \textsl{Hit:} 15 (2d10 + 4) slashing damage, or 9 (1d10 + 4) slashing damage in Small or Medium form.

\DndMonsterAction{Change Shape}
The oni magically polymorphs into a Small or Medium humanoid, into a Large giant, or back into its true form. Other than its size, its statistics are the same in each form. The only equipment that is transformed is its glaive, which shrinks so that it can be wielded in humanoid form. If the oni dies, it reverts to its true form, and its glaive reverts to its normal size.

\end{DndMonster}



\begin{DndMonster}[float=t]{Orc}
\DndMonsterType{Medium humanoid (orc), chaotic evil}
\DndMonsterBasics[
armor-class = {13 (\textit{hide armor})},
hit-points = {\DndDice{2d8 + 6}},
speed = {30 ft.}
]
\DndMonsterAbilityScores[
 str = 16,
 dex = 12,
 con = 16,
 int = 7,
 wis = 11,
 cha = 10
]
\DndMonsterDetails[
skills = {Intimidation +2},
senses = {darkvision 60 ft., passive perception 10},
languages = {Common, Orc},
challenge = 1/2,
]
\DndMonsterAction{Aggressive}
As a bonus action, the orc can move up to its speed toward a hostile creature that it can see.

\par
\DndMonsterSection{Actions}
\DndMonsterAction{Greataxe}
\textsl{Melee Weapon Attack:} +5 to hit, reach 5 ft., one target. \textsl{Hit:} 9 (1d12 + 3) slashing damage.

\DndMonsterAction{Javelin}
\textsl{Melee or Ranged Weapon Attack:} +5 to hit, reach 5 ft. or range 30/120 ft., one target. \textsl{Hit:} 6 (1d6 + 3) piercing damage.

\end{DndMonster}



\begin{DndMonster}[float=t]{Owlbear}
\DndMonsterType{Large monstrosity, unaligned}
\DndMonsterBasics[
armor-class = {13 (natural armor)},
hit-points = {\DndDice{7d10 + 21}},
speed = {40 ft.}
]
\DndMonsterAbilityScores[
 str = 20,
 dex = 12,
 con = 17,
 int = 3,
 wis = 12,
 cha = 7
]
\DndMonsterDetails[
skills = {Perception +3},
senses = {darkvision 60 ft., passive perception 13},
challenge = 3,
]
\DndMonsterAction{Keen Sight and Smell}
The owlbear has advantage on Wisdom (Perception) checks that rely on sight or smell.

\par
\DndMonsterSection{Actions}
\DndMonsterAction{Multiattack}
The owlbear makes two attacks: one with its beak and one with its claws.

\DndMonsterAction{Beak}
\textsl{Melee Weapon Attack:} +7 to hit, reach 5 ft., one creature. \textsl{Hit:} 10 (1d10 + 5) piercing damage.

\DndMonsterAction{Claws}
\textsl{Melee Weapon Attack:} +7 to hit, reach 5 ft., one target. \textsl{Hit:} 14 (2d8 + 5) slashing damage.

\end{DndMonster}



\begin{DndMonster}[float=t]{Percival de Rolo}
\DndMonsterType{Medium humanoid (human), U}
\DndMonsterBasics[
armor-class = {18 (studded leather)},
hit-points = {\DndDice{20d10 + 40 + 40}},
speed = {30 ft.}
]
\DndMonsterAbilityScores[
 str = 12,
 dex = 22,
 con = 14,
 int = 20,
 wis = 16,
 cha = 14
]
\DndMonsterDetails[
saving-throws = {Str +7, Con +8},
skills = {Acrobatics +12, History +11, Perception +9, Persuasion +8},
damage-resistances = {lightning},
languages = {Celestial, Common, Elvish},
challenge = 18,
]
\DndMonsterAction{Action Surge (2/resting)}
Once on his turn, Percy can take one additional action.

\DndMonsterAction{Cabal's Ruin}
Cabal's Ruin grants Percy Advantage and Disadvantage on saving throws against spells or magical effects. Additionally, Cabal's Ruin has 10 charges. When Percy hits with a weapon attack, he can expend up to 10 charges to have the attack deal 1d6 extra lightning damage per charge spent.

\DndMonsterAction{Indomitable (3/resting)}
Percy can reroll a failed saving throws.

\DndMonsterAction{Misfire}
When Percy rolls a 2 or lower on the d20 when making an attack with Animus or Bad News, the attack misses and he can't use the weapon again until he spends an action to repair it. On a misfire with Animus, he takes 7 (2d6) psychic damage.

\DndMonsterAction{No Mercy}
Percy's weapon attacks score a critical hit on a roll of 19-20.

\DndMonsterAction{Ranged Master}
Percy's ranged weapon attacks have a +2 bonus to hit (included in attacks), ignore cover and cover, and have no range penalty. Additionally, Percy can take a −5 penalty to any ranged attack roll to gain a +10 bonus to the damage roll.

\DndMonsterAction{Special Equipment}
Percy wears a \textit{ring of lightning resistance} and Cabal's Ruin. He wields Animus and Bad News, which are both +1 weapons.

\DndMonsterAction{Quick Draw}
Percy adds his proficiency bonus to initiative rolls, giving him a +12 to initiative rolls.

\DndMonsterAction{Toughness}
Percy has an additional 40 hit points.

\par
\DndMonsterSection{Actions}
\DndMonsterAction{Multiattack}
Percy makes four ranged attacks with Animus, or two ranged attacks with Bad News. When Percy hits with an attack with either weapon, he can choose one of the following additional effects:


\begin{itemize}
\item{Deadeye.}
Percy gains Advantage and Disadvantage on his next attack roll this turn made against the same target.

\item{Disarm.}
The target must make a DC 20 Strength saving throws. On a failure, the target drops one item it is holding of Percy's choice.

\end{itemize}

\DndMonsterAction{Animus}
\textsl{Ranged Weapon Attack:} +15 to hit, range 320 ft., one target. \textsl{Hit:} 11 (1d10 + 6) piercing damage and 3 (1d6) psychic damage.

\DndMonsterAction{Bad News}
\textsl{Ranged Weapon Attack:} +15 to hit, range 800 ft., one target. \textsl{Hit:} 19 (2d12 + 6) piercing damage.

\par
\DndMonsterSection{Reactions}
\DndMonsterAction{Absorb Magic (Recharges after a resting)}
When Percy is targeted by an enemy's spell, he can absorb a portion of the spell's energy into Cabal's Ruin. The spell affects him normally, but the cloak gains a number of charges equal to the level of the spell. When he does so, he also gains resistance to one type of damage dealt by the spell until the end of his next turn.

\end{DndMonster}



\begin{DndMonster}[float=t]{Pike Trickfoot}
\DndMonsterType{Small humanoid (gnome), U}
\DndMonsterBasics[
armor-class = {23 (plate of the dawnmartyr)},
hit-points = {\DndDice{20d8 + 80}},
speed = {25 ft.}
]
\DndMonsterAbilityScores[
 str = 19,
 dex = 12,
 con = 18,
 int = 14,
 wis = 20,
 cha = 14
]
\DndMonsterDetails[
saving-throws = {Wis +11, Cha +8},
skills = {Athletics +10, Perception +11, Persuasion +8, Religion +8, Stealth +7},
damage-resistances = {fire; bludgeoning, piercing, slashing from nonmagical attacks},
senses = {darkvision 60 ft., passive perception 21},
languages = {Common, Dwarvish, Gnomish, Undercommon},
challenge = 17,
]
\DndMonsterAction{Gnome Cunning}
Pike has Advantage and Disadvantage on Intelligence, Wisdom, and Charisma saving throws against magic.

\DndMonsterAction{Special Equipment}
Pike wields a \textit{mace of disruption} and wears \textit{gauntlets of ogre power}, \textit{boots of speed}, and the Plate of the Dawnmartyr.

\DndMonsterAction{Dawnmartyr's Grace (1/resting)}
When Pike is reduced to 0 hit points while wearing the Plate of the Dawnmartyr, she can choose to have a blast of healing flame surround her. She regains 1 hit point and is lifted to her feet. Each enemy within 10 feet of her takes 7 (2d6) fire damage or radiant damage (her choice).

\DndMonsterAction{Spellcasting}
Pike is a 20th-level spellcaster. Her spellcasting ability is Wisdom (spell save DC 19, +11 to hit with spell attacks). She has the following cleric spells prepared:
\begin{DndMonsterSpells}
  \DndMonsterSpellLevel{light, sacred flame, spare the dying, thaumaturgy}
   \DndMonsterSpellLevel[1][4]{command, cure wounds, guiding bolt, healing word, sanctuary, shield of faith}
   \DndMonsterSpellLevel[2][3]{aid, lesser restoration, spiritual weapon}
   \DndMonsterSpellLevel[3][3]{glyph of warding, mass healing word, protection from energy, revivify, speak with dead}
   \DndMonsterSpellLevel[4][3]{death ward, freedom of movement, guardian of faith}
   \DndMonsterSpellLevel[5][3]{flame strike, greater restoration, hallow, mass cure wounds, scrying}
   \DndMonsterSpellLevel[6][2]{heal, heroes' feast, planar ally}
   \DndMonsterSpellLevel[7][2]{resurrection}
   \DndMonsterSpellLevel[8][1]{holy aura}
\end{DndMonsterSpells}
\par
\DndMonsterSection{Actions}
\DndMonsterAction{Mace of Disruption}
\textsl{Melee Weapon Attack:} +10 to hit, reach 5 ft., one target. \textsl{Hit:} 17 (3d8 + 4) bludgeoning damage, plus an extra 7 (2d6) radiant damage if the target is a fiend or an undead. If the target has 25 hit points or fewer after taking this damage, it must succeed on a DC 15 Wisdom saving throws or be destroyed. On a successful save, the creature becomes frightened of Pike until the end of her next turn.

\DndMonsterAction{Guiding Bolt (5th Level)}
\textsl{Ranged Spell Attack:} +11 to hit, range 120 ft., one target. \textsl{Hit:} 28 (8d6) radiant damage, and the next attack roll made against the target before the end of Pike's next turn has Advantage and Disadvantage.

\par
\DndMonsterSection{Bonus Actions}
\DndMonsterAction{Battle Priest (5/resting)}
If Pike uses the Attack action, she can make one weapon attack as a bonus action.

\DndMonsterAction{\textit{Healing Word} (4th Level)}
Pike casts \textit{healing word}, restoring 16 (4d4 + 6) hit points to herself or another creature she can see within 60 feet of her.

\par
\DndMonsterSection{Reactions}
\DndMonsterAction{Dawnmartyr's Rebuke}
When a creature within 5 feet of Pike hits her with a melee attack, she can deal 7 (2d6) fire or radiant damage to it (her choice).

\DndMonsterAction{Everlight's Guidance (3/resting)}
When Pike or a creature she can see within 30 feet of her makes an attack roll, she can grant a +10 bonus to the roll. She can grant this bonus after the roll is made but before it hits or misses.

\end{DndMonster}



\begin{DndMonster}[float=t]{Pixie}
\DndMonsterType{Tiny fey, neutral good}
\DndMonsterBasics[
armor-class = {15},
hit-points = {\DndDice{1d4 - 1}},
speed = {10 ft., fly 30 ft.}
]
\DndMonsterAbilityScores[
 str = 2,
 dex = 20,
 con = 8,
 int = 10,
 wis = 14,
 cha = 15
]
\DndMonsterDetails[
skills = {Perception +4, Stealth +7},
languages = {Sylvan},
challenge = 1/4,
]
\DndMonsterAction{Magic Resistance}
The pixie has advantage on saving throws against spells and other magical effects.

\DndMonsterAction{Innate Spellcasting}
The pixie's innate spellcasting ability is Charisma (spell save DC 12). It can innately cast the following spells, requiring only its pixie dust as a component:
\begin{DndMonsterSpells}
  \DndInnateSpellLevel{druidcraft}
  \DndInnateSpellLevel[1e]{confusion, dancing lights, detect evil and good, detect thoughts, dispel magic, entangle, fly, phantasmal force, polymorph, sleep}
\end{DndMonsterSpells}
\par
\DndMonsterSection{Actions}
\DndMonsterAction{Superior Invisibility}
The pixie magically turns invisible until its concentration ends (as if concentration on a spell). Any equipment the pixie wears or carries is invisible with it.

\end{DndMonster}



\begin{DndMonster}[float=t]{Plainscow}
\DndMonsterType{Large beast, U}
\DndMonsterBasics[
armor-class = {13 (natural armor)},
hit-points = {\DndDice{5d10 + 15}},
speed = {40 ft.}
]
\DndMonsterAbilityScores[
 str = 18,
 dex = 8,
 con = 16,
 int = 2,
 wis = 10,
 cha = 6
]
\DndMonsterDetails[
skills = {Insight +2},
damage-resistances = {; bludgeoning, piercing, slashing from nonmagical attacks},
challenge = 2,
]
\DndMonsterAction{Trampling Charge}
If the plainscow moves at least 20 feet straight toward a creature and then hits it with a hooves attack on the same turn, that target must succeed on a DC 14 Strength saving throws or be knocked prone. If the target is prone, the plainscow can make another attack with its hooves as a bonus action.

\par
\DndMonsterSection{Actions}
\DndMonsterAction{Hooves}
\textsl{Melee Weapon Attack:} +6 to hit, reach 5 ft., one target. \textsl{Hit:} 11 (2d6 + 4) bludgeoning damage.

\end{DndMonster}



\begin{DndMonster}[float=t]{Platinum Golem}
\DndMonsterType{Large construct, U}
\DndMonsterBasics[
armor-class = {21 (natural armor)},
hit-points = {\DndDice{18d10 + 90}},
speed = {30 ft.}
]
\DndMonsterAbilityScores[
 str = 22,
 dex = 10,
 con = 21,
 int = 3,
 wis = 12,
 cha = 1
]
\DndMonsterDetails[
damage-immunities = {acid, cold, fire, poison, psychic; bludgeoning, piercing, slashing from nonmagical attacks that aren't adamantine},
senses = {darkvision 120 ft., passive perception 11},
languages = {understands the languages of its creator but can't speak},
challenge = 16,
]
\DndMonsterAction{Divine Custodian}
When the golem hits a creature with an opportunity attack, that creature's speed becomes 0 for the rest of the turn.

\DndMonsterAction{Immutable Form}
The golem is immune to any spell or effect that would alter its form.

\DndMonsterAction{Magic Resistance}
The golem has Advantage and Disadvantage on saving throws against spells and other magical effects.

\DndMonsterAction{Magic Weapons}
The golem's weapon attacks are magical.

\par
\DndMonsterSection{Actions}
\DndMonsterAction{Multiattack}
The golem makes two melee attacks.

\DndMonsterAction{Slam}
\textsl{Melee Weapon Attack:} +11 to hit, reach 5 ft., one target. \textsl{Hit:} 19 (3d8 + 6) piercing damage.

\DndMonsterAction{Mace}
\textsl{Melee Weapon Attack:} +11 to hit, reach 10 ft., one target. \textsl{Hit:} 22 (3d10 + 6) bludgeoning damage.

\DndMonsterAction{Radiant Blast (Recharge 6)}
The golem emits a burst of radiant light in a 20-foot radius. Each creature in the area must make a DC 18 Constitution saving throws. On a failed save, a creature takes 18 (4d8) radiant damage and is blinded until the end of its next turn. On a success, the creature takes half as much damage and is not blinded.

\par
\DndMonsterSection{Reactions}
\DndMonsterAction{Divine Judgment}
When a creature within 10 feet of the golem makes an attack against a target other than the golem, the golem can make a mace attack against that creature.

\end{DndMonster}



\begin{DndMonster}[float=t]{Quasit}
\DndMonsterType{Tiny fiend (demon), chaotic evil}
\DndMonsterBasics[
armor-class = {13},
hit-points = {\DndDice{3d4}},
speed = {40 ft.}
]
\DndMonsterAbilityScores[
 str = 5,
 dex = 17,
 con = 10,
 int = 7,
 wis = 10,
 cha = 10
]
\DndMonsterDetails[
skills = {Stealth +5},
damage-resistances = {cold, fire, lightning; bludgeoning, piercing, slashing from nonmagical attacks},
damage-immunities = {poison},
senses = {darkvision 120 ft., passive perception 10},
languages = {Abyssal, Common},
challenge = 1,
]
\DndMonsterAction{Shapechanger}
The quasit can use its action to polymorph into a beast form that resembles a bat (speed 10 feet fly 40 ft.), a centipede (40 ft., climb 40 ft.), or a toad (40 ft., swim 40 ft.), or back into its true form. Its statistics are the same in each form, except for the speed changes noted. Any equipment it is wearing or carrying isn't transformed. It reverts to its true form if it dies.

\DndMonsterAction{Magic Resistance}
The quasit has advantage on saving throws against spells and other magical effects.

\par
\DndMonsterSection{Actions}
\DndMonsterAction{Claw (Bite in Beast Form)}
\textsl{Melee Weapon Attack:} +4 to hit, reach 5 ft., one target. \textsl{Hit:} 5 (1d4 + 3) piercing damage, and the target must succeed on a DC 10 Constitution saving throw or take 5 (2d4) poison damage and become poisoned for 1 minute. The target can repeat the saving throw at the end of each of its turns, ending the effect on itself on a success.

\DndMonsterAction{Scare (1/Day)}
One creature of the quasit's choice within 20 feet of it must succeed on a DC 10 Wisdom saving throw or be frightened for 1 minute. The target can repeat the saving throw at the end of each of its turns, with disadvantage if the quasit is within line of sight, ending the effect on itself on a success.

\DndMonsterAction{Invisibility}
The quasit magically turns invisible until it attacks or uses Scare, or until its concentration ends (as if concentration on a spell). Any equipment the quasit wears or carries is invisible with it.

\end{DndMonster}



\begin{DndMonster}[float*=b,width=\textwidth + 8pt]{Ravager Slaughter Lord}
\begin{multicols}{2}
\DndMonsterType{Large humanoid (any), chaotic evil}
\DndMonsterBasics[
armor-class = {17 (unarmored defense)},
hit-points = {\DndDice{15d10 + 75}},
speed = {30 ft.}
]
\DndMonsterAbilityScores[
 str = 22,
 dex = 14,
 con = 20,
 int = 12,
 wis = 16,
 cha = 16
]
\DndMonsterDetails[
saving-throws = {Str +10, Con +9, Wis +7},
skills = {Intimidation +11, Religion +5},
senses = {darkvision 60 ft., passive perception 13},
languages = {Abyssal, Common, one other language},
challenge = 9,
]
\DndMonsterAction{Legendary Resistance (2/Day)}
If the Slaughter Lord fails a saving throws, it can choose to succeed instead.

\DndMonsterAction{Unarmored Defense}
While the Slaughter Lord is wearing no armor and not wielding a \textit{shield}, its AC includes its Constitution modifier.

\DndMonsterAction{Innate Spellcasting}
The Slaughter Lord's innate spellcasting ability is Wisdom (spell save DC 15, +7 to hit with spell attacks). It can innately cast the following spells, requiring no material components:
\begin{DndMonsterSpells}
  \DndInnateSpellLevel{thaumaturgy}
  \DndInnateSpellLevel[3e]{flame strike, spirit guardians}
  \DndInnateSpellLevel[1e]{control weather, divine word, fire storm}
\end{DndMonsterSpells}
\par
\DndMonsterSection{Actions}
\DndMonsterAction{Multiattack}
The Slaughter Lord makes four attacks with its dual greatswords or three attacks with its spear.

\DndMonsterAction{Dual Greatswords}
\textsl{Melee Weapon Attack:} +10 to hit, reach 5 ft., one target. \textsl{Hit:} 13 (2d6 + 6) piercing damage plus 3 (1d6) fire damage.

\DndMonsterAction{Spear}
\textsl{Melee or Ranged Weapon Attack:} +10 to hit, range 20/60 ft., one target. \textsl{Hit:} 13 (2d6 + 6) piercing damage plus 3 (1d6) fire damage, or 15 (2d8 + 6) piercing damage plus 3 (1d6) fire damage if used with both hands to make a melee attack.

\par
\DndMonsterSection{Bonus Actions}
\DndMonsterAction{Aggressive}
The Slaughter Lord can move up to its speed toward a hostile creature it can see.

\par
\DndMonsterSection{Legendary Actions}
Ravager Slaughter Lord can take 3 legendary actions, choosing from the options below. Only one legendary action can be used at a time and only at the end of another creature's turn. Ravager Slaughter Lord regains spent legendary actions at the start of its turn.

\begin{DndMonsterLegendaryActions}
\DndMonsterLegendaryAction{Attack}{The Slaughter Lord makes a spear or dual greatswords attack.
}
\DndMonsterLegendaryAction{Move}{The Slaughter Lord moves up to half its speed.
}
\DndMonsterLegendaryAction{Cast a Spell (Costs 3 Actions)}{The Slaughter Lord cast a spell.
}
\end{DndMonsterLegendaryActions}
\end{multicols}
\end{DndMonster}



\begin{DndMonster}[float=t]{Ravager Stabby-Stabber}
\DndMonsterType{Small humanoid (any), chaotic evil}
\DndMonsterBasics[
armor-class = {15 (studded leather)},
hit-points = {\DndDice{8d6 + 8}},
speed = {30 ft.}
]
\DndMonsterAbilityScores[
 str = 10,
 dex = 16,
 con = 12,
 int = 10,
 wis = 8,
 cha = 7
]
\DndMonsterDetails[
saving-throws = {Wis +1},
senses = {darkvision 60 ft., passive perception 9},
languages = {Common, one other language},
challenge = 3,
]
\DndMonsterAction{Stabby Frenzy}
When the stabby-stabber hits with a melee attack, it can attack again and continue making additional attacks until it misses. The stabby-stabber cannot move between the attacks made with this trait.

\par
\DndMonsterSection{Actions}
\DndMonsterAction{Shortsword}
\textsl{Melee Weapon Attack:} +5 to hit, reach 5 ft., one target. \textsl{Hit:} 6 (1d6 + 3) piercing damage.

\DndMonsterAction{Dagger}
\textsl{Melee or Ranged Weapon Attack:} +5 to hit, reach 5 ft., range 20/60 ft., one target. \textsl{Hit:} 5 (1d4 + 3) piercing damage.

\par
\DndMonsterSection{Bonus Actions}
\DndMonsterAction{Nimble Escape}
The stabby-stabber can use the Disengage or Hide action.

\end{DndMonster}



\begin{DndMonster}[float=t]{Remnant Chosen}
\DndMonsterType{Medium humanoid (any), U}
\DndMonsterBasics[
armor-class = {12},
hit-points = {\DndDice{18d8 + 36}},
speed = {30 ft.}
]
\DndMonsterAbilityScores[
 str = 10,
 dex = 14,
 con = 14,
 int = 15,
 wis = 16,
 cha = 20
]
\DndMonsterDetails[
saving-throws = {Cha +9, Wis +7},
skills = {Arcana +10, Deception +9},
damage-resistances = {damage from spells},
damage-immunities = {necrotic},
senses = {truesight 60 ft., passive perception 13},
languages = {Abyssal, Common, Infernal},
challenge = 12,
]
\DndMonsterAction{Magic Resistance}
The chosen has Advantage and Disadvantage on saving throws against spells and other magical effects.

\DndMonsterAction{Withering Spells}
When the chosen casts a spell of 1st level or higher, the next Withered Hand attack it makes before the end of its next turn deals extra damage equal to 1d6 per level of the spell cast.

\DndMonsterAction{Spellcasting}
The chosen is an 18th-level spellcaster. Its spellcasting ability is Charisma (spell save DC 17, +9 to hit with spell attacks). It has the following sorcerer spells prepared:
\begin{DndMonsterSpells}
  \DndMonsterSpellLevel{chill touch, dancing lights, mage hand, message, shocking grasp}
   \DndMonsterSpellLevel[1][4]{charm person, mage armor, shield}
   \DndMonsterSpellLevel[2][3]{blindness/deafness, detect thoughts}
   \DndMonsterSpellLevel[3][3]{counterspell, fly, hypnotic pattern}
   \DndMonsterSpellLevel[4][3]{greater invisibility}
   \DndMonsterSpellLevel[5][3]{dominate person, seeming}
   \DndMonsterSpellLevel[6][1]{eyebite}
   \DndMonsterSpellLevel[7][1]{finger of death}
   \DndMonsterSpellLevel[8][1]{power word stun}
\end{DndMonsterSpells}
\par
\DndMonsterSection{Actions}
\DndMonsterAction{Withering Touch}
\textsl{Melee Spell Attack:} +9 to hit, reach 5 ft., one target. \textsl{Hit:} 17 (5d6) force damage plus extra damage from the Withering Touch trait. If this damage reduces a creature to 0 hit points, the creature and everything it is wearing and carrying, except magic items, are reduced to a pile of fine gray dust. The creature can be restored to life only by means of a \textit{true resurrection} or a \textit{wish} spell.

\par
\DndMonsterSection{Bonus Actions}
\DndMonsterAction{Inescapable Sight}
The Chosen selects a creature or object under the effect of an illusion spell of 4th level or lower. One illusion of the Chosen's choice affecting the target is dispelled.

\end{DndMonster}



\begin{DndMonster}[float=t]{Remnant Cultist}
\DndMonsterType{Medium humanoid (any), U}
\DndMonsterBasics[
armor-class = {13},
hit-points = {\DndDice{11d8 + 11}},
speed = {30 ft.}
]
\DndMonsterAbilityScores[
 str = 10,
 dex = 16,
 con = 12,
 int = 18,
 wis = 8,
 cha = 8
]
\DndMonsterDetails[
saving-throws = {Wis +2},
skills = {Deception +5, Stealth +6},
damage-resistances = {necrotic, psychic},
languages = {Common, Infernal},
challenge = 7,
]
\DndMonsterAction{Unknowable Secrets}
Any attempt to form a mental link with the cultist, cast \textit{scrying} on the cultist, or cast \textit{speak with dead} on the cultist automatically fails, and the creature that initiated the attempt takes 21 (6d6) psychic damage.

\DndMonsterAction{One-Eyed}
The cultist has Advantage and Disadvantage on any attack roll made against a target more than 30 feet away.

\DndMonsterAction{Spellsteal}
If the cultist successfully uses \textit{counterspell} to negate another creature's spell, it can cast the countered spell once before the end of its next turn using one of its spell slots of the same level or higher, and requiring no material components.

\DndMonsterAction{Spellcasting}
The remnant cultist is an 11th-level spellcaster. Its spellcasting ability is Intelligence (spell save DC 15, +7 to hit with spell attacks). The cultist has the following wizard spells prepared:
\begin{DndMonsterSpells}
  \DndMonsterSpellLevel{chill touch, message, minor illusion, prestidigitation, ray of frost}
   \DndMonsterSpellLevel[1][4]{detect magic, Tasha's hideous laughter, mage armor, shield}
   \DndMonsterSpellLevel[2][3]{detect thoughts, suggestion}
   \DndMonsterSpellLevel[3][3]{counterspell, fear, vampiric touch}
   \DndMonsterSpellLevel[4][3]{greater invisibility, phantasmal killer}
   \DndMonsterSpellLevel[5][2]{dream, mislead}
   \DndMonsterSpellLevel[6][1]{circle of death}
\end{DndMonsterSpells}
\par
\DndMonsterSection{Actions}
\DndMonsterAction{Dagger of Wounding}
\textsl{Melee or Ranged Weapon Attack:} +6 to hit, reach 5 ft., range 20/60 ft., one target. \textsl{Hit:} 5 (1d4 + 3) piercing damage. Hit points lost to this weapon's damage can be regained only through a resting, rather than by regeneration, magic, or any other means.

Once per turn when the cultist hits a creature with an attack using this weapon, it can wound its target. At the start of each of the wounded creature's turns, it takes 2 (1d4) necrotic damage for each time it's been wounded by this weapon. It can then make a DC 15 Constitution saving throws, ending the effect of all such wounds on itself on a success. Alternatively, the wounded creature, or a creature within 5 feet of it, can use an action to make a DC 15 Wisdom (Medicine) check, ending the effect of such wounds on it on a success.

\end{DndMonster}



\begin{DndMonster}[float=t]{Remorhaz}
\DndMonsterType{Huge monstrosity, unaligned}
\DndMonsterBasics[
armor-class = {17 (natural armor)},
hit-points = {\DndDice{17d12 + 85}},
speed = {30 ft., burrow 20 ft.}
]
\DndMonsterAbilityScores[
 str = 24,
 dex = 13,
 con = 21,
 int = 4,
 wis = 10,
 cha = 5
]
\DndMonsterDetails[
damage-immunities = {cold, fire},
senses = {darkvision 60 ft., tremorsense 60 ft., passive perception 10},
challenge = 11,
]
\DndMonsterAction{Heated Body}
A creature that touches the remorhaz or hits it with a melee attack while within 5 feet of it takes 10 (3d6) fire damage.

\par
\DndMonsterSection{Actions}
\DndMonsterAction{Bite}
\textsl{Melee Weapon Attack:} +11 to hit, reach 10 ft., one target. \textsl{Hit:} 40 (6d10 + 7) piercing damage plus 10 (3d6) fire damage. If the target is a creature, it is grappled (escape DC 17). Until this grapple ends, the target is restrained, and the remorhaz can't bite another target.

\DndMonsterAction{Swallow}
The remorhaz makes one bite attack against a Medium or smaller creature it is grappling. If the attack hits, that creature takes the bite's damage and is swallowed, and the grapple ends. While swallowed, the creature is blinded and restrained, it has total cover against attacks and other effects outside the remorhaz, and it takes 21 (6d6) acid damage at the start of each of the remorhaz's turns.

If the remorhaz takes 30 damage or more on a single turn from a creature inside it, the remorhaz must succeed on a DC 15 Constitution saving throw at the end of that turn or regurgitate all swallowed creatures, which fall prone in a space within 10 feet of the remorhaz. If the remorhaz dies, a swallowed creature is no longer restrained by it and can escape from the corpse using 15 feet of movement, exiting prone.

\end{DndMonster}



\begin{DndMonster}[float=t]{Rivermaw Brawler}
\DndMonsterType{Medium humanoid (any), U}
\DndMonsterBasics[
armor-class = {15 (unarmored defense)},
hit-points = {\DndDice{10d8 + 20}},
speed = {40 ft.}
]
\DndMonsterAbilityScores[
 str = 17,
 dex = 14,
 con = 15,
 int = 14,
 wis = 16,
 cha = 8
]
\DndMonsterDetails[
saving-throws = {Str +5, Wis +5},
skills = {Acrobatics +4, Athletics +5, Perception +5},
languages = {Common, Giant},
challenge = 4,
]
\DndMonsterAction{Unarmored Defense}
While the brawler is wearing no armor and wielding no \textit{shield}, its AC includes its Wisdom modifier.

\par
\DndMonsterSection{Actions}
\DndMonsterAction{Multiattack}
The brawler makes three unarmed strikes.

\DndMonsterAction{Unarmed Strike}
\textsl{Melee Weapon Attack:} +5 to hit, reach 5 ft., one target. \textsl{Hit:} 7 (1d8 + 3) bludgeoning damage. If the target is a Large or smaller creature, the brawler can choose one of the following additional effects:


\begin{itemize}
\item{Bruise.}
The target must succeed on a DC 13 Constitution saving throws. On a failure, its speed is reduced by 5 feet until it receives at least 1 point of magical healing. A creature can use its action to make a DC 15 Wisdom (Medicine) check to heal an affected creature within 5 feet of it (including itself), restoring any lost speed on a success.

\item{Disarm.}
The target must succeed on a DC 13 Strength saving throws or drop one item it is holding. The brawler can choose which item is dropped, and can pick it up.

\item{Suplex.}
The target must succeed on a DC 13 Strength saving throws or be knocked prone.

\end{itemize}

\par
\DndMonsterSection{Reactions}
\DndMonsterAction{Clever Fighter}
When the brawler hits with an attack or is targeted by an attack, it can add its Intelligence modifier (+2) to its damage rolls or AC until the end of its turn.

\DndMonsterAction{Endurance (Recharges after a resting)}
When the brawler takes damage, it can reduce the damage by 8 (1d12 + 2).

\end{DndMonster}



\begin{DndMonster}[float=t]{Rivermaw Stormborn}
\DndMonsterType{Medium humanoid (any), U}
\DndMonsterBasics[
armor-class = {13 (\textit{hide armor})},
hit-points = {\DndDice{12d8 + 48}},
speed = {30 ft.}
]
\DndMonsterAbilityScores[
 str = 18,
 dex = 12,
 con = 18,
 int = 11,
 wis = 10,
 cha = 12
]
\DndMonsterDetails[
saving-throws = {Con +7, Cha +4},
skills = {Athletics +7, Intimidation +4},
damage-resistances = {lightning},
languages = {Common, Giant},
challenge = 5,
]
\DndMonsterAction{Reckless}
At the start of its turn, the stormborn can gain Advantage and Disadvantage on all melee weapon attack rolls during that turn, but attack rolls against it have Advantage and Disadvantage until the start of its next turn.

\DndMonsterAction{Innate Spellcasting}
The stormborn's innate spellcasting ability is Charisma (spell save DC 12). The stormborn can innately cast the following spells, requiring no material components:
\begin{DndMonsterSpells}
  \DndInnateSpellLevel[1]{call lightning}
  \DndInnateSpellLevel[3]{lightning bolt}
\end{DndMonsterSpells}
\par
\DndMonsterSection{Actions}
\DndMonsterAction{Multiattack}
The stormborn makes two greataxe attacks. If it is concentrating on \textit{call lightning}, it can also call down a lightning bolt.

\DndMonsterAction{Greataxe}
\textsl{Melee Weapon Attack:} +7 to hit, reach 5 ft., one target. \textsl{Hit:} 10 (1d12 + 4) slashing damage plus 3 (1d6) lightning damage.

\par
\DndMonsterSection{Reactions}
\DndMonsterAction{Endurance (Recharges after a resting)}
When the stormborn takes damage, it can reduce the damage by 10 (1d12 + 4).

\end{DndMonster}


\clearpage

\begin{DndMonster}[float=t]{Guardian Naga}
\DndMonsterType{Large monstrosity, lawful good}
\DndMonsterBasics[
armor-class = {18 (natural armor)},
hit-points = {\DndDice{15d10 + 45}},
speed = {40 ft.}
]
\DndMonsterAbilityScores[
 str = 19,
 dex = 18,
 con = 16,
 int = 16,
 wis = 19,
 cha = 18
]
\DndMonsterDetails[
saving-throws = {Dex +8, Con +7, Int +7, Wis +8, Cha +8},
damage-immunities = {poison},
senses = {darkvision 60 ft., passive perception 14},
languages = {Celestial, Common},
challenge = 10,
]
\DndMonsterAction{Rejuvenation}
If it dies, the naga returns to life in 1d6 days and regains all its hit points. Only a \textit{wish} spell can prevent this trait from functioning.

\DndMonsterAction{Spellcasting}
The naga is an 11th-level spellcaster. Its spellcasting ability is Wisdom (spell save DC 16, +8 to hit with spell attacks), and it needs only verbal components to cast its spells. It has the following cleric spells prepared:
\begin{DndMonsterSpells}
  \DndMonsterSpellLevel{mending, sacred flame, thaumaturgy}
   \DndMonsterSpellLevel[1][4]{command, cure wounds, shield of faith}
   \DndMonsterSpellLevel[2][3]{calm emotions, hold person}
   \DndMonsterSpellLevel[3][3]{bestow curse, clairvoyance}
   \DndMonsterSpellLevel[4][3]{banishment, freedom of movement}
   \DndMonsterSpellLevel[5][2]{flame strike, geas}
   \DndMonsterSpellLevel[6][1]{true seeing}
\end{DndMonsterSpells}
\par
\DndMonsterSection{Actions}
\DndMonsterAction{Bite}
\textsl{Melee Weapon Attack:} +8 to hit, reach 10 ft., one creature. \textsl{Hit:} 8 (1d8 + 4) piercing damage, and the target must make a DC 15 Constitution saving throw, taking 45 (10d8) poison damage on a failed save, or half as much damage on a successful one.

\DndMonsterAction{Spit Poison}
\textsl{Ranged Weapon Attack:} +8 to hit, range 15/30 ft., one creature. \textsl{Hit:} The target must make a DC 15 Constitution saving throw, taking 45 (10d8) poison damage on a failed save, or half as much damage on a successful one.

\end{DndMonster}



\begin{DndMonster}[float=t]{Sahuagin}
\DndMonsterType{Medium humanoid (sahuagin), lawful evil}
\DndMonsterBasics[
armor-class = {12 (natural armor)},
hit-points = {\DndDice{4d8 + 4}},
speed = {30 ft., swim 40 ft.}
]
\DndMonsterAbilityScores[
 str = 13,
 dex = 11,
 con = 12,
 int = 12,
 wis = 13,
 cha = 9
]
\DndMonsterDetails[
skills = {Perception +5},
senses = {darkvision 120 ft., passive perception 15},
languages = {Sahuagin},
challenge = 1/2,
]
\DndMonsterAction{Blood Frenzy}
The sahuagin has advantage on melee attack rolls against any creature that doesn't have all its hit points.

\DndMonsterAction{Limited Amphibiousness}
The sahuagin can breathe air and water, but it needs to be submerged at least once every 4 hours to avoid suffocating.

\DndMonsterAction{Shark Telepathy}
The sahuagin can magically command any shark within 120 feet of it, using a limited telepathy.

\par
\DndMonsterSection{Actions}
\DndMonsterAction{Multiattack}
The sahuagin makes two melee attacks: one with its bite and one with its claws or spear.

\DndMonsterAction{Bite}
\textsl{Melee Weapon Attack:} +3 to hit, reach 5 ft., one target. \textsl{Hit:} 3 (1d4 + 1) piercing damage.

\DndMonsterAction{Claws}
\textsl{Melee Weapon Attack:} +3 to hit, reach 5 ft., one target. \textsl{Hit:} 3 (1d4 + 1) slashing damage.

\DndMonsterAction{Spear}
\textsl{Melee or Ranged Weapon Attack:} +3 to hit, reach 5 ft. or range 20/60 ft., one target. \textsl{Hit:} 4 (1d6 + 1) piercing damage, or 5 (1d8 + 1) piercing damage if used with two hands to make a melee attack.

\end{DndMonster}



\begin{DndMonster}[float=t]{Sahuagin Baron}
\DndMonsterType{Large humanoid (sahuagin), lawful evil}
\DndMonsterBasics[
armor-class = {16 (\textit{breastplate})},
hit-points = {\DndDice{9d10 + 27}},
speed = {30 ft., swim 50 ft.}
]
\DndMonsterAbilityScores[
 str = 19,
 dex = 15,
 con = 16,
 int = 14,
 wis = 13,
 cha = 17
]
\DndMonsterDetails[
saving-throws = {Dex +5, Con +6, Int +5, Wis +4},
skills = {Perception +7},
senses = {darkvision 120 ft., passive perception 17},
languages = {Sahuagin},
challenge = 5,
]
\DndMonsterAction{Blood Frenzy}
The sahuagin has advantage on melee attack rolls against any creature that doesn't have all its hit points.

\DndMonsterAction{Limited Amphibiousness}
The sahuagin can breathe air and water, but it needs to be submerged at least once every 4 hours to avoid suffocating.

\DndMonsterAction{Shark Telepathy}
The sahuagin can magically command any shark within 120 feet of it, using a limited telepathy.

\par
\DndMonsterSection{Actions}
\DndMonsterAction{Multiattack}
The sahuagin makes three attacks: one with his bite and two with his claws or trident.

\DndMonsterAction{Bite}
\textsl{Melee Weapon Attack:} +7 to hit, reach 5 ft., one target. \textsl{Hit:} 9 (2d4 + 4) piercing damage.

\DndMonsterAction{Claws}
\textsl{Melee Weapon Attack:} +7 to hit, reach 5 ft., one target. \textsl{Hit:} 11 (2d6 + 4) slashing damage.

\DndMonsterAction{Trident}
\textsl{Melee or Ranged Weapon Attack:} +7 to hit, reach 5 ft. or range 20/60 ft., one target. \textsl{Hit:} 11 (2d6 + 4) piercing damage, or 13 (2d8 + 4) piercing damage if used with two hands to make a melee attack.

\end{DndMonster}



\begin{DndMonster}[float=t]{Salamander}
\DndMonsterType{Large elemental, neutral evil}
\DndMonsterBasics[
armor-class = {15 (natural armor)},
hit-points = {\DndDice{12d10 + 24}},
speed = {30 ft.}
]
\DndMonsterAbilityScores[
 str = 18,
 dex = 14,
 con = 15,
 int = 11,
 wis = 10,
 cha = 12
]
\DndMonsterDetails[
damage-vulnerabilities = {cold},
damage-resistances = {; bludgeoning, piercing, slashing from nonmagical attacks},
damage-immunities = {fire},
senses = {darkvision 60 ft., passive perception 10},
languages = {Ignan},
challenge = 5,
]
\DndMonsterAction{Heated Body}
A creature that touches the salamander or hits it with a melee attack while within 5 feet of it takes 7 (2d6) fire damage.

\DndMonsterAction{Heated Weapons}
Any metal melee weapon the salamander wields deals an extra 3 (1d6) fire damage on a hit (included in the attack).

\par
\DndMonsterSection{Actions}
\DndMonsterAction{Multiattack}
The salamander makes two attacks: one with its spear and one with its tail.

\DndMonsterAction{Spear}
\textsl{Melee or Ranged Weapon Attack:} +7 to hit, reach 5 ft. or range 20/60 ft., one target. \textsl{Hit:} 11 (2d6 + 4) piercing damage, or 13 (2d8 + 4) piercing damage if used with two hands to make a melee attack, plus 3 (1d6) fire damage.

\DndMonsterAction{Tail}
\textsl{Melee Weapon Attack:} +7 to hit, reach 10 ft., one target. \textsl{Hit:} 11 (2d6 + 4) bludgeoning damage plus 7 (2d6) fire damage, and the target is grappled (escape DC 14). Until this grapple ends, the target is restrained, the salamander can automatically hit the target with its tail, and the salamander can't make tail attacks against other targets.

\end{DndMonster}



\begin{DndMonster}[float=t]{Satyr}
\DndMonsterType{Medium fey, chaotic neutral}
\DndMonsterBasics[
armor-class = {14 (\textit{leather armor})},
hit-points = {\DndDice{7d8}},
speed = {40 ft.}
]
\DndMonsterAbilityScores[
 str = 12,
 dex = 16,
 con = 11,
 int = 12,
 wis = 10,
 cha = 14
]
\DndMonsterDetails[
skills = {Perception +2, Performance +6, Stealth +5},
languages = {Common, Elvish, Sylvan},
challenge = 1/2,
]
\DndMonsterAction{Magic Resistance}
The satyr has advantage on saving throws against spells and other magical effects.

\par
\DndMonsterSection{Actions}
\DndMonsterAction{Ram}
\textsl{Melee Weapon Attack:} +3 to hit, reach 5 ft., one target. \textsl{Hit:} 6 (2d4 + 1) bludgeoning damage.

\DndMonsterAction{Shortsword}
\textsl{Melee Weapon Attack:} +5 to hit, reach 5 ft., one target. \textsl{Hit:} 6 (1d6 + 3) piercing damage.

\DndMonsterAction{Shortbow}
\textsl{Ranged Weapon Attack:} +5 to hit, range 80/320 ft., one target. \textsl{Hit:} 6 (1d6 + 3) piercing damage.

\end{DndMonster}



\begin{DndMonster}[float=t]{Scanlan Shorthalt}
\DndMonsterType{Small humanoid (gnome), U}
\DndMonsterBasics[
armor-class = {16 (\textit{+2 Studded Leather Armor of Acid Resistance})},
hit-points = {\DndDice{20d8 + 60 + 40}},
speed = {25 ft.}
]
\DndMonsterAbilityScores[
 str = 13,
 dex = 11,
 con = 16,
 int = 16,
 wis = 7,
 cha = 22
]
\DndMonsterDetails[
saving-throws = {Str +3, Dex +7, Con +5, Int +5, Wis +0, Cha +13},
skills = {Acrobatics +5, Arcana +8, Athletics +6, Deception +16, Intimidation +11, Investigation +13, Performance +16, Persuasion +16},
damage-resistances = {acid},
senses = {darkvision 60 ft., passive perception 8},
languages = {Common, Gnomish, Marquesian},
challenge = 15,
]
\DndMonsterAction{Gnome Cunning}
Scanlan has Advantage and Disadvantage on Intelligence, Wisdom, and Charisma saving throws against magic.

\DndMonsterAction{Focused}
Scanlan has Advantage and Disadvantage on Constitution saving throws made to maintain concentration on spells.

\DndMonsterAction{Special Equipment}
Scanlan wields Mythcarver, wears \textit{+2 Studded Leather Armor of Acid Resistance} and a ring of protection with a +2 bonus, and carries a \textit{hat of disguise} and a \textit{wand of magic missiles}.

\DndMonsterAction{Toughness}
Scanlan has an additional 40 hit points.

\DndMonsterAction{Spellcasting}
Scanlan is a 20th-level spellcaster. His spellcasting ability is Charisma (spell save DC 19, +11 to hit with spell attacks). He has the following bard spells prepared:
\begin{DndMonsterSpells}
  \DndMonsterSpellLevel{mage hand, message, minor illusion, vicious mockery}
   \DndMonsterSpellLevel[1][4]{healing word, thunderwave}
   \DndMonsterSpellLevel[2][3]{blindness/deafness, hold person, invisibility, suggestion}
   \DndMonsterSpellLevel[3][3]{counterspell, lightning bolt, stinking cloud}
   \DndMonsterSpellLevel[4][3]{dimension door, polymorph, Otiluke's Resilient Sphere}
   \DndMonsterSpellLevel[5][3]{bigby's hand, dominate person, modify memory, seeming}
   \DndMonsterSpellLevel[6][2]{eyebite}
   \DndMonsterSpellLevel[7][2]{Mordenkainen's Magnificent Mansion, reverse gravity}
   \DndMonsterSpellLevel[8][1]{dominate monster}
\end{DndMonsterSpells}
\par
\DndMonsterSection{Actions}
\DndMonsterAction{Mythcarver}
\textsl{Melee Weapon Attack:} +9 to hit, reach 5 ft., one target. \textsl{Hit:} 7 (1d8 + 3) piercing damage and 3 (1d6) force damage.

\par
\DndMonsterSection{Bonus Actions}
\DndMonsterAction{Arcane Fist (8th Level)}
\textsl{Melee Spell Attack:} +11 to hit, reach 5 ft., one target. \textsl{Hit:} 45 (10d8) force damage. Scanlan can use this bonus action only while concentrating on the \textit{Bigby's Hand} spell.

\DndMonsterAction{Bardic Inspiration (6/resting)}
Scanlan grants a Bardic Inspiration die to one creature other than himself within 60 feet of him who can hear him. Once within the next 10 minutes, the creature can roll a d12 and add the number rolled to one ability check, attack roll, or saving throws it makes.  When Scanlan uses this ability, he gains Advantage and Disadvantage on attack rolls with Mythcarver until the end of his turn.

\DndMonsterAction{\textit{Healing Word} (5th Level)}
Scanlan casts \textit{healing word}, restoring 18 (5d4 + 6) hit points to himself or another creature he can see within 60 feet of him.

\par
\DndMonsterSection{Reactions}
\DndMonsterAction{Cutting Words}
When a creature Scanlan can see within 60 feet of him makes an attack roll, an ability check, or a damage roll, he can expend a use of Bardic Inspiration to roll a d12 and subtract the number rolled from the creature's roll. The creature also has Advantage and Disadvantage on its next saving throws.

\end{DndMonster}



\begin{DndMonster}[float=t]{Scout}
\DndMonsterType{Medium humanoid (any race), any alignment}
\DndMonsterBasics[
armor-class = {13 (\textit{leather armor})},
hit-points = {\DndDice{3d8 + 3}},
speed = {30 ft.}
]
\DndMonsterAbilityScores[
 str = 11,
 dex = 14,
 con = 12,
 int = 11,
 wis = 13,
 cha = 11
]
\DndMonsterDetails[
skills = {Nature +4, Perception +5, Stealth +6, Survival +5},
languages = {any one language (usually Common)},
challenge = 1/2,
]
\DndMonsterAction{Keen Hearing and Sight}
The scout has advantage on Wisdom (Perception) checks that rely on hearing or sight.

\par
\DndMonsterSection{Actions}
\DndMonsterAction{Multiattack}
The scout makes two melee attacks or two ranged attacks.

\DndMonsterAction{Shortsword}
\textsl{Melee Weapon Attack:} +4 to hit, reach 5 ft., one target. \textsl{Hit:} 5 (1d6 + 2) piercing damage.

\DndMonsterAction{Longbow}
\textsl{Ranged Weapon Attack:} +4 to hit, ranged 150/600 ft., one target. \textsl{Hit:} 6 (1d8 + 2) piercing damage.

\end{DndMonster}



\begin{DndMonster}[float*=b,width=\textwidth + 8pt]{Sea Hag}
\begin{multicols}{2}
\DndMonsterType{Medium fey, chaotic evil}
\DndMonsterBasics[
armor-class = {14 (natural armor)},
hit-points = {\DndDice{7d8 + 21}},
speed = {30 ft., swim 40 ft.}
]
\DndMonsterAbilityScores[
 str = 16,
 dex = 13,
 con = 16,
 int = 12,
 wis = 12,
 cha = 13
]
\DndMonsterDetails[
senses = {darkvision 60 ft., passive perception 11},
languages = {Aquan, Common, Giant},
challenge = {'cr': '2', 'coven': '4'},
]
\DndMonsterAction{Amphibious}
The hag can breathe air and water.

\DndMonsterAction{Horrific Appearance}
Any humanoid that starts its turn within 30 feet of the hag and can see the hag's true form must make a DC 11 Wisdom saving throw. On a failed save, the creature is frightened for 1 minute. A creature can repeat the saving throw at the end of each of its turns, with disadvantage if the hag is within line of sight, ending the effect on itself on a success. If a creature's saving throw is successful or the effect ends for it, the creature is immune to the hag's Horrific Appearance for the next 24 hours.

Unless the target is Surprise or the revelation of the hag's true form is sudden, the target can avert its eyes and avoid making the initial saving throw. Until the start of its next turn, a creature that averts its eyes has disadvantage on attack rolls against the hag.

\par
\DndMonsterSection{Actions}
\DndMonsterAction{Claws}
\textsl{Melee Weapon Attack:} +5 to hit, reach 5 ft., one target. \textsl{Hit:} 10 (2d6 + 3) slashing damage.

\DndMonsterAction{Death Glare}
The hag targets one frightened creature she can see within 30 feet of her. If the target can see the hag, it must succeed on a DC 11 Wisdom saving throw against this magic or drop to 0 hit points.

\DndMonsterAction{Illusory Appearance}
The hag covers herself and anything she is wearing or carrying with a magical illusion that makes her look like an ugly creature of her general size and humanoid shape. The effect ends if the hag takes a bonus action to end it or if she dies.

The changes wrought by this effect fail to hold up to physical inspection. For example, the hag could appear to have no claws, but someone touching her hand might feel the claws. Otherwise, a creature must take an action to visually inspect the illusion and succeed on a DC 16 Intelligence (Investigation) check to discern that the hag is disguised.

\end{multicols}
\end{DndMonster}



\begin{DndMonster}[float=t]{Shambling Mound}
\DndMonsterType{Large plant, unaligned}
\DndMonsterBasics[
armor-class = {15 (natural armor)},
hit-points = {\DndDice{16d10 + 48}},
speed = {20 ft., swim 20 ft.}
]
\DndMonsterAbilityScores[
 str = 18,
 dex = 8,
 con = 16,
 int = 5,
 wis = 10,
 cha = 5
]
\DndMonsterDetails[
skills = {Stealth +2},
damage-resistances = {cold, fire},
damage-immunities = {lightning},
senses = {blindsight 60 ft. (blind beyond this radius), passive perception 10},
challenge = 5,
]
\DndMonsterAction{Lightning Absorption}
Whenever the shambling mound is subjected to lightning damage, it takes no damage and regains a number of hit points equal to the lightning damage dealt.

\par
\DndMonsterSection{Actions}
\DndMonsterAction{Multiattack}
The shambling mound makes two slam attacks. If both attacks hit a Medium or smaller target, the target is grappled (escape DC 14), and the shambling mound uses its Engulf on it.

\DndMonsterAction{Slam}
\textsl{Melee Weapon Attack:} +7 to hit, reach 5 ft., one target. \textsl{Hit:} 13 (2d8 + 4) bludgeoning damage.

\DndMonsterAction{Engulf}
The shambling mound engulfs a Medium or smaller creature grappled by it. The engulfed target is blinded, restrained, and unable to breathe, and it must succeed on a DC 14 Constitution saving throw at the start of each of the mound's turns or take 13 (2d8 + 4) bludgeoning damage. If the mound moves, the engulfed target moves with it. The mound can have only one creature engulfed at a time.

\end{DndMonster}



\begin{DndMonster}[float=t]{Siren}
\DndMonsterType{Medium fey, chaotic good}
\DndMonsterBasics[
armor-class = {14},
hit-points = {\DndDice{7d8 + 7}},
speed = {30 ft., swim 30 ft.}
]
\DndMonsterAbilityScores[
 str = 10,
 dex = 18,
 con = 12,
 int = 13,
 wis = 14,
 cha = 16
]
\DndMonsterDetails[
skills = {Medicine +4, Nature +3, Stealth +6, Survival +4},
senses = {darkvision 60 ft., passive perception 12},
languages = {Common, Elvish, Sylvan},
challenge = 3,
]
\DndMonsterAction{Amphibious}
Siren can breathe air and water.

\DndMonsterAction{Magic Resistance}
Siren has advantage on saving throws against spells and other magical effects.

\DndMonsterAction{Innate Spellcasting}
Siren's innate spellcasting ability is Charisma (spell save DC 13). She can innately cast the following spells, requiring no material components:
\begin{DndMonsterSpells}
  \DndInnateSpellLevel[1e]{charm person, fog cloud, greater invisibility, polymorph}
\end{DndMonsterSpells}
\par
\DndMonsterSection{Actions}
\DndMonsterAction{Shortsword}
\textsl{Melee Weapon Attack:} +6 to hit, reach 5 ft., one target. \textsl{Hit:} 7 (1d6 + 4) piercing damage.

\DndMonsterAction{Stupefying Touch}
Siren touches one creature she can see within 5 feet of her. The creature must succeed on a DC 13 Intelligence saving throw or take 13 (3d6 + 3) psychic damage and be stunned until the start of Siren's next turn.

\end{DndMonster}



\begin{DndMonster}[float=t]{Skeleton}
\DndMonsterType{Medium undead, lawful evil}
\DndMonsterBasics[
armor-class = {13 (armor scraps)},
hit-points = {\DndDice{2d8 + 4}},
speed = {30 ft.}
]
\DndMonsterAbilityScores[
 str = 10,
 dex = 14,
 con = 15,
 int = 6,
 wis = 8,
 cha = 5
]
\DndMonsterDetails[
damage-vulnerabilities = {bludgeoning},
damage-immunities = {poison},
senses = {darkvision 60 ft., passive perception 9},
languages = {understands all languages it spoke in life but can't speak},
challenge = 1/4,
]
\par
\DndMonsterSection{Actions}
\DndMonsterAction{Shortsword}
\textsl{Melee Weapon Attack:} +4 to hit, reach 5 ft., one target. \textsl{Hit:} 5 (1d6 + 2) piercing damage.

\DndMonsterAction{Shortbow}
\textsl{Ranged Weapon Attack:} +4 to hit, range 80/320 ft., one target. \textsl{Hit:} 5 (1d6 + 2) piercing damage.

\end{DndMonster}



\begin{DndMonster}[float=t]{Specter}
\DndMonsterType{Medium undead, chaotic evil}
\DndMonsterBasics[
armor-class = {12},
hit-points = {\DndDice{5d8}},
speed = {0 ft., fly 50 ft. \textit{(hover)}}
]
\DndMonsterAbilityScores[
 str = 1,
 dex = 14,
 con = 11,
 int = 10,
 wis = 10,
 cha = 11
]
\DndMonsterDetails[
damage-resistances = {acid, cold, fire, lightning, thunder; bludgeoning, piercing, slashing from nonmagical attacks},
damage-immunities = {necrotic, poison},
senses = {darkvision 60 ft., passive perception 10},
languages = {understands all languages it knew in life but can't speak},
challenge = 1,
]
\DndMonsterAction{Incorporeal Movement}
The specter can move through other creatures and objects as if they were difficult terrain. It takes 5 (1d10) force damage if it ends its turn inside an object.

\DndMonsterAction{Sunlight Sensitivity}
While in sunlight, the specter has disadvantage on attack rolls, as well as on Wisdom (Perception) checks that rely on sight.

\par
\DndMonsterSection{Actions}
\DndMonsterAction{Life Drain}
\textsl{Melee Spell Attack:} +4 to hit, reach 5 ft., one creature. \textsl{Hit:} 10 (3d6) necrotic damage. The target must succeed on a DC 10 Constitution saving throw or its hit point maximum is reduced by an amount equal to the damage taken. This reduction lasts until the creature finishes a long rest. The target dies if this effect reduces its hit point maximum to 0.

\end{DndMonster}



\begin{DndMonster}[float=t]{Spirit Naga}
\DndMonsterType{Large monstrosity, chaotic evil}
\DndMonsterBasics[
armor-class = {15 (natural armor)},
hit-points = {\DndDice{10d10 + 20}},
speed = {40 ft.}
]
\DndMonsterAbilityScores[
 str = 18,
 dex = 17,
 con = 14,
 int = 16,
 wis = 15,
 cha = 16
]
\DndMonsterDetails[
saving-throws = {Dex +6, Con +5, Wis +5, Cha +6},
damage-immunities = {poison},
senses = {darkvision 60 ft., passive perception 12},
languages = {Abyssal, Common},
challenge = 8,
]
\DndMonsterAction{Rejuvenation}
If it dies, the naga returns to life in 1d6 days and regains all its hit points. Only a \textit{wish} spell can prevent this trait from functioning.

\DndMonsterAction{Spellcasting}
The naga is a 10th-level spellcaster. Its spellcasting ability is Intelligence (spell save DC 14, +6 to hit with spell attacks), and it needs only verbal components to cast its spells. It has the following wizard spells prepared:
\begin{DndMonsterSpells}
  \DndMonsterSpellLevel{mage hand, minor illusion, ray of frost}
   \DndMonsterSpellLevel[1][4]{charm person, detect magic, sleep}
   \DndMonsterSpellLevel[2][3]{detect thoughts, hold person}
   \DndMonsterSpellLevel[3][3]{lightning bolt, water breathing}
   \DndMonsterSpellLevel[4][3]{blight, dimension door}
   \DndMonsterSpellLevel[5][2]{dominate person}
\end{DndMonsterSpells}
\par
\DndMonsterSection{Actions}
\DndMonsterAction{Bite}
\textsl{Melee Weapon Attack:} +7 to hit, reach 10 ft., one creature. \textsl{Hit:} 7 (1d6 + 4) piercing damage, and the target must make a DC 13 Constitution saving throw, taking 31 (7d8) poison damage on a failed save, or half as much damage on a successful one.

\end{DndMonster}



\begin{DndMonster}[float=t]{Sprite}
\DndMonsterType{Tiny fey, neutral good}
\DndMonsterBasics[
armor-class = {15 (\textit{leather armor})},
hit-points = {\DndDice{1d4}},
speed = {10 ft., fly 40 ft.}
]
\DndMonsterAbilityScores[
 str = 3,
 dex = 18,
 con = 10,
 int = 14,
 wis = 13,
 cha = 11
]
\DndMonsterDetails[
skills = {Perception +3, Stealth +8},
languages = {Common, Elvish, Sylvan},
challenge = 1/4,
]
\par
\DndMonsterSection{Actions}
\DndMonsterAction{Longsword}
\textsl{Melee Weapon Attack:} +2 to hit, reach 5 ft., one target. \textsl{Hit:} 1 slashing damage.

\DndMonsterAction{Shortbow}
\textsl{Ranged Weapon Attack:} +6 to hit, range 40/160 ft., one target. \textsl{Hit:} 1 piercing damage, and the target must succeed on a DC 10 Constitution saving throw or become poisoned for 1 minute. If its saving throw result is 5 or lower, the poisoned target falls unconscious for the same duration, or until it takes damage or another creature takes an action to shake it awake.

\DndMonsterAction{Heart Sight}
The sprite touches a creature and magically knows the creature's current emotional state. If the target fails a DC 10 Charisma saving throw, the sprite also knows the creature's alignment. Celestials, fiends, and undead automatically fail the saving throw.

\DndMonsterAction{Invisibility}
The sprite magically turns invisible until it attacks or casts a spell, or until its concentration ends (as if concentration on a spell). Any equipment the sprite wears or carries is invisible with it.

\end{DndMonster}



\begin{DndMonster}[float=t]{Stone Giant}
\DndMonsterType{Huge giant, neutral}
\DndMonsterBasics[
armor-class = {17 (natural armor)},
hit-points = {\DndDice{11d12 + 55}},
speed = {40 ft.}
]
\DndMonsterAbilityScores[
 str = 23,
 dex = 15,
 con = 20,
 int = 10,
 wis = 12,
 cha = 9
]
\DndMonsterDetails[
saving-throws = {Dex +5, Con +8, Wis +4},
skills = {Athletics +12, Perception +4},
senses = {darkvision 60 ft., passive perception 14},
languages = {Giant},
challenge = 7,
]
\DndMonsterAction{Stone Camouflage}
The giant has advantage on Dexterity (Stealth) checks made to hide in rocky terrain.

\par
\DndMonsterSection{Actions}
\DndMonsterAction{Multiattack}
The giant makes two greatclub attacks.

\DndMonsterAction{Greatclub}
\textsl{Melee Weapon Attack:} +9 to hit, reach 15 ft., one target. \textsl{Hit:} 19 (3d8 + 6) bludgeoning damage.

\DndMonsterAction{Rock}
\textsl{Ranged Weapon Attack:} +9 to hit, range 60/240 ft., one target. \textsl{Hit:} 28 (4d10 + 6) bludgeoning damage. If the target is a creature, it must succeed on a DC 17 Strength saving throw or be knocked prone.

\par
\DndMonsterSection{Reactions}
\DndMonsterAction{Rock Catching}
If a rock or similar object is hurled at the giant, the giant can, with a successful DC 10 Dexterity saving throw, catch the missile and take no bludgeoning damage from it.

\end{DndMonster}



\begin{DndMonster}[float=t]{Storm Giant}
\DndMonsterType{Huge giant, chaotic good}
\DndMonsterBasics[
armor-class = {16 (\textit{scale mail})},
hit-points = {\DndDice{20d12 + 100}},
speed = {50 ft., swim 50 ft.}
]
\DndMonsterAbilityScores[
 str = 29,
 dex = 14,
 con = 20,
 int = 16,
 wis = 18,
 cha = 18
]
\DndMonsterDetails[
saving-throws = {Str +14, Con +10, Wis +9, Cha +9},
skills = {Arcana +8, Athletics +14, History +8, Perception +9},
damage-resistances = {cold},
damage-immunities = {lightning, thunder},
languages = {Common, Giant},
challenge = 13,
]
\DndMonsterAction{Amphibious}
The giant can breathe air and water.

\DndMonsterAction{Innate Spellcasting}
The giant's innate spellcasting ability is Charisma (spell save DC 17). It can innately cast the following spells, requiring no material components:
\begin{DndMonsterSpells}
  \DndInnateSpellLevel{detect magic, feather fall, levitate, light}
  \DndInnateSpellLevel[3e]{control weather, water breathing}
\end{DndMonsterSpells}
\par
\DndMonsterSection{Actions}
\DndMonsterAction{Multiattack}
The giant makes two greatsword attacks.

\DndMonsterAction{Greatsword}
\textsl{Melee Weapon Attack:} +14 to hit, reach 10 ft., one target. \textsl{Hit:} 30 (6d6 + 9) slashing damage.

\DndMonsterAction{Rock}
\textsl{Ranged Weapon Attack:} +14 to hit, range 60/240 ft., one target. \textsl{Hit:} 35 (4d12 + 9) bludgeoning damage.

\DndMonsterAction{Lightning Strike (Recharge 5\textendash 6)}
The giant hurls a magical lightning bolt at a point it can see within 500 feet of it. Each creature within 10 feet of that point must make a DC 17 Dexterity saving throw, taking 54 (12d8) lightning damage on a failed save, or half as much damage on a successful one.

\end{DndMonster}



\begin{DndMonster}[float=t]{Swashbuckler}
\DndMonsterType{Medium humanoid (any race), any non-lawful alignment}
\DndMonsterBasics[
armor-class = {17 (\textit{leather armor})},
hit-points = {\DndDice{12d8 + 12}},
speed = {30 ft.}
]
\DndMonsterAbilityScores[
 str = 12,
 dex = 18,
 con = 12,
 int = 14,
 wis = 11,
 cha = 15
]
\DndMonsterDetails[
skills = {Acrobatics +8, Athletics +5, Persuasion +6},
languages = {any one language (usually Common)},
challenge = 3,
]
\DndMonsterAction{Lightfooted}
The swashbuckler can take the Dash or Disengage action as a bonus action on each of its turns.

\DndMonsterAction{Suave Defense}
While the swashbuckler is wearing light or no armor and wielding no shield, its AC includes its Charisma modifier.

\par
\DndMonsterSection{Actions}
\DndMonsterAction{Multiattack}
The swashbuckler makes three attacks: one with a dagger and two with its rapier.

\DndMonsterAction{Dagger}
\textsl{Melee or Ranged Weapon Attack:} +6 to hit, reach 5 ft. or range 20/60 ft., one target. \textsl{Hit:} 6 (1d4 + 4) piercing damage.

\DndMonsterAction{Rapier}
\textsl{Melee Weapon Attack:} +6 to hit, reach 5 ft., one target. \textsl{Hit:} 8 (1d8 + 4) piercing damage.

\end{DndMonster}



\begin{DndMonster}[float=t]{Taryon Darrington}
\DndMonsterType{Medium humanoid (human), U}
\DndMonsterBasics[
armor-class = {15 (\textit{+1 breastplate})},
hit-points = {\DndDice{15d8 + 30 + 30}},
speed = {40 ft., fly 50 ft. \textit{(with \textit{broom of flying})}}
]
\DndMonsterAbilityScores[
 str = 12,
 dex = 11,
 con = 14,
 int = 18,
 wis = 8,
 cha = 15
]
\DndMonsterDetails[
saving-throws = {Con +7, Int +9},
skills = {Arcana +9, History +9, Perception +4, Survival +4},
damage-resistances = {acid, poison},
languages = {Common, Draconic, Gnomish, Sylvan, Undercommon},
challenge = 9,
]
\DndMonsterAction{Alchemical Defense}
Tary's experimental concoctions have granted him resistance to acid and poison damage, and immunity to the poisoned condition.

\DndMonsterAction{Doty}
\textbf{Doty X} moves and acts on Tary's turn, and Tary can use a bonus action to command him to use the Dash, Dodge, Disengage, or Attack action. He can also command him to Attack as part of his Multitask action (see below). Doty can take only one action on each of Tary's turns.

\DndMonsterAction{Focused}
Tary has Advantage and Disadvantage on Constitution saving throws made to maintain concentration on spells.

\DndMonsterAction{Special Equipment}
Tary wears a \textit{helm of brilliance}, a \textit{robe of useful items}, and a \textit{ring of the ram}. He carries a \textit{rod of mercurial form}, as well as a \textit{bag of holding}, 1d4 + 4 beads of force, a \textit{broom of flying}, and a \textit{chime of opening}.

\DndMonsterAction{Toughness}
Tary has an additional 30 hit points.

\DndMonsterAction{Spellcasting}
Tary is a 15th-level spellcaster. His spellcasting ability is Intelligence (spell save DC 16, +8 to hit with spell attacks). He has the following spells prepared:
\begin{DndMonsterSpells}
  \DndMonsterSpellLevel{mending, prestidigitation, ray of frost}
   \DndMonsterSpellLevel[1][4]{detect magic, heroism, identify, shield}
   \DndMonsterSpellLevel[2][3]{alter self, branding smite, see invisibility, warding bond}
   \DndMonsterSpellLevel[3][3]{dispel magic, fly, haste, magic circle, revivify}
   \DndMonsterSpellLevel[4][2]{banishment, fabricate, fire shield}
\end{DndMonsterSpells}
\par
\DndMonsterSection{Actions}
\DndMonsterAction{Multitask}
Tary makes two weapon attacks or casts a spell with a casting time of 1 action, and can also command \textbf{Doty X} to use his Empowered Fist or Repair action.

\DndMonsterAction{Mercurial Longsword}
\textsl{Melee Weapon Attack:} +5 to hit, reach 5 ft., one target. \textsl{Hit:} 6 (1d10 + 1) slashing damage.

\DndMonsterAction{\textit{Bead of Force}}
Tary throws a \textit{bead of force} up to 60 feet. Each creature within a 10-foot radius of where the bead landed must succeed on a DC 15 Dexterity saving throws or take 12 (5d4) force damage. A sphere of transparent force then encloses the area for 1 minute. A creature that fails the save and is completely within the area is trapped inside the sphere.

\par
\DndMonsterSection{Reactions}
\DndMonsterAction{Brilliant Intellect}
When Tary or another creature he can see within 30 feet of him makes an ability check or a saving throws, he can use his reaction to add +4 to the roll.

\end{DndMonster}



\begin{DndMonster}[float=t]{Treant}
\DndMonsterType{Huge plant, chaotic good}
\DndMonsterBasics[
armor-class = {16 (natural armor)},
hit-points = {\DndDice{12d12 + 60}},
speed = {30 ft.}
]
\DndMonsterAbilityScores[
 str = 23,
 dex = 8,
 con = 21,
 int = 12,
 wis = 16,
 cha = 12
]
\DndMonsterDetails[
damage-vulnerabilities = {fire},
damage-resistances = {bludgeoning, piercing},
languages = {Common, Druidic, Elvish, Sylvan},
challenge = 9,
]
\DndMonsterAction{False Appearance}
While the treant remains motionless, it is indistinguishable from a normal tree.

\DndMonsterAction{Siege Monster}
The treant deals double damage to objects and structures.

\par
\DndMonsterSection{Actions}
\DndMonsterAction{Multiattack}
The treant makes two slam attacks.

\DndMonsterAction{Slam}
\textsl{Melee Weapon Attack:} +10 to hit, reach 5 ft., one target. \textsl{Hit:} 16 (3d6 + 6) bludgeoning damage.

\DndMonsterAction{Rock}
\textsl{Ranged Weapon Attack:} +10 to hit, range 60/180 ft., one target. \textsl{Hit:} 28 (4d10 + 6) bludgeoning damage.

\DndMonsterAction{Animate Trees (1/Day)}
The treant magically animates one or two trees it can see within 60 feet of it. These trees have the same statistics as a \textbf{treant}, except they have Intelligence and Charisma scores of 1, they can't speak, and they have only the Slam action option. An animated tree acts as an ally of the treant. The tree remains animate for 1 day or until it dies; until the treant dies or is more than 120 feet from the tree; or until the treant takes a bonus action to turn it back into an inanimate tree. The tree then takes root if possible.

\end{DndMonster}



\begin{DndMonster}[float=t]{Trinket}
\DndMonsterType{Large beast (brown bear), U}
\DndMonsterBasics[
armor-class = {18 (\textit{plate barding})},
hit-points = {\DndDice{6d10 + 18 + 10}},
speed = {40 ft., climb 30 ft.}
]
\DndMonsterAbilityScores[
 str = 19,
 dex = 10,
 con = 16,
 int = 5,
 wis = 13,
 cha = 7
]
\DndMonsterDetails[
skills = {Perception +10},
languages = {understands Common but can't speak it},
challenge = 5,
]
\DndMonsterAction{Chosen Companion}
Trinket has an additional 10 hit points, and has a +6 bonus to attack and damage rolls and to Wisdom (Perception) checks.

\DndMonsterAction{Empowered Attacks}
Trinket's attacks are treated as magical for the purpose of overcoming resistance and immunity to damage from nonmagical attacks.

\DndMonsterAction{Keen Smell}
Trinket has Advantage and Disadvantage on Wisdom (Perception) checks that rely on smell.

\DndMonsterAction{Raven's Slumber}
If Trinket is reduced to 0 hit points and is within 100 feet of Vex, he is immediately drawn into the demiplane within \textbf{Vex'ahlia} \textit{raven's slumber} and stabilized.

\par
\DndMonsterSection{Actions}
\DndMonsterAction{Multiattack}
Trinket makes two attacks: one with his bite and one with his claws.

\DndMonsterAction{Bite}
\textsl{Melee Weapon Attack:} +13 to hit, reach 5 ft., one target. \textsl{Hit:} 14 (1d8 + 10) piercing damage.

\DndMonsterAction{Claws}
\textsl{Melee Weapon Attack:} +13 to hit, reach 5 ft., one target. \textsl{Hit:} 17 (2d6 + 10) slashing damage.

\end{DndMonster}


\clearpage

\begin{DndMonster}[float=t]{Troll}
\DndMonsterType{Large giant, chaotic evil}
\DndMonsterBasics[
armor-class = {15 (natural armor)},
hit-points = {\DndDice{8d10 + 40}},
speed = {30 ft.}
]
\DndMonsterAbilityScores[
 str = 18,
 dex = 13,
 con = 20,
 int = 7,
 wis = 9,
 cha = 7
]
\DndMonsterDetails[
skills = {Perception +2},
senses = {darkvision 60 ft., passive perception 12},
languages = {Giant},
challenge = 5,
]
\DndMonsterAction{Keen Smell}
The troll has advantage on Wisdom (Perception) checks that rely on smell.

\DndMonsterAction{Regeneration}
The troll regains 10 hit points at the start of its turn. If the troll takes acid or fire damage, this trait doesn't function at the start of the troll's next turn. The troll dies only if it starts its turn with 0 hit points and doesn't regenerate.

\par
\DndMonsterSection{Actions}
\DndMonsterAction{Multiattack}
The troll makes three attacks: one with its bite and two with its claws.

\DndMonsterAction{Bite}
\textsl{Melee Weapon Attack:} +7 to hit, reach 5 ft., one target. \textsl{Hit:} 7 (1d6 + 4) piercing damage.

\DndMonsterAction{Claw}
\textsl{Melee Weapon Attack:} +7 to hit, reach 5 ft., one target. \textsl{Hit:} 11 (2d6 + 4) slashing damage.

\end{DndMonster}



\begin{DndMonster}[float*=b,width=\textwidth + 8pt]{Unicorn}
\begin{multicols}{2}
\DndMonsterType{Large celestial, lawful good}
\DndMonsterBasics[
armor-class = {12},
hit-points = {\DndDice{9d10 + 18}},
speed = {50 ft.}
]
\DndMonsterAbilityScores[
 str = 18,
 dex = 14,
 con = 15,
 int = 11,
 wis = 17,
 cha = 16
]
\DndMonsterDetails[
damage-immunities = {poison},
senses = {darkvision 60 ft., passive perception 13},
languages = {Celestial, Elvish, Sylvan, telepathy 60 ft.},
challenge = 5,
]
\DndMonsterAction{Charge}
If the unicorn moves at least 20 feet straight toward a target and then hits it with a horn attack on the same turn, the target takes an extra 9 (2d8) piercing damage. If the target is a creature, it must succeed on a DC 15 Strength saving throw or be knocked prone.

\DndMonsterAction{Magic Resistance}
The unicorn has advantage on saving throws against spells and other magical effects.

\DndMonsterAction{Magic Weapons}
The unicorn's weapon attacks are magical.

\DndMonsterAction{Innate Spellcasting}
The unicorn's innate spellcasting ability is Charisma (spell save DC 14). The unicorn can innately cast the following spells, requiring no components:
\begin{DndMonsterSpells}
  \DndInnateSpellLevel{detect evil and good, druidcraft, pass without trace}
  \DndInnateSpellLevel[1e]{calm emotions, dispel evil and good, entangle}
\end{DndMonsterSpells}
\par
\DndMonsterSection{Actions}
\DndMonsterAction{Multiattack}
The unicorn makes two attacks: one with its hooves and one with its horn.

\DndMonsterAction{Hooves}
\textsl{Melee Weapon Attack:} +7 to hit, reach 5 ft., one target. \textsl{Hit:} 11 (2d6 + 4) bludgeoning damage.

\DndMonsterAction{Horn}
\textsl{Melee Weapon Attack:} +7 to hit, reach 5 ft., one target. \textsl{Hit:} 8 (1d8 + 4) piercing damage.

\DndMonsterAction{Healing Touch (3/Day)}
The unicorn touches another creature with its horn. The target magically regains 11 (2d8 + 2) hit points. In addition, the touch removes all diseases and neutralizes all poisons afflicting the target.

\DndMonsterAction{Teleport (1/Day)}
The unicorn magically teleports itself and up to three willing creatures it can see within 5 feet of it, along with any equipment they are wearing or carrying, to a location the unicorn is familiar with, up to 1 mile away.

\par
\DndMonsterSection{Legendary Actions}
Unicorn can take 3 legendary actions, choosing from the options below. Only one legendary action can be used at a time and only at the end of another creature's turn. Unicorn regains spent legendary actions at the start of its turn.

\begin{DndMonsterLegendaryActions}
\DndMonsterLegendaryAction{Hooves}{The unicorn makes one attack with its hooves.
}
\DndMonsterLegendaryAction{Shimmering Shield (Costs 2 Actions)}{The unicorn creates a shimmering, magical field around itself or another creature it can see within 60 feet of it. The target gains a +2 bonus to AC until the end of the unicorn's next turn.
}
\DndMonsterLegendaryAction{Heal Self (Costs 3 Actions)}{The unicorn magically regains 11 (2d8 + 2) hit points.
}
\end{DndMonsterLegendaryActions}
\par \DndMonsterSection{Regional Effects}
Transformed by the creature's celestial presence, the domain of a unicorn might include any of the following magical effects:


\begin{itemize}
\item Open flames of a non magical nature are extinguished within the unicorn's domain. Torches and campfires refuse to burn, but closed lanterns are unaffected.
\item Creatures native to the unicorn's domain have an easier time hiding; they have advantage on all Dexterity (Stealth) checks made to hide.
\item When a good-aligned creature casts a spell or uses a magical effect that causes another good-aligned creature to regain hit points, the target regains the maximum number of hit points possible for the spell or effect.
\item Curses affecting any good-aligned creature are suppressed.
\end{itemize}

If the unicorn dies, these effects end immediately.

\end{multicols}
\end{DndMonster}



\begin{DndMonster}[float*=b,width=\textwidth + 8pt]{Vampire}
\begin{multicols}{2}
\DndMonsterType{Medium undead (shapechanger), lawful evil}
\DndMonsterBasics[
armor-class = {16 (natural armor)},
hit-points = {\DndDice{17d8 + 68}},
speed = {30 ft.}
]
\DndMonsterAbilityScores[
 str = 18,
 dex = 18,
 con = 18,
 int = 17,
 wis = 15,
 cha = 18
]
\DndMonsterDetails[
saving-throws = {Dex +9, Wis +7, Cha +9},
skills = {Perception +7, Stealth +9},
damage-resistances = {necrotic; bludgeoning, piercing, slashing from nonmagical attacks},
senses = {darkvision 120 ft., passive perception 17},
languages = {the languages it knew in life},
challenge = 13,
]
\DndMonsterAction{Shapechanger}
If the vampire isn't in sunlight or running water, it can use its action to polymorph into a Tiny bat or a Medium cloud of mist, or back into its true form.

While in bat form, the vampire can't speak, its walking speed is 5 feet, and it has a flying speed of 30 feet. Its statistics, other than its size and speed, are unchanged. Anything it is wearing transforms with it, but nothing it is carrying does. It reverts to its true form if it dies.

While in mist form, the vampire can't take any actions, speak, or manipulate objects. It is weightless, has a flying speed of 20 feet, can hover, and can enter a hostile creature's space and stop there. In addition, if air can pass through a space, the mist can do so without squeezing, and it can't pass through water. It has advantage on Strength, Dexterity, and Constitution saving throws, and it is immune to all nonmagical damage, except the damage it takes from sunlight.

\DndMonsterAction{Legendary Resistance (3/Day)}
If the vampire fails a saving throw, it can choose to succeed instead.

\DndMonsterAction{Misty Escape}
When it drops to 0 hit points outside its resting place, the vampire transforms into a cloud of mist (as in the Shapechanger trait) instead of falling unconscious, provided that it isn't in sunlight or running water. If it can't transform, it is destroyed.

While it has 0 hit points in mist form, it can't revert to its vampire form, and it must reach its resting place within 2 hours or be destroyed. Once in its resting place, it reverts to its vampire form. It is then paralyzed until it regains at least 1 hit point. After spending 1 hour in its resting place with 0 hit points, it regains 1 hit point.

\DndMonsterAction{Regeneration}
The vampire regains 20 hit points at the start of its turn if it has at least 1 hit point and isn't in sunlight or running water. If the vampire takes radiant damage or damage from holy water, this trait doesn't function at the start of the vampire's next turn.

\DndMonsterAction{Spider Climb}
The vampire can climb difficult surfaces, including upside down on ceilings, without needing to make an ability check.

\DndMonsterAction{Vampire Weaknesses}
The vampire has the following flaws:

\textit{Forbiddance.} The vampire can't enter a residence without an invitation from one of the occupants.

\textit{Harmed by Running Water.} The vampire takes 20 acid damage if it ends its turn in running water.

\textit{Stake to the Heart.} If a piercing weapon made of wood is driven into the vampire's heart while the vampire is incapacitated in its resting place, the vampire is paralyzed until the stake is removed.

\textit{Sunlight Hypersensitivity.} The vampire takes 20 radiant damage when it starts its turn in sunlight. While in sunlight, it has disadvantage on attack rolls and ability checks.

\par
\DndMonsterSection{Actions}
\DndMonsterAction{Multiattack (Vampire Form Only)}
The vampire makes two attacks, only one of which can be a bite attack.

\DndMonsterAction{Unarmed Strike (Vampire Form Only)}
\textsl{Melee Weapon Attack:} +9 to hit, reach 5 ft., one creature. \textsl{Hit:} 8 (1d8 + 4) bludgeoning damage. Instead of dealing damage, the vampire can grapple the target (escape DC 18).

\DndMonsterAction{Bite (Bat or Vampire Form Only)}
\textsl{Melee Weapon Attack:} +9 to hit, reach 5 ft., one willing creature, or a creature that is grappled by the vampire, incapacitated, or restrained. \textsl{Hit:} 7 (1d6 + 4) piercing damage plus 10 (3d6) necrotic damage. The target's hit point maximum is reduced by an amount equal to the necrotic damage taken, and the vampire regains hit points equal to that amount. The reduction lasts until the target finishes a long rest. The target dies if this effect reduces its hit point maximum to 0. A humanoid slain in this way and then buried in the ground rises the following night as a \textbf{vampire spawn} under the vampire's control.

\DndMonsterAction{Charm}
The vampire targets one humanoid it can see within 30 feet of it. If the target can see the vampire, the target must succeed on a DC 17 Wisdom saving throw against this magic or be charmed by the vampire. The charmed target regards the vampire as a trusted friend to be heeded and protected. Although the target isn't under the vampire's control, it takes the vampire's requests or actions in the most favorable way it can, and it is a willing target for the vampire's bite attack.

Each time the vampire or the vampire's companions do anything harmful to the target, it can repeat the saving throw, ending the effect on itself on a success. Otherwise, the effect lasts 24 hours or until the vampire is destroyed, is on a different plane of existence than the target, or takes a bonus action to end the effect.

\DndMonsterAction{Children of the Night (1/Day)}
The vampire magically calls 2d4 swarms of \textbf{swarm of bats} or \textbf{swarm of rats}, provided that the sun isn't up. While outdoors, the vampire can call 3d6 \textbf{wolf} instead. The called creatures arrive in 1d4 rounds, acting as allies of the vampire and obeying its spoken commands. The beasts remain for 1 hour, until the vampire dies, or until the vampire dismisses them as a bonus action.

\par
\DndMonsterSection{Legendary Actions}
Vampire can take 3 legendary actions, choosing from the options below. Only one legendary action can be used at a time and only at the end of another creature's turn. Vampire regains spent legendary actions at the start of its turn.

\begin{DndMonsterLegendaryActions}
\DndMonsterLegendaryAction{Move}{The vampire moves up to its speed without provoking opportunity attacks.
}
\DndMonsterLegendaryAction{Unarmed Strike}{The vampire makes one unarmed strike.
}
\DndMonsterLegendaryAction{Bite (Costs 2 Actions)}{The vampire makes one bite attack.
}
\end{DndMonsterLegendaryActions}
\par \DndMonsterSection{Regional Effects}
The region surrounding a vampire's lair is warped by the creature's unnatural presence, creating any of the following effects:


\begin{itemize}
\item There's a noticeable increase in the populations of bats, rats, and wolves in the region.
\item Plants within 500 feet of the lair wither, and their stems and branches become twisted and thorny.
\item Shadows cast within 500 feet of the lair seem abnormally gaunt and sometimes move as though alive.
\item A creeping fog clings to the ground within 500 feet of the vampire's lair. The fog occasionally takes eerie forms, such as grasping claws and writhing serpents.
\end{itemize}

If the vampire is destroyed, these effects end after 2d6 days.

\end{multicols}
\end{DndMonster}



\begin{DndMonster}[float=t]{Veteran}
\DndMonsterType{Medium humanoid (any race), any alignment}
\DndMonsterBasics[
armor-class = {17 (\textit{splint armor})},
hit-points = {\DndDice{9d8 + 18}},
speed = {30 ft.}
]
\DndMonsterAbilityScores[
 str = 16,
 dex = 13,
 con = 14,
 int = 10,
 wis = 11,
 cha = 10
]
\DndMonsterDetails[
skills = {Athletics +5, Perception +2},
languages = {any one language (usually Common)},
challenge = 3,
]
\par
\DndMonsterSection{Actions}
\DndMonsterAction{Multiattack}
The veteran makes two longsword attacks. If it has a shortsword drawn, it can also make a shortsword attack.

\DndMonsterAction{Longsword}
\textsl{Melee Weapon Attack:} +5 to hit, reach 5 ft., one target. \textsl{Hit:} 7 (1d8 + 3) slashing damage, or 8 (1d10 + 3) slashing damage if used with two hands.

\DndMonsterAction{Shortsword}
\textsl{Melee Weapon Attack:} +5 to hit, reach 5 ft., one target. \textsl{Hit:} 6 (1d6 + 3) piercing damage.

\DndMonsterAction{Heavy Crossbow}
\textsl{Ranged Weapon Attack:} +3 to hit, range 100/400 ft., one target. \textsl{Hit:} 6 (1d10 + 1) piercing damage.

\end{DndMonster}



\begin{DndMonster}[float=t]{Vex'ahlia}
\DndMonsterType{Medium humanoid (half-elf), U}
\DndMonsterBasics[
armor-class = {21 (\textit{+2 Studded White Dragon Leather Armor of Cold Resistance})},
hit-points = {\DndDice{20d10 + 20}},
speed = {30 ft., fly 50 ft. \textit{(with \textit{broom of flying})}}
]
\DndMonsterAbilityScores[
 str = 7,
 dex = 20,
 con = 12,
 int = 14,
 wis = 16,
 cha = 17
]
\DndMonsterDetails[
saving-throws = {Str +6, Dex +13, Con +3, Int +4, Wis +5, Cha +5},
skills = {Acrobatics +17, Athletics +4, Deception +9, Insight +9, Investigation +8, Perception +15, Persuasion +9, Stealth +17, Survival +9},
damage-resistances = {cold},
senses = {darkvision 60 ft., passive perception 25},
languages = {Abyssal, Common, Draconic, Elvish, Thieves' Cant, Undercommon},
challenge = 18,
]
\DndMonsterAction{Assassinate}
During her first turn, Vex has Advantage and Disadvantage on attack rolls against any creature that hasn't taken a turn. Any hit Vex scores against a surprised creature is a critical hit.

\DndMonsterAction{Evasion}
If Vex is subjected to an effect that allows her to make a Dexterity saving throws to take only half damage, she instead takes no damage if she succeeds on the saving throws, and only half damage if she fails.

\DndMonsterAction{Fey Ancestry}
Vex has Advantage and Disadvantage on saving throws against being charmed, and magic can't put her to sleep.

\DndMonsterAction{Ranged Master}
Vex's ranged weapon attacks have a +2 bonus to hit (included in attacks), and ignore cover and cover, and have no range penalty. Additionally, Vex can take a −5 penalty to any ranged attack roll to gain a +10 bonus to the damage roll.

\DndMonsterAction{Sneak Attack (1/Turn)}
Vex deals an extra 14 (4d6) damage when she hits a target with a weapon attack and has Advantage and Disadvantage on the attack roll, or when the target is within 5 feet of an ally of Vex that isn't incapacitated and she doesn't have Advantage and Disadvantage on the attack roll.

\DndMonsterAction{Special Equipment}
Vex wears \textit{+2 Studded White Dragon Leather Armor of Cold Resistance} and a ring of protection with a +2 bonus. She carries a \textit{broom of flying} and a \textit{raven's slumber}, and wields Fenthras.

\DndMonsterAction{Trinket}
\textbf{Trinket} moves and acts on Vex's turn, and Vex can use a bonus action to command him to take the Dash, Dodge, Disengage, or Attack action. She can also command him to Attack as part of her Multiattack action (see below). \textbf{Trinket} can take only one action on each of Vex's turns.

\DndMonsterAction{Spellcasting}
Vex is a 13th-level spellcaster. Her spell-casting ability is Wisdom (spell save DC 17, +9 to hit with spell attacks). She has the following ranger spells prepared:
\begin{DndMonsterSpells}
   \DndMonsterSpellLevel[1][4]{cure wounds, hunter's mark}
   \DndMonsterSpellLevel[2][3]{pass without trace, protection from poison}
   \DndMonsterSpellLevel[3][3]{lightning arrow, nondetection}
   \DndMonsterSpellLevel[4][1]{freedom of movement, grasping vine}
\end{DndMonsterSpells}
\par
\DndMonsterSection{Actions}
\DndMonsterAction{Multiattack}
Vex makes two ranged attacks, and can command \textbf{Trinket} to move up to his speed and use his Multiattack action.

\DndMonsterAction{Fenthras}
\textsl{Ranged Weapon Attack:} +16 to hit, range 600 ft., one target. \textsl{Hit:} 22 (1d8 + 18) piercing damage and 3 (1d6) lightning damage, and Vex can use the following ability:

\DndMonsterAction{Bramble Shot (2/resting)}
The attack deals an extra 18 (4d8) piercing damage, and the target must succeed on a DC 17 Strength saving throws or be restrained. The target can repeat the saving throws at the end of each of its turns, ending the effect on itself on a success.

\par
\DndMonsterSection{Bonus Actions}
\DndMonsterAction{Cunning Action}
Vex takes the Dash, Disengage, or Hide action.

\DndMonsterAction{Hunter's Mark}
Vex casts \textit{hunter's mark} to mark a creature she can see within 90 feet of her. For 1 hour, Vex's weapon attacks deal an extra 1d6 damage to the marked creature, and she has Advantage and Disadvantage on Wisdom (Perception) and Wisdom (Survival) checks made to find it.

\par
\DndMonsterSection{Reactions}
\DndMonsterAction{Uncanny Dodge}
When an attacker that Vex can see hits her with an attack, she takes half damage from the attack.

\end{DndMonster}



\begin{DndMonster}[float=t]{Vos'skyriss Serpentfolk}
\DndMonsterType{Large monstrosity, neutral evil}
\DndMonsterBasics[
armor-class = {14 (natural armor)},
hit-points = {\DndDice{8d10 + 16}},
speed = {40 ft., swim 40 ft.}
]
\DndMonsterAbilityScores[
 str = 16,
 dex = 15,
 con = 15,
 int = 12,
 wis = 10,
 cha = 14
]
\DndMonsterDetails[
skills = {Perception +2, Stealth +4},
damage-immunities = {poison},
senses = {darkvision 60 ft., passive perception 12},
languages = {Common, Draconic},
challenge = 3,
]
\DndMonsterAction{Hold Breath}
The serpentfolk can hold its breath for 30 minutes.

\DndMonsterAction{Innate Spellcasting}
The serpentfolk's innate spellcasting ability is Charisma (spell save DC 12). It can innately cast the following spells, requiring no material components:
\begin{DndMonsterSpells}
  \DndInnateSpellLevel{animal friendship, speak with animals}
  \DndInnateSpellLevel[1]{pass without trace}
  \DndInnateSpellLevel[2e]{sleep, suggestion}
\end{DndMonsterSpells}
\par
\DndMonsterSection{Actions}
\DndMonsterAction{Multiattack}
The serpentfolk makes two scimitar attacks.

\DndMonsterAction{Bite}
\textsl{Melee Weapon Attack:} +5 to hit, reach 5 ft., one creature. \textsl{Hit:} 6 (1d6 + 3) piercing damage plus 7 (2d6) poison damage.

\DndMonsterAction{Scimitar}
\textsl{Melee Weapon Attack:} +5 to hit, reach 5 ft., one target. \textsl{Hit:} 6 (1d6 + 3) slashing damage.

\DndMonsterAction{Constrict}
\textsl{Melee Weapon Attack:} +5 to hit, reach 5 ft., one creature. \textsl{Hit:} 7 (1d8 + 3) bludgeoning damage, and the target is grappled (escape DC 13). Until this grapple ends, the creature is restrained, and the serpentfolk can't constrict another target.

\end{DndMonster}



\begin{DndMonster}[float=t]{Vos'skyriss Serpentfolk Ghost}
\DndMonsterType{Large undead, neutral evil}
\DndMonsterBasics[
armor-class = {13},
hit-points = {\DndDice{10d10}},
speed = {0 ft., fly 40 ft. \textit{(hover)}}
]
\DndMonsterAbilityScores[
 str = 7,
 dex = 16,
 con = 10,
 int = 12,
 wis = 10,
 cha = 16
]
\DndMonsterDetails[
skills = {Perception +2, Stealth +7},
damage-resistances = {acid, fire, lightning, thunder; bludgeoning, piercing, slashing from nonmagical attacks},
damage-immunities = {cold, necrotic, poison},
senses = {darkvision 60 ft., passive perception 12},
languages = {Common, Draconic},
challenge = 4,
]
\DndMonsterAction{Ethereal Sight}
The ghost can see 60 feet into the Ethereal Plane when it is on the Material Plane, and vice versa.

\DndMonsterAction{Incorporeal Movement}
The ghost can move through other creatures and objects as if they were difficult terrain. It takes 5 (1d10) force damage if it ends its turn inside an object.

\par
\DndMonsterSection{Actions}
\DndMonsterAction{Multiattack}
The ghost makes two spear attacks.

\DndMonsterAction{Spear}
\textsl{Melee or Ranged Weapon Attack:} +5 to hit, reach 5 ft. or range 20/60 ft., one target. \textsl{Hit:} 6 (1d6 + 3) piercing damage plus 7 (2d6) necrotic damage. If the ghost makes a ranged attack, its spear rematerializes in its hands after the attack is resolved.

\DndMonsterAction{Etherealness}
The ghost enters the Ethereal Plane from the Material Plane, or vice versa. It is visible on the Material Plane while it is in the Border Ethereal, and vice versa, yet it can't affect or be affected by anything on the other plane.

\DndMonsterAction{Terrifying Hiss}
The ghost emits an uncanny, rasping hiss. Each non-undead creature within 30 feet of the ghost that can hear it must succeed on a DC 13 Wisdom saving throws or be frightened for 1 minute. A frightened target can repeat the saving throws at the end of each of its turns, ending the condition on itself on a success. If a target's saving throws is successful or the effect ends for it, the target is immune to this ghost's Terrifying Hiss for the next 24 hours.

\end{DndMonster}



\begin{DndMonster}[float=t]{Water Elemental}
\DndMonsterType{Large elemental, neutral}
\DndMonsterBasics[
armor-class = {14 (natural armor)},
hit-points = {\DndDice{12d10 + 48}},
speed = {30 ft., swim 90 ft.}
]
\DndMonsterAbilityScores[
 str = 18,
 dex = 14,
 con = 18,
 int = 5,
 wis = 10,
 cha = 8
]
\DndMonsterDetails[
damage-resistances = {acid; bludgeoning, piercing, slashing from nonmagical attacks},
damage-immunities = {poison},
senses = {darkvision 60 ft., passive perception 10},
languages = {Aquan},
challenge = 5,
]
\DndMonsterAction{Water Form}
The elemental can enter a hostile creature's space and stop there. It can move through a space as narrow as 1 inch wide without squeezing.

\DndMonsterAction{Freeze}
If the elemental takes cold damage, it partially freezes; its speed is reduced by 20 feet until the end of its next turn.

\par
\DndMonsterSection{Actions}
\DndMonsterAction{Multiattack}
The elemental makes two slam attacks.

\DndMonsterAction{Slam}
\textsl{Melee Weapon Attack:} +7 to hit, reach 5 ft., one target. \textsl{Hit:} 13 (2d8 + 4) bludgeoning damage.

\DndMonsterAction{Whelm (Recharge 4\textendash 6)}
Each creature in the elemental's space must make a DC 15 Strength saving throw. On a failure, a target takes 13 (2d8 + 4) bludgeoning damage. If it is Large or smaller, it is also grappled (escape DC 14). Until this grapple ends, the target is restrained and unable to breathe unless it can breathe water. If the saving throw is successful, the target is pushed out of the elemental's space.

The elemental can grapple one Large creature or up to two Medium or smaller creatures at one time. At the start of each of the elemental's turns, each target grappled by it takes 13 (2d8 + 4) bludgeoning damage. A creature within 5 feet of the elemental can pull a creature or object out of it by taking an action to make a DC 14 Strength check and succeeding.

\end{DndMonster}



\begin{DndMonster}[float=t]{Werebear}
\DndMonsterType{Medium humanoid (human, shapechanger), neutral good}
\DndMonsterBasics[
armor-class = {10 in humanoid form (11 (natural armor) in bear or hybrid form)},
hit-points = {\DndDice{18d8 + 54}},
speed = {{'number': 30, 'condition': '(40 ft., climb 30 ft. in bear or hybrid form)'} ft.}
]
\DndMonsterAbilityScores[
 str = 19,
 dex = 10,
 con = 17,
 int = 11,
 wis = 12,
 cha = 12
]
\DndMonsterDetails[
skills = {Perception +7},
damage-immunities = {; bludgeoning, piercing, slashing from nonmagical attacks that aren't silvered},
languages = {Common (can't speak in bear form)},
challenge = 5,
]
\DndMonsterAction{Shapechanger}
The werebear can use its action to polymorph into a Large bear-humanoid hybrid or into a Large bear, or back into its true form, which is humanoid. Its statistics, other than its size and AC, are the same in each form. Any equipment it is wearing or carrying isn't transformed. It reverts to its true form if it dies.

\DndMonsterAction{Keen Smell}
The werebear has advantage on Wisdom (Perception) checks that rely on smell.

\par
\DndMonsterSection{Actions}
\DndMonsterAction{Multiattack}
In bear form, the werebear makes two claw attacks. In humanoid form, it makes two greataxe attacks. In hybrid form, it can attack like a bear or a humanoid.

\DndMonsterAction{Bite (Bear or Hybrid Form Only)}
\textsl{Melee Weapon Attack:} +7 to hit, reach 5 ft., one target. \textsl{Hit:} 15 (2d10 + 4) piercing damage. If the target is a humanoid, it must succeed on a DC 14 Constitution saving throw or be cursed with werebear lycanthropy.

\DndMonsterAction{Claw (Bear or Hybrid Form Only)}
\textsl{Melee Weapon Attack:} +7 to hit, reach 5 ft., one target. \textsl{Hit:} 13 (2d8 + 4) slashing damage.

\DndMonsterAction{Greataxe (Humanoid or Hybrid Form Only)}
\textsl{Melee Weapon Attack:} +7 to hit, reach 5 ft., one target. \textsl{Hit:} 10 (1d12 + 4) slashing damage.

\end{DndMonster}



\begin{DndMonster}[float=t]{Werewolf}
\DndMonsterType{Medium humanoid (human, shapechanger), chaotic evil}
\DndMonsterBasics[
armor-class = {11 in humanoid form (12 (natural armor) in wolf or hybrid form)},
hit-points = {\DndDice{9d8 + 18}},
speed = {{'number': 30, 'condition': '(40 ft. in wolf form)'} ft.}
]
\DndMonsterAbilityScores[
 str = 15,
 dex = 13,
 con = 14,
 int = 10,
 wis = 11,
 cha = 10
]
\DndMonsterDetails[
skills = {Perception +4, Stealth +3},
damage-immunities = {; bludgeoning, piercing, slashing from nonmagical attacks that aren't silvered},
languages = {Common (can't speak in wolf form)},
challenge = 3,
]
\DndMonsterAction{Shapechanger}
The werewolf can use its action to polymorph into a wolf-humanoid hybrid or into a wolf, or back into its true form, which is humanoid. Its statistics, other than its AC, are the same in each form. Any equipment it is wearing or carrying isn't transformed. It reverts to its true form if it dies.

\DndMonsterAction{Keen Hearing and Smell}
The werewolf has advantage on Wisdom (Perception) checks that rely on hearing or smell.

\par
\DndMonsterSection{Actions}
\DndMonsterAction{Multiattack (Humanoid or Hybrid Form Only)}
The werewolf makes two attacks: two with its spear (humanoid form) or one with its bite and one with its claws (hybrid form).

\DndMonsterAction{Bite (Wolf or Hybrid Form Only)}
\textsl{Melee Weapon Attack:} +4 to hit, reach 5 ft., one target. \textsl{Hit:} 6 (1d8 + 2) piercing damage. If the target is a humanoid, it must succeed on a DC 12 Constitution saving throw or be cursed with werewolf lycanthropy.

\DndMonsterAction{Claws (Hybrid Form Only)}
\textsl{Melee Weapon Attack:} +4 to hit, reach 5 ft., one creature. \textsl{Hit:} 7 (2d4 + 2) slashing damage.

\DndMonsterAction{Spear (Humanoid Form Only)}
\textsl{Melee or Ranged Weapon Attack:} +4 to hit, reach 5 ft. or range 20/60 ft., one creature. \textsl{Hit:} 5 (1d6 + 2) piercing damage, or 6 (1d8 + 2) piercing damage if used with two hands to make a melee attack.

\end{DndMonster}



\begin{DndMonster}[float=t]{Wight}
\DndMonsterType{Medium undead, neutral evil}
\DndMonsterBasics[
armor-class = {14 (studded leather)},
hit-points = {\DndDice{6d8 + 18}},
speed = {30 ft.}
]
\DndMonsterAbilityScores[
 str = 15,
 dex = 14,
 con = 16,
 int = 10,
 wis = 13,
 cha = 15
]
\DndMonsterDetails[
skills = {Perception +3, Stealth +4},
damage-resistances = {necrotic; bludgeoning, piercing, slashing from nonmagical attacks that aren't silvered},
damage-immunities = {poison},
senses = {darkvision 60 ft., passive perception 13},
languages = {the languages it knew in life},
challenge = 3,
]
\DndMonsterAction{Sunlight Sensitivity}
While in sunlight, the wight has disadvantage on attack rolls, as well as on Wisdom (Perception) checks that rely on sight.

\par
\DndMonsterSection{Actions}
\DndMonsterAction{Multiattack}
The wight makes two longsword attacks or two longbow attacks. It can use its Life Drain in place of one longsword attack.

\DndMonsterAction{Life Drain}
\textsl{Melee Weapon Attack:} +4 to hit, reach 5 ft., one creature. \textsl{Hit:} 5 (1d6 + 2) necrotic damage. The target must succeed on a DC 13 Constitution saving throw or its hit point maximum is reduced by an amount equal to the damage taken. This reduction lasts until the target finishes a long rest. The target dies if this effect reduces its hit point maximum to 0.

A humanoid slain by this attack rises 24 hours later as a \textbf{zombie} under the wight's control, unless the humanoid is restored to life or its body is destroyed. The wight can have no more than twelve zombies under its control at one time.

\DndMonsterAction{Longsword}
\textsl{Melee Weapon Attack:} +4 to hit, reach 5 ft., one target. \textsl{Hit:} 6 (1d8 + 2) slashing damage, or 7 (1d10 + 2) slashing damage if used with two hands.

\DndMonsterAction{Longbow}
\textsl{Ranged Weapon Attack:} +4 to hit, range 150/600 ft., one target. \textsl{Hit:} 6 (1d8 + 2) piercing damage.

\end{DndMonster}



\begin{DndMonster}[float=t]{Worg}
\DndMonsterType{Large monstrosity, neutral evil}
\DndMonsterBasics[
armor-class = {13 (natural armor)},
hit-points = {\DndDice{4d10 + 4}},
speed = {50 ft.}
]
\DndMonsterAbilityScores[
 str = 16,
 dex = 13,
 con = 13,
 int = 7,
 wis = 11,
 cha = 8
]
\DndMonsterDetails[
skills = {Perception +4},
senses = {darkvision 60 ft., passive perception 14},
languages = {Goblin, Worg},
challenge = 1/2,
]
\DndMonsterAction{Keen Hearing and Smell}
The worg has advantage on Wisdom (Perception) checks that rely on hearing or smell.

\par
\DndMonsterSection{Actions}
\DndMonsterAction{Bite}
\textsl{Melee Weapon Attack:} +5 to hit, reach 5 ft., one target. \textsl{Hit:} 10 (2d6 + 3) piercing damage. If the target is a creature, it must succeed on a DC 13 Strength saving throw or be knocked prone.

\end{DndMonster}



\begin{DndMonster}[float=t]{Wraith}
\DndMonsterType{Medium undead, neutral evil}
\DndMonsterBasics[
armor-class = {13},
hit-points = {\DndDice{9d8 + 27}},
speed = {0 ft., fly 60 ft. \textit{(hover)}}
]
\DndMonsterAbilityScores[
 str = 6,
 dex = 16,
 con = 16,
 int = 12,
 wis = 14,
 cha = 15
]
\DndMonsterDetails[
damage-resistances = {acid, cold, fire, lightning, thunder; bludgeoning, piercing, slashing from nonmagical attacks that aren't silvered},
damage-immunities = {necrotic, poison},
senses = {darkvision 60 ft., passive perception 12},
languages = {the languages it knew in life},
challenge = 5,
]
\DndMonsterAction{Incorporeal Movement}
The wraith can move through other creatures and objects as if they were difficult terrain. It takes 5 (1d10) force damage if it ends its turn inside an object.

\DndMonsterAction{Sunlight Sensitivity}
While in sunlight, the wraith has disadvantage on attack rolls, as well as on Wisdom (Perception) checks that rely on sight.

\par
\DndMonsterSection{Actions}
\DndMonsterAction{Life Drain}
\textsl{Melee Weapon Attack:} +6 to hit, reach 5 ft., one creature. \textsl{Hit:} 21 (4d8 + 3) necrotic damage. The target must succeed on a DC 14 Constitution saving throw or its hit point maximum is reduced by an amount equal to the damage taken. This reduction lasts until the target finishes a long rest. The target dies if this effect reduces its hit point maximum to 0.

\DndMonsterAction{Create Specter}
The wraith targets a humanoid within 10 feet of it that has been dead for no longer than 1 minute and died violently. The target's spirit rises as a \textbf{specter} in the space of its corpse or in the nearest unoccupied space. The \textbf{specter} is under the wraith's control. The wraith can have no more than seven specters under its control at one time.

\end{DndMonster}



\begin{DndMonster}[float=t]{Wraithroot Tree}
\DndMonsterType{Huge plant, neutral evil}
\DndMonsterBasics[
armor-class = {16 (natural armor)},
hit-points = {\DndDice{20d12 + 140}},
speed = {30 ft.}
]
\DndMonsterAbilityScores[
 str = 23,
 dex = 8,
 con = 24,
 int = 12,
 wis = 16,
 cha = 12
]
\DndMonsterDetails[
skills = {Perception +8},
damage-vulnerabilities = {fire},
damage-resistances = {bludgeoning, piercing},
languages = {Common, Druidic, Elvish, Sylvan},
challenge = 14,
]
\DndMonsterAction{False Appearance}
While the wraithroot tree remains motionless, it is indistinguishable from a normal tree.

\DndMonsterAction{Wraith Shift}
On its turn, the wraithroot tree can move through other creatures and objects as if they were difficult terrain. The tree does not provoke opportunity attack while moving, and must finish its movement in a space unoccupied by any Large or larger objects or creatures. Any Medium or smaller creature in the tree's space when it finishes moving is pushed to the nearest unoccupied space of the creature's choice. Medium or smaller objects in the tree's space are pushed to the nearest unoccupied spaces of the tree's choice.

\par
\DndMonsterSection{Actions}
\DndMonsterAction{Multiattack}
The wraithroot tree makes two slam attacks.

\DndMonsterAction{Slam}
\textsl{Melee Weapon Attack:} +11 to hit, reach 10 ft., one target. \textsl{Hit:} 16 (3d6 + 6) bludgeoning damage.

\DndMonsterAction{Rock}
\textsl{Ranged Weapon Attack:} +11 to hit, range 60/180 ft., one target. \textsl{Hit:} 28 (4d10 + 6) bludgeoning damage.

\DndMonsterAction{Wraithstorm (Recharge 5\textendash 6)}
The wraithroot tree pulls life energy from the area around it, leaving creatures temporarily frail. Each creature that is not a construct or undead within 30 feet of the tree must succeed on a DC 16 Constitution saving throws or have its speed halved and gain vulnerability to bludgeoning damage until the end of the tree's next turn.

\end{DndMonster}



\begin{DndMonster}[float=t]{Wyvern}
\DndMonsterType{Large dragon, unaligned}
\DndMonsterBasics[
armor-class = {13 (natural armor)},
hit-points = {\DndDice{13d10 + 39}},
speed = {20 ft., fly 80 ft.}
]
\DndMonsterAbilityScores[
 str = 19,
 dex = 10,
 con = 16,
 int = 5,
 wis = 12,
 cha = 6
]
\DndMonsterDetails[
skills = {Perception +4},
senses = {darkvision 60 ft., passive perception 14},
challenge = 6,
]
\par
\DndMonsterSection{Actions}
\DndMonsterAction{Multiattack}
The wyvern makes two attacks: one with its bite and one with its stinger. While flying, it can use its claws in place of one other attack.

\DndMonsterAction{Bite}
\textsl{Melee Weapon Attack:} +7 to hit, reach 10 ft., one creature. \textsl{Hit:} 11 (2d6 + 4) piercing damage.

\DndMonsterAction{Claws}
\textsl{Melee Weapon Attack:} +7 to hit, reach 5 ft., one target. \textsl{Hit:} 13 (2d8 + 4) slashing damage.

\DndMonsterAction{Stinger}
\textsl{Melee Weapon Attack:} +7 to hit, reach 10 ft., one creature. \textsl{Hit:} 11 (2d6 + 4) piercing damage. The target must make a DC 15 Constitution saving throw, taking 24 (7d6) poison damage on a failed save, or half as much damage on a successful one.

\end{DndMonster}



\begin{DndMonster}[float=t]{Young Green Dragon}
\DndMonsterType{Large dragon, lawful evil}
\DndMonsterBasics[
armor-class = {18 (natural armor)},
hit-points = {\DndDice{16d10 + 48}},
speed = {40 ft., fly 80 ft., swim 40 ft.}
]
\DndMonsterAbilityScores[
 str = 19,
 dex = 12,
 con = 17,
 int = 16,
 wis = 13,
 cha = 15
]
\DndMonsterDetails[
saving-throws = {Dex +4, Con +6, Wis +4, Cha +5},
skills = {Deception +5, Perception +7, Stealth +4},
damage-immunities = {poison},
senses = {blindsight 30 ft., darkvision 120 ft., passive perception 17},
languages = {Common, Draconic},
challenge = 8,
]
\DndMonsterAction{Amphibious}
The dragon can breathe air and water.

\par
\DndMonsterSection{Actions}
\DndMonsterAction{Multiattack}
The dragon makes three attacks: one with its bite and two with its claws.

\DndMonsterAction{Bite}
\textsl{Melee Weapon Attack:} +7 to hit, reach 10 ft., one target. \textsl{Hit:} 15 (2d10 + 4) piercing damage plus 7 (2d6) poison damage.

\DndMonsterAction{Claw}
\textsl{Melee Weapon Attack:} +7 to hit, reach 5 ft., one target. \textsl{Hit:} 11 (2d6 + 4) slashing damage.

\DndMonsterAction{Poison Breath (Recharge 5\textendash 6)}
The dragon exhales poisonous gas in a 30-foot cone. Each creature in that area must make a DC 14 Constitution saving throw, taking 42 (12d6) poison damage on a failed save, or half as much damage on a successful one.

\end{DndMonster}



\begin{DndMonster}[float=t]{Young Magma Landshark}
\DndMonsterType{Large elemental, U}
\DndMonsterBasics[
armor-class = {16 (natural armor)},
hit-points = {\DndDice{10d10 + 50}},
speed = {40 ft., burrow 40 ft., swim 40 ft. \textit{(lava only)}}
]
\DndMonsterAbilityScores[
 str = 20,
 dex = 11,
 con = 21,
 int = 2,
 wis = 10,
 cha = 5
]
\DndMonsterDetails[
skills = {Perception +4},
damage-immunities = {fire},
senses = {darkvision 60 ft., tremorsense 120 ft., passive perception 14},
challenge = 9,
]
\DndMonsterAction{Searing Presence}
Any creature that starts its turn within 10 feet of the landshark takes 11 (2d10) fire damage.

\DndMonsterAction{Standing Leap}
The landshark's jumping is up to 40 feet and its jumping is up to 20 feet, with or without a running start.

\DndMonsterAction{Illumination}
The landshark sheds bright light in a 15-foot radius and dim light for an additional 15 feet.

\par
\DndMonsterSection{Actions}
\DndMonsterAction{Bite}
\textsl{Melee Weapon Attack:} +9 to hit, reach 5 ft., one target. \textsl{Hit:} 31 (4d12 + 5) piercing damage, and 14 (4d6) fire damage.

\DndMonsterAction{Deadly Leap}
If the landshark jumping at least 15 feet as part of its movement, it can then use this action to land on its feet in a space that contains one or more other creatures. Each of those creatures must succeed on a DC 16 Strength or Dexterity saving throws (target's choice) or be knocked prone and take 19 (4d6 + 5) bludgeoning damage plus 14 (4d6) fire damage. On a successful save, the creature takes only half the damage, isn't knocked prone, and is pushed 5 feet out of the landshark's space into an unoccupied space of the creature's choice. If no unoccupied space is within range, the creature instead falls prone in the land-shark's space.

\end{DndMonster}



\begin{DndMonster}[float=t]{Zombie}
\DndMonsterType{Medium undead, neutral evil}
\DndMonsterBasics[
armor-class = {8},
hit-points = {\DndDice{3d8 + 9}},
speed = {20 ft.}
]
\DndMonsterAbilityScores[
 str = 13,
 dex = 6,
 con = 16,
 int = 3,
 wis = 6,
 cha = 5
]
\DndMonsterDetails[
saving-throws = {Wis +0},
damage-immunities = {poison},
senses = {darkvision 60 ft., passive perception 8},
languages = {understands all languages it spoke in life but can't speak},
challenge = 1/4,
]
\DndMonsterAction{Undead Fortitude}
If damage reduces the zombie to 0 hit points, it must make a Constitution saving throw with a DC of 5 + the damage taken, unless the damage is radiant or from a critical hit. On a success, the zombie drops to 1 hit point instead.

\par
\DndMonsterSection{Actions}
\DndMonsterAction{Slam}
\textsl{Melee Weapon Attack:} +3 to hit, reach 5 ft., one target. \textsl{Hit:} 4 (1d6 + 1) bludgeoning damage.

\end{DndMonster}


\end{document}
